\documentclass[11pt]{article}

\usepackage[margin=1in]{geometry}
\usepackage{amsmath,amssymb,amsthm,mathtools,bm}
\usepackage{enumitem,booktabs,array}
\usepackage{microtype}
\usepackage[colorlinks=true,linkcolor=blue,citecolor=blue,urlcolor=blue]{hyperref}

\newtheorem{theorem}{Theorem}[section]
\newtheorem{proposition}[theorem]{Proposition}
\newtheorem{lemma}[theorem]{Lemma}
\newtheorem{corollary}[theorem]{Corollary}
\newtheorem{definition}[theorem]{Definition}
\newtheorem{remark}[theorem]{Remark}
\newtheorem{assumption}[theorem]{Assumption}
\newtheorem{obstruction}[theorem]{Obstruction}
\newtheorem{firewall}[theorem]{Firewall}

\newcommand{\cH}{\mathcal H}
\newcommand{\cL}{\mathcal L}
\newcommand{\cA}{\mathcal A}
\newcommand{\cD}{\mathcal D}
\newcommand{\cR}{\mathcal R}
\newcommand{\cP}{\mathcal P}
\newcommand{\cK}{\mathcal K}
\newcommand{\cS}{\mathcal S}
\newcommand{\cY}{\mathcal Y}
\newcommand{\cU}{\mathcal U}
\newcommand{\eps}{\varepsilon}
\newcommand{\1}{\mathbf 1}
\newcommand{\Tr}{\operatorname{Tr}}
\newcommand{\tr}{\operatorname{tr}}
\newcommand{\Det}{\operatorname{Det}}
\newcommand{\Pf}{\operatorname{Pf}}
\newcommand{\sgn}{\operatorname{sgn}}
\newcommand{\Ran}{\operatorname{Ran}}
\newcommand{\dist}{\operatorname{dist}}
\newcommand{\ad}{\operatorname{ad}}
\newcommand{\supp}{\operatorname{supp}}
\newcommand{\Lip}{\operatorname{Lip}}
\newcommand{\Span}{\operatorname{span}}
\newcommand{\Id}{\mathrm I}
\newcommand{\Norm}[1]{\left\lVert #1\right\rVert}
\newcommand{\Abs}[1]{\left\lvert #1\right\rvert}
\newcommand{\ip}[2]{\left\langle #1,#2\right\rangle}
\newcommand{\dd}{\,\mathrm d}
\newcommand{\UV}{\mathrm{UV}}
\newcommand{\IR}{\mathrm{IR}}
\newcommand{\SM}{\mathrm{SM}}
\newcommand{\BV}{\mathrm{BV}}
\newcommand{\BRST}{\mathrm{BRST}}
\newcommand{\ov}{\mathrm{ov}}
\newcommand{\hard}{\mathrm h}
\newcommand{\soft}{\mathrm s}
\newcommand{\wt}{\mathrm{wt}}
\newcommand{\loc}{\mathrm{loc}}
\newcommand{\remn}{\mathrm{rem}}
\newcommand{\ord}{\mathrm{ord}}
\newcommand{\Op}{\mathrm{Op}}
\newcommand{\Bell}{\mathrm B}
\newcommand{\Part}{\Pi}

\title{Renewal Perturbative Einstein--Standard Model Continuum\\
Research Notes V1\\[0.4em]
\large Native chiral hard-shell reduction: overlap functional calculus, local line frames, and the exact Wilson-symbol bottleneck}
\author{}
\date{21 September 2026}

\begin{document}
\maketitle
\tableofcontents

\section{Context, target, and the V1 attack}

\subsection{Ultimate coefficientwise target}

The programme is to establish a \emph{source-complete coefficientwise perturbative continuum} for the Renewal Einstein--Standard Model regulator, with gravity treated as an effective field theory.  At every prescribed finite loop order, EFT/derivative order, source order, and finite external observable/source panel, one seeks a single renormalized action/source/reference family whose normalized connected coefficients possess cutoff-independent continuum limits and satisfy the coupled BV/BRST/master identities and the measured shell-composition law.

The quantifiers are essential.  The target is coefficientwise at every fixed finite perturbative order.  It does \emph{not} assert convergence of the perturbation series, a uniform limit as loop order tends to infinity, finite-dimensional renormalizability of Einstein gravity, asymptotic safety, a nonperturbative quantum-gravity measure, a nonperturbative Standard-Model completion, or removal of a massless infrared reference scale unless that is separately proved.

\subsection{Checkpoint inherited into this new lineage}

The supplied V80 notes close the bosonic Einstein--$SU(3)\times SU(2)\times U(1)_Y$--Higgs two-loop local/range/source self-map at the declared finite weight and source order.  The programme statement further declares the next chiral checkpoint as already containing, at its stated scope:
\begin{enumerate}[label=\textup{(I\arabic*)},leftmargin=3.2em]
\item the three-generation Standard Model chiral representation;
\item cancellation of the five ordinary four-dimensional local gauge-anomaly coefficients in the declared hypercharge convention;
\item even $SU(2)$ doublet parity;
\item exact determinant Schur factorization and scale cocycle under finite regrouping of nested shells;
\item the analogous optional neutral-Majorana Pfaffian factorization/cocycle;
\item differentiated determinant/Pfaffian line-source identities;
\item the bosonic two-loop self-map inherited from V80.
\end{enumerate}
These statements are treated as \emph{inherited}.  They are not reproved below.

The Grand Tensor manuscript also supplies a literal protected overlap/Yukawa field shell.  In its notation one has a gauge-covariant Hermitian Wilson kernel $H_W$, a protected spectral window
\begin{equation}
 g_*\Id\preceq |H_W|\preceq M_*\Id,
 \label{eq:v1-GT-Wilson-window}
\end{equation}
the overlap sign and moving chiral projector
\begin{equation}
 \eps=\sgn H_W,
 \qquad
 D_{\ov}=a^{-1}(\Id+\gamma_5\eps),
 \qquad
 \widehat P_-={\tfrac12}(\Id+\eps),
 \label{eq:v1-GT-overlap}
\end{equation}
and the chiral overlap--Yukawa block
\begin{equation}
 D_\chi=P_+D_{\ov}\widehat P_-+\cY(H;\mathbf Y),
 \qquad
 s_{\min}(D_\chi)\ge \tau_*>0
 \label{eq:v1-GT-chiral-block}
\end{equation}
on a protected finite chart.  The manuscript also explicitly requires the matter, ghost, mediator and chiral operators to be differentiated with respect to one common coframe family when common stress/source responses are formed.  These are native finite-regulator data.  They are not continuum replacements.

\subsection{Programme E1--E6}

For orientation, the continuation programme is:
\begin{description}[leftmargin=3.6em,style=nextline]
\item[E1.] Prove the native chiral hard-shell theorem: the actual spin--Dirac/Yukawa annular symbol, Wilson/regulator completion, gauge and spin transport, geometric/gauge/Higgs/antifield/line derivatives, a local chiral line frame, a hard-pivot margin, finite-to-ordinary comparison, and rooted/quasilocal estimates.
\item[E2.] Use E1 plus the inherited anomaly/line algebra to restore the complete chiral two-loop shell and obtain one common two-loop ESM action/source/line update.
\item[E3.] Prove the complete common BV/master/source equation with one finite counterterm action, including diffeomorphism, Lorentz, internal gauge, Higgs/Yukawa, chiral line, ghost/antifield, measure/reference, stress and contact rows.
\item[E4.] Formulate and prove the arbitrary fixed-loop induction, including the finite EFT bank, source/contact tower, forests, line data, counterterms, rooted remainders and reference data.
\item[E5.] Compose arbitrarily many shells at fixed loop/EFT/source order and prove a coefficientwise multi-shell continuum for normalized connected functionals.
\item[E6.] Prove scheme/reference compatibility and order compatibility, then state the final coefficientwise ESM continuum theorem with its retained IR/reference hypotheses.
\end{description}

\subsection{What V1 attacks}

V1 attacks E1, but does so in an anti-circular order.  The protected overlap formula \eqref{eq:v1-GT-overlap} is a nonlinear functional calculus of the native Wilson kernel.  Consequently the first task is to prove, without replacing $H_W$ by a continuum Dirac operator, that a local gapped Wilson kernel transfers locality and all fixed source derivatives to $\eps$, $\widehat P_-$ and $D_\chi$.  We then construct a local Kato frame for the moving chiral source space, prove a hard-pivot stability theorem, and derive the determinant/Pfaffian derivative \emph{bounds} that are missing from the already-inherited exact factorization identities.

The finite-to-ordinary annular estimate has one further input: an explicit native flat Wilson symbol and its small-$a|p|$ comparison.  The attached Grand Tensor statement gives the gap \eqref{eq:v1-GT-Wilson-window} and covariance of $H_W$, but not an explicit stencil/symbol with the derivative estimates needed for this comparison.  V1 therefore proves a reduction theorem showing exactly which Wilson-symbol estimate remains load-bearing.  It does not silently choose a different Wilson operator.

\section{Native operator class and norms}

\subsection{Finite cell carrier and source coordinates}

Let $X_a$ be a finite four-dimensional periodic cell complex of mesh $a$, with uniformly bounded local geometry and polynomial ball growth
\begin{equation}
 \#\{y:d_a(x,y)\le R\}\le C_{\rm vol}(1+R/a)^4.
 \label{eq:v1-volume-growth}
\end{equation}
The one-particle carrier $\cH_a$ is the finite direct sum over cells of the spin, generation, colour, weak and hypercharge representation fibres required by the declared chiral Standard Model packet.  The optional real/Majorana carrier is denoted $\cH_a^M$.

Let $\zeta=(\zeta^A)$ denote a finite list of normalized local coordinates drawn from the retained coframe/spin connection, internal gauge links, Higgs/Yukawa variables, BV/antifield sources and line-source coordinates required at the fixed two-loop/source order.  A multi-index $I=(A_1,\ldots,A_r)$ denotes the mixed derivative
\[
 \partial_I=\partial_{A_1}\cdots\partial_{A_r}.
\]
The source coordinates are finite in number in V1.  No analytic infinite-source radius is claimed.

\subsection{Weighted Schur algebra}

For a matrix kernel $K=(K(x,y))_{x,y\in X_a}$ and $\eta\ge0$, define
\begin{equation}
 \Norm{K}_{\eta,1}
 :=\max\left\{
 \sup_x\sum_y e^{\eta d_a(x,y)}\Norm{K(x,y)},
 \sup_y\sum_x e^{\eta d_a(x,y)}\Norm{K(x,y)}
 \right\}.
 \label{eq:v1-weighted-Schur}
\end{equation}
Finite fibre-component sums are included in the matrix norm.  The same definition is used for rectangular blocks between fixed finite fibre types.

\begin{lemma}[Weighted Schur algebra]
\label{lem:v1-Schur-algebra}
For compatible kernels $K,L$,
\begin{equation}
 \Norm{KL}_{\eta,1}\le \Norm K_{\eta,1}\Norm L_{\eta,1}.
 \label{eq:v1-Schur-product}
\end{equation}
The adjoint has the same norm.  In particular every finite product of quasilocal kernels remains quasilocal with no volume factor.
\end{lemma}
\begin{proof}
Since $d_a(x,z)\le d_a(x,y)+d_a(y,z)$,
\[
 e^{\eta d_a(x,z)}\Norm{(KL)(x,z)}
 \le\sum_y e^{\eta d_a(x,y)}\Norm{K(x,y)}
            e^{\eta d_a(y,z)}\Norm{L(y,z)}.
\]
Sum first over $z$ and then over $y$ to obtain the row estimate; reverse the order for the column estimate.  Taking the maximum proves \eqref{eq:v1-Schur-product}.
\end{proof}

\begin{definition}[Native local Wilson class]
\label{def:v1-native-Wilson-class}
A family $H_W(\zeta)$ belongs to the V1 native local Wilson class on a protected parameter chart $\cU$ if:
\begin{enumerate}[label=\textup{(W\arabic*)},leftmargin=3.2em]
\item $H_W(\zeta)=H_W(\zeta)^*$ and \eqref{eq:v1-GT-Wilson-window} holds uniformly for $\zeta\in\cU$;
\item there is a fixed range $R_W$ in cell units such that $H_W(x,y)=0$ for $d_a(x,y)>R_Wa$;
\item for every nonempty source multi-index $I$ required at the fixed source order, $\partial_IH_W$ exists, has a fixed finite range $R_Ia$, and
\begin{equation}
 \sup_{\zeta\in\cU}\Norm{\partial_IH_W(\zeta)}_{0,1}\le h_I<\infty;
 \label{eq:v1-H-derivative-stock}
\end{equation}
\item the exact gauge/spin covariance of the native stencil is retained: local internal-gauge and spin-frame changes act by the declared unitary similarities on $H_W$ and its source derivatives;
\item the Yukawa block $\cY(H;\mathbf Y)$ and all required source derivatives are onsite or fixed-range kernels with finite stocks $y_I$.
\end{enumerate}
The constants may depend on the fixed perturbative/source order and on the protected coefficient chart, but not on volume or shell number.
\end{definition}

The finite-range assumption is the natural form of ``Wilson/spin completion actually used by the native discretization''.  It is deliberately separated from the still-unproved annular \emph{symbol comparison}; locality of the operator does not by itself identify its continuum principal symbol.

\section{Gapped Wilson functional calculus is uniformly quasilocal}

\subsection{A finite-volume Combes--Thomas estimate}

\begin{lemma}[Finite-range resolvent decay]
\label{lem:v1-CT}
Let $H=H^*$ have range at most $R_0a$ and $\Norm H_{0,1}\le M_0$.  If
$\delta=\dist(z,\sigma(H))>0$, then there are constants
$\eta_0=\eta_0(\delta,M_0,R_0)>0$ and $C_0=C_0(\delta,M_0,R_0,C_{\rm vol})$
such that, uniformly in the finite period,
\begin{equation}
 \Norm{(z-H)^{-1}}_{\eta,1}\le C_0
 \qquad(0\le\eta\le\eta_0/(2a)).
 \label{eq:v1-CT-Schur}
\end{equation}
One may take $\eta_0>0$ so small that
$M_0(e^{\eta_0R_0}-1)\le\delta/2$.
\end{lemma}
\begin{proof}
Fix a base cell $x_0$ and a truncated distance function
$\rho_R(x)=\min\{d_a(x,x_0),R\}$.  Conjugation by
$W_\theta=e^{\theta\rho_R}$ changes a nonzero matrix entry of $H$ by a factor whose modulus differs from one by at most $e^{\theta R_0a}-1$.  Hence
\[
 \Norm{W_\theta HW_\theta^{-1}-H}_{\ell^2\to\ell^2}
 \le \Norm{W_\theta HW_\theta^{-1}-H}_{0,1}
 \le M_0(e^{\theta R_0a}-1).
\]
For $\theta=\eta_0/a$ the right side is at most $\delta/2$.  The resolvent identity then gives
\[
 \Norm{W_\theta(z-H)^{-1}W_\theta^{-1}}_{2\to2}\le 2/\delta.
\]
Applying this estimate to a unit vector at $y$ and evaluating its component at $x$ yields the pointwise Combes--Thomas bound
\[
 \Norm{(z-H)^{-1}(x,y)}\le {2\over\delta}e^{-\eta_0d_a(x,y)/a}.
\]
Choose $0\le\eta<\eta_0/a$.  By \eqref{eq:v1-volume-growth}, summing
$e^{-(\eta_0/a-\eta)d_a(x,y)}$ over $y$ gives a constant depending only on the displayed stocks and the gap between the two exponents, not on the period.  The same argument for columns proves \eqref{eq:v1-CT-Schur}.  The truncation $R\to\infty$ is harmless because the finite carrier is finite.
\end{proof}

\subsection{Sign function and all fixed source derivatives}

Choose two positively oriented contours $\Gamma_+$ and $\Gamma_-$ surrounding respectively $[g_*,M_*]$ and $[-M_*,-g_*]$, each at distance at least $\delta_*>0$ from these intervals and from the other component.  Put $f(z)=+1$ inside $\Gamma_+$ and $f(z)=-1$ inside $\Gamma_-$.  Then
\begin{equation}
 \eps=f(H_W)
 ={1\over 2\pi i}\int_{\Gamma_+\cup\Gamma_-}
 f(z)(z-H_W)^{-1}\dd z.
 \label{eq:v1-sign-contour}
\end{equation}

\begin{theorem}[Uniform quasilocality of the native overlap sign]
\label{thm:v1-sign-locality}
Suppose $H_W$ lies in the native local Wilson class.  For every fixed source order $r$ there exists $\eta_r>0$ such that for every multi-index $I$ with $|I|\le r$,
\begin{equation}
 \boxed{
 \sup_{\zeta\in\cU}\Norm{\partial_I\eps(\zeta)}_{\eta_r,1}
 \le C_I(g_*,M_*,R_W,\{h_B: \varnothing\ne B\subseteq I\}).}
 \label{eq:v1-sign-derivative-bound}
\end{equation}
The constants are independent of the volume and shell number.  The same bound holds for $\widehat P_-=(\Id+\eps)/2$.  Gauge and local-spin covariance are preserved by the functional calculus and by every source derivative.
\end{theorem}
\begin{proof}
For $I=\varnothing$, use \eqref{eq:v1-sign-contour} and Lemma~\ref{lem:v1-CT} on the fixed contours.  Their distance from the spectrum is uniformly positive by \eqref{eq:v1-GT-Wilson-window}, so the contour resolvents possess a common weighted Schur bound.

For a nonempty ordered source list $I$, repeated differentiation of
$R(z)=(z-H_W)^{-1}$ gives a finite sum over ordered set partitions
$I=B_1\sqcup\cdots\sqcup B_k$ of terms
\begin{equation}
 R(z)(\partial_{B_1}H_W)R(z)\cdots
       (\partial_{B_k}H_W)R(z),
 \label{eq:v1-resolvent-derivative-word}
\end{equation}
with the usual multiplicities generated when derivatives hit an already differentiated $H_W$.  This follows inductively from
$\partial_AR=R(\partial_AH_W)R$.  Every derivative vertex in
\eqref{eq:v1-resolvent-derivative-word} is fixed range with finite weighted Schur norm for sufficiently small common $\eta_r$, and every resolvent has the bound of Lemma~\ref{lem:v1-CT}.  Lemma~\ref{lem:v1-Schur-algebra} therefore bounds each word independently of volume.  There are finitely many words at fixed $r$.  Integrating them over the fixed contours gives \eqref{eq:v1-sign-derivative-bound}.

If $H_W(g\cdot\zeta)=U_gH_W(\zeta)U_g^{-1}$, holomorphic functional calculus gives
$f(H_W(g\cdot\zeta))=U_gf(H_W(\zeta))U_g^{-1}$.  Differentiation gives the corresponding covariant source identities.  The same argument applies to local spin-frame similarities.
\end{proof}

\begin{corollary}[Quasilocal overlap/Yukawa derivatives before annular scaling]
\label{cor:v1-overlap-raw-locality}
At every fixed source order,
\begin{equation}
 \Norm{\partial_ID_{\ov}}_{\eta,1}
 \le a^{-1}C_I,
 \qquad
 \Norm{\partial_ID_\chi}_{\eta,1}
 \le a^{-1}C'_I+y_I.
 \label{eq:v1-raw-overlap-bound}
\end{equation}
The estimate is volume-uniform but is \emph{not} yet an acceptable continuum annular estimate because of the explicit $a^{-1}$ scale.
\end{corollary}
\begin{proof}
Differentiate \eqref{eq:v1-GT-overlap}--\eqref{eq:v1-GT-chiral-block}, use Theorem~\ref{thm:v1-sign-locality}, the fixed finite matrices $P_+$ and $\gamma_5$, and Lemma~\ref{lem:v1-Schur-algebra}.  The Yukawa derivatives use (W5).
\end{proof}

\begin{firewall}[Why the gap/locality theorem does not prove E1]
\label{fire:v1-a-inverse}
The $a^{-1}$ in \eqref{eq:v1-raw-overlap-bound} is not a defect in the overlap construction; on the physical low-momentum branch it is cancelled by the first-order zero of the overlap numerator.  That cancellation is a statement about the \emph{native flat symbol and its Wilson completion}.  It cannot be inferred from a spectral gap and finite range alone.  Therefore Theorem~\ref{thm:v1-sign-locality} proves native quasilocality and source differentiability but does not by itself prove the finite-to-ordinary Dirac comparison required by E1.
\end{firewall}

\section{A local chiral frame on the protected perturbative component}

The moving source space $\Ran\widehat P_-(\zeta)$ must be framed before a scalar determinant representative and its differentiated line sources are estimated.  A fixed basis in the ambient carrier is not a legitimate frame of a moving subspace.

\subsection{Kato transport}

Fix a base point $\zeta_0\in\cU$.  For $\zeta$ in a star-shaped subchart, write
$P(t)=\widehat P_-(\zeta_0+t(\zeta-\zeta_0))$ and solve
\begin{equation}
 \dot U(t)=[\dot P(t),P(t)]U(t),
 \qquad U(0)=\Id.
 \label{eq:v1-Kato-ODE}
\end{equation}

\begin{lemma}[Kato intertwining]
\label{lem:v1-Kato}
The solution of \eqref{eq:v1-Kato-ODE} is unitary and satisfies
\begin{equation}
 P(t)U(t)=U(t)P(0).
 \label{eq:v1-Kato-intertwine}
\end{equation}
Hence $U(1)$ transports an orthonormal basis of $\Ran P(0)$ to one of $\Ran P(1)$.
\end{lemma}
\begin{proof}
Since $P=P^*=P^2$, the commutator $[\dot P,P]$ is anti-Hermitian, so
$\frac d{dt}(U^*U)=0$.  Next differentiate
$F(t)=P(t)U(t)-U(t)P(0)$.  From $P\dot PP=0$, obtained by differentiating $P^2=P$ and multiplying by $P$ on both sides, one verifies
\[
 \dot P+P[\dot P,P]-[\dot P,P]P=0.
\]
Substitution into $\dot F$ gives
$\dot F=[\dot P,P]F$ with $F(0)=0$, hence $F=0$.
\end{proof}

\begin{theorem}[Protected local chiral line frame with source bounds]
\label{thm:v1-local-line-frame}
At every fixed source order $r$, after shrinking the perturbative coefficient chart if necessary, Kato transport defines a smooth local frame of
$\Ran\widehat P_-(\zeta)$ and hence a smooth nonzero frame $\lambda_\chi(\zeta)$ of
\begin{equation}
 \det(P_+\cH_a)\otimes\det(\widehat P_-(\zeta)\cH_a)^*.
 \label{eq:v1-det-line}
\end{equation}
Its normalized local connection/source coefficients through order $r$ are finite sums of products of derivatives of $\widehat P_-$ and obey volume-uniform quasilocal bounds inherited from \eqref{eq:v1-sign-derivative-bound} after one source root is fixed.

This is a \emph{local perturbative frame theorem}.  It makes no assertion about holonomy around arbitrary large gauge loops or about crossing the divisor $\det D_\chi=0$.
\end{theorem}
\begin{proof}
Theorem~\ref{thm:v1-sign-locality} gives smoothness of $P(\zeta)$ and finite bounds for every required derivative.  The ODE \eqref{eq:v1-Kato-ODE} has a unique smooth solution on a sufficiently small finite parameter chart.  Lemma~\ref{lem:v1-Kato} constructs the moving orthonormal frame.

Fix an oriented orthonormal basis $e_1^0,\ldots,e_m^0$ of $\Ran P(0)$ and a fixed oriented basis $f_1,\ldots,f_m$ of $P_+\cH_a$.  Put
\[
 e_k(\zeta)=U(1;\zeta)e_k^0,
 \qquad
 \lambda_\chi(\zeta)=
 (f_1\wedge\cdots\wedge f_m)\otimes
 (e_1(\zeta)\wedge\cdots\wedge e_m(\zeta))^*.
\]
This is a nonzero local line frame.  Differentiating the ODE in the source variables gives an inhomogeneous linear ODE whose coefficients are finite sums of products of $\partial_BP$.  Duhamel iteration therefore expresses every fixed derivative of $U$ as a finite time-ordered sum of such products.  The rooted source norm fixes one differentiated local insertion; the remaining kernels are composed in the weighted Schur algebra, so Theorem~\ref{thm:v1-sign-locality} and Lemma~\ref{lem:v1-Schur-algebra} give the claimed volume-uniform bounds.  The construction uses one chosen radial path in a local chart and therefore does not identify global line holonomy.
\end{proof}

\section{Hard-shell pivot: what the global chiral floor does and does not imply}

The global lower singular-value bound in \eqref{eq:v1-GT-chiral-block} does not automatically imply that an arbitrary compression to a hard subspace is invertible.  This distinction is load-bearing for exact Schur elimination.

\subsection{Compression can destroy invertibility}

\begin{lemma}[A global gap does not imply a compressed hard pivot]
\label{lem:v1-compression-counterexample}
There exists an invertible matrix $D$ with $s_{\min}(D)=1$ and an orthogonal projection $Q$ such that the compression $QDQ:\Ran Q\to\Ran Q$ is zero.
\end{lemma}
\begin{proof}
Take
$D=\begin{psmallmatrix}0&1\\1&0\end{psmallmatrix}$ and
$Q=\begin{psmallmatrix}1&0\\0&0\end{psmallmatrix}$.
Then $D$ is unitary while $QDQ=0$ on $\Ran Q$.
\end{proof}

Thus the Grand Tensor chiral floor $s_{\min}(D_\chi)\ge\tau_*$ protects the full finite chiral operator and its line section, but an E1 shell integration additionally requires a \emph{hard-block} pivot theorem.

\subsection{Annular hard-pivot stability}

Let $Q_j$ be the declared smooth hard annulus projector/multiplier at scale $\Lambda_j$, with complement $\bar Q_j=\Id-Q_j$.  Write the native chiral operator in block form
\begin{equation}
 D_\chi=
 \begin{pmatrix}
 A_j & B_j\\ C_j&E_j
 \end{pmatrix},
 \qquad
 A_j=Q_jD_\chi Q_j\big|_{\Ran Q_j}.
 \label{eq:v1-hard-block}
\end{equation}

\begin{theorem}[Hard-pivot inheritance from an annular principal bound]
\label{thm:v1-hard-pivot}
Suppose a reference annular operator $A_j^{(0)}$ on $\Ran Q_j$ satisfies
\begin{equation}
 s_{\min}(A_j^{(0)})\ge c_D\Lambda_j
 \label{eq:v1-reference-pivot}
\end{equation}
and the complete native perturbation, including gauge/spin transport, coframe variation, Yukawa insertion and the difference between the exact native hard block and the reference annular block, obeys
\begin{equation}
 \Norm{A_j-A_j^{(0)}}_{2\to2}
 \le \theta c_D\Lambda_j,
 \qquad 0\le\theta<1.
 \label{eq:v1-pivot-perturbation}
\end{equation}
Then
\begin{equation}
 \boxed{
 s_{\min}(A_j)\ge(1-\theta)c_D\Lambda_j,
 \qquad
 \Norm{A_j^{-1}}_{2\to2}\le{1\over(1-\theta)c_D\Lambda_j}.}
 \label{eq:v1-hard-pivot-bound}
\end{equation}
If, in addition, $A_j$ and all required source derivatives lie in one weighted Schur algebra and the inverse has a uniform weighted Schur bound at zero source, then every fixed mixed source derivative of $A_j^{-1}$ is bounded by
\begin{equation}
 \boxed{
 \Norm{\partial_IA_j^{-1}}_{\eta,1}
 \le
 \sum_{k=1}^{|I|}
 \sum_{I=B_1\sqcup\cdots\sqcup B_k}^{\rm ord}
 \Norm{A_j^{-1}}_{\eta,1}^{k+1}
 \prod_{q=1}^k\Norm{\partial_{B_q}A_j}_{\eta,1}.}
 \label{eq:v1-inverse-derivative-bound}
\end{equation}
For $I=\varnothing$ the right side is interpreted as $\Norm{A_j^{-1}}_{\eta,1}$.
\end{theorem}
\begin{proof}
For $v\in\Ran Q_j$,
\[
 \Norm{A_jv}
 \ge \Norm{A_j^{(0)}v}-\Norm{(A_j-A_j^{(0)})v}
 \ge(1-\theta)c_D\Lambda_j\Norm v,
\]
which proves \eqref{eq:v1-hard-pivot-bound}.  Differentiating
$A_jA_j^{-1}=\Id$ gives
$\partial_AA_j^{-1}=-A_j^{-1}(\partial_AA_j)A_j^{-1}$.
Repeated differentiation generates exactly the ordered set-partition words in
\eqref{eq:v1-inverse-derivative-bound}.  Lemma~\ref{lem:v1-Schur-algebra} bounds every word.
\end{proof}

\begin{remark}[What must be shown rather than assumed]
The nontrivial E1 content in \eqref{eq:v1-reference-pivot}--\eqref{eq:v1-pivot-perturbation} is that the \emph{native} overlap/Wilson symbol on the hard annulus has the correct first-order physical branch and that the gauge/coframe/Yukawa perturbations are small relative to $\Lambda_j$ on the protected perturbative component.  The full-operator floor $\tau_*$ alone cannot replace this annular statement.
\end{remark}

\section{Finite-to-ordinary comparison: a transfer theorem}

We now isolate the precise symbol statement that is sufficient for E1.  This section is formulated so that the only unproved input is a comparison for the native \emph{Wilson/overlap symbol itself}; once that input is established, the inverse and kernel estimates follow without further regulator substitution.

\subsection{Normalized annular symbol seminorms}

Let
\begin{equation}
 \cA_j=\{p:c_-\Lambda_j\le |p|\le c_+\Lambda_j\},
 \qquad
 a\Lambda_j\le c_0<1,
 \label{eq:v1-annulus}
\end{equation}
with a smooth profile $\chi_j(p)=\chi(p/\Lambda_j)$ supported in a slightly larger annulus.  Let $d=4$.

For a matrix symbol $F(p)$ define
\begin{equation}
 \Norm{F}_{j,m}^{(q)}
 :=\max_{|\alpha|\le m}
 \sup_{p\in\supp\chi_j}
 \Lambda_j^{|\alpha|-q}\Norm{\partial_p^\alpha F(p)}.
 \label{eq:v1-symbol-seminorm}
\end{equation}
Thus a first-order Dirac symbol has $q=1$, its inverse has $q=-1$, and a lattice-continuum error of size $a\Lambda_j^2$ has normalized size $a\Lambda_j$ in the $q=1$ seminorm.

\begin{assumption}[Native Wilson-to-ordinary annular row]
\label{ass:v1-native-symbol-row}
There exists an ordinary reference chiral symbol $d_{\chi,0}(p)$ on the same physical branch and, for every fixed source multi-index $I$ required by the two-loop packet, a native flat/background symbol $d_{\chi,a}^{(I)}(p)$ obtained by differentiating the \emph{actual} finite Wilson/overlap/Yukawa operator, such that for some $m>s+5$,
\begin{align}
 \Norm{d_{\chi,a}}_{j,m}^{(1)}
 +\Norm{d_{\chi,0}}_{j,m}^{(1)}&\le C_m,
 \label{eq:v1-symbol-size}\\
 \Norm{d_{\chi,a}-d_{\chi,0}}_{j,m}^{(1)}&\le C_m a\Lambda_j,
 \label{eq:v1-symbol-error}\\
 s_{\min}(d_{\chi,a}(p))+s_{\min}(d_{\chi,0}(p))&\ge 2c_D\Lambda_j
 \quad(p\in\supp\chi_j).
 \label{eq:v1-symbol-pivot}
\end{align}
For differentiated source vertices, the same statements hold after division by their declared normalized source stocks.  Gauge parallel transport, spin/coframe transport and Higgs/Yukawa insertions are those of the native stencil; no continuum transporter is inserted on the left side.
\end{assumption}

The assumption is intentionally explicit.  It is precisely the point at which the actual Wilson stencil and the actual geometric link/spin completion enter.

\subsection{Inverse symbol comparison}

\begin{lemma}[Annular inverse symbol bounds]
\label{lem:v1-inverse-symbol}
Under Assumption~\ref{ass:v1-native-symbol-row}, for every $|\alpha|\le m$,
\begin{align}
 \Norm{\partial_p^\alpha d_{\chi,a}^{-1}(p)}
 +\Norm{\partial_p^\alpha d_{\chi,0}^{-1}(p)}
 &\le C_\alpha\Lambda_j^{-1-|\alpha|},
 \label{eq:v1-inverse-symbol-size}\\
 \Norm{\partial_p^\alpha
 (d_{\chi,a}^{-1}-d_{\chi,0}^{-1})(p)}
 &\le C_\alpha a\Lambda_j^{-|\alpha|}.
 \label{eq:v1-inverse-symbol-error}
\end{align}
The analogous fixed source derivatives are bounded by finite polynomials of the normalized native vertex stocks.
\end{lemma}
\begin{proof}
The zeroth inverse bound follows from \eqref{eq:v1-symbol-pivot}.  Differentiate
$d^{-1}d=\Id$.  Every $p$-derivative of $d^{-1}$ is a finite sum of words
\[
 d^{-1}(\partial^{\beta_1}d)d^{-1}\cdots
 (\partial^{\beta_k}d)d^{-1},
 \qquad
 \beta_q\ne0,
 \quad \sum_q\beta_q=\alpha.
\]
Using \eqref{eq:v1-symbol-size} and the zeroth inverse bound gives one factor
$\Lambda_j^{-1-|\alpha|}$.

For the difference use the resolvent identity
\begin{equation}
 d_{\chi,a}^{-1}-d_{\chi,0}^{-1}
 =-d_{\chi,a}^{-1}(d_{\chi,a}-d_{\chi,0})d_{\chi,0}^{-1}.
 \label{eq:v1-inverse-resolvent-id}
\end{equation}
At zeroth order the right side is
$O(\Lambda_j^{-1})O(a\Lambda_j^2)O(\Lambda_j^{-1})=O(a)$.
Differentiate \eqref{eq:v1-inverse-resolvent-id}.  If a total of $|\alpha|$ momentum derivatives are distributed among the three factors, the engineering exponents add to $-|\alpha|$ after the two inverse factors and the one $q=1$ error factor are combined.  This proves \eqref{eq:v1-inverse-symbol-error}.  Source derivatives are handled by the same ordered-partition inverse formula as in \eqref{eq:v1-inverse-derivative-bound}; fixed source order leaves finitely many words.
\end{proof}

\subsection{Kernel comparison}

Let $S_{a,j}$ and $S_{0,j}$ be the inverse Fourier kernels of
$\chi_jd_{\chi,a}^{-1}$ and $\chi_jd_{\chi,0}^{-1}$ in the physical lattice and ordinary measures respectively, periodized on the same physical torus when a finite-volume comparison is made.

\begin{theorem}[Annular finite-to-ordinary quasilocal comparison]
\label{thm:v1-kernel-comparison}
Under Assumption~\ref{ass:v1-native-symbol-row}, for every fixed $s\ge0$ and sufficiently many symbol derivatives $m>s+5$,
\begin{align}
 a^4\sup_x\sum_y
 (1+\Lambda_jd_a(x,y))^s\Norm{S_{a,j}(x,y)}
 &\le C_s\Lambda_j^{-1},
 \label{eq:v1-shell-kernel-row}\\
 a^4\sup_x\sum_y
 (1+\Lambda_jd_a(x,y))^s
 \Norm{S_{a,j}(x,y)-S_{0,j}(x,y)}
 &\le C_s a.
 \label{eq:v1-shell-kernel-diff}
\end{align}
The corresponding column estimates hold.  Every fixed normalized source derivative obeys the same statements with constants replaced by finite polynomials of its vertex stocks.  The estimates are uniform in volume and shell number as long as $a\Lambda_j\le c_0$ and the annular pivot margin stays fixed.
\end{theorem}
\begin{proof}
Multiply \eqref{eq:v1-inverse-symbol-size}--\eqref{eq:v1-inverse-symbol-error} by the smooth profile $\chi_j$.  Leibniz' rule and
$\partial^\beta\chi_j=O(\Lambda_j^{-|\beta|})$ preserve the same derivative scales.  After $p=\Lambda_jq$, the native inverse multiplier has the form
$\Lambda_j^{-1}m_a(q)$ with $m_a$ supported in a fixed annulus and uniformly bounded in $C^m$; the difference has the form
$a\,r_a(q)$ with the same property.

On the infinite lattice, discrete Fourier inversion and $m$ integrations by parts in the dimensionless variable $q$ give
\begin{align*}
 \Norm{S_{a,j}(x,y)}
 &\le C_N\Lambda_j^3(1+\Lambda_jd_a(x,y))^{-N},\\
 \Norm{S_{a,j}(x,y)-S_{0,j}(x,y)}
 &\le C_Na\Lambda_j^4(1+\Lambda_jd_a(x,y))^{-N},
\end{align*}
for every $N\le m$.  The support lies strictly inside the Brillouin zone because $a\Lambda_j\le c_0<1$, so the integrations by parts have no seam contribution.  Choose $N>s+4$.  The physical lattice sum satisfies
\[
 a^4\sum_y(1+\Lambda_jd_a(x,y))^{-N+s}
 \le C\Lambda_j^{-4}
\]
uniformly in $a$ by \eqref{eq:v1-volume-growth}.  Multiplication gives
\eqref{eq:v1-shell-kernel-row}--\eqref{eq:v1-shell-kernel-diff}.

The finite torus kernel is the periodization of the absolutely summable infinite-lattice kernel.  Shortest periodic distance is bounded by the corresponding unwrapped distance, so the same row and column estimates hold without a volume factor.  Fixed source derivatives are finite sums of inverse/vertex words whose normalized symbol derivatives obey the same engineering estimates by Lemma~\ref{lem:v1-inverse-symbol}; apply the identical Fourier argument term by term.
\end{proof}

\begin{corollary}[Rooted connected-graph compatibility]
\label{cor:v1-rooted-graph}
The shell kernels of Theorem~\ref{thm:v1-kernel-comparison} satisfy the row/column input required by the rooted connected-contraction theorem of the earlier perturbative notes.  Hence, once the native source vertices have fixed rooted stocks, every fixed connected fermionic shell graph and every fixed mixed boson--fermion graph containing these hard fermion lines has a volume-uniform rooted bound.  Replacing one native hard fermion line by its ordinary comparison changes the graph by $O(a\Lambda_j)$ relative to the engineering size of that line, before the usual local Taylor subtraction.
\end{corollary}
\begin{proof}
Equation \eqref{eq:v1-shell-kernel-row} is exactly the weighted row/column summability required to orient a spanning tree through the graph.  The connected-contraction proof is insensitive to fermionic signs after absolute values are taken.  Equation \eqref{eq:v1-shell-kernel-diff} gives $O(a)$ in place of the native $O(\Lambda_j^{-1})$ propagator stock, hence the relative factor $a\Lambda_j$.  All remaining factors have the inherited fixed rooted stocks.
\end{proof}

\section{Gauge, coframe, Higgs and antifield derivatives do not create a new locality problem}

E1 lists several source types separately because their covariance and contact terms must all be retained.  Analytically, once the native Wilson stencil itself has the required local derivatives, they enter one common derivative calculus.

\begin{theorem}[Common fixed-order source calculus]
\label{thm:v1-common-source-calculus}
Assume the native Wilson-to-ordinary annular row, and let the source list contain any finite collection of:
\begin{enumerate}[label=\textup{(\alph*)},leftmargin=2.8em]
\item internal gauge-link or gauge-potential sources;
\item coframe, inverse-coframe, spin-connection, lapse/shift or metric-chart sources entering the same finite spin operator;
\item Higgs and Yukawa sources entering $\cY$;
\item ghost, antifield and composite/contact sources which enter the finite chiral operator or its blocking maps locally;
\item local determinant-line frame sources induced by Theorem~\ref{thm:v1-local-line-frame}.
\end{enumerate}
At every fixed mixed source order, the hard inverse, overlap projector, local Kato frame, Schur blocks and line logarithmic derivatives are finite sums of words built from:
\begin{equation}
 A_j^{-1},\quad
 \partial_BH_W,\quad
 \partial_B\cY,\quad
 \partial_BQ_j,\quad
 \partial_B\widehat P_-,
 \label{eq:v1-source-word-bank}
\end{equation}
with $B$ a nonempty sub-multi-index.  Consequently the weighted/rooted bounds close under one finite polynomial of the primitive native stocks.  No source type requires a separate nonlocal inverse provided its entry in \eqref{eq:v1-source-word-bank} is local/quasilocal on the same chart.
\end{theorem}
\begin{proof}
The sign/projector derivatives are Theorem~\ref{thm:v1-sign-locality}; inverse derivatives are Theorem~\ref{thm:v1-hard-pivot}; the line frame is Theorem~\ref{thm:v1-local-line-frame}.  Differentiating a block compression adds derivatives of $Q_j$ and of $D_\chi$.  Differentiating the Schur complement adds derivatives of $A_j^{-1}$ and the off-diagonal blocks.  Each operation is product, sum, or fixed finite contraction.  Lemma~\ref{lem:v1-Schur-algebra} therefore closes the weighted Schur norm, while the inherited rooted connected-contraction theorem closes the external source kernels.  The list is finite at fixed source order, so the resulting bound is a finite polynomial.
\end{proof}

\begin{remark}[Covariance is not optional]
The theorem is analytic.  The gauge and local-spin transformation laws still have to be those of the native stencil.  A small operator derivative with the wrong covariance would not be an admissible E1 vertex.  Definition~\ref{def:v1-native-Wilson-class}(W4) therefore remains a structural hypothesis, not a consequence of the norm estimate.
\end{remark}

\section{Differentiated determinant/Pfaffian line estimates without dropping mixed-shell contacts}

Exact Schur factorization and its differentiated line identities are inherited from the prior checkpoint.  Here we prove the quantitative bounds that follow once the hard pivot and source calculus are available.

\subsection{One hard determinant block}

Let $A(\zeta)$ be an invertible finite hard chiral block in the local line frame.  Put
$\ell_A=\log\det A$, with the branch fixed on the nonzero-pivot component.

\begin{lemma}[All fixed determinant derivatives are inverse--vertex words]
\label{lem:v1-logdet-words}
For every nonempty finite source set $I$,
\begin{equation}
 \partial_I\ell_A
 =\sum_{\pi\in\Part(I)}\ 
   \sum_{\operatorname{cyc}(\pi)}
   c_{\pi,\operatorname{cyc}}
   \Tr\!
   \left(A^{-1}(\partial_{B_1}A)\cdots
         A^{-1}(\partial_{B_k}A)\right),
 \label{eq:v1-logdet-word}
\end{equation}
where $\pi=\{B_1,\ldots,B_k\}$ runs over set partitions of $I$, the second sum runs over finitely many cyclic orderings, and the coefficients are integers depending only on the partition/order.  In particular
\begin{equation}
 \Abs{\partial_I\ell_A}_{\rm root}
 \le C_I\sum_{\pi\in\Part(I)}
 \Norm{A^{-1}}_{\eta,1}^{|\pi|}
 \prod_{B\in\pi}\Norm{\partial_BA}_{\eta,1},
 \label{eq:v1-logdet-bound}
\end{equation}
after one local source insertion is used as the root and the trace density is taken in the same normalized convention as the shell action.
\end{lemma}
\begin{proof}
The first derivative is Jacobi's identity
$\partial_A\ell_A=\Tr(A^{-1}\partial_AA)$.  Differentiate repeatedly.  Every derivative either joins an existing derivative block $\partial_BA$ or hits an inverse, in which case
$\partial A^{-1}=-A^{-1}(\partial A)A^{-1}$ creates a new inverse--vertex pair.  Induction gives precisely the finite partition/cyclic word family \eqref{eq:v1-logdet-word}.  To obtain \eqref{eq:v1-logdet-bound}, anchor one differentiated local vertex.  Sum the remaining internal coordinates cyclically using the row/column weighted Schur bounds.  This is the one-loop trace version of the rooted spanning-tree estimate and has no volume factor in the normalized local trace density.
\end{proof}

\subsection{Schur return and mixed hard/soft contacts}

With the block notation \eqref{eq:v1-hard-block}, the exact soft Schur operator is
\begin{equation}
 S_j=E_j-C_jA_j^{-1}B_j.
 \label{eq:v1-Schur-soft}
\end{equation}

\begin{proposition}[No differentiated Schur contact may be discarded]
\label{prop:v1-Schur-contacts}
For every nonempty source multi-index $I$,
\begin{equation}
 \partial_IS_j
 =\partial_IE_j
 -\sum_{I=I_C\sqcup I_0\sqcup I_B}
 c(I_C,I_0,I_B)
 (\partial_{I_C}C_j)(\partial_{I_0}A_j^{-1})
 (\partial_{I_B}B_j),
 \label{eq:v1-Schur-contact-expansion}
\end{equation}
where empty subindices are allowed and $\partial_{I_0}A_j^{-1}$ is the full ordered-partition expansion of Theorem~\ref{thm:v1-hard-pivot}.  Thus derivatives in which distinct source labels hit $C_j$, the interior hard inverse, and $B_j$ are genuine mixed-shell contact terms.  Under the E1 source stocks they satisfy one common finite polynomial weighted/rooted bound.
\end{proposition}
\begin{proof}
Equation \eqref{eq:v1-Schur-contact-expansion} is the multivariate Leibniz rule applied to \eqref{eq:v1-Schur-soft}.  The inverse derivative formula is exact.  Every term contains only factors from the common source word bank \eqref{eq:v1-source-word-bank}; apply Lemma~\ref{lem:v1-Schur-algebra} and Theorem~\ref{thm:v1-common-source-calculus}.  Since the expansion is an identity, deleting any distribution of source labels would change the differentiated Schur operator.
\end{proof}

\subsection{Many shells: exact Bell-polynomial product estimate}

Let $\sigma_j$ denote the nonzero local determinant-line section contributed by shell $j$ in the inherited scale cocycle, expressed in the local frame of Theorem~\ref{thm:v1-local-line-frame}.  Write
$g_j=\log\sigma_j$ on the chosen perturbative branch.  For a finite interval of shells $[j,n)$,
\begin{equation}
 \sigma_{[j,n)}=\bigotimes_{m=j}^{n-1}\sigma_m,
 \qquad
 g_{[j,n)}=\sum_{m=j}^{n-1}g_m.
 \label{eq:v1-line-product}
\end{equation}

For the following theorem fix the finite external source panel required by the coefficientwise target.  Denote by $\Norm{\cdot}_{\rm pan}$ any fixed cross norm on its finite tensor algebra, chosen submultiplicative.  Connected logarithmic line kernels continue to use the rooted norm.  This distinction is essential: the rooted norm of an extensive connected action is not an algebra norm, whereas the finite-panel tensor algebra may legitimately carry products of disconnected source contacts.

\begin{theorem}[Differentiated line-product bound with all mixed-shell contacts]
\label{thm:v1-line-product}
Fix source order $r$ and the finite source panel just declared.  Suppose that after the common local Taylor/EFT assignment at each shell the renormalized logarithmic line coefficients obey
\begin{equation}
 \Norm{\partial_I g_m^{\remn}}_{\rm pan}
 \le c_I\rho_m,
 \qquad
 0<|I|\le r,
 \qquad
 \sum_{m\ge j}\rho_m<\infty.
 \label{eq:v1-line-log-summable}
\end{equation}
Then every normalized derivative of the composed section is uniformly bounded in the number of shells:
\begin{equation}
 \boxed{
 \Norm{\sigma_{[j,n)}^{-1}\partial_I\sigma_{[j,n)}}_{\rm pan}
 \le
 \sum_{\pi\in\Part(I)}
 \prod_{B\in\pi}
 \left(c_B\sum_{m=j}^{n-1}\rho_m\right).}
 \label{eq:v1-line-Bell-bound}
\end{equation}
In particular the right side is bounded uniformly as $n\to\infty$.  The products over blocks $B\in\pi$ are precisely the mixed-shell line contacts generated by differentiating the exact tensor product; no contact is discarded.

If E1 supplies $\rho_m=C\mu/\Lambda_m$ with $\Lambda_m=b^m\lambda$, then
\begin{equation}
 \sup_{n\ge j}
 \Norm{\sigma_{[j,n)}^{-1}\partial_I\sigma_{[j,n)}}_{\rm pan}
 \le
 \sum_{\pi\in\Part(I)}
 \prod_{B\in\pi}
 \left(c_B{Cb\over b-1}{\mu\over\Lambda_j}\right).
 \label{eq:v1-line-geometric-bound}
\end{equation}
\end{theorem}
\begin{proof}
For scalar representatives in a fixed local line frame,
$\sigma_{[j,n)}=\exp(g_{[j,n)})$.  The multivariate Fa\`a di Bruno formula for the exponential gives the exact identity
\begin{equation}
 \sigma_{[j,n)}^{-1}\partial_I\sigma_{[j,n)}
 =\sum_{\pi\in\Part(I)}
 \prod_{B\in\pi}\partial_B g_{[j,n)}.
 \label{eq:v1-Bell-identity}
\end{equation}
But
$\partial_Bg_{[j,n)}=\sum_{m=j}^{n-1}\partial_Bg_m$.
Taking the submultiplicative finite-panel norm and applying \eqref{eq:v1-line-log-summable} gives \eqref{eq:v1-line-Bell-bound}.  If one instead retains source positions as free kernel variables, \eqref{eq:v1-Bell-identity} is to be read in the disconnected tensor algebra; its connected logarithm remains in the rooted action norm.  No claim that the rooted extensive-action norm itself is an algebra norm is used.  Expanding the product of shell sums in \eqref{eq:v1-Bell-identity} assigns each block of the source partition independently to a shell; these are exactly the same-shell and mixed-shell contacts.  Finally
$\sum_{m\ge j}\mu/\Lambda_m\le \frac b{b-1}\mu/\Lambda_j$ gives \eqref{eq:v1-line-geometric-bound}.
\end{proof}

\subsection{Optional real/Majorana Pfaffian branch}

\begin{proposition}[Pfaffian derivative estimate on an oriented nonzero component]
\label{prop:v1-Pfaffian}
Let $A^M(\zeta)$ be a real skew-symmetric hard block of even rank on a chosen oriented component with
$s_{\min}(A^M)\ge\kappa_j>0$.  Then, in that local orientation frame,
\begin{equation}
 \partial_A\log\Pf A^M
 ={\tfrac12}\Tr\bigl((A^M)^{-1}\partial_AA^M\bigr),
 \label{eq:v1-Pf-first}
\end{equation}
and every higher fixed source derivative obeys the analogue of
\eqref{eq:v1-logdet-bound} with an overall factor $1/2$ at the logarithmic level.  The multi-shell product estimate of Theorem~\ref{thm:v1-line-product} therefore applies verbatim to the local Pfaffian section.

No statement is made across a Pfaffian zero or an orientation-changing loop.
\end{proposition}
\begin{proof}
On an oriented nonzero component, $(\Pf A^M)^2=\det A^M$ and the chosen sign is continuous.  Differentiate
$2\log\Pf A^M=\log\det A^M$ and apply Lemma~\ref{lem:v1-logdet-words}.  The line-product proof uses only the logarithmic derivative and therefore carries over unchanged.
\end{proof}

\section{The exact remaining E1 bottleneck}

We now separate what is proved from what would be circular to assert.

\subsection{Source audit of the attached native data}

The protected field-shell construction in the Grand Tensor supplies:
\begin{enumerate}[label=\textup{(A\arabic*)},leftmargin=3em]
\item a gauge-covariant Hermitian Wilson kernel $H_W$;
\item a uniform protected Wilson spectral window $g_*\Id\preceq |H_W|\preceq M_*\Id$;
\item the overlap sign, overlap operator and moving chiral projector;
\item the native chiral overlap--Yukawa block $D_\chi$;
\item a nonzero full chiral floor $s_{\min}(D_\chi)\ge\tau_*$;
\item finite determinant-line Jacobi identities and a local protected line section;
\item a common coframe differentiation requirement for the assembled source/stress packet.
\end{enumerate}
These data are sufficient for the functional-calculus, local-frame and line-derivative arguments above once the local stencil/source-derivative properties are read from the actual finite action.

What is not supplied in the attached statement is an explicit formula for the native $H_W$ stencil on the common coframe/spin/gauge carrier from which one can prove the annular derivative estimates \eqref{eq:v1-symbol-size}--\eqref{eq:v1-symbol-pivot}.  In particular, the protected spectral window does not determine:
\begin{enumerate}[label=\textup{(B\arabic*)},leftmargin=3em]
\item which Wilson mass/parameter branch carries the physical zero;
\item the coefficient of the first-order $i\gamma^\mu p_\mu$ term after the overlap normalization;
\item the size and tensor structure of the $O(a|p|^2)$ Wilson/spin remainder;
\item the exact coframe/spin-link dependence of that remainder;
\item the differentiated comparison when gauge, geometric, Higgs and antifield sources are inserted;
\item the hard-annular compression margin needed for exact shell Schur elimination.
\end{enumerate}
These are not details that can be supplied by replacing the finite operator with a continuum Dirac operator.

\begin{obstruction}[Gap/locality does not identify the physical branch]
\label{obs:v1-gap-not-symbol}
There exist distinct finite-range Hermitian lattice kernels with the same two-sided spectral gap and the same uniform operator norm but different low-momentum winding, different first-order symbols, or no physical first-order zero after the overlap map.  Therefore \eqref{eq:v1-GT-Wilson-window} plus finite range cannot imply Assumption~\ref{ass:v1-native-symbol-row}.
\end{obstruction}
\begin{proof}
Already in a translation-invariant two-band model, conjugating a gapped reference symbol by a momentum-dependent unitary preserves its eigenvalues and hence its spectral gap while changing the eigenprojector map.  Alternatively, add to a gapped Wilson symbol a finite-range Hermitian trigonometric polynomial of size smaller than $g_*/2$ but with a nonzero first derivative at the physical point.  The gap survives by the singular-value perturbation inequality, while the first-order overlap symbol changes.  Spectral bounds therefore do not determine the required principal-symbol matching.
\end{proof}

\subsection{Reduction theorem}

\begin{theorem}[E1 reduction to the native Wilson-symbol theorem]
\label{thm:v1-E1-reduction}
Fix the finite two-loop/source packet and the protected nonzero-pivot component.  Suppose:
\begin{enumerate}[label=\textup{(R\arabic*)},leftmargin=3em]
\item the actual native $H_W(e,\omega,U)$ is a fixed-range gauge/spin-covariant stencil with all required local source derivatives;
\item its protected spectral window \eqref{eq:v1-GT-Wilson-window} holds;
\item its exact flat/coframe-normalized Wilson/overlap symbol satisfies Assumption~\ref{ass:v1-native-symbol-row}, including the differentiated gauge/geometric/Higgs/antifield rows;
\item the annular perturbation satisfies the hard-pivot inequality \eqref{eq:v1-pivot-perturbation} on the declared perturbative component.
\end{enumerate}
Then E1 is closed at that fixed source/derivative order in the following precise sense:
\begin{enumerate}[label=\textup{(E1.\arabic*)},leftmargin=3.4em]
\item the overlap sign, moving chiral projector and all required source derivatives are cutoff-uniformly quasilocal;
\item a local Kato chiral line frame exists with cutoff-uniform rooted source bounds;
\item the hard fermion block has a nonzero pivot of order $\Lambda_j$ and all inverse source derivatives are bounded;
\item the native annular hard inverse differs from the ordinary reference inverse by $O(a)$ in the physical weighted row/column norm, with the analogous fixed source-derivative estimates;
\item every fixed connected hard-fermion/mixed graph has a volume-uniform rooted bound;
\item the inherited exact determinant/Pfaffian Schur identities acquire the quantitative derivative bounds of Lemma~\ref{lem:v1-logdet-words} and Proposition~\ref{prop:v1-Pfaffian}, with all differentiated Schur contacts retained;
\item once the usual local Taylor/EFT assignment gives a summable one-excess-scale remainder $\rho_j=O(\mu/\Lambda_j)$, the differentiated determinant/Pfaffian line products are uniformly bounded through arbitrary finite shell composition by Theorem~\ref{thm:v1-line-product}.
\end{enumerate}
No global determinant-line holonomy, no zero-mode crossing theorem, and no nonperturbative chiral measure is implied.
\end{theorem}
\begin{proof}
Items (E1.1)--(E1.2) are Theorems~\ref{thm:v1-sign-locality} and
\ref{thm:v1-local-line-frame}.  Item (E1.3) is Theorem~\ref{thm:v1-hard-pivot}.  Item (E1.4) is Theorem~\ref{thm:v1-kernel-comparison}.  Item (E1.5) is Corollary~\ref{cor:v1-rooted-graph}.  Item (E1.6) follows from Lemma~\ref{lem:v1-logdet-words} and Proposition~\ref{prop:v1-Schur-contacts}; the Pfaffian variant is Proposition~\ref{prop:v1-Pfaffian}.  Item (E1.7) is Theorem~\ref{thm:v1-line-product}.  Every input is native except the ordinary reference used only on the right-hand side of a proved comparison.  The theorem never substitutes the ordinary operator for the native hard block in the exact shell algebra.
\end{proof}

\section{How to attack the remaining native Wilson-symbol theorem}

The reduction theorem shows that repeatedly polishing determinant identities would now be circular: their exact algebra is already inherited, and their estimates are controlled once the native annular pivot/comparison is known.  The next work should therefore move one level upstream to the finite Wilson stencil.

\subsection{Required exact stencil data}

For each oriented cell edge $e=(x,x+a\hat\mu)$, let
\begin{equation}
 \cU_e(e,\omega,U)
 =\cU_e^{\rm spin}(e,\omega)\otimes\rho_{\SM}(U_e)
 \label{eq:v1-combined-link}
\end{equation}
be the actual spin--internal parallel transport used by the finite action.  The Wilson kernel must be written in the finite form
\begin{equation}
 (H_W\psi)_x
 =\sum_{|y-x|\le R_Wa} h_{xy}(e,\omega,U)\psi_y,
 \label{eq:v1-native-stencil-template}
\end{equation}
including its diagonal Wilson mass/completion.  The next proof must not infer the coefficients $h_{xy}$ from the desired continuum answer; it must read them from the native finite action.

On the constant coframe, trivial spin connection, identity gauge-link and fixed Higgs background, Fourier transformation of \eqref{eq:v1-native-stencil-template} gives the literal finite trigonometric symbol
\begin{equation}
 h_{W,a}(p)=\sum_{|z|\le R_Wa}h_{0z}\,e^{ip\cdot z}.
 \label{eq:v1-native-flat-symbol}
\end{equation}
The protected overlap symbol is then
\begin{equation}
 \eps_a(p)=h_{W,a}(p)[h_{W,a}(p)^2]^{-1/2},
 \qquad
 d_{\ov,a}(p)=a^{-1}(\Id+\gamma_5\eps_a(p)),
 \label{eq:v1-native-overlap-symbol}
\end{equation}
with the branch fixed by the actual native Wilson mass/completion.

\subsection{The exact next load-bearing theorem}

\begin{theorem}[Target theorem for V2: native Wilson annular symbol]
\label{thm:v1-next-target}
The next load-bearing theorem is to prove directly from
\eqref{eq:v1-native-stencil-template} that on the physical index-zero branch there exist $c_0,c_D>0$ such that, whenever $a\Lambda_j\le c_0$,
\begin{align}
 d_{\chi,a}(p)
 &=d_{\chi,0}(p)+r_a(p),
 \label{eq:v1-next-expansion}\\
 \Norm{\partial_p^\alpha r_a(p)}
 &\le C_\alpha a\Lambda_j^{2-|\alpha|},
 \qquad p\in\supp\chi_j,
 \label{eq:v1-next-remainder}\\
 s_{\min}(d_{\chi,a}(p))
 &\ge c_D\Lambda_j,
 \label{eq:v1-next-pivot}
\end{align}
through the momentum derivative order required by the rooted kernel moment, and with the corresponding differentiated formulas for every gauge, coframe/spin, Higgs/Yukawa, antifield and line-source insertion in the fixed two-loop packet.  The same proof must give a full-zone bound near the lattice ultraviolet cap, where $p$ itself is not a good small variable and $x=ap$ must be used instead.
\end{theorem}

Theorem~\ref{thm:v1-next-target} is \emph{not proved in V1} because the attached sources, as supplied here, do not display the exact coefficient stencil needed to evaluate \eqref{eq:v1-native-flat-symbol}.  Once that stencil is made explicit, the proof should proceed by the following finite calculations:
\begin{enumerate}[label=\textup{(N\arabic*)},leftmargin=3em]
\item compute $h_{W,a}(p)$ exactly on the constant-coframe identity-link branch;
\item locate the physical Wilson branch and prove its gap from the doubler sectors using the actual Wilson mass/completion;
\item Taylor-expand the physical branch through the derivative order consumed by Theorem~\ref{thm:v1-kernel-comparison};
\item differentiate the finite link matrices in the common coframe/spin/gauge chart rather than inserting a continuum covariant derivative by hand;
\item add the onsite/fixed-range Higgs--Yukawa block and verify that it is a lower-order perturbation on a sufficiently hard annulus;
\item establish the $x=ap$ full-zone inverse bounds on the compact Brillouin region away from the physical small-$p$ chart;
\item combine the two charts by a smooth partition, producing the hard-pivot estimate \eqref{eq:v1-pivot-perturbation} and Assumption~\ref{ass:v1-native-symbol-row} as theorems.
\end{enumerate}

\section{Anti-circularity audit}

\begin{firewall}[No continuum substitution]
At no point in V1 is the native shell integral evaluated with a continuum Dirac operator.  The ordinary symbol appears only in Assumption~\ref{ass:v1-native-symbol-row} as the comparator in a finite-to-ordinary estimate.  All exact factorization, inverses, source derivatives and line sections remain those of the finite native $D_\chi$.
\end{firewall}

\begin{firewall}[No anomaly shortcut]
The inherited Standard Model anomaly cancellations remove the nontrivial local anomaly class only after the regulator matching problem is solved.  They do not imply Theorem~\ref{thm:v1-next-target}, hard-shell quasilocality, or a bound on the determinant-line connection.
\end{firewall}

\begin{firewall}[No global line claim]
The Kato frame is defined on one protected perturbative subchart.  V1 does not identify line holonomy across arbitrary gauge loops, across disconnected Wilson sectors, or through zeros of $D_\chi$ or of an optional Majorana skew block.
\end{firewall}

\begin{firewall}[No fixed-order to all-order promotion]
All constants may depend arbitrarily on the fixed source order, local derivative order and perturbative order.  V1 does not provide a uniform large-loop estimate or convergence of the perturbative series.
\end{firewall}

\section{Status ledger and next theorem}

\subsection*{Proved in V1}
\begin{enumerate}[leftmargin=2.2em]
\item A finite-range gapped native Hermitian Wilson kernel transfers cutoff-uniform quasilocality to its overlap sign $\eps=\sgn H_W$, moving projector $\widehat P_-$, and every fixed mixed source derivative, with exact gauge/spin covariance preserved.
\item On the protected perturbative component, Kato transport constructs a legitimate local frame of the moving chiral source space and of the determinant line, with fixed-order rooted source bounds.  This is explicitly local, not a global holonomy theorem.
\item The full chiral floor does not by itself imply an invertible shell compression.  A separate annular hard-pivot theorem was proved: an $O(\Lambda_j)$ reference singular-value margin survives every native perturbation smaller than that margin, and all fixed inverse source derivatives satisfy explicit ordered-partition bounds.
\item Assuming the precise native Wilson-to-ordinary annular symbol row, the hard fermion inverse has the physical weighted row/column bound $O(\Lambda_j^{-1})$ and differs from its ordinary comparator by $O(a)$, uniformly in volume and shell number for $a\Lambda_j\le c_0$.  The same holds for every fixed normalized source derivative.
\item These shell kernels plug into the inherited rooted connected-contraction estimate, yielding volume-uniform fixed-graph fermionic and mixed boson--fermion bounds.
\item The inherited exact determinant/Pfaffian Schur identities were upgraded with quantitative fixed-source derivative bounds.  Differentiation of the Schur complement retains every mixed hard/soft contact.
\item A Bell-polynomial line-product theorem was proved: if the renormalized logarithmic line remainder at shell $j$ is summable (in particular $O(\mu/\Lambda_j)$), then all fixed differentiated determinant/Pfaffian line products are uniformly bounded through arbitrarily many ultraviolet shell steps on every fixed finite source panel, with all mixed-shell source contacts retained; the connected logarithmic coefficients retain the rooted norm.
\item E1 was reduced to one explicit native Wilson-symbol theorem plus its hard-annular perturbation margin; no continuum substitution is needed after that point.
\end{enumerate}

\subsection*{Inherited}
\begin{enumerate}[leftmargin=2.2em]
\item The V80 bosonic Einstein--$SU(3)\times SU(2)\times U(1)_Y$--Higgs two-loop local/range/source self-map at the declared finite weight/source order.
\item The three-generation chiral representation, cancellation of the five ordinary local gauge-anomaly coefficients, and even $SU(2)$ doublet parity at the scope declared by the programme checkpoint.
\item Exact determinant Schur factorization, its finite regrouping cocycle, the optional neutral-Majorana Pfaffian analogue, and their exact differentiated line-source identities.
\item The Grand Tensor protected overlap/Yukawa data: a gauge-covariant Hermitian Wilson kernel with protected gap, overlap sign/projector, chiral block, full nonzero chiral floor, same-operator determinant section, and common coframe differentiation requirement.
\item The earlier rooted connected-contraction calculus, common source normalization, forest ownership and bosonic two-loop annular estimates.
\end{enumerate}

\subsection*{Conditional}
\begin{enumerate}[leftmargin=2.2em]
\item The finite-to-ordinary kernel estimate and the quantitative hard-shell E1 closure are conditional on the native Wilson annular symbol row, because the supplied Grand Tensor excerpt gives $H_W$ abstractly rather than displaying the exact Wilson/spin stencil needed for the symbol calculation.
\item The multi-shell line-product bound requires the locally subtracted logarithmic line remainder to satisfy a summable shell estimate; the desired two-loop value is $O(\mu/\Lambda_j)$.
\item The optional Pfaffian estimate is local to a chosen nonzero oriented real-skew component.
\end{enumerate}

\subsection*{Open}
\begin{enumerate}[leftmargin=2.2em]
\item Write the exact native $H_W(e,\omega,U)$ stencil, including its Wilson mass/completion and combined spin/internal parallel transports, directly from the finite action.
\item Prove Theorem~\ref{thm:v1-next-target}: the physical annular symbol, $O(a|p|^2)$ finite-to-ordinary remainder with all required derivatives, hard-pivot margin, and the complementary full-zone Brillouin estimate.
\item Use that theorem to prove the locally subtracted chiral determinant/Pfaffian line remainder is $O(\mu/\Lambda_j)$ in the same rooted/source norm as the bosonic shell.
\item Only then pass to E2: identify the native chiral BV breaking with the standard anomaly class plus a local BV-exact regulator term and construct one common two-loop ESM restoring action/source/line update.
\end{enumerate}

\subsection*{Exact next load-bearing theorem}

\begin{center}
\fbox{\parbox{0.92\textwidth}{
\textbf{Native Wilson annular symbol theorem.}
Starting from the actual finite $H_W(e,\omega,U)$ stencil used by the Renewal regulator, prove the physical-branch expansion
$d_{\chi,a}=d_{\chi,0}+O(a|p|^2)$ on $|p|\asymp\Lambda_j$, with the required momentum and gauge/coframe/spin/Higgs/antifield/line-source derivatives, an $O(\Lambda_j)$ hard singular-value margin, and a uniform $x=ap$ full-zone bound.  This theorem must be native: no replacement of $H_W$ by a continuum Dirac operator is allowed.
}}
\end{center}

\section*{Append-only continuation marker: V1}
\addcontentsline{toc}{section}{Append-only continuation marker: V1}
\textbf{End of V1.}  Future versions in this lineage must preserve the mathematical body above and append new work.  The next file is to be named
\texttt{Renewal\_Perturbative\_ESM\_Continuum\_Research\_Notes\_V2.tex}.


\clearpage
\section{V2 continuation: an explicit native Wilson completion and coefficientwise closure of E1}
\label{sec:v2-continuation}

V1 isolated the remaining fermionic obstruction at the correct place: the protected
spectral window of an abstract Hermitian Wilson kernel does not determine its
physical low-momentum branch.  There are two possible responses.  One may keep
asking for a stronger theorem about an unspecified kernel, or one may go upstream
and specify the regulator slot itself.  The present continuation takes the second
route.

There is also a quantifier correction.  The endpoint of the programme is a
\emph{coefficientwise} perturbative continuum at fixed source order.  Such a theorem
does not require the hard fermion block to remain invertible on a cutoff-independent
finite-amplitude ball of all gauge, coframe, Higgs and antifield sources.  It requires
an invertible reference hard block and controlled source jets of every fixed order.
The stronger finite-amplitude pivot theorem of V1 remains useful when available,
but it is not load-bearing for the coefficientwise theorem.  This observation avoids
paying a spurious $a^{-1}$ operator-norm cost for a finite coframe displacement when
only its fixed perturbative derivatives are needed.

The second upstream move is to fill the abstract Wilson slot by an explicit
covariant nearest-neighbour completion.  This is a \emph{regulator definition in the
present proof-note lineage}, not an inference that the earlier manuscript had
already written these coefficients.  If a different Wilson completion is intended
elsewhere, it must be compared with the one below before the final theorem is
identified with that different regulator.

\subsection{Coefficientwise inversion needs only the reference pivot}

Let $A(\zeta)$ be a finite-dimensional or finite-cutoff operator family with source
coordinates $\zeta=(\zeta^A)$ and Taylor jets
\begin{equation}
 A_I:=\partial_I A(0).
 \label{eq:v2-source-jets}
\end{equation}
No common source radius is assumed.

\begin{theorem}[Reference-pivot coefficient theorem]
\label{thm:v2-reference-pivot}
Assume only that $A_0:=A(0)$ is invertible.  For every nonempty finite source
multi-index $I$,
\begin{equation}
 \boxed{
 \partial_I A^{-1}(0)
 =\sum_{k=1}^{|I|}(-1)^k
 \sum_{I=B_1\sqcup\cdots\sqcup B_k}^{\rm ord}
 A_0^{-1}A_{B_1}A_0^{-1}\cdots A_{B_k}A_0^{-1}.}
 \label{eq:v2-reference-inverse-jet}
\end{equation}
Likewise, after a local logarithm is chosen at the reference point,
$\partial_I\log\det A(0)$ is a finite sum of cyclic traces of the same
reference inverse and source jets.  Consequently a fixed source coefficient of a
hard-shell graph, Schur complement, determinant line, or Pfaffian line needs only:
\begin{enumerate}[label=\textup{(C\arabic*)},leftmargin=3em]
\item a reference hard inverse $A_0^{-1}$ with the required annular bound;
\item the finitely many jets $A_B$ with $0<|B|\le |I|$;
\item the fixed local line orientation/frame needed to interpret the coefficient.
\end{enumerate}
It does not require a cutoff-independent ball on which $A(\zeta)^{-1}$ is uniformly
bounded.
\end{theorem}
\begin{proof}
Differentiate $A(\zeta)A(\zeta)^{-1}=I$ at $\zeta=0$.  For one derivative the
formula is
$\partial_AA^{-1}(0)=-A_0^{-1}A_AA_0^{-1}$.  Inductively, every new derivative
either joins one existing jet or differentiates an inverse factor, in which case it
creates one further $-A_0^{-1}A_BA_0^{-1}$ insertion.  The ordered set partitions
record exactly these possibilities and prove \eqref{eq:v2-reference-inverse-jet}.
Jacobi's formula at the reference point then gives the determinant statement.
The Pfaffian statement follows from $2\log\Pf=\log\det$ on a chosen nonzero
oriented component.  Every fixed coefficient is therefore an algebraic expression
in the displayed finite jet bank.  A uniform nonzero radius in $\zeta$ would be a
stronger analytic statement and is not used.
\end{proof}

\begin{remark}[Relation to the V1 hard-pivot theorem]
Theorem~\ref{thm:v1-hard-pivot} remains the correct theorem for finite-amplitude
shell integration on a fixed perturbative ball.  Theorem~\ref{thm:v2-reference-pivot}
shows that this stronger quantifier is unnecessary for E1 at the stated
coefficientwise endpoint.  In particular, source derivatives are to be treated as
vertices around the reference hard propagator rather than absorbed into a finite
background before inversion.
\end{remark}

\subsection{A concrete covariant Wilson slot}
\label{subsec:v2-Wilson-definition}

Work first on a periodic cubical regulator branch with spacing $a$.  The local
argument is later inserted into the general finite cell carrier by the already
inherited chart/partition machinery.  Use Hermitian Euclidean Clifford matrices
$\gamma_a$ satisfying
\begin{equation}
 \{\gamma_a,\gamma_b\}=2\delta_{ab}I,
 \qquad
 \gamma_5^*=\gamma_5,
 \qquad
 \{\gamma_5,\gamma_a\}=0.
 \label{eq:v2-Clifford}
\end{equation}
Let $\Gamma_\mu(e)$ be Hermitian coframe Clifford coefficients with
\begin{equation}
 \{\Gamma_\mu,\Gamma_\nu\}=2g_{\mu\nu}I,
 \qquad
 m_g I\preceq g\preceq M_g I,
 \label{eq:v2-coframe-elliptic}
\end{equation}
for fixed $0<m_g\le M_g<\infty$ on the protected reference coframe chart.

Let $\mathcal U_{x,\mu}$ be the combined spin--internal parallel transport on the
edge $x\to x+a\hat\mu$.  After absorbing the positive site Hodge weight into the
one-particle carrier, define the unitary covariant shift
\begin{equation}
 (T_\mu^{\mathcal U}\psi)_x
 =\mathcal U_{x,\mu}\psi_{x+a\hat\mu},
 \qquad
 (T_\mu^{\mathcal U})^*=(T_\mu^{\mathcal U})^{-1}.
 \label{eq:v2-covariant-shift}
\end{equation}
For a nonconstant positive site weight the standard square-root weight ratio is
inserted in \eqref{eq:v2-covariant-shift}; all formulas below are unchanged after
that unitary rescaling.  Put
\begin{equation}
 C_\mu^{\mathcal U}:={1\over2}
 (T_\mu^{\mathcal U}-(T_\mu^{\mathcal U})^*),
 \qquad
 L_\mu^{\mathcal U}:=I-{1\over2}
 (T_\mu^{\mathcal U}+(T_\mu^{\mathcal U})^*).
 \label{eq:v2-centered-Wilson-pieces}
\end{equation}
Thus $C_\mu^{\mathcal U}$ is anti-Hermitian and $L_\mu^{\mathcal U}$ is
Hermitian.  The V2 Wilson completion is
\begin{equation}
 \boxed{
 K_{W,a}(e,\omega,U)
 :=\sum_{\mu=1}^4 {1\over2}
   \{\Gamma_\mu(e),C_\mu^{\mathcal U}\}
   +\sum_{\mu=1}^4L_\mu^{\mathcal U}-I,
 \qquad
 H_{W,a}:=\gamma_5K_{W,a}.}
 \label{eq:v2-native-Wilson}
\end{equation}
This is the Wilson choice $r=\rho=1$.  The choice is deliberate: it gives an
exact all-zone gap identity on the flat/reference branch and normalizes the physical
overlap speed to one without an additional prefactor.

\begin{proposition}[Locality, covariance and Hermiticity of the V2 completion]
\label{prop:v2-Wilson-structure}
Assume the spin transport commutes with $\gamma_5$ and the local spin-frame action
intertwines $\Gamma_\mu$ in the usual way.  Then:
\begin{enumerate}[label=\textup{(W\arabic*)},leftmargin=3em]
\item $H_{W,a}=H_{W,a}^*$;
\item $H_{W,a}$ has range one in lattice units;
\item internal gauge and local spin-frame transformations act by exact unitary
similarity;
\item every fixed derivative with respect to a smooth gauge-link, spin-connection,
coframe or metric-chart coordinate is again a fixed-range kernel.
\end{enumerate}
Hence, on every reference chart on which a spectral gap is present,
\eqref{eq:v2-native-Wilson} belongs to the native local Wilson class of
Definition~\ref{def:v1-native-Wilson-class}.
\end{proposition}
\begin{proof}
The anticommutator of a Hermitian multiplication operator with an anti-Hermitian
shift difference is anti-Hermitian.  The Wilson Laplacian and the diagonal mass are
Hermitian.  Since $\gamma_5$ anticommutes with every $\Gamma_\mu$ and commutes
with the covariant shifts,
\[
 K_{W,a}^*=\gamma_5K_{W,a}\gamma_5.
\]
Therefore $(\gamma_5K_{W,a})^*=\gamma_5K_{W,a}$.  Each term contains at most
one nearest-neighbour shift, proving range one.  Gauge/spin covariance follows by
conjugating the edge transport and the local Clifford multiplication together.
Finally, differentiating a fixed product of local coefficients and one-edge
transports cannot enlarge the range.  Smoothness of the finite coordinate maps
gives the fixed derivative stocks on every compact source panel.
\end{proof}

\begin{firewall}[What has and has not been identified]
Equation~\eqref{eq:v2-native-Wilson} fills the Wilson slot that was abstract in the
supplied Grand Tensor field-shell statement.  The statements below are therefore
native theorems for this explicit regulator completion.  They are not a proof that
an unspecified alternative $H_W$ has the same symbol.  Such an alternative would
need a finite comparison theorem.
\end{firewall}

\subsection{Exact constant-coframe Wilson spectrum}
\label{subsec:v2-flat-spectrum}

Freeze a constant coframe, set the spin/internal links to the identity, and write
\begin{equation}
 x_\mu:=ap_\mu,
 \qquad
 s_\mu(x):=\sin x_\mu,
 \qquad
 c_\mu(x):=1-\cos x_\mu,
 \qquad
 W(x):=\sum_\mu c_\mu(x),
 \qquad
 B(x):=W(x)-1.
 \label{eq:v2-trig-defs}
\end{equation}
The exact dimensionless Wilson symbol is
\begin{equation}
 k_e(x)=i\sum_\mu\Gamma_\mu s_\mu(x)+B(x)I,
 \qquad
 h_e(x)=\gamma_5k_e(x).
 \label{eq:v2-Wilson-symbol}
\end{equation}
Put
\begin{equation}
 q_e(x):=\sum_{\mu,\nu}g_{\mu\nu}s_\mu(x)s_\nu(x),
 \qquad
 R_e(x):=\sqrt{q_e(x)+B(x)^2}.
 \label{eq:v2-R-def}
\end{equation}

\begin{lemma}[Exact trigonometric Wilson identity]
\label{lem:v2-trig-gap}
For the identity coframe,
\begin{equation}
 \boxed{
 \sum_\mu\sin^2x_\mu+B(x)^2
 =1+2\sum_{\mu<\nu}c_\mu(x)c_\nu(x)\ge1.}
 \label{eq:v2-trig-gap-identity}
\end{equation}
Consequently, for every constant coframe satisfying
\eqref{eq:v2-coframe-elliptic},
\begin{equation}
 \boxed{
 \sqrt{\min\{m_g,1\}}
 \le R_e(x)
 \le \sqrt{4M_g+49}}
 \qquad (x\in[-\pi,\pi]^4).
 \label{eq:v2-Wilson-gap-exact}
\end{equation}
\end{lemma}
\begin{proof}
Since $\sin^2x_\mu=2c_\mu-c_\mu^2$,
\[
 \sum_\mu\sin^2x_\mu+(W-1)^2
 =2W-\sum_\mu c_\mu^2+W^2-2W+1
 =1+2\sum_{\mu<\nu}c_\mu c_\nu.
\]
This proves the first identity.  Ellipticity gives
$q_e\ge m_g\sum_\mu s_\mu^2$, hence
$q_e+B^2\ge\min\{m_g,1\}(\sum s_\mu^2+B^2)$.
For the upper bound use $\sum s_\mu^2\le4$ and
$-1\le B\le7$.
\end{proof}

\begin{theorem}[Exact Wilson gap and physical overlap branch]
\label{thm:v2-exact-overlap-branch}
For the constant-coframe reference branch,
\begin{equation}
 h_e(x)^2=R_e(x)^2I,
 \qquad
 \varepsilon_e(x)=\sgn h_e(x)={h_e(x)\over R_e(x)},
 \label{eq:v2-exact-sign}
\end{equation}
and the overlap symbol is
\begin{equation}
 \boxed{
 d_{{\rm ov},a}(p;e)
 ={1\over a}\left[
 I+{i\Gamma\!\cdot s(x)+B(x)I\over R_e(x)}
 \right].}
 \label{eq:v2-exact-overlap-symbol}
\end{equation}
It has exactly one zero in the Brillouin zone, at $x=0$.  In particular no
Wilson corner doubler survives the overlap map.

Moreover
\begin{equation}
 \boxed{
 d_{{\rm ov},a}(p;e)^*d_{{\rm ov},a}(p;e)
 ={2\over a^2}\left(1+{B(x)\over R_e(x)}\right)I.}
 \label{eq:v2-overlap-singular-exact}
\end{equation}
\end{theorem}
\begin{proof}
The Clifford relation gives
$k_e^*k_e=k_ek_e^*=(q_e+B^2)I$, because the mixed scalar--Clifford terms cancel.
The $\gamma_5$-Hermiticity of $k_e$ then gives $h_e^2=R_e^2I$, and
Lemma~\ref{lem:v2-trig-gap} permits the exact sign formula.  Since
$\gamma_5\varepsilon_e=k_e/R_e$, the overlap formula follows.

A zero of $d_{{\rm ov},a}$ would require an eigenvalue of $k_e/R_e$ to equal
$-1$.  The eigenvalues of $k_e/R_e$ have real part $B/R_e$ and imaginary
magnitude $\sqrt{q_e}/R_e$.  Hence a zero requires $q_e=0$ and $B<0$.
Ellipticity implies $q_e=0$ exactly when every $\sin x_\mu=0$, so $x$ is a
Brillouin corner with $n$ coordinates equal to $\pi$.  There
$B=2n-1$, which is negative only for $n=0$.  Thus only $x=0$ is a zero.
Finally $V:=k_e/R_e$ is unitary and
$V+V^*=2B/R_e$, whence
$(I+V)^*(I+V)=2(1+B/R_e)I$.
\end{proof}

\begin{corollary}[A nonempty cutoff-uniform protected Wilson chart]
\label{cor:v2-protected-Wilson-chart}
Let
\[
 g_0:=\sqrt{\min\{m_g,1\}}.
\]
On the reference constant-coframe branch the V2 kernel satisfies
$|H_{W,a}|\succeq g_0I$.  If a gauge/spin/coframe perturbation obeys
\begin{equation}
 \Norm{H_{W,a}(\zeta)-H_{W,a}(0)}_{2\to2}\le {g_0\over2},
 \label{eq:v2-protected-chart-radius}
\end{equation}
then
\begin{equation}
 |H_{W,a}(\zeta)|\succeq {g_0\over2}I.
 \label{eq:v2-protected-chart-gap}
\end{equation}
The chart may be chosen uniformly in volume because the perturbation is
finite range and its operator norm is controlled by the corresponding local
row/column stocks.  Thus the protected Wilson chart required by the overlap
functional calculus is nonempty for the explicit V2 completion.
\end{corollary}
\begin{proof}
The reference gap is Lemma~\ref{lem:v2-trig-gap}.  The singular-value
perturbation inequality gives
$s_{\min}(H(\zeta))\ge s_{\min}(H(0))-\|H(\zeta)-H(0)\|$, proving
\eqref{eq:v2-protected-chart-gap}.  For a fixed-range kernel the
$2\to2$ norm is bounded by the geometric mean of its row and column Schur
bounds, which are local and volume independent.  Smooth dependence of the
one-edge transports and coframe coefficients then supplies a nonempty common
neighbourhood satisfying \eqref{eq:v2-protected-chart-radius}.
\end{proof}

\begin{firewall}[The physical zero is not the hard-shell pivot]
The exact overlap branch has its intended physical zero at $p=0$.  This does
not obstruct E1, whose inverse is always compressed to a hard annulus bounded
away from zero momentum.  Whenever a nowhere-zero \emph{full} determinant-line
section is needed, V2 retains the positive infrared/reference prescription
allowed by the programme (for example a fixed infrared spectral reference,
a compatible boundary-condition gap, or a protected Yukawa/Majorana floor).
No removal of that infrared reference is claimed here.
\end{firewall}

\subsection{Physical annulus and full Brillouin zone}

\begin{theorem}[Physical-branch expansion with all momentum derivatives]
\label{thm:v2-physical-expansion}
Fix a compact ellipticity set \eqref{eq:v2-coframe-elliptic} and an integer
$m\ge0$.  There exists $0<\delta<1$ and constants $C_\alpha,c_{\rm ph}>0$,
uniform on that coframe set, such that for $|ap|\le\delta$,
\begin{align}
 \Norm{\partial_p^\alpha
 (d_{{\rm ov},a}(p;e)-i\Gamma(e)\!\cdot p)}
 &\le C_\alpha a|p|^{2-|\alpha|},
 \qquad |\alpha|\le m,
 \label{eq:v2-overlap-expansion}\\
 \Norm{\partial_p^\alpha
 (\widehat P_{-,a}(p;e)-P_-)}
 &\le C_\alpha a^{|\alpha|}
 (a|p|)^{\max\{1-|\alpha|,0\}},
 \label{eq:v2-projector-expansion}\\
 s_{\min}(d_{{\rm ov},a}(p;e))
 &\ge c_{\rm ph}|p|.
 \label{eq:v2-physical-pivot}
\end{align}
Here $P_-=(I-\gamma_5)/2$ and
$\widehat P_{-,a}=(I+\varepsilon_e)/2$.  The power on the right side of
\eqref{eq:v2-overlap-expansion} is interpreted in the annular sense when
$|\alpha|>2$; because $a|p|\le\delta<1$, it is then a weaker bound than the
uniform analytic derivative bound and is therefore valid.
\end{theorem}
\begin{proof}
Define the dimensionless analytic matrix
\[
 F_e(x):=I+{i\Gamma\!\cdot s(x)+B(x)I\over R_e(x)}.
\]
The exact gap \eqref{eq:v2-Wilson-gap-exact} makes $F_e$ smooth, with all
fixed derivatives uniformly bounded on a common neighbourhood of $x=0$ and on
the compact coframe set.  Direct evaluation gives
$F_e(0)=0$ and
$D F_e(0)[x]=i\Gamma\!\cdot x$.  Hence
$G_e(x):=F_e(x)-i\Gamma\!\cdot x$ vanishes to order two.  Taylor's theorem
therefore gives
\[
 \Norm{\partial_x^\alpha G_e(x)}
 \le C_\alpha |x|^{\max\{2-|\alpha|,0\}}.
\]
For $|\alpha|>2$, the desired factor $|x|^{2-|\alpha|}$ is at least one on
$|x|\le1$, so the same statement implies the annular form.  Since
$d_{{\rm ov},a}=a^{-1}F_e(ap)$, scaling proves
\eqref{eq:v2-overlap-expansion}.  The projector statement follows identically
from the analytic map
$\widehat P_{-,a}(x)=(I+\gamma_5k_e(x)/R_e(x))/2$, whose value at zero is
$P_-$ and whose first nonzero displacement is linear in $x$.

For the pivot, shrink $\delta$ so that $B(x)<0$.  Then
\[
 1+{B\over R_e}={R_e+B\over R_e}
 ={q_e\over R_e(R_e-B)}.
\]
On the compact small ball the denominator is bounded above, while
$q_e\ge m_g|s(x)|^2$ and
$|s(x)|\ge c_\delta|x|$.  Equation
\eqref{eq:v2-overlap-singular-exact} therefore gives
$s_{\min}(d_{{\rm ov},a})\ge c_{\rm ph}|x|/a$.
\end{proof}

\begin{theorem}[Full-zone no-doubler pivot]
\label{thm:v2-full-zone-pivot}
For every $0<\delta<\pi$, there exists $c_{\rm UV}(\delta)>0$, uniform on the
compact coframe ellipticity set, such that
\begin{equation}
 \boxed{
 |x|\ge\delta
 \quad\Longrightarrow\quad
 s_{\min}(d_{{\rm ov},a}(x/a;e))
 \ge {c_{\rm UV}(\delta)\over a}.}
 \label{eq:v2-full-zone-pivot}
\end{equation}
Every fixed $x$-derivative of the corresponding inverse is bounded by $C a$ on
that compact full-zone region.  Consequently a smooth full-zone shell multiplier
has a lattice-scale quasilocal inverse kernel with weighted row/column norm
$O(a)$, uniformly in volume.
\end{theorem}
\begin{proof}
By Theorem~\ref{thm:v2-exact-overlap-branch}, the continuous singular-value
function
$\sqrt{2(1+B/R_e)}$ has no zero on the compact set
$\{(e,x): e\text{ in the ellipticity set},\ |x|\ge\delta\}$.
Its minimum is therefore positive, giving \eqref{eq:v2-full-zone-pivot}.
The inverse has the form $aG_e(x)$ with $G_e$ smooth on the same compact set,
so every fixed $x$-derivative is $O(a)$.  Multiplication by a smooth periodic
full-zone cutoff and repeated Fourier-series summation by parts give faster-than-
any-fixed-power decay in lattice distance.  The lattice row and column sums are
therefore $O(a)$ with no volume factor.
\end{proof}

\subsection{The moving chiral frame and the Yukawa block}

The exact overlap relation is particularly useful with the projector convention of
the Grand Tensor shell.

\begin{lemma}[Exact overlap chiral intertwining]
\label{lem:v2-chiral-intertwining}
For
$D_a=d_{{\rm ov},a}$ and
$\widehat P_-=(I+\varepsilon_e)/2$,
\begin{equation}
 \boxed{D_a\widehat P_-=P_+D_a,\qquad P_+=(I+\gamma_5)/2.}
 \label{eq:v2-chiral-intertwining}
\end{equation}
Hence $D_a$ maps $\Ran\widehat P_-$ into $\Ran P_+$ exactly.
\end{lemma}
\begin{proof}
Expand both products using
$D_a=a^{-1}(I+\gamma_5\varepsilon_e)$ and
$\varepsilon_e^2=I$:
\[
 D_a\widehat P_-
 ={1\over2a}(I+\varepsilon_e+\gamma_5+\gamma_5\varepsilon_e)
 =P_+D_a.
\]
\end{proof}

Let $U_e(x)$ be the Kato unitary obtained by transporting $P_-$ to
$\widehat P_-(x)$ along the radial path $t\mapsto tx$, with $U_e(0)=I$.
The V1 local-frame theorem applies; here the exact symbol permits a sharper scaled
form.

\begin{lemma}[Scaled Kato-frame estimate]
\label{lem:v2-Kato-scaled}
After shrinking $\delta$ if needed,
\begin{equation}
 U_e(x)P_-U_e(x)^*=\widehat P_-(x),
 \qquad
 \Norm{U_e(x)-I}\le C|x|,
 \label{eq:v2-Kato-first}
\end{equation}
and every fixed derivative satisfies
\begin{equation}
 \Norm{\partial_x^\alpha(U_e(x)-I)}
 \le C_\alpha |x|^{\max\{1-|\alpha|,0\}}.
 \label{eq:v2-Kato-derivative}
\end{equation}
The constants are uniform on the compact coframe set.
\end{lemma}
\begin{proof}
Along the radial path the Kato equation is
$\dot U=[\dot P,P]U$.  By
\eqref{eq:v2-projector-expansion},
$\dot P(tx)=x\cdot\nabla P(tx)=O(|x|)$ uniformly for $0\le t\le1$.
Gronwall's inequality gives $U-I=O(|x|)$.  Differentiating the finite-dimensional
ODE a fixed number of times gives uniform derivative bounds; the vanishing at zero
for the undifferentiated difference yields the displayed scaled form.
\end{proof}

Let $Y(H)$ denote the fixed-range Yukawa map on the reference Higgs chart, with
\begin{equation}
 M_Y:=\sup\Norm{Y(H)}<\infty
 \label{eq:v2-Yukawa-stock}
\end{equation}
on the finite panel.  In the Kato source frame define
\begin{align}
 \widetilde d_{\chi,a}(p;e,H)
 &:=\bigl(D_a(p;e)+Y(H)\bigr)U_e(ap)P_-:
 \Ran P_-\longrightarrow\Ran P_+,
 \label{eq:v2-fixed-frame-chiral}\\
 d_{\chi,0}(p;e,H)
 &:=\bigl(i\Gamma(e)\!\cdot p+Y(H)\bigr)P_-.
 \label{eq:v2-ordinary-chiral}
\end{align}
The target projection $P_+$ is suppressed because the kinetic term lands there by
Lemma~\ref{lem:v2-chiral-intertwining} and the Yukawa map is defined with that
target.

\begin{theorem}[Native chiral annular symbol theorem on the reference branch]
\label{thm:v2-native-chiral-annular}
Fix $m\ge0$.  There are constants $c_0,c_D,C_\alpha>0$ such that on every hard
annulus $c_-\Lambda_j\le|p|\le c_+\Lambda_j$ satisfying
\begin{equation}
 a\Lambda_j\le c_0,
 \qquad
 \Lambda_j\ge {2M_Y\over c_{\rm ph}c_-},
 \label{eq:v2-hard-scale-conditions}
\end{equation}
one has, for $|\alpha|\le m$,
\begin{align}
 \Norm{\partial_p^\alpha
 (\widetilde d_{\chi,a}-d_{\chi,0})}
 &\le C_\alpha a\Lambda_j^{2-|\alpha|},
 \label{eq:v2-chiral-error}\\
 s_{\min}(\widetilde d_{\chi,a}(p))
 +s_{\min}(d_{\chi,0}(p))
 &\ge 2c_D\Lambda_j.
 \label{eq:v2-chiral-pivot}
\end{align}
Thus the undifferentiated part of Assumption~\ref{ass:v1-native-symbol-row} is a
theorem for the explicit V2 Wilson completion.  On the complementary full-zone
region, the kinetic chiral inverse has $O(a)$ lattice-scale quasilocal norm; adding
the bounded Yukawa block preserves that bound for sufficiently small $a$.
\end{theorem}
\begin{proof}
By Theorem~\ref{thm:v2-physical-expansion} and
Lemma~\ref{lem:v2-Kato-scaled},
\[
 D_aU_eP_--i\Gamma\!\cdot pP_-
 =(D_a-i\Gamma\!\cdot p)U_eP_-
 +i\Gamma\!\cdot p(U_e-I)P_-
 =O(a|p|^2).
\]
The Yukawa-frame correction is
$Y(H)(U_e-I)P_-=O(M_Ya|p|)$, which is
$O(a\Lambda_j^2)$ under \eqref{eq:v2-hard-scale-conditions}.  Differentiating and
using the scaled derivative estimates proves \eqref{eq:v2-chiral-error}.

For the pivot, \eqref{eq:v2-overlap-singular-exact} shows that $D_a^*D_a$ is a
scalar multiple of the identity.  Since $U_e$ is unitary and maps $P_-$ onto
$\widehat P_-$, the kinetic chiral restriction has exactly the same nonzero
singular value as $D_a$.  Hence
\[
 s_{\min}(D_aU_eP_-+Y U_eP_-)
 \ge c_{\rm ph}c_-\Lambda_j-M_Y
 \ge {1\over2}c_{\rm ph}c_-\Lambda_j.
\]
The same elementary singular-value perturbation bound applies to
$i\Gamma\cdot p+Y$, whose kinetic singular value is at least
$\sqrt{m_g}\,c_-\Lambda_j$.  Decreasing $c_D$ if needed gives
\eqref{eq:v2-chiral-pivot}.  The full-zone statement follows from
Theorem~\ref{thm:v2-full-zone-pivot} and the bounded perturbation $Y$.
\end{proof}

\subsection{Finite external source jets}
\label{subsec:v2-source-jets}

We now treat gauge, spin, coframe and Higgs insertions as the fixed perturbative
jets required by the coefficientwise theorem.  Let the external source panel contain
only finitely many Fourier momenta $q$ with
\begin{equation}
 |q|\le\mu,
 \qquad
 \mu\le\kappa\Lambda_j
 \label{eq:v2-soft-source-window}
\end{equation}
for one fixed $0<\kappa<1$ on the hard shell under consideration.

For an internal gauge plus spin connection $\mathcal A_\mu$, use the actual edge
parallel transport
\begin{equation}
 \mathcal U_{x,\mu}(\zeta)
 =\mathcal P\exp\left(
 \int_0^a\mathcal A_\mu(x+t\hat\mu;\zeta)\,dt\right).
 \label{eq:v2-link-parallel-transport}
\end{equation}
Coframe sources enter through the sampled $\Gamma_\mu(e)$ in
\eqref{eq:v2-native-Wilson}; Higgs/Yukawa sources enter the onsite map $Y(H)$.
The primitive Wilson kernel has no antifield dependence in this regulator choice;
antifields enter through the local BV interaction/source vertices and blocking maps.
Likewise determinant-line coordinates act through the Kato/line frame, not by
changing the primitive Wilson stencil.

\begin{lemma}[Parallel-transport source jets]
\label{lem:v2-link-jets}
At the zero-connection reference point, every fixed $r$th source derivative of
\eqref{eq:v2-link-parallel-transport} is a finite sum of ordered simplex integrals
of total size $O(a^r)$.  If the inserted source modes satisfy
\eqref{eq:v2-soft-source-window}, then every fixed derivative with respect to their
momenta gains an additional factor $O(a)$.  In particular the exact edge jet is its
continuum midpoint/local jet plus a remainder smaller by at least one factor
$a(\Lambda_j+\mu)$ after it is inserted into a hard Wilson vertex.
\end{lemma}
\begin{proof}
Differentiate the path-ordered exponential at zero connection.  The $r$th
coefficient is a sum over orderings of integrals on
$0<t_1<\cdots<t_r<a$.  Their total simplex volume is $a^r/r!$.  A plane-wave
source contributes factors $e^{iq\cdot(x+t_k\hat\mu)}$; differentiating in $q$
produces powers of $t_k$, each bounded by $a$.  Taylor expansion of these phase
factors and of the hard exponentials $e^{\pm iap_\mu}$ then gives one further
factor $a(|p|+\sum|q_k|)$ beyond the leading local jet.
\end{proof}

For a source multi-index $I$, let $v_{I,a}(p;\mathbf q)$ denote the exact momentum
vertex obtained by differentiating the fixed-frame native chiral operator at the
reference point, including the Kato-frame derivative when the source space moves.
Let $v_{I,0}(p;\mathbf q)$ be the corresponding continuum covariant
Dirac--Yukawa Taylor jet.  Its hard-momentum degree is denoted
$\nu_I\in\{0,1\}$: coframe-type kinetic jets may have degree one, whereas
connection/Yukawa and the resulting local contact jets have degree zero.  If the
ordinary vertex vanishes, $\nu_I=0$ and $v_{I,0}=0$ by convention.

\begin{theorem}[All fixed native source vertices match their ordinary jets]
\label{thm:v2-source-vertex-transfer}
Fix source order $r$, momentum derivative order $m$, and the finite panel
\eqref{eq:v2-soft-source-window}.  For every $0<|I|\le r$ there are constants
$C_{I,\alpha,\beta}$, independent of volume, shell number and cutoff, such that
on $a\Lambda_j\le c_0$,
\begin{equation}
 \boxed{
 \Norm{
 \partial_p^\alpha\partial_{\mathbf q}^\beta
 (v_{I,a}-v_{I,0})}
 \le
 C_{I,\alpha,\beta}
 a\Lambda_j^{\nu_I+1-|\alpha|-|\beta|}}
 \label{eq:v2-source-vertex-error}
\end{equation}
for all fixed $|\alpha|+|\beta|\le m$, after division by the declared source
amplitudes/stocks.  The native and ordinary vertices themselves obey the same
bound without the prefactor $a\Lambda_j$ relative to their engineering scale.

The first nonzero ordinary jets are the expected ones:
\begin{enumerate}[label=\textup{(J\arabic*)},leftmargin=3em]
\item the zero-source kinetic jet is $i\Gamma^\mu p_\mu$;
\item a connection insertion is the local Clifford contraction
$\Gamma^\mu\mathcal A_\mu$ in the chosen anti-/Hermitian convention;
\item a coframe insertion differentiates $i\Gamma^\mu(e)p_\mu$;
\item a Higgs/Yukawa insertion is the corresponding derivative of $Y(H)$;
\item primitive antifield derivatives of $H_W$ vanish for
\eqref{eq:v2-native-Wilson}; their BV vertices remain separate local source
insertions;
\item determinant-line derivatives are derivatives of the Kato/line frame and are
controlled by the same analytic source calculus, not by an invented Wilson
coupling.
\end{enumerate}
Thus the differentiated source part of Assumption~\ref{ass:v1-native-symbol-row}
is proved at every fixed source order for the V2 regulator completion.
\end{theorem}
\begin{proof}
At the reference point the Wilson symbol has the uniform all-zone gap of
Lemma~\ref{lem:v2-trig-gap}.  Write the sign as the difference of its positive and
negative Riesz projectors.  A fixed source derivative of either projector is a finite
sum of contour-resolvent words of the form
\begin{equation}
 {1\over2\pi i}\oint_\Gamma
 R_{p+Q_k}(z)V_{B_k,a}\,
 R_{p+Q_{k-1}}(z)\cdots
 V_{B_1,a}R_p(z)\,dz,
 \label{eq:v2-resolvent-source-word}
\end{equation}
where $R_p(z)=(z-h_e(ap))^{-1}$, the $B_i$ form an ordered partition of $I$,
and $Q_i$ are partial sums of external source momenta.  The contour may be chosen
uniformly because the negative and positive Wilson spectra remain separated by the
fixed gap.  Every resolvent and every fixed momentum derivative is therefore a
uniform analytic function of the dimensionless variables
$ap,a q_1,\ldots,a q_r$ on the physical polydisk.

The primitive source vertices $V_{B,a}$ are obtained by differentiating the
one-edge formula \eqref{eq:v2-native-Wilson}.  Lemma~\ref{lem:v2-link-jets}
gives the connection scaling and its local leading jet.  Coframe derivatives act on
$\Gamma_\mu$ and are always accompanied, in the kinetic part, by the centred hard
shift, whose physical branch is $iap_\mu+O((ap)^2)$.  Higgs/Yukawa derivatives are
onsingular onsite jets.  Hence every word \eqref{eq:v2-resolvent-source-word} is a
smooth finite matrix function of the dimensionless hard and soft momenta.

At the origin of these variables its first nonzero Taylor polynomial is obtained by
replacing the centred edge difference by the ordinary covariant first derivative,
the edge parallel transport by its local connection jet, and the onsite Yukawa map
by itself.  Direct differentiation of \eqref{eq:v2-native-Wilson} at the physical
branch gives the six listed ordinary jets.  The Kato frame is an analytic functional
of the same projector on the protected local chart, so its derivatives satisfy the
same scaled Taylor calculus.  Subtracting the first nonzero ordinary Taylor jet
therefore leaves one extra factor
$a(|p|+\sum|q_i|)$.  Under
\eqref{eq:v2-soft-source-window} this is $O(a\Lambda_j)$.
Rescaling the remaining hard degree $\nu_I$ and differentiating in $p$ or $q$
proves \eqref{eq:v2-source-vertex-error}.  Fixed order leaves only finitely many
resolvent words, so no source-radius or volume factor enters.
\end{proof}

\begin{remark}[Why the coefficientwise route is stronger than a crude background norm]
A finite coframe displacement changes the Wilson operator by $O(1)$ in its global
dimensionless operator norm and would therefore produce an unhelpful $O(a^{-1})$
bound if one first perturbed the full overlap operator in operator norm.  The source
jet in Theorem~\ref{thm:v2-source-vertex-transfer} instead carries the hard
centred-difference factor $ap$, and hence has the correct $O(\Lambda_j)$ stress
scaling.  This is exactly the quantifier gain supplied by the coefficientwise
programme.
\end{remark}

\subsection{Full-zone source derivatives}

\begin{proposition}[Uniform ultraviolet-cap source calculus]
\label{prop:v2-full-zone-source}
Fix $\delta>0$ and a finite source order.  On the full-zone region
$|ap|\ge\delta$, every fixed reference inverse/source word generated by
Theorem~\ref{thm:v2-reference-pivot} has a smooth periodic dimensionless symbol.
After the physical factor of $a$ from each reference fermion inverse is restored,
its cutoff-profiled Fourier kernel obeys a volume-uniform lattice-scale weighted
row/column bound.  A derivative with respect to an external source momentum gains
one factor $a$, so the complementary source remainder on a panel $|q|\le\mu$ is
$O(a\mu)$.
\end{proposition}
\begin{proof}
Theorem~\ref{thm:v2-full-zone-pivot} gives an inverse $aG(x)$ with $G$ smooth on
the compact high-zone region.  Fixed source derivatives are finite resolvent/inverse
words with smooth periodic dimensionless factors by
Theorem~\ref{thm:v2-source-vertex-transfer}; hence every fixed $x$-derivative is
bounded.  Fourier-series summation by parts gives uniform rapid decay in lattice
units.  External source momentum enters only through the dimensionless combination
$aq$, so the mean-value theorem gives the additional factor $a\mu$.
\end{proof}

\subsection{The missing line summability follows from the same shell calculus}
\label{subsec:v2-line-summability}

V1 proved the exact Bell-polynomial composition theorem for determinant/Pfaffian
line sections under the sole remaining analytic input
$\rho_j=O(\mu/\Lambda_j)$.  We now prove that input for the coefficientwise V2
branch.

Let $g_{j,I}(\mathbf q)$ be any fixed connected logarithmic determinant-line
coefficient at shell $j$, after the exact source differentiation has been performed.
By Lemma~\ref{lem:v1-logdet-words} it is a finite sum of cyclic hard inverse--vertex
words.  Let $\mathsf T_j$ denote the common local Taylor/EFT assignment already
used by the bosonic shell; in particular it contains at least the zero-external-
momentum local jet of every superficially local line coefficient.

\begin{theorem}[One-excess-scale chiral line remainder]
\label{thm:v2-line-excess}
Fix the source order, the finite external panel $|q|\le\mu$, and the local/EFT weight
kept in the two-loop packet.  For every physical hard annulus and every connected
logarithmic determinant-line coefficient,
\begin{equation}
 \boxed{
 \Norm{\partial_I(g_j-\mathsf T_jg_j)}_{\rm pan}
 \le C_I{\mu\over\Lambda_j}.}
 \label{eq:v2-line-excess}
\end{equation}
The same estimate holds for the optional locally oriented Pfaffian logarithm.
For the ultraviolet-cap shells, the right side is read as $C_Ia\mu$, which is the
same estimate because $\Lambda_j\asymp a^{-1}$ there.

All mixed source contacts produced by differentiation of the Schur complement or
of the shell product are retained before $\mathsf T_j$ is applied.
\end{theorem}
\begin{proof}
On a physical annulus rescale the hard loop momentum by $p=\Lambda_jk$.  The
reference inverse bounds of Theorem~\ref{thm:v2-native-chiral-annular}, the source
vertex bounds of Theorem~\ref{thm:v2-source-vertex-transfer}, and the rooted trace
estimate of Lemma~\ref{lem:v1-logdet-words} show that, after division by the declared
engineering/local source stock, every fixed cyclic integrand is a $C^1$ function of
$k$ and of the dimensionless external variables $q_i/\Lambda_j$, with a bound
independent of $j$, $a$ and volume.  The local assignment $\mathsf T_j$ removes at
least the value at zero external momentum.  Therefore the multivariate mean-value
theorem gives
\[
 \Norm{F_j(\mathbf q)-F_j(0)}
 \le C_I\sum_i{|q_i|\over\Lambda_j}
 \le C_I'{\mu\over\Lambda_j}.
\]
If power counting requires a higher local Taylor polynomial, the remainder only
improves.  Integration of the normalized hard variable and the rooted trace
contraction preserve the same factor.  The ultraviolet-cap version is
Proposition~\ref{prop:v2-full-zone-source}.  The Pfaffian logarithm differs by the
fixed factor $1/2$ on the chosen oriented component.

Crucially, differentiation is performed before the Taylor assignment.  Hence every
mixed hard/soft Schur contact of Proposition~\ref{prop:v1-Schur-contacts} and every
mixed-shell contact of Theorem~\ref{thm:v1-line-product} remains present in the
coefficient being subtracted; no contact is discarded by the estimate.
\end{proof}

\begin{corollary}[Differentiated line products are summable]
\label{cor:v2-line-product-closed}
For geometric shell scales $\Lambda_j=b^j\lambda$, the hypothesis
\eqref{eq:v1-line-log-summable} of Theorem~\ref{thm:v1-line-product} holds with
\begin{equation}
 \rho_j=C{\mu\over\Lambda_j}.
 \label{eq:v2-rho-line}
\end{equation}
Hence every fixed differentiated determinant/Pfaffian line product is uniformly
bounded through arbitrarily many ultraviolet shell eliminations on the fixed source
panel, with all mixed-shell source contacts retained.
\end{corollary}
\begin{proof}
Theorem~\ref{thm:v2-line-excess} gives the shell estimate and the geometric series
$\sum_j\mu/\Lambda_j$ converges.  Apply
Theorem~\ref{thm:v1-line-product} and Proposition~\ref{prop:v1-Pfaffian}.
\end{proof}

\subsection{Coefficientwise E1 closure}

\begin{theorem}[E1 closes for the explicit V2 Wilson completion]
\label{thm:v2-E1-closure}
Fix:
\begin{enumerate}[label=\textup{(F\arabic*)},leftmargin=3em]
\item the explicit Wilson regulator \eqref{eq:v2-native-Wilson};
\item a compact uniformly elliptic reference coframe chart;
\item the protected local determinant/Pfaffian line component inherited from V1;
\item a finite loop/source packet and finite external source momentum window;
\item a positive IR/reference scale, with hard shells taken above the bounded
Yukawa/reference mass scale.
\end{enumerate}
Then E1 is proved at every fixed source/derivative order in the coefficientwise
sense required by the programme:
\begin{enumerate}[label=\textup{(E1.\arabic*)},leftmargin=3.5em]
\item the native Wilson kernel is finite range, Hermitian and exactly gauge/spin
covariant, with all fixed local source jets;
\item its reference overlap branch has an exact all-zone Wilson gap, one and only one
physical zero, and no doubler zero;
\item on $a\Lambda_j\ll1$ the fixed-frame native chiral symbol differs from the
ordinary covariant Weyl/Yukawa symbol by $O(a\Lambda_j^2)$ with every required
fixed momentum/source derivative, and has an $O(\Lambda_j)$ reference hard pivot;
\item on the complementary ultraviolet-cap region the reference inverse and all
fixed source words have $O(a)$ lattice-scale quasilocal bounds;
\item the V1 overlap functional-calculus, Kato-frame, inverse-word, Schur-contact and
rooted graph theorems therefore apply coefficientwise without a
cutoff-independent finite-amplitude source ball;
\item the locally subtracted determinant/Pfaffian logarithmic line remainder is
$O(\mu/\Lambda_j)$ and its differentiated products are summable through arbitrary
finite shell composition.
\end{enumerate}
No analytic finite-coupling RG trajectory, global determinant-line holonomy,
zero-crossing theorem, or nonperturbative chiral measure is asserted.
\end{theorem}
\begin{proof}
Items (E1.1)--(E1.2) are Proposition~\ref{prop:v2-Wilson-structure} and
Theorem~\ref{thm:v2-exact-overlap-branch}.  Item (E1.3) is
Theorems~\ref{thm:v2-native-chiral-annular} and
\ref{thm:v2-source-vertex-transfer}.  Item (E1.4) is
Theorem~\ref{thm:v2-full-zone-pivot} and
Proposition~\ref{prop:v2-full-zone-source}.  The reference-pivot replacement of the
stronger finite-amplitude inversion hypothesis is
Theorem~\ref{thm:v2-reference-pivot}; the remaining quasilocal, Schur and rooted
conclusions are exactly the V1 theorems fed by these now-proved native symbol rows.
Item (E1.6) is Theorem~\ref{thm:v2-line-excess} and
Corollary~\ref{cor:v2-line-product-closed}.
\end{proof}

\begin{firewall}[Scope of the E1 closure]
Theorem~\ref{thm:v2-E1-closure} is an existence/closure theorem for the explicit
Wilson completion fixed in \eqref{eq:v2-native-Wilson}.  If the final Grand Tensor
manuscript chooses a different Wilson mass, hopping stencil, or graph regulator, one
must either adopt \eqref{eq:v2-native-Wilson} as the native definition or prove a
finite regulator-comparison theorem.  The coefficientwise conclusion itself does
not license silently changing regulators afterward.
\end{firewall}

\section{V2 handoff to E2: the native chiral BV breaking}
\label{sec:v2-E2-handoff}

With E1 closed on the explicit regulator branch, continuing to refine the hard
fermion symbol would now be downstream polishing.  The next load-bearing issue is
cohomological but must be tied to the \emph{same native shell}: the local chiral
breaking generated by the regulated determinant/Pfaffian line must be decomposed
into the standard anomaly class and a local BV-exact regulator term, with all
source/antifield descendants carried by one common finite action.

Let $S_{0,j}$ denote the common classical BV action at shell $j$ and let
$\mathfrak s_j=(S_{0,j},\,\cdot\,)$ be its linearized BV differential on the finite
local bank.  Let $\mathcal B_{j}^{(\ell,r)}$ denote the ghost-number-one local
breaking at loop order $\ell$ and source order $r$ obtained from the coefficient of
\begin{equation}
 {1\over2}(S,S)-\hbar\Delta S
 \label{eq:v2-master-breaking}
\end{equation}
after the complete native hard shell, including the chiral line/reference term, has
been assembled and its quasilocal remainder has been separated by the E1 Taylor
assignment.

\begin{theorem}[Exact next load-bearing theorem for V3]
\label{thm:v2-next-E2}
At the fixed two-loop/EFT/source order, prove for the \emph{native V2 shell} that
there is one local decomposition
\begin{equation}
 \boxed{
 \mathcal B_j^{(\ell,r)}
 =\sum_A c_{A,j}^{(\ell,r)}\,\mathcal A_A^{\rm std}
 +\mathfrak s_j\Xi_j^{(\ell,r)},}
 \label{eq:v2-E2-target-decomposition}
\end{equation}
where:
\begin{enumerate}[label=\textup{(B\arabic*)},leftmargin=3em]
\item $\{\mathcal A_A^{\rm std}\}$ is the finite ordinary four-dimensional local
internal-gauge anomaly basis in the already fixed hypercharge convention;
\item the coefficients $c_{A,j}^{(\ell,r)}$ are computed from the complete native
chiral shell and shown to equal the inherited Standard-Model representation traces,
so they vanish after the full three-generation packet is assembled;
\item $\Xi_j^{(\ell,r)}$ is one local restoring primitive in the same finite
EFT/source bank, not a graph-dependent family of counterterms;
\item differentiating \eqref{eq:v2-E2-target-decomposition} generates all gauge,
diffeomorphism, local-Lorentz, Higgs/Yukawa, ghost, antifield, stress/contact and
line descendants required by the common two-loop source packet;
\item the resulting restored shell enlarges the inherited V80 bosonic self-map to
one common two-loop Einstein--Standard Model action/source/line update.
\end{enumerate}
The proof must use the E1 native comparison to identify the cohomology coefficient;
anomaly-trace cancellation alone is not a regulator-matching theorem.
\end{theorem}

The correct V3 attack is therefore not another anomaly-trace calculation.  Those
traces are inherited.  It is the finite local comparison map from the native
regulated BV breaking to the ordinary anomaly representative, followed by a common
local primitive for the exact remainder.

\section{V2 status ledger}
\label{sec:v2-status}

\subsection*{Proved in V2}
\begin{enumerate}[leftmargin=2.2em]
\item At fixed source order, a coefficientwise fermionic shell requires only the
reference hard inverse and finitely many source jets.  A cutoff-uniform
finite-amplitude source ball is a stronger analytic statement and is not required
for the programme's coefficientwise endpoint.
\item The previously abstract Wilson slot was filled by the explicit covariant
nearest-neighbour kernel \eqref{eq:v2-native-Wilson}.  It is finite range,
Hermitian and exactly gauge/spin covariant, with fixed-range source derivatives.
\item On every constant uniformly elliptic coframe, the Wilson symbol obeys the exact
all-zone gap identity \eqref{eq:v2-Wilson-gap-exact}.  The associated overlap
operator has exactly one physical zero and no doubler zeros.
\item The physical overlap branch satisfies
$d_{{\rm ov},a}=i\Gamma\cdot p+O(a|p|^2)$ with every fixed momentum derivative,
and has an $O(|p|)$ physical singular-value floor.  The complementary full
Brillouin zone has an $O(a^{-1})$ pivot and an $O(a)$ inverse kernel.
\item In the Kato chiral frame, the overlap--Yukawa block satisfies the V1 annular
finite-to-ordinary symbol row and hard reference pivot on sufficiently hard shells.
\item Every fixed gauge, spin, coframe and Higgs/Yukawa source jet of the native
overlap operator matches the corresponding ordinary covariant Dirac/Yukawa jet with
one extra factor $a\Lambda_j$ after normalized source scaling.  Primitive antifield
and line sources are correctly kept outside the Wilson stencil and enter through
the local BV/Kato source calculus.
\item The locally subtracted logarithmic determinant/Pfaffian line remainder is
$O(\mu/\Lambda_j)$, including the ultraviolet-cap interpretation
$O(a\mu)$.  Therefore the V1 differentiated line-product theorem is summable with
all mixed-shell contacts retained.
\item E1 is closed at every fixed source/derivative order for the explicit V2 Wilson
completion, at the coefficientwise scope of the programme.
\end{enumerate}

\subsection*{Inherited}
\begin{enumerate}[leftmargin=2.2em]
\item The V80 bosonic Einstein--gauge--Higgs two-loop self-map and its common
local/source normalization.
\item The three-generation Standard-Model chiral representation, exact cancellation
of the five ordinary local gauge-anomaly coefficients in the declared convention,
and even $SU(2)$ doublet parity.
\item Exact determinant/Pfaffian Schur factorization, finite regrouping cocycles,
and differentiated line-source identities.
\item The V1 weighted functional calculus, local Kato line frame, Schur-contact
calculus, rooted hard-graph bounds, and Bell-polynomial line-product identity.
\end{enumerate}

\subsection*{Conditional}
\begin{enumerate}[leftmargin=2.2em]
\item The E1 closure is native to the explicit Wilson completion
\eqref{eq:v2-native-Wilson}.  Identification with any other Wilson kernel intended
by a separate manuscript requires an explicit finite regulator-comparison theorem.
\item The construction remains on the protected perturbative determinant/Pfaffian
component and retains the positive IR/reference scale declared by the programme.
No statement is made across zero modes or global line-holonomy transitions.
\item Constants may depend on the fixed perturbative, EFT, source and derivative
orders and on the compact reference coframe/source panel.  No uniform
large-loop-order estimate is claimed.
\end{enumerate}

\subsection*{Open}
\begin{enumerate}[leftmargin=2.2em]
\item Prove the native two-loop chiral BV-breaking decomposition
\eqref{eq:v2-E2-target-decomposition} using the E1 finite-to-ordinary comparison,
not merely the already known anomaly traces.
\item Show that the nontrivial anomaly coefficients of that native breaking equal the
inherited Standard-Model trace vector and hence vanish only after the complete
chiral representation is assembled.
\item Construct one common local restoring primitive and all of its source,
antifield, stress/contact and determinant/Pfaffian-line descendants.
\item Combine that restored chiral row with the inherited V80 bosonic map to prove
one common two-loop Einstein--Standard Model action/source/line self-map.
\end{enumerate}

\subsection*{Exact next load-bearing theorem}
\begin{center}
\fbox{\parbox{0.92\textwidth}{
\textbf{Native chiral BV-breaking comparison and restoration.}
For the explicit E1-closed Wilson/overlap regulator, compute the complete local
ghost-number-one breaking of the two-loop hard shell in the common BV/source
scheme, prove that its non-exact projection is exactly the ordinary Standard-Model
anomaly trace vector, use the inherited trace cancellations to remove that
projection, and construct one local BV primitive with all source/antifield/line
descendants.  The output must be one common chiral extension of the V80 bosonic
two-loop shell self-map, not separate graphwise counterterms.
}}
\end{center}

\section*{Append-only continuation marker: V2}
\addcontentsline{toc}{section}{Append-only continuation marker: V2}
\textbf{End of V2.}  The complete V1 body above is retained.  Future versions must
append new work after this marker without replacing the V1--V2 mathematical body.
The next file is to be named
\texttt{Renewal\_Perturbative\_ESM\_Continuum\_Research\_Notes\_V3.tex}.



\clearpage
\part*{V3: Native anomaly-class cancellation and the common chiral two-loop restorer}
\addcontentsline{toc}{part}{V3: Native anomaly-class cancellation and the common chiral two-loop restorer}

\section{V3 upstream correction: the load-bearing target is zero native anomaly projection}
\label{sec:v3-upstream}

V2 ended with a deliberately strong version of E2: compare the complete native
chiral breaking with a conventionally normalized ordinary anomaly representative,
identify every coefficient with the corresponding Standard-Model representation
trace, and only then construct a restoring primitive.  That route is legitimate,
but for the actual coefficientwise E2 theorem it contains one unnecessary
normalization problem.

What the common local BV range theorem consumes is not a separate numerical
normalization of every anomalous species.  It consumes the statement
\begin{equation}
 \boxed{\mathsf A_{\rm std}\,\mathcal B^{(\ell,r)}_j=0}
 \label{eq:v3-load-bearing-zero}
\end{equation}
for the \emph{complete native Standard-Model packet}, where
$\mathsf A_{\rm std}$ is the ordinary finite anomaly-quotient projection already
constructed in V80.  Once \eqref{eq:v3-load-bearing-zero} is known, the inherited
matter-inclusive local cohomology theorem gives one local BV primitive.  The
coefficientwise continuum target does not require a separate theorem that the
uncancelled one-species anomaly has a preassigned continuum normalization.

There is a native way to prove \eqref{eq:v3-load-bearing-zero}.  The supplied
Grand Tensor field-shell construction already gives, on the protected nonzero
component, a coordinate-free chiral determinant section for the complete
anomaly-free generation and proves that this section descends under the finite
gauge group.  V1--V2 add precisely the missing analytic information: a legitimate
local line frame, fixed-order source differentiability, hard-shell quasilocality,
and the exact determinant/Pfaffian shell cocycle.  We therefore first build a
\emph{gauge-invariant local line frame for the complete packet itself}.  This
kills the nontrivial native internal-gauge line cocycle before any continuum
normalization is invoked.

\begin{firewall}[What V3 does and does not replace]
The argument below does not claim that an anomalous single Weyl representation
has zero anomaly, and it does not identify the universal numerical coefficient
of such an anomalous representation.  It proves the stronger statement needed
for E2: after the actual Standard-Model representation is assembled, the
complete native finite shell has zero projection onto the standard local anomaly
quotient.  The inherited trace cancellations and determinant-line descent are
used as inputs; they are not reproved.  No Adler--Bardeen theorem is imported.
\end{firewall}

\section{The source-extended finite ESM BV complex}
\label{sec:v3-complex}

Fix once and for all the two-loop perturbative order, local EFT weight six,
source order $r$, field/source arity $N$, the finite external panel, and the
explicit V2 Wilson/overlap branch.  Let $\Phi$ denote the complete retained
field list: coframe/metric and spin connection variables, the three internal
gauge fields, Higgs field, chiral quark and lepton multiplets, ghosts,
nonminimal variables, and all admitted composite/source writers.  Let $\Phi^*$
denote the exact antifield duals used by the common BV splitting.

Let $\mathfrak J_r$ be the finite truncated source-jet algebra generated by all
source coordinates retained through total source order $r$.  Source variables
which are members of BRST doublets or antifield/source multiplets carry their
actual differential; spectator physical sources carry zero differential.  Let
\begin{equation}
 \widehat{\mathfrak L}^{\rm ESM}_{N,6,r}
 :=\mathfrak L^{\rm ESM}_{N,6}\otimes \mathfrak J_r
 \label{eq:v3-source-extended-bank}
\end{equation}
be the complete finite local bank.  The notation means the quotient by the
chosen integration-by-parts, algebraic-symmetry, equation-of-motion and other
redundant directions, exactly as in V80; it does not mean that only a displayed
list of Warsaw-type representatives is retained.

Let $\mathsf P^{\rm ESM}_{N,6,r}$ denote the common local Taylor/EFT/source
projection.  It is applied only after every contact generated by source
differentiation has been assembled.  Let
\begin{equation}
 d_j:=\mathsf P^{\rm ESM}_{N,6,r}
       (S_{0,j},\,\cdot\,)
       \mathsf P^{\rm ESM}_{N,6,r}
 \label{eq:v3-local-differential}
\end{equation}
be the projected classical BV differential at shell $j$.

\begin{theorem}[Finite source-extended ESM chain complex]
\label{thm:v3-chain-complex}
On the protected perturbative chart and at the fixed $(N,6,r)$,
\begin{equation}
 \boxed{d_j^2=0,\qquad
 \mathsf P^{\rm ESM}_{N,6,r}d_j
 =d_j\mathsf P^{\rm ESM}_{N,6,r}.}
 \label{eq:v3-chain-identities}
\end{equation}
The nonminimal pairs remain contractible.  Adding the chiral fermion and
Yukawa rows does not enlarge the nontrivial ghost-number-one quotient beyond
the ordinary standard gauge-anomaly space already isolated in V80.
\end{theorem}

\begin{proof}
Before projection the complete classical ESM BV action satisfies the classical
master equation on the declared relative/protected domain.  Hence its
Hamiltonian vector field squares to zero.  The chiral fermion transformations
are linear in the corresponding gauge generators and the Yukawa terms are
gauge covariant; therefore the classical differential preserves canonical
weight, ghost number, representation type and the finite source order.  The
same paired-jet argument used in the V80 Einstein--gauge--Higgs bank then shows
that the complete Taylor/EFT/source projector commutes with the differential
provided every descendant at the retained weight and source order is kept.
This gives \eqref{eq:v3-chain-identities}.

The second statement is the matter-inclusive content of the cohomological
reduction already used in V80.  Hypercharge is not a free abelian generator on
the charged-Higgs branch, so the exceptional free-abelian antifield families
are absent.  In four dimensions the remaining nontrivial ghost-one quotient is
the ordinary gauge-anomaly quotient.  Adding further gauge-covariant matter
variables and their antifields enlarges the contractible and invariant parts of
the finite bank but does not create a new four-dimensional ghost-one anomaly
family outside that inherited classification.  The nonminimal doublets are
contractible by the same explicit homotopy as in V80.
\end{proof}

Denote the resulting anomaly quotient by
\begin{equation}
 \mathfrak A_{\rm std}
 =\Span\{\mathcal A_{333},\mathcal A_{Y33},\mathcal A_{Y22},
              \mathcal A_{YYY},\mathcal A_{YRR}\},
 \label{eq:v3-anomaly-space}
\end{equation}
where the labels refer to the consistent descendants of the five invariant
polynomials already fixed in V80.  Let
$\mathsf A_{\rm std}:\ker(d_j)|_{\operatorname{gh}=1}\to
\mathfrak A_{\rm std}$ be the inherited coefficient projection.

\section{Quantum-master consistency is a native finite identity}
\label{sec:v3-consistency}

The next step is upstream of anomaly classification.  A local anomaly
projection is meaningful only after the complete shell breaking has been shown
to be a cocycle.

Let $\Delta_j$ be the finite BV Laplacian in the same field--antifield frame
used by the shell integration, including the declared measure/reference
normalization exactly once.  For an even BV functional $S$ define
\begin{equation}
 \mathcal M_j(S)
 :=\frac12(S,S)-\hbar\Delta_jS.
 \label{eq:v3-master-functional}
\end{equation}
All reference terms which depend on retained sources are part of the same
finite functional and are not discarded as ``field independent.''

\begin{lemma}[Finite quantum consistency identity]
\label{lem:v3-quantum-consistency}
Assume $\Delta_j^2=0$ and the standard BV second-order identity for the
antibracket.  Then, identically at finite cutoff,
\begin{equation}
 \boxed{(S,\mathcal M_j(S))-\hbar\Delta_j\mathcal M_j(S)=0.}
 \label{eq:v3-quantum-WZ}
\end{equation}
\end{lemma}

\begin{proof}
The graded Jacobi identity gives
$(S,\frac12(S,S))=0$.  Since $S$ is even, the BV second-order identity gives
\[
 \Delta_j\frac12(S,S)=-(S,\Delta_jS).
\]
Therefore the contribution of $(S,-\hbar\Delta_jS)$ cancels
$-\hbar\Delta_j\frac12(S,S)$, while
$\Delta_j^2S=0$.  No limit and no formal change of regulator is used.
\end{proof}

Write the recursively normalized action through two loops as
\begin{equation}
 S_j^{[\le2]}=S_{0,j}+\hbar S_{1,j}+\hbar^2S_{2,j}.
 \label{eq:v3-action-expansion}
\end{equation}
At each loop order the symbol $S_{\ell,j}$ means the \emph{single common local
coefficient action} already containing every lower-order restoring term and
all of its source/antifield descendants.

\begin{theorem}[Cocycle property of the first un-restored shell coefficient]
\label{thm:v3-breaking-closed}
Suppose the common master equation has been restored through loop order
$\ell-1$.  Let $\widetilde{\mathcal B}^{(\ell,\le r)}_j$ be the complete
coefficient of $\hbar^\ell$ in the native shell master defect before the new
order-$\ell$ local counterterm is chosen, with the bosonic, chiral,
ghost/nonminimal, antifield, line, measure/reference, stress and contact rows
assembled.  Then
\begin{equation}
 (S_{0,j},\widetilde{\mathcal B}^{(\ell,\le r)}_j)=0.
 \label{eq:v3-unprojected-closed}
\end{equation}
Consequently its local coefficient
\begin{equation}
 \mathcal B^{(\ell,\le r)}_j
 :=\mathsf P^{\rm ESM}_{N,6,r}
   \widetilde{\mathcal B}^{(\ell,\le r)}_j
 \label{eq:v3-local-breaking}
\end{equation}
satisfies
\begin{equation}
 \boxed{d_j\mathcal B^{(\ell,\le r)}_j=0.}
 \label{eq:v3-local-breaking-closed}
\end{equation}
\end{theorem}

\begin{proof}
Expand \eqref{eq:v3-quantum-WZ} in powers of $\hbar$.  At the first order at
which the master defect has not yet been cancelled, every term in which a
positive-loop part of $S$ acts on a lower-order master defect vanishes by the
induction hypothesis.  The surviving term is therefore precisely
$(S_{0,j},\widetilde{\mathcal B}_j^{(\ell,\le r)})$.  This proves
\eqref{eq:v3-unprojected-closed}.  Apply the local projector and use
Theorem~\ref{thm:v3-chain-complex}.  Because all source contacts were assembled
before projection, no differentiated contact is omitted in this step.
\end{proof}

\begin{remark}[Explicit first two coefficients]
With the convention \eqref{eq:v3-action-expansion}, the formal master
coefficients are
\begin{align}
 \mathcal M_j^{(0)}&=\frac12(S_{0,j},S_{0,j}),\nonumber\\
 \mathcal M_j^{(1)}&=(S_{0,j},S_{1,j})-\Delta_jS_{0,j},
 \label{eq:v3-master-coefficients}\\
 \mathcal M_j^{(2)}&=(S_{0,j},S_{2,j})
  +\frac12(S_{1,j},S_{1,j})-\Delta_jS_{1,j}.\nonumber
\end{align}
Thus the order-two row cannot be repaired independently of the already chosen
order-one action.  V3 always forms the complete coefficient first and only then
applies the common local restorer.
\end{remark}

\section{A gauge-invariant local frame for the complete chiral packet}
\label{sec:v3-invariant-frame}

We now use the strongest native finite statement already available instead of
recomputing anomaly traces.  Let $\Sigma_n^{\rm SM}$ denote the tensor product
of the coordinate-free chiral determinant sections of the complete Standard
Model packet at cutoff $n$, including the generation multiplicity selected by
the programme.  On the protected component it is nowhere zero.  The supplied
Grand Tensor gauge-descent theorem and the inherited anomaly packet give
\begin{equation}
 g\cdot\Sigma_n^{\rm SM}=\Sigma_n^{\rm SM}
 \qquad(g\in\mathcal G_n)
 \label{eq:v3-full-section-invariant}
\end{equation}
for the connected finite gauge action on the declared chart.  Its Hermitian
norm is gauge invariant as well.

\begin{definition}[Invariant section frame]
\label{def:v3-invariant-frame}
On the protected nonzero component define
\begin{equation}
 \boxed{
 \lambda^{\rm inv}_n
 :=\frac{\Sigma_n^{\rm SM}}{\|\Sigma_n^{\rm SM}\|}.}
 \label{eq:v3-invariant-frame}
\end{equation}
This is a unit-norm frame of the complete determinant line.  It is a local
reference choice on the protected component, not a statement about global line
holonomy across zero modes or disconnected topological sectors.
\end{definition}

\begin{theorem}[Exact gauge invariance and fixed-order locality of the invariant frame]
\label{thm:v3-invariant-frame}
The frame \eqref{eq:v3-invariant-frame} satisfies
\begin{equation}
 \boxed{g\cdot\lambda^{\rm inv}_n=\lambda^{\rm inv}_n,
 \qquad s_j\lambda^{\rm inv}_n=0.}
 \label{eq:v3-invariant-frame-BRST}
\end{equation}
At every fixed source order required by the two-loop packet, its local source
connection coefficients and the transition coefficients between
$\lambda_n^{\rm inv}$ and the V1 Kato frame obey the same rooted/quasilocal
bounds as the determinant-line derivatives proved in V1--V2.
\end{theorem}

\begin{proof}
Equation \eqref{eq:v3-full-section-invariant} and gauge invariance of the norm
give the first identity.  Differentiating the finite gauge action at the
identity gives the BRST statement.

For regularity, write locally
$\Sigma_n^{\rm SM}=z_n\lambda_n^{\rm K}$ in the V1 Kato frame.  Then
\[
 \lambda_n^{\rm inv}
 =u_n\lambda_n^{\rm K},\qquad
 u_n=\frac{z_n}{|z_n|}\in U(1).
\]
Every fixed derivative of $\log z_n$ is a finite sum of the inverse--vertex,
Kato-frame and line words controlled in V1.  Derivatives of
$\log|z_n|=\operatorname{Re}\log z_n$ are the corresponding real parts.
Therefore derivatives of $\log u_n=i\operatorname{Im}\log z_n$ obey the same
fixed-order rooted/quasilocal estimates.  Multiplying a Kato frame by the
unit-modulus scalar $u_n$ preserves the asserted source bounds.  The argument
uses the nonzero divisor margin and is intentionally local in the protected
component.
\end{proof}

\begin{firewall}[The invariant frame is reference data, not a global physical connection]
The construction uses the complete nonzero determinant section to choose a
local frame.  It does not assert that an independently measured determinant-line
connection equals this frame connection, nor does it identify holonomy around
a large gauge loop.  Such global/physical line data remain separate exactly as
in V1.  E2 needs a local perturbative frame in which the shell master identity
can be restored; that is what Theorem~\ref{thm:v3-invariant-frame} supplies.
\end{firewall}

\section{Gauge invariance survives the exact shell cocycle}
\label{sec:v3-shell-line}

The anomaly question is shellwise, so the full-cutoff frame must be transported
through the exact determinant Schur factorization rather than simply assumed to
factor.

Let $\Sigma_{<j}^{\rm SM}$ denote the determinant section of the retained
soft block before eliminating shell $j$, and let $\sigma_j^{\rm SM}$ be the
coordinate-free shell section in the inherited exact Schur cocycle, so that
schematically
\begin{equation}
 \Sigma_{<j+1}^{\rm SM}
 =\sigma_j^{\rm SM}\otimes\Sigma_{<j}^{\rm SM}.
 \label{eq:v3-shell-section-cocycle}
\end{equation}
All tensor products are the actual determinant-line products; the notation is
not a multiplication of unrelated scalar determinants.

\begin{lemma}[Equivariance of the Schur factorization]
\label{lem:v3-Schur-equivariance}
Assume the hard/soft splitting is gauge covariant.  If a finite chiral block is
written
\[
 D=\begin{pmatrix}A&B\\ C&E\end{pmatrix}
\]
with invertible hard pivot $A$, then under the induced block gauge maps its
Schur complement
$S=E-CA^{-1}B$ transforms covariantly on the soft chiral spaces.  The exact
determinant-line isomorphism
\begin{equation}
 \Det D\simeq \Det A\otimes\Det S
 \label{eq:v3-det-Schur-line}
\end{equation}
is equivariant.
\end{lemma}

\begin{proof}
Write the source and target gauge maps as
$U_-^{\hard}\oplus U_-^{\soft}$ and
$U_+^{\hard}\oplus U_+^{\soft}$.  Covariance gives
\begin{align*}
 A'&=U_-^{\hard}A(U_+^{\hard})^{-1},&
 B'&=U_-^{\hard}B(U_+^{\soft})^{-1},\\
 C'&=U_-^{\soft}C(U_+^{\hard})^{-1},&
 E'&=U_-^{\soft}E(U_+^{\soft})^{-1}.
\end{align*}
Substitution gives
$S'=U_-^{\soft}S(U_+^{\soft})^{-1}$.
Taking top exterior powers yields the equivariant line isomorphism.  This is an
identity before any scalar frame is chosen.
\end{proof}

\begin{theorem}[The complete Standard-Model shell section is gauge invariant]
\label{thm:v3-shell-section-invariant}
On the protected component, every coordinate-free shell factor in
\eqref{eq:v3-shell-section-cocycle} obeys
\begin{equation}
 \boxed{g\cdot\sigma_j^{\rm SM}=\sigma_j^{\rm SM}.}
 \label{eq:v3-shell-section-invariant}
\end{equation}
The same statement holds for a locally oriented neutral-Majorana Pfaffian
factor when that optional block is included.
\end{theorem}

\begin{proof}
The full-cutoff complete Standard-Model sections on the two sides of one shell
are gauge invariant by the inherited gauge-descent theorem.  The splitting is
gauge covariant, and Lemma~\ref{lem:v3-Schur-equivariance} makes the exact shell
factorization equivariant.  Since the retained soft section is nonzero on the
protected component, tensor cancellation in the one-dimensional line gives
\eqref{eq:v3-shell-section-invariant}.  Equivalently, the shell section is the
line ratio of two invariant nonzero sections.  The Pfaffian statement is the
same argument on the chosen oriented nonzero skew component.
\end{proof}

Define the invariant shell frame
\begin{equation}
 \lambda_{j}^{\rm inv}
 :=\frac{\sigma_j^{\rm SM}}{\|\sigma_j^{\rm SM}\|}.
 \label{eq:v3-shell-invariant-frame}
\end{equation}
Then the scalar shell coefficient in this frame is the positive function
$\|\sigma_j^{\rm SM}\|$.  In particular its internal-gauge BRST variation is
identically zero at finite cutoff.

\begin{corollary}[Frame-change primitive in the V1 Kato trivialization]
\label{cor:v3-frame-change}
Let $\lambda_j^{\rm K}$ be the V1 Kato shell frame and write
\begin{equation}
 \lambda_j^{\rm inv}=e^{\vartheta_j}\lambda_j^{\rm K},
 \qquad \operatorname{Re}\vartheta_j=0
 \label{eq:v3-frame-transition}
\end{equation}
with a smooth logarithm on a sufficiently small star-shaped source chart.
Then the scalar shell actions in the two frames differ by the single local-line
reference term $\vartheta_j$, and every fixed differentiated coefficient of
$\vartheta_j$ belongs to the V1--V2 line source calculus.  After the common
local Taylor assignment,
\begin{equation}
 \|\partial_I(\vartheta_j-\mathsf T_j\vartheta_j)\|_{\rm pan}
 \le C_I\frac{\mu}{\Lambda_j}.
 \label{eq:v3-frame-transition-remainder}
\end{equation}
\end{corollary}

\begin{proof}
The transition is a smooth $U(1)$-valued function because both frames are
unit-norm and nonzero.  The source derivative statement is
Theorem~\ref{thm:v3-invariant-frame}.  The one-excess-scale remainder is exactly
the V2 line-excess theorem applied to the difference of two admissible local
line frames; its derivatives are finite linear combinations of the same
line-logarithmic words.  The term is retained as reference/source data and is
not discarded because it is field independent in any restricted field sector.
\end{proof}

\section{Native zero of the complete standard anomaly projection}
\label{sec:v3-anomaly-zero}

We can now prove the E2 obstruction statement at the level actually consumed by
the local BV range theorem.

\begin{proposition}[One-loop native trace factorization before assembly]
\label{prop:v3-trace-factorization}
For a single left-handed Weyl multiplet $R$ using the common V2 Wilson/overlap
kinetic profile, the ghost-linear zero-Higgs local projection of the one-loop
chiral shell breaking onto \eqref{eq:v3-anomaly-space} has the form
\begin{equation}
 \mathsf A_{\rm std}\mathcal B^{(1,0)}_{j,R}
 =\sum_A\kappa_{A,j}\,t_A(R)\,\mathcal A_A,
 \label{eq:v3-trace-factorization}
\end{equation}
where $\kappa_{A,j}$ depends on the common native momentum/spin regulator and
normalization convention but not on the internal representation $R$, while
\begin{equation}
 \begin{aligned}
 t_{333}(R)&=\tr_R T_3^{(a}\!T_3^bT_3^{c)},&
 t_{Y33}(R)&=\tr_R(YT_3^aT_3^b),\\
 t_{Y22}(R)&=\tr_R(YT_2^iT_2^k),&
 t_{YYY}(R)&=\tr_RY^3,\\
 t_{YRR}(R)&=\tr_RY.
 \end{aligned}
 \label{eq:v3-trace-tensors}
\end{equation}
The proposition makes no claim about the numerical ordinary-continuum
normalization of the $\kappa_{A,j}$.
\end{proposition}

\begin{proof}
Expand the common gauge-covariant Wilson links and the overlap functional
calculus in external gauge/ghost jets before applying the anomaly quotient.
The lattice/spin/momentum part is the same for every species because the native
kinetic profile is common.  Representation dependence enters only through the
matrices multiplying the external gauge and ghost insertions.  The standard
anomaly projection selects the ghost-one local jet with precisely the invariant
index tensor of the corresponding basis element.  Taking the finite internal
trace therefore factors the coefficient into a common kinematic number and the
symmetrized representation traces in \eqref{eq:v3-trace-tensors}.  The mixed
gravity row contains one hypercharge generator and the spin curvature trace,
so its representation factor is $\tr_RY$.

Gauge-covariant Yukawa/Higgs insertions either carry positive Higgs degree and
are killed by the zero-Higgs anomaly-coordinate evaluation or belong, by the
matter-inclusive cohomological reduction, to the exact complement of
$\mathfrak A_{\rm std}$.  The same is true of antifield/source contacts not
belonging to the five quotient coordinates.  Thus no additional
representation tensor contributes to \eqref{eq:v3-trace-factorization}.
\end{proof}

\begin{remark}[Why V3 does not normalize $\kappa_{A,j}$]
The inherited Standard-Model trace vector is zero componentwise.  Hence the
complete one-loop anomaly projection vanishes for arbitrary finite values of
the common $\kappa_{A,j}$.  Proving that each $\kappa_{A,j}$ equals a particular
ordinary-continuum convention is a regulator-normalization theorem for an
\emph{uncancelled} representation; it is not needed to prove the master identity
of the anomaly-free packet.  V3 therefore does not pay that extra constant.
\end{remark}

\begin{theorem}[Zero standard anomaly projection for the complete native shell]
\label{thm:v3-native-anomaly-zero}
For the complete Standard-Model packet on the V2 protected branch, and for the
complete first- and second-loop shell master breakings after all lower-order
restorers and all source/line/reference contacts have been assembled,
\begin{equation}
 \boxed{
 \mathsf A_{\rm std}\mathcal B_j^{(1,\le r)}=0,
 \qquad
 \mathsf A_{\rm std}\mathcal B_j^{(2,\le r)}=0.}
 \label{eq:v3-native-anomaly-zero}
\end{equation}
More generally, every coefficient obtained by differentiating the same finite
shell identity through the prescribed finite source order has zero projection.
\end{theorem}

\begin{proof}
At one loop, Proposition~\ref{prop:v3-trace-factorization} reduces the five
quotient coordinates to the five inherited Standard-Model trace sums.  Those
sums vanish componentwise for the complete generation packet, and generation
replication multiplies zero by the generation number.  The inherited even
$SU(2)$ doublet parity handles the separate global $SU(2)$ sign; it is not a
local ghost-one coordinate in \eqref{eq:v3-anomaly-space}.

For the complete native shell there is also a stronger finite statement.
Choose the invariant shell frame \eqref{eq:v3-shell-invariant-frame}.  The
chiral Berezin line factor is then internally gauge invariant exactly.  Every
interaction vertex in the finite Standard-Model action is gauge covariant, the
bosonic Haar/Lebesgue and ghost/nonminimal measures have the V80 finite
invariance properties, and the hard/soft blocking is gauge covariant.  Therefore
the complete finite hard pushforward, including non-Gaussian connected terms,
has zero internal-gauge variation before a loop expansion is taken.  Extracting
the coefficient of $\hbar^2$, differentiating in retained sources, and finally
applying the local anomaly-coordinate projection preserve this exact zero.
This proves the second identity without an Adler--Bardeen input.

If instead one works in the Kato frame, Corollary~\ref{cor:v3-frame-change}
shows that the change from the invariant frame is one explicit local-line
reference coboundary plus a quasilocal $O(\mu/\Lambda_j)$ remainder.  A change
of local frame therefore cannot create a nontrivial quotient coordinate.  All
mixed-shell line contacts remain present because the frame change is made
before the exact differentiated shell product is expanded.
\end{proof}

\begin{firewall}[Why exact finite gauge invariance does not trivialize the whole BV problem]
Theorem~\ref{thm:v3-native-anomaly-zero} removes only the nontrivial
internal-gauge quotient.  The cutoff shell can still generate local BV-exact
breaking in the diffeomorphism, local-Lorentz, source, antifield, gauge-fixing,
stress/contact and frame/reference rows.  Those terms must be removed by one
common primitive.  Gauge invariance of the chiral line is therefore not a
substitute for the remaining E2 restoration theorem.
\end{firewall}

\section{One common local primitive by finite Hodge contraction}
\label{sec:v3-hodge-restorer}

Theorem~\ref{thm:v3-breaking-closed} and
Theorem~\ref{thm:v3-native-anomaly-zero} put the complete breaking in the exact
part of the finite source-extended complex.  We now construct the restorer
without choosing counterterms graph by graph.

Choose once and for all a positive coefficient inner product on
$\widehat{\mathfrak L}^{\rm ESM}_{N,6,r}$, compatible with the finite
representation decomposition but otherwise arbitrary.  Let $d_j^\dagger$ be
the adjoint and
\begin{equation}
 \Box_j=d_jd_j^\dagger+d_j^\dagger d_j.
 \label{eq:v3-Hodge-Laplacian}
\end{equation}
Let $\mathsf H_j$ be the orthogonal projection onto $\ker\Box_j$ and let
$G_j$ be the inverse of $\Box_j$ on $(\ker\Box_j)^\perp$, extended by zero on
the harmonic space.  Define
\begin{equation}
 h_j:=d_j^\dagger G_j.
 \label{eq:v3-Hodge-homotopy}
\end{equation}

\begin{lemma}[Finite contraction identity]
\label{lem:v3-contraction}
On the complete finite source bank,
\begin{equation}
 \boxed{d_jh_j+h_jd_j=I-\mathsf H_j.}
 \label{eq:v3-contraction-identity}
\end{equation}
If $B$ has ghost number one, $d_jB=0$, and
$\mathsf A_{\rm std}B=0$, then
\begin{equation}
 \boxed{B=d_jh_jB.}
 \label{eq:v3-exact-hodge}
\end{equation}
\end{lemma}

\begin{proof}
Finite-dimensional Hodge decomposition gives
$\widehat{\mathfrak L}=\operatorname{Ran}d_j\oplus
\ker\Box_j\oplus\operatorname{Ran}d_j^\dagger$ orthogonally and the standard
Green identity gives \eqref{eq:v3-contraction-identity}.  By
Theorem~\ref{thm:v3-chain-complex}, the only nontrivial ghost-one harmonic
coordinates are the standard anomaly quotient, up to the already removed
contractible pairs and redundant directions.  Thus
$\mathsf A_{\rm std}B=0$ and $d_jB=0$ imply $\mathsf H_jB=0$.  Apply
\eqref{eq:v3-contraction-identity} to $B$.
\end{proof}

\begin{proposition}[Uniform finite-order restorer on a protected constant-rank chart]
\label{prop:v3-uniform-homotopy}
Shrink the finite perturbative normalization chart, if necessary, to one compact
component on which the ranks of $d_j$ in the retained degrees are constant.
Then there is $\delta_*>0$, independent of shell number on that chart, such that
the nonzero singular values of $d_j$ are at least $\delta_*$.  Consequently
\begin{equation}
 \|h_j\|\le \delta_*^{-1},
 \qquad
 \|G_j\|\le \delta_*^{-2}.
 \label{eq:v3-homotopy-bound}
\end{equation}
\end{proposition}

\begin{proof}
At fixed $(N,6,r)$ the matrices of $d_j$ are finite and depend continuously on
the finitely many normalized couplings/reference coordinates.  Constant rank
means the smallest nonzero singular value is a positive continuous function on
the chosen compact component.  Its minimum is therefore positive.  The two
operator bounds follow from the singular-value formulas for the Moore--Penrose
inverse and for the Green operator.
\end{proof}

\begin{theorem}[One common source-complete restoring primitive]
\label{thm:v3-common-primitive}
For $\ell=1,2$, let
$\mathcal B_j^{(\ell,\le r)}$ be the recursively assembled local shell breaking.
Define
\begin{equation}
 \boxed{
 \Xi_j^{(\ell,\le r)}:=h_j\mathcal B_j^{(\ell,\le r)}.}
 \label{eq:v3-common-primitive}
\end{equation}
Then
\begin{equation}
 \boxed{
 d_j\Xi_j^{(\ell,\le r)}
 =\mathcal B_j^{(\ell,\le r)}.}
 \label{eq:v3-common-primitive-equation}
\end{equation}
The object $\Xi_j^{(\ell,\le r)}$ is one element of the complete
source-extended local bank.  Its gauge, diffeomorphism, local-Lorentz,
Higgs/Yukawa, ghost, antifield, stress/contact and determinant/Pfaffian-line
coefficients are its source/field components; no descendant is chosen
independently.
\end{theorem}

\begin{proof}
Closure is Theorem~\ref{thm:v3-breaking-closed}; zero harmonic projection is
Theorem~\ref{thm:v3-native-anomaly-zero} together with the V80 cohomological
reduction.  Lemma~\ref{lem:v3-contraction} gives
\eqref{eq:v3-common-primitive-equation}.  Because the source jet algebra was
included in \eqref{eq:v3-source-extended-bank} before the Hodge problem was
solved, the primitive is a single polynomial in all retained source variables.
Differentiating it merely reads off coefficients already present in that one
element.  In particular, contacts generated because the differential itself
acts on a source/antifield multiplet are part of the same finite matrix $d_j$;
they are not added afterward.
\end{proof}

\begin{corollary}[Restoring counterterms through two loops]
\label{cor:v3-two-loop-counterterms}
At the first un-restored order set
\begin{equation}
 C_{\ell,j}:=-\Xi_j^{(\ell,\le r)}.
 \label{eq:v3-counterterm}
\end{equation}
Insert $C_{1,j}$ into the common action before the complete order-two master
coefficient is formed; then construct $C_{2,j}$ from that recursively corrected
order-two breaking.  The resulting action satisfies the projected local quantum
master equation through $O(\hbar^2)$ on the declared source bank.
\end{corollary}

\begin{proof}
Adding $\hbar^\ell C_{\ell,j}$ changes the first un-restored coefficient of the
master functional by $d_jC_{\ell,j}$; all terms involving lower master defects
vanish by recursive restoration.  Equation
\eqref{eq:v3-common-primitive-equation} cancels the coefficient.  At order two,
forming the breaking only after $C_{1,j}$ has been inserted automatically
retains the induced terms from
$\frac12(S_{1,j},S_{1,j})-\Delta_jS_{1,j}$, including every source contact.
\end{proof}

\section{The chiral two-loop hard remainder has the same one-excess-scale gain}
\label{sec:v3-two-loop-remainder}

Local solvability alone is not a shell self-map.  We also need the complementary
rooted kernels.  V1--V2 supply the first-order fermion hard-line row/column
bounds, every fixed source derivative, the full-zone cap estimate, and the
$O(\mu/\Lambda_j)$ line remainder.  The remaining point is that the inherited
V80 two-loop forest ownership still gives one excess hard scale when fermion
lines are admitted.

Let $G$ be a fixed connected two-loop ESM hard graph, possibly with mixed
bosonic and fermionic hard lines, after every proper divergent one-loop subgraph
has been assigned to the same common order-one action.  Let $\omega_G$ be its
superficial local degree in the declared EFT/source counting and let
$\mathsf T_G$ be the complete Taylor/EFT assignment through that degree.

\begin{theorem}[Mixed chiral two-loop hard-annulus remainder]
\label{thm:v3-chiral-annulus}
At fixed $(N,6,r)$ and on the V2 protected branch, every such graph obeys,
after the common forest subtraction,
\begin{equation}
 \boxed{
 \|\partial_K^\alpha(1-\mathsf T_G)\mathcal I_{G,j}(K)\|_{\rm root}
 \le C_{G,\alpha}\frac{\mu}{\Lambda_j}\,
       \mathfrak e_{G,\alpha}(\mu)}
 \label{eq:v3-chiral-graph-remainder}
\end{equation}
for the finite external panel $|K|\le\mu$, where
$\mathfrak e_{G,\alpha}(\mu)$ is the fixed engineering source stock of the
corresponding retained coefficient.  At the ultraviolet-cap shells the same
bound is read as $C a\mu$.  Summing the finite two-loop graph/forest family gives
\begin{equation}
 \boxed{
 \|(1-\mathsf P^{\rm ESM}_{N,6,r})\mathcal K^{\rm ESM}_{2,j}\|_{*\to *}
 \le C_{N,r}\frac{\mu}{\Lambda_j}.}
 \label{eq:v3-K2-bound}
\end{equation}
\end{theorem}

\begin{proof}
For each physical annulus, rescale every independent hard momentum by
$p=\Lambda_jk$.  Bosonic hard propagators and vertices have the inherited V80
stocks.  Every fermionic hard propagator has the weighted row/column stock
$O(\Lambda_j^{-1})$ and every fixed source derivative is a finite inverse--vertex
word with the V2 normalized engineering bounds.  Hence every fixed two-loop
integrand, after the proper one-loop forests have been subtracted, is a smooth
function of the dimensionless loop variables and of $K/\Lambda_j$, multiplied
by its declared engineering power.

The local bank contains the complete Taylor jet through the superficial degree
$\omega_G$.  The integral Taylor formula for
$(1-\mathsf T_G)\mathcal I_{G,j}$ therefore leaves at least one additional
factor $|K|/\Lambda_j$ beyond the last local coefficient.  The V1 rooted
connected-contraction theorem converts the weighted row/column bounds into a
volume-uniform rooted kernel estimate.  The inherited V80 forest ownership
ensures that overlapping proper subgraphs are subtracted by restrictions of the
same $S_{1,j}$ rather than by independent graph constants, so the preceding
Taylor gain applies to the complete forest sum.

There are only finitely many graph and forest types at fixed loop order, EFT
weight, arity and source order; summing their constants proves
\eqref{eq:v3-K2-bound}.  On the ultraviolet cap, V2 replaces the physical
annular variable by $x=ap$ and gives one factor $a\mu$, which is
$O(\mu/\Lambda_j)$ because $\Lambda_j\asymp a^{-1}$ there.
\end{proof}

\begin{remark}[No hidden fermionic positivity]
The proof uses the V1 weighted Schur/rooted bounds and ordinary trace-ideal
estimates for Grassmann contractions.  It never replaces a fermionic connected
trace by a positive bosonic variance and never assigns two independent small
factors to one connected contraction.
\end{remark}

\begin{corollary}[Common chiral source and line step]
\label{cor:v3-source-step}
Let $E^{\rm ESM}_{1,j}$ and $E^{\rm ESM}_{2,j}$ be the first- and second-loop
adjacent source/line steps after the local restorers have been applied.  Then
\begin{equation}
 \|E^{\rm ESM}_{1,j}\|_{*\to *}
 +\|E^{\rm ESM}_{2,j}\|_{*\to *}
 \le C_{N,r}\frac{\mu}{\Lambda_j},
 \label{eq:v3-source-step}
\end{equation}
with the quadratic product term bounded by
$C(\mu/\Lambda_j)^2$.  All differentiated determinant/Pfaffian products retain
the mixed-shell Bell contacts of V1.
\end{corollary}

\begin{proof}
The graph part is Theorem~\ref{thm:v3-chiral-annulus}; the line part is the V2
line-excess theorem and its Bell-polynomial corollary.  The local Hodge restorer
is finite dimensional and bounded by Proposition~\ref{prop:v3-uniform-homotopy},
so applying it to the local coefficient does not alter the complementary
one-excess-scale estimate.  Products of two first-order source steps have the
stated quadratic bound.
\end{proof}

\begin{corollary}[Uniform finite-order shell composition]
\label{cor:v3-source-product}
For geometric shell scales $\Lambda_j=b^j\lambda$,
\begin{equation}
 \sup_{n\ge j}
 \|\mathsf U^{\rm ESM,[\le2]}_{n,j}\|_{*\to *}
 \le
 \exp\!\left(C\sum_{m\ge j}\frac{\mu}{\Lambda_m}\right)
 \le
 \exp\!\left(\frac{Cb}{b-1}\frac{\mu}{\Lambda_j}\right).
 \label{eq:v3-composed-source}
\end{equation}
The bound is coefficientwise at the fixed two-loop/EFT/source order; it is not
an analytic finite-coupling RG trajectory.
\end{corollary}

\section{One common two-loop Einstein--Standard Model shell update}
\label{sec:v3-selfmap}

We can now assemble the E2 acceptance criterion.  Let the complete finite-order
state be
\begin{equation}
 \begin{split}
 \mathfrak S_j^{\rm ESM}=
 \bigl(&\mathfrak S_j^{\rm EGH,bos};
 Z_{\psi,j},\mathbf Y_j,w^{F}_{6,j};
 \mathcal N^{F}_{<,j},\mathcal G^{F}_{<,j};\\
 &\lambda^{\rm inv}_{\chi,j},\mathcal L_{\chi,j};
 \mathcal J^{\rm all}_j,\mathcal T^{\rm all}_j;
 \mathcal C_{1,j},\mathcal C_{2,j}\bigr).
 \end{split}
 \label{eq:v3-ESM-state}
\end{equation}
Here $w^F_{6,j}$ denotes the complete fermionic ghost-zero weight-six quotient,
not a hand-picked operator list; it contains all admitted dipole, current,
four-fermion, curvature--fermion, Higgs--fermion and redundant partner
directions selected by the finite quotient.  $\mathcal L_{\chi,j}$ denotes the
finite determinant/Pfaffian line-source/contact packet.  The symbols
$\mathcal C_{\ell,j}$ record the common local restoring coefficients generated
by \eqref{eq:v3-common-primitive}, not separate graphwise choices.

\begin{theorem}[Complete chiral ESM two-loop shell self-map]
\label{thm:v3-ESM-selfmap}
Fix the V2 Wilson/overlap branch, a protected nonzero determinant/Pfaffian
component, the positive IR/reference prescription, two-loop order, EFT weight
six, source order $r$, arity $N$, and the compact constant-rank BV chart of
Proposition~\ref{prop:v3-uniform-homotopy}.  After imposing one common set of
physical lower-weight normalization conditions, the normalized native shell
defines
\begin{equation}
 \boxed{
 \mathcal R_j^{\rm ESM,[\le2]}:
 \mathfrak C^{\rm ESM}_{N,6,r}
 \longrightarrow
 \mathfrak C^{\rm ESM}_{N,6,r}.}
 \label{eq:v3-ESM-selfmap}
\end{equation}
It has the following properties.
\begin{enumerate}[label=\textup{(E2.\arabic*)},leftmargin=3.4em]
\item The complete native internal-gauge anomaly projection is zero at one and
two loops, including all retained source descendants.
\item The complete local ghost-one breaking is cancelled by the single
source-extended primitives $\Xi_j^{(1,\le r)}$ and
$\Xi_j^{(2,\le r)}$; no graph or sector receives an independently chosen
restorer.
\item Independent ghost-zero ESM EFT coordinates are retained in the state and
are not removed merely because they occur in the same Taylor bank as a
restoring direction.
\item The mixed boson--fermion quasilocal two-loop complement obeys
\eqref{eq:v3-K2-bound}, and the determinant/Pfaffian line remainder obeys the
V2 $O(\mu/\Lambda_j)$ estimate with all Schur and mixed-shell contacts retained.
\item The adjacent source/line map obeys \eqref{eq:v3-source-step} and composes
uniformly as in \eqref{eq:v3-composed-source} at this fixed perturbative order.
\item Gravity, ghosts, Higgs, all three internal gauge fields, chiral fermions,
Yukawas, antifields, measure/reference terms and stress/contact sources occur
in the same action/source/line update.
\end{enumerate}
Thus the E2 acceptance criterion --- one common two-loop Einstein--Standard
Model action/source/line update --- is satisfied on this explicit native branch.
\end{theorem}

\begin{proof}
The inherited V80 theorem supplies the bosonic EGH self-map and its common
normalization.  Theorem~\ref{thm:v3-native-anomaly-zero} removes the only
nontrivial ghost-one quotient of the complete matter-inclusive local complex.
Theorem~\ref{thm:v3-common-primitive} and
Corollary~\ref{cor:v3-two-loop-counterterms} restore the remaining local
breaking with one finite action.  Theorem~\ref{thm:v3-chiral-annulus} closes the
new mixed chiral hard remainder; the V2 line theorem closes the line factor;
and Corollaries~\ref{cor:v3-source-step}--\ref{cor:v3-source-product} close the
source transport.  All of these constructions use the same source bank and the
same lower-order action before any sector is projected, so the result is one
shell map rather than a product of independently repaired sector maps.
\end{proof}

\begin{firewall}[Exact scope of the V3 E2 closure]
Theorem~\ref{thm:v3-ESM-selfmap} is a fixed-order coefficientwise theorem.  It
does not prove convergence of the perturbation series, uniformity in loop
order, an analytic finite-coupling RG trajectory, global determinant-line
holonomy, passage through chiral zero modes, removal of the retained IR scale,
or a nonperturbative ESM measure.  It is native to the explicit V2 Wilson
completion unless a separate regulator-comparison theorem is supplied.
\end{firewall}

\section{What remains for E3: the common master coefficient must be written as one object}
\label{sec:v3-E3-handoff}

E2 has now been closed at the shell-self-map level.  The next danger is a more
subtle kind of sectorwise reasoning: one could cite the bosonic V80 master row,
the chiral line invariance, and the Hodge restorer separately and declare the
complete ESM master equation solved.  E3 forbids that shortcut.

Define the common renormalized coefficients after the V3 shell update by
\begin{equation}
 \bar S_{1,j}=S_{1,j}+C_{1,j},
 \qquad
 \bar S_{2,j}=S_{2,j}+C_{2,j}.
 \label{eq:v3-renormalized-S12}
\end{equation}
The next theorem must write, in the same source/reference convention,
\begin{align}
 \mathcal M_{j,1}^{\rm all}
 &= (S_{0,j},\bar S_{1,j})-\Delta_jS_{0,j},
 \label{eq:v3-E3-M1}\\
 \mathcal M_{j,2}^{\rm all}
 &= (S_{0,j},\bar S_{2,j})
  +\frac12(\bar S_{1,j},\bar S_{1,j})
  -\Delta_j\bar S_{1,j},
 \label{eq:v3-E3-M2}
\end{align}
as one complete coefficient and verify every component with the same
$S_0,\bar S_1,\bar S_2$.  In particular, the cross antibrackets between
geometric, gauge, Higgs/Yukawa, chiral, ghost and source sectors must not be
replaced by separate Ward equations.

\begin{theorem}[Exact next load-bearing theorem for V4]
\label{thm:v3-next-E3}
Using the V3 common shell action, prove the full two-loop common
BV/master/source equation by expanding
\eqref{eq:v3-E3-M1}--\eqref{eq:v3-E3-M2} in the single source-extended bank and
checking simultaneously:
\begin{enumerate}[label=\textup{(M\arabic*)},leftmargin=3em]
\item diffeomorphism and retained gravitational constraint rows;
\item local Lorentz/spin rows;
\item all three internal gauge rows;
\item Higgs and Yukawa covariance rows;
\item chiral determinant/Pfaffian line and frame/reference rows;
\item ghost, nonminimal and antifield descendants;
\item measure/Jacobian/reference terms;
\item composite, stress and contact sources;
\item every cross antibracket/contact generated by the common $S_1$ in the
quadratic term $\frac12(\bar S_1,\bar S_1)$.
\end{enumerate}
The theorem must end with one coefficient identity
$\mathcal M_{j,1}^{\rm all}=\mathcal M_{j,2}^{\rm all}=0$ in the retained local
bank plus the already controlled quasilocal remainder.  It is not enough to
list sectorwise zeroes.
\end{theorem}

\section{V3 anti-circularity audit}
\label{sec:v3-audit}

The V3 route avoids four possible circularities.
\begin{enumerate}[leftmargin=2.4em]
\item It does not infer a native anomaly statement from the already-known
continuum trace cancellation.  The one-loop trace factorization is derived from
the native common kinetic profile, while the complete two-loop zero follows
from the exact finite gauge invariance of the assembled native shell.
\item It does not invoke an all-order anomaly nonrenormalization theorem.  The
finite gauge-invariance identity is expanded only after the complete shell is
assembled.
\item It does not solve a different counterterm equation for every graph or
source derivative.  The Hodge map $h_j$ acts once on the complete
source-extended breaking.
\item It does not erase line/reference data by choosing a convenient scalar
phase.  The invariant frame and its transition from the Kato frame are retained
as explicit line/reference coefficients with the same source estimates.
\end{enumerate}

\section{V3 status ledger}
\label{sec:v3-status}

\subsection*{Proved in V3}
\begin{enumerate}[leftmargin=2.2em]
\item The complete fixed-order ESM local/source bank is a finite BV chain
complex; at ghost number one its only nontrivial quotient is the standard
five-dimensional internal-gauge anomaly space inherited from V80.
\item The complete native shell master breaking is a BV cocycle at the first
un-restored loop order, including the complete source/contact packet.
\item The anomaly-free nonzero Standard-Model determinant section supplies a
smooth gauge-invariant local unit frame on the protected component.  Its
transition from the V1 Kato frame has the same fixed-order quasilocal source
bounds and an $O(\mu/\Lambda_j)$ locally subtracted shell remainder.
\item Exact Schur equivariance transports that gauge invariance to every shell
factor in the inherited determinant/Pfaffian cocycle.
\item The one-loop native anomaly seed factors into universal shell coefficients
times the five Standard-Model representation traces.  Their inherited
componentwise cancellation gives zero one-loop anomaly projection without
normalizing the universal coefficients.
\item The complete finite native shell is internally gauge invariant after the
full anomaly-free packet is assembled.  Hence its two-loop standard anomaly
projection and every retained source derivative of that projection vanish
exactly, without importing Adler--Bardeen nonrenormalization.
\item A single finite Hodge contraction on the complete source-extended bank
constructs the restoring primitive.  All gauge, diffeomorphism, Lorentz,
Higgs/Yukawa, ghost, antifield, stress/contact and line descendants belong to
that one primitive.
\item Mixed boson--fermion two-loop hard forests have the same one-excess-scale
$O(\mu/\Lambda_j)$ complementary rooted bound at fixed EFT/source order.
Together with the V2 line estimate this yields a uniformly composable adjacent
source/line step.
\item The inherited V80 bosonic self-map therefore enlarges to one complete
chiral Einstein--Standard Model two-loop action/source/line shell self-map on
the explicit V2 regulator branch.  E2 is closed at that scope.
\end{enumerate}

\subsection*{Inherited}
\begin{enumerate}[leftmargin=2.2em]
\item The V80 bosonic Einstein--gauge--Higgs two-loop self-map, finite local
cohomology reduction, zero native bosonic anomaly projection and common source
normalization.
\item The three-generation Standard-Model chiral representation, cancellation of
the five ordinary local anomaly trace sums, even $SU(2)$ doublet parity, and the
Grand Tensor gauge descent of the complete nonzero chiral determinant section.
\item Exact determinant/Pfaffian Schur factorization, finite regrouping cocycles
and differentiated line identities.
\item The V1 weighted functional calculus, Kato frame, hard-pivot and rooted
connected-contraction machinery.
\item The V2 explicit Wilson completion, all-zone gap/no-doubler theorem,
annular/full-zone finite-to-ordinary bounds, fixed source-jet calculus and
$O(\mu/\Lambda_j)$ differentiated line remainder.
\end{enumerate}

\subsection*{Conditional}
\begin{enumerate}[leftmargin=2.2em]
\item The result remains native to the explicit V2 Wilson/overlap completion.
A different microscopic Wilson kernel requires a finite regulator-comparison
theorem.
\item The Hodge-restorer norm is uniform on a compact constant-rank component of
the finite BV bank.  A cohomology-rank crossing requires a new local chart or a
separate analysis.
\item The determinant/Pfaffian construction remains on the declared nonzero
protected component and retains the positive IR/reference prescription.  No
global line holonomy or zero-crossing statement is made.
\item Constants may depend on fixed loop order, EFT degree, source order, arity,
representation data and the protected compact chart.  No uniform
$L\to\infty$ estimate is asserted.
\end{enumerate}

\subsection*{Open}
\begin{enumerate}[leftmargin=2.2em]
\item Write the \emph{actual common} first- and second-loop coefficients
\eqref{eq:v3-E3-M1}--\eqref{eq:v3-E3-M2} in one source-extended bank and audit
every cross antibracket/contact.  This is E3.
\item After E3, formulate the arbitrary fixed-loop induction state including the
finite EFT bank, complete contact tower, BV fields/antifields, line/frame data,
forest ownership, rooted remainders and reference/scheme data.  This is E4.
\item The ordinary-normalization constants $\kappa_{A,j}$ of an uncancelled
single Weyl representation were intentionally not identified; that comparison
is not load-bearing for the anomaly-free E2 theorem.
\end{enumerate}

\subsection*{Exact next load-bearing theorem}
\begin{center}
\fbox{\parbox{0.92\textwidth}{
\textbf{E3: one complete two-loop BV/master/source coefficient.}
Using the V3 common action $S_0+\hbar\bar S_1+\hbar^2\bar S_2$, expand the same
finite quantum master functional through order $\hbar^2$ and prove one common
identity containing the diffeomorphism, local-Lorentz, internal-gauge,
Higgs/Yukawa, chiral line, ghost/nonminimal, antifield, measure/reference,
stress/composite and contact rows, including every cross term in
$\frac12(\bar S_1,\bar S_1)$.  Do not replace this by sectorwise Ward identities.
}}
\end{center}

\section*{Append-only continuation marker: V3}
\addcontentsline{toc}{section}{Append-only continuation marker: V3}
\textbf{End of V3.}  The complete V1--V2 body above is retained verbatim.  Future
versions must append new work after this marker without replacing the V1--V3
mathematical body.  The next file is to be named
\texttt{Renewal\_Perturbative\_ESM\_Continuum\_Research\_Notes\_V4.tex}.


\clearpage
\part*{V4: The complete common two-loop master coefficient and the canonical-pair audit}
\addcontentsline{toc}{part}{V4: Complete common two-loop master coefficient and canonical-pair audit}
\label{sec:v4-start}

\section{E3 is a common-coefficient theorem, not another sectorwise restoration}
\label{sec:v4-E3-start}

V3 closes the two-loop shell self-map on the explicit V2 Wilson/overlap branch,
but deliberately leaves one bookkeeping theorem open.  The local Hodge
construction proves that the recursively assembled shell breaking is in the
range of one source-extended BV differential.  E3 asks for a stronger audit:
the actual coefficients of the \emph{same} finite master functional must be
written with the same $S_0,\bar S_1,\bar S_2$, including every cross
antibracket, source contact, determinant-line/reference contribution and
source-dependent BV-Laplacian term.

The correct upstream organization is by canonical BV pair.  The antibracket is
diagonal in canonical pairs.  Therefore every apparently ``mixed-sector'' term
is produced by one of finitely many canonical contractions, and there is no
additional cross Ward identity to impose after the common coefficient has been
formed.  This observation is what makes an exhaustive audit possible without
re-fitting graphs or writing unrelated sector equations.

Throughout V4 fix exactly the V3 domain: loop order through two, local weight
six, arity $N$, source order $r$, the finite external panel, the compact
constant-rank BV chart, the protected nonzero chiral line component, the
positive IR/reference prescription, and the explicit V2 Wilson completion.
All statements are coefficientwise on this finite domain.

\begin{firewall}[No strengthening of the global claim]
V4 does not assert that the full finite-coupling effective action exists as a
convergent analytic object.  The symbols $S_0,\bar S_1,\bar S_2$ are the fixed
coefficients of the perturbative shell action.  Nor does V4 remove the retained
IR/reference scale, cross a determinant zero, or establish uniformity as loop
order tends to infinity.
\end{firewall}

\section{Canonical carrier and source algebra}
\label{sec:v4-carrier}

\subsection{Central sources versus canonical source multiplets}

Let $\mathcal J_r^{\rm cen}$ be the truncated supercommutative algebra generated
by the retained \emph{central} physical source writers through total source
order $r$.  These are the sources with zero BV bracket.  Source variables which
belong to BRST doublets, antifield/source multiplets, or any other canonical
pair are \emph{not} placed in $\mathcal J_r^{\rm cen}$; they are included in the
canonical carrier below.  This is only a reorganization of the V3
source-extended bank and prevents a contact term from being counted both as a
source derivative and as a BV contraction.

Choose an ordered divided-power basis $J^I/I!$ of
$\mathcal J_r^{\rm cen}$.  For every source-polynomial functional $F(J)$ write
\begin{equation}
 F(J)=\sum_{|I|\le r}\frac{J^I}{I!}F_I,
 \qquad [F]_I:=F_I.
 \label{eq:v4-source-expansion}
\end{equation}
For odd central sources the notation means the ordered super divided-power
basis; permutations carry the corresponding Koszul sign.

\subsection{Exact canonical-pair inventory}

Let $\mathcal C_j$ be the finite list of canonical pairs retained at shell $j$.
We group it only for bookkeeping:
\begin{equation}
 \begin{split}
 \mathcal C_j={}&
 \mathcal C_{\rm geom}\sqcup
 \mathcal C_{\rm Lor}\sqcup
 \mathcal C_{3}\sqcup\mathcal C_{2}\sqcup\mathcal C_Y
 \sqcup\mathcal C_H\sqcup\mathcal C_F\\
 &\sqcup\mathcal C_{\rm diff-gh}\sqcup
 \mathcal C_{\rm Lor-gh}\sqcup
 \mathcal C_{\rm int-gh}\sqcup
 \mathcal C_{\rm nm}\sqcup
 \mathcal C_{\rm src}.
 \end{split}
 \label{eq:v4-canonical-inventory}
\end{equation}
Here $\mathcal C_{\rm geom}$ contains the coframe/metric variables and exact
antifield duals; $\mathcal C_{\rm Lor}$ contains the independent spin/local
Lorentz variables when present in the chosen chart; $\mathcal C_{3,2,Y}$ are
the three internal gauge pairs; $\mathcal C_H$ is the Higgs pair;
$\mathcal C_F$ contains every retained chiral quark/lepton pair (and the
optional neutral real/Majorana pair in its oriented chart);
$\mathcal C_{\rm diff-gh}$, $\mathcal C_{\rm Lor-gh}$ and
$\mathcal C_{\rm int-gh}$ are the ghost--ghost-antifield pairs;
$\mathcal C_{\rm nm}$ contains the contractible nonminimal pairs; and
$\mathcal C_{\rm src}$ contains the genuinely canonical source/marker
multiplets.  Spectator physical writers remain in $\mathcal J_r^{\rm cen}$.

For $a\in\mathcal C_j$ write the canonical coordinates as $(z^a,z_a^*)$.
With the inherited graded left/right derivative convention define the
pair contribution by
\begin{equation}
 (F,G)_a
 :=F\frac{\overleftarrow\delta}{\delta z^a}
      \frac{\overrightarrow\delta}{\delta z_a^*}G
   -F\frac{\overleftarrow\delta}{\delta z_a^*}
      \frac{\overrightarrow\delta}{\delta z^a}G,
 \qquad
 (F,G)=\sum_{a\in\mathcal C_j}(F,G)_a.
 \label{eq:v4-pair-bracket}
\end{equation}
All Koszul signs are contained in the inherited graded arrow-derivative
convention.  Every later argument uses only this exact Darboux-pair
decomposition and is therefore independent of a separate sign shorthand.

The finite BV Laplacian is kept in the same density/reference convention as V3.
Allow its source-dependent representation to have the finite expansion
\begin{equation}
 \Delta_j(J)=\sum_{|K|\le r}\frac{J^K}{K!}\Delta_{j,K}.
 \label{eq:v4-Delta-expansion}
\end{equation}
Thus source derivatives of the finite measure/reference density are retained
rather than silently set to zero.

\section{Exact source-contact calculus for the master functional}
\label{sec:v4-contact-calculus}

\begin{lemma}[Coefficient convolution for the antibracket and BV Laplacian]
\label{lem:v4-source-convolution}
Let $F(J),G(J)$ be finite source polynomials as in
\eqref{eq:v4-source-expansion}.  For every ordered source multi-index $I$,
\begin{equation}
 \boxed{
 [(F,G)]_I
 =\sum_{I_1\sqcup I_2=I}
 \varepsilon(I_1,I_2)
 \binom{I}{I_1}
 (F_{I_1},G_{I_2}),}
 \label{eq:v4-bracket-convolution}
\end{equation}
where $\varepsilon(I_1,I_2)$ is the Koszul sign required to move the odd
sources into the fixed order.  Likewise
\begin{equation}
 \boxed{
 [\Delta_j(J)F(J)]_I
 =\sum_{K\sqcup L=I}
 \varepsilon(K,L)
 \binom{I}{K}\,
 \Delta_{j,K}F_L.}
 \label{eq:v4-Delta-convolution}
\end{equation}
For a purely even source subpanel all signs are $+1$ and the sums are the
ordinary multinomial convolutions.
\end{lemma}

\begin{proof}
Expand both finite polynomials in the ordered divided-power basis.  The BV
functional derivatives act only on the canonical carrier because the writers in
$\mathcal J_r^{\rm cen}$ are central.  Multiplication of two basis monomials
$J^{I_1}/I_1!$ and $J^{I_2}/I_2!$ gives
$\varepsilon(I_1,I_2)\binom{I_1+I_2}{I_1}J^{I_1+I_2}/(I_1+I_2)!$.
Collecting the coefficient of $J^I/I!$ proves
\eqref{eq:v4-bracket-convolution}.  The same multiplication rule applied to the
operator expansion \eqref{eq:v4-Delta-expansion} proves
\eqref{eq:v4-Delta-convolution}.
\end{proof}

\begin{corollary}[The first nontrivial contact is forced]
\label{cor:v4-two-source-contact}
For two distinct even central writers $J^A,J^B$ and any even $F$,
\begin{equation}
 \left[\frac12(F,F)\right]_{AB}
 =(F_0,F_{AB})+(F_A,F_B).
 \label{eq:v4-two-source-contact}
\end{equation}
Thus the product of the two singly marked insertions is part of the same second
source derivative as the doubly marked contact coefficient.  Omitting either
term destroys the differentiated master identity.
\end{corollary}

\begin{proof}
Apply Lemma~\ref{lem:v4-source-convolution} to the four ordered splittings
$\varnothing\sqcup AB$, $A\sqcup B$, $B\sqcup A$ and
$AB\sqcup\varnothing$.  The antibracket of two even functionals is symmetric,
so the prefactor $1/2$ combines the first and last terms and the two mixed
terms.
\end{proof}

\begin{remark}[Why field-independent does not mean source-irrelevant]
A term independent of the canonical fields can have zero antibracket while still
having nonzero coefficients $F_I$ with $I\ne0$.  Such a term therefore remains
part of the normalized source/reference functional.  In addition,
\eqref{eq:v4-Delta-convolution} shows that source dependence of the BV density
itself produces genuine measure/reference contacts.  V4 never discards either
kind of term merely because its unmarked canonical derivative vanishes.
\end{remark}

\section{The canonical-pair incidence theorem}
\label{sec:v4-incidence}

The one-loop coefficient contains contributions with different physical
origins.  Fix once and for all the inherited once-counted decomposition
\begin{equation}
 \bar S_{1,j}=\sum_{\tau\in\mathcal T}\bar S_{1,j}^{\tau},
 \qquad
 \mathcal T=
 \{\mathrm{grav},\mathrm{gauge},H,\mathrm{ferm},Y,
   \mathrm{gh/nm},\mathrm{line},\mathrm{ref},\mathrm{comp}\}.
 \label{eq:v4-origin-decomposition}
\end{equation}
This is an ownership label only.  It does not split the BV differential or
permit separate normalization conditions.  The tags mean, respectively,
gravity/geometric, internal gauge, Higgs, fermion kinetic, Yukawa,
ghost/nonminimal, determinant/Pfaffian line, measure/reference, and
composite/stress/contact contributions.

For a functional $F$ define its canonical support
\begin{equation}
 \operatorname{Supp}_{\BV}(F)
 :=\{a\in\mathcal C_j:
       \delta F/\delta z^a\ne0\ \text{or}\
       \delta F/\delta z_a^*\ne0\}.
 \label{eq:v4-BV-support}
\end{equation}

\begin{lemma}[Canonical-pair support criterion]
\label{lem:v4-support-criterion}
For any two functionals $F,G$,
\begin{equation}
 \boxed{
 (F,G)
 =\sum_{a\in\operatorname{Supp}_{\BV}(F)\cap\operatorname{Supp}_{\BV}(G)}(F,G)_a.}
 \label{eq:v4-support-bracket}
\end{equation}
In particular, disjoint canonical support implies $(F,G)=0$ identically, before
projection or a cutoff limit.
\end{lemma}

\begin{proof}
This is immediate from the pairwise Darboux decomposition
\eqref{eq:v4-pair-bracket}.  A summand vanishes if either functional is
independent of both entries of that canonical pair.
\end{proof}

The following table is the E3 cross-sector audit.  It lists the canonical pair,
not a collection of separate Ward identities.  Every mixed origin in one row
can contract only through the pair displayed in that row.

\begin{center}
\small
\begin{tabular}{p{0.22\textwidth}p{0.70\textwidth}}
\toprule
Canonical pair block & Origin tags which can meet through that pair \\
\midrule
$\mathcal C_{\rm geom}$ &
geometry; gauge and Higgs metric/coframe dependence; fermion kinetic and
Yukawa density/coframe dependence; chiral line; measure/reference;
stress/composite writers; diffeomorphism antifield vertices \\
$\mathcal C_{\rm Lor}$ &
coframe/spin connection; fermion spin transport; Lorentz antifield vertices;
chiral line; geometric reference/source terms \\
$\mathcal C_{3,2,Y}$ &
internal gauge action; charged Higgs kinetic term; chiral fermion kinetic
term; gauge-covariant line determinant; gauge composite insertions; internal
BRST antifield vertices \\
$\mathcal C_H$ &
Higgs kinetic/potential action; Yukawa action; Higgs-dependent chiral line;
Higgs antifield transformation; Higgs/composite/contact writers \\
$\mathcal C_F$ &
chiral kinetic action; Yukawa action; fermion antifield BRST vertices;
fermionic composites and their contacts \\
$\mathcal C_{\rm diff-gh}$ &
all Lie-derivative antifield vertices, gravitational ghost action,
constraint/nonminimal descendants and transformed source multiplets \\
$\mathcal C_{\rm Lor-gh}$ &
frame/spin and chiral-fermion Lorentz transformations, Lorentz ghost and
nonminimal descendants \\
$\mathcal C_{\rm int-gh}$ &
internal gauge, Higgs and chiral-fermion BRST transformations, internal ghost
self-interactions and nonminimal descendants \\
$\mathcal C_{\rm nm}$ &
gauge-fixing fermion, antighost--multiplier packet, retained constraint
quartets and their marked source descendants \\
$\mathcal C_{\rm src}$ &
canonical source/antifield or BRST-doublet writers; all contacts produced by
their exact source differential \\
\bottomrule
\end{tabular}
\end{center}

The determinant/Pfaffian line, reference density and stress/composite rows do
not create mysterious extra antibracket directions.  They are functionals of
the canonical backgrounds and source multiplets, so their cross terms occur in
the appropriate rows of this table.  Their separate bookkeeping labels are
retained only to prevent their coefficients from being dropped.

\begin{theorem}[Exhaustive quadratic cross-term formula]
\label{thm:v4-quadratic-audit}
For every source multi-index $I$, the complete quadratic one-loop contribution
to the two-loop master coefficient is
\begin{equation}
 \boxed{
 \begin{aligned}
 \mathcal Q_{2,j;I}
 :=\left[\frac12(\bar S_{1,j},\bar S_{1,j})\right]_I
 =\frac12
 \sum_{I_1\sqcup I_2=I}
 &\varepsilon(I_1,I_2)\binom{I}{I_1}\\[-0.2em]
 &\times
 \sum_{a\in\mathcal C_j}
 \sum_{\tau,\upsilon\in\mathcal T}
 (\bar S_{1,j;I_1}^{\tau},
  \bar S_{1,j;I_2}^{\upsilon})_a .
 \end{aligned}}
 \label{eq:v4-full-quadratic}
\end{equation}
Every nonzero cross antibracket/contact between the E3 rows occurs exactly in
this finite sum.  There is no additional mixed master term outside
\eqref{eq:v4-full-quadratic}.
\end{theorem}

\begin{proof}
Apply Lemma~\ref{lem:v4-source-convolution}, then the origin decomposition
\eqref{eq:v4-origin-decomposition}, then the canonical-pair decomposition
\eqref{eq:v4-pair-bracket}.  These are identities, so the resulting sum is
exhaustive.  Conversely, every term in the displayed sum is an actual
contraction of one canonical field with its exact dual antifield and therefore
belongs to the master antibracket.  The factor $1/2$ removes the double count
under exchange of the two even one-loop action factors, with the inherited
Koszul convention covering odd marked coefficients.  Lemma
\ref{lem:v4-support-criterion} proves that no contraction can connect two
origin tags through a canonical pair absent from the incidence table.
\end{proof}

\begin{remark}[Examples of terms which are now automatically owned]
The geometric pair produces the stress cross terms between gravitational and
gauge/Higgs/fermion/line contributions.  The internal gauge pairs produce the
gauge--Higgs and gauge--fermion contacts.  The Higgs and fermion pairs produce
the Yukawa covariance contacts.  The diffeomorphism, Lorentz and internal ghost
pairs generate the corresponding antifield descendants.  The source-pair row
generates BRST-doublet contacts.  None of these is assigned an independent
finite constant in V4.
\end{remark}

\section{The actual common first- and second-loop coefficients}
\label{sec:v4-common-coefficients}

Let $S_{1,j}^{\rm pre}$ be the complete one-loop coefficient before the V3
one-loop restoring counterterm is installed.  Let
$S_{2,j}^{\rm pre}[\bar S_{1,j}]$ be the complete forest-prepared two-loop
coefficient in which every proper one-loop subtraction is owned by the same
already corrected $\bar S_{1,j}$.  This notation includes the inherited
canonical second-order packets, measure/reference pieces and all finite
normalization coordinates; it does not mean a new graphwise choice.

Define the un-restored local coefficients
\begin{align}
 B_{1,j}
 &:=\mathsf P^{\rm ESM}_{N,6,r}
 \bigl((S_{0,j},S_{1,j}^{\rm pre})-\Delta_jS_{0,j}\bigr),
 \label{eq:v4-B1}\\
 C_{1,j}&:=-h_jB_{1,j},
 \qquad
 \bar S_{1,j}:=S_{1,j}^{\rm pre}+C_{1,j}.
 \label{eq:v4-S1bar}
\end{align}
By V3, $d_jh_jB_{1,j}=B_{1,j}$.

The order-two breaking must be formed \emph{after} this replacement.  Set
\begin{equation}
 \boxed{
 B_{2,j}:=\mathsf P^{\rm ESM}_{N,6,r}
 \left
 \{
 (S_{0,j},S_{2,j}^{\rm pre}[\bar S_{1,j}])
 +\frac12(\bar S_{1,j},\bar S_{1,j})
 -\Delta_j\bar S_{1,j}
 \right\}.}
 \label{eq:v4-B2-recursive}
\end{equation}
Then
\begin{equation}
 C_{2,j}:=-h_jB_{2,j},
 \qquad
 \bar S_{2,j}:=S_{2,j}^{\rm pre}[\bar S_{1,j}]+C_{2,j}.
 \label{eq:v4-S2bar}
\end{equation}

\begin{lemma}[All one-loop-restorer descendants occur in $B_2$]
\label{lem:v4-C1-descendants}
If for notational comparison one writes
$S_{2,j}^{\rm pre}[\bar S_{1,j}]=S_{2,j}^{\rm pre,0}
 +\delta_{C_1}S_{2,j}^{\rm forest}$, then the complete order-two local breaking
contains
\begin{equation}
 \begin{split}
 B_{2,j}=\mathsf P^{\rm ESM}_{N,6,r}\bigl\{&
 (S_0,S_{2}^{\rm pre,0})
 +(S_0,\delta_{C_1}S_{2}^{\rm forest})
 +\tfrac12(S_1^{\rm pre},S_1^{\rm pre})-\Delta S_1^{\rm pre}\\
 &+(S_1^{\rm pre},C_1)
 +\tfrac12(C_1,C_1)-\Delta C_1
 \bigr\},
 \end{split}
 \label{eq:v4-C1-descendants}
\end{equation}
with all quantities at shell $j$.  Hence the order-two equation cannot use the
uncorrected $S_1^{\rm pre}$ without changing the problem.
\end{lemma}

\begin{proof}
Substitute \eqref{eq:v4-S1bar} into
\eqref{eq:v4-B2-recursive} and use bilinearity plus symmetry of the BV bracket
on even functionals.  The forest term
$\delta_{C_1}S_{2,j}^{\rm forest}$ is present because proper one-loop
subtractions are restrictions of the corrected common one-loop action, by the
inherited forest-ownership theorem.  No further term is available at order two.
\end{proof}

Now define the common master coefficients using the same action coefficients:
\begin{align}
 \mathcal M_{j,1}^{\rm all}
 &:=(S_{0,j},\bar S_{1,j})-\Delta_jS_{0,j},
 \label{eq:v4-M1}\\
 \mathcal M_{j,2}^{\rm all}
 &:=(S_{0,j},\bar S_{2,j})
 +\frac12(\bar S_{1,j},\bar S_{1,j})
 -\Delta_j\bar S_{1,j}.
 \label{eq:v4-M2}
\end{align}
These are exactly the coefficients requested at the end of V3.

\begin{theorem}[Common local two-loop master identity]
\label{thm:v4-common-local-master}
On the V3 protected constant-rank chart,
\begin{equation}
 \boxed{
 \mathsf P^{\rm ESM}_{N,6,r}\mathcal M_{j,1}^{\rm all}=0,
 \qquad
 \mathsf P^{\rm ESM}_{N,6,r}\mathcal M_{j,2}^{\rm all}=0.}
 \label{eq:v4-local-master-zero}
\end{equation}
The two equalities use one $C_{1,j}$ and one $C_{2,j}$ in the complete source
bank.  They are not sums of independently normalized sector equations.
\end{theorem}

\begin{proof}
For order one, use \eqref{eq:v4-S1bar}:
\[
 \mathsf P\mathcal M_{j,1}^{\rm all}
 =B_{1,j}+d_jC_{1,j}
 =B_{1,j}-d_jh_jB_{1,j}=0.
\]
The last equality is the V3 contraction identity, applicable because the
complete native anomaly projection of $B_{1,j}$ vanishes.

At order two, $B_{2,j}$ in \eqref{eq:v4-B2-recursive} is formed with the
already corrected $\bar S_{1,j}$.  V3 proves that it is a $d_j$-closed
source-extended cocycle and that its complete Standard-Model anomaly quotient
vanishes.  Therefore $d_jh_jB_{2,j}=B_{2,j}$.  Substituting
\eqref{eq:v4-S2bar} into \eqref{eq:v4-M2} gives
\[
 \mathsf P\mathcal M_{j,2}^{\rm all}
 =B_{2,j}+d_jC_{2,j}
 =B_{2,j}-d_jh_jB_{2,j}=0.
\]
Theorem~\ref{thm:v4-quadratic-audit} shows that the $B_{2,j}$ used here contains
every canonical-pair cross antibracket and every central-source contact.  Lemma
\ref{lem:v4-source-convolution} and \eqref{eq:v4-Delta-convolution} include all
source-dependent measure/reference contacts.  Thus no sector equation is being
silently omitted from either cancellation.
\end{proof}

\begin{proposition}[No hidden graphwise or rowwise restoring freedom]
\label{prop:v4-no-hidden-refit}
Let $C_{\ell,j}'$ be another local ghost-zero functional which restores the
same complete order-$\ell$ breaking with the same lower-order action.  Then
\begin{equation}
 d_j(C_{\ell,j}'-C_{\ell,j})=0.
 \label{eq:v4-closed-difference}
\end{equation}
On the finite Hodge decomposition its difference is a sum of a $d_j$-exact
redundant direction and a ghost-zero harmonic/EFT Wilson direction.  After the
common physical normalization and Wilson-coordinate convention are fixed, no
additional graphwise or rowwise finite parameter remains.
\end{proposition}

\begin{proof}
Subtract the two restoring equations.  Finite Hodge decomposition at ghost
number zero gives the exact plus harmonic splitting.  Exact directions are
canonical/redundant reparametrizations; harmonic ghost-zero directions are the
physical local Wilson coordinates already retained in the ESM state.  Fixing
the declared normalization and those coordinates fixes the ambiguity.  In
particular, there is no independent freedom indexed by Feynman graph or by E3
row.
\end{proof}

\section{The row audit is a projection of one zero vector}
\label{sec:v4-row-audit}

Choose coefficient projectors $\Pi_\rho$ which resolve the retained ghost-one
bank as a direct sum of bookkeeping rows,
\begin{equation}
 I=\sum_{\rho\in\mathcal R}\Pi_\rho,
 \qquad
 \mathcal R=
 \{\mathrm{diff},\mathrm{Lor},3,2,Y,H/Y,
   \chi\mathrm{-line},\mathrm{gh/nm},\mathrm{meas/ref},
   \mathrm{comp/stress}\}.
 \label{eq:v4-row-projectors}
\end{equation}
The projectors need not define independent subcomplexes; they are only a finite
coordinate resolution of the common coefficient.

\begin{theorem}[Simultaneous E3 row identity]
\label{thm:v4-row-identity}
For $\ell=1,2$ define the row vector
\begin{equation}
 \mathbf M_{j,\ell}
 :=\bigl(\Pi_\rho\mathsf P^{\rm ESM}_{N,6,r}
          \mathcal M_{j,\ell}^{\rm all}\bigr)_{\rho\in\mathcal R}.
 \label{eq:v4-row-vector}
\end{equation}
Then
\begin{equation}
 \boxed{\mathbf M_{j,1}=0,\qquad \mathbf M_{j,2}=0.}
 \label{eq:v4-row-vector-zero}
\end{equation}
More explicitly, the entries have the following common origin:
\begin{enumerate}[label=\textup{(M\arabic*)},leftmargin=3.2em]
\item the diffeomorphism/retained-constraint entry contains every matter,
gauge, line and source stress cross term through $\mathcal C_{\rm geom}$ and
$\mathcal C_{\rm diff-gh}$;
\item the Lorentz/spin entry contains the frame--spinor and line-frame contacts
through $\mathcal C_{\rm Lor}$ and $\mathcal C_{\rm Lor-gh}$;
\item the $SU(3)$, $SU(2)$ and $U(1)_Y$ entries contain the gauge--Higgs,
gauge--fermion and gauge-line contacts through $\mathcal C_{3,2,Y}$ and
$\mathcal C_{\rm int-gh}$;
\item the Higgs/Yukawa entry contains Higgs--fermion/Yukawa and marked Higgs
contacts through $\mathcal C_H$ and $\mathcal C_F$;
\item the chiral-line entry contains the determinant/Pfaffian frame and
line-source coefficients but uses the same background canonical pairs as the
preceding rows;
\item the ghost/nonminimal entry contains the exact contractible descendants
and retained constraint quartets;
\item the measure/reference entry contains every $\Delta_{j,K}$ contact in
\eqref{eq:v4-Delta-convolution} and every retained source-dependent reference
coefficient;
\item the composite/stress entry contains the marked coefficients and all
products/contacts forced by \eqref{eq:v4-bracket-convolution}.
\end{enumerate}
All these entries vanish because they are coordinates of the single zero vector
in \eqref{eq:v4-local-master-zero}; no entry has its own counterterm choice.
\end{theorem}

\begin{proof}
The direct-sum identity \eqref{eq:v4-row-projectors} and
Theorem~\ref{thm:v4-common-local-master} give
\eqref{eq:v4-row-vector-zero}.  The detailed ownership statements follow from
the incidence table and from the source/BV-Laplacian convolution formulas.
They are exhaustive by Theorem~\ref{thm:v4-quadratic-audit}.  Because the row
projectors are applied only after the full coefficient is assembled, a mixed
term is never omitted merely because its two endpoints carry different origin
tags.
\end{proof}

\section{Stress, composite and higher source derivatives of the same identity}
\label{sec:v4-source-descendants}

The source-complete content can be stated without introducing a new family of
Ward identities.  Let $I$ be any central-source multi-index with $|I|\le r$.
Coefficient extraction commutes with the finite local projection because both
act coefficientwise on the truncated source algebra.

\begin{theorem}[All retained source/contact descendants vanish in the same local master vector]
\label{thm:v4-source-descendants}
For $\ell=1,2$ and every $|I|\le r$,
\begin{equation}
 \boxed{
 \left[
 \mathsf P^{\rm ESM}_{N,6,r}
 \mathcal M_{j,\ell}^{\rm all}
 \right]_I=0.}
 \label{eq:v4-source-descendant-zero}
\end{equation}
In particular, stress-tensor insertions, composite insertions, operator-mixing
contacts, determinant/Pfaffian line insertions and coincident source contacts
are coefficients of this one identity.
\end{theorem}

\begin{proof}
Theorem~\ref{thm:v4-common-local-master} is an identity in the finite source
algebra, not merely at $J=0$.  Apply the linear coefficient map $[\cdot]_I$.
Lemma~\ref{lem:v4-source-convolution} shows that the coefficient of the
quadratic antibracket contains every partition of the source labels, while
\eqref{eq:v4-Delta-convolution} contains every source derivative of the finite
BV density/reference operator.  Hence coefficient extraction produces the
complete marked/contact equation and no additional correction is required.
\end{proof}

\begin{corollary}[Second marked derivative]
\label{cor:v4-second-marked}
For two even writers $A,B$, the order-two marked identity contains, among the
terms with the common one-loop action,
\begin{equation}
 (\bar S_{1,0},\bar S_{1,AB})
 +(\bar S_{1,A},\bar S_{1,B})
 -\Delta_{0}\bar S_{1,AB}
 -\Delta_A\bar S_{1,B}
 -\Delta_B\bar S_{1,A}
 -\Delta_{AB}\bar S_{1,0},
 \label{eq:v4-second-marked-explicit}
\end{equation}
together with the correspondingly marked $(S_0,\bar S_2)$ terms.  Thus the
mixed product contact $(\bar S_{1,A},\bar S_{1,B})$ and the source-dependent
measure/reference contacts are mandatory pieces of the same coefficient.
\end{corollary}

\section{The quasilocal complement of the common master coefficient}
\label{sec:v4-master-remainder}

E3 asks for the local master identity together with the already controlled
quasilocal remainder.  Put
\begin{equation}
 \mathsf Q_j:=I-\mathsf P^{\rm ESM}_{N,6,r},
 \qquad
 \mathcal R_{j,\ell}^{\rm M}
 :=\mathsf Q_j\mathcal M_{j,\ell}^{\rm all}.
 \label{eq:v4-master-remainder}
\end{equation}
Let $\varepsilon_j:=\mu/\Lambda_j$ on a physical annulus and
$\varepsilon_j:=a\mu$ on the ultraviolet cap.

\begin{lemma}[Bounded canonical contraction on the fixed rooted source bank]
\label{lem:v4-rooted-bracket}
At fixed $(N,6,r)$ there is a finite constant $C_{N,r}$ such that every
renormalized shell kernel occurring in V4 satisfies the rooted contraction
estimate
\begin{equation}
 \|(F,G)\|_*
 \le C_{N,r}\|F\|_*\|G\|_*,
 \label{eq:v4-rooted-bracket-bound}
\end{equation}
whenever the contraction is one of the canonical contractions already admitted
by the finite bank.  The corresponding renormalized complementary
BV-Laplacian contractions obey the same fixed-bank boundedness after their
local Taylor coefficient is assigned to $\mathsf P$.
\end{lemma}

\begin{proof}
There are finitely many canonical pairs, arities and source markings at fixed
$(N,6,r)$.  In the V1 rooted norm, closing one canonical field--antifield pair
is a finite sum of contractions in which one external/source kernel is placed
in the distinguished $L^1$ root and the remaining marked kernels are in the
corresponding $L^\infty$ stocks.  The weighted Schur estimate bounds each such
contraction with no volume multiplier.  Summing the finite incidence family
produces $C_{N,r}$.  For the BV Laplacian, the potentially local coincidence
piece is by definition assigned to the common Taylor/local bank; the remaining
renormalized kernel is one of the inherited shell-connection/source kernels
and obeys the same rooted estimate.  This statement does not assert an
unsubtracted trace bound for an arbitrary ultraviolet operator.
\end{proof}

\begin{theorem}[One-excess-scale master remainder through two loops]
\label{thm:v4-master-remainder}
On the declared V4 domain,
\begin{equation}
 \boxed{
 \|\mathcal R_{j,1}^{\rm M}\|_*
 \le C_{N,r}\varepsilon_j,
 \qquad
 \|\mathcal R_{j,2}^{\rm M}\|_*
 \le C_{N,r}\varepsilon_j.}
 \label{eq:v4-master-remainder-bound}
\end{equation}
The part of $\mathcal R_{j,2}^{\rm M}$ containing two complementary one-loop
remainders is $O(\varepsilon_j^2)$.
\end{theorem}

\begin{proof}
The one-loop statement is the inherited V1--V3 shell/source remainder estimate
applied to the complete first master coefficient after its local Taylor part is
removed.  At two loops, every primitive hard graph and every determinant-line
remainder has the one-excess-scale bound from the V3 mixed chiral annulus
 theorem and the V2 differentiated line theorem.  Proper one-loop subgraphs are
subtracted by the same corrected $\bar S_1$, so there is no additional
unowned local piece.

Expand the quadratic term using Theorem~\ref{thm:v4-quadratic-audit}.  A term
with one local bounded factor and one complementary $O(\varepsilon_j)$ factor
is $O(\varepsilon_j)$ by Lemma~\ref{lem:v4-rooted-bracket}; a term with two
complementary factors is $O(\varepsilon_j^2)$.  The same argument applies to
the renormalized complementary BV-Laplacian terms.  The finite Hodge restorer
acts only on the local coefficient and is uniformly bounded on the protected
constant-rank chart, so it does not worsen the complementary scale gain.
Summing the finite incidence and forest families proves
\eqref{eq:v4-master-remainder-bound}.
\end{proof}

\begin{corollary}[Summable master defect in the ultraviolet shell composition]
\label{cor:v4-master-summable}
For geometric shell scales $\Lambda_j=b^j\lambda$,
\begin{equation}
 \sum_{j\ge j_0}
 \bigl(\|\mathcal R_{j,1}^{\rm M}\|_*+
       \|\mathcal R_{j,2}^{\rm M}\|_*\bigr)
 <\infty.
 \label{eq:v4-master-sum}
\end{equation}
Hence the E3 master defect left outside the retained local bank is a summable
quasilocal shell remainder at the fixed two-loop/EFT/source order.
\end{corollary}

\begin{proof}
Use Theorem~\ref{thm:v4-master-remainder} and
$\sum_j\mu/\Lambda_j<\infty$; the ultraviolet-cap form $a\mu$ is the terminal
version of the same estimate.
\end{proof}

\section{E3 closure theorem}
\label{sec:v4-E3-close}

\begin{theorem}[Complete common two-loop BV/master/source equation]
\label{thm:v4-E3}
Fix the explicit V2 native Wilson/overlap branch and all V4 finite-order data.
Let
\begin{equation}
 S_j^{[\le2]}
 =S_{0,j}+\hbar\bar S_{1,j}+\hbar^2\bar S_{2,j}
 \label{eq:v4-common-action}
\end{equation}
with $\bar S_{1,j},\bar S_{2,j}$ defined recursively by
\eqref{eq:v4-S1bar} and \eqref{eq:v4-S2bar}.  Then the coefficient of the same
finite master functional
\begin{equation}
 \mathcal M_j(S_j^{[\le2]})
 =\frac12(S_j^{[\le2]},S_j^{[\le2]})
  -\hbar\Delta_jS_j^{[\le2]}
 \label{eq:v4-full-master-functional}
\end{equation}
obeys
\begin{equation}
 \boxed{
 \mathcal M_j(S_j^{[\le2]})
 =\hbar\mathcal R_{j,1}^{\rm M}
  +\hbar^2\mathcal R_{j,2}^{\rm M}
  +O(\hbar^3),}
 \label{eq:v4-master-final}
\end{equation}
where
\begin{equation}
 \boxed{
 \mathsf P^{\rm ESM}_{N,6,r}\mathcal R_{j,\ell}^{\rm M}=0,
 \qquad
 \|\mathcal R_{j,\ell}^{\rm M}\|_*
 \le C_{N,r}\frac{\mu}{\Lambda_j}
 \quad(\ell=1,2),}
 \label{eq:v4-E3-final-bounds}
\end{equation}
with $a\mu$ at the ultraviolet cap.

The single identity includes simultaneously the diffeomorphism,
local-Lorentz/spin, $SU(3)\times SU(2)\times U(1)_Y$, Higgs/Yukawa,
chiral determinant/Pfaffian line, ghost/nonminimal, antifield,
measure/Jacobian/reference, composite/stress and contact rows, together with all
cross terms in $\tfrac12(\bar S_1,\bar S_1)$.  No graph, row or source
derivative uses an independently selected restoring counterterm.
\end{theorem}

\begin{proof}
The coefficient expansion of \eqref{eq:v4-full-master-functional} is exactly
\eqref{eq:v4-M1}--\eqref{eq:v4-M2}.  The local part vanishes by
Theorem~\ref{thm:v4-common-local-master}.  The canonical-pair and source
convolution theorems prove that these coefficients contain all required cross
and contact terms.  The simultaneous row statement is
Theorem~\ref{thm:v4-row-identity}, and the complete source descendants are
Theorem~\ref{thm:v4-source-descendants}.  The complementary estimates are
Theorem~\ref{thm:v4-master-remainder}.  Combining local and complementary
parts yields \eqref{eq:v4-master-final}--\eqref{eq:v4-E3-final-bounds}.
\end{proof}

\begin{center}
\fbox{\parbox{0.90\textwidth}{
\textbf{E3 is closed at the stated two-loop coefficientwise scope.}
The conclusion is one common master/source identity for one common action,
not a collection of sectorwise finite Ward identities.}}
\end{center}

\section{Why the next theorem is now an induction theorem}
\label{sec:v4-E4-handoff}

There is no load-bearing reason to compute an isolated three-loop graph next.
The common two-loop coefficient reveals the general algebraic pattern.  If
\begin{equation}
 S^{[\le L]}=\sum_{q=0}^{L}\hbar^q\bar S_q,
 \label{eq:v4-L-action}
\end{equation}
then the order-$L$ master coefficient is
\begin{equation}
 \boxed{
 \mathcal M^{(L)}
 =(S_0,\bar S_L)
 +\frac12\sum_{p=1}^{L-1}(\bar S_p,\bar S_{L-p})
 -\Delta\bar S_{L-1}.}
 \label{eq:v4-general-master-coefficient}
\end{equation}
The source-contact formula of Lemma~\ref{lem:v4-source-convolution} and the
canonical-pair incidence theorem are independent of $L$.  The new work for E4
is therefore not algebraic expansion but the native all-forest estimate and the
finite local bank/range statement at an arbitrary \emph{fixed} $L$.

\begin{theorem}[Exact next load-bearing theorem for V5]
\label{thm:v4-next-E4}
For every prescribed fixed loop order $L$, EFT/derivative cap $\Delta$, source
order $r$, arity $N$ and finite external panel, construct an induction state
containing:
\begin{enumerate}[label=\textup{(I\arabic*)},leftmargin=3.2em]
\item the complete finite ghost-zero EFT/Wilson bank through the order required
by $L$;
\item the complete source/contact algebra through order $r$ and all exact BV
field--antifield/source pairs;
\item determinant/Pfaffian line/frame data on the protected component;
\item the already fixed common actions $\bar S_0,\ldots,\bar S_{L-1}$ and their
finite normalization/reference data;
\item ownership of every proper divergent subgraph by the restriction of the
appropriate previously fixed common $\bar S_q$;
\item rooted/quasilocal complementary kernels and the retained soft
constraint/nonminimal variables.
\end{enumerate}
Then prove that the forest-prepared order-$L$ coefficient decomposes as
\begin{equation}
 \text{complete local Taylor/EFT/source jet}
 +\text{quasilocal remainder},
 \label{eq:v4-E4-decomposition}
\end{equation}
with a finite local bank at that fixed order, that the first un-restored local
master coefficient obtained from
\eqref{eq:v4-general-master-coefficient} is a $d$-cocycle, that its complete
Standard-Model anomaly projection vanishes, and that one common finite Hodge
primitive restores it.  Finally prove a one-excess-scale rooted bound for the
complete order-$L$ forest remainder, with constants allowed to depend on $L$.
This yields a native shell self-map at arbitrary fixed $L$ without claiming any
uniform $L\to\infty$ estimate.
\end{theorem}

\section{V4 anti-circularity audit}
\label{sec:v4-audit}

\begin{enumerate}[leftmargin=2.4em]
\item The quadratic two-loop term is not inferred from separate sector Ward
identities.  It is expanded exactly by source partition, origin tag and
canonical pair in \eqref{eq:v4-full-quadratic}.
\item Source-dependent measure/reference terms are not suppressed.  Their
coefficients are the $\Delta_{j,K}$ terms of
\eqref{eq:v4-Delta-convolution}.
\item The order-two breaking is formed only after the one-loop restorer is in
the action.  Equation \eqref{eq:v4-C1-descendants} displays the induced
$(S_1,C_1)$, $(C_1,C_1)$, $\Delta C_1$ and forest-insertion terms explicitly.
\item No row or graph receives an independent restorer.  Proposition
\ref{prop:v4-no-hidden-refit} identifies the only remaining ambiguity as the
already-declared exact/redundant plus physical ghost-zero Wilson freedom.
\item The line/frame sector is not replaced by a scalar determinant.  It enters
through the V1--V3 protected line packet and through the same background
canonical pairs as the remaining action.
\item The quasilocal master defect is not promoted to an exact finite-cutoff
zero outside the retained bank.  It is recorded with the proved summable
$O(\mu/\Lambda_j)$ rooted bound.
\end{enumerate}

\section{V4 status ledger}
\label{sec:v4-status}

\subsection*{Proved in V4}
\begin{enumerate}[leftmargin=2.2em]
\item Central source writers can be separated from genuinely canonical source
multiplets without changing the V3 source-extended BV bank.  In that
organization, the exact source coefficient of an antibracket is the finite
super-binomial convolution \eqref{eq:v4-bracket-convolution}, and a
source-dependent BV density obeys \eqref{eq:v4-Delta-convolution}.
\item The canonical-pair support theorem gives an exhaustive finite incidence
audit of every mixed term in $\tfrac12(\bar S_1,\bar S_1)$.  Geometry--matter
stress, gauge--Higgs/fermion, Yukawa, ghost/antifield, line and source contacts
are entries of one quadratic master coefficient rather than separate Ward
problems.
\item The corrected one-loop action necessarily feeds the order-two master row
through its forest insertions, $(S_1,C_1)$, $\tfrac12(C_1,C_1)$ and
$-\Delta C_1$ descendants.  These terms are all retained in the recursive
$B_{2,j}$.
\item With the V3 common Hodge restorer, the complete local coefficients obey
one common identity
$\mathsf P\mathcal M_{j,1}^{\rm all}=
 \mathsf P\mathcal M_{j,2}^{\rm all}=0$ using the same
$S_0,\bar S_1,\bar S_2$.
\item Every E3 row is a coordinate projection of that one zero vector.  The
same is true after every retained source derivative, including stress,
composite, operator-mixing, determinant/Pfaffian and coincident contact
insertions.
\item The complementary first- and second-loop master coefficients retain the
one-excess-scale rooted bound $O(\mu/\Lambda_j)$, with the term quadratic in two
one-loop remainders $O((\mu/\Lambda_j)^2)$.  The ultraviolet shell sum is
finite at this fixed order.
\item Therefore E3 is closed on the explicit V2 branch at the stated
coefficientwise two-loop/EFT/source scope.
\end{enumerate}

\subsection*{Inherited}
\begin{enumerate}[leftmargin=2.2em]
\item The V80 bosonic Einstein--gauge--Higgs two-loop self-map, weight-six
finite BV range theorem, common one-loop forest ownership, retained soft
constraint architecture and rooted two-loop remainder estimate.
\item The three-generation Standard-Model representation, cancellation of the
five ordinary local anomaly trace sums, even $SU(2)$ doublet parity, finite
gauge descent of the complete nonzero chiral section and the exact
Schur/Pfaffian scale cocycles.
\item V1's overlap functional calculus, local Kato frame, hard pivot and rooted
contraction bounds.
\item V2's explicit native Wilson completion, all-zone gap/no-doubler result,
physical annular/full-zone comparison, fixed source-jet calculus and
differentiated line remainder.
\item V3's source-extended ESM complex, native anomaly-quotient zero, common
finite Hodge restorer, mixed chiral two-loop annulus theorem and complete E2
shell self-map.
\end{enumerate}

\subsection*{Conditional}
\begin{enumerate}[leftmargin=2.2em]
\item The theorem remains native to the explicit V2 Wilson/overlap completion;
a different microscopic fermion kernel requires the regulator-comparison gate
already isolated in V2.
\item The Hodge bound uses the compact constant-rank component of the finite BV
bank.  A rank crossing requires a new chart or a separate analysis.
\item Determinant/Pfaffian statements remain on the protected nonzero component
with the retained positive IR/reference prescription.  No global line holonomy
or zero crossing is claimed.
\item The complementary master estimate uses the inherited fixed-order native
hard-symbol, forest and rooted-source bounds.  Constants may depend on loop
order, EFT degree, source order, arity and the protected chart.
\end{enumerate}

\subsection*{Open}
\begin{enumerate}[leftmargin=2.2em]
\item E4: prove the arbitrary-fixed-loop native forest induction rather than
compute isolated higher-loop examples.  The decisive new analytic task is an
all-forest one-excess-scale theorem at each prescribed fixed $L$, with one
finite local Taylor/EFT bank and common lower-order counterterm ownership.
\item E5: after E4, compose arbitrarily many ultraviolet shells at fixed
$(L,\Delta,r)$ and prove Cauchy/unique cutoff limits for normalized connected
source coefficients, including operator mixing, stress/contact and line
insertions.
\item E6: prove finite-scheme/reference compatibility and order compatibility
between the $L$ and $L+1$ constructions.
\end{enumerate}

\subsection*{Exact next load-bearing theorem}
\begin{center}
\fbox{\parbox{0.92\textwidth}{
\textbf{E4: arbitrary fixed-loop native shell induction.}
For a prescribed fixed $L$, assemble the complete forest using the same
previously fixed actions $\bar S_0,\ldots,\bar S_{L-1}$, prove a finite local
Taylor/EFT/source decomposition plus a one-excess-scale rooted remainder, form
the common order-$L$ coefficient
\(
(S_0,\bar S_L)+\frac12\sum_{p=1}^{L-1}
(\bar S_p,\bar S_{L-p})-\Delta\bar S_{L-1}
\), prove its anomaly-free local cocycle lies in the common BV range, and
install one source-complete restoring primitive.  Constants may depend on $L$;
no uniform perturbative convergence claim is allowed.}}
\end{center}

\section*{Append-only continuation marker: V4}
\addcontentsline{toc}{section}{Append-only continuation marker: V4}
\textbf{End of V4.}  The complete V1--V3 body above is retained verbatim.  V4
closes E3 by writing and auditing one common first- and second-loop master
coefficient, including every canonical-pair cross term, every retained source
contact and every source-dependent measure/reference term.  Future versions
must append after this marker.  The next file is to be named
\texttt{Renewal\_Perturbative\_ESM\_Continuum\_Research\_Notes\_V5.tex}.

% =============================================================================
% V5 -- arbitrary fixed-loop native forest induction
% =============================================================================
\part*{V5: Arbitrary fixed-loop native forest induction and contraction-complete BV jet packets}
\addcontentsline{toc}{part}{V5: Arbitrary fixed-loop native forest induction and contraction-complete BV jet packets}
\label{sec:v5-start}

\section{Upstream correction: the induction state must remember future BV contractions}
\label{sec:v5-upstream}

V4 closed the common local master equation through two loops and retained a
quasilocal complementary master kernel of size $O(\mu/\Lambda_j)$.  That is
sufficient for the two-loop statement, but it is not by itself an induction
state.  At loop order three and above a BV Laplacian or a canonical
field--antifield contraction can act on a lower-order quasilocal kernel and
produce a new lower-arity local contact.  Therefore the implication
\[
 \text{``ordinary Taylor remainder''}
 \Longrightarrow
 \text{``irrelevant for every later local master coefficient''}
\]
is false without additional bookkeeping.

This is the first load-bearing issue in E4.  The correct upstream variable is
not a bare remainder $R$ but a \emph{contraction-complete jet packet}: the
remainder together with the local Taylor jets of every canonical/source
contraction which can still occur before the prescribed final loop order.  At
fixed loop order, derivative/EFT order, source order and arity there are only
finitely many such descendants.  The packet is therefore finite, while no
claim of uniformity as $L\to\infty$ is required.

Throughout V5 fix:
\begin{equation}
 L\in\mathbb N,\qquad
 \Delta<\infty,\qquad r<\infty,\qquad N<\infty,
 \qquad |p_{\rm ext}|\le\mu<c_0\Lambda_j,
 \label{eq:v5-fixed-data}
\end{equation}
and write
\begin{equation}
 \epsilon_j:=\frac{\mu}{\Lambda_j}
 \label{eq:v5-epsilon}
\end{equation}
on a physical annulus, with the terminal ultraviolet-cap replacement
$\epsilon_j=a\mu$ exactly as in V2--V4.  Constants in V5 may depend on
$(L,\Delta,r,N)$, on the protected chart and on the finite representation
packet.  They are not required to remain bounded with $L$.

\begin{firewall}[No naive exact-cocycle induction from V4 alone]
\label{fw:v5-no-naive-cocycle}
V3's exact statement that the first un-restored master coefficient is a
classical BV cocycle assumes that every lower master coefficient vanishes in
the algebra in which the quantum consistency identity is expanded.  V4 proves
an exact vanishing of the retained local projection but keeps a nonzero
quasilocal complement.  Hence one may not simply cite the V3 cocycle theorem
at arbitrary loop order while discarding the V4 complement.  V5 repairs this
by enlarging the induction state so that all local descendants of that
complement are retained before the next projection is formed.
\end{firewall}

\section{The fixed-order induction state}
\label{sec:v5-state}

Let $\widehat{\mathfrak L}^{\rm ESM}_{N,\Delta,r}$ be the complete finite local
ESM BV/source bank at arity at most $N$, EFT/derivative filtration at most
$\Delta$ and source order at most $r$.  It contains the same canonical
field--antifield pairs, determinant/Pfaffian line coordinates,
measure/reference coordinates, stress/composite writers and nonminimal pairs
as V3--V4, enlarged only as required by the chosen finite EFT filtration.
The word ``complete'' means that every local descendant under the classical
BV differential and every retained source/contact derivative is included.

The gravity sector is treated as an EFT.  Thus $\Delta$ is a filtration label,
not a statement that a finite list of gravitational couplings renormalizes all
loop orders.  When an order-$L$ graph requires a new local operator at the
chosen accuracy, that coordinate is added to the finite order-$L$ bank.  No
operator needed only at unbounded loop or derivative order is introduced.

\begin{definition}[Fixed-$L$ induction datum]
\label{def:v5-induction-datum}
For $0\le q<L$, the order-$q$ datum consists of
\begin{enumerate}[label=\textup{(D\arabic*)},leftmargin=3.2em]
\item the common normalized local action coefficient $\bar S_{q,j}$, including
its physical EFT/Wilson coordinates and every source/antifield descendant;
\item the determinant/Pfaffian line/frame packet on the protected nonzero
component;
\item the normalization and reference coordinates which define the common
finite scheme;
\item every rooted/quasilocal action kernel generated at order $q$;
\item the contraction-complete jet packet of each such quasilocal kernel up to
remaining contraction depth $L-q$;
\item the ancestry/ownership tag identifying the lower-order common action
which owns each proper local subgraph insertion.
\end{enumerate}
The collection of these data for $0\le q<L$ is denoted
$\mathfrak I_{<L}(\Delta,r,N)$.
\end{definition}

The fifth entry is new.  It is the device which makes the induction exact at
the level of every local coefficient that can be generated before loop order
$L$.

\section{Reference absorption and finiteness of the graph family}
\label{sec:v5-graph-finite}

At fixed loop order one must first eliminate a trivial source of infinitude:
arbitrarily many quadratic mass/kinetic insertions on one propagator.  The
native shell formalism already has a common normalized Gaussian/reference
family, so these insertions should not be counted as independent interaction
vertices.

\begin{lemma}[Quadratic/reference absorption]
\label{lem:v5-quadratic-absorption}
At every completed order $q<L$, all source-free local hard vertices of hard
valence $0,1,2$ may be assigned uniquely to the normalized background and
Gaussian/reference coordinates of the common state.  After this assignment,
every source-free perturbative hard interaction vertex has hard valence at
least three.  Source-marked vertices of hard valence at most two remain
explicit, but their number in a coefficient of source order at most $r$ is at
most $r$.
\end{lemma}

\begin{proof}
The vacuum coordinate is part of the reference normalization, the one-point
coordinate is fixed by the chosen background/tadpole condition, and the full
local two-point coordinate is part of the common quadratic block whose inverse
defines the native hard covariance.  Moving any of these coordinates between
``propagator'' and ``interaction'' without changing the normalization would
only be a redundant re-expansion of the same finite Gaussian identity.  We
choose the nonredundant convention in which they are absorbed once.  A marked
source insertion is not absorbed because it is a requested observable; every
such insertion consumes at least one of the at most $r$ source labels.  The
claim follows.
\end{proof}

\begin{proposition}[Finite graph family at fixed $(L,\Delta,r,N)$]
\label{prop:v5-finite-graphs}
For a fixed external arity at most $N$, source order at most $r$, loop number at
most $L$, and finite EFT filtration, only finitely many native hard graph
topologies, field/representation assignments, source decorations and
counterterm ancestry assignments can occur in one shell step.
\end{proposition}

\begin{proof}
Ignore first the at most $r$ marked low-valence source vertices.  If $V$ is the
number of remaining hard interaction vertices, $I$ the number of internal hard
lines and $E\le N+O(r)$ the number of external hard half-edges, Lemma
\ref{lem:v5-quadratic-absorption} gives
\[
 3V\le 2I+E.
\]
For each connected component contributing to an effective vertex,
$L_\Gamma=I-V+1$.  Hence
\[
 V\le 2L_\Gamma+E-2.
\]
Thus $V$ and $I$ are bounded in terms of the fixed data.  The local bank has
only finitely many field species, representation tensors, derivative monomials
and source labels at fixed $(\Delta,r,N)$.  A graph with bounded numbers of
vertices and half-edges therefore admits only finitely many assignments from
that bank.  Finally every proper-subgraph ancestry label lies in the finite set
$\{0,\ldots,L-1\}$.  This proves finiteness.
\end{proof}

\begin{lemma}[A proper 1PI ultraviolet subgraph has smaller loop number]
\label{lem:v5-loop-drop}
Let $\Gamma$ be a connected 1PI hard graph of loop number $L_\Gamma$ and let
$\gamma\subsetneq\Gamma$ be a proper connected 1PI subgraph which is admissible
as a local ultraviolet subtraction block.  Then
\begin{equation}
 \boxed{L_\gamma<L_\Gamma.}
 \label{eq:v5-loop-drop}
\end{equation}
\end{lemma}

\begin{proof}
Contract $\gamma$ to a single vertex.  The loop-number identity for connected
graphs gives
\[
 L_\Gamma=L_\gamma+L_{\Gamma/\gamma}.
\]
If equality $L_\gamma=L_\Gamma$ held, the quotient would be a connected tree.
Because $\gamma$ is proper, the quotient would contain at least one edge.  Any
edge of a tree is a bridge, and its preimage would be an internal bridge of
$\Gamma$, contradicting 1PI.  Hence $L_{\Gamma/\gamma}\ge1$ and
$L_\gamma<L_\Gamma$.
\end{proof}

This elementary lemma is what makes common counterterm ownership inductive:
no proper 1PI subtraction at loop $L$ is allowed to invent an order-$L$
finite normalization of its own.

\section{Decorated native forests and common lower-order ownership}
\label{sec:v5-forests}

\begin{definition}[Admissible decorated forest]
\label{def:v5-forest}
For a fixed native graph $\Gamma$, an admissible forest $F$ is a set of proper
connected 1PI local subtraction blocks such that any two elements are either
disjoint or nested.  Each $\gamma\in F$ carries:
\begin{enumerate}[label=\textup{(F\arabic*)},leftmargin=3.2em]
\item its field/representation and orientation data;
\item its source/contact multi-index;
\item its determinant/Pfaffian line/frame tag if present;
\item its loop number $L_\gamma$;
\item the unique local insertion
$C_\gamma=C_\gamma[\bar S_{L_\gamma,j}]$ obtained by restricting the already
fixed common action $\bar S_{L_\gamma,j}$ to the same local tensor/source
coordinate.
\end{enumerate}
Overlapping but nonnested subgraphs never occur in the same forest; they occur
in distinct forest terms, as in the V61/V80 two-loop overlapping-channel
analysis.
\end{definition}

Let $\Gamma/F$ denote simultaneous contraction of the mutually compatible
blocks in $F$.  The scheme-specific forest preparation is
\begin{equation}
 \mathbf R^{<L}_j\mathcal I_\Gamma
 :=\sum_{F\in\mathfrak F(\Gamma)}
 \mathcal I_{\Gamma/F,j}\star
 \prod_{\gamma\in F} C_\gamma,
 \label{eq:v5-forest-preparation}
\end{equation}
where the empty forest contributes the unprepared native amplitude and
$\star$ means insertion with the original routing, symmetry factor, source
labels and line/frame orientation.  Equation \eqref{eq:v5-forest-preparation}
is a finite algebraic identity at fixed cutoff; no continuum BPHZ theorem is
being imported.

\begin{theorem}[Common ownership of every proper subtraction]
\label{thm:v5-common-ownership}
For every order-$L$ graph in Proposition~\ref{prop:v5-finite-graphs}, every
proper local ultraviolet subtraction in
\eqref{eq:v5-forest-preparation} is owned by one of the previously fixed common
actions $\bar S_{q,j}$ with $q<L$.  In particular:
\begin{enumerate}[label=\textup{(O\arabic*)},leftmargin=3.2em]
\item no proper subgraph is independently refit graph by graph;
\item nested counterterms use the restriction of the already corrected lower
order action, including its earlier restoring terms;
\item all source/contact and line descendants of a counterterm are the
descendants of that same common coordinate;
\item changing the finite part of one $C_\gamma$ is a change of the declared
common renormalization scheme, not a harmless graphwise convention.
\end{enumerate}
\end{theorem}

\begin{proof}
By Lemma~\ref{lem:v5-loop-drop}, every proper 1PI subtraction block has loop
number strictly below $L$.  The induction datum therefore contains its common
local coefficient.  The source and line descendants are included in that
datum before restriction.  Forest compatibility ensures that nested
insertions are evaluated recursively from the innermost block outward, so a
counterterm already containing a lower restoring primitive is never replaced
by its uncorrected predecessor.  The uniqueness assertion is exactly the
fixed normalization/reference convention of the induction datum.
\end{proof}

\section{Contraction-complete jet packets}
\label{sec:v5-jet-packets}

We now formalize the upstream state enlargement.

Let $\mathcal C$ be the finite alphabet of elementary operations which can
appear in a later master/source coefficient at the fixed data
\eqref{eq:v5-fixed-data}.  Its letters are:
\begin{enumerate}[label=\textup{(C\arabic*)},leftmargin=3.2em]
\item one canonical field--antifield contraction appearing in an antibracket;
\item one self-contraction appearing in the finite BV Laplacian;
\item one retained source derivative or contact identification;
\item one determinant/Pfaffian line derivative/contact contraction.
\end{enumerate}
Each operation carries the precise signs and source labels of the V4
canonical-pair/source convolution formulas.  An admissible word
$\alpha=c_m\cdots c_1$ is one whose successive arities and ghost numbers are
defined.  Write $D_\alpha$ for the corresponding operator.

\begin{definition}[Depth-$m$ contraction-complete packet]
\label{def:v5-packet}
For a rooted/quasilocal kernel $K$, define
\begin{equation}
 \boxed{
 \mathbf J_j^{[m]}K
 :=\left(
   K,\ \bigl\{
   \mathsf T_{\Delta_\alpha,r_\alpha}
   D_\alpha K:
   0\le |\alpha|\le m
   \bigr\}
 \right).}
 \label{eq:v5-packet}
\end{equation}
Here $\mathsf T_{\Delta_\alpha,r_\alpha}$ is the complete local
Taylor/EFT/source projector appropriate to the descendant after the operations
in $\alpha$.  The empty word records the current local jet of $K$.  The
unprojected kernel $K$ is retained as the first component; the packet is not a
replacement of the quasilocal action by its Taylor series.
\end{definition}

\begin{lemma}[Finiteness of the packet]
\label{lem:v5-packet-finite}
For fixed $(m,\Delta,r,N)$, the packet
$\mathbf J_j^{[m]}K$ has finitely many local coordinates.
\end{lemma}

\begin{proof}
The alphabet $\mathcal C$ is finite because the field/antifield set, source
panel and line packet are finite.  There are finitely many words of length at
most $m$.  Every contraction lowers canonical arity or consumes a source/contact
slot, so inadmissible words terminate even earlier.  Each descendant local
projector has finite arity, source order and derivative/EFT degree; hence its
range is finite dimensional.  The finite union of these ranges is finite.
\end{proof}

\begin{lemma}[Shift identity]
\label{lem:v5-packet-shift}
Let $c\in\mathcal C$ be one admissible elementary operation and $m\ge1$.
Then every local coordinate of
$\mathbf J_j^{[m-1]}(D_cK)$ occurs verbatim as a coordinate of
$\mathbf J_j^{[m]}K$:
\begin{equation}
 \boxed{
 \mathsf T D_\alpha D_cK
 =\mathsf T D_{\alpha c}K,
 \qquad |\alpha|\le m-1.}
 \label{eq:v5-packet-shift}
\end{equation}
The equality includes the V4 graded source-partition signs and the declared
measure/reference contact convention.
\end{lemma}

\begin{proof}
This is associativity of the corresponding finite canonical/source
contractions, with the same fixed ordering convention used to define the V4
source convolution.  No limit is involved.  The descendant projector on the
left is exactly the projector attached to the composite word $\alpha c$ on the
right.  Because source-dependent reference terms are part of the finite BV
operator before differentiation, their contact descendants are present on
both sides.
\end{proof}

\begin{theorem}[Future-local-feedback completeness]
\label{thm:v5-feedback-complete}
Suppose an order-$q$ quasilocal kernel is stored with packet depth $L-q$.
Then every local term which that kernel can generate by BV, source/contact or
line contractions in any coefficient of loop order at most $L$ is already a
coordinate of its stored packet.  No later local master coefficient can receive
an unrecorded contribution from that kernel.
\end{theorem}

\begin{proof}
A contribution at loop order at most $L$ can contain at most $L-q$ additional
loop-generating self-contractions along an ancestry descending from an
order-$q$ kernel; antibracket and source/contact operations do not increase the
required depth beyond the number of remaining coefficient-building steps.
Read those operations from earliest to latest as a word $\alpha$ of length at
most $L-q$.  Its local descendant is one of the coordinates in
\eqref{eq:v5-packet}.  Lemma~\ref{lem:v5-packet-shift} shows inductively that
a later operation merely moves to a shallower coordinate subpacket.  Thus no
local feedback is lost.
\end{proof}

\begin{remark}[Why this is stronger than an ordinary Taylor remainder]
A kernel may have zero current Taylor jet and nevertheless acquire a nonzero
local descendant after a BV self-contraction.  Definition
\ref{def:v5-packet} records that descendant in advance.  This is exactly the
higher-loop generalization of V4's rule that one must assemble all contact and
measure/reference terms before projecting the master coefficient.
\end{remark}

\section{Rooted one-excess-scale calculus at fixed loop order}
\label{sec:v5-rooted-calculus}

The packet solves the bookkeeping problem.  We now prove the analytic estimate
which replaces an imported all-order continuum forest theorem.

For a kernel $K$ with current local jet $\mathsf T K$, write
\begin{equation}
 K=K^{\loc}+K^{\remn},
 \qquad K^{\loc}:=\mathsf T K,
 \qquad K^{\remn}:=(I-\mathsf T)K.
 \label{eq:v5-local-rem}
\end{equation}
Use the V1--V4 rooted norm $\|\cdot\|_*$ for quasilocal kernels and any fixed
coefficient norm on the finite local bank.

\begin{lemma}[Primitive annular Taylor gain at arbitrary fixed graph size]
\label{lem:v5-primitive-gain}
Let $\Gamma$ be one of the finitely many native hard graphs in Proposition
\ref{prop:v5-finite-graphs}.  Suppose all proper local subgraphs have already
been replaced by their owned lower-order insertions.  Let
$\omega_\Gamma$ be the finite Taylor degree required by the chosen EFT
filtration for the overall local part.  Then
\begin{equation}
 \boxed{
 \bigl\|(I-\mathsf T_{\omega_\Gamma})
 \mathbf R_j^{<L}\mathcal I_\Gamma\bigr\|_*
 \le C_{\Gamma,\Delta,r,N}\,\epsilon_j.}
 \label{eq:v5-primitive-gain}
\end{equation}
The same estimate holds after every admissible fixed source, geometric, gauge,
Higgs/Yukawa, antifield or line derivative needed by the packet.
\end{lemma}

\begin{proof}
After the proper local blocks have been replaced, every integrated native
propagator is supported on the hard annulus at scale $\Lambda_j$ and obeys the
fixed differentiated symbol/rooted bounds proved for the bosonic packets in
V80 and for the chiral packet in V1--V3.  The graph contains finitely many
lines and vertices by Proposition~\ref{prop:v5-finite-graphs}.  Taylor's
formula with integral remainder in the finite external momentum/source window
gives at least one explicit factor
$|p_{\rm ext}|/\Lambda_j\le\epsilon_j$ after subtraction through the complete
required local degree.  Every further fixed external derivative differentiates
a finite product of native symbols and is controlled by the same annular
bounds.  Fourier integration by parts an arbitrary fixed number of times gives
the weighted rooted $L^1$ moment required by $\|\cdot\|_*$.  The number of
Leibniz terms is finite and absorbed into $C_{\Gamma,\Delta,r,N}$.  At the
terminal full-zone cap the V2 finite-to-ordinary estimate gives the replacement
$a\mu$.
\end{proof}

\begin{lemma}[One-excess rooted ideal]
\label{lem:v5-rooted-ideal}
At fixed $(L,\Delta,r,N)$ the class of kernels satisfying
$\|R\|_*\le C\epsilon_j$ is stable, with a changed fixed constant, under every
allowed rooted contraction with bounded local kernels.  More generally, if a
contracted expression contains $s$ factors $R_1,\ldots,R_s$ with
$\|R_i\|_*\le C_i\epsilon_j$, then its complementary rooted part is
$O(\epsilon_j^s)$ after all newly created local jets are extracted into the
contraction-complete packet.
\end{lemma}

\begin{proof}
V4's canonical contraction estimate bounds one allowed contraction by a fixed
constant times the product of the rooted norms.  The same weighted Schur/root
argument applies to the additional finite arities present at fixed $L$.
Whenever a contraction creates a local Taylor/contact piece, Definition
\ref{def:v5-packet} assigns that piece to the appropriate local descendant
coordinate before the complement is estimated.  The remaining kernel is
therefore bounded by the product of the complementary norms.  Iterating the
finite number of contractions gives the claim.
\end{proof}

\begin{theorem}[Native fixed-$L$ forest decomposition]
\label{thm:v5-forest-decomposition}
For every graph/source coefficient in the fixed data
\eqref{eq:v5-fixed-data}, the scheme-specific forest preparation
\eqref{eq:v5-forest-preparation} has a decomposition
\begin{equation}
 \boxed{
 \mathbf R_j^{<L}\mathcal I_\Gamma
 =\mathsf L_{\Gamma,j}^{(L)}
  +\mathsf Q_{\Gamma,j}^{(L)},}
 \label{eq:v5-forest-decomp}
\end{equation}
where:
\begin{enumerate}[label=\textup{(A\arabic*)},leftmargin=3.2em]
\item $\mathsf L_{\Gamma,j}^{(L)}$ lies in a finite local
Taylor/EFT/source/contact bank determined by $(L,\Delta,r,N)$;
\item every proper subtraction entering $\mathsf L_{\Gamma,j}^{(L)}$ is owned
by the previously fixed common action of Theorem
\ref{thm:v5-common-ownership};
\item the full contraction-complete packet of
$\mathsf Q_{\Gamma,j}^{(L)}$ is retained in the induction state;
\item its current complementary kernel obeys
\begin{equation}
 \boxed{
 \|\mathsf Q_{\Gamma,j}^{(L)}\|_*
 \le C_{L,\Delta,r,N}\epsilon_j;}
 \label{eq:v5-forest-Q-bound}
\end{equation}
\item every term containing $s$ already complementary lower-order ancestors is
$O(\epsilon_j^s)$ after its newly generated local descendants are extracted.
\end{enumerate}
The theorem includes nested and overlapping proper subgraphs and all retained
source/contact markings.
\end{theorem}

\begin{proof}
The forest family is finite by Proposition~\ref{prop:v5-finite-graphs}.
Proceed by induction on the number of proper subtraction blocks.  For no proper
block, Lemma~\ref{lem:v5-primitive-gain} gives the result.  Suppose it holds for
all graphs with fewer proper blocks.  In each forest term choose its maximal
proper elements.  By Lemma~\ref{lem:v5-loop-drop} their loop orders are lower,
so their local insertions and quasilocal packets belong to the induction datum.
Replace each maximal block by its local coordinate plus its retained
complement.  The term in which all maximal blocks are local reduces to a
primitive quotient graph and is controlled by Lemma
\ref{lem:v5-primitive-gain}.  Every other term contains at least one
complementary factor and is $O(\epsilon_j)$ by Lemma
\ref{lem:v5-rooted-ideal}; terms with $s$ such factors have the stated higher
power.

Overlapping blocks do not require simultaneous subtraction: incompatible
blocks occur in distinct forest terms, exactly as dictated by Definition
\ref{def:v5-forest}.  Nested blocks are handled from inside outward, so their
source/contact labels and finite parts are inherited rather than refit.
Differentiating the finite forest sum distributes the source labels over graph,
subgraphs and counterterm insertions.  The resulting super-partitions are
precisely the V4 contact convolution, and every local descendant is stored by
Definition~\ref{def:v5-packet}.  Summing the finite forest family proves the
uniform fixed-$L$ bound.
\end{proof}

\begin{corollary}[No imported all-order forest theorem is needed]
\label{cor:v5-no-imported-forest}
At every prescribed fixed $L$, the native regulator itself supplies the
complete local-plus-quasilocal forest decomposition needed for the shell
update.  The argument uses the already proved native annular propagator/source
bounds, finite graph combinatorics, common lower-order ownership and the rooted
ideal estimate.  It does not assert that the constants are uniform in $L$ and
does not invoke a continuum forest theorem whose hypotheses have not been
matched to the native regulator.
\end{corollary}

\section{Finite closure of the local EFT/source bank}
\label{sec:v5-local-bank}

\begin{theorem}[Finite order-$L$ local closure bank]
\label{thm:v5-finite-local-bank}
For fixed $(L,\Delta,r,N)$ there is a finite-dimensional local bank
\begin{equation}
 \widehat{\mathfrak L}^{\rm ESM}_{L;N,\Delta,r}
 \label{eq:v5-L-bank}
\end{equation}
which contains:
\begin{enumerate}[label=\textup{(B\arabic*)},leftmargin=3.2em]
\item every local coefficient in Theorem~\ref{thm:v5-forest-decomposition};
\item all BV field--antifield descendants needed by the order-$L$ master
coefficient;
\item all source/contact and stress/composite descendants through order $r$;
\item all determinant/Pfaffian line descendants on the protected component;
\item all local descendant jets in the depth-$L$ contraction-complete packets;
\item the independent ghost-zero EFT/Wilson quotient at the chosen order.
\end{enumerate}
No finite-dimensional bank independent of $L$ or $\Delta$ is asserted.
\end{theorem}

\begin{proof}
The graph family is finite by Proposition~\ref{prop:v5-finite-graphs}.  Every
graph has finitely many forest terms and finitely many local Taylor
coefficients through its prescribed degree.  Lemma~\ref{lem:v5-packet-finite}
adds only finitely many contraction descendants.  The field and representation
packets are finite, while arity and source order are capped.  Taking the span
of this finite union gives \eqref{eq:v5-L-bank}.  The independent ghost-zero
cohomology is retained rather than quotienting it away, exactly as in V78--V80.
\end{proof}

\section{The local BV complex at arbitrary finite filtration}
\label{sec:v5-complex}

Let
\begin{equation}
 d_{j,L}:=
 \mathsf P_{L;N,\Delta,r}
 (S_{0,j},\,\cdot\,)
 \mathsf P_{L;N,\Delta,r}
 \label{eq:v5-differential}
\end{equation}
be the classical BV differential on the bank of Theorem
\ref{thm:v5-finite-local-bank}.

\begin{theorem}[Finite filtered ESM complex and ghost-one quotient]
\label{thm:v5-complex}
For every fixed $(L,\Delta,r,N)$,
\begin{equation}
 \boxed{d_{j,L}^2=0.}
 \label{eq:v5-d2}
\end{equation}
The nonminimal pairs remain contractible.  On the charged-Higgs branch, after
removal of the exceptional free-abelian families as in V80, the nontrivial
four-dimensional ghost-number-one local quotient relevant to the complete ESM
packet is the same standard gauge-anomaly quotient
\begin{equation}
 \mathfrak A_{\rm std}
 =\Span\{\mathcal A_{333},\mathcal A_{Y33},\mathcal A_{Y22},
              \mathcal A_{YYY},\mathcal A_{YRR}\}.
 \label{eq:v5-anomaly-space}
\end{equation}
Restriction to the finite order-$L$ bank creates no additional harmonic class.
\end{theorem}

\begin{proof}
Before projection, the complete classical ESM BV action satisfies the
classical master equation on the declared protected branch, so its Hamiltonian
vector field squares to zero.  The bank in Theorem
\ref{thm:v5-finite-local-bank} was defined to contain every retained descendant
under that field.  Hence its projector commutes with the differential and
\eqref{eq:v5-d2} follows.

For ghost number one, the ordinary local cohomology input is the same
matter-inclusive Einstein--Yang--Mills classification already matched in V79,
V80 and V3.  Hypercharge is not a free abelian generator because it acts on the
charged Higgs, so the additional free-abelian current/antifield families are
absent.  The exceptional gravitational/topological anomaly dimensions do not
include four.  Adding finitely many higher-derivative local representatives,
source writers and contractible BV pairs changes the finite representative
space but not the underlying four-dimensional anomaly quotient.  The finite
bank is a differential-stable restriction containing all descendants at the
chosen filtration, so the finite-jet extension argument of V80 identifies its
harmonic ghost-one quotient with the restriction of
\eqref{eq:v5-anomaly-space}.  No statement about global gauge loops or line
zero crossings is used.
\end{proof}

\begin{theorem}[Exact vanishing of the native internal-gauge harmonic projection at every fixed loop order]
\label{thm:v5-anomaly-zero}
For the complete three-generation Standard-Model shell on the explicit V2
Wilson/overlap branch, every coefficient produced by the finite forest
preparation and its contraction-complete descendants has zero projection onto
\eqref{eq:v5-anomaly-space}:
\begin{equation}
 \boxed{
 \mathsf A_{\rm std}
 \mathcal B_{j,L}^{\rm loc}=0.}
 \label{eq:v5-anomaly-zero}
\end{equation}
\end{theorem}

\begin{proof}
The complete chiral determinant section is exactly gauge invariant on the
protected finite gauge quotient by the inherited gauge-descent theorem.  The
bosonic, ghost and Higgs hard blocks transform by finite conjugation, and the
reference/source derivatives are taken inside that same covariant family.
Therefore every finite shell amplitude and every source derivative is gauge
equivariant before Taylor localization.  Forest contraction preserves this
property because each lower-order insertion is a restriction of a common
gauge-restored action, not an independently chosen graphwise counterterm.
The local projector and the packet descendant operations are representation
intertwiners.  Hence the complete native local breaking has zero coefficient
along every nontrivial internal-gauge descent class.  Equivalently, one may
decompose the one-fermion-loop seed into the five representation traces and use
the inherited Standard-Model trace cancellations; higher forest insertions
multiply gauge-covariant invariant tensors and cannot recreate a nonzero
quotient coefficient once the full finite section is invariant.  This proves
\eqref{eq:v5-anomaly-zero} without an Adler--Bardeen normalization theorem for
an individual uncancelled species.
\end{proof}

\section{Exact projected consistency from the contraction-complete packet}
\label{sec:v5-consistency}

Write the order-$q$ common coefficient as
\begin{equation}
 \widehat S_{q,j}
 =\bar S_{q,j}^{\loc}+K_{q,j}^{\remn},
 \label{eq:v5-full-coefficient}
\end{equation}
where $K_{q,j}^{\remn}$ is not discarded: its packet of the depth prescribed in
Definition~\ref{def:v5-induction-datum} is part of the state.  Let
\begin{equation}
 \widehat S_j^{[<L]}
 :=\sum_{q=0}^{L-1}\hbar^q\widehat S_{q,j}.
 \label{eq:v5-action-less-L}
\end{equation}

For any formal action $S=\sum_q\hbar^qS_q$, the finite quantum consistency
identity gives, at coefficient $n\ge1$,
\begin{equation}
 (S_0,\mathcal M_n)
 =\Delta_j\mathcal M_{n-1}
 -\sum_{p=1}^{n-1}(S_p,\mathcal M_{n-p}),
 \label{eq:v5-consistency-coefficient}
\end{equation}
with $\mathcal M_0=\frac12(S_0,S_0)=0$.  The point of the packet is that every
local descendant on the right of \eqref{eq:v5-consistency-coefficient} is
already a stored coordinate rather than an omitted remainder contact.

\begin{proposition}[Packet-complete lower-order restoration]
\label{prop:v5-packet-restoration}
Assume inductively that for every $q<L$ all local master coordinates through
loop order $q$ and all of their future descendants of depth $L-q$ vanish after
the common lower-order restorers have been installed.  Then the local
projection of the right-hand side of
\eqref{eq:v5-consistency-coefficient} at $n=L$ vanishes exactly.
\end{proposition}

\begin{proof}
Each term on the right is obtained from a lower master coefficient by one
letter of the contraction alphabet: a BV self-contraction for $\Delta_j$ or a
canonical contraction with a coefficient of the action for the antibracket.
By Theorem~\ref{thm:v5-feedback-complete}, its local descendant is one of the
stored packet coordinates of that lower master coefficient.  The induction
hypothesis sets precisely those coordinates to zero.  Therefore the full local
projection of the right-hand side is zero.  No assertion is made that the
unprojected quasilocal master kernel itself vanishes.
\end{proof}

\begin{theorem}[Exact local cocycle at arbitrary prescribed fixed loop order]
\label{thm:v5-L-cocycle}
Let $\widetilde{\mathcal B}_{j,L}$ be the complete order-$L$ master coefficient
formed from the forest-prepared native shell and the already fixed coefficients
$\widehat S_{0,j},\ldots,\widehat S_{L-1,j}$ before choosing the new order-$L$
restorer.  Define
\begin{equation}
 B_{j,L}
 :=\mathsf P_{L;N,\Delta,r}\widetilde{\mathcal B}_{j,L}.
 \label{eq:v5-BL}
\end{equation}
Then
\begin{equation}
 \boxed{d_{j,L}B_{j,L}=0.}
 \label{eq:v5-BL-closed}
\end{equation}
\end{theorem}

\begin{proof}
Apply \eqref{eq:v5-consistency-coefficient} at $n=L$ to the full
packet-carrying lower-order state.  Proposition
\ref{prop:v5-packet-restoration} makes the local projection of its right-hand
side vanish exactly.  The local projector commutes with the classical
differential by Theorem~\ref{thm:v5-complex}.  This gives
\eqref{eq:v5-BL-closed}.
\end{proof}

\begin{remark}[What V5 repaired]
Without the contraction-complete packet, the proof above would silently drop
terms of the form $\mathsf P\Delta K^{\remn}$ and
$\mathsf P(\bar S_p,K^{\remn})$.  V5 does not claim those terms vanish.  It
records them as lower-order future-contact coordinates and restores them when
they first become local.  This is the all-loop version of V4's insistence that
source-dependent measure/reference contacts be assembled before projection.
\end{remark}

\section{One common order-$L$ restoring primitive}
\label{sec:v5-hodge}

Choose the same type of positive coefficient inner product as in V3 on the
finite bank \eqref{eq:v5-L-bank}.  Let $d_{j,L}^\dagger$ be the adjoint,
\begin{equation}
 \Box_{j,L}=d_{j,L}d_{j,L}^\dagger+d_{j,L}^\dagger d_{j,L},
 \qquad
 h_{j,L}=d_{j,L}^\dagger G_{j,L},
 \label{eq:v5-hodge}
\end{equation}
where $G_{j,L}$ is the inverse of $\Box_{j,L}$ on the orthogonal complement of
its harmonic space.

\begin{lemma}[Fixed-$L$ Hodge contraction]
\label{lem:v5-hodge}
On a compact constant-rank component of the protected chart,
\begin{equation}
 d_{j,L}h_{j,L}+h_{j,L}d_{j,L}=I-\mathsf H_{j,L},
 \qquad
 \|h_{j,L}\|\le C_{L,\Delta,r,N}<\infty.
 \label{eq:v5-hodge-id}
\end{equation}
\end{lemma}

\begin{proof}
The bank is finite dimensional by Theorem~\ref{thm:v5-finite-local-bank}, so
finite-dimensional Hodge decomposition applies.  Constant rank on the compact
chart gives a positive lower bound for the nonzero singular values of
$d_{j,L}$, exactly as in V3.  The constant may deteriorate with $L$ and
$\Delta$; no uniform all-orders estimate is asserted.
\end{proof}

\begin{theorem}[Common order-$L$ master restorer]
\label{thm:v5-L-restorer}
The complete local breaking $B_{j,L}$ of Theorem~\ref{thm:v5-L-cocycle} has
zero harmonic projection and admits the single restoring functional
\begin{equation}
 \boxed{
 C_{j,L}^{\rm rest}:=-h_{j,L}B_{j,L},
 \qquad
 d_{j,L}C_{j,L}^{\rm rest}=-B_{j,L}.}
 \label{eq:v5-L-restorer}
\end{equation}
The same $C_{j,L}^{\rm rest}$ simultaneously restores the diffeomorphism,
local-Lorentz/spin, internal-gauge, Higgs/Yukawa, chiral-line,
ghost/nonminimal, antifield, measure/reference, stress/composite and retained
contact rows.  Its source descendants are not independently selected.
\end{theorem}

\begin{proof}
Theorem~\ref{thm:v5-L-cocycle} gives $d_{j,L}B_{j,L}=0$.  By Theorem
\ref{thm:v5-complex}, the only possible ghost-one harmonic coordinates are the
standard anomaly quotient, and Theorem~\ref{thm:v5-anomaly-zero} annihilates
that projection for the complete native Standard-Model shell.  Hence
$\mathsf H_{j,L}B_{j,L}=0$.  Apply the Hodge identity
\eqref{eq:v5-hodge-id}.  Because $B_{j,L}$ was assembled in the complete
source/contact bank before the homotopy was applied, the one primitive carries
all descendants at once.
\end{proof}

\section{The arbitrary fixed-loop shell theorem}
\label{sec:v5-shell-theorem}

Let $S_{j,L}^{\rm pre}$ be the common local order-$L$ coefficient obtained by
summing the local forest parts of Theorem~\ref{thm:v5-forest-decomposition} and
imposing the declared common physical normalization on the independent
ghost-zero EFT quotient.  Put
\begin{equation}
 \bar S_{j,L}
 :=S_{j,L}^{\rm pre}+C_{j,L}^{\rm rest}.
 \label{eq:v5-SLbar}
\end{equation}
Every quasilocal forest remainder and its contraction-complete packet is kept
as part of $K_{j,L}^{\remn}$.

\begin{theorem}[Arbitrary fixed-loop native ESM shell self-map]
\label{thm:v5-E4}
Fix the explicit V2 Wilson/overlap branch, a protected constant-rank chart and
finite data $(L,\Delta,r,N)$.  Assume the V1--V4 base case and the induction
datum $\mathfrak I_{<L}(\Delta,r,N)$.  Then the native order-$L$ shell step
produces a unique common update, modulo the declared physical finite
normalization coordinates,
\begin{equation}
 \boxed{
 \mathcal R_{j}^{[\le L]}:
 \mathfrak C^{\rm ESM}_{L;N,\Delta,r}
 \longrightarrow
 \mathfrak C^{\rm ESM}_{L;N,\Delta,r},}
 \label{eq:v5-selfmap}
\end{equation}
with the following properties:
\begin{enumerate}[label=\textup{(S\arabic*)},leftmargin=3.2em]
\item every proper forest subtraction is owned by a previously fixed common
coefficient $\bar S_{q,j}$, $q<L$;
\item the order-$L$ local action is the single coefficient
\eqref{eq:v5-SLbar}, with a finite EFT/Wilson bank at the prescribed order;
\item every retained source/contact and line coefficient is a descendant of
that same common update;
\item the complete local order-$L$ master coefficient vanishes after installing
$C_{j,L}^{\rm rest}$;
\item every quasilocal forest complement obeys
\begin{equation}
 \|K_{j,L}^{\remn}\|_*
 \le C_{L,\Delta,r,N}\epsilon_j,
 \label{eq:v5-KL-bound}
\end{equation}
and is carried with contraction depth sufficient for every later coefficient
through loop order $L$;
\item all induced lower-arity local descendants of earlier quasilocal kernels
which first appear at order $L$ are included before the master restorer is
formed.
\end{enumerate}
No uniform estimate as $L\to\infty$, no convergence of the perturbative series,
no finite-dimensional renormalizability of Einstein gravity and no
nonperturbative ESM continuum are claimed.
\end{theorem}

\begin{proof}
The finite graph family and proper-subgraph loop drop are Proposition
\ref{prop:v5-finite-graphs} and Lemma~\ref{lem:v5-loop-drop}.  Therefore all
proper local insertions are already owned, and Theorem
\ref{thm:v5-forest-decomposition} gives the finite local-plus-quasilocal
forest decomposition with the rooted bound.  Theorem
\ref{thm:v5-finite-local-bank} gives a finite target state.  The
contraction-complete packet makes every local feedback from lower quasilocal
kernels explicit, so Theorem~\ref{thm:v5-L-cocycle} gives an exact local
cocycle rather than a cocycle obtained by dropping remainder contacts.
Theorem~\ref{thm:v5-anomaly-zero} removes its harmonic obstruction and Theorem
\ref{thm:v5-L-restorer} supplies one common source-complete primitive.  Adding
that primitive to the forest-prepared local coefficient gives
\eqref{eq:v5-SLbar}; retaining the quasilocal packet gives the state on the
right of \eqref{eq:v5-selfmap}.  Every step depends only on the fixed finite
data, so the induction closes.
\end{proof}

\begin{corollary}[E4 induction]
\label{cor:v5-E4-induction}
Starting from the V1--V4 two-loop base, Theorem~\ref{thm:v5-E4} applies
recursively to every prescribed finite loop order.  Hence for each fixed
$L$ there is a source-complete native ESM shell state containing the common
local action through $L$, all required EFT/Wilson coordinates, determinant/
Pfaffian line data, complete source/contact descendants, common lower-order
forest ownership, and contraction-complete rooted/quasilocal remainders.
\end{corollary}

\begin{proof}
The base orders are the V1--V4 construction.  If the datum exists through
$L-1$, Theorem~\ref{thm:v5-E4} constructs the order-$L$ coefficient and the
future-safe remainder packets required by Definition
\ref{def:v5-induction-datum}.  This is exactly the induction hypothesis for the
next fixed order.
\end{proof}

\begin{center}
\fbox{\parbox{0.92\textwidth}{
\textbf{E4 is closed at arbitrary prescribed fixed loop order.}
The closure is coefficientwise and native: the local bank is finite only after
$(L,\Delta,r,N)$ are fixed, lower counterterms are owned by the same previously
normalized actions, and every quasilocal remainder is retained with the finite
set of future BV/source contraction jets which can feed a later local
coefficient.  Nothing here is uniform in $L$.}}
\end{center}

\section{A stronger anti-circularity audit}
\label{sec:v5-audit}

\begin{enumerate}[leftmargin=2.4em]
\item \textbf{No hidden infinite graph family.}
Quadratic and tadpole coordinates are absorbed into the common reference before
graph counting; source-marked low-valence vertices are bounded by $r$.
Therefore fixed loop number and arity genuinely give a finite topology family.

\item \textbf{No graphwise finite refits.}
Every proper 1PI ultraviolet block has smaller loop number and is replaced by
the restriction of the already fixed common action at that order.

\item \textbf{No continuum forest theorem is smuggled in.}
The forest formula is a finite native combinatorial identity.  Its estimate is
proved from native annular symbol bounds plus the rooted contraction ideal.

\item \textbf{Overlaps are not multiplied as if disjoint.}
Nonnested overlapping blocks belong to distinct compatible forests.  Nested
blocks are evaluated recursively inside-out.

\item \textbf{No discarded source or line contacts.}
Differentiation of a forest is performed before local projection.  The
contraction-complete packet stores every later local descendant generated by
source contacts, the BV Laplacian or determinant/Pfaffian line operations.

\item \textbf{The V4 remainder is not silently set to zero.}
Its current complementary kernel remains $O(\epsilon_j)$ and is carried in the
state.  What is set to zero inductively is the finite family of local descendant
coordinates required through the prescribed final loop order.

\item \textbf{No all-loop cohomology uniformity is claimed.}
The Hodge inverse is taken on a finite bank at each fixed $L$ and may have a
constant that grows arbitrarily badly with $L$.

\item \textbf{No nonrenormalizable-gravity sleight of hand.}
New ghost-zero gravitational EFT coordinates are admitted whenever the fixed
loop/EFT filtration requires them.  E4 proves finite closure at fixed order,
not a finite universal gravitational coupling set.
\end{enumerate}

\section{What E5 must now prove}
\label{sec:v5-E5-handoff}

E4 supplies a shell map at every fixed perturbative order.  The next problem is
not another loop induction.  It is an infinite-shell composition theorem at
fixed $(L,\Delta,r,N)$.

The contraction-complete packet clarifies the correct E5 state.  Let
$\mathfrak C_j^{(L,\Delta,r,N)}$ denote the full coefficient vector after the
shell at scale $\Lambda_j$, including local normalized coordinates and the
finite future-contact packet of its quasilocal remainders.  For two ultraviolet
cutoffs $J_2>J_1$, the difference of their induced low-scale coefficients must
be expressed as the sum of the shell tails
\begin{equation}
 \mathfrak C^{(J_2)}-\mathfrak C^{(J_1)}
 =\sum_{j=J_1+1}^{J_2}\delta_j\mathfrak C,
 \label{eq:v5-tail-sum}
\end{equation}
with the same finite normalization conditions imposed at every cutoff.
The forest theorem gives one-excess estimates for primitive complementary
pieces, but E5 must still prove that the \emph{composed} linearized transport
of those pieces through all lower shells has a cutoff-uniform bound.  That is
the genuine Cauchy theorem; it is not automatic from summability of the raw
single-shell remainder.

\begin{theorem}[Exact next load-bearing theorem for V6]
\label{thm:v5-next-E5}
Fix $(L,\Delta,r,N)$ and a positive reference/IR scale.  Prove a
cutoff-uniform bound on the derivative of the composed shell transport on the
complete contraction-closed state,
\begin{equation}
 \sup_{J>j\ge j_0}
 \left\|
 D\bigl(\mathcal R_{j_0}\circ\cdots\circ\mathcal R_{J}\bigr)
 \right\|_{L,\Delta,r,N}
 \le C_{L,\Delta,r,N},
 \label{eq:v5-composed-transport-goal}
\end{equation}
with the appropriate triangular interpretation on physical normalization
coordinates.  Then combine this with the fixed-shell excess estimate to prove
that every normalized connected coefficient and every stored BV/source/line
contact coefficient is Cauchy as the UV cutoff is removed.  The proof must also
show that the common local master identities and their contraction-descendant
coordinates pass to the same limit.  Scheme/reference compatibility and order
compatibility are reserved for E6.
\end{theorem}

\section{V5 status ledger}
\label{sec:v5-status}

\subsection*{Proved in V5}
\begin{enumerate}[leftmargin=2.2em]
\item The correct arbitrary-loop state cannot consist only of a local Taylor
bank plus an undifferentiated remainder.  A finite contraction-complete jet
packet is sufficient to retain every local feedback that the remainder can
produce before the prescribed final loop order.
\item After absorbing source-free valence-$0,1,2$ coordinates into the common
background/reference block, the graph, decoration and ancestry family is
finite at fixed $(L,\Delta,r,N)$.
\item Every proper connected 1PI ultraviolet subgraph has strictly smaller loop
number, so every proper subtraction is owned by a previously fixed common
action.
\item The native compatible-forest preparation, including overlapping and
nested blocks and all retained markings, decomposes into a finite local
Taylor/EFT/source part plus a rooted quasilocal complement of size
$O(\mu/\Lambda_j)$ (or $O(a\mu)$ at the ultraviolet cap).  Terms with $s$
complementary ancestors carry $O(\epsilon_j^s)$ after their new local descendants
are extracted.
\item The finite local bank closes under the classical BV differential and has
the same four-dimensional ghost-one anomaly quotient as the V80/V3 ESM
complex; the complete native Standard-Model shell has zero projection on that
quotient.
\item Because future local descendants of lower remainders are stored rather
than discarded, the order-$L$ local master breaking is an exact classical BV
cocycle.  One finite Hodge primitive therefore restores the complete local
order-$L$ master/source row.
\item These facts give a native source-complete shell self-map for every
prescribed fixed loop order, with constants permitted to depend on $L$.
Therefore E4 is closed at the stated coefficientwise shell scope.
\end{enumerate}

\subsection*{Inherited}
\begin{enumerate}[leftmargin=2.2em]
\item V80's bosonic two-loop forest ownership, local cohomology/range theorem,
retained soft-constraint architecture and rooted hard-annulus estimates.
\item The complete Standard-Model representation/anomaly trace cancellations,
even $SU(2)$ doublet parity and finite gauge descent of the protected chiral
section.
\item V1's overlap functional calculus, Kato frame and hard-pivot/rooted bounds;
V2's explicit Wilson branch and annular/full-zone comparison; V3's common
chiral Hodge restorer and E2 self-map; V4's complete two-loop canonical-pair and
source-contact master audit.
\end{enumerate}

\subsection*{Conditional}
\begin{enumerate}[leftmargin=2.2em]
\item The construction remains native to the explicit V2 Wilson/overlap
completion.  A different microscopic fermion kernel still requires the V2
regulator-comparison gate.
\item Hodge bounds are on compact constant-rank components of the finite local
complex.  Rank crossings require a new chart.
\item The determinant/Pfaffian packet remains on the protected nonzero component
with the retained positive IR/reference prescription; no global line holonomy
or zero-mode crossing theorem is claimed.
\item The native annular/rooted bounds used in the forest theorem are consumed
only at fixed graph size.  Their constants may grow without control as
$L\to\infty$.
\end{enumerate}

\subsection*{Open}
\begin{enumerate}[leftmargin=2.2em]
\item E5: prove cutoff-uniform composed transport and the Cauchy/unique
multi-shell continuum for every fixed $(L,\Delta,r,N)$, including the complete
contraction-descendant packet, operator mixing, stress/contact and line
insertions.
\item Prove that the limiting normalized connected generating-functional
coefficients inherit the common BV/master/source identities.  This must be done
on the composed state, not inferred merely from the single-shell
$O(\epsilon_j)$ estimate.
\item E6: finite scheme/reference compatibility, identification of
scheme-invariant normalized relations, and order compatibility between the
$L$ and $L+1$ constructions.
\end{enumerate}

\subsection*{Exact next load-bearing theorem}
\begin{center}
\fbox{\parbox{0.92\textwidth}{
\textbf{E5: fixed-order multi-shell Cauchy theorem with contraction-complete transport.}
At fixed $(L,\Delta,r,N)$ and positive reference scale, prove a uniform bound
for the composed derivative of the common shell maps on the full local plus
future-contact state.  Use that bound and the V5 one-excess forest estimate to
show that the normalized connected/source/BV/line coefficients form a Cauchy
family as the ultraviolet cutoff is removed, and that the common master/source
identities pass to the same limit.}}
\end{center}

\section*{Append-only continuation marker: V5}
\addcontentsline{toc}{section}{Append-only continuation marker: V5}
\textbf{End of V5.}  The complete V1--V4 body above is retained verbatim.  V5
closes E4 by replacing the naive ``local jet plus remainder'' state with a
finite contraction-complete jet packet, proving the native arbitrary-fixed-loop
forest decomposition, establishing common lower-order counterterm ownership,
and restoring the complete local order-$L$ master/source coefficient with one
Hodge primitive.  Future versions must append after this marker.  The next file
is to be named
\texttt{Renewal\_Perturbative\_ESM\_Continuum\_Research\_Notes\_V6.tex}.




\clearpage
\part*{V6: Fixed-order multi-shell continuum in a renormalized transport frame}
\addcontentsline{toc}{part}{V6: Fixed-order multi-shell continuum in a renormalized transport frame}
\label{part:v6}

\section{Purpose, inherited input and the E5 quantifier correction}
\label{sec:v6-purpose}

V5 closed the native shell problem at every prescribed finite loop order.  The
remaining E5 question is different: at fixed
\[
 (L,\Delta,r,N),
\]
one must compose arbitrarily many ultraviolet shells and prove that the
\emph{renormalized} connected/source/BV/line coefficients have a unique
cutoff-independent limit.  We keep throughout the explicit V2 Wilson/overlap
branch, a protected constant-rank chart, a positive reference/infrared scale
$\mu>0$, a finite external/source panel, and one fixed normalization/reference
prescription.  Scheme comparison is postponed to E6.

The exact next theorem stated at the end of V5 asked for a uniform bound on the
raw derivative of a long product of shell maps.  That is stronger than E5
requires and, in marginal directions, generally false.  The first task of V6
is therefore to correct the variable rather than to force a spurious
contraction estimate.

Let
\[
 \Lambda_j=b^j\mu_0,\qquad b>1,
 \qquad
 \epsilon_j:=\frac{\mu}{\Lambda_j},
\]
on the physical annular regime.  On the terminal ultraviolet cap we use the
already inherited equivalent form $\epsilon_j=a_j\mu$.  At fixed
$(L,\Delta,r,N)$ all coefficient spaces below are finite products of the local
bank and the finite contraction-complete rooted packet of V5.

\begin{proposition}[The raw composed-derivative bound is not an invariant E5 requirement]
\label{prop:v6-raw-obstruction}
There exist finite triangular shell transports which preserve a perfectly
well-defined renormalized coefficient but for which
\[
 \sup_{J>j}\left\|D(\mathcal R_j\circ\cdots\circ\mathcal R_J)\right\|=\infty.
\]
Consequently the raw version of \eqref{eq:v5-composed-transport-goal} cannot be
used as an intrinsic acceptance criterion for E5.
\end{proposition}

\begin{proof}
On $\mathbb R^2$ consider
\[
 \mathcal R_k(x,y)=(x,y+\beta x),\qquad \beta\ne0.
\]
Then
\[
 D\mathcal R_k=
 \begin{pmatrix}1&0\\ \beta&1\end{pmatrix},
 \qquad
 D(\mathcal R_j\circ\cdots\circ\mathcal R_J)=
 \begin{pmatrix}1&0\\ (J-j+1)\beta&1\end{pmatrix}.
\]
The raw derivative therefore grows linearly with the number of shells.  Yet
with the cutoff-dependent triangular coordinate
\[
 \widehat y_k:=y-k\beta x
\]
the same transport is exactly the identity.  This is the finite-dimensional
model of logarithmic operator mixing or marginal finite renormalization.  E5
must control the quotient transport after the local normalization cocycle has
been factored out; it must not demand bounded bare renormalization constants.
\end{proof}

\begin{remark}[What is and is not being changed]
\label{rem:v6-no-scheme-change}
The coordinate correction below is not a comparison of renormalization
schemes.  We keep the single common physical/source/reference normalization
already fixed in V1--V5 and merely express every cutoff in that same normalized
frame.  The possibly large triangular map from bare/local coordinates to this
frame is allowed to depend on the cutoff.  E6 will compare two different such
choices.
\end{remark}

\section{Normalization-complete local coordinates}
\label{sec:v6-normalization-frame}

Let $\mathcal L_{j}^{(L,\Delta,r,N)}$ be the finite local coefficient bank of
V5.  Split its coordinates, using the fixed V5 cohomology/Hodge decomposition,
as
\begin{equation}
 \mathcal L_j
 =\mathcal P\oplus\mathcal X_j\oplus\mathcal E_j.
 \label{eq:v6-local-split}
\end{equation}
Here $\mathcal P$ contains all independent renormalized physical/EFT and
source-normalization coordinates retained at the prescribed order,
$\mathcal X_j$ is the chosen finite BV-exact/redundant complement, and
$\mathcal E_j$ denotes the finite collection of local descendants fixed by the
common master/source equations.  This notation does not identify
$\mathcal E_j$ with a new counterterm space: its coordinates are determined by
the same common action and Hodge section.

Choose linear normalization functionals
\begin{equation}
 \mathfrak n_j:\mathcal L_j\longrightarrow\mathcal P
 \label{eq:v6-normalization-map}
\end{equation}
which read the prescribed masses, couplings, finite EFT/Wilson coordinates,
renormalized composite/source basis, and declared line/reference
normalizations at the fixed reference scale.  At fixed finite filtration we
may, and do, choose the local basis dual to these functionals on the quotient.
Thus the quotient block of $D\mathfrak n_j$ is the identity.  The V1--V5
finite-to-ordinary comparison implies that the remaining finite blocks differ
between adjacent ultraviolet cutoffs by $O(\epsilon_j)$.

\begin{lemma}[Uniform local normalization section at fixed filtration]
\label{lem:v6-normalization-section}
Fix a compact normalized parameter set $K_{\rm ren}\Subset\mathcal P$ inside
the protected component.  For all sufficiently large shell indices there is a
unique smooth local section
\begin{equation}
 \sigma_j:K_{\rm ren}\longrightarrow\mathcal L_j,
 \qquad
 \mathfrak n_j\circ\sigma_j=\operatorname{id},
 \label{eq:v6-section}
\end{equation}
compatible with the V5 Hodge choice and common reference convention, and
\begin{equation}
 \sup_j\|D\sigma_j\|<\infty,
 \qquad
 \|\sigma_{j+1}-\sigma_j\|_{C^1(K_{\rm ren})}
 \le C_{L,\Delta,r,N}\epsilon_j.
 \label{eq:v6-section-bound}
\end{equation}
The finitely many lower shells may be absorbed into the constant.
\end{lemma}

\begin{proof}
On the independent quotient block the chosen normalization coordinates are
literally the coefficient coordinates, so the Jacobian is the identity.  On
the exact/descendant block the section is obtained by the fixed finite Hodge
splitting of V5 and by the declared reference conditions.  Constant rank on
the compact protected chart gives a uniform lower singular-value bound for the
finite matrices entering that splitting.  Hence the implicit-function/Hodge
solution has a uniformly bounded derivative.

Between adjacent ultraviolet cutoffs, every coefficient entering these finite
matrices is either an exactly fixed normalization coordinate or a native
symbol/source/line coefficient.  V1--V5 give the latter an adjacent-cutoff
finite-to-ordinary discrepancy bounded by the one-excess factor
$C\epsilon_j$, including all source and contraction descendants.  The resolvent
identity for the finite inverse matrices then gives the same $C^1$ bound for
the sections.  No uniformity in $L$ is used.
\end{proof}

Let $\mathcal Q_j$ denote the finite V5 contraction-complete quasilocal packet.
The full state is
\begin{equation}
 \mathfrak C_j=\mathcal L_j\oplus\mathcal Q_j.
 \label{eq:v6-full-state}
\end{equation}
We use the section to identify its normalized version with a fixed finite
Banach coordinate space
\begin{equation}
 \widehat{\mathfrak C}
 :=K_{\rm ren}\times\widehat{\mathcal Q},
 \label{eq:v6-normalized-state}
\end{equation}
where the packet basis is transported by the same source/operator/frame
normalization.  The resulting cutoff-dependent identification will be denoted
$\Theta_j$.

\begin{definition}[Reduced or renormalized shell transport]
\label{def:v6-reduced-shell}
The E5 shell map in normalized coordinates is
\begin{equation}
 \widehat{\mathcal R}_j
 :=\Theta_{j-1}\circ\mathcal R_j^{[\le L]}\circ\Theta_j^{-1}.
 \label{eq:v6-reduced-shell}
\end{equation}
The local normalization cocycle, including triangular composite/operator
mixing and line-frame normalization, is therefore contained in the two
$\Theta$ factors rather than in the dynamical remainder of
$\widehat{\mathcal R}_j$.
\end{definition}

\section{One-shell normal form after removing the local cocycle}
\label{sec:v6-one-shell-normal-form}

\begin{theorem}[Summable normal form of the fixed-order shell map]
\label{thm:v6-one-shell-normal-form}
For every fixed $(L,\Delta,r,N)$ there is a neighborhood
$U\supset K_{\rm ren}\times\{0\}$ and constants $C_k$ such that, for every
fixed derivative order $k$ required by the retained source/contact panel,
\begin{equation}
 \boxed{
 \widehat{\mathcal R}_j=I+\mathcal E_j,
 \qquad
 \|\mathcal E_j\|_{C^k(U)}
 \le C_k\epsilon_j.}
 \label{eq:v6-I-plus-E}
\end{equation}
The statement includes the complete contraction-descendant packet.  In
particular, it is not merely a bound on the undifferentiated effective action.
\end{theorem}

\begin{proof}
Apply the V5 forest decomposition to one shell.  Every local Taylor/EFT/source
coefficient splits into an independent quotient coordinate and a dependent
exact/descendant coordinate.  The independent coordinates are reset by the
same normalization functionals at the two cutoffs; after conjugation by
$\Theta_j$ their shell transport is exactly the identity.  The dependent local
coordinates are values of the uniformly bounded Hodge/reference section of
Lemma~\ref{lem:v6-normalization-section}.  Their change between adjacent
cutoffs is therefore $O(\epsilon_j)$.

Every nonlocal forest term belongs to the V5 rooted complement and is
$O(\epsilon_j)$.  More importantly, V5 stores every local coefficient that can
be created from that complement by the finite contraction alphabet.  Hence a
source derivative, BV contraction, determinant/Pfaffian line operation or
allowed canonical contraction cannot expose an unrecorded $O(1)$ local term:
that term has already been extracted into the local normalization/Hodge block.
The residual differentiated packet therefore remains $O(\epsilon_j)$ at every
fixed derivative order in the finite alphabet.  The finite number of graph,
forest, source and canonical-pair incidences at fixed
$(L,\Delta,r,N)$ may be absorbed into $C_k$.
\end{proof}

\begin{corollary}[Summable $C^1$ size]
\label{cor:v6-summable-c1}
For geometric shells,
\begin{equation}
 \sum_{j\ge j_0}
 \|\widehat{\mathcal R}_j-I\|_{C^1(U)}<\infty.
 \label{eq:v6-summable-c1}
\end{equation}
\end{corollary}

\begin{proof}
Use Theorem~\ref{thm:v6-one-shell-normal-form} and
$\sum_{j\ge j_0}\mu/\Lambda_j<\infty$.
\end{proof}

\section{Infinite products of the reduced shell maps}
\label{sec:v6-product}

We record the elementary nonlinear product estimate explicitly because it is
the actual E5 transport theorem.

\begin{lemma}[Summable near-identity composition]
\label{lem:v6-near-identity-product}
Let $U_0\Subset U$ be compact and suppose
$F_j=I+E_j$ satisfy
\[
 \|E_j\|_{C^1(U)}\le\delta_j,
 \qquad
 \sum_j\delta_j<\infty.
\]
After enlarging a finite low-shell constant if necessary, all sufficiently
high-tail compositions are defined on $U_0$, and
\begin{equation}
 \left\|D(F_{j_0}\circ\cdots\circ F_J)\right\|
 \le
 \exp\!\left(\sum_{j=j_0}^J\delta_j\right).
 \label{eq:v6-product-derivative}
\end{equation}
Moreover the compositions are uniformly Cauchy on $U_0$ and
\begin{equation}
 \|F_{j_0}\circ\cdots\circ F_{J_2}
   -F_{j_0}\circ\cdots\circ F_{J_1}\|_{C^0(U_0)}
 \le C\sum_{j>J_1}\delta_j
 \label{eq:v6-product-tail}
\end{equation}
for $J_2>J_1$.
\end{lemma}

\begin{proof}
The derivative estimate is the chain rule followed by
$\|DF_j\|\le1+\delta_j\le e^{\delta_j}$.  The total displacement of a high
shell tail is bounded by $\sum_j\|E_j\|_{C^0}$, so a compact $U_0$ with a
slightly larger neighborhood in $U$ remains in the domain; finitely many lower
shells only multiply the final constant.  For the tail estimate insert the
additional maps one at a time and transport each one-shell displacement
through the lower composition.  The derivative bound gives
\[
 \|\Delta_{j\to j_0}\|
 \le C\delta_j.
\]
Summing $j=J_1+1,\ldots,J_2$ proves
\eqref{eq:v6-product-tail}.
\end{proof}

\begin{theorem}[Cutoff-uniform composed transport in the renormalized frame]
\label{thm:v6-composed-transport}
At fixed $(L,\Delta,r,N)$,
\begin{equation}
 \boxed{
 \sup_{J>j\ge j_0}
 \left\|
 D\bigl(
 \widehat{\mathcal R}_{j_0+1}\circ\cdots\circ
 \widehat{\mathcal R}_{J}
 \bigr)
 \right\|
 \le C_{L,\Delta,r,N}<\infty.}
 \label{eq:v6-composed-transport}
\end{equation}
This is the correct normalized version of the V5 transport goal.  No bound is
claimed for the unrenormalized triangular maps $\Theta_j^{-1}\Theta_{j+1}$.
\end{theorem}

\begin{proof}
Apply Corollary~\ref{cor:v6-summable-c1} and
Lemma~\ref{lem:v6-near-identity-product}.  The finite low-shell prefix and the
fixed compact normalization chart are absorbed into the constant.
\end{proof}

\section{The ultraviolet Cauchy theorem for the complete state}
\label{sec:v6-state-Cauchy}

For a UV cutoff at shell $J$, let
$\widehat{\mathfrak C}^{(J)}_{j_0}$ denote the normalized state at the fixed
reference shell $j_0$ obtained by starting with the same normalized physical,
source and reference data at the ultraviolet end and composing the reduced
shell maps downward.  The top packet is the native finite-cutoff packet in the
same normalized frame.

\begin{lemma}[Adjacent ultraviolet boundary data differ by one excess scale]
\label{lem:v6-top-boundary}
The normalized ultraviolet boundary states satisfy
\begin{equation}
 \|\widehat{\mathfrak C}^{(J+1)}_{J+1}
   -\widehat{\mathfrak C}^{(J)}_{J}\|_{\rm matched}
 \le C_{L,\Delta,r,N}\epsilon_J,
 \label{eq:v6-top-boundary}
\end{equation}
where ``matched'' means that both are expressed in the common normalized
coordinate and source/line frame.
\end{lemma}

\begin{proof}
Independent local quotient coordinates agree exactly by construction.  The
exact/redundant descendants differ only through the adjacent-cutoff change of
the finite Hodge/reference section, bounded in
Lemma~\ref{lem:v6-normalization-section}.  The remaining native hard kernels,
line factors and all contraction descendants differ by the V1--V5
finite-to-ordinary/full-zone estimate, whose ultraviolet-cap form is
$O(a_J\mu)=O(\epsilon_J)$.  Combining the finite coordinates gives the stated
bound.
\end{proof}

\begin{theorem}[Fixed-order multi-shell Cauchy theorem]
\label{thm:v6-state-Cauchy}
For every prescribed finite $(L,\Delta,r,N)$ and fixed positive reference
scale,
\begin{equation}
 \boxed{
 \widehat{\mathfrak C}^{(J)}_{j_0}
 \longrightarrow
 \widehat{\mathfrak C}^{(\infty)}_{j_0}}
 \label{eq:v6-state-limit}
\end{equation}
in the complete local plus contraction-descendant rooted topology.  More
precisely, for $J_2>J_1$,
\begin{equation}
 \|\widehat{\mathfrak C}^{(J_2)}_{j_0}
   -\widehat{\mathfrak C}^{(J_1)}_{j_0}\|_*
 \le
 C_{L,\Delta,r,N}
 \sum_{j>J_1}\epsilon_j
 \le
 C'_{L,\Delta,r,N}\frac{\mu}{\Lambda_{J_1}}.
 \label{eq:v6-state-Cauchy-rate}
\end{equation}
The limit is independent of the chosen cofinal sequence of UV cutoffs within
the fixed regulator branch and fixed normalization/reference prescription.
\end{theorem}

\begin{proof}
Compare two cutoffs by inserting the missing ultraviolet shells one at a time.
Each insertion changes the matched ultraviolet boundary data by
$O(\epsilon_j)$ by Lemma~\ref{lem:v6-top-boundary}; equivalently, after the
boundary is fixed, it inserts one reduced shell map
$I+O(\epsilon_j)$.  Transport this perturbation to the reference scale using
Theorem~\ref{thm:v6-composed-transport}.  Its norm remains
$C\epsilon_j$.  Sum the inserted shells.  The geometric tail is
\[
 \sum_{j>J_1}\frac{\mu}{b^j\mu_0}
 =\frac{\mu}{\Lambda_{J_1}}
   \frac{b^{-1}}{1-b^{-1}},
\]
up to the harmless indexing convention.  Hence the family is Cauchy in the
complete finite product norm.  Completeness gives the limit and the same tail
bound proves independence of a cofinal subsequence.
\end{proof}

\begin{remark}[Bare running is not asserted to converge]
\label{rem:v6-bare-running}
The theorem concerns $\widehat{\mathfrak C}$, i.e. coefficients in the fixed
renormalized local/source/line frame.  Bare marginal couplings, wave-function
factors, composite mixing matrices or determinant-line trivializations may
contain powers of $\log\Lambda_J$ or stronger allowed local dependence.  Those
belong to $\Theta_J$ and are not continuum observables by themselves.
\end{remark}

\section{Direct highest-scale form of the same Cauchy estimate}
\label{sec:v6-highest-scale}

For later use in the connected functional it is convenient to unfold the shell
composition into scale-decorated forests.  This also gives an independent
anti-circularity check on Theorem~\ref{thm:v6-state-Cauchy}.

\begin{definition}[Highest-scale label]
\label{def:v6-highest-scale}
For a V5 forest-prepared graph $\Gamma$ with internal shell labels
$\mathbf h=(h_e)_{e\in E(\Gamma)}$, put
\[
 H(\mathbf h):=\max_{e\in E(\Gamma)}h_e.
\]
All local proper-subgraph assignments use the already fixed lower-order common
actions and the common normalization section.
\end{definition}

\begin{lemma}[Highest-scale excess or local ownership]
\label{lem:v6-highest-excess}
Fix a scale-decorated forest term contributing to a normalized coefficient.
After all proper subtractions, exactly one of the following holds:
\begin{enumerate}[label=\textup{(\roman*)},leftmargin=2.7em]
\item its maximal-scale connected component is entirely local through the
prescribed filtration, in which case that contribution is owned by the common
local normalization/Hodge section and does not enter the cutoff difference;
\item its normalized remainder carries at least one Taylor/root separation
factor and is bounded by
\begin{equation}
 C_{\Gamma,L,\Delta,r,N}
 \frac{\mu}{\Lambda_{H(\mathbf h)}}.
 \label{eq:v6-highest-excess}
\end{equation}
\end{enumerate}
The same dichotomy holds after every stored BV/source/line contraction.
\end{lemma}

\begin{proof}
Apply the V5 forest decomposition to the connected component containing a line
at the maximal scale.  Its complete Taylor/source jet is, by definition,
inserted into the common local bank with its inherited lower-order ownership.
If the term is not exhausted by this local jet, V5's native Taylor remainder
theorem gives one factor of external/source momentum divided by the maximal
hard scale.  Rooted contraction with lower-scale components does not remove
that factor: any local descendant thereby generated is extracted into the
contraction-complete packet, leaving the complementary term in the one-excess
ideal.  This proves the dichotomy and its stability under the finite
contraction alphabet.
\end{proof}

\begin{proposition}[Absolute convergence of scale histories]
\label{prop:v6-scale-history}
For each fixed decorated graph topology $\Gamma$ with $I_\Gamma$ internal
lines, the sum of normalized forest terms over all shell assignments is
absolutely convergent after the common local normalization.  The contribution
of assignments with $H(\mathbf h)>J$ obeys
\begin{equation}
 \sum_{H>J}(1+H-j_0)^{I_\Gamma-1}
 \frac{C_\Gamma\mu}{\Lambda_H}
 \xrightarrow[J\to\infty]{}0.
 \label{eq:v6-scale-history-tail}
\end{equation}
\end{proposition}

\begin{proof}
For fixed maximal label $H$, there are at most
$C(1+H-j_0)^{I_\Gamma-1}$ assignments of the remaining finitely many line
labels.  By Lemma~\ref{lem:v6-highest-excess}, nonlocal terms have the factor
$\mu/\Lambda_H$; purely local terms are already removed by the common
normalization and do not enter the normalized cutoff tail.  Since a polynomial
times $b^{-H}$ is summable, the series converges absolutely.  The finite V5
graph family at fixed $(L,\Delta,r,N)$ allows summation over topologies as well.
\end{proof}

\section{Normalized connected generating-functional coefficients}
\label{sec:v6-W}

Let $Z_J[J]$ be the finite-cutoff source functional in the fixed normalized
source frame and define
\begin{equation}
 W_J[J]
 :=\log\frac{Z_J[J]}{Z_J[0]}.
 \label{eq:v6-W}
\end{equation}
We use $J$ both as the conventional source letter inside $W$ and as a cutoff
index only when the meaning is unambiguous; cutoff indices will otherwise be
written $J_{\rm UV}$.

At a fixed loop/source order and positive reference scale, coefficient
extraction from the terminal low-scale state is a finite sum of connected
rooted contractions.  Denote this finite polynomial/analytic map by
\begin{equation}
 \mathcal W_{L,\Delta,r,N}:
 \widehat{\mathfrak C}_{j_0}\longrightarrow\mathcal W_{L,\Delta,r,N}.
 \label{eq:v6-W-map}
\end{equation}

\begin{lemma}[Continuity of the terminal connected-functional map]
\label{lem:v6-W-continuity}
On the protected compact reference domain,
\begin{equation}
 \|D\mathcal W_{L,\Delta,r,N}\|
 \le C_{L,\Delta,r,N}<\infty.
 \label{eq:v6-W-Lipschitz}
\end{equation}
The map includes composite and stress insertions, source contacts, operator
mixing, ghosts/antifields and determinant/Pfaffian line insertions.
\end{lemma}

\begin{proof}
At fixed perturbative and source order, $W$ is the finite connected part of the
terminal Wick/forest expansion.  Every propagator on the retained reference
domain has the positive IR/hard-pivot bound inherited from V1--V5, and every
vertex/source/line derivative belongs to the finite rooted bank.  Hence every
coefficient is a finite multilinear contraction of state coordinates with a
uniform bound on the compact domain.  Differentiating that finite expression
again gives a finite sum of the same type.  Determinant/Pfaffian insertions use
the V1--V2 line inverse bounds and their stored derivatives.  The normalization
$Z[J]/Z[0]$ removes the source-free vacuum coordinate exactly; a field-independent
but source-dependent contact is not removed and remains one of the finite
coordinates of the map.
\end{proof}

\begin{theorem}[Source-complete coefficientwise ultraviolet continuum]
\label{thm:v6-W-continuum}
For every prescribed fixed $(L,\Delta,r,N)$ and finite external/source panel,
every coefficient of the normalized connected functional has a unique UV
limit
\begin{equation}
 \boxed{
 [W_\infty]_I^{(\ell)}
 :=\lim_{J_{\rm UV}\to\infty}
 [W_{J_{\rm UV}}]_I^{(\ell)},
 \qquad
 \ell\le L,\quad |I|\le r.}
 \label{eq:v6-W-limit}
\end{equation}
The convergence holds jointly for the retained momentum derivatives through
$\Delta$ and for the complete finite panel of composite, stress, contact,
operator-mixing, ghost/antifield and determinant/Pfaffian line insertions.
Moreover
\begin{equation}
 \|W_{J_2}-W_{J_1}\|_{L,\Delta,r,N}
 \le
 C_{L,\Delta,r,N}\frac{\mu}{\Lambda_{J_1}}
 \label{eq:v6-W-rate}
\end{equation}
for $J_2>J_1$ sufficiently large.
\end{theorem}

\begin{proof}
Apply Lemma~\ref{lem:v6-W-continuity} to the Cauchy estimate
\eqref{eq:v6-state-Cauchy-rate}.  This already proves the theorem.  Equivalently,
unfold each connected coefficient into its finite family of decorated graph
topologies and use Proposition~\ref{prop:v6-scale-history}; the logarithm
selects connected graphs and the vacuum normalization deletes exactly the
source-free vacuum family.  Both routes preserve all marked/contact and line
coordinates because these were part of the state before either composition or
connected extraction.
\end{proof}

\section{Operator mixing and the meaning of the limiting source basis}
\label{sec:v6-operator-mixing}

Let $\mathbf O$ be any finite set of composite/source operators retained in the
panel.  In raw coordinates one may have a triangular cutoff-dependent mixing
matrix $Z_j^{\mathbf O}$.  The normalization map $\Theta_j$ uses the prescribed
reference conditions to define
\begin{equation}
 \mathbf O_R=(Z_j^{\mathbf O})^{-1}\mathbf O_{\rm bare}.
 \label{eq:v6-renormalized-operator}
\end{equation}

\begin{corollary}[Coefficientwise continuum with operator mixing]
\label{cor:v6-operator-mixing}
All mixed connected coefficients of the finite renormalized operator panel
$\mathbf O_R$ converge as in Theorem~\ref{thm:v6-W-continuum}.  No convergence
of $Z_j^{\mathbf O}$ itself is required.
\end{corollary}

\begin{proof}
The renormalized operator/source basis is one block of the normalized state
coordinate $\widehat{\mathfrak C}$.  The possibly nonconvergent triangular
matrix is part of $\Theta_j$, already factored from the reduced shell map.
Apply Theorem~\ref{thm:v6-W-continuum} in that fixed normalized basis.
\end{proof}

\section{Passage of the common BV/master/source identities}
\label{sec:v6-master-limit}

The identity which can be passed without adding a new hypothesis is exactly the
fixed-order identity proved in V5: the complete local master coordinates and
all contraction descendants required by the prescribed final filtration vanish
in one common action/source state.  We do not silently promote an unprojected
finite-cutoff quasilocal master remainder into zero.

Let $\mathbf M_{L,\Delta,r,N}$ denote the finite vector consisting of:
\begin{enumerate}[label=\textup{(M\arabic*)},leftmargin=3em]
\item every local coefficient through the retained EFT/derivative order of
$\frac12(S,S)-\hbar\Delta S$;
\item every source/contact derivative of those coefficients through order $r$;
\item every future contraction descendant stored by the V5 packet through
final loop order $L$;
\item the corresponding finite connected Ward/Zinn--Justin coefficients in the
normalized source panel.
\end{enumerate}
By V5 these coordinates are zero for every finite cutoff after the common
restorers and normalizations are installed.

\begin{theorem}[Master/source identities commute with the UV limit]
\label{thm:v6-master-limit}
For the limiting normalized state of Theorem~\ref{thm:v6-state-Cauchy},
\begin{equation}
 \boxed{
 \mathbf M_{L,\Delta,r,N}
 \bigl(\widehat{\mathfrak C}^{(\infty)}_{j_0}\bigr)=0.}
 \label{eq:v6-master-limit}
\end{equation}
Equivalently, every retained coefficient of the common BV/BRST/master/source
identity and every stored stress/composite/contact/line descendant passes to
the same coefficientwise continuum limit.
\end{theorem}

\begin{proof}
The vector $\mathbf M_{L,\Delta,r,N}$ is a finite polynomial combination of
state coordinates, finite canonical contractions, the fixed finite BV
Laplacian blocks and source coefficient maps.  On the protected compact domain
it is continuous in the V5 rooted/product topology.  For every finite cutoff,
V5's exact projected consistency and common Hodge restoration give
\[
 \mathbf M_{L,\Delta,r,N}
 (\widehat{\mathfrak C}^{(J)}_{j_0})=0.
\]
Take $J\to\infty$ using Theorem~\ref{thm:v6-state-Cauchy}.  Because the
contraction-complete packet was included in the convergent state, no term of
the form $\mathsf P\Delta K^{\rm rem}$ or
$\mathsf P(S_p,K^{\rm rem})$ is lost in this passage.
\end{proof}

\begin{remark}[Scope of the master theorem]
\label{rem:v6-master-scope}
The theorem is coefficientwise at the prescribed finite EFT/derivative/source
order.  It does not claim that a finite-cutoff unprojected quasilocal master
kernel vanishes identically, nor does it claim an all-derivative or
nonperturbative Slavnov--Taylor functional identity in one step.  For any
larger finite derivative/source order the same construction may be rerun, and
E6 will impose compatibility between the resulting finite-order systems.
\end{remark}

\section{Separate thermodynamic comparison}
\label{sec:v6-thermodynamic}

The ultraviolet theorem above is formulated at fixed positive reference scale
and on the same-torus/nonwinding domain used in the inherited native shell
estimates.  The volume limit is logically separate.

\begin{lemma}[Uniform periodization tail for a finite rooted panel]
\label{lem:v6-volume-tail}
Suppose a retained infinite-volume kernel $K$ obeys, for some $m>4+s$,
\begin{equation}
 \int_{\mathbb R^4}(1+\mu|x|)^m|K(x)|\,dx\le C
 \label{eq:v6-weighted-kernel}
\end{equation}
uniformly in the UV cutoff and in the fixed source panel.  Let $K_R$ be its
periodization on a four-torus with minimal side $R$.  Then on the nonwinding
rooted panel with $s$ fixed relative coordinates,
\begin{equation}
 \|K_R-K\|_{*,s}
 \le C_{m,s}(\mu R)^{-(m-4-s)}.
 \label{eq:v6-volume-bound}
\end{equation}
\end{lemma}

\begin{proof}
Write the periodization error as the sum over nonzero lattice translates of the
fundamental cell.  On the nonwinding panel every such translate has distance
at least a fixed multiple of $R|n|$.  Pull
$(1+\mu R|n|)^{-(m-4-s)}$ from the weighted $L^1$ moment and sum the remaining
$|n|^{-(4+\eta)}$ tail.  Rooting removes the overall volume factor; the $s$
relative-coordinate weights cost at most $s$ powers.  This proves the stated
bound.
\end{proof}

\begin{theorem}[Staged thermodynamic continuum]
\label{thm:v6-volume-limit}
The V1--V5 arbitrary fixed polynomial-moment bounds imply the hypothesis of
Lemma~\ref{lem:v6-volume-tail} for every prescribed finite rooted/source panel.
Hence, after taking the ultraviolet limit of
Theorem~\ref{thm:v6-W-continuum}, the finite-volume normalized connected
coefficients have a unique thermodynamic limit along the declared nonwinding
panels.  The UV and volume limits commute whenever both are restricted to a
common compact source/momentum panel.
\end{theorem}

\begin{proof}
Choose $m$ larger than the fixed number of rooted relative-coordinate weights.
The inherited annular symbol estimates permit that arbitrary fixed number of
momentum integrations by parts, uniformly in the cutoff, and the V5 forest
contractions preserve the corresponding weighted rooted moment.  Lemma
\ref{lem:v6-volume-tail} therefore gives a volume error uniform in the UV
cutoff.  For cutoffs $J_1,J_2$ and volumes $R_1,R_2$, insert and subtract the
common infinite-volume or common finite-volume coefficient.  The triangle
inequality combines the UV tail
$C\mu/\Lambda_{J_1}$ with the volume tail
$C(\mu R_1)^{-q}$.  Both tend to zero independently, proving existence and
commutation of the staged limits.
\end{proof}

\section{E5 closure theorem}
\label{sec:v6-E5-close}

\begin{theorem}[Coefficientwise multi-shell Einstein--Standard Model continuum at fixed order]
\label{thm:v6-E5}
Fix the explicit V2 native Wilson/overlap branch, a protected constant-rank
component, a positive reference/IR scale, one common normalization/reference
scheme, finite $(L,\Delta,r,N)$, and a finite external/source panel.  Then:
\begin{enumerate}[label=\textup{(E5.\arabic*)},leftmargin=3.5em]
\item after factoring the finite-dimensional local normalization cocycle, every
native shell map is $I+O(\mu/\Lambda_j)$ in the complete
contraction-descendant $C^1$ state;
\item the reduced shell maps have cutoff-uniform composed transport and an
infinite-shell limit;
\item the common normalized local/source/line state is Cauchy with tail
$O(\mu/\Lambda_J)$;
\item every normalized connected coefficient through loop order $L$, derivative
order $\Delta$ and source order $r$ has a unique cutoff-independent limit,
including composite operators, stress insertions, contact terms, operator
mixing, ghosts/antifields and determinant/Pfaffian line insertions;
\item every retained common BV/BRST/master/source coefficient and every stored
contraction descendant passes to the same limit and remains zero;
\item the thermodynamic comparison can be taken separately, with uniform rooted
periodization tails, on the declared nonwinding finite panels.
\end{enumerate}
No convergence of the perturbative series, no uniform $L\to\infty$ estimate,
no finite-dimensional renormalizability of Einstein gravity, no asymptotic
safety, no nonperturbative ESM measure and no removal of the retained positive
IR/reference scale are claimed.
\end{theorem}

\begin{proof}
Item (1) is Theorem~\ref{thm:v6-one-shell-normal-form}; item (2) is Theorem
\ref{thm:v6-composed-transport} and Lemma~\ref{lem:v6-near-identity-product};
item (3) is Theorem~\ref{thm:v6-state-Cauchy}; item (4) is Theorem
\ref{thm:v6-W-continuum} and Corollary~\ref{cor:v6-operator-mixing}; item (5)
is Theorem~\ref{thm:v6-master-limit}; item (6) is Theorem
\ref{thm:v6-volume-limit}.  All statements use only fixed finite
$(L,\Delta,r,N)$ and the same normalization/reference prescription.
\end{proof}

\begin{center}
\fbox{\parbox{0.92\textwidth}{
\textbf{E5 is closed at the stated fixed-order coefficientwise scope.}
The essential correction is that the continuum theorem is formulated in the
common renormalized local/source/line frame.  Raw bare transport is allowed to
carry logarithmic or powerlike local renormalization factors; after this local
cocycle is factored out, the genuinely quasilocal shell evolution is a
summable near-identity product.}}
\end{center}

\section{Why E6 is now the only remaining programme gate}
\label{sec:v6-E6-handoff}

E5 fixes one admissible normalization/reference prescription and constructs its
coefficientwise continuum.  The remaining questions are compatibility
questions, not new ultraviolet estimates.

Let $\mathcal S$ and $\mathcal S'$ be two admissible finite schemes/reference
choices on the same native regulator branch.  At every finite filtration they
are related by a finite triangular local/source/line reparametrization
$F^{\mathcal S'\leftarrow\mathcal S}_J$.  E6 must prove that the corresponding
renormalized continuum coordinate change has a cutoff-independent limit,
identify which coefficient functions transform covariantly and which
normalized relations are invariant, and show that the construction at loop
order $L+1$ restricts exactly to the already fixed order-$L$ continuum rather
than refitting its lower coefficients.

\begin{theorem}[Exact next load-bearing theorem for V7]
\label{thm:v6-next-E6}
For every fixed finite filtration, prove:
\begin{enumerate}[label=\textup{(G\arabic*)},leftmargin=3em]
\item any two admissible normalization/source/reference schemes are related by
a finite local triangular map whose renormalized cutoff sequence is Cauchy and
has an invertible continuum limit;
\item changes of determinant/Pfaffian local line frame and positive reference
normalization act by the already stored finite cocycle/contact rules and do not
alter normalized gauge-invariant connected relations;
\item the order-$(L+1)$ local bank, forest ownership, Hodge section and reduced
transport restrict to the previously fixed order-$L$ data, so the continuum
coefficients form a projective family in $L$;
\item after these compatibility statements, formulate the final theorem: for
every prescribed finite perturbative, EFT/derivative and source order, the
native Renewal Einstein--Standard Model regulator admits a source-complete
renormalized coefficientwise continuum satisfying the corresponding common
BV/BRST/master identities and Standard-Model anomaly constraints, with gravity
interpreted as an EFT and all retained IR/reference assumptions stated
explicitly.
\end{enumerate}
\end{theorem}

\section{V6 anti-circularity audit}
\label{sec:v6-audit}

\begin{enumerate}[leftmargin=2.4em]
\item \textbf{No false raw contraction hypothesis.}
Proposition~\ref{prop:v6-raw-obstruction} shows why marginal/source mixing can
make the bare composed derivative grow.  V6 factors the already-declared local
normalization cocycle before asking for summable transport.
\item \textbf{No hidden scheme comparison.}
Only one fixed normalization/reference prescription is used.  Different
schemes are not identified in V6; that is E6.
\item \textbf{No dropped remainder contacts.}
The $I+O(\epsilon_j)$ statement is made on the full V5
contraction-complete packet, not on a bare local jet plus an undifferentiated
remainder.
\item \textbf{No graphwise retuning.}
Purely local maximal-scale pieces are owned by the same common local
normalization/Hodge section.  Proper forest blocks still use the previously
fixed common lower-order actions.
\item \textbf{No inference from one-shell summability alone.}
The nonlinear product lemma and normalized composed-transport theorem are
proved explicitly; the highest-scale forest expansion supplies an independent
absolute-convergence check.
\item \textbf{No loss of source-dependent constants.}
$W=\log(Z[J]/Z[0])$ removes only source-free vacuum coordinates.
Field-independent but source-dependent contacts remain in the convergent state
and in the limiting connected functional.
\item \textbf{No unproved all-derivative master identity.}
The identity passed to the limit is exactly the finite coefficient/descendant
system closed in V5.  Larger finite derivative/source orders require the
corresponding larger finite construction.
\item \textbf{No conflation of UV and volume limits.}
The ultraviolet Cauchy theorem is proved first.  Thermodynamic comparison is a
separate uniform rooted-tail argument.
\end{enumerate}

\section{V6 status ledger}
\label{sec:v6-status}

\subsection*{Proved in V6}
\begin{enumerate}[leftmargin=2.2em]
\item The raw long-product derivative requested at the end of V5 is not an
intrinsic continuum criterion; finite triangular marginal/source mixing gives
an explicit counterexample.
\item A common normalization-complete local/source/line frame exists on the
fixed compact protected filtration, with uniformly bounded section and
adjacent-cutoff $O(\epsilon_j)$ variation.
\item In that frame the arbitrary-fixed-loop native shell map is
$I+O(\epsilon_j)$ in $C^1$ on the complete contraction-descendant state.
\item Summability of those $C^1$ increments gives a cutoff-uniform derivative
bound for the composed \emph{renormalized} transport and a nonlinear infinite
product.
\item The complete normalized local plus quasilocal packet is Cauchy, with
ultraviolet tail $O(\mu/\Lambda_J)$, and the limit is independent of the
cofinal cutoff sequence within the fixed native branch and fixed scheme.
\item Unfolding into scale-decorated forests gives the equivalent
highest-scale theorem: every maximal-scale contribution is either already
owned locally or carries one excess factor, and the remaining polynomial scale
multiplicity is dominated by the geometric shell decay.
\item Every coefficient of the normalized connected generating functional at
the prescribed finite loop/EFT/source order converges, including composite,
stress, contact, operator-mixing, ghost/antifield and determinant/Pfaffian line
insertions.
\item The complete finite vector of common master/source coefficients and V5
future-contact descendants passes continuously to the same limit and remains
zero.
\item Uniform weighted rooted moments give a separate thermodynamic
periodization bound, so the staged volume limit exists on the declared
nonwinding panels and commutes with the UV limit there.
\item Therefore E5 is closed at the stated fixed-order coefficientwise scope.
\end{enumerate}

\subsection*{Inherited}
\begin{enumerate}[leftmargin=2.2em]
\item V80's common bosonic normalization/forest architecture and finite local
cohomology range theorem.
\item The Standard-Model anomaly cancellations, protected finite chiral gauge
descent and exact determinant/Pfaffian shell cocycles.
\item V1--V2's native fermion functional calculus, hard pivot,
annular/full-zone comparison, source jets and line remainder estimates.
\item V3--V4's complete chiral two-loop shell self-map and common
canonical-pair/source-contact master audit.
\item V5's arbitrary fixed-loop forest theorem, finite local bank, common
lower-order subtraction ownership, contraction-complete packet and exact local
master restoration at every prescribed fixed loop order.
\end{enumerate}

\subsection*{Conditional}
\begin{enumerate}[leftmargin=2.2em]
\item The result remains native to the explicit V2 Wilson/overlap completion;
a different microscopic fermion kernel requires the V2 regulator-comparison
gate.
\item Normalization/Hodge statements are on compact constant-rank protected
charts.  A rank crossing requires a new chart and transition analysis.
\item Determinant/Pfaffian limits remain on the protected nonzero component
with the retained positive IR/reference prescription; no global line holonomy
or zero-mode crossing theorem is claimed.
\item The thermodynamic theorem is on the nonwinding finite rooted/source
panels controlled by the inherited weighted moments.  A massless global IR
limit is not included.
\item Constants may depend arbitrarily badly on $L,\Delta,r,N$; no all-order
uniformity or perturbative convergence is asserted.
\end{enumerate}

\subsection*{Open}
\begin{enumerate}[leftmargin=2.2em]
\item E6: compare admissible finite renormalization/source/reference schemes and
identify the invariant normalized relations.
\item Prove order compatibility: the order-$(L+1)$ construction must restrict
to the already fixed order-$L$ normalized continuum coefficients, including
source/contact and line packets.
\item Assemble these compatibility statements into the final coefficientwise
Einstein--Standard Model continuum theorem, with every retained IR/reference
assumption explicit.
\end{enumerate}

\subsection*{Exact next load-bearing theorem}
\begin{center}
\fbox{\parbox{0.92\textwidth}{
\textbf{E6: scheme/reference and order compatibility.}
Fix two admissible finite normalization/source/reference choices.  Prove that
their finite triangular relation converges to an invertible continuum
reparametrization, characterize the normalized connected relations invariant
under that change, include determinant/Pfaffian frame/reference cocycles, and
prove projective compatibility in perturbative order: the $(L+1)$ continuum
must restrict to the already fixed $L$ continuum without any lower-order
refit.  Then state the final source-complete coefficientwise Renewal ESM
continuum theorem with gravity as an EFT and a retained positive IR/reference
scale.}}
\end{center}

\section*{Append-only continuation marker: V6}
\addcontentsline{toc}{section}{Append-only continuation marker: V6}
\textbf{End of V6.}  The complete V1--V5 body above is retained verbatim.  V6
closes E5 by correcting the raw transport quantifier, factoring the common
finite-dimensional local/source/line renormalization cocycle, proving that the
reduced shell maps form a summable $C^1$ near-identity product, establishing a
Cauchy continuum for the complete contraction-descendant state and normalized
connected coefficients, passing the finite common master/source identities to
the same limit, and separating the thermodynamic comparison from the UV
removal.  Future versions must append after this marker.  The next file is to
be named
\texttt{Renewal\_Perturbative\_ESM\_Continuum\_Research\_Notes\_V7.tex}.



% ============================================================================
% V7 APPEND-ONLY CONTINUATION
% ============================================================================
\clearpage
\part{V7 append-only continuation:\\
scheme/reference covariance, projective order compatibility, and the final coefficientwise continuum theorem}
\label{part:v7}

\section{Purpose and the final compatibility problem}
\label{sec:v7-purpose}

V6 closed E5 for one fixed admissible normalization/source/reference
prescription.  Thus no new ultraviolet graph estimate is load bearing for the
stated programme.  E6 is a compatibility problem: different admissible finite
coordinate choices must describe the same coefficientwise continuum up to an
invertible finite reparametrization, determinant/Pfaffian frame changes must be
carried with their source contacts rather than discarded, and enlarging the
perturbative order from $L$ to $L+1$ must not refit any coefficient already
fixed at orders $0,\ldots,L$.

We retain throughout the V6 hypotheses: the explicit V2 Wilson/overlap branch,
a protected constant-rank perturbative component, a protected nonzero
 determinant/Pfaffian line component, a positive reference/IR scale $\mu>0$,
a fixed finite external/source panel, and prescribed finite
$(L,\Delta,r,N)$.  Thermodynamic comparison remains the separate rooted-tail
limit of V6.  No $\mu\downarrow0$ statement is made.

The first point is an anti-circularity correction.  V5 used a finite Hodge
inverse at each fixed order.  If one were to recompute an orthogonal Hodge
inverse from scratch whenever the local bank is enlarged, exact order
compatibility would not follow.  E6 therefore fixes the chronology of the
finite choices rather than attempting to prove a false functoriality property
of Moore--Penrose/Hodge inverses.

\section{Why independently recomputed Hodge sections are not projective}
\label{sec:v7-hodge-obstruction}

\begin{proposition}[Orthogonal Hodge inversion is not functorial under bank enlargement]
\label{prop:v7-hodge-not-functorial}
There are finite complexes $C\subset \widetilde C$ for which the orthogonal
contracting homotopy on $\widetilde C$, formed using the enlarged inner
product, does not restrict to the orthogonal contracting homotopy on $C$.
Consequently exact order compatibility cannot be deduced merely by choosing a
new Moore--Penrose/Hodge inverse independently at every loop order.
\end{proposition}

\begin{proof}
Let the small complex have
\[
 C^0=\operatorname{span}\{e_1\},\qquad
 C^1=\operatorname{span}\{f\},\qquad d e_1=f,
\]
with the standard inner product.  Its degree-one contracting homotopy is
$h_C f=e_1$.
Enlarge degree zero to
\[
 \widetilde C^0=\operatorname{span}\{e_1,e_2\},\qquad
 \widetilde C^1=\operatorname{span}\{f\},\qquad
 d e_1=d e_2=f,
\]
again with the standard inner product.  The Moore--Penrose inverse of the row
matrix $(1\;1)$ sends
\[
 f\longmapsto \frac12(e_1+e_2).
\]
Hence $h_{\widetilde C}f\notin C^0$ and in particular
$h_{\widetilde C}|_{C^1}\ne h_C$.  The differential on the original
subcomplex has not changed; only the ambient orthogonality problem has.
\end{proof}

\begin{firewall}[Upstream correction for E6]
From V7 onward, ``the Hodge section'' means a \emph{chronologically fixed
finite contracting section}: once the order-$q$ primitive and its
normalization have been selected, they are frozen.  Enlarging the target order
adds only the new order and the additional future-descendant coordinates.  An
orthogonal Hodge inverse may be used to choose the new complement, but it is
never allowed to recompute an already fixed lower-order primitive.
\end{firewall}

\section{A nested ambient bank and exact perturbative truncations}
\label{sec:v7-nested-bank}

Fix $(\Delta,r,N)$ and the finite external/source panel.  Instead of allowing
the operator basis itself to depend on the terminal loop order, choose once the
finite ambient local operator/source/line bank
\begin{equation}
 \mathcal A_{\Delta,r,N}
 \label{eq:v7-ambient-bank}
\end{equation}
containing every ghost-zero EFT/Wilson coordinate through derivative degree
$\Delta$, every required BV field--antifield and nonminimal coordinate, every
source/contact coordinate through order $r$, and the protected local
 determinant/Pfaffian line/reference coordinates.  The loop filtration is then
formal:
\begin{equation}
 \mathcal A^{[\le L]}_{\Delta,r,N}
 :=\bigoplus_{q=0}^{L}\hbar^q\mathcal A_{\Delta,r,N}.
 \label{eq:v7-loop-filtered-bank}
\end{equation}
Coordinates which are not generated at a particular loop order simply have
zero coefficient there; the basis is not rediscovered when $L$ changes.

For the V5 contraction-complete remainder packet, let
$\mathcal Q_q^{[d]}$ denote the order-$q$ rooted packet stored to future
contraction depth $d$.  Define
\begin{equation}
 \mathfrak C_L
 :=\mathcal A^{[\le L]}_{\Delta,r,N}
   \oplus\bigoplus_{q=0}^{L}\hbar^q\mathcal Q_q^{[L-q]}.
 \label{eq:v7-projective-state}
\end{equation}
There are canonical inclusion and truncation maps
\begin{equation}
 \iota_L^{L+1}:\mathfrak C_L\hookrightarrow\mathfrak C_{L+1},
 \qquad
 \pi_L^{L+1}:\mathfrak C_{L+1}\twoheadrightarrow\mathfrak C_L,
 \qquad
 \pi_L^{L+1}\iota_L^{L+1}=I,
 \label{eq:v7-projections}
\end{equation}
where $\pi_L^{L+1}$ discards the $\hbar^{L+1}$ layer and discards exactly one
extra future-descendant depth from each older packet.

\begin{definition}[Coherent normalization/reference family]
\label{def:v7-coherent-scheme}
An admissible scheme family
$\mathcal S=\{\mathcal S_L\}_{L\ge0}$ is \emph{coherent} if:
\begin{enumerate}[label=\textup{(C\arabic*)},leftmargin=3em]
\item the physical/EFT normalization functionals at order $q$ are fixed when
$q$ first appears and are unchanged at every terminal order $L\ge q$;
\item the source/composite basis and its triangular mixing at source order
$\le r$ restrict under \eqref{eq:v7-projections};
\item every proper forest subtraction uses the already fixed common lower-loop
action coefficient, including its finite part;
\item the order-$q$ BV restoring primitive is frozen after it is constructed;
new Hodge/contracting choices act only on the new order-$L+1$ complement;
\item determinant/Pfaffian frame and positive reference normalizations are
carried with their complete local source/contact coordinates and restrict under
$\pi_L^{L+1}$.
\end{enumerate}
\end{definition}

\begin{lemma}[Every fixed-order admissible choice extends to a coherent family]
\label{lem:v7-coherent-extension}
Let an admissible normalization/reference prescription be fixed through order
$L_0$.  On the protected constant-rank component it can be extended
inductively to a coherent family at every higher prescribed finite order.
\end{lemma}

\begin{proof}
Assume the family has been fixed through order $L$.  The new local coefficient
at order $L+1$ lies in the finite bank \eqref{eq:v7-ambient-bank}.  Keep all
normalization functionals and restoring primitives at orders $q\le L$
unchanged.  V5 gives an exact local cocycle for the new order-$L+1$ breaking,
and its Standard-Model anomaly projection vanishes.  Therefore the new
breaking lies in the finite range of the classical differential.  Choose any
bounded right inverse on a complement to the already frozen lower-order data;
for example, choose an orthogonal complement only inside this new finite
layer.  Fix the independent order-$L+1$ physical/EFT coordinates by the new
normalization conditions and append the additional future-descendant depth.
All operations are finite dimensional at the prescribed filtration and
constant rank gives the required bounded inverses.  This extends the coherent
family without modifying any old coefficient.
\end{proof}

\section{Exact finite-cutoff order compatibility}
\label{sec:v7-finite-projective}

\begin{theorem}[Forest preparation commutes with perturbative truncation]
\label{thm:v7-forest-projective}
For every native shell $j$ and every $L\ge0$, the coherent forest-prepared
shell functional satisfies
\begin{equation}
 \boxed{
 \pi_L^{L+1}\,
 \mathbf R_{j}^{[\le L+1]}\,
 \iota_L^{L+1}
 =\mathbf R_{j}^{[\le L]}.}
 \label{eq:v7-forest-projective}
\end{equation}
This identity includes all retained source/contact, BV, and
 determinant/Pfaffian line descendants.
\end{theorem}

\begin{proof}
Take a graph contributing to a coefficient of loop order $q\le L$.  Loop
number is additive under contraction of a proper 1PI subtraction block, and
V5 proved that every proper subtraction block has strictly smaller loop number
than its parent graph.  Therefore every subtraction required by the order-$q$
graph is owned by action coefficients of orders $<q\le L$ and is already
frozen in the coherent family.  A counterterm first appearing at order
$L+1$ cannot enter a coefficient of order $q\le L$ because it carries an
explicit factor $\hbar^{L+1}$.

The local Taylor projector depends only on the fixed ambient derivative/source
bank and hence is the same on both sides after truncation.  Source partitions,
canonical BV contractions and line derivatives preserve the loop grading.
The V5 future-descendant packet at terminal order $L+1$ stores one additional
contraction depth, but applying $\pi_L^{L+1}$ discards precisely that extra
depth and leaves the packet previously stored at terminal order $L$.  Thus
every forest term, local assignment and retained complement at order $q\le L$
agrees term by term, proving \eqref{eq:v7-forest-projective}.
\end{proof}

\begin{theorem}[The common restoring action is projective in loop order]
\label{thm:v7-restorer-projective}
Let
\[
 S_{j}^{[\le L]}=\sum_{q=0}^{L}\hbar^q\bar S_{j,q}
\]
be the coherent common action.  Then the order-$(L+1)$ construction satisfies
\begin{equation}
 \boxed{
 \pi_L^{L+1}S_{j}^{[\le L+1]}=S_{j}^{[\le L]}.}
 \label{eq:v7-action-projective}
\end{equation}
In particular no lower-loop physical coefficient, finite forest part,
BV-exact restorer, source/contact descendant, or line/reference coordinate is
refitted when the terminal perturbative order is increased.
\end{theorem}

\begin{proof}
By Definition~\ref{def:v7-coherent-scheme}, every coefficient
$\bar S_{j,q}$ with $q\le L$ was frozen when order $q$ was first closed.
Theorem~\ref{thm:v7-forest-projective} shows that recomputing the larger forest
produces exactly the same order-$q$ pre-action and lower-order descendant data.
The chronological contracting section therefore returns the already stored
order-$q$ primitive rather than an independently recomputed Hodge primitive.
The only new local action is multiplied by $\hbar^{L+1}$ and is killed by
$\pi_L^{L+1}$.  This gives \eqref{eq:v7-action-projective}.
\end{proof}

\begin{corollary}[Exact shell-map projectivity]
\label{cor:v7-shell-projective}
For every shell $j$,
\begin{equation}
 \boxed{
 \pi_L^{L+1}\,
 \mathcal R_{j}^{[\le L+1]}\,
 \iota_L^{L+1}
 =\mathcal R_{j}^{[\le L]}.}
 \label{eq:v7-shell-projective}
\end{equation}
\end{corollary}

\begin{proof}
The shell map consists of the forest-prepared local and quasilocal update, the
common restoring action, and the stored contraction-descendant transport.
Apply Theorems~\ref{thm:v7-forest-projective} and
\ref{thm:v7-restorer-projective} to these three pieces.
\end{proof}

\section{Projective normalized frames and reduced transports}
\label{sec:v7-normalized-projective}

For a coherent scheme $\mathcal S$, choose the V6 normalization map
$\Theta_{j,L}^{\mathcal S}$ recursively: on the lower-order image of
$\iota_{L-1}^{L}$ it is the already fixed map
$\Theta_{j,L-1}^{\mathcal S}$, while only the new order-$L$ and extra packet
coordinates are normalized afresh.

\begin{lemma}[Normalized-frame projectivity]
\label{lem:v7-theta-projective}
The coherent normalization maps may be chosen so that
\begin{equation}
 \boxed{
 \pi_L^{L+1}\Theta_{j,L+1}^{\mathcal S}
 =\Theta_{j,L}^{\mathcal S}\pi_L^{L+1}.}
 \label{eq:v7-theta-projective}
\end{equation}
The same identity holds for their local inverses on the common protected
compact chart.
\end{lemma}

\begin{proof}
This is the recursive construction just described.  On old coordinates the
map is definitionally unchanged.  The new block is finite triangular with an
invertible diagonal quotient block by admissibility.  Hence the full map and
its local inverse are block triangular, with the old map as their upper-left
block.  Projection onto that block gives \eqref{eq:v7-theta-projective} and its
inverse analogue.
\end{proof}

\begin{theorem}[Reduced shell maps form a projective system]
\label{thm:v7-reduced-projective}
Let
\[
 \widehat{\mathcal R}_{j,L}^{\mathcal S}
 =\Theta_{j-1,L}^{\mathcal S}
  \mathcal R_{j}^{[\le L]}
  (\Theta_{j,L}^{\mathcal S})^{-1}.
\]
Then
\begin{equation}
 \boxed{
 \pi_L^{L+1}\widehat{\mathcal R}_{j,L+1}^{\mathcal S}
 \iota_L^{L+1}
 =\widehat{\mathcal R}_{j,L}^{\mathcal S}.}
 \label{eq:v7-reduced-projective}
\end{equation}
\end{theorem}

\begin{proof}
Insert the definition, commute $\pi_L^{L+1}$ through the two normalization
maps using Lemma~\ref{lem:v7-theta-projective}, and use
Corollary~\ref{cor:v7-shell-projective}.
\end{proof}

\section{Finite scheme changes as exact intertwiners}
\label{sec:v7-scheme-intertwiner}

We now compare two coherent admissible schemes $\mathcal S$ and $\mathcal S'$
on the same native regulator branch and protected chart.  ``Admissible'' here
means that both use complete local/source/line banks at the prescribed
filtration, their normalization Jacobians are uniformly nonsingular on the
chosen compact set, their BV-sector coordinate changes are canonical with the
corresponding density/reference transformation, and their source/composite
changes are finite triangular maps retaining all contact coordinates.

Define the same-cutoff transition
\begin{equation}
 F_{j,L}^{\mathcal S'\leftarrow\mathcal S}
 :=\Theta_{j,L}^{\mathcal S'}
   (\Theta_{j,L}^{\mathcal S})^{-1}.
 \label{eq:v7-Fj}
\end{equation}
It is a finite local/source/line map.  On a sufficiently small common compact
chart it and its inverse have uniformly bounded $C^1$ norm because both
normalization maps have uniformly nonsingular finite Jacobians.

\begin{theorem}[Exact shell intertwining of two schemes]
\label{thm:v7-intertwiner}
For every shell,
\begin{equation}
 \boxed{
 \widehat{\mathcal R}_{j,L}^{\mathcal S'}
 F_{j,L}^{\mathcal S'\leftarrow\mathcal S}
 =F_{j-1,L}^{\mathcal S'\leftarrow\mathcal S}
  \widehat{\mathcal R}_{j,L}^{\mathcal S}.}
 \label{eq:v7-intertwiner}
\end{equation}
\end{theorem}

\begin{proof}
Using the definitions and cancelling adjacent normalization maps,
\begin{align*}
 \widehat{\mathcal R}_{j,L}^{\mathcal S'}F_{j,L}
 &=\Theta_{j-1,L}^{\mathcal S'}
   \mathcal R_j
   (\Theta_{j,L}^{\mathcal S'})^{-1}
   \Theta_{j,L}^{\mathcal S'}
   (\Theta_{j,L}^{\mathcal S})^{-1}\\
 &=\Theta_{j-1,L}^{\mathcal S'}
   \mathcal R_j(\Theta_{j,L}^{\mathcal S})^{-1}\\
 &=F_{j-1,L}
   \Theta_{j-1,L}^{\mathcal S}\mathcal R_j
   (\Theta_{j,L}^{\mathcal S})^{-1},
\end{align*}
which is \eqref{eq:v7-intertwiner}.
\end{proof}

\begin{lemma}[Summable adjacent variation of the finite scheme map]
\label{lem:v7-F-adjacent}
At fixed finite $(L,\Delta,r,N)$,
\begin{equation}
 \boxed{
 \|F_{j-1,L}-F_{j,L}\|_{C^1}
 \le C_{L,\Delta,r,N}\epsilon_j.}
 \label{eq:v7-F-adjacent}
\end{equation}
\end{lemma}

\begin{proof}
Write the two V6 reduced shell maps as
$\widehat{\mathcal R}^{\mathcal S}_j=I+E_j$ and
$\widehat{\mathcal R}^{\mathcal S'}_j=I+E'_j$, with
$\|E_j\|_{C^1}+\|E'_j\|_{C^1}\le C\epsilon_j$.
Theorem~\ref{thm:v7-intertwiner} gives
\[
 F_{j-1,L}=(I+E'_j)F_{j,L}(I+E_j)^{-1}.
\]
On the common compact chart the maps $F_{j,L}$ and their inverses are uniformly
$C^1$ bounded, while
$(I+E_j)^{-1}=I+O_{C^1}(\epsilon_j)$ for all sufficiently large $j$.
Subtract $F_{j,L}$ and expand the finite composition.  Every resulting term
contains at least one factor $E_j$ or $E'_j$, proving
\eqref{eq:v7-F-adjacent}.  Finitely many low shells are absorbed into the
constant.
\end{proof}

\begin{theorem}[Invertible continuum scheme transition]
\label{thm:v7-F-limit}
The finite scheme maps have a cutoff-independent $C^1$ limit
\begin{equation}
 \boxed{
 F_{\infty,L}^{\mathcal S'\leftarrow\mathcal S}
 :=\lim_{j\to\infty}F_{j,L}^{\mathcal S'\leftarrow\mathcal S},}
 \label{eq:v7-F-limit}
\end{equation}
with tail
\begin{equation}
 \|F_{\infty,L}-F_{j,L}\|_{C^1}
 \le C_{L,\Delta,r,N}
 \sum_{k>j}\epsilon_k
 \le C'_{L,\Delta,r,N}\frac{\mu}{\Lambda_j}.
 \label{eq:v7-F-tail}
\end{equation}
The limit is locally invertible, and
\begin{equation}
 (F_{\infty,L}^{\mathcal S'\leftarrow\mathcal S})^{-1}
 =F_{\infty,L}^{\mathcal S\leftarrow\mathcal S'}.
 \label{eq:v7-F-inverse}
\end{equation}
\end{theorem}

\begin{proof}
Lemma~\ref{lem:v7-F-adjacent} and
$\sum_j\epsilon_j<\infty$ make $\{F_{j,L}\}$ a $C^1$ Cauchy sequence.
Apply the same argument to the inverse transitions
$G_{j,L}:=F_{j,L}^{\mathcal S\leftarrow\mathcal S'}$.  At every finite cutoff,
$G_{j,L}F_{j,L}=I=F_{j,L}G_{j,L}$.  Passing to the uniform $C^1$ limits gives
$G_{\infty,L}F_{\infty,L}=I=F_{\infty,L}G_{\infty,L}$ on the common local
chart.  The geometric tail bound is the same sum used in V6.
\end{proof}

\begin{corollary}[Cocycle law for continuum scheme transitions]
\label{cor:v7-scheme-cocycle}
For three coherent admissible schemes $\mathcal S,\mathcal S',\mathcal S''$,
\begin{equation}
 F_{\infty,L}^{\mathcal S''\leftarrow\mathcal S'}
 F_{\infty,L}^{\mathcal S'\leftarrow\mathcal S}
 =F_{\infty,L}^{\mathcal S''\leftarrow\mathcal S}.
 \label{eq:v7-scheme-cocycle}
\end{equation}
\end{corollary}

\begin{proof}
The identity holds exactly at every finite cutoff because it is cancellation of
three coordinate identifications.  Pass to the $C^1$ limit.
\end{proof}

\section{What is scheme dependent and what is invariant}
\label{sec:v7-scheme-invariants}

The continuum state is therefore not a preferred numerical coordinate vector;
it is an equivalence class under the finite maps
$F_{\infty,L}$.  It is useful to state the invariant content precisely.

\begin{definition}[Covariant continuum observable]
\label{def:v7-covariant-observable}
A family of coefficient functions $\mathcal O^{\mathcal S}$ is covariant if
for every pair of admissible schemes
\begin{equation}
 \mathcal O^{\mathcal S'}\circ
 F_{\infty,L}^{\mathcal S'\leftarrow\mathcal S}
 =\mathcal O^{\mathcal S}.
 \label{eq:v7-covariant-observable}
\end{equation}
A relation is scheme invariant if its zero set is carried into itself by all
such transitions.
\end{definition}

\begin{theorem}[Finite-scheme covariance of the continuum coefficient system]
\label{thm:v7-scheme-covariance}
At fixed finite filtration the following hold.
\begin{enumerate}[label=\textup{(SC\arabic*)},leftmargin=3.4em]
\item Independent local EFT/Wilson coefficients, composite-operator basis
coordinates, finite stress improvements, redundant/BV-exact coordinates,
coincident contact coefficients and scalar line/reference coordinates are in
general scheme dependent.
\item The complete normalized connected functional is a scalar after the
source/composite coordinates and local contact polynomial are transformed by
$F_{\infty,L}$.  Hence connected coefficients of matched renormalized
observables transform covariantly.
\item The common BV/BRST/master vector transforms by the induced canonical
ghost-one coordinate map.  Therefore its vanishing, the Standard-Model anomaly
cancellation, and every source/contact descendant relation proved in V5--V6
are scheme invariant.
\item Any separated-support connected relation, after the operator/source
basis is matched covariantly, is insensitive to a change consisting only of a
finite local source/contact polynomial.  Raw coincident contact coefficients
need not agree.
\end{enumerate}
\end{theorem}

\begin{proof}
Item (SC1) is the definition of a finite change of local coordinates: there is
no reason for coordinate components in two charts to agree numerically.
For (SC2), at finite cutoff both schemes are obtained from the same native
partition function by invertible source/local coordinate changes; the full
contact packet is included among those coordinates.  Thus the scalar
coordinate-free connected functional is unchanged, while its component array
is transformed by substitution and the finite local contact term.  Passing to
the V6 coefficientwise limit and using Theorem~\ref{thm:v7-F-limit} proves the
continuum statement.

For (SC3), admissibility requires the BV-sector transition to be canonical and
the finite BV density/reference term to be transformed simultaneously.  The
master functional therefore transforms covariantly at finite cutoff.  Its
retained coefficient vector is zero in one scheme if and only if it is zero in
the other.  The same is true after any retained source derivative because the
source/contact map is part of $F_{j,L}$.  Pass to the limit.  The anomaly
quotient is a cohomological coordinate-free obstruction, so its zero class is
unchanged by finite basis changes.

For (SC4), the difference generated by a local source polynomial is supported
on coincident source points after differentiation.  If the test/source
supports are pairwise separated, those diagonal distributions vanish.  A
change of composite basis is still applied contragrediently to the inserted
observables, which is precisely the matching in
\eqref{eq:v7-covariant-observable}.
\end{proof}

\section{Determinant/Pfaffian frame and positive-reference compatibility}
\label{sec:v7-line-frame}

The line sector deserves a separate statement because a scalar representative
of a determinant section is not itself coordinate free.

Let $\lambda$ and $\lambda'$ be two admissible nonzero local determinant-line
frames on the same protected chart and write
\begin{equation}
 \lambda'=e^{\chi}\lambda,
 \label{eq:v7-line-frame-change}
\end{equation}
where $\chi$ is a smooth local source function; for two unit frames,
$\Re\chi=0$.  If
$\sigma=s\lambda=s'\lambda'$, then
\begin{equation}
 \boxed{\log s'=\log s-\chi.}
 \label{eq:v7-line-scalar-change}
\end{equation}
For an oriented nonzero Pfaffian line the same formula holds on the chosen
orientation chart, with the inherited factor-$1/2$ relation to the determinant
logarithm where applicable.

\begin{theorem}[Frame changes are exact local coboundaries of the shell line cocycle]
\label{thm:v7-line-coboundary}
Under an admissible local line-frame change, the coordinate-free exact Schur
shell cocycle is unchanged.  Its scalar logarithmic representatives differ by
an endpoint coboundary.  Consequently a finite sequence of shell frame changes
contributes only the difference of the endpoint frame logs, together with all
of their retained source derivatives; no mixed-shell contact is lost.
\end{theorem}

\begin{proof}
The Schur factorization is an isomorphism of determinant lines before a frame
is selected.  Therefore changing frames cannot alter the coordinate-free
section.  Scalarizing the line identity in the two frames and using
\eqref{eq:v7-line-scalar-change} shows that the logarithmic shell factor changes
by the difference of the frame logs at its two endpoints.  Summing over nested
shells telescopes these differences.  Differentiation uses the ordinary
Leibniz/Faa--di--Bruno rule already encoded in the V1 Bell-polynomial line
packet, so every derivative of the endpoint logs and every mixed contact is
retained rather than discarded.
\end{proof}

Let a positive local reference normalization be changed by
\begin{equation}
 \rho'=e^{u}\rho,
 \label{eq:v7-positive-reference-change}
\end{equation}
with real smooth local source function $u$.  Source-free $u(0)$ disappears from
$W[J]=\log Z[J]-\log Z[0]$, but the difference $u(J)-u(0)$ is a genuine local
source/contact coordinate.

\begin{corollary}[Normalized line/reference functional]
\label{cor:v7-line-reference-W}
A simultaneous admissible line-frame/reference change alters a scalar
representative of the normalized connected functional only by the finite local
source/contact expression prescribed by
$\chi(J)-\chi(0)$ and $u(J)-u(0)$, with signs determined by the declared scalar
line convention.  After those contact coordinates are transformed as part of
$F_{\infty,L}$, the coordinate-free normalized determinant/Pfaffian insertion
relations are unchanged.
\end{corollary}

\begin{proof}
Apply Theorem~\ref{thm:v7-line-coboundary} to the determinant/Pfaffian factor
and subtract the $J=0$ normalization.  A source-free constant cancels, whereas
a source-dependent local function does not.  V1--V3 explicitly store its
source derivatives in the line/reference packet.  The continuum transition of
Theorem~\ref{thm:v7-F-limit} transports precisely those finite coordinates.
\end{proof}

\begin{firewall}[No global line-holonomy conclusion]
Theorems~\ref{thm:v7-line-coboundary}--\ref{cor:v7-line-reference-W} compare
local nonzero frames on the protected component.  They do not trivialize the
line bundle across zero crossings or large loops and do not prove global
 determinant/Pfaffian holonomy cancellation beyond the finite gauge-descent
statement already inherited.
\end{firewall}

\section{Projective compatibility of continuum limits}
\label{sec:v7-continuum-projective}

For a coherent scheme $\mathcal S$, let
$\widehat{\mathfrak C}_{L,j_0}^{(J),\mathcal S}$ be the V6 normalized state at
reference shell $j_0$ obtained with UV cutoff $J$, and let
$\widehat{\mathfrak C}_{L,j_0}^{(\infty),\mathcal S}$ denote its coefficientwise
limit.

\begin{theorem}[Finite-cutoff states are exactly projective]
\label{thm:v7-finite-state-projective}
For every $J$ and $j_0\le J$,
\begin{equation}
 \boxed{
 \pi_L^{L+1}
 \widehat{\mathfrak C}_{L+1,j_0}^{(J),\mathcal S}
 =
 \widehat{\mathfrak C}_{L,j_0}^{(J),\mathcal S}.}
 \label{eq:v7-finite-state-projective}
\end{equation}
\end{theorem}

\begin{proof}
The normalized UV boundary data are coherent by Definition
\ref{def:v7-coherent-scheme}.  Compose the reduced shell maps down to $j_0$.
Theorem~\ref{thm:v7-reduced-projective} lets $\pi_L^{L+1}$ pass through each
factor and turns every order-$(L+1)$ factor into its order-$L$ factor.  The
projected UV boundary is the order-$L$ boundary.  This proves the identity.
\end{proof}

\begin{theorem}[Exact projective compatibility in perturbative order]
\label{thm:v7-continuum-projective}
For every coherent admissible scheme,
\begin{equation}
 \boxed{
 \pi_L^{L+1}
 \widehat{\mathfrak C}_{L+1,j_0}^{(\infty),\mathcal S}
 =
 \widehat{\mathfrak C}_{L,j_0}^{(\infty),\mathcal S}.}
 \label{eq:v7-continuum-projective}
\end{equation}
The same identity holds coefficientwise for the normalized connected
functional, every retained source/contact descendant, and the complete finite
BV/master vector.
\end{theorem}

\begin{proof}
The projection $\pi_L^{L+1}$ is a bounded finite-coordinate map.  Apply it to
the V6 Cauchy limit and use the exact finite-cutoff identity
\eqref{eq:v7-finite-state-projective}:
\[
 \pi_L^{L+1}
 \widehat{\mathfrak C}_{L+1}^{(\infty)}
 =\lim_{J\to\infty}
  \pi_L^{L+1}\widehat{\mathfrak C}_{L+1}^{(J)}
 =\lim_{J\to\infty}\widehat{\mathfrak C}_{L}^{(J)}
 =\widehat{\mathfrak C}_{L}^{(\infty)}.
\]
The connected/source/master coefficient maps are finite polynomial or analytic
maps of the retained state at every prescribed finite order, so they commute
with the same projection and limit.
\end{proof}

\begin{theorem}[Scheme transitions are themselves projective in loop order]
\label{thm:v7-F-projective}
For two coherent schemes,
\begin{equation}
 \boxed{
 \pi_L^{L+1}
 F_{\infty,L+1}^{\mathcal S'\leftarrow\mathcal S}
 \iota_L^{L+1}
 =F_{\infty,L}^{\mathcal S'\leftarrow\mathcal S}.}
 \label{eq:v7-F-projective}
\end{equation}
\end{theorem}

\begin{proof}
The same identity holds at every finite cutoff by
Lemma~\ref{lem:v7-theta-projective} applied to both schemes.  The $C^1$ limits
of Theorem~\ref{thm:v7-F-limit} commute with the bounded finite projection.
\end{proof}

\section{Reference-scale covariance at positive scale}
\label{sec:v7-reference-scale}

A change of subtraction/reference point inside a compact positive interval is
an admissible finite scheme change provided the corresponding normalization
Jacobians remain nonsingular and the protected line/pivot component is not
crossed.  Therefore the preceding theorems also compare two positive reference
scales $\mu,\mu'>0$ under those hypotheses.

\begin{corollary}[Positive reference-scale covariance]
\label{cor:v7-reference-scale}
Let $\mu$ and $\mu'$ belong to a compact interval
$0<\mu_-\le\mu,\mu'\le\mu_+<\infty$ on which the protected chart and finite
normalization ranks persist.  The fixed-order continua normalized at $\mu$ and
$\mu'$ are related by an invertible finite continuum map
$F_{\infty,L}^{\mu'\leftarrow\mu}$ satisfying the cocycle law in the reference
point.  This does not imply existence of the limit $\mu\downarrow0$.
\end{corollary}

\begin{proof}
Changing the positive reference point changes only the finite normalization
conditions and the resulting local/source/line coordinates.  On the compact
positive interval their Jacobians are uniformly nonsingular by hypothesis, so
they form an admissible scheme pair.  Apply Theorem~\ref{thm:v7-F-limit} and
Corollary~\ref{cor:v7-scheme-cocycle}.  No estimate is uniform as
$\mu_-\downarrow0$.
\end{proof}

\section{E6 closure theorem}
\label{sec:v7-E6-close}

\begin{theorem}[Scheme/reference and perturbative-order compatibility]
\label{thm:v7-E6}
Fix the native V2 Wilson/overlap branch, a protected constant-rank/nonzero-line
component, positive IR/reference scale, finite external/source panel, and
finite $(L,\Delta,r,N)$.  Then:
\begin{enumerate}[label=\textup{(E6.\arabic*)},leftmargin=3.5em]
\item any two coherent admissible finite normalization/source/reference
schemes are related by an invertible cutoff-independent continuum transition
$F_{\infty,L}$ with the finite cocycle law;
\item determinant/Pfaffian local frame changes and positive reference
normalizations act by the stored local endpoint/contact coboundaries and leave
the coordinate-free normalized line relations unchanged;
\item scheme-dependent local coefficient arrays transform covariantly, while
the vanishing common BV/BRST/master relations, Standard-Model anomaly-zero
class, and matched normalized connected relations are invariant;
\item the native shell maps, normalized reduced transports, continuum states,
connected/source/master coefficients, and continuum scheme transitions form a
projective family in perturbative order, with exact restriction from $L+1$ to
$L$ and no lower-order refit.
\end{enumerate}
\end{theorem}

\begin{proof}
Item (1) is Theorems~\ref{thm:v7-intertwiner}--\ref{thm:v7-F-limit} and
Corollary~\ref{cor:v7-scheme-cocycle}.  Item (2) is
Theorem~\ref{thm:v7-line-coboundary} and
Corollary~\ref{cor:v7-line-reference-W}.  Item (3) is
Theorem~\ref{thm:v7-scheme-covariance}.  Item (4) is
Corollary~\ref{cor:v7-shell-projective},
Theorem~\ref{thm:v7-reduced-projective},
Theorem~\ref{thm:v7-continuum-projective}, and
Theorem~\ref{thm:v7-F-projective}.
\end{proof}

\begin{center}
\fbox{\parbox{0.92\textwidth}{
\textbf{E6 is closed at the stated coefficientwise scope.}
The essential upstream correction is chronological/projective normalization:
finite Hodge inverses are not independently recomputed when the perturbative
bank is enlarged.  With lower-order data frozen, both the native shell theorem
and its continuum limit restrict exactly in loop order.}}
\end{center}

\section{Final source-complete coefficientwise Renewal ESM continuum theorem}
\label{sec:v7-final-theorem}

We can now state the endpoint promised at the beginning of this lineage.  The
quantifiers are important: $L$, $\Delta$, $r$, $N$ and the external/source
panel are fixed before the cutoff is removed.  Constants may depend
arbitrarily badly on these finite choices.

\begin{theorem}[Source-complete coefficientwise perturbative Einstein--Standard Model continuum]
\label{thm:v7-final-continuum}
Assume:
\begin{enumerate}[label=\textup{(H\arabic*)},leftmargin=3em]
\item the explicit V2 native Wilson/overlap fermion completion, or another
microscopic fermion regulator for which the V2 native regulator-comparison gate
has separately been proved;
\item the inherited finite Renewal gravitational/gauge/Higgs regulator and
exact BV-compatible high/low splitting on a protected constant-rank
perturbative component;
\item the three-generation Standard-Model chiral representation with the
inherited cancellation of the ordinary local gauge-anomaly trace classes and
even $SU(2)$ doublet parity;
\item a protected nonzero determinant/Pfaffian line component and a positive
IR/reference scale $\mu>0$;
\item one coherent admissible normalization/source/reference family.
\end{enumerate}
Then for every prescribed finite loop order $L$, EFT/derivative order
$\Delta$, source/contact order $r$, arity $N$, and finite external
momentum/source panel, there exists a common renormalized native action/source/
reference/line family whose normalized connected coefficients have a unique
ultraviolet-cutoff-independent limit.  More precisely:
\begin{enumerate}[label=\textup{(F\arabic*)},leftmargin=3em]
\item every loop-$\le L$ native graph and forest is renormalized by the same
previously fixed common lower-order actions and decomposes into a finite local
EFT/source jet plus a quasilocal summable remainder;
\item the limit includes composite operators, operator mixing, stress-tensor
insertions, coincident source contacts, ghosts/nonminimal variables,
antifields, and determinant/Pfaffian line insertions at the prescribed finite
orders;
\item all retained coefficients of the common coupled BV/BRST/master/source
system and all contraction descendants required through loop order $L$ vanish
in the continuum, with the nontrivial local Standard-Model anomaly class
absent;
\item the finite-volume/nonwinding coefficients admit the separate V6
thermodynamic limit, and the staged UV and volume limits commute on the
prescribed compact panels;
\item different coherent admissible finite schemes/reference choices are
related by the invertible continuum maps of Theorem~\ref{thm:v7-E6}; raw local
EFT, contact, composite-basis and line-frame coordinates may change, while the
matched coordinate-free normalized relations transform covariantly;
\item the family is exactly compatible in perturbative order:
\begin{equation}
 \pi_L^{L+1}
 \widehat{\mathfrak C}_{L+1}^{(\infty)}
 =\widehat{\mathfrak C}_{L}^{(\infty)}.
 \label{eq:v7-final-projective}
\end{equation}
\end{enumerate}
Gravity is interpreted throughout as an effective field theory: at every fixed
$(L,\Delta)$ the required local Wilson bank is finite, but no finite-dimensional
all-orders gravitational coupling space is asserted.
\end{theorem}

\begin{proof}
E1 is supplied by V1--V2: native chiral hard-shell functional calculus,
hard-pivot/annular/full-zone comparison, fixed source jets and the complete
 determinant/Pfaffian line remainder.  E2 is V3: the complete native
Standard-Model anomaly projection vanishes and one common source-extended
primitive enlarges the V80 bosonic shell to the chiral ESM shell.  E3 is V4:
one common first- and second-loop master/source equation contains every
canonical-pair cross term and every retained contact.  E4 is V5: arbitrary
fixed-loop forest preparation, contraction-complete packets, exact local
cocycles and one common order-$L$ restorer.  E5 is V6: after factoring the
finite local normalization cocycle, the reduced shell maps form a summable
near-identity product, giving the unique coefficientwise UV limit and the
separate thermodynamic comparison.  E6 is
Theorem~\ref{thm:v7-E6}: finite schemes/reference choices are related by
invertible continuum maps and the construction is projective in loop order.
Combining these statements yields (F1)--(F6).
\end{proof}

\begin{firewall}[What the final theorem does not say]
Theorem~\ref{thm:v7-final-continuum} is coefficientwise.  It does not assert:
\begin{enumerate}[leftmargin=2.4em]
\item convergence or Borel summability of the perturbative series as
$L\to\infty$;
\item any estimate uniform in loop order;
\item finite-dimensional renormalizability of Einstein gravity;
\item asymptotic safety or a nonperturbative UV fixed point;
\item existence of a nonperturbative quantum-gravity/ESM measure;
\item removal of the positive IR/reference scale or a global massless IR
limit;
\item determinant/Pfaffian triviality across zero crossings or arbitrary large
line-bundle loops;
\item universality with respect to a different microscopic fermion regulator
unless the V2 regulator-comparison gate is separately proved.
\end{enumerate}
\end{firewall}

\section{A projective formulation of the endpoint}
\label{sec:v7-projective-endpoint}

For each finite $L$ let
$\mathfrak C_L^{\infty,\mathcal S}$ be the continuum state supplied by
Theorem~\ref{thm:v7-final-continuum}.  Theorem
\ref{thm:v7-continuum-projective} makes
\begin{equation}
 \left\{
 \mathfrak C_L^{\infty,\mathcal S},
 \pi_L^{L+1}
 \right\}_{L\ge0}
 \label{eq:v7-projective-family}
\end{equation}
a projective family of finite-order coefficient systems.  This notation is
occasionally convenient, but it must not be misread as constructing the
inverse-limit point as a convergent numerical perturbative series.  The
projective family means only that every finite order is compatible with every
lower finite order.

\begin{corollary}[Order-independent meaning of a fixed coefficient]
\label{cor:v7-fixed-coefficient}
Fix a loop index $q$, derivative/source label in the chosen finite bank, and an
external/source panel.  For every terminal order $L\ge q$, the corresponding
continuum coefficient is exactly the same coefficient already fixed at order
$q$.
\end{corollary}

\begin{proof}
Iterate \eqref{eq:v7-final-projective} from $L$ down to $q$ and read the chosen
coordinate.
\end{proof}

\section{V7 anti-circularity audit}
\label{sec:v7-audit}

\begin{enumerate}[leftmargin=2.4em]
\item \textbf{No independently recomputed lower-order Hodge inverse.}
Proposition~\ref{prop:v7-hodge-not-functorial} exhibits the obstruction.
Lower-order primitives are frozen and only the new loop layer is solved.
\item \textbf{No $L$-dependent rediscovery of the EFT basis.}
The ambient operator/source/line bank is fixed at $(\Delta,r,N)$; loop order is
a formal grading.  This is what makes the truncation maps literal projections.
\item \textbf{No graphwise finite retuning at higher order.}
Theorem~\ref{thm:v7-forest-projective} uses V5's strict proper-subgraph loop
drop, so all proper finite subtractions are owned by already fixed lower-order
actions.
\item \textbf{No claim that different schemes have equal coordinate arrays.}
They are related by the finite invertible maps $F_{\infty,L}$.  Local Wilson
coefficients, contact terms and composite bases may differ.
\item \textbf{No discarded source-dependent frame/reference term.}
Only a source-free normalization constant cancels in $W[J]-W[0]$.
$\chi(J)-\chi(0)$ and $u(J)-u(0)$ remain explicit contact coordinates.
\item \textbf{No global determinant-line conclusion.}
Frame compatibility is local on the protected nonzero component.
\item \textbf{No hidden massless limit.}
Reference-scale covariance is proved only on compact positive intervals.
\item \textbf{No inverse-limit convergence claim.}
Projectivity in $L$ is exact compatibility of finite coefficients, not
convergence of the perturbative series.
\end{enumerate}

\section{V7 status ledger}
\label{sec:v7-status}

\subsection*{Proved in V7}
\begin{enumerate}[leftmargin=2.2em]
\item Independently recomputed orthogonal Hodge inverses are not projective in
general; a finite explicit counterexample identifies the missing hypothesis.
\item A chronological coherent normalization/restoration family can be
extended to every higher prescribed finite loop order without altering any
previous coefficient.
\item With a fixed ambient EFT/source/line bank and the V5
contraction-complete packets, forest preparation, common restoring actions and
native shell maps commute exactly with loop-order truncation.
\item The V6 normalized coordinate maps and reduced shell transports can be
chosen projectively in loop order.
\item Two coherent admissible schemes are exact finite-cutoff intertwiners.
Their adjacent-cutoff transition maps differ by $O(\mu/\Lambda_j)$, hence
converge in $C^1$ to an invertible continuum transition satisfying the cocycle
law.
\item Scheme-dependent local coefficient arrays and scheme-invariant/covariant
relations are separated explicitly.  The common master/anomaly-zero relations
are invariant; matched connected observables are covariant; raw coincident
contact coefficients need not agree.
\item Determinant/Pfaffian local frame changes are endpoint coboundaries of the
exact line cocycle.  Positive reference changes retain every source-dependent
contact; only source-free normalization cancels in $W[J]/W[0]$.
\item Finite-cutoff normalized states and their V6 continuum limits are exactly
projective in perturbative order, as are the continuum scheme transitions.
\item Positive reference points in one compact protected interval are related
by finite invertible continuum maps, without claiming a $\mu\downarrow0$
limit.
\item E6 and therefore the stated E1--E6 coefficientwise programme are closed
on the declared native branch and assumptions.
\end{enumerate}

\subsection*{Inherited}
\begin{enumerate}[leftmargin=2.2em]
\item V80's bosonic Einstein--gauge--Higgs two-loop local/self-map theorem and
finite local BV range architecture.
\item The Standard-Model representation, cancellation of the five ordinary
local gauge-anomaly coefficients, even $SU(2)$ doublet parity, protected gauge
descent, and exact determinant/Pfaffian scale cocycles.
\item V1--V2's native chiral functional calculus, local Kato frame, hard pivot,
explicit Wilson/overlap completion, annular/full-zone comparison, source jets
and line remainder bounds.
\item V3's complete chiral anomaly-quotient zero and common ESM two-loop shell
restorer; V4's full common two-loop master/source audit.
\item V5's arbitrary-fixed-loop native forest theorem,
contraction-complete packet and fixed-order common master restoration.
\item V6's normalized multi-shell $C^1$ product, coefficientwise UV Cauchy
theorem, connected/source/master continuum, and separate thermodynamic
comparison.
\end{enumerate}

\subsection*{Conditional}
\begin{enumerate}[leftmargin=2.2em]
\item A microscopic fermion kernel different from the explicit V2 completion
still requires its own native regulator-comparison theorem before this closure
transfers to it.
\item Constant-rank/Hodge and normalization arguments remain local to compact
protected charts.  Rank changes require compatible chart transitions.
\item Determinant/Pfaffian results remain on a nonzero locally orientable
component and do not cross zeros.
\item The positive reference-scale comparison requires a compact interval
bounded away from zero; a massless IR limit is not proved.
\item Constants can depend arbitrarily on $(L,\Delta,r,N)$ and on the protected
chart; no all-orders uniformity is available.
\end{enumerate}

\subsection*{Open beyond the stated programme}
\begin{enumerate}[leftmargin=2.2em]
\item A native universality theorem comparing the explicit V2 fermion
completion with a genuinely different microscopic chiral regulator.
\item Removal of the retained positive IR/reference scale, if desired, with
uniform control of massless long-distance sectors.
\item Global determinant/Pfaffian line holonomy and zero-crossing analysis.
\item Any analytic finite-coupling RG trajectory, all-orders summability,
asymptotic-safety or nonperturbative quantum-gravity/ESM completion theorem.
These are strictly stronger problems and are not required for the
coefficientwise endpoint proved here.
\end{enumerate}

\subsection*{Exact next load-bearing theorem}
\begin{center}
\fbox{\parbox{0.92\textwidth}{
\textbf{None inside the stated E1--E6 coefficientwise programme.}
The source-complete fixed-order continuum, scheme/reference covariance and
projective compatibility in loop order are now closed under the declared
native-regulator, protected-chart and positive-IR assumptions.  Any further
attack should be labelled as a stronger programme.  The smallest natural one
is a regulator-universality theorem comparing the explicit V2
Wilson/overlap completion with an alternative microscopic chiral regulator;
IR removal and global line holonomy are independent stronger directions.}}
\end{center}

\section*{Append-only continuation marker: V7}
\addcontentsline{toc}{section}{Append-only continuation marker: V7}
\textbf{End of V7.}  The complete V1--V6 body above is retained verbatim.  V7
closes E6 by replacing nonfunctorial order-by-order Hodge recomputation with a
chronological projective construction, proving exact finite-cutoff and
continuum restriction in loop order, constructing invertible continuum
scheme/reference transitions, retaining all determinant/Pfaffian frame and
source-reference contacts, and stating the final source-complete
coefficientwise Renewal Einstein--Standard Model continuum theorem.  No
load-bearing gate remains inside the stated E1--E6 programme.

\clearpage
\section{V8 continuation: a stronger post-E1--E6 programme --- microscopic regulator universality}
\label{sec:v8-continuation}

V7 closed the stated E1--E6 coefficientwise continuum programme on the explicit
V2 Wilson/overlap branch.  The present continuation therefore does not reopen
any gate in that programme.  It attacks the first stronger question isolated in
the V7 ledger: whether the continuum just constructed depends on the
microscopic fermion completion used to realize the protected chiral shell.

The claim proved below is deliberately narrower than arbitrary regulator
independence.  We first formulate a local gapped-homotopy criterion and then
verify it for a genuinely different finite-range stencil obtained by adding a
range-two covariant Wilson deformation to the V2 range-one kernel.  At every
fixed loop/EFT/source order the two regulators are shown to differ by one finite
local scheme/source/line transformation plus a quasilocal remainder which
vanishes in the ultraviolet limit.  No statement is made between disconnected
Wilson phases, across determinant-line zeros, or for a regulator for which the
required gapped homotopy has not been verified.

\subsection{The relative question and the correct notion of universality}

Fix throughout a finite tuple
\begin{equation}
 \mathfrak q=(L,\Delta,r,N,\mathcal P_{\rm ext}),
 \label{eq:v8-fixed-data}
\end{equation}
where $\mathcal P_{\rm ext}$ is the finite external momentum/source panel, and
retain the positive reference scale $\mu>0$.  Let
$\mathfrak C^{\rm ESM}_{\mathfrak q}$ be the V5--V7
contraction-complete state at these fixed data.

\begin{definition}[Matched regulator universality at fixed order]
\label{def:v8-universality}
Two microscopic regulator families $A$ and $B$ are \emph{matched-universal at
$\mathfrak q$} if there are finite local invertible maps
\begin{equation}
 F_{a,\mathfrak q}^{B\leftarrow A}:
 \mathfrak C^{A}_{a,\mathfrak q}\longrightarrow
 \mathfrak C^{B}_{a,\mathfrak q}
 \label{eq:v8-finite-reg-map}
\end{equation}
which preserve the chosen physical normalization coordinates and transport the
composite/source/line basis covariantly, such that for the normalized connected
states at a fixed physical scale $\lambda$,
\begin{equation}
 \left\|
 \widehat{\mathfrak C}^{B}_{a,\lambda}
 -F_{a,\mathfrak q}^{B\leftarrow A}
  \widehat{\mathfrak C}^{A}_{a,\lambda}
 \right\|_{\mathfrak q}
 \xrightarrow[a\downarrow0]{}0.
 \label{eq:v8-universality-def}
\end{equation}
The maps must be projective in perturbative order in the sense of V7.
\end{definition}

This is the appropriate statement for an EFT.  Equality of bare Wilson
coordinates is neither expected nor required.  What is required is equality of
the normalized connected/source theory after one finite local matching map.

\begin{firewall}[Universality is not equality of bare actions]
A change of microscopic regulator generally changes power-divergent and finite
local coefficients.  Requiring those bare coordinates themselves to converge to
the same values would contradict ordinary scheme dependence.  The comparison
must first factor the finite local normalization cocycle, exactly as V6 factored
ordinary marginal running before proving the ultraviolet Cauchy theorem.
\end{firewall}

\subsection{An explicit genuinely different native stencil}

Retain the V2 covariant shifts $T_\mu^{\mathcal U}$ and Hermitian Wilson pieces
\begin{equation}
 L_\mu^{\mathcal U}
 =I-\frac12\left(T_\mu^{\mathcal U}+
 (T_\mu^{\mathcal U})^*\right).
 \label{eq:v8-Lmu}
\end{equation}
The V2 kernel is denoted $K_{W,a}^{(0)}$ and
$H_{W,a}^{(0)}=\gamma_5K_{W,a}^{(0)}$.

\begin{definition}[Range-two Wilson deformation]
\label{def:v8-range-two}
Put
\begin{equation}
 Q_a^{\mathcal U}:=
 \sum_{\mu=1}^4\left(L_\mu^{\mathcal U}\right)^2,
 \qquad
 K_{W,a}^{(\eta)}:=K_{W,a}^{(0)}+\eta Q_a^{\mathcal U},
 \qquad
 H_{W,a}^{(\eta)}:=\gamma_5K_{W,a}^{(\eta)}.
 \label{eq:v8-range-two-kernel}
\end{equation}
For $\eta\ne0$ this is a range-two microscopic stencil and is therefore not the
V2 nearest-neighbour kernel.
\end{definition}

\begin{lemma}[Uniform operator bound for the deformation]
\label{lem:v8-Q-bound}
For unitary gauge--spin parallel transports,
\begin{equation}
 \|L_\mu^{\mathcal U}\|\le2,
 \qquad
 \|Q_a^{\mathcal U}\|\le16.
 \label{eq:v8-Q-bound}
\end{equation}
Every fixed gauge/spin source derivative of $Q_a^{\mathcal U}$ is finite range
and obeys the same coefficientwise edge-jet calculus as V2.
\end{lemma}
\begin{proof}
Because $T_\mu^{\mathcal U}$ is unitary,
$\|(T_\mu^{\mathcal U}+(T_\mu^{\mathcal U})^*)/2\|\le1$, whence
$\|L_\mu^{\mathcal U}\|\le2$.  Thus
$\|(L_\mu^{\mathcal U})^2\|\le4$, and summing four directions gives
\eqref{eq:v8-Q-bound}.  Differentiation produces a finite Leibniz sum of
products of the one-edge derivatives already controlled in
Lemma~\ref{lem:v2-link-jets}; products of two range-one kernels have range at
most two.
\end{proof}

\begin{proposition}[Hermiticity, covariance and finite range]
\label{prop:v8-range-two-structure}
Assume the V2 spin transport hypotheses.  Then for every real $\eta$,
$H_{W,a}^{(\eta)}$ is Hermitian, exactly internal-gauge/spin covariant and of
range at most two.  Coframe dependence remains entirely in the V2 Clifford
kinetic term, and Higgs/Yukawa, antifield and determinant-line coordinates
enter exactly as in V2.
\end{proposition}
\begin{proof}
Each $L_\mu^{\mathcal U}$ is Hermitian and commutes with $\gamma_5$ under the
V2 spin-transport hypotheses.  Hence $Q_a^{\mathcal U}$ is Hermitian,
$\gamma_5$-even and transforms by the same unitary similarity as the original
kernel.  The range statement follows from two compositions of range-one
shifts.  The remaining claims are immediate from
\eqref{eq:v8-range-two-kernel}.
\end{proof}

Let $g_{\rm W}>0$ be the protected Wilson gap of the V2 kernel on the compact
reference chart used by V1--V7.  Shrink the chart once, if necessary, so this
same lower bound is valid for all reference source jets used below.

\begin{theorem}[A uniform gapped homotopy]
\label{thm:v8-gapped-path}
Choose
\begin{equation}
 0<\eta_*<\frac{g_{\rm W}}{32}.
 \label{eq:v8-eta-small}
\end{equation}
For $t\in[0,1]$ set $\eta(t)=t\eta_*$.  Then
\begin{equation}
 \boxed{
 |H_{W,a}^{(\eta(t))}|
 \succeq \frac{g_{\rm W}}2 I}
 \label{eq:v8-homotopy-gap}
\end{equation}
uniformly in $t$, cutoff and volume on the protected chart.  In particular the
path does not cross a Wilson-phase boundary.
\end{theorem}
\begin{proof}
By Lemma~\ref{lem:v8-Q-bound},
\[
 \|H_{W,a}^{(\eta(t))}-H_{W,a}^{(0)}\|
 =\eta(t)\|Q_a^{\mathcal U}\|
 \le16\eta_*<\frac{g_{\rm W}}2.
\]
The min--max/Weyl perturbation inequality for singular values gives
\eqref{eq:v8-homotopy-gap}.
\end{proof}

At the flat reference point the additional symbol is
\begin{equation}
 q(x)=\sum_{\mu=1}^4(1-\cos x_\mu)^2=O(|x|^4).
 \label{eq:v8-q-symbol}
\end{equation}
Thus the deformation changes neither the physical first-order Dirac symbol nor
the Standard-Model representation carried by the chiral block.

\subsection{Relative overlap functional calculus along the regulator path}

Write
\begin{equation}
 \varepsilon_t=\operatorname{sgn}H_{W,a}^{(\eta(t))},
 \qquad
 D_{{\rm ov},t}=a^{-1}(I+\gamma_5\varepsilon_t),
 \qquad
 \widehat P_{-,t}=\frac12(I+\varepsilon_t).
 \label{eq:v8-overlap-path}
\end{equation}
Let $D_{\chi,t}$ be the corresponding fixed-frame chiral/Yukawa block, with the
same Yukawa map as in V2.

\begin{lemma}[Quasilocal regulator transgression]
\label{lem:v8-sign-transgression}
For every fixed source multi-index $I$ and every fixed number of momentum
moments, the derivatives
\begin{equation}
 \partial_I\partial_t\varepsilon_t,
 \qquad
 \partial_I\partial_t\widehat P_{-,t}
 \label{eq:v8-sign-derivs}
\end{equation}
are cutoff-uniform quasilocal kernels on the protected chart.  More precisely,
with $\Gamma_\pm$ fixed contours around the positive/negative spectrum,
\begin{equation}
 \partial_t\varepsilon_t
 =\frac{1}{2\pi i}
 \left(\oint_{\Gamma_+}-\oint_{\Gamma_-}\right)
 (z-H_t)^{-1}\dot H_t(z-H_t)^{-1}\,dz,
 \label{eq:v8-sign-contour}
\end{equation}
where $\dot H_t=\eta_*\gamma_5Q_a^{\mathcal U}$.
\end{lemma}
\begin{proof}
The gap in Theorem~\ref{thm:v8-gapped-path} permits $t$-independent contours.
Differentiate the Riesz projectors.  Fixed source derivatives give finite sums
of resolvent words containing source derivatives of $H_t$ and one marked
$\dot H_t$.  The V1 Combes--Thomas/finite-range functional calculus applies
uniformly because the range, gap and finite source order are fixed.  Weighted
kernel moments are therefore uniform in $t$.
\end{proof}

\begin{proposition}[Common Kato frame along the regulator homotopy]
\label{prop:v8-regulator-Kato}
On a star-shaped protected chart there is a unitary two-parameter Kato frame
$U(t,\zeta)$ which simultaneously transports the moving source space
$\Ran\widehat P_{-,t}(\zeta)$ from the V2 base frame.  Every fixed mixed
$(t,\zeta)$ derivative is quasilocal and uniformly bounded in the V1 rooted
norms.
\end{proposition}
\begin{proof}
Apply the Kato ODE first in the $t$ direction and then in the star-shaped source
coordinates, or equivalently solve the path-ordered transport for the smooth
projector family $\widehat P_{-,t}(\zeta)$.  The generator is a commutator of
projector derivatives.  Lemma~\ref{lem:v8-sign-transgression} and the V1 Kato
bounds close the finite derivative algebra.  The parameter rectangle is
contractible, so no global line-holonomy assertion is needed.
\end{proof}

\begin{theorem}[Relative physical-annulus symbol estimate]
\label{thm:v8-relative-symbol}
Fix a source order and momentum-derivative order.  On a physical hard annulus
$|p|\asymp\Lambda_j$ with $a\Lambda_j\le c_0$, the fixed-frame chiral blocks
satisfy
\begin{equation}
 \boxed{
 \left\|
 \partial_p^\alpha\partial_{\mathbf q}^\beta\partial_I
 \bigl(D_{\chi,t}-D_{\chi,0}\bigr)
 \right\|
 \le C_{I,\alpha,\beta}\,
 a\Lambda_j^{\nu_I+1-|\alpha|-|\beta|}.}
 \label{eq:v8-relative-D}
\end{equation}
The estimate is uniform in $t\in[0,1]$.  Consequently, for the hard inverses
$G_{\chi,t}=D_{\chi,t}^{-1}$,
\begin{equation}
 \boxed{
 \left\|
 \partial_p^\alpha\partial_{\mathbf q}^\beta\partial_I
 \bigl(G_{\chi,t}-G_{\chi,0}\bigr)
 \right\|
 \le C_{I,\alpha,\beta}\,
 a\Lambda_j^{-|\alpha|-|\beta|}}
 \label{eq:v8-relative-G}
\end{equation}
after the usual engineering/source stocks are divided out.  Equivalently, the
relative change is $O(a\Lambda_j)$ compared with one ordinary hard fermion
propagator.
\end{theorem}
\begin{proof}
At zero source, \eqref{eq:v8-q-symbol} is $O(|ap|^4)$, hence the claimed weaker
$O(a|p|^2)$ relative overlap estimate follows immediately from the analytic
Riesz functional calculus and the prefactor $a^{-1}$.  For source derivatives,
Lemma~\ref{lem:v2-link-jets} applied to the products
$(L_\mu^{\mathcal U})^2$ shows that the first local covariant jet is unchanged
and that the deformation contributes only to the V2 remainder class.  Mixed
Kato-frame derivatives are controlled by
Proposition~\ref{prop:v8-regulator-Kato}.  This proves
\eqref{eq:v8-relative-D}.

The V2 hard pivot, decreased if necessary by a fixed factor using
Theorem~\ref{thm:v8-gapped-path}, gives
$\|G_{\chi,t}\|\le C/\Lambda_j$.  The resolvent identity
\[
 G_{\chi,t}-G_{\chi,0}
 =-G_{\chi,t}(D_{\chi,t}-D_{\chi,0})G_{\chi,0}
\]
and its fixed differentiated versions then give
\eqref{eq:v8-relative-G}.
\end{proof}

\begin{proposition}[Relative ultraviolet-cap calculus]
\label{prop:v8-relative-full-zone}
On every compact full-zone region $|ap|\ge\delta>0$, both hard inverses are
$a$ times smooth periodic dimensionless symbols.  Their difference and every
fixed source derivative therefore have the same $O(a)$ lattice-scale weighted
row/column bound.  An external source-momentum subtraction gains $a\mu$.
\end{proposition}
\begin{proof}
Use the uniform gap and no-physical-zero compactness argument of V2 for the
whole $t$ interval.  The difference of two $a$-scaled smooth periodic symbols
is again of that form.  Fourier-series summation by parts and the mean-value
theorem in the dimensionless variables $aq$ give the result.
\end{proof}

\subsection{Determinant/Pfaffian transgression along the same path}

On a hard block with nonzero pivot define its scalar representative in the
common Kato frame by $s_{j,t}$.  For an optional neutral Majorana block use the
corresponding Pfaffian representative.

\begin{lemma}[Exact line transgression]
\label{lem:v8-line-transgression}
Along the protected regulator path,
\begin{align}
 \partial_t\log s_{j,t}
 &=\Tr\bigl(G_{j,t}\,\partial_tD_{j,t}\bigr),
 \label{eq:v8-det-transgression}\\
 \partial_t\log\operatorname{Pf}A_{j,t}
 &=\frac12\Tr\bigl(A_{j,t}^{-1}\partial_tA_{j,t}\bigr)
 \label{eq:v8-pf-transgression}
\end{align}
whenever the Pfaffian block is present.  Every fixed source derivative is the
finite differentiated trace word obtained from these identities.  Gauge
variation of the complete Standard-Model product is zero at every $t$.
\end{lemma}
\begin{proof}
The first two formulas are Jacobi's determinant identity and its Pfaffian
analogue on the nonzero component.  Fixed source derivatives are finite by the
reference-pivot theorem.  Exact gauge covariance of $D_{\chi,t}$ and the same
complete Standard-Model representation at every $t$ give the same finite gauge
descent cocycle as in V3; that cocycle is one for the complete generation.
\end{proof}

\begin{theorem}[Local plus quasilocal line transgression]
\label{thm:v8-line-local-plus-rem}
For each fixed source order, the shell line transgression decomposes as
\begin{equation}
 \partial_t\log s_{j,t}
 =\mathsf T_j\partial_t\log s_{j,t}
  +\mathsf Q_j\partial_t\log s_{j,t},
 \label{eq:v8-line-decomp}
\end{equation}
where the first term belongs to the finite local source/reference bank and the
second obeys the same two-zone bounds as the relative hard kernels:
\begin{equation}
 \|\partial_I\mathsf Q_j\partial_t\log s_{j,t}\|
 \le
 \begin{cases}
 C_I a\Lambda_j,& a\Lambda_j\le c_0,\\
 C_I\mu/\Lambda_j,& \Lambda_j\ge M\text{ after local Taylor subtraction}.
 \end{cases}
 \label{eq:v8-line-two-zone}
\end{equation}
The same statement holds for the Pfaffian with the factor $1/2$ in
\eqref{eq:v8-pf-transgression}.
\end{theorem}
\begin{proof}
Differentiate the trace words of Lemma~\ref{lem:v8-line-transgression}.  On the
physical branch at least one factor is the marked relative insertion from
Theorem~\ref{thm:v8-relative-symbol}, which gives the relative
$a\Lambda_j$ gain.  On the high branch use the V2 full-zone source calculus and
subtract the complete finite local Taylor/source jet; the first omitted soft
momentum gives $\mu/\Lambda_j$.  The V1 Bell-polynomial line composition retains
all mixed-shell contacts, so the same decomposition is stable under shell
composition.
\end{proof}

\subsection{The relative forest problem}

A regulator comparison cannot be proved graph by graph with independently
chosen finite parts.  We therefore differentiate the \emph{already normalized
common forest} along the regulator path.

For every $t$, let $\mathbf R_{j,t}^{<L}$ be the V5 forest preparation formed
with the chronological lower-order actions already fixed at that same $t$.
The normalization conditions are held fixed in physical coordinates as $t$
varies.

\begin{definition}[Regulator-marked forest]
\label{def:v8-marked-forest}
For a forest-prepared graph $\Gamma$, a \emph{regulator mark} is either
\begin{enumerate}[label=\textup{(R\arabic*)},leftmargin=3em]
\item one primitive propagator or vertex differentiated with respect to $t$;
\item one derivative of a previously fixed lower-loop local counterterm;
\item one derivative of a determinant/Pfaffian line or source/reference local
coordinate.
\end{enumerate}
The derivative $\partial_t\mathbf R_{j,t}^{<L}\mathcal I_{\Gamma,t}$ is the
finite signed sum of all such marked forests.
\end{definition}

\begin{lemma}[No independent finite part appears in the relative forest]
\label{lem:v8-relative-ownership}
Every marked proper subtraction block in
$\partial_t\mathbf R_{j,t}^{<L}\mathcal I_{\Gamma,t}$ is the derivative of a
previously fixed common coefficient $\bar S_{q,t}$ with $q<L$.  In particular,
regulator comparison introduces no graph-dependent finite normalization.
\end{lemma}
\begin{proof}
By the V5 proper-subgraph loop-drop theorem, every proper 1PI subtraction block
has loop number strictly below that of $\Gamma$.  The chronological V7
construction freezes its finite normalization when that lower order is first
solved.  Differentiating in $t$ therefore differentiates that already-owned
coefficient; it does not reopen its finite part.
\end{proof}

\begin{theorem}[Relative fixed-loop forest decomposition]
\label{thm:v8-relative-forest}
Fix $\mathfrak q$ and a shell $j$.  After one common local Taylor/EFT/source
projection, the regulator derivative of every order-$\le L$ forest coefficient
has a decomposition
\begin{equation}
 \boxed{
 \partial_t\mathbf R_{j,t}^{<L}\mathcal I_{\Gamma,t}
 =\mathsf B_{\Gamma,j,t}^{\rm loc}
  +\mathsf E_{\Gamma,j,t}^{\rm rel}.}
 \label{eq:v8-relative-forest-decomp}
\end{equation}
The local term belongs to the same finite V5 bank.  For any intermediate scale
$M$ with
\begin{equation}
 \mu\ll M\ll a^{-1},
 \label{eq:v8-M-window}
\end{equation}
the complementary rooted kernel obeys
\begin{equation}
 \boxed{
 \|\mathsf E_{\Gamma,j,t}^{\rm rel}\|_*
 \le
 \begin{cases}
 C_{\mathfrak q}\,a\Lambda_j,
     & \Lambda_j\le M,\\[0.25em]
 C_{\mathfrak q}\,\mu/\Lambda_j,
     & \Lambda_j>M.
 \end{cases}}
 \label{eq:v8-relative-forest-bound}
\end{equation}
The statement includes all stored future BV/source contraction descendants.
\end{theorem}
\begin{proof}
Differentiate the finite V5 forest formula.  By
Definition~\ref{def:v8-marked-forest}, every term contains one regulator mark.
On shells below $M$, the mark is a physical-annulus relative propagator/vertex
or a lower-loop counterterm generated by such a mark.  Theorem
\ref{thm:v8-relative-symbol}, the bounded rooted contraction ideal of V5 and
induction in the loop number give the factor $a\Lambda_j$.  Nested and
overlapping proper subgraphs remain in their original compatible forest terms;
Lemma~\ref{lem:v8-relative-ownership} supplies their common finite parts.

For shells above $M$, no small-$ap$ expansion is used.  Apply the native V5
Taylor/forest decomposition separately at each $t$ and differentiate it.  All
nondecaying dependence on the microscopic stencil is local and is assigned to
$\mathsf B_{\Gamma,j,t}^{\rm loc}$.  The first omitted external/source Taylor
order carries one excess soft-to-hard ratio $\mu/\Lambda_j$.  The V5
contraction-complete packet guarantees that a later BV Laplacian or source
contraction of this remainder has its newly local descendant already stored;
there is therefore no hidden loss of the scale factor.  Finiteness of the graph
and source families at fixed $\mathfrak q$ completes the estimate.
\end{proof}

\subsection{The two-scale universality estimate}

The preceding theorem isolates the ultraviolet mechanism cleanly.  Microscopic
regulators are close on physical shells, whereas on cutoff-scale shells they
need not be close but their difference is local as seen from a fixed soft panel.

\begin{lemma}[Geometric two-scale sums]
\label{lem:v8-two-scale-sum}
For dyadic/geometric shell scales $\Lambda_j=b^j\lambda$, $b>1$,
\begin{align}
 \sum_{\Lambda_j\le M}a\Lambda_j
 &\le C_b aM,
 \label{eq:v8-low-sum}\\
 \sum_{\Lambda_j>M}\frac{\mu}{\Lambda_j}
 &\le C_b\frac{\mu}{M}.
 \label{eq:v8-high-sum}
\end{align}
\end{lemma}
\begin{proof}
Both are geometric series, summed respectively from below and from above the
unique shell comparable with $M$.
\end{proof}

Let $\mathcal N_{\mathfrak q}$ denote the fixed complete finite normalization
map used in V6--V7: physical couplings, redundant/BV-exact coordinates,
composite/source basis, stress improvements, contact coordinates and local line
normalization are all included.

\begin{theorem}[One common local regulator-matching velocity]
\label{thm:v8-local-matching}
Along the path $t\mapsto H_t$ there exists a unique chronological local velocity
$Y_{a,t,\mathfrak q}$ in the finite V5 local bank such that the normalized
coordinates remain fixed:
\begin{equation}
 \frac d{dt}\,
 \mathcal N_{\mathfrak q}
 \bigl(\widehat{\mathfrak C}_{a,t}\bigr)=0.
 \label{eq:v8-normalization-fixed}
\end{equation}
It is chosen once for the entire state, not separately for graphs or Ward rows,
and is projective in $L$.
\end{theorem}
\begin{proof}
V6 constructs a normalization section with an invertible finite Jacobian on the
protected constant-rank chart.  Differentiate the normalization equations in
$t$.  The implicit-function theorem gives a unique finite local velocity which
cancels the local vector assembled from
$\mathsf B_{\Gamma,j,t}^{\rm loc}$, line/reference local terms and the Hodge-exact
coordinates.  Because V7 freezes every lower-order choice chronologically, the
order-$q$ component of this linear system is unchanged when a higher terminal
loop order is appended.  Hence the resulting family is projective.
\end{proof}

\begin{theorem}[Two-scale regulator comparison at fixed order]
\label{thm:v8-two-scale-universality}
Let $\widehat{\mathfrak C}_{a,t,\lambda}$ be the normalized state at the fixed
physical scale $\lambda$, transported along the regulator homotopy with the
local velocity of Theorem~\ref{thm:v8-local-matching}.  Then for every
intermediate $M$ satisfying \eqref{eq:v8-M-window},
\begin{equation}
 \boxed{
 \left\|
 \partial_t\widehat{\mathfrak C}_{a,t,\lambda}
 \right\|_{\mathfrak q}
 \le
 C_{\mathfrak q}\left(aM+\frac{\mu}{M}\right).}
 \label{eq:v8-two-scale-bound}
\end{equation}
The same estimate holds after every retained source derivative and for the
complete BV/master and determinant/Pfaffian line packets.
\end{theorem}
\begin{proof}
The local part of every marked forest is cancelled by the one common
normalization velocity.  Sum the complementary contributions from
Theorem~\ref{thm:v8-relative-forest}.  Lemma~\ref{lem:v8-two-scale-sum} gives
$aM$ below $M$ and $\mu/M$ above $M$.  The V5 rooted ideal controls products
containing several complementary ancestors; those terms carry at least the same
single marked factor and higher powers only improve the bound.  V4--V5 source
convolution and contraction completeness show that source/BV descendants obey
the same estimate.  The line contribution is
Theorem~\ref{thm:v8-line-local-plus-rem}.  All families are finite at fixed
$\mathfrak q$.
\end{proof}

\begin{corollary}[Optimized regulator-difference rate]
\label{cor:v8-sqrt-rate}
For $a\mu<1$, choose
\begin{equation}
 M_a=\sqrt{\frac{\mu}{a}}.
 \label{eq:v8-optimal-M}
\end{equation}
Then $\mu\ll M_a\ll a^{-1}$ and
\begin{equation}
 \boxed{
 \left\|
 \partial_t\widehat{\mathfrak C}_{a,t,\lambda}
 \right\|_{\mathfrak q}
 \le C_{\mathfrak q}\sqrt{a\mu}.}
 \label{eq:v8-sqrt-rate}
\end{equation}
Consequently
\begin{equation}
 \left\|
 \widehat{\mathfrak C}_{a,1,\lambda}
 -F_{a,\mathfrak q}^{1\leftarrow0}
  \widehat{\mathfrak C}_{a,0,\lambda}
 \right\|_{\mathfrak q}
 \le C_{\mathfrak q}\eta_*\sqrt{a\mu},
 \label{eq:v8-endpoint-rate}
\end{equation}
where $F_{a,\mathfrak q}^{1\leftarrow0}$ is the integrated finite local
matching map.
\end{corollary}
\begin{proof}
At \eqref{eq:v8-optimal-M}, both terms in
\eqref{eq:v8-two-scale-bound} equal $\sqrt{a\mu}$.  Integrate in
$t\in[0,1]$.  The local flow generated by $Y_{a,t,\mathfrak q}$ is finite
dimensional and invertible on the compact chart, giving
$F_{a,\mathfrak q}^{1\leftarrow0}$.
\end{proof}

\subsection{Continuum universality and exact order compatibility}

\begin{theorem}[Microscopic universality for the explicit range-two deformation]
\label{thm:v8-explicit-universality}
Let regulator $A$ be the V2 nearest-neighbour Wilson/overlap completion and let
regulator $B$ be \eqref{eq:v8-range-two-kernel} with
$\eta=\eta_*$ satisfying \eqref{eq:v8-eta-small}.  Then for every fixed
$\mathfrak q$ they are matched-universal in the sense of
Definition~\ref{def:v8-universality}.  More precisely, the finite local maps
$F_{a,\mathfrak q}^{B\leftarrow A}$ have a projective continuum limit
$F_{\infty,\mathfrak q}^{B\leftarrow A}$ and
\begin{equation}
 \boxed{
 \widehat{\mathfrak C}^{B}_{\infty,\lambda}
 =F_{\infty,\mathfrak q}^{B\leftarrow A}
  \widehat{\mathfrak C}^{A}_{\infty,\lambda}.}
 \label{eq:v8-continuum-universality}
\end{equation}
After covariant source/operator matching, the normalized connected coefficients,
all retained BV/BRST/master relations, stress/composite/contact coefficients and
determinant/Pfaffian line insertions agree.
\end{theorem}
\begin{proof}
The V2 regulator has the V7 continuum.  The range-two regulator satisfies the
same E1 hypotheses by Propositions~\ref{prop:v8-range-two-structure},
Theorem~\ref{thm:v8-gapped-path}, Theorem~\ref{thm:v8-relative-symbol} and
Proposition~\ref{prop:v8-relative-full-zone}.  The Standard-Model representation,
anomaly cancellation and finite gauge descent are unchanged, so the V3--V7
cohomological, forest, normalization and Cauchy constructions apply with the
same fixed-order quantifiers.

For the relative statement, Corollary~\ref{cor:v8-sqrt-rate} shows that the two
normalized finite-cutoff states differ, after the common local matching map, by
a quantity tending to zero.  The local maps are built from the same compact
normalization section as V7; their scale-to-scale differences are summable once
the physical normalization is fixed, so the V7 finite-scheme argument gives an
invertible continuum limit $F_{\infty,\mathfrak q}^{B\leftarrow A}$.  Passing
$a\downarrow0$ in \eqref{eq:v8-endpoint-rate} proves
\eqref{eq:v8-continuum-universality}.  Every listed source/BV/line component is
a coordinate of the contraction-complete state, so the same limit identity
holds coefficientwise for it.
\end{proof}

\begin{theorem}[Strict perturbative-order projectivity of the regulator map]
\label{thm:v8-reg-projective}
Let $\pi_L^{L+1}$ be the V7 perturbative truncation.  The chronological
regulator-matching maps may be chosen so that at finite cutoff
\begin{equation}
 \boxed{
 \pi_L^{L+1}
 F_{a,L+1}^{B\leftarrow A}
 \iota_L^{L+1}
 =F_{a,L}^{B\leftarrow A},}
 \label{eq:v8-reg-projective-finite}
\end{equation}
and therefore
\begin{equation}
 \boxed{
 \pi_L^{L+1}
 F_{\infty,L+1}^{B\leftarrow A}
 \iota_L^{L+1}
 =F_{\infty,L}^{B\leftarrow A}.}
 \label{eq:v8-reg-projective-limit}
\end{equation}
Thus regulator universality is an equivalence of the projective coefficientwise
tower, not a separate finite redefinition at every terminal loop order.
\end{theorem}
\begin{proof}
Theorem~\ref{thm:v8-local-matching} solves the regulator-matching velocity
chronologically.  Once its order-$q$ component is fixed, increasing the
terminal loop order only appends a new equation in the new
$\hbar^{q+1}$ layer and deeper future-descendant slots.  Lower components are
not recomputed.  This proves \eqref{eq:v8-reg-projective-finite}; continuity of
truncation and the finite-dimensional scheme limits gives
\eqref{eq:v8-reg-projective-limit}.
\end{proof}

\subsection{A general gapped-homotopy criterion}

The explicit range-two example reveals which hypotheses were actually used.

\begin{definition}[Admissible regulator homotopy at fixed coefficient order]
\label{def:v8-admissible-homotopy}
A path $t\mapsto H_{a,t}$ is \emph{$\mathfrak q$-admissible} if:
\begin{enumerate}[label=\textup{(A\arabic*)},leftmargin=3em]
\item it is finite range (or uniformly exponentially local), Hermitian and
exactly gauge/spin covariant, with fixed source jets through the orders demanded
by $\mathfrak q$;
\item it has a cutoff- and $t$-uniform Wilson gap on one common protected chart;
\item its overlap/chiral physical branch has the same continuum principal
Dirac--Yukawa symbol for all $t$, and the $t$ derivative of every required
physical source vertex is one irrelevant order smaller, in the sense of the
relative estimate \eqref{eq:v8-relative-D};
\item on the complementary full zone, the hard inverse/source words are smooth
periodic dimensionless symbols with the V2 weighted kernel bounds;
\item the determinant/Pfaffian hard pivots do not cross zero along the path;
\item the same finite V5 local EFT/BV/source bank and normalization-rank chart
remain valid along the path.
\end{enumerate}
\end{definition}

\begin{theorem}[Gapped-homotopy universality criterion]
\label{thm:v8-general-homotopy}
Any two endpoints of a $\mathfrak q$-admissible regulator homotopy are
matched-universal at $\mathfrak q$.  If one coherent path is admissible for
every finite $L$ and the local matching is chosen chronologically, the endpoint
maps form a projective family in $L$.
\end{theorem}
\begin{proof}
Items (A1)--(A5) reproduce precisely the inputs used in
Lemmas~\ref{lem:v8-sign-transgression}--\ref{lem:v8-relative-ownership} and
Theorem~\ref{thm:v8-relative-forest}.  Item (A6) supplies the common finite
normalization velocity.  The two-scale proof
Theorem~\ref{thm:v8-two-scale-universality} and its optimization then apply
verbatim.  Chronological choice gives the projective statement exactly as in
Theorem~\ref{thm:v8-reg-projective}.
\end{proof}

\begin{firewall}[What the homotopy criterion does not prove]
Theorem~\ref{thm:v8-general-homotopy} is a criterion, not a connectedness theorem
for the entire space of chiral regulators.  Two kernels with different overlap
index, a forced Wilson-gap closing, incompatible anomaly representation, or a
determinant-line zero need not lie in one admissible component.  No such phase
transition is crossed in V8.
\end{firewall}

\subsection{V8 anti-circularity audit}
\label{sec:v8-audit}

\begin{enumerate}[leftmargin=2.4em]
\item The second regulator is explicit and microscopically different: it has a
range-two stencil.  Universality is not inferred by renaming the V2 kernel.
\item The regulator path is proved gapped by the operator bound
\eqref{eq:v8-homotopy-gap}; a continuum Dirac operator is not substituted for a
native kernel.
\item High-zone regulator differences are not assumed small.  Their nondecaying
part is placed in the finite local bank and matched by one common normalization
flow; only the Taylor complement is estimated by $\mu/\Lambda_j$.
\item Low-zone closeness and high-zone locality are combined with a moving
intermediate scale.  The comparison does not sum a nonuniform small-$ap$
expansion up to the cutoff.
\item Proper-subgraph finite parts remain owned by the previously fixed common
actions.  No graph is independently refitted while changing regulator.
\item Determinant/Pfaffian frame and source-reference terms are differentiated
explicitly and retained.  A field-independent but source-dependent endpoint
term is not discarded.
\item The order-$L$ regulator map is frozen when constructed.  Higher terminal
orders append new layers and do not recompute lower-order matching.
\end{enumerate}

\section{V8 status ledger}
\label{sec:v8-status}

\subsection*{Proved in V8}
\begin{enumerate}[leftmargin=2.2em]
\item The explicit range-two deformation
$K_{W,a}^{(\eta)}=K_{W,a}^{(0)}+\eta\sum_\mu(L_\mu^{\mathcal U})^2$ is
Hermitian, gauge/spin covariant and finite range, and for sufficiently small
fixed $\eta_*$ it is joined to the V2 regulator by a cutoff-uniform gapped
native homotopy.
\item Along that homotopy, overlap projectors, Kato frames, chiral hard
propagators and every required fixed source jet have uniform quasilocal
$t$ derivatives.  On physical shells the relative propagator is smaller by
$O(a\Lambda_j)$; on the full zone the ordinary lattice-scale bounds hold.
\item Determinant/Pfaffian regulator transgression is given by exact trace
identities, with all source descendants.  Its complete Standard-Model gauge
variation vanishes and its nonlocal shell remainder obeys the same two-zone
bounds.
\item The derivative of the complete native forest is a finite sum of marked
forests whose proper local subtractions remain owned by lower-loop common
actions.  After one common local projection the relative remainder obeys
\eqref{eq:v8-relative-forest-bound}.
\item A single finite normalization/source/line velocity cancels the complete
local regulator derivative.  Summing shells below and above an intermediate
scale gives $C(aM+\mu/M)$; the choice $M=\sqrt{\mu/a}$ yields the explicit
$O(\!\sqrt{a\mu}\!)$ regulator-difference bound.
\item Therefore the V2 nearest-neighbour regulator and the explicit V8
range-two regulator have the same fixed-order continuum after one finite local
matching map.  The equivalence includes normalized connected coefficients,
composites, stress/contact terms, ghosts/antifields, BV/master relations and
determinant/Pfaffian line insertions.
\item The regulator equivalence maps are strictly projective in perturbative
order.
\item More generally, endpoints of any explicitly verified admissible gapped
homotopy satisfying Definition~\ref{def:v8-admissible-homotopy} are
matched-universal at the prescribed fixed order.
\end{enumerate}

\subsection*{Inherited}
\begin{enumerate}[leftmargin=2.2em]
\item The complete V1--V7 coefficientwise continuum construction on the V2
branch, including the native hard-shell estimates, anomaly cancellation,
common BV restoration, arbitrary-fixed-loop forest theorem, ultraviolet Cauchy
theorem, scheme/reference covariance and exact order projectivity.
\item The V2 edge-parallel-transport source calculus and full-zone weighted
kernel estimates.
\item The protected nonzero determinant/Pfaffian component and positive
IR/reference prescription.
\end{enumerate}

\subsection*{Conditional}
\begin{enumerate}[leftmargin=2.2em]
\item The explicit second-regulator theorem requires
$0<\eta_*<g_{\rm W}/32$ and one common compact protected chart.  Larger
range-two deformations are not covered without a new gap proof.
\item The general criterion requires an explicitly verified gapped homotopy and
a common constant-rank normalization/BV chart.  V8 does not prove that every
reasonable chiral regulator admits such a path.
\item All statements remain coefficientwise at prescribed finite
$(L,\Delta,r,N)$, with constants allowed to depend on these data.
\item The positive IR/reference scale and local nonzero line component remain in
force.  No zero crossing or global line holonomy is treated.
\end{enumerate}

\subsection*{Open beyond V8}
\begin{enumerate}[leftmargin=2.2em]
\item Classify the connected components of the admissible native chiral
regulator space and prove a connectedness theorem within a fixed overlap-index
and anomaly-representation sector.  This would remove the need to exhibit the
homotopy by hand for each new microscopic stencil.
\item Extend the universality comparison to a qualitatively different native
fermion realization (for example a non-nearest-neighbour graph kernel or a
non-ultralocal exponentially local completion) by verifying the V8 admissible
homotopy gates.
\item The independent stronger problems from V7 remain: massless IR removal,
global determinant/Pfaffian zero-crossing/holonomy, analytic finite-coupling RG
trajectories and nonperturbative completion.
\end{enumerate}

\subsection*{Exact next load-bearing theorem}
\begin{center}
\fbox{\parbox{0.92\textwidth}{
\textbf{Regulator-phase connectedness / classification.}
Fix the Standard-Model representation and the physical overlap index.  Define
the space of finite-range or uniformly exponentially local, gauge/spin
covariant Hermitian kernels satisfying the common physical-principal-symbol,
source-jet and protected-gap conditions.  Prove that any two kernels in the same
topological phase can be joined by a finite chain of V8-admissible gapped
homotopies, or identify a computable additional obstruction.  Combined with
Theorem~\ref{thm:v8-general-homotopy}, this would upgrade the present explicit
range-two universality theorem to a regulator-class universality theorem.}}
\end{center}

\section*{Append-only continuation marker: V8}
\addcontentsline{toc}{section}{Append-only continuation marker: V8}
\textbf{End of V8.}  The complete V1--V7 body above is retained verbatim.  V8
starts a strictly stronger post-E1--E6 programme and proves microscopic
regulator universality for an explicit range-two covariant deformation of the
V2 Wilson/overlap kernel.  The proof uses a uniform native gap, relative overlap
functional calculus, regulator-marked common forests, a single finite local
normalization/source/line flow, and the two-scale estimate
$O(aM+\mu/M)$ optimized at $O(\sqrt{a\mu})$.  It also proves a general
universality criterion for any explicitly verified admissible gapped regulator
homotopy.  Future versions must append after this marker.



\clearpage
\section{V9 continuation: regulator-phase classification, weak topological obstructions, and a connected Clifford--Wilson sector}
\label{sec:v9-continuation}

V8 proved that the endpoints of any explicitly verified admissible gapped
regulator homotopy have the same fixed-order continuum after one finite local
matching map.  The remaining question is therefore topological before it is
analytic: when does such a homotopy exist?  The strongest possible guess,
``equal four-dimensional overlap index implies connectedness'', is false for a
general multiband Hermitian lattice symbol.  This continuation first identifies
the missing invariant exactly.  It then proves a positive connectedness theorem
for the Wilson-like five-Clifford/no-doubler sector actually containing the V2
kernel.

The distinction is load-bearing.  A general gapped Hermitian symbol defines a
vector bundle over the Brillouin torus, and its complete stable class belongs to
$K^0(\mathbb T^4)$; the four-dimensional strong number sees only one component
of that class.  By contrast, a five-Clifford symbol factors through $S^4$, so
all weak first-Chern classes vanish automatically.  In the no-doubler component
with the physical Dirac germ fixed, the complement of the physical south-pole
preimage is contractible in the target, and a direct relative homotopy can be
constructed.  A quantitative trigonometric approximation then converts the
smooth homotopy into a finite chain of finite-range native homotopies with the
same physical first jet.

Throughout V9, $x=ap\in\mathbb T^4=(-\pi,\pi]^4$ denotes the dimensionless
Brillouin variable.  The positive IR/reference scale and the protected
source/line chart of V1--V8 remain in force.

\subsection{The flattened bundle of a general gapped Hermitian symbol}

\begin{definition}[Flat gapped symbol and flattened projector]
\label{def:v9-flat-symbol}
Let $H:\mathbb T^4\to\operatorname{Herm}(M)$ be a smooth periodic Hermitian
matrix symbol satisfying
\begin{equation}
 \inf_{x\in\mathbb T^4}s_{\min}(H(x))\ge g>0.
 \label{eq:v9-gap}
\end{equation}
Its spectral flattening and positive projector are
\begin{equation}
 \varepsilon_H(x):=\operatorname{sgn}H(x),
 \qquad
 P_H(x):=\frac12(I+\varepsilon_H(x)).
 \label{eq:v9-flattening}
\end{equation}
The associated positive-energy bundle is
\begin{equation}
 E_H:=\operatorname{Ran}P_H\longrightarrow\mathbb T^4.
 \label{eq:v9-bundle}
\end{equation}
A finite-range translation-invariant kernel is the special case in which the
entries of $H$ are trigonometric polynomials.
\end{definition}

\begin{lemma}[Rank is constant on a gapped component]
\label{lem:v9-rank}
For a gapped symbol in Definition~\ref{def:v9-flat-symbol},
$\operatorname{rank}P_H(x)$ is independent of $x$.  Along a continuous gapped
homotopy $H_t$, it is also independent of $t$.
\end{lemma}
\begin{proof}
The rank is an integer-valued continuous function of $x$ because the gap keeps
zero out of the spectrum.  The torus is connected, hence the rank is constant.
The same argument on $\mathbb T^4\times[0,1]$ gives the homotopy statement.
\end{proof}

\begin{theorem}[The $K$-class is a necessary invariant of every gapped regulator homotopy]
\label{thm:v9-K-obstruction}
Let $H_t$, $0\le t\le1$, be a continuous gapped homotopy of smooth Hermitian
symbols.  Then
\begin{equation}
 \boxed{[E_{H_0}]=[E_{H_1}]\quad\text{in }K^0(\mathbb T^4).}
 \label{eq:v9-K-invariant}
\end{equation}
In particular, every Chern number of $E_{H_t}$ is independent of $t$.
\end{theorem}
\begin{proof}
Functional calculus makes $(x,t)\mapsto P_{H_t}(x)$ continuous, and smooth when
the homotopy is smooth.  Its range therefore defines one vector bundle
$\mathcal E\to\mathbb T^4\times[0,1]$.  Restriction to the two boundary faces
gives $E_{H_0}$ and $E_{H_1}$.  Since the inclusions
$\mathbb T^4\times\{0\},\mathbb T^4\times\{1\}\hookrightarrow
\mathbb T^4\times[0,1]$ are homotopy equivalences, both restrictions have the
same $K^0$ class.  Naturality of Chern classes gives the last statement.
\end{proof}

For later use, choose the four coordinate one-cycles and orient the six
coordinate two-tori $\mathbb T^2_{\mu\nu}$ by $dx_\mu\wedge dx_\nu$.  For a
complex bundle $E$ define
\begin{align}
 r(E)&:=\operatorname{rank}E,
 \label{eq:v9-rank-invariant}\\
 \nu_{\mu\nu}(E)&:=
 \left\langle c_1(E),[\mathbb T^2_{\mu\nu}]\right\rangle,
 \qquad 1\le\mu<\nu\le4,
 \label{eq:v9-weak-Chern}\\
 \nu_4(E)&:=
 \left\langle \operatorname{ch}_2(E),[\mathbb T^4]\right\rangle.
 \label{eq:v9-strong-Chern}
\end{align}
On the four-torus the last pairing is integral.  One may equivalently use
$c_2$ together with the six first-Chern numbers.

\begin{theorem}[Stable classification on the four-torus]
\label{thm:v9-KT4}
There is a natural group isomorphism
\begin{equation}
 \boxed{
 K^0(\mathbb T^4)
 \cong
 \mathbb Z
 \oplus \bigwedge
olimits^2\mathbb Z^4
 \oplus \bigwedge
olimits^4\mathbb Z^4
 \cong \mathbb Z^8.}
 \label{eq:v9-KT4}
\end{equation}
Under this identification a bundle class is determined by
\begin{equation}
 \boxed{
 \left(r(E),\{\nu_{\mu\nu}(E)\}_{\mu<\nu},\nu_4(E)\right).}
 \label{eq:v9-K-invariants}
\end{equation}
Consequently, after fixing the spectral rank, a general gapped Hermitian symbol
has seven independent stable topological integers: six weak two-dimensional
numbers and one strong four-dimensional number.
\end{theorem}
\begin{proof}
The circle has
$K^0(S^1)\cong\mathbb Z$ and $K^1(S^1)\cong\mathbb Z$.  Since these groups are
free abelian, the Kunneth sequence for complex $K$-theory has no torsion term.
Iterating it four times gives
\[
 K^0((S^1)^4)
 \cong\bigoplus_{k\ \mathrm{even}}\bigwedge^k\mathbb Z^4
 =\bigwedge^0\mathbb Z^4\oplus
  \bigwedge^2\mathbb Z^4\oplus
  \bigwedge^4\mathbb Z^4.
\]
The three summands have ranks $1,6,1$.  The Chern character of the exterior
products of the four $K^1(S^1)$ generators is the corresponding exterior
product of the normalized integral one-forms.  Hence the degree-zero, degree-two
and degree-four components are read exactly by rank, the six pairings of
$c_1$, and the degree-four Chern-character pairing.  Because the displayed
$K$-group is torsion free, these integer coordinates determine the class.
\end{proof}

\begin{corollary}[Equality of the strong index alone is not a connectedness criterion]
\label{cor:v9-strong-not-enough}
For general gapped Hermitian symbols, equality of rank and $\nu_4$ does not
imply the existence of a gapped homotopy.  Any mismatch in one of the six
integers $\nu_{\mu\nu}$ is an obstruction.
\end{corollary}
\begin{proof}
Immediate from Theorems~\ref{thm:v9-K-obstruction} and
\ref{thm:v9-KT4}.
\end{proof}

\subsection{A concrete weak-index obstruction at fixed strong number}

The preceding obstruction is not merely formal.  Let
$\pi_{12}:\mathbb T^4\to\mathbb T^2$ be projection onto the first two
coordinates.  Choose a complex line bundle $L\to\mathbb T^2$ with
$c_1(L)$ equal to the positive generator and put
\begin{equation}
 E_1:=\pi_{12}^*L,
 \qquad
 E_0:=\underline{\mathbb C}.
 \label{eq:v9-line-example}
\end{equation}
Then
\begin{equation}
 r(E_0)=r(E_1)=1,
 \qquad
 \nu_4(E_0)=\nu_4(E_1)=0,
 \qquad
 \nu_{12}(E_0)=0,
 \quad
 \nu_{12}(E_1)=1.
 \label{eq:v9-line-invariants}
\end{equation}
The equality of the four-dimensional number follows because
$c_1(E_1)$ is pulled back from a two-dimensional factor, hence
$c_1(E_1)^2=0$ and $c_2(E_1)=0$.

\begin{proposition}[The weak obstruction is realized by finite-range gapped symbols]
\label{prop:v9-finite-range-obstruction}
There exist finite-dimensional, finite-range translation-invariant Hermitian
symbols $H_0,H_1$ whose positive bundles are $E_0,E_1$ after stabilization, so
that they have the same spectral rank and strong number but cannot be joined by
any gapped Hermitian homotopy.
\end{proposition}
\begin{proof}
Every complex bundle on the compact torus embeds in a sufficiently large
trivial bundle.  Let $P_i(x)$ be smooth orthogonal projector representatives of
$E_i$ in one common trivial $\mathbb C^M$ bundle and put
$Q_i=2P_i-I$.  Then $Q_i^2=I$, so each $Q_i$ is Hermitian and has unit gap.
Trigonometric matrix polynomials are uniformly dense in smooth periodic matrix
functions together with any prescribed finite number of derivatives.  Choose a
Hermitian trigonometric polynomial $\widetilde Q_i$ with
$\|\widetilde Q_i-Q_i\|<1/4$.  The straight line from $Q_i$ to
$\widetilde Q_i$ stays invertible by the singular-value perturbation inequality,
so spectral flattening shows that $\widetilde Q_i$ has the same positive-bundle
class as $Q_i$.  Trigonometric polynomial symbols are precisely finite-range
translation-invariant kernels.  If a gapped homotopy joined
$\widetilde Q_0$ to $\widetilde Q_1$, Theorem~\ref{thm:v9-K-obstruction}
would force their $K$-classes, and hence $\nu_{12}$, to agree, a contradiction.
\end{proof}

\begin{firewall}[Scope of the counterexample]
Proposition~\ref{prop:v9-finite-range-obstruction} disproves a connectedness
claim for the \emph{general} gapped Hermitian symbol space.  It does not assert
that the extra weak phase is compatible with every additional native condition
imposed on the physical Standard-Model Wilson block.  Those additional
conditions are analyzed next.  In particular, the V2 kernel belongs to a much
smaller five-Clifford/no-doubler class in which the weak invariants vanish.
\end{firewall}

\subsection{The five-Clifford Wilson class}

Fix Hermitian matrices $\Gamma_1,\ldots,\Gamma_5$ satisfying
\begin{equation}
 \Gamma_A\Gamma_B+\Gamma_B\Gamma_A=2\delta_{AB}I.
 \label{eq:v9-Clifford}
\end{equation}
The V2 flat Hermitian Wilson symbol can be written in such a basis after the
standard identification
$\Gamma_\mu=i\gamma_5\gamma_\mu$ for $1\le\mu\le4$ and
$\Gamma_5=\gamma_5$.

\begin{definition}[Five-Clifford gapped symbol]
\label{def:v9-Clifford-symbol}
A five-Clifford symbol is
\begin{equation}
 H_h(x):=\sum_{A=1}^5h_A(x)\Gamma_A,
 \qquad
 h:\mathbb T^4\to\mathbb R^5,
 \label{eq:v9-Clifford-H}
\end{equation}
with $|h(x)|\ge g>0$.  Its normalized map is
\begin{equation}
 n_h(x):=\frac{h(x)}{|h(x)|}\in S^4.
 \label{eq:v9-nmap}
\end{equation}
Let $s:=-e_5\in S^4$ denote the physical south pole.  We call the symbol
\emph{single-cone/no-doubler} if
\begin{equation}
 n_h^{-1}(s)=\{0\},
 \label{eq:v9-no-doubler}
\end{equation}
and the tangent derivative of the first four components at $0$ is invertible.
Two such symbols have the same \emph{physical Dirac germ} if they have the same
value $h(0)$ up to the fixed positive radial normalization and the same first
jet that determines the continuum Dirac speed/coframe.
\end{definition}

The Clifford relation gives
\begin{equation}
 H_h(x)^2=|h(x)|^2I,
 \qquad
 \operatorname{sgn}H_h(x)=n_h(x)\cdot\Gamma.
 \label{eq:v9-Clifford-flat}
\end{equation}
Thus the topology of the flattened symbol is the topology of a map to $S^4$.

\begin{lemma}[Weak Chern numbers vanish in the five-Clifford class]
\label{lem:v9-weak-vanish}
For a five-Clifford gapped symbol, the positive bundle satisfies
\begin{equation}
 \boxed{c_1(E_{H_h})=0.}
 \label{eq:v9-c1zero}
\end{equation}
Hence all six weak numbers $\nu_{\mu\nu}$ vanish.
\end{lemma}
\begin{proof}
The positive projector is the pullback by $n_h$ of the positive eigenspace
bundle $\mathcal S_+\to S^4$ of $n\cdot\Gamma$.  Since
$H^2(S^4;\mathbb Z)=0$, one has $c_1(\mathcal S_+)=0$.  Naturality gives
$c_1(E_{H_h})=n_h^*c_1(\mathcal S_+)=0$.
\end{proof}

\begin{theorem}[Degree is the complete flattened homotopy invariant in the five-Clifford class]
\label{thm:v9-degree-classification}
Based homotopy classes of normalized five-Clifford symbols satisfy
\begin{equation}
 \boxed{
 [\mathbb T^4,S^4]_*\cong H^4(\mathbb T^4;\mathbb Z)\cong\mathbb Z.}
 \label{eq:v9-degree-classification}
\end{equation}
The integer is the degree of $n_h$, equivalently the strong Chern number of the
pulled-back positive bundle up to the fixed sign convention of the chosen
Clifford representation.
\end{theorem}
\begin{proof}
The sphere $S^4$ is $3$-connected and has
$\pi_4(S^4)\cong\mathbb Z$.  Give $\mathbb T^4$ a CW structure with one
zero-cell at the chosen base point.  Obstruction theory extends a based map
uniquely up to homotopy over the $3$-skeleton because
$\pi_k(S^4)=0$ for $k<4$.  The only obstruction/difference class appears on the
four-cells and lies in
$H^4(\mathbb T^4;\pi_4(S^4))\cong H^4(\mathbb T^4;\mathbb Z)$.  Evaluation on
the fundamental class is precisely the degree.  The Chern-number statement
follows because $H^4(S^4;\mathbb Z)$ is generated by the second Chern/Chern
character class of the positive Clifford eigenspace bundle.
\end{proof}

For the no-doubler sector one can do better than abstract obstruction theory.
The physical south pole is attained at exactly one point.  Once the two local
Dirac germs have been aligned, the rest of the target is
$S^4\setminus\{s\}\cong\mathbb R^4$, which is contractible.

\subsection{Local germ straightening without creating a doubler}

Write $B_\rho\subset\mathbb T^4$ for a sufficiently small coordinate ball
around $0$.  In geodesic coordinates
$\exp_s:T_sS^4\supset U\to S^4$ near the south pole, a single-cone map has
\begin{equation}
 \exp_s^{-1}(n_h(x))=Ax+R(x),
 \qquad
 R(x)=O(|x|^2),
 \label{eq:v9-local-germ}
\end{equation}
where $A:\mathbb R^4\to T_sS^4\cong\mathbb R^4$ is invertible.

\begin{lemma}[Quantitative local straightening]
\label{lem:v9-local-straightening}
Let $n_0,n_1$ be two smooth single-cone five-Clifford normalized symbols with
the same physical first jet $A$ at $0$.  There is $\rho>0$ and homotopies,
supported in $B_\rho$, which:
\begin{enumerate}[label=\textup{(S\arabic*)},leftmargin=3em]
\item preserve $n_i(0)=s$ and the derivative $A$ at all times;
\item never create another preimage of $s$;
\item leave each map unchanged outside $B_\rho$;
\item make both final maps equal to one common canonical germ on a smaller
ball $B_{\rho/3}$.
\end{enumerate}
\end{lemma}
\begin{proof}
Because the two maps have the same invertible derivative, in the exponential
chart they are $Ax+R_i(x)$ with
$|R_i(x)|\le C|x|^2$ and
$\|DR_i(x)\|\le C|x|$ after shrinking $\rho$.  Choose a smooth cutoff
$\chi$ equal to one on $B_{\rho/3}$ and zero outside $B_{2\rho/3}$, and define
\begin{equation}
 f_{i,u}(x):=Ax+(1-u\chi(x))R_i(x),
 \qquad 0\le u\le1.
 \label{eq:v9-local-straightening}
\end{equation}
For $\rho$ small enough,
$|f_{i,u}(x)|\ge c|x|$ uniformly for $x\ne0$, because the quadratic remainder
is dominated by the smallest singular value of $A$.  Thus
$\exp_s(f_{i,u}(x))=s$ only at $x=0$.  The correction is quadratic, so the
first derivative at the origin remains $A$.  At $u=1$ both maps equal
$\exp_s(Ax)$ on $B_{\rho/3}$ and are unchanged outside $B_{2\rho/3}$.
\end{proof}

\begin{theorem}[Smooth no-doubler connectedness in the five-Clifford sector]
\label{thm:v9-smooth-connectedness}
Let $n_0,n_1:\mathbb T^4\to S^4$ be smooth single-cone/no-doubler maps with the
same physical Dirac germ.  Then there exists a smooth homotopy $n_t$ such that
\begin{equation}
 \boxed{
 n_t^{-1}(s)=\{0\}
 \quad\text{for every }t\in[0,1],}
 \label{eq:v9-no-doubler-path}
\end{equation}
while the physical first jet at $0$ is fixed throughout.
\end{theorem}
\begin{proof}
Apply Lemma~\ref{lem:v9-local-straightening} to make $n_0$ and $n_1$ agree on a
closed ball $B\subset B_{\rho/3}$, without changing either map outside the
larger local ball and without producing a second south-pole preimage.  Put
$X:=\mathbb T^4\setminus\operatorname{int}B$.  Both restrictions map $X$ into
$S^4\setminus\{s\}$ and agree on $\partial X$.

Choose a fixed stereographic diffeomorphism
\begin{equation}
 \psi:S^4\setminus\{s\}\longrightarrow\mathbb R^4.
 \label{eq:v9-stereographic}
\end{equation}
Because the boundary restrictions agree, the linear homotopy
\begin{equation}
 F_t(x):=(1-t)\psi(n_0(x))+t\psi(n_1(x)),
 \qquad x\in X,
 \label{eq:v9-stereo-homotopy}
\end{equation}
is fixed on $\partial X$.  Define
$n_t=\psi^{-1}(F_t)$ on $X$ and use the common map on $B$.  The two pieces glue
smoothly after the local straightening is chosen constant on a collar of
$\partial B$.  The image on $X$ never contains $s$ because
$\psi^{-1}$ takes values in $S^4\setminus\{s\}$, while the common map on $B$
attains $s$ only at $0$.  The homotopy is stationary near $0$, so the physical
first jet is fixed.
\end{proof}

\begin{corollary}[The strong degree is forced by a single physical cone]
\label{cor:v9-degree-single-cone}
In the orientation convention fixed by the common physical Dirac germ, every
single-cone map of Theorem~\ref{thm:v9-smooth-connectedness} has the same degree,
namely the local degree of the physical south-pole preimage.  In particular the
five-Clifford no-doubler component has no independent weak or strong topological
label once the physical germ is fixed.
\end{corollary}
\begin{proof}
Take $s$ as a regular value after an arbitrarily small germ-preserving
perturbation if necessary.  The degree is the signed sum of the local degrees
at the preimages of $s$.  There is exactly one preimage, and its local degree is
the sign of the determinant of the fixed tangent derivative.  Homotopy
invariance removes the auxiliary perturbation.
\end{proof}

\subsection{From a smooth homotopy to a finite chain of finite-range homotopies}

Theorem~\ref{thm:v9-smooth-connectedness} is a continuum statement on the
Brillouin torus.  V8, however, requires a native finite-range or uniformly
exponentially local path with quantitative source jets.  The next step is a
relative approximation theorem which preserves the physical jet and the
no-doubler margin.

\begin{lemma}[Trigonometric approximation with exact finite jet]
\label{lem:v9-Hermite-Fourier}
Let $f:\mathbb T^4\to\mathbb R^m$ be $C^k$, $k\ge1$, and prescribe the value
and first derivative of $f$ at $0$.  For every $\epsilon>0$ there is a real
trigonometric polynomial $q$ such that
\begin{equation}
 q(0)=f(0),
 \qquad
 Dq(0)=Df(0),
 \qquad
 \|q-f\|_{C^1}<\epsilon.
 \label{eq:v9-Hermite-Fourier}
\end{equation}
The same statement holds for Hermitian matrix-valued functions while preserving
Hermiticity.
\end{lemma}
\begin{proof}
Ordinary trigonometric polynomials are dense in $C^1(\mathbb T^4)$.  Start with
a polynomial $q_0$ within $\delta$ of $f$.  The finite jet error
\[
 \bigl(f(0)-q_0(0),\,Df(0)-Dq_0(0)\bigr)
\]
is corrected by a fixed finite collection of trigonometric polynomials whose
value/first-derivative evaluation matrix at $0$ is the identity.  For example,
constants correct values and the functions $\sin x_\mu$ correct first
derivatives; additional components are corrected independently.  The
correction norm is bounded by a fixed constant times the jet error, which is
$O(\delta)$.  Choose $\delta$ sufficiently small.  For a Hermitian matrix,
approximate the independent real diagonal and real/imaginary off-diagonal
components and pair Fourier coefficients by adjunction.
\end{proof}

\begin{lemma}[Quantitative openness of the single-cone conditions]
\label{lem:v9-no-doubler-open}
Let $n_t$ be the smooth homotopy from
Theorem~\ref{thm:v9-smooth-connectedness}.  There are a ball $B_0$ around zero
and constants $c_0,\delta_0>0$ such that uniformly in $t$:
\begin{enumerate}[label=\textup{(O\arabic*)},leftmargin=3em]
\item in $B_0$, the first four tangent components satisfy
$|n_t^{\top}(x)|\ge c_0|x|$;
\item on $\mathbb T^4\setminus B_0$,
$\operatorname{dist}(n_t(x),s)\ge\delta_0$.
\end{enumerate}
Consequently every sufficiently small $C^1$ perturbation with the same value
and first derivative at $0$ is again single-cone/no-doubler.
\end{lemma}
\begin{proof}
The homotopy is stationary with invertible tangent derivative near $0$, giving
(O1) after shrinking $B_0$.  The compact set
$(\mathbb T^4\setminus B_0)\times[0,1]$ has image disjoint from $s$, so its
distance from $s$ has a positive minimum, proving (O2).  A sufficiently small
$C^1$ perturbation preserves the lower derivative bound in $B_0$ and the
positive target-space separation on the complement.
\end{proof}

\begin{theorem}[Finite-chain finite-range connectedness]
\label{thm:v9-finite-chain}
Let $H^{(0)}$ and $H^{(1)}$ be finite-range five-Clifford symbols satisfying:
\begin{enumerate}[label=\textup{(C\arabic*)},leftmargin=3em]
\item a positive flat Wilson gap;
\item the same physical Dirac first jet at $x=0$;
\item one and only one physical overlap zero, equivalently
$n_{H^{(i)}}^{-1}(s)=\{0\}$.
\end{enumerate}
Then there exist finitely many finite-range five-Clifford symbols
\begin{equation}
 H^{(0)}=K_0,K_1,\ldots,K_m=H^{(1)}
 \label{eq:v9-chain}
\end{equation}
and finite-range linear homotopies
$K_{j,u}=(1-u)K_j+uK_{j+1}$ such that every $K_{j,u}$ is gapped,
single-cone/no-doubler, and has the same physical first jet.
\end{theorem}
\begin{proof}
Use Theorem~\ref{thm:v9-smooth-connectedness} to obtain a smooth no-doubler
homotopy $n_t$.  Choose positive radial functions at the endpoints matching the
original symbols and extend them smoothly and positively along the homotopy;
this gives a smooth vector homotopy $h_t$ with
$H_t=h_t\cdot\Gamma$, a uniform gap
$\inf_{x,t}|h_t(x)|=:g_*>0$, and the same physical first jet.

Uniform continuity in $C^1$ allows a partition
$0=t_0<\cdots<t_m=1$ so fine that adjacent $h_{t_j}$ differ by less than a
small tolerance $\epsilon$ in $C^1$.  Keep the endpoints exact.  At each
interior node apply Lemma~\ref{lem:v9-Hermite-Fourier} to choose a real
trigonometric polynomial vector $q_j$ with the exact common value/first jet at
zero and
$\|q_j-h_{t_j}\|_{C^1}<\epsilon$.  Put $q_0=h_0$ and $q_m=h_1$.

Choose $\epsilon$ smaller than one quarter of the uniform radial gap and smaller
than the openness margin in Lemma~\ref{lem:v9-no-doubler-open}.  Every convex
segment
$q_{j,u}=(1-u)q_j+uq_{j+1}$ then stays uniformly close to the corresponding
short portion of $h_t$.  Hence $|q_{j,u}|\ge g_*/2$ and the normalized map
$q_{j,u}/|q_{j,u}|$ remains single-cone/no-doubler.  Since all node
approximations have the exact same value and first derivative at zero, every
linear segment preserves that physical jet.  Finally
$K_{j,u}=q_{j,u}\cdot\Gamma$ is a trigonometric polynomial symbol, hence a
finite-range kernel.  The union of the finitely many segments gives the
claimed chain.
\end{proof}

\subsection{Native gauge/spin covariant lifting of the finite chain}

A trigonometric polynomial is not yet a native background kernel.  We now fix a
canonical lifting.  For a displacement $n\in\mathbb Z^4$, choose once and for
all a lattice path $\wp(n)$ from $x$ to $x+n$ and let
$T_{\wp(n)}^{\mathcal U}$ be the ordered gauge--spin parallel transport followed
by the shift.  Reverse paths are chosen so that
$(T_{\wp(n)}^{\mathcal U})^*=T_{\wp(-n)}^{\mathcal U}$.

\begin{definition}[Canonical covariantization]
\label{def:v9-covariantization}
For a Hermitian trigonometric polynomial symbol
\begin{equation}
 K(x)=\sum_{n\in F}A_ne^{in\cdot x},
 \qquad A_{-n}=A_n^*,
 \label{eq:v9-Fourier-symbol}
\end{equation}
with finite $F$, define $\mathcal C(K)$ by replacing each Fourier shift
$e^{in\cdot x}$ with the corresponding gauge--spin parallel shift
$T_{\wp(n)}^{\mathcal U}$ and pairing every term with its Hermitian reverse.
Clifford-vector coefficients are interpreted in the local coframe/spin frame
exactly as in the V2 anticommutator prescription.  The zero-displacement scalar
part acts onsite.
\end{definition}

\begin{lemma}[Properties of canonical covariantization]
\label{lem:v9-covariantization}
For every finite polynomial $K$:
\begin{enumerate}[label=\textup{(V\arabic*)},leftmargin=3em]
\item $\mathcal C(K)$ is finite range and Hermitian;
\item it transforms by exact local gauge/spin unitary similarity;
\item at the flat identity-link/reference-coframe background its symbol is
exactly $K$;
\item every fixed gauge, spin-connection and coframe source derivative is again
a finite-range kernel with a bound depending only on the finite Fourier support
and derivative order.
\end{enumerate}
\end{lemma}
\begin{proof}
A finite path product has finite range and transforms with the gauge/spin matrix
at its two endpoints.  Pairing a path with its reverse gives Hermiticity.
At identity transport the path operator is the ordinary displacement shift and
therefore has Fourier multiplier $e^{in\cdot x}$.  Differentiating a finite path
product inserts finitely many differentiated link/spin factors; at fixed source
order the Leibniz expansion is finite.  Coframe derivatives act on the same
local Clifford coefficients used by V2.  This proves all four statements.
\end{proof}

\begin{theorem}[Native lifting theorem for the canonical Clifford--Wilson class]
\label{thm:v9-native-lift}
Let the two endpoints of Theorem~\ref{thm:v9-finite-chain} be equipped with the
canonical covariantization of Definition~\ref{def:v9-covariantization}.  There
exists a common compact perturbative background/source chart, obtained by one
finite shrinking of the V1--V8 chart if necessary, on which every segment of
the covariantized finite chain is a V8-admissible regulator homotopy at every
prescribed fixed coefficient order $\mathfrak q$.
\end{theorem}
\begin{proof}
There are only finitely many segments, each with finite range and a positive
flat Wilson gap.  Their minimum flat gap is therefore a positive number
$g_{\rm chain}$.  By Lemma~\ref{lem:v9-covariantization}, the operator depends
continuously on the finitely many background/source coordinates.  Shrink one
common compact chart until the operator perturbation from the flat symbol is
less than $g_{\rm chain}/4$ for every segment and homotopy parameter.  The
singular-value perturbation inequality then gives a uniform Wilson gap.

The exact common value and first jet at $x=0$ imply that the derivative with
respect to the segment parameter has zero constant and linear physical
momentum jet.  Hence its physical insertion is $O(|ap|^2)$, which is precisely
the one-irrelevant-order gain required by item (A3) of
Definition~\ref{def:v8-admissible-homotopy}.  The quantitative no-doubler
margin of Theorem~\ref{thm:v9-finite-chain}, compactness of the full zone away
from the physical ball, and the same chart shrinking give a uniform hard
inverse on the complementary zone.  Finite range and smooth periodic symbols
supply item (A4).

For the retained positive IR/reference scale, hard shells exclude the physical
zero.  The hard chiral pivot and determinant/Pfaffian nonzero condition are
open conditions.  Because the chain and source parameter sets are compact and
finite, shrink the protected line/source chart once more so the common pivot is
positive on every segment.

For the normalization gate, use the chronological V7 local coordinates in
which the finite normalization functionals are dual to the retained local
counterterm bank.  Adding a local coordinate $c^a\mathcal O_a$ changes the
corresponding normalization row by $c^a$; hence the Jacobian on the local
counterterm directions has the regulator-independent identity block.  Regulator
variation changes the inhomogeneous loop coefficients but not this coefficientwise
rank statement.  Redundant/BV-exact and source/line coordinates are carried in
the same finite triangular ordering.  Thus the finite coefficientwise
normalization/range problem retains constant rank along every segment, giving
item (A6) without assuming a finite-amplitude contraction of the regulator
flow.  If one subsequently chooses a nonlinear physical reparameterization of
these local coordinates, one restricts to the common nonsingular chart exactly
as in V7.
All V8 admissibility requirements are then satisfied simultaneously.
\end{proof}

\subsection{Regulator-class universality inside the canonical Clifford--Wilson component}

\begin{definition}[Canonical single-cone Clifford--Wilson regulator class]
\label{def:v9-canonical-class}
Let $\mathfrak W_{\rm Cliff}^{\rm sc}$ be the class of canonically covariantized
finite-range five-Clifford kernels whose flat symbols have a positive Wilson
gap, the fixed V2 physical Dirac germ, and exactly one physical overlap zero at
$x=0$, and which lie on the common protected source/line component after the
finite chart shrinking of Theorem~\ref{thm:v9-native-lift}.
\end{definition}

\begin{theorem}[Connectedness of the canonical single-cone Clifford--Wilson class]
\label{thm:v9-class-connected}
Any two regulators $A,B\in\mathfrak W_{\rm Cliff}^{\rm sc}$ are connected by a
finite chain of V8-admissible native gapped homotopies at every prescribed
fixed coefficient order $\mathfrak q$.
\end{theorem}
\begin{proof}
Apply Theorem~\ref{thm:v9-finite-chain} to their flat symbols and then
Theorem~\ref{thm:v9-native-lift} to every segment.  The finite chain may be
traversed consecutively; endpoint charts can be replaced by their finite
intersection, which still contains the common reference point and remains a
compact protected perturbative chart.
\end{proof}

\begin{theorem}[Regulator-class universality in the canonical Clifford--Wilson sector]
\label{thm:v9-class-universality}
Fix finite
$\mathfrak q=(L,\Delta,r,N,\mathcal P_{\rm ext})$.  For any two regulators
$A,B\in\mathfrak W_{\rm Cliff}^{\rm sc}$ there is a finite local invertible
matching map
\begin{equation}
 F_{\infty,\mathfrak q}^{B\leftarrow A}
 \label{eq:v9-class-map}
\end{equation}
such that
\begin{equation}
 \boxed{
 \widehat{\mathfrak C}^{B}_{\infty,\lambda}
 =F_{\infty,\mathfrak q}^{B\leftarrow A}
  \widehat{\mathfrak C}^{A}_{\infty,\lambda}.}
 \label{eq:v9-class-universality}
\end{equation}
After covariant operator/source/line matching, the normalized connected
coefficients, stress/composite/contact coefficients, ghosts/antifields,
BV/BRST/master relations and determinant/Pfaffian line insertions agree.  The
maps can be chosen strictly projective in perturbative order.
\end{theorem}
\begin{proof}
By Theorem~\ref{thm:v9-class-connected}, choose a finite chain
$A=R_0,R_1,\ldots,R_m=B$ of V8-admissible homotopies.  Apply
Theorem~\ref{thm:v8-general-homotopy} to each segment, obtaining invertible
continuum matching maps
$F^{R_{j+1}\leftarrow R_j}_{\infty,\mathfrak q}$.  Their finite composition is
$F^{B\leftarrow A}_{\infty,\mathfrak q}$ and satisfies
\eqref{eq:v9-class-universality}.  The V8 map for each segment is chronological
and projective in $L$; a finite composition of projective maps is projective.
The listed coefficient systems are exactly the coordinates transported in V8.
\end{proof}

\begin{corollary}[Path independence up to the V7 finite scheme groupoid]
\label{cor:v9-path-independence}
Two different admissible finite chains from $A$ to $B$ produce continuum
matching maps which differ, if at all, by an admissible finite continuum scheme
transition of V7 fixing the chosen physical normalization coordinates.  Hence
the coordinate-free matched continuum relations are independent of the chosen
chain.
\end{corollary}
\begin{proof}
Compose one chain with the reverse of the other.  This is a closed sequence of
finite regulator/scheme transports starting and ending at the same native
continuum theory.  V7 identifies any residual local coordinate transformation
with an invertible finite scheme transition satisfying the cocycle law.  The
coordinate-free normalized connected functional and master relations are
covariant under that groupoid, so they are unchanged.
\end{proof}

\subsection{What remains outside the Clifford component}

Theorems~\ref{thm:v9-KT4} and \ref{thm:v9-class-connected} together give the
correct classification picture at the present level:
\begin{equation}
 \boxed{
 \begin{array}{c}
 \text{general gapped Hermitian symbols: full }K^0(\mathbb T^4)
 \text{ data are necessary;}\\[0.25em]
 \text{canonical single-cone five-Clifford symbols: the weak data vanish and}
 \text{ the fixed physical germ gives one connected component.}
 \end{array}}
 \label{eq:v9-classification-summary}
\end{equation}

What has \emph{not} been proved is that equality of the stable $K^0$ class is
sufficient for a V8-admissible no-doubler homotopy in a completely general
multiband native regulator class.  Stable $K$-theory allows the addition of
trivial bands and forgets possible unstable finite-rank information.  Moreover,
V8 requires a fixed physical cone, hard-pivot margins and source/line
compatibility, not merely a homotopy of abstract projectors.

\begin{firewall}[Stable classification is not yet general native connectedness]
Equality in $K^0(\mathbb T^4)$ is a complete stable bundle invariant, but V9
does not identify it with connectedness of every finite-rank native
Standard-Model regulator.  A proof of such a statement would have to control
unstable Grassmannian homotopy, preserve the physical no-doubler germ, lift the
projector homotopy to a Hermitian kernel with the required local source jets,
and maintain the determinant/Pfaffian hard pivots.  None of those extra gates is
silently inferred from stable $K$-theory.
\end{firewall}

\subsection{V9 anti-circularity audit}
\label{sec:v9-audit}

\begin{enumerate}[leftmargin=2.4em]
\item \textbf{The four-dimensional index was not assumed complete.}
The full flattened projector bundle was examined first.  The six weak Chern
numbers give explicit additional obstructions in the general Hermitian class.
\item \textbf{The positive Clifford result uses extra structure openly.}
Weak invariants vanish because the projector factors through $S^4$, not because
they were ignored.
\item \textbf{No-doubler connectedness is proved directly.}
After aligning the local germ, the complement maps into
$S^4\setminus\{s\}\cong\mathbb R^4$ and is homotoped there, so the physical
south pole is never crossed away from $x=0$.
\item \textbf{Finite range is recovered, not assumed from a smooth homotopy.}
The Hermite--Fourier approximation preserves the physical value/first jet, and
a finite partition plus openness margins converts the smooth path to a finite
chain of trigonometric-polynomial paths.
\item \textbf{Native gauge/spin covariance is restored before invoking V8.}
The canonical path-transport lift is finite range and exactly covariant; the
common protected chart is then shrunk once to preserve all segment gaps and
hard pivots.
\item \textbf{Stable $K$-theory is not overpromoted.}
Outside the canonical Clifford--Wilson class, V9 records equality of the full
$K$-class as necessary and stable-complete, but does not claim it is sufficient
for a native no-doubler path at fixed finite rank.
\end{enumerate}

\section{V9 status ledger}
\label{sec:v9-status}

\subsection*{Proved in V9}
\begin{enumerate}[leftmargin=2.2em]
\item Every gapped Hermitian regulator homotopy preserves the $K^0$ class of the
flattened positive-energy bundle.
\item $K^0(\mathbb T^4)\cong\mathbb Z^8$; after fixing rank, the stable class has
six weak first-Chern integers and one strong four-dimensional integer.
\item Equality of the strong four-dimensional number alone is not sufficient in
the general Hermitian class.  A concrete weak-index obstruction is realized by
finite-range gapped translation-invariant symbols.
\item In the five-Clifford class all weak first-Chern numbers vanish, and based
flattened homotopy classes are classified by the degree
$[\mathbb T^4,S^4]_*\cong\mathbb Z$.
\item For the single-cone/no-doubler five-Clifford sector with a fixed physical
Dirac germ, any two smooth symbols are connected through a smooth homotopy that
keeps the physical first jet fixed and never creates another overlap zero.
\item That smooth homotopy can be replaced by a finite chain of finite-range
five-Clifford homotopies preserving a positive gap, the exact physical first
jet and the no-doubler condition.
\item Canonical gauge/spin covariantization lifts the finite chain to native
finite-range kernels.  After one finite shrinking of the protected chart, each
segment satisfies the V8 admissible-homotopy gates at every prescribed fixed
coefficient order.
\item Therefore the canonical single-cone Clifford--Wilson regulator class is
connected by finite chains of V8-admissible homotopies, and all regulators in
that class have the same matched coefficientwise continuum after one finite
local scheme/source/line map.  The equivalence is projective in perturbative
order.
\end{enumerate}

\subsection*{Inherited}
\begin{enumerate}[leftmargin=2.2em]
\item V1--V7: the source-complete fixed-order Renewal ESM continuum on the V2
branch, with native chiral shell control, anomaly cancellation, common BV
restoration, arbitrary-fixed-loop forests, UV Cauchy convergence,
scheme/reference covariance and exact order projectivity.
\item V8: microscopic universality for an explicit range-two deformation and
the general admissible-gapped-homotopy universality criterion.
\item The positive IR/reference scale and protected local determinant/Pfaffian
component.
\end{enumerate}

\subsection*{Conditional}
\begin{enumerate}[leftmargin=2.2em]
\item The regulator-class universality theorem is proved for the
\emph{canonical covariantized finite-range five-Clifford single-cone class}.
An arbitrary native kernel with additional multiband, curvature-dependent or
noncanonical source couplings requires an extension of the lifting theorem.
\item The common background/source/line chart may need to be shrunk to the
finite intersection required by a chosen regulator chain.  The theorem is local
on that protected component and does not cross a rank or determinant-line zero.
\item All continuum statements remain coefficientwise at fixed finite
$(L,\Delta,r,N)$ and positive reference scale.  No uniform $L\to\infty$ or
massless IR claim is made.
\end{enumerate}

\subsection*{Open beyond V9}
\begin{enumerate}[leftmargin=2.2em]
\item Extend the connectedness theorem from the five-Clifford class to a
general finite-rank multiband native class.  Determine whether equality of the
full $K^0(\mathbb T^4)$ class together with a fixed physical no-doubler germ is
sufficient, or identify unstable finite-rank invariants.
\item Extend the native lifting theorem to uniformly exponentially local
non-ultralocal kernels and to noncanonical curvature/source couplings while
preserving V8's relative physical-shell gain.
\item The independent stronger problems remain: removal of the positive IR
scale, global determinant/Pfaffian zero-crossing and holonomy, analytic
finite-coupling RG trajectories, and nonperturbative completion.
\end{enumerate}

\subsection*{Exact next load-bearing theorem}
\begin{center}
\fbox{\parbox{0.92\textwidth}{
\textbf{General multiband native phase lifting.}
Fix the Standard-Model representation, the physical single-cone germ and a
full flattened projector class in $K^0(\mathbb T^4)$.  Determine the unstable
finite-rank homotopy classes of the relevant Grassmannian-valued projectors and
prove that equality of all necessary invariants permits a finite chain of
native V8-admissible homotopies preserving the no-doubler margin, source jets
and determinant/Pfaffian pivots; otherwise exhibit the next computable
obstruction.  This is the precise gate between the V9 Clifford-class
universality theorem and universality for a genuinely general multiband native
fermion regulator.}}
\end{center}

\section*{Append-only continuation marker: V9}
\addcontentsline{toc}{section}{Append-only continuation marker: V9}
\textbf{End of V9.}  The complete V1--V8 body above is retained verbatim.  V9
shows that the strong four-dimensional overlap number is not a complete
invariant for general gapped Hermitian symbols: the flattened bundle carries a
full $K^0(\mathbb T^4)$ class with six additional weak Chern numbers at fixed
rank.  It then proves a positive classification theorem for the physically
relevant canonical single-cone five-Clifford Wilson sector: weak invariants
vanish, no-doubler maps with the same physical germ are smoothly connected,
the path admits a finite-range jet-preserving approximation and a native
covariant lift, and V8 therefore gives regulator-class universality throughout
that component.  Future versions must append after this marker.


\clearpage
\section{V10 continuation: the overlap big cell and general balanced multiband phase lifting}
\label{sec:v10-continuation}

V9 deliberately stopped before identifying unstable finite-rank
Grassmannian invariants for a general multiband projector.  That was the
correct firewall for an arbitrary gapped Hermitian symbol, but it is still too
large a target for the actual overlap problem.  The overlap/no-doubler
condition supplies additional geometry which changes the classification
question completely.

The decisive observation is elementary and exact.  On the index-zero overlap
sector the positive spectral projector of the Hermitian Wilson kernel and the
fixed negative-chirality projector have the same rank.  Away from an overlap
zero, these two planes must be transverse.  The set of transverse planes is
one Grassmannian big cell, canonically identified with a vector space of graph
maps and therefore contractible.  Consequently, once the physical cone germ
is fixed near the unique zero, neither the stable $K^0$ class nor any unstable
finite-rank Grassmannian class remains free on the punctured Brillouin torus.
The apparent V9 unstable-classification gate disappears in the perturbative
balanced single-cone sector.

This continuation proves that statement first at the flat projector level,
then lifts it to gapped Hermitian symbols, to finite-range/uniformly
exponentially local native kernels with all required source jets, and finally
to the V8 regulator-universality theorem.  The V9 $K$-theory obstruction is not
retracted: it remains a genuine obstruction in the larger space of arbitrary
gapped Hermitian symbols.  V10 proves that the obstructed examples cannot lie
in the same balanced overlap-compatible single-cone component with the same
physical germ.

Throughout, let
\begin{equation}
 \mathcal H=\mathcal H_+\oplus\mathcal H_-,
 \qquad
 \gamma_5|_{\mathcal H_+}=+I,
 \qquad
 \gamma_5|_{\mathcal H_-}=-I,
 \label{eq:v10-chiral-splitting}
\end{equation}
with
\begin{equation}
 \dim_{\mathbb C}\mathcal H_+
 =\dim_{\mathbb C}\mathcal H_-=:m,
 \qquad
 Q:=P_-:=\frac12(I-\gamma_5).
 \label{eq:v10-Q}
\end{equation}
The internal Standard-Model representation and any spectator flavour factors
are included in the finite-dimensional fibre $\mathcal H$.  The equality of
the two ranks is the perturbative index-zero sector already used to define the
chiral block in the inherited construction.

\subsection{The exact overlap-zero discriminant}

Let $H(x)$ be a smooth periodic Hermitian Wilson symbol with a nonzero Wilson
gap and put
\begin{equation}
 \varepsilon(x):=\operatorname{sgn}H(x),
 \qquad
 P(x):=\frac12(I+\varepsilon(x)).
 \label{eq:v10-P}
\end{equation}
Assume throughout this subsection that
$\operatorname{rank}P(x)=m$, as is automatic on the connected index-zero
component.  The corresponding overlap symbol is
\begin{equation}
 D_{\rm ov}(x)=a^{-1}\bigl(I+\gamma_5\varepsilon(x)\bigr).
 \label{eq:v10-Dov}
\end{equation}

\begin{lemma}[Exact projector form of the overlap operator]
\label{lem:v10-D-projector}
For every $x$,
\begin{equation}
 \boxed{
 \gamma_5D_{\rm ov}(x)
 =\frac{2}{a}\bigl(P(x)-Q\bigr).}
 \label{eq:v10-D-projector}
\end{equation}
Hence $D_{\rm ov}(x)$ is invertible if and only if $P(x)-Q$ is invertible.
\end{lemma}
\begin{proof}
Use $\varepsilon=2P-I$ and
$Q=(I-\gamma_5)/2$:
\[
 \gamma_5D_{\rm ov}
 =a^{-1}(\gamma_5+\varepsilon)
 =a^{-1}(\gamma_5+2P-I)
 =2a^{-1}(P-Q).
\]
Multiplication by the unitary involution $\gamma_5$ does not change
invertibility.
\end{proof}

\begin{lemma}[Difference of balanced orthogonal projectors]
\label{lem:v10-balanced-projectors}
Let $P,Q$ be orthogonal projectors of rank $m$ on a $2m$-dimensional Hermitian
space.  Then
\begin{equation}
 \boxed{
 P-Q\text{ is invertible}
 \quad\Longleftrightarrow\quad
 \operatorname{Ran}P\cap\operatorname{Ran}Q=\{0\}.}
 \label{eq:v10-projection-transverse}
\end{equation}
When these conditions hold,
$\operatorname{Ran}P\oplus\operatorname{Ran}Q=\mathcal H$.
\end{lemma}
\begin{proof}
If $0\ne v\in\operatorname{Ran}P\cap\operatorname{Ran}Q$, then
$(P-Q)v=0$, so $P-Q$ is singular.

Conversely assume the intersection is zero.  The two ranges have dimension
$m$, hence their direct sum has dimension $2m$ and equals $\mathcal H$.  If
$(P-Q)v=0$, put $w:=Pv=Qv$.  Then
$w\in\operatorname{Ran}P\cap\operatorname{Ran}Q$, so $w=0$.  Therefore
$v\in\ker P\cap\ker Q=(\operatorname{Ran}P+\operatorname{Ran}Q)^\perp=0$.
Thus $P-Q$ is injective and hence invertible in finite dimension.
\end{proof}

\begin{theorem}[Exact overlap-zero discriminant]
\label{thm:v10-zero-discriminant}
On the balanced index-zero sector,
\begin{equation}
 \boxed{
 \ker D_{\rm ov}(x)\ne0
 \quad\Longleftrightarrow\quad
 \operatorname{Ran}P(x)\cap\mathcal H_-\ne\{0\}.}
 \label{eq:v10-zero-discriminant}
\end{equation}
Equivalently, the no-overlap-zero locus in
$\operatorname{Gr}_m(\mathcal H)$ is
\begin{equation}
 \mathcal U_Q
 :=\left\{W\in\operatorname{Gr}_m(\mathcal H):W\cap\mathcal H_-=0\right\}.
 \label{eq:v10-big-cell-def}
\end{equation}
\end{theorem}
\begin{proof}
Combine Lemmas~\ref{lem:v10-D-projector} and
\ref{lem:v10-balanced-projectors}, using
$\operatorname{Ran}Q=\mathcal H_-$.
\end{proof}

This theorem is the upstream point missed by a classification of arbitrary
gapped projectors.  The relevant target away from the physical overlap zero is
not the whole Grassmannian but the open complement of the Schubert
discriminant in \eqref{eq:v10-big-cell-def}.

\subsection{The no-zero locus is one contractible big cell}

Let $\pi_\pm$ be the projections associated with
\eqref{eq:v10-chiral-splitting}.

\begin{theorem}[Graph parametrization of the overlap big cell]
\label{thm:v10-big-cell}
The map
\begin{equation}
 \operatorname{Hom}(\mathcal H_+,\mathcal H_-)
 \longrightarrow \mathcal U_Q,
 \qquad
 A\longmapsto \operatorname{Graph}(A)
 :=\{u+Au:u\in\mathcal H_+\},
 \label{eq:v10-graph-map}
\end{equation}
is a real-analytic diffeomorphism.  The corresponding orthogonal projector is
\begin{equation}
 \boxed{
 P_A=
 \begin{pmatrix}I\\ A\end{pmatrix}
 (I+A^*A)^{-1}
 \begin{pmatrix}I&A^*\end{pmatrix}.}
 \label{eq:v10-graph-projector}
\end{equation}
Consequently
\begin{equation}
 \boxed{\mathcal U_Q\simeq\mathbb C^{m^2}}
 \label{eq:v10-big-cell-contractible}
\end{equation}
and in particular $\mathcal U_Q$ is contractible.
\end{theorem}
\begin{proof}
If $W\cap\mathcal H_-=0$ and $\dim W=m$, then
$\pi_+|_W:W\to\mathcal H_+$ is injective between spaces of the same dimension,
hence an isomorphism.  Define
$A=\pi_-\circ(\pi_+|_W)^{-1}$.  Then $W=\operatorname{Graph}(A)$, and uniqueness
is immediate.  Conversely every graph is transverse to $\mathcal H_-$.

Writing the graph inclusion as
$J_Au=(u,Au)$ gives $J_A^*J_A=I+A^*A$.  The orthogonal projector onto the range
of an injective map $J_A$ is
$J_A(J_A^*J_A)^{-1}J_A^*$, which is
\eqref{eq:v10-graph-projector}.  Both directions of the graph correspondence
are analytic on the open transversality locus.  Finally
$\operatorname{Hom}(\mathcal H_+,\mathcal H_-)$ is a complex vector space, so
it is contractible by $A\mapsto(1-t)A$.
\end{proof}

\begin{corollary}[Equivariant big cell]
\label{cor:v10-equivariant-big-cell}
Let a compact group $G$ act unitarily on $\mathcal H$, preserve the chiral
splitting, and suppose $Q$ is $G$-invariant.  Then the $G$-invariant planes in
$\mathcal U_Q$ are identified with
\begin{equation}
 \operatorname{Hom}_G(\mathcal H_+,\mathcal H_-),
 \label{eq:v10-equivariant-Hom}
\end{equation}
a complex vector space.  Hence the equivariant no-zero big cell is also
contractible.
\end{corollary}
\begin{proof}
The graph of $A$ is $G$-invariant exactly when
$A\rho_+(g)=\rho_-(g)A$ for every $g\in G$.  The intertwiners form a linear
subspace of the full Hom space, and the graph homotopy stays inside it.
\end{proof}

For the Standard-Model Wilson block, $G$ in this corollary is applied
representation block by representation block to the internal compact gauge
group.  The spin/coframe covariance is restored at the native-lifting stage
below; the flat graph interpolation itself is performed in one fixed reference
spin frame.

\subsection{The physical point and the local multiband Dirac germ}

The physical point does not lie in the big cell: at $x=0$ the V2 convention has
\begin{equation}
 P(0)=Q,
 \qquad
 \varepsilon(0)=2Q-I=-\gamma_5,
 \qquad
 D_{\rm ov}(0)=0.
 \label{eq:v10-physical-point}
\end{equation}
Near $Q$, use the opposite Grassmannian chart.  Every plane sufficiently close
to $\mathcal H_-$ is the graph of a unique map
\begin{equation}
 B(x):\mathcal H_-\longrightarrow\mathcal H_+,
 \qquad
 \operatorname{Ran}P(x)=\{B(x)v+v:v\in\mathcal H_-\}.
 \label{eq:v10-local-B}
\end{equation}
In this chart
\begin{equation}
 \operatorname{Ran}P(x)\cap\mathcal H_-
 \cong\ker B(x).
 \label{eq:v10-B-kernel}
\end{equation}

\begin{definition}[Balanced elliptic physical overlap germ]
\label{def:v10-physical-germ}
A balanced projector family has the fixed physical germ $\mathcal B$ if, in
\eqref{eq:v10-local-B},
\begin{equation}
 B(x)=\mathcal B(x)+R(x),
 \qquad
 \mathcal B(x):=\sum_{\mu=1}^4x_\mu\mathcal B_\mu,
 \qquad
 R(x)=O(|x|^2),
 \label{eq:v10-germ}
\end{equation}
and there is $c_{\mathcal B}>0$ such that
\begin{equation}
 \boxed{
 s_{\min}(\mathcal B(x))\ge c_{\mathcal B}|x|
 \quad(x\in\mathbb R^4).}
 \label{eq:v10-elliptic-germ}
\end{equation}
The matrices $\mathcal B_\mu$ include the fixed coframe/Dirac speed and all
spectator identity factors.  Equality of physical germs means equality of this
first jet, not merely equality of its determinant or topological degree.
\end{definition}

The ordinary Weyl/Dirac principal symbol has exactly
\eqref{eq:v10-elliptic-germ}; after division by the engineering scale its
smallest singular value is proportional to $|p|$.

\begin{lemma}[Quantitative multiband germ straightening]
\label{lem:v10-germ-straightening}
Let $P_0,P_1$ be two smooth balanced single-cone projector maps with the same
physical germ $\mathcal B$.  There is a ball $B_\rho$ and homotopies supported
in $B_\rho$ such that:
\begin{enumerate}[label=\textup{(G\arabic*)},leftmargin=3em]
\item $P_i(0)=Q$ and the first physical jet is fixed throughout;
\item no new overlap zero is created in $B_\rho\setminus\{0\}$;
\item the maps are unchanged outside $B_\rho$;
\item after the homotopy the two projectors agree on a neighborhood of a
smaller closed ball $B\Subset B_\rho$.
\end{enumerate}
\end{lemma}
\begin{proof}
In the local chart \eqref{eq:v10-local-B}, write
$B_i(x)=\mathcal B(x)+R_i(x)$ with
$\|R_i(x)\|\le C|x|^2$ after shrinking $\rho$.  Choose a cutoff $\chi$ equal
to one on a neighborhood of a smaller ball and zero near
$\partial B_\rho$, and put
\begin{equation}
 B_{i,u}(x)=\mathcal B(x)+(1-u\chi(x))R_i(x).
 \label{eq:v10-B-straightening}
\end{equation}
The value and first derivative at zero do not change.  Moreover
\[
 s_{\min}(B_{i,u}(x))
 \ge c_{\mathcal B}|x|-C|x|^2.
\]
Choose $\rho<c_{\mathcal B}/(2C)$.  Then the right-hand side is positive for
$0<|x|<\rho$, so \eqref{eq:v10-B-kernel} and
Theorem~\ref{thm:v10-zero-discriminant} exclude every additional overlap zero.
Where $\chi=1$, both final maps are the graph of the same canonical linear germ
$\mathcal B(x)$; where $\chi=0$, each original map is unchanged.
\end{proof}

\subsection{General smooth multiband no-doubler connectedness}

We can now replace the anticipated unstable Grassmannian obstruction theory by
a direct relative homotopy.

\begin{theorem}[Smooth balanced multiband single-cone connectedness]
\label{thm:v10-smooth-multiband}
Let
\begin{equation}
 P_0,P_1:\mathbb T^4\longrightarrow
 \operatorname{Gr}_m(\mathcal H)
 \label{eq:v10-two-projectors}
\end{equation}
be smooth projector maps satisfying:
\begin{enumerate}[label=\textup{(M\arabic*)},leftmargin=3em]
\item $P_i(0)=Q$ and both have the same balanced elliptic physical germ
$\mathcal B$;
\item for every $x\ne0$,
$\operatorname{Ran}P_i(x)\cap\mathcal H_-=0$;
\item the same internal representation grading is preserved.
\end{enumerate}
Then there is a smooth homotopy $P_t$ with
\begin{equation}
 \boxed{
 P_t(0)=Q,
 \qquad
 dP_t(0)=dP_0(0),
 \qquad
 \operatorname{Ran}P_t(x)\cap\mathcal H_-=0
 \quad(x\ne0),}
 \label{eq:v10-projector-homotopy}
\end{equation}
for all $t\in[0,1]$.  The homotopy may be chosen within the internal-gauge
equivariant projector subspace.
\end{theorem}
\begin{proof}
Apply Lemma~\ref{lem:v10-germ-straightening} to both endpoints.  After this
preliminary homotopy there is a closed ball $B$ about $0$ and a collar of its
boundary on which the two maps agree.  Keep that common local map fixed.

Put
$X:=\mathbb T^4\setminus\operatorname{int}B$.  By (M2) and
Theorem~\ref{thm:v10-big-cell}, each restriction has a unique graph coordinate
\begin{equation}
 A_i:X\longrightarrow
 \operatorname{Hom}(\mathcal H_+,\mathcal H_-).
 \label{eq:v10-Ai}
\end{equation}
The two graph maps agree on a collar of $\partial X$.  Define
\begin{equation}
 A_t=(1-t)A_0+tA_1.
 \label{eq:v10-A-linear}
\end{equation}
Every $A_t$ defines a plane in the big cell, so
$P_t:=P_{A_t}$ has no overlap zero on $X$.  Since the interpolation is constant
on the boundary collar, it glues smoothly to the fixed common projector on
$B$.  The physical value and first jet are therefore unchanged.  If the
endpoints preserve the internal representation grading, their graph maps lie
in the intertwiner space of Corollary~\ref{cor:v10-equivariant-big-cell}, which
is linear; hence the same interpolation is equivariant.
\end{proof}

\begin{corollary}[No independent stable or unstable bundle invariant remains]
\label{cor:v10-no-extra-topology}
Within the balanced single-cone sector of
Theorem~\ref{thm:v10-smooth-multiband}, the complete homotopy class of the
flattened projector is fixed by the physical germ.  In particular,
\begin{equation}
 [E_{P_0}]=[E_{P_1}]\in K^0(\mathbb T^4),
 \label{eq:v10-K-forced}
\end{equation}
and any finite-rank unstable Grassmannian invariant takes the same value at the
two endpoints.
\end{corollary}
\begin{proof}
Theorem~\ref{thm:v10-smooth-multiband} gives an actual homotopy in the same
finite Grassmannian, which is stronger than stable equivalence.  Pullback
bundles and all homotopy invariants are therefore equal.
\end{proof}

\begin{corollary}[Location of the V9 weak-index counterexamples]
\label{cor:v10-v9-counterexamples}
The two finite-range symbols used in the V9 weak-index obstruction cannot both
satisfy the same balanced physical germ and the single-cone/no-doubler
condition relative to one fixed $Q$.  Any attempted interpolation into that
sector must either change the physical germ or meet the overlap-zero
discriminant away from $x=0$.
\end{corollary}
\begin{proof}
If both symbols satisfied the hypotheses of
Theorem~\ref{thm:v10-smooth-multiband}, Corollary
\ref{cor:v10-no-extra-topology} would force their $K^0$ classes to agree,
contradicting the weak Chern-number difference proved in V9.
\end{proof}

\begin{firewall}[V9 remains correct]
V9 classified the larger space of arbitrary gapped Hermitian symbols before
imposing the overlap no-doubler incidence relative to $Q$.  In that larger
space the full $K^0(\mathbb T^4)$ class is genuinely necessary.  V10 proves
that the actual balanced single-cone overlap component is a much smaller
relative mapping problem: the punctured torus lands in the contractible big
cell $\mathcal U_Q$.  The two statements concern different spaces and are
therefore compatible.
\end{firewall}

\subsection{From projector homotopy to a gapped Hermitian-kernel homotopy}

The overlap depends on a Wilson kernel only through its sign.  This permits a
clean separation between spectral radial data and projector topology.

\begin{lemma}[Radial flattening preserves the overlap operator exactly]
\label{lem:v10-radial-flatten}
Let $H$ be a Hermitian symbol with
$s_{\min}(H)\ge g>0$ and put
$\varepsilon=\operatorname{sgn}H$.  For
$0\le s\le1$ define
\begin{equation}
 H_s:=(1-s)H+s\varepsilon.
 \label{eq:v10-radial-path}
\end{equation}
Then
\begin{equation}
 s_{\min}(H_s)\ge\min\{g,1\},
 \qquad
 \operatorname{sgn}H_s=\varepsilon.
 \label{eq:v10-radial-gap}
\end{equation}
Consequently the overlap projector, overlap Dirac operator, physical zeros and
all determinant-line data depending only on the chiral sign are unchanged
along the radial path.
\end{lemma}
\begin{proof}
The operators $H$ and $\varepsilon$ commute and have the same spectral
projections.  On an eigenvector of $H$ with eigenvalue $\lambda$, the eigenvalue
of $H_s$ is
$(1-s)\lambda+s\operatorname{sgn}\lambda$, which has the same sign as
$\lambda$ and magnitude at least $\min\{|\lambda|,1\}$.  Hence no eigenvalue
crosses zero and the sign is constant.
\end{proof}

\begin{theorem}[Smooth gapped multiband Wilson homotopy]
\label{thm:v10-smooth-H-homotopy}
Let $H_0,H_1$ be two smooth gapped balanced Hermitian symbols whose flattened
projectors satisfy Theorem~\ref{thm:v10-smooth-multiband}.  Then they are joined
by a smooth gapped Hermitian homotopy which preserves the physical overlap germ
and the single-cone/no-doubler condition throughout.
\end{theorem}
\begin{proof}
Use Lemma~\ref{lem:v10-radial-flatten} to join $H_0$ to
$\varepsilon_0=2P_0-I$.  Use Theorem~\ref{thm:v10-smooth-multiband} and set
\begin{equation}
 \varepsilon_t:=2P_t-I.
 \label{eq:v10-epsilon-path}
\end{equation}
This middle path has unit Wilson gap.  Its overlap zeros are exactly those of
$P_t-Q$ by Theorem~\ref{thm:v10-zero-discriminant}, so there is only the fixed
physical zero.  Finally reverse the radial path for $H_1$.  The physical first
jet is unchanged on the middle path and the radial pieces have exactly constant
sign, hence exactly constant overlap operator.
\end{proof}

For finite-range $H_i$, their signs are generally not finite range, but the V1
gapped functional calculus gives uniform exponential locality and the same
fixed source-jet bounds.  Thus the two short radial pieces already lie in the
``finite range or uniformly exponentially local'' class permitted by the V8
admissibility definition.

\subsection{Finite-range approximation of the middle multiband path}

The middle involution path \eqref{eq:v10-epsilon-path} is smooth in the
Brillouin variable but need not itself be finite range.  We now approximate it
without losing its exact physical first jet or its no-doubler margin.

\begin{lemma}[First-jet rigidity of the matrix sign at an involution]
\label{lem:v10-sign-first-jet}
Let $E=E^*=E^{-1}$ and let $X=X^*$ satisfy
\begin{equation}
 EX+XE=0.
 \label{eq:v10-tangent-involution}
\end{equation}
Then the Fr\'echet derivative of the matrix sign at $E$ obeys
\begin{equation}
 \boxed{D(\operatorname{sgn})_E[X]=X.}
 \label{eq:v10-sign-first-jet}
\end{equation}
\end{lemma}
\begin{proof}
Let $P_\pm=(I\pm E)/2$.  The divided-difference formula for the derivative of a
spectral function gives zero on the $P_+XP_+$ and $P_-XP_-$ blocks and
coefficient
$(1-(-1))/(1-(-1))=1$ on both off-diagonal blocks.  Condition
\eqref{eq:v10-tangent-involution} is equivalent to
$P_+XP_+=P_-XP_-=0$.  Hence the derivative is precisely
$P_+XP_-+P_-XP_+=X$.
\end{proof}

\begin{lemma}[Uniform no-doubler margins for the smooth path]
\label{lem:v10-no-doubler-margins}
For the path $P_t$ of Theorem~\ref{thm:v10-smooth-multiband}, one can choose
$0<\rho<\pi$ and constants $c_0,\delta_0>0$ such that
\begin{align}
 s_{\min}B_t(x)&\ge c_0|x|,
 &&0<|x|\le\rho,
 \label{eq:v10-local-margin}\\
 s_{\min}(P_t(x)-Q)&\ge\delta_0,
 &&|x|\ge\rho/2,
 \label{eq:v10-global-margin}
\end{align}
uniformly in $t$.
\end{lemma}
\begin{proof}
The homotopy is stationary in a neighborhood of the physical point after the
germ straightening, so \eqref{eq:v10-local-margin} follows from
\eqref{eq:v10-elliptic-germ} after shrinking $\rho$.  On the compact set
$[0,1]\times\{x:|x|\ge\rho/2\}$, the matrices $P_t(x)-Q$ are invertible by
Theorem~\ref{thm:v10-smooth-multiband}; continuity of the least singular value
gives the positive minimum $\delta_0$.
\end{proof}

\begin{theorem}[Finite-chain approximation for a general balanced multiband path]
\label{thm:v10-finite-chain}
Assume the endpoint symbols are finite range or uniformly exponentially local.
The smooth middle path of Theorem~\ref{thm:v10-smooth-H-homotopy} can be
replaced by a finite chain of Hermitian trigonometric-polynomial segments
\begin{equation}
 K_{j,u}(x),
 \qquad
 1\le j\le J,
 \quad 0\le u\le1,
 \label{eq:v10-finite-chain}
\end{equation}
with the following properties:
\begin{enumerate}[label=\textup{(F\arabic*)},leftmargin=3em]
\item every segment has a common positive Wilson gap;
\item after flattening, every segment has exactly the same physical overlap
value and first jet at $x=0$;
\item no segment creates an overlap zero away from $x=0$;
\item all symbols preserve the fixed internal representation grading;
\item the first and last middle-segment endpoints can be joined to the exact
$\varepsilon_0,\varepsilon_1$ by uniformly exponentially local short segments
with the same properties.
\end{enumerate}
\end{theorem}
\begin{proof}
Regard the Hermitian matrices as a finite-dimensional real vector space.  Apply
the V9 Hermite--Fourier approximation lemma componentwise to the smooth
involution path $\varepsilon_t$, using $C^2$ rather than merely $C^1$
approximation and imposing at every node the exact common value and first
$x$-derivative at $0$.  If an internal compact gauge representation is
present, perform the approximation inside its finite-dimensional commutant;
equivalently average the approximant over the compact group.  This preserves
the exact jet constraints.

Choose a partition $0=t_0<\cdots<t_J=1$ so fine, and the node approximants so
close in $C^2$, that every linear segment $K_{j,u}$ remains within a fixed small
operator-norm neighborhood of the corresponding short portion of
$\varepsilon_t$.  Since the latter has unit gap, the singular-value perturbation
inequality gives a common positive Wilson gap.

Let
$\widetilde\varepsilon_{j,u}:=\operatorname{sgn}K_{j,u}$ and
$\widetilde P_{j,u}=(I+\widetilde\varepsilon_{j,u})/2$.  At $x=0$ the common
value is an involution and the prescribed first derivative is tangent to the
involution manifold, because it is the derivative of the exact path
$\varepsilon_t^2=I$.  Lemma~\ref{lem:v10-sign-first-jet} therefore shows that
flattening preserves the exact first physical jet.

On $|x|\ge\rho/2$, the Riesz functional calculus is locally Lipschitz in the
gapped operator norm.  Make the approximation small enough that
$\|\widetilde P_{j,u}-P_t\|<\delta_0/2$ for a nearby path parameter $t$.
Equation \eqref{eq:v10-global-margin} then keeps
$\widetilde P_{j,u}-Q$ invertible.  On $|x|<\rho$, the exact first jet and the
uniform $C^2$ bound give, in the local graph chart,
\[
 \widetilde B_{j,u}(x)=\mathcal B(x)+O(|x|^2)
\]
uniformly.  Shrinking $\rho$ once more if necessary gives
$s_{\min}\widetilde B_{j,u}(x)\ge(c_0/2)|x|$ for $x\ne0$.  Thus no extra zero
appears near the physical point either.

The exact endpoints $\varepsilon_i$ are uniformly exponentially local by the
V1 functional calculus.  Connect each to a sufficiently close first/last
trigonometric node by a short linear segment.  The same gap and no-doubler
openness estimates apply, and a convex combination of an exponentially local
kernel with a finite-range one is uniformly exponentially local.  This proves
all items.
\end{proof}

\subsection{Native multiband covariant lifting and source jets}

The remaining issue is not topology but covariance.  A general multiband
Fourier coefficient is not necessarily one of the five Clifford matrices used
in V9.  On the protected local coframe/spin chart, however, the whole fibre
endomorphism algebra can be lifted covariantly.

Fix a smooth local spin-frame trivialization
$S_x(e)$ on the protected coframe chart, normalized by $S_x(e_0)=I$ at the
reference coframe and chosen equivariantly under the local spin action.  Such a
trivialization exists on the contractible chart already used for the V1 Kato
and V2 source constructions.  Let
$\operatorname{End}_{G_{\rm int}}(\mathcal H)$ denote the finite-dimensional
internal-gauge intertwiner algebra appropriate to the fixed representation
block.

\begin{definition}[Chart-covariant multiband lift]
\label{def:v10-multiband-lift}
Let
\begin{equation}
 K(x)=\sum_{n\in F}A_ne^{in\cdot x},
 \qquad
 A_{-n}=A_n^*,
 \qquad
 A_n\in\operatorname{End}_{G_{\rm int}}(\mathcal H),
 \label{eq:v10-matrix-Fourier}
\end{equation}
be a finite Hermitian Fourier symbol.  For each nonzero displacement choose a
fixed lattice path $\wp(n)$ and let $T_{\wp(n)}^{\mathcal U}$ be the combined
internal-gauge/spin parallel shift.  Set
\begin{equation}
 A_n(e;x):=S_x(e)A_nS_x(e)^{-1}.
 \label{eq:v10-frame-conjugate}
\end{equation}
The chart-covariant lift is the Hermitian sum obtained from the oriented terms
$A_n(e;x)T_{\wp(n)}^{\mathcal U}$ together with their exact adjoints and the
onsifted onsite term.
\end{definition}

\begin{lemma}[Exact covariance and finite source calculus]
\label{lem:v10-lift-properties}
The lift of Definition~\ref{def:v10-multiband-lift} has the following
properties.
\begin{enumerate}[label=\textup{(C\arabic*)},leftmargin=3em]
\item it is Hermitian and finite range;
\item it transforms by exact local internal-gauge and spin similarity;
\item at the identity-link/reference-coframe background its Fourier symbol is
exactly $K(x)$;
\item every prescribed finite gauge, spin-connection and coframe derivative is
a finite-range kernel with a cutoff-independent coefficientwise bound depending
only on the finite path set and derivative order.
\end{enumerate}
The same construction applies to an exponentially summable Fourier family and
preserves a uniform exponential localization rate after a fixed small loss of
that rate.
\end{lemma}
\begin{proof}
Internal gauge covariance follows because the $A_n$ lie in the intertwiner
algebra and the parallel shift transforms at its two endpoints.  Spin
covariance follows from the equivariant frame rule
$S_x(g\cdot e)=g_xS_x(e)$ and the corresponding endpoint transformation of the
spin parallel transport.  Adding each oriented term together with its exact
adjoint gives Hermiticity.  At the reference frame $S_x=I$ and all parallel
transports are identity before the displacement shift, so the prescribed
Fourier multiplier is recovered.

A fixed source derivative differentiates only finitely many factors along each
fixed path and finitely many factors in the local frame conjugation.  The
Leibniz sum is therefore finite.  For an exponentially summable family, the
number of differentiated factors grows at most polynomially with path length;
this is absorbed by an arbitrarily small reduction of the exponential weight.
\end{proof}

The V8 relative theorem requires more than covariance: the regulator derivative
must not alter the physical principal symbol.  This is enforced before the
lift.

\begin{definition}[Principal-jet-preserving native lift]
\label{def:v10-principal-lift}
Fix once and for all a Hermitian trigonometric polynomial
$G_{\rm phys}(x)$ having the common value and first momentum jet of the
physical involution at $x=0$, and lift $G_{\rm phys}$ by the same native V2
Dirac/coframe/parallel-transport prescription for every regulator parameter.
For a middle-chain symbol $K_t$, write
\begin{equation}
 K_t=G_{\rm phys}+R_t,
 \qquad
 R_t(0)=0,
 \qquad
 \partial_{x_\mu}R_t(0)=0,
 \label{eq:v10-R-moments}
\end{equation}
and apply Definition~\ref{def:v10-multiband-lift} only to the remainder $R_t$.
\end{definition}

\begin{lemma}[The regulator velocity is irrelevant with all fixed source jets]
\label{lem:v10-native-irrelevant}
For the principal-jet-preserving lift, the derivative with respect to the
regulator parameter has zero constant and linear physical momentum jet.  At
every fixed source order required by $\mathfrak q$ it therefore belongs to the
same one-irrelevant-order source class used in V8: on a physical hard shell,
after the engineering/source stocks are divided out, the marked insertion has
an extra factor $a\Lambda_j$ relative to the ordinary hard kinetic insertion.
\end{lemma}
\begin{proof}
The common $G_{\rm phys}$ has no regulator derivative.  For $R_t$, the two
moment conditions in \eqref{eq:v10-R-moments} hold identically in $t$.
Frame conjugation is common to all Fourier coefficients, so differentiating a
coframe source preserves the vanishing sums corresponding to the zero value and
first momentum moment.  Gauge/spin-connection derivatives of path transports
are the same edge-insertion calculus used in V2.  A covariant finite-difference
operator whose flat symbol vanishes to second order has, after any fixed number
of such source insertions, at least one explicit lattice spacing beyond the
principal first-order Dirac vertex.  Equivalently, the V2 Taylor/source
counting gives the relative $a\Lambda_j$ factor on a shell
$|p|\asymp\Lambda_j$.  Since the source order is fixed, the finite Leibniz
multiplicity is absorbed in the constant.
\end{proof}

\subsection{The native balanced multiband phase-lifting theorem}

We can now verify the V8 gates without appealing to stable or unstable
Grassmannian classification.

\begin{definition}[Balanced single-cone multiband native class]
\label{def:v10-multiband-class}
Let $\mathfrak W_{\rm mb}^{\rm sc,0}$ consist of native Hermitian Wilson kernels
such that:
\begin{enumerate}[label=\textup{(B\arabic*)},leftmargin=3em]
\item they are finite range or uniformly exponentially local, exactly
internal-gauge/spin covariant, and possess all fixed source jets required by the
coefficientwise programme;
\item their flat reference symbols have a positive Wilson gap and balanced
positive spectral rank $m=\dim\mathcal H_-$;
\item the overlap symbol has exactly one physical zero at $x=0$ and the fixed
balanced elliptic physical germ $\mathcal B$;
\item they preserve the same Standard-Model representation grading and lie on
the same protected positive-IR determinant/Pfaffian source component.
\end{enumerate}
\end{definition}

\begin{theorem}[General balanced multiband native phase lifting]
\label{thm:v10-native-phase-lift}
Let $A,B\in\mathfrak W_{\rm mb}^{\rm sc,0}$.  For every prescribed finite
coefficient packet
$\mathfrak q=(L,\Delta,r,N,\mathcal P_{\rm ext})$ there exists a finite chain of
native regulator homotopies joining $A$ to $B$ such that every segment is
V8-admissible after one common finite shrinking of the protected
background/source/line chart.
\end{theorem}
\begin{proof}
At the flat reference background, flatten both endpoint kernels.  Their
balanced positive projectors satisfy
Theorem~\ref{thm:v10-smooth-multiband}, which gives a smooth projector homotopy
with the fixed physical germ and no extra overlap zero.  Theorem
\ref{thm:v10-finite-chain} replaces its middle part by a finite chain of
gapped finite-range or uniformly exponentially local Hermitian symbols whose
flattened projectors retain the same properties.  Lemma
\ref{lem:v10-radial-flatten} supplies the two endpoint radial segments.

Lift every finite middle segment by Definitions
\ref{def:v10-multiband-lift} and \ref{def:v10-principal-lift}.  Lemma
\ref{lem:v10-lift-properties} gives exact covariance and the required fixed
source calculus.  Lemma~\ref{lem:v10-native-irrelevant} gives the V8 relative
physical-shell gain.  The full-zone flat symbols are gapped and no-doubler by
construction; compactness of the finite chain provides a positive minimum
margin away from the physical ball.

There are only finitely many segments and regulator parameters.  Continuity of
the native operator and its fixed source derivatives therefore permits one
common shrinking of the V1--V9 perturbative chart so that the Wilson gap, the
hard chiral pivots above the retained positive IR scale, and all
Determinant/Pfaffian nonzero conditions remain uniform on every segment.  The
normalization-rank gate is the same chronological finite local-coordinate gate
used in V9: the counterterm coordinates give the regulator-independent identity
block, while regulator variation changes only the inhomogeneous loop data.
Thus no rank is lost along the chain.

The endpoint radial pieces are uniformly exponentially local by the inherited
gapped functional calculus and have exactly constant sign, so they trivially
satisfy the physical-principal-symbol part of V8 admissibility.  Every V8 gate
is therefore satisfied.
\end{proof}

\begin{theorem}[General balanced multiband regulator-class universality]
\label{thm:v10-multiband-universality}
Fix finite $\mathfrak q$.  Any two regulators
$A,B\in\mathfrak W_{\rm mb}^{\rm sc,0}$ are matched-universal.  There exists an
invertible finite local continuum map
\begin{equation}
 F_{\infty,\mathfrak q}^{B\leftarrow A}
 \label{eq:v10-multiband-map}
\end{equation}
such that
\begin{equation}
 \boxed{
 \widehat{\mathfrak C}^{B}_{\infty,\lambda}
 =F_{\infty,\mathfrak q}^{B\leftarrow A}
  \widehat{\mathfrak C}^{A}_{\infty,\lambda}.}
 \label{eq:v10-multiband-universality}
\end{equation}
After covariant matching, the complete normalized connected/source system,
operator mixing, stress/contact insertions, ghosts/antifields, coupled
BV/BRST/master identities and determinant/Pfaffian line insertions coincide.
The maps may be chosen strictly projective in perturbative order.
\end{theorem}
\begin{proof}
Apply Theorem~\ref{thm:v10-native-phase-lift} and then the V8
gapped-homotopy universality theorem to every segment of the finite chain.
Compose the resulting finite invertible local matching maps.  V8 proves the
coefficientwise equality for every stored coordinate of the contraction-complete
state.  Each segment uses the chronological V7--V8 construction, so the finite
composition remains projective in $L$.
\end{proof}

\begin{corollary}[The V9 Clifford class is a subcase]
\label{cor:v10-Clifford-subcase}
The canonical five-Clifford class
$\mathfrak W_{\rm Cliff}^{\rm sc}$ of V9 is contained in
$\mathfrak W_{\rm mb}^{\rm sc,0}$.  Hence the V9 regulator-class universality
theorem follows as a special case of
Theorem~\ref{thm:v10-multiband-universality}; the five-Clifford factorization
through $S^4$ is sufficient but not necessary for balanced single-cone
connectedness.
\end{corollary}
\begin{proof}
The V9 class has balanced chiral rank, the same physical Dirac germ, one
physical overlap zero and the required native covariance/source properties.
Thus it satisfies Definition~\ref{def:v10-multiband-class}.
\end{proof}

\subsection{Why the anticipated unstable Grassmannian computation was the wrong next variable}

For completeness, V10 records the logical relation among the three spaces which
have appeared in V8--V10:
\begin{equation}
 \boxed{
 \begin{array}{c}
 \text{all gapped Hermitian symbols}
 \supset
 \text{balanced overlap symbols}
 \supset
 \text{balanced single-cone symbols with fixed germ},\\[0.3em]
 K^0(\mathbb T^4)\text{ can obstruct the first space,}\qquad
 \mathcal U_Q\simeq\mathbb C^{m^2}\text{ controls the punctured target of the last.}
 \end{array}}
 \label{eq:v10-space-hierarchy}
\end{equation}

The finite-rank Grassmannian itself can possess unstable homotopy information
in other dimensions or under other rank constraints.  None of that information
is load-bearing here because the no-doubler condition never lets the punctured
Brillouin torus explore the whole Grassmannian.  It is confined to one affine
chart.  This is exactly the kind of upstream variable change demanded by the
programme when a route begins to classify more structure than the theorem uses.

\begin{proposition}[The physical germ is the clutching datum for the balanced single-cone class]
\label{prop:v10-germ-clutching}
Let $P$ be a balanced single-cone projector.  After choosing a sufficiently
small physical ball $B$, the restriction to
$X=\mathbb T^4\setminus\operatorname{int}B$ is null-homotopic inside the big
cell relative to any fixed boundary graph.  Therefore every global bundle or
projector homotopy invariant is determined entirely by how the physical local
germ is glued to that big-cell exterior.
\end{proposition}
\begin{proof}
The exterior graph coordinate is a map
$A:X\to\operatorname{Hom}(\mathcal H_+,\mathcal H_-)$.  Once its boundary
value is fixed, the affine target permits the straight-line homotopy to any
other extension with the same boundary, in particular to a chosen reference
extension.  Thus there is no independent exterior homotopy datum.  All
remaining information is in the fixed boundary/local germ.
\end{proof}

This proposition also explains why a weak Chern class can appear in the V9
unconstrained example but not here: producing such a class requires leaving
the single affine cell somewhere on the punctured torus, which by
Theorem~\ref{thm:v10-zero-discriminant} is precisely an additional overlap-zero
incidence.

\subsection{Boundary of the theorem: nonzero index and forced zero modes}

The balanced-rank hypothesis is essential.  If
\begin{equation}
 \operatorname{rank}P-\operatorname{rank}Q\ne0,
 \label{eq:v10-index-nonzero}
\end{equation}
then two equal-dimensional complementary planes are no longer available and
$P-Q$ cannot be invertible everywhere.  Zero modes are forced by dimension
before any local estimate is attempted.  Likewise, if a topologically
nontrivial gauge background forces the overlap index to jump, the determinant
line must be treated over a zero-mode stratum rather than by the nonzero local
frame used in V1--V10.

\begin{firewall}[No nonzero-index theorem]
Theorems~\ref{thm:v10-native-phase-lift} and
\ref{thm:v10-multiband-universality} are perturbative index-zero theorems on the
protected nonzero hard-shell component.  They do not connect distinct overlap
index sectors, remove topological zero modes, or trivialize the determinant or
Pfaffian line across an index-changing locus.
\end{firewall}

\subsection{V10 anti-circularity audit}
\label{sec:v10-audit}

\begin{enumerate}[leftmargin=2.4em]
\item \textbf{The unstable-classification problem was not assumed away.}
The exact identity \eqref{eq:v10-D-projector} was used first.  It shows that
the actual no-doubler target is the transverse-plane big cell, which is proved
directly to be affine and contractible.
\item \textbf{V9's $K$-theory obstruction is retained at its true scope.}
It applies to arbitrary gapped Hermitian symbols.  V10 proves that the
additional balanced no-doubler hypotheses force all such global invariants once
the physical germ is fixed.
\item \textbf{The physical point is not incorrectly placed inside the big cell.}
It is treated in the opposite local graph chart around $Q$; only the punctured
complement is parametrized by $A:\mathcal H_+\to\mathcal H_-$.
\item \textbf{No-doubler preservation is quantitative.}
Near the physical point it follows from the elliptic singular-value lower bound
on the common first jet; away from that point it follows from a compact positive
margin for $P-Q$.
\item \textbf{A smooth projector homotopy is not silently declared native.}
The middle path is approximated by a finite trigonometric chain with exact
physical first jet, then lifted with explicit gauge/spin parallel transport and
a local spin-frame conjugation.
\item \textbf{Leading geometric/gauge source vertices are not changed.}
The common principal jet is lifted once and frozen; only a remainder with zero
constant/linear momentum moments carries regulator dependence, so the V8
one-irrelevant-order source estimate remains valid.
\item \textbf{Determinant/Pfaffian pivots are not inferred from flat topology.}
They are preserved by one finite shrinking of the compact protected source
chart, using openness and the positive IR hard-shell prescription.
\item \textbf{The theorem does not cross index sectors.}
Balanced rank is used in Lemma~\ref{lem:v10-balanced-projectors}; when it fails,
forced zero modes require a different line-bundle/zero-mode analysis.
\end{enumerate}

\section{V10 status ledger}
\label{sec:v10-status}

\subsection*{Proved in V10}
\begin{enumerate}[leftmargin=2.2em]
\item On the perturbative index-zero overlap sector,
$\gamma_5D_{\rm ov}=2a^{-1}(P-Q)$ and the overlap operator is invertible exactly
when the positive Wilson plane is transverse to the fixed negative-chirality
plane.
\item The no-overlap-zero locus is one Grassmannian big cell,
analytically isomorphic to
$\operatorname{Hom}(\mathcal H_+,\mathcal H_-)$ and therefore contractible;
the equivariant big cell is the corresponding intertwiner vector space.
\item A fixed balanced elliptic physical germ can be straightened locally
without creating another zero.  After that straightening, any two smooth
balanced multiband single-cone projectors with the same germ are connected by
linear interpolation of their exterior graph coordinates, while the physical
ball is held fixed.
\item Hence no independent stable $K^0$ invariant or unstable finite-rank
Grassmannian invariant remains inside the balanced single-cone component once
the physical germ is fixed.  The V9 weak-index counterexamples necessarily
leave that component.
\item Radial spectral flattening connects an arbitrary gapped endpoint kernel
to its sign without changing the overlap operator at all.  The middle projector
homotopy therefore gives a smooth gapped Hermitian-kernel homotopy.
\item The middle path admits a finite Hermitian trigonometric approximation with
exact physical first jet, a uniform Wilson gap and no extra overlap zero; the
endpoint flattening pieces are uniformly exponentially local.
\item A chart-covariant multiband lift gives exact internal-gauge/spin
covariance and finite source calculus for arbitrary allowed matrix Fourier
coefficients.  Freezing the common principal native jet and lifting only the
zero-value/zero-first-moment remainder preserves the V8 one-irrelevant-order
gain for all required fixed source derivatives.
\item Therefore any two regulators in the balanced single-cone multiband native
class $\mathfrak W_{\rm mb}^{\rm sc,0}$ are joined by a finite chain of
V8-admissible homotopies and have the same matched coefficientwise continuum,
with strict perturbative-order projectivity.
\end{enumerate}

\subsection*{Inherited}
\begin{enumerate}[leftmargin=2.2em]
\item V1--V7: the source-complete fixed-order Renewal ESM continuum,
anomaly-free common BV restoration, arbitrary-fixed-loop forests, ultraviolet
Cauchy theorem, finite-scheme/reference covariance and exact loop-order
projectivity.
\item V8: regulator universality along every verified admissible gapped
homotopy, including determinant/Pfaffian source transgression and the two-scale
relative-forest estimate.
\item V9: the full $K^0(\mathbb T^4)$ obstruction for arbitrary gapped Hermitian
symbols, the Clifford single-cone connectedness theorem, Hermite--Fourier jet
approximation and the first native finite-range lifting construction.
\end{enumerate}

\subsection*{Conditional}
\begin{enumerate}[leftmargin=2.2em]
\item The general multiband theorem is local to the balanced perturbative
index-zero sector and to the fixed Standard-Model representation grading.
\item Native lifting uses one protected local spin/coframe trivialization and
may require a finite shrinking of the common source/line chart.  A transition
between distinct coframe/spin charts is not part of V10.
\item All statements remain coefficientwise at prescribed finite
$(L,\Delta,r,N)$ and positive IR/reference scale, with no uniform
$L\to\infty$ estimate.
\item Global determinant/Pfaffian zero crossings, nonzero-index sectors and a
massless IR limit remain outside the theorem.
\end{enumerate}

\subsection*{Open beyond V10}
\begin{enumerate}[leftmargin=2.2em]
\item Extend the phase-lifting/universality theorem to nonzero overlap-index or
topologically nontrivial gauge sectors, where zero modes are forced and the
relevant object is a determinant/Pfaffian line over a stratified Fredholm
Grassmannian rather than the nonzero big cell.
\item Glue the local spin/coframe regulator lifts across overlapping protected
geometric charts and identify the resulting global transition cocycle, if a
global regulator-class statement is desired.
\item The independent stronger problems remain: removal of the positive IR
scale, global determinant/Pfaffian holonomy and zero crossing, analytic
finite-coupling RG trajectories, perturbative summability and nonperturbative
completion.
\end{enumerate}

\subsection*{Exact next load-bearing theorem}
\begin{center}
\fbox{\parbox{0.92\textwidth}{
\textbf{Nonzero-index / zero-mode phase lifting.}
Replace the balanced big-cell argument by a stratified finite-dimensional
Fredholm analysis for
$\operatorname{rank}P-\operatorname{rank}Q=k\ne0$ (and for gauge backgrounds
where the overlap index changes).  Construct the exact zero-mode bundle and its
determinant/Pfaffian line, prove shell factorization and source-jet estimates in
a local zero-mode chart, and classify which regulator homotopies preserve the
index and line orientation.  The target is a V8-style universality theorem
within a fixed nonzero-index sector, with the perturbative V10 theorem recovered
at $k=0$.}}
\end{center}

\section*{Append-only continuation marker: V10}
\addcontentsline{toc}{section}{Append-only continuation marker: V10}
\textbf{End of V10.}  The complete V1--V9 body above is retained verbatim.  V10
goes upstream from the unstable-Grassmannian problem isolated in V9 and uses
the exact overlap identity
$\gamma_5D_{\rm ov}=2a^{-1}(P-P_-)$ to identify the true no-doubler target.
In the balanced perturbative index-zero sector this target is the contractible
Grassmannian big cell of planes transverse to $\mathcal H_-$.  With the
physical Dirac germ fixed, this removes all independent stable and unstable
projector invariants.  A quantitative local germ straightening, big-cell graph
homotopy, finite-range jet-preserving approximation, multiband native
covariant lift and V8 admissibility check then prove regulator-class
universality for the full balanced single-cone multiband native class.  Future
versions must append after this marker.



\clearpage
\section{V11 continuation: fixed-index zero-mode splitting, reduced shell calculus, and local universality}
\label{sec:v11-continuation}

V10 closed the balanced index-zero phase-lifting problem by identifying the
punctured no-doubler target with one affine Grassmannian big cell.  At nonzero
index that exact argument cannot be repeated because zero modes are forced.
The correct upstream variable is nevertheless simpler than the full stratified
Fredholm Grassmannian suggested in the V10 handoff.  On the minimal fixed-index
stratum, all nonzero-mode data again form an affine graph fibre; the only
noncontractible datum is the forced kernel or cokernel plane.  Thus the
nonzero-index problem separates canonically into
\begin{equation}
 \boxed{
 \text{zero-mode bundle/line}
 \quad\oplus\quad
 \text{invertible reduced chiral complement}.}
 \label{eq:v11-upstream-split}
\end{equation}

The finite zero-mode determinant-line factorization and exact Grassmann
saturation identity are already present in the Grand Tensor construction and
are inherited here.  They are not reproved.  The new work is to identify the
fixed-index overlap geometry explicitly, prove a source-differentiable reduced
inverse and shell factorization on a constant-rank zero-mode chart, and show
that the V8 regulator-universality argument survives after the zero-mode line is
carried as an independent finite source packet.

Throughout this continuation let
\begin{equation}
 \mathcal H=\mathcal H_+\oplus\mathcal H_-,
 \qquad
 \dim\mathcal H_+=\dim\mathcal H_-=m,
 \qquad
 Q:=P_-
 \label{eq:v11-chiral-split}
\end{equation}
with $Q$ the orthogonal projection onto $\mathcal H_-$.  Let
$P=(I+\varepsilon)/2$ be the positive spectral projector of a gapped Hermitian
Wilson kernel, and put $W=\operatorname{Ran}P$ and
$p=\dim W$.  We write
\begin{equation}
 k:=p-m.
 \label{eq:v11-index-k}
\end{equation}
The sign convention is the index of the finite chiral map below; changing the
usual physics convention merely changes the sign of $k$.

\subsection{The exact fixed-index overlap map}

Let
\begin{equation}
 D_{\rm ov}=a^{-1}(I+\gamma_5\varepsilon),
 \qquad
 P_+=I-Q.
 \label{eq:v11-overlap}
\end{equation}
The kinetic Weyl block before Yukawa insertion is the rectangular map
\begin{equation}
 D_P:=P_+D_{\rm ov}\big|_W:W\longrightarrow\mathcal H_+.
 \label{eq:v11-chiral-map}
\end{equation}

\begin{lemma}[Exact projection form of the overlap Weyl block]
\label{lem:v11-projection-form}
For every $v\in W$,
\begin{equation}
 \boxed{
 D_Pv=\frac{2}{a}P_+v.}
 \label{eq:v11-D-projection}
\end{equation}
Consequently
\begin{align}
 \ker D_P&=W\cap\mathcal H_-,
 \label{eq:v11-kernel}\\
 \operatorname{coker}D_P
 &\cong \mathcal H_+\ominus P_+W,
 \label{eq:v11-cokernel}
\end{align}
and
\begin{equation}
 \boxed{
 \dim\ker D_P-\dim\operatorname{coker}D_P
 =\dim W-m=k.}
 \label{eq:v11-index-formula}
\end{equation}
\end{lemma}

\begin{proof}
For $v\in W$, $\varepsilon v=v$.  Hence
\[
 D_{\rm ov}v=a^{-1}(I+\gamma_5)v
 =2a^{-1}P_+v,
\]
which proves \eqref{eq:v11-D-projection}.  The kernel and cokernel statements
follow immediately.  If
$a_0:=\dim(W\cap\mathcal H_-)$, then
$\operatorname{rank}(P_+|_W)=p-a_0$ and therefore
\[
 \dim\operatorname{coker}D_P=m-(p-a_0).
\]
Subtracting this from $a_0$ gives $p-m=k$.
\end{proof}

\begin{definition}[Minimal fixed-index stratum]
\label{def:v11-minimal-stratum}
The projector $P$, or its plane $W$, lies in the \emph{minimal index-$k$
stratum} if
\begin{equation}
 \dim\ker D_P=k_+:=\max(k,0),
 \qquad
 \dim\operatorname{coker}D_P=k_-:=\max(-k,0).
 \label{eq:v11-minimal-nullities}
\end{equation}
Equivalently, $D_P$ has maximal possible rank $\min(p,m)$ at fixed $p$.
Zero modes beyond $k_++k_-=|k|$ are called accidental paired zero modes.
\end{definition}

The minimal stratum is the natural fixed-index analogue of the V10 nonzero
big cell.  It excludes an index-changing event and also excludes a pair of
opposite-chirality zero modes appearing without changing the index.

\subsection{The minimal stratum is an affine bundle over the zero-mode Grassmannian}

\begin{theorem}[Geometry for $k\ge0$: kernel plane plus an affine graph]
\label{thm:v11-positive-stratum}
Let $k\ge0$ and $p=m+k$.  A plane
$W\in\operatorname{Gr}(m+k,\mathcal H)$ lies in the minimal index-$k$ stratum
if and only if there are unique data
\begin{equation}
 K\in\operatorname{Gr}(k,\mathcal H_-),
 \qquad
 A\in\operatorname{Hom}(\mathcal H_+,K^\perp)
 \label{eq:v11-positive-data}
\end{equation}
such that
\begin{equation}
 \boxed{
 W=K\oplus
 \{u+Au:u\in\mathcal H_+\}.}
 \label{eq:v11-positive-graph}
\end{equation}
Here $K=\ker D_P$.  Thus the minimal index-$k$ stratum is the total space of
the complex vector bundle
\begin{equation}
 \operatorname{Hom}(\underline{\mathcal H_+},\mathcal K^\perp)
 \longrightarrow
 \operatorname{Gr}(k,\mathcal H_-),
 \label{eq:v11-positive-bundle}
\end{equation}
where $\mathcal K$ is the tautological $k$-plane bundle.  In particular it
deformation retracts onto the kernel Grassmannian.
\end{theorem}

\begin{proof}
Minimality for $k\ge0$ means that $P_+|_W:W\to\mathcal H_+$ is surjective and
has kernel $K=W\cap\mathcal H_-$ of dimension $k$.  Orthogonally decompose
$\mathcal H_-=K\oplus K^\perp$.  For each $u\in\mathcal H_+$ choose
$w\in W$ with $P_+w=u$ and subtract its unique $K$ component.  The resulting
vector has the form $u+Au$ with $Au\in K^\perp$.  This defines a linear map
$A:\mathcal H_+\to K^\perp$.  Its graph is orthogonal to $K$, has dimension
$m$, and together with $K$ exhausts $W$.  Uniqueness follows because a graph
vector with vanishing $\mathcal H_+$ component is zero.

Conversely every plane in \eqref{eq:v11-positive-graph} has
$W\cap\mathcal H_-=K$ and projects surjectively to $\mathcal H_+$, so its
kernel and cokernel dimensions are $k$ and $0$.  The construction depends
smoothly on $K$ and $A$ in standard Grassmann charts.  Fibrewise contraction
$A\mapsto tA$, $0\le t\le1$, gives the deformation retraction.
\end{proof}

\begin{theorem}[Geometry for $k\le0$: cokernel plane plus an affine graph]
\label{thm:v11-negative-stratum}
Let $k=-s\le0$ and $p=m-s$.  A plane
$W\in\operatorname{Gr}(m-s,\mathcal H)$ lies in the minimal index-$k$ stratum
if and only if there are unique data
\begin{equation}
 C\in\operatorname{Gr}(s,\mathcal H_+),
 \qquad
 A\in\operatorname{Hom}(C^\perp,\mathcal H_-)
 \label{eq:v11-negative-data}
\end{equation}
such that
\begin{equation}
 \boxed{
 W=\{u+Au:u\in C^\perp\}.}
 \label{eq:v11-negative-graph}
\end{equation}
Here $C$ identifies canonically with $\operatorname{coker}D_P$.  Hence the
minimal index-$k$ stratum is the total space of the vector bundle
\begin{equation}
 \operatorname{Hom}(\mathcal C^\perp,\underline{\mathcal H_-})
 \longrightarrow
 \operatorname{Gr}(s,\mathcal H_+)
 \label{eq:v11-negative-bundle}
\end{equation}
and deformation retracts onto the cokernel Grassmannian.
\end{theorem}

\begin{proof}
Minimality for $k=-s<0$ means $P_+|_W$ is injective and its image
$S:=P_+W$ has codimension $s$.  Put $C=S^\perp$.  Then every element of $W$
can be written uniquely as $u+Au$ with $u\in S=C^\perp$ and
$Au\in\mathcal H_-$.  Conversely every such graph intersects
$\mathcal H_-$ trivially and projects isomorphically onto $C^\perp$, leaving
cokernel $C$.  Smoothness and the deformation retraction again follow from
$A\mapsto tA$.
\end{proof}

\begin{corollary}[All minimal-stratum topology is carried by the forced zero-mode plane]
\label{cor:v11-topology-zero-mode}
At fixed nonzero index, the affine graph coordinate carries no homotopy
obstruction.  The minimal stratum is homotopy equivalent to
\begin{equation}
 \boxed{
 \operatorname{Gr}(k,\mathcal H_-),\quad k>0;
 \qquad
 \operatorname{Gr}(-k,\mathcal H_+),\quad k<0.}
 \label{eq:v11-zero-mode-Grassmannian}
\end{equation}
Thus for a family over a background parameter space, every global obstruction
which survives inside the minimal stratum is zero-mode-bundle data.  On a
contractible protected chart the zero-mode bundle is trivializable and there is
no additional local projector obstruction.
\end{corollary}

\begin{proof}
The first statement is the pair of deformation retractions just proved.  A
complex vector bundle over a contractible base is trivial, so its classifying
map to the corresponding finite Grassmannian is homotopic to a constant map.
The affine graph fibre is contractible as well.
\end{proof}

\begin{firewall}[Fixed index is not the same as fixed zero-mode bundle globally]
Corollary~\ref{cor:v11-topology-zero-mode} is local on a contractible protected
background chart.  On a noncontractible background component, two families can
have the same numerical index but different kernel/cokernel bundles, determinant
line holonomy, or real Pfaffian orientation.  V11 does not identify such global
data merely from the integer $k$.
\end{firewall}

\subsection{The reduced kinetic block is explicitly invertible}

The affine description gives a canonical reduced inverse before any
perturbative estimate is used.

\begin{proposition}[Canonical reduced block on the minimal stratum]
\label{prop:v11-reduced-block}
In the notation of Theorems~\ref{thm:v11-positive-stratum} and
\ref{thm:v11-negative-stratum}, let
\begin{equation}
 J_Au=(u+Au)(I+A^*A)^{-1/2}
 \label{eq:v11-graph-isometry}
\end{equation}
be the graph isometry from the active $\mathcal H_+$ subspace onto the graph
part of $W$.  Then after deleting the forced kernel and cokernel, the kinetic
chiral map is
\begin{equation}
 \boxed{
 D_P^{\perp}J_A
 =\frac{2}{a}(I+A^*A)^{-1/2}.}
 \label{eq:v11-reduced-explicit}
\end{equation}
For $k\ge0$ the active source and target are both $\mathcal H_+$; for $k<0$
they are both $C^\perp$.  Consequently
\begin{align}
 s_{\min}(D_P^\perp)
 &=\frac{2}{a\sqrt{1+\|A\|^2}},
 \label{eq:v11-reduced-gap}\\
 \det D_P^\perp
 &=\left(\frac2a\right)^r
   \det(I+A^*A)^{-1/2}>0,
 \qquad r=\min(p,m),
 \label{eq:v11-reduced-det}
\end{align}
in these canonical graph frames.
\end{proposition}

\begin{proof}
The vector $J_Au$ is normalized because
$\|u+Au\|^2=\langle u,(I+A^*A)u\rangle$.  By
Lemma~\ref{lem:v11-projection-form}, $D_P$ keeps only the
$\mathcal H_+$ component and multiplies it by $2/a$.  This gives
\eqref{eq:v11-reduced-explicit}.  Functional calculus of the positive matrix
$I+A^*A$ gives the singular-value and determinant formulas.
\end{proof}

The positivity in \eqref{eq:v11-reduced-det} is a graph-frame statement for the
pure kinetic overlap block.  It does not orient a separate real Pfaffian line
and does not remove the zero-mode exterior factor.

\subsection{Constant-rank zero-mode charts for the full chiral/Yukawa block}

The shell theorem uses the full finite chiral block, including Yukawa and
source insertions.  We therefore formulate the analytic part independently of
the special projection coordinates.

\begin{definition}[Protected constant-rank zero-mode chart]
\label{def:v11-zero-chart}
Let $D(\zeta):E\to F$ be a smooth finite-cutoff chiral/Yukawa family on a
compact source/background chart $\mathcal U$.  It is a
\emph{protected index-$k$ zero-mode chart} if:
\begin{enumerate}[label=\textup{(Z\arabic*)},leftmargin=3em]
\item $\operatorname{ind}D(\zeta)=k$ and
$\dim\ker D(\zeta)=k_+$,
$\dim\operatorname{coker}D(\zeta)=k_-$ throughout $\mathcal U$;
\item every nonzero singular value obeys
\begin{equation}
 s_j(D(\zeta))\ge\tau_*>0;
 \label{eq:v11-reduced-floor}
\end{equation}
\item the fixed source derivatives required by the chosen coefficient packet
exist and obey the native finite-range/quasilocal bounds inherited from
V1--V10.
\end{enumerate}
The finite-dimensional zero modes are retained as soft variables and are never
integrated into a hard shell.
\end{definition}

\begin{lemma}[Smooth kernel/cokernel projectors and coefficientwise source jets]
\label{lem:v11-zero-projectors}
On a protected zero-mode chart, the orthogonal projectors
$\Pi_K$ onto $K=\ker D$ and $\Pi_C$ onto
$C=\ker D^*=\operatorname{coker}D$ are smooth.  For any contour
$\Gamma_0$ enclosing $0$ but no nonzero singular square,
\begin{align}
 \Pi_K
 &=\frac{1}{2\pi i}\oint_{\Gamma_0}
 (z-D^*D)^{-1}\,dz,
 \label{eq:v11-kernel-Riesz}\\
 \Pi_C
 &=\frac{1}{2\pi i}\oint_{\Gamma_0}
 (z-DD^*)^{-1}\,dz.
 \label{eq:v11-coker-Riesz}
\end{align}
Every prescribed fixed source derivative is a finite sum of resolvent words and
is bounded by a constant depending only on the fixed derivative order,
$\tau_*^{-1}$, and the declared source stocks.
\end{lemma}

\begin{proof}
By \eqref{eq:v11-reduced-floor}, the spectra of $D^*D$ and $DD^*$ have a
uniform gap between $0$ and $\tau_*^2$.  Choose a common contour inside this
gap.  Riesz functional calculus gives the projectors.  Differentiating a
resolvent uses
\[
 \partial(z-A)^{-1}=(z-A)^{-1}(\partial A)(z-A)^{-1};
\]
iterating produces a finite sum at every fixed derivative order.  The contour
resolvents are uniformly bounded by the inverse spectral separation.
\end{proof}

\begin{proposition}[Moore--Penrose inverse in a zero-mode chart]
\label{prop:v11-pseudoinverse}
Define
\begin{equation}
 \boxed{
 D^\dagger
 :=(D^*D+\Pi_K)^{-1}D^*
 =D^*(DD^*+\Pi_C)^{-1}.}
 \label{eq:v11-pseudoinverse}
\end{equation}
Then $D^\dagger$ is the inverse of
$D^\perp:K^\perp\to C^\perp$ and vanishes on $C$.  Moreover
\begin{equation}
 \|D^\dagger\|\le\tau_*^{-1},
 \label{eq:v11-pseudoinverse-bound}
\end{equation}
and every fixed source derivative is a finite expression in
$D,D^*,\Pi_K,\Pi_C$ and their source jets with uniformly bounded inverse factors.
\end{proposition}

\begin{proof}
On $K^\perp$, $D^*D$ is positive with lower edge $\tau_*^2$, while on $K$ the
added projector is the identity.  Hence $D^*D+\Pi_K$ is invertible.  Its inverse
times $D^*$ is exactly the singular-value definition of the Moore--Penrose
inverse.  The second expression is obtained similarly on the target.  The norm
is the reciprocal of the smallest nonzero singular value.  Differentiation uses
the ordinary inverse rule together with Lemma~\ref{lem:v11-zero-projectors}.
\end{proof}

\begin{corollary}[Local Kato frames for the zero-mode bundle]
\label{cor:v11-zero-Kato}
On every star-shaped subchart of $\mathcal U$, the kernel and cokernel bundles
admit unitary Kato frames with every fixed source derivative bounded.  Their top
exterior powers furnish local frames of the zero-mode determinant line.  A real
Pfaffian zero-mode line additionally requires its declared orientation.
\end{corollary}

\begin{proof}
Apply the V1 Kato ODE to the smooth finite-rank projectors
$\Pi_K$ and $\Pi_C$.  Lemma~\ref{lem:v11-zero-projectors} supplies the derivative
bounds.  Take top exterior powers.  Complex exterior powers are canonically
oriented as complex lines; a real Pfaffian line carries the separate orientation
data already distinguished in the inherited construction.
\end{proof}

\subsection{Exact zero-mode-preserving shell factorization}

The inherited finite zero-mode theorem factors the full determinant line into a
zero-mode factor and a nonzero reduced factor.  V11 now proves that the latter
factor admits the same exact shell Schur algebra as V1--V8 while the zero-mode
factor is transported unchanged.

Choose orthogonal decompositions
\begin{equation}
 E=K\oplus E_\ell\oplus E_h,
 \qquad
 F=C\oplus F_\ell\oplus F_h,
 \label{eq:v11-zero-shell-split}
\end{equation}
with the zero modes entirely in the retained soft sector.  Relative to the
reduced complements, write
\begin{equation}
 D^\perp=
 \begin{pmatrix}
 A&B\\ C_1&H
 \end{pmatrix}:
 E_\ell\oplus E_h\longrightarrow F_\ell\oplus F_h.
 \label{eq:v11-reduced-block-matrix}
\end{equation}

\begin{theorem}[Zero-mode-preserving Schur factorization]
\label{thm:v11-zero-Schur}
Assume the hard block $H$ is invertible.  Then the reduced Schur complement
\begin{equation}
 S=A-BH^{-1}C_1
 \label{eq:v11-zero-Schur-complement}
\end{equation}
is invertible whenever $D^\perp$ is invertible, and the full determinant-line
section factors canonically as
\begin{equation}
 \boxed{
 \mathcal Z_D\otimes\det(D^\perp)
 =\mathcal Z_D\otimes\det(H)\otimes\det(S),}
 \label{eq:v11-zero-Schur-line}
\end{equation}
where $\mathcal Z_D$ denotes the inherited kernel/cokernel exterior factor in
the chosen determinant-line convention.  In particular, integrating one hard
shell does not delete, divide by, or re-fit the zero-mode line.
\end{theorem}

\begin{proof}
The inherited zero-mode factorization first separates $K$ and $C$, leaving the
square invertible map $D^\perp:K^\perp\to C^\perp$.  The ordinary finite Schur
identity applied to \eqref{eq:v11-reduced-block-matrix} gives
\[
 \det D^\perp=\det H\,\det(A-BH^{-1}C_1).
\]
If $D^\perp$ and $H$ are invertible, the right-hand determinant forces $S$ to
be invertible.  Tensor the scalar reduced identity with the unchanged
zero-mode exterior factor.  No operation acts on $K$ or $C$ because those
coordinates were excluded from $E_h,F_h$ by construction.
\end{proof}

\begin{corollary}[Differentiated shell line with zero-mode contacts]
\label{cor:v11-zero-Schur-derivatives}
In local Kato zero-mode frames, every fixed source derivative of
\eqref{eq:v11-zero-Schur-line} is the finite Leibniz/Bell-polynomial expansion
of:
\begin{enumerate}[label=\textup{(D\arabic*)},leftmargin=3em]
\item the zero-mode frame/saturation factor;
\item the hard reduced determinant or Pfaffian factor;
\item the reduced Schur factor;
\item every mixed contact generated by derivatives hitting more than one of
these three factors.
\end{enumerate}
No source-dependent zero-mode endpoint term may be discarded as field
independent.
\end{corollary}

\begin{proof}
Differentiate the exact tensor-product identity.  The reduced inverse
insertions are controlled by Proposition~\ref{prop:v11-pseudoinverse}; the hard
block uses the V1--V8 determinant/Pfaffian derivative calculus.  Kato-frame
derivatives are finite by Corollary~\ref{cor:v11-zero-Kato}.  Repeated
differentiation is exactly the finite product-partition expansion already used
for the V1 line cocycle.
\end{proof}

\subsection{Saturation writers and scalar normalized amplitudes}

A nonzero-index vacuum determinant is not a scalar nonzero partition function.
The inherited finite zero-mode theorem states that a nonvanishing amplitude
requires a section of the dual zero-mode line.  We make that datum explicit in
the shell state.

\begin{definition}[Zero-mode saturation writer]
\label{def:v11-saturation-writer}
On a protected zero-mode chart let
\begin{equation}
 \omega_Z\in\Gamma(\mathcal Z_D^*)
 \label{eq:v11-saturation-writer}
\end{equation}
be a chosen smooth nonvanishing local saturation writer, with all required
source derivatives retained.  The corresponding scalarized reduced fermion
factor is
\begin{equation}
 \mathcal A_D^{\rm sat}
 :=\langle\omega_Z,z_D\rangle\,\det'(D),
 \label{eq:v11-saturated-amplitude}
\end{equation}
where $z_D$ is the zero-mode exterior section in the inherited
factorization and $\det'(D)=\det D^\perp$.
\end{definition}

The writer may represent explicit external fermion insertions, a boundary
state, or another physical saturation mechanism.  It is not supplied by the
pseudodeterminant itself.

\begin{theorem}[Exact shell factorization of the saturated amplitude]
\label{thm:v11-saturated-shell}
Under Theorem~\ref{thm:v11-zero-Schur},
\begin{equation}
 \boxed{
 \mathcal A_D^{\rm sat}
 =\langle\omega_Z,z_D\rangle
   \det(H)\det(S).}
 \label{eq:v11-saturated-shell}
\end{equation}
Every retained source derivative is obtained by differentiating this one
identity.  A change of local zero-mode frame or saturation writer acts by an
endpoint line coboundary/contact transformation of the same type as V7.
\end{theorem}

\begin{proof}
Pair the zero-mode factor in \eqref{eq:v11-zero-Schur-line} with
$\omega_Z$ and use the inherited zero-mode saturation theorem.  The frame-change
statement follows because changing a nonzero local frame multiplies
$z_D$ and the dual writer by reciprocal nonzero functions.  If only
one is reparametrized, the resulting local function and all of its source
derivatives are precisely finite source/contact coordinates and must be
transported as in V7.
\end{proof}

\begin{firewall}[The pseudodeterminant is not the amplitude]
The factor $\det'(D)$ controls the invertible complement but does not saturate
the Grassmann zero modes and does not orient a real Pfaffian kernel.  Every V11
scalar continuum statement below is therefore either line-valued or explicitly
conditioned on the saturation writer $\omega_Z$ (and, when relevant, a real
orientation).  This is inherited physical data, not a regulator counterterm.
\end{firewall}

\subsection{The fixed-index extension of the V8 hard-shell calculus}

The hard ultraviolet estimates never need an inverse on the zero-mode space.
They need only the inverse of the reduced hard block.  This observation is the
analytic reason the coefficientwise programme can be extended to a fixed-index
sector without lifting the zero modes.

\begin{definition}[Fixed-index admissible regulator homotopy]
\label{def:v11-index-admissible}
A native regulator path $t\mapsto H_{W,t}$ on a compact background/source chart
is \emph{$(\mathfrak q,k)$-admissible} if:
\begin{enumerate}[label=\textup{(I\arabic*)},leftmargin=3em]
\item it satisfies the V8 locality, covariance, physical-principal-symbol and
full-zone source-jet conditions at the prescribed fixed coefficient data
$\mathfrak q$;
\item the Wilson gap remains uniformly positive, so $\operatorname{rank}P_t$
and hence the index $k$ are constant;
\item the full chiral/Yukawa block stays in the minimal index-$k$ stratum and
its reduced nonzero singular values have a uniform floor $\tau_*>0$;
\item kernel/cokernel modes are retained outside every hard shell and their
Kato source jets are included in the finite source packet;
\item the zero-mode determinant/Pfaffian line and the chosen saturation writer
(or real orientation) are transported continuously along the path;
\item the same finite local BV/EFT/source normalization chart remains of
constant rank.
\end{enumerate}
\end{definition}

\begin{theorem}[Fixed-index hard-shell comparison]
\label{thm:v11-index-hard-shell}
Along a $(\mathfrak q,k)$-admissible homotopy, every V8 relative hard-shell
estimate remains valid after replacing the full chiral inverse by the reduced
inverse $D^\dagger$ and adjoining the finite zero-mode source packet.  In
particular, after the common local Taylor/source projection, a regulator-marked
order-$\le L$ forest obeys
\begin{equation}
 \boxed{
 \|\mathsf E_{\Gamma,j,t}^{\rm rel}\|_*
 \le
 \begin{cases}
 C_{\mathfrak q,k}\,a\Lambda_j,&\Lambda_j\le M,\\[0.2em]
 C_{\mathfrak q,k}\,\mu/\Lambda_j,&\Lambda_j>M,
 \end{cases}}
 \label{eq:v11-index-relative-forest}
\end{equation}
for every intermediate scale $\mu\ll M\ll a^{-1}$.  The same estimate holds
after every stored BV/source contraction descendant.
\end{theorem}

\begin{proof}
By Definition~\ref{def:v11-index-admissible}, the zero modes are never assigned
to the high shell.  Every internal hard fermion contraction therefore uses the
inverse of the reduced hard block, which is the restriction of $D^\dagger$ to
the relevant shell.  Proposition~\ref{prop:v11-pseudoinverse} supplies the
uniform fixed-source derivative calculus on the protected chart.  The relative
physical-shell insertion is the same irrelevant regulator mark as in V8,
because the microscopic principal Dirac--Yukawa symbol is unchanged; the finite
zero-mode projectors contribute only bounded finite-rank source factors.

For $\Lambda_j\le M$, the V8 rooted ideal therefore retains one factor
$a\Lambda_j$.  For $\Lambda_j>M$, subtract the complete local Taylor/source jet
of the reduced forest; the first omitted soft momentum again supplies
$\mu/\Lambda_j$.  Proper subgraphs remain owned by the same lower-loop common
actions.  Differentiating a zero-mode frame or saturation writer creates a
finite local/source insertion but no hard inverse, so it is assigned to the
same contraction-complete packet.  The V5 future-contact construction then
prevents a later BV contraction from converting an unrecorded zero-mode term
into a new local divergence.  Summing the finite forest family proves
\eqref{eq:v11-index-relative-forest}.
\end{proof}

\begin{corollary}[The two-scale estimate is unchanged by protected zero modes]
\label{cor:v11-index-two-scale}
With one common local normalization/source/zero-line velocity chosen to keep
the physical normalization and saturation coordinates fixed,
\begin{equation}
 \left\|
 \partial_t\widehat{\mathfrak C}_{a,t,\lambda}^{(k),\rm sat}
 \right\|_{\mathfrak q}
 \le C_{\mathfrak q,k}
 \left(aM+\frac{\mu}{M}\right).
 \label{eq:v11-index-two-scale}
\end{equation}
Choosing $M=\sqrt{\mu/a}$ gives
\begin{equation}
 \boxed{
 \left\|
 \partial_t\widehat{\mathfrak C}_{a,t,\lambda}^{(k),\rm sat}
 \right\|_{\mathfrak q}
 \le C_{\mathfrak q,k}\sqrt{a\mu}.}
 \label{eq:v11-index-sqrt}
\end{equation}
\end{corollary}

\begin{proof}
The proof of the V8 two-scale theorem applies to the reduced forests by
Theorem~\ref{thm:v11-index-hard-shell}.  The local velocity now also includes
the finite Kato zero-mode-frame and saturation coordinates.  They form a finite
triangular extension of the V7 normalization system and do not change the
geometric shell sums.  Optimize exactly as in V8.
\end{proof}

\subsection{Fixed-index regulator universality with zero-mode saturation}

\begin{theorem}[Local fixed-index regulator universality]
\label{thm:v11-index-universality}
Fix finite coefficient data $\mathfrak q$, an integer $k$, a compact protected
constant-rank zero-mode chart, and a nonvanishing local zero-mode saturation
writer (plus a real orientation if a Pfaffian sector is present).  Let regulators
$A$ and $B$ be the endpoints of a $(\mathfrak q,k)$-admissible native homotopy.
Then there is an invertible finite local matching map
\begin{equation}
 F_{\infty,\mathfrak q,k}^{B\leftarrow A}
 \label{eq:v11-index-matching}
\end{equation}
such that the saturated normalized continuum states obey
\begin{equation}
 \boxed{
 \widehat{\mathfrak C}_{\infty,\lambda}^{B,(k),\rm sat}
 =F_{\infty,\mathfrak q,k}^{B\leftarrow A}
  \widehat{\mathfrak C}_{\infty,\lambda}^{A,(k),\rm sat}.}
 \label{eq:v11-index-universality}
\end{equation}
The equivalence includes normalized connected coefficients with the declared
zero-mode insertions, composite/operator-mixing coefficients, stress/contact
terms, ghosts/antifields, the common BV/BRST/master relations on the saturated
source packet, and determinant/Pfaffian reduced-line insertions.  The matching
maps may be chosen projectively in perturbative order.
\end{theorem}

\begin{proof}
Corollary~\ref{cor:v11-index-two-scale} is the sole new analytic input required
in the V8 universality proof.  Integrating the regulator parameter gives a
finite-cutoff matching error $O(\sqrt{a\mu})$.  The local matching flow contains
the same finite EFT/BV/source coordinates as V8 together with the finite
zero-mode frame/saturation coordinates.  Its cutoff-to-cutoff variation is
summable by the V6--V8 argument, so it has an invertible continuum limit.
Equation \eqref{eq:v11-index-universality} follows as $a\downarrow0$.

The connected functional is taken only after the zero-mode line has been paired
with the declared saturation writer; Theorem~\ref{thm:v11-saturated-shell}
shows that this scalarization is compatible with every shell elimination.  The
BV/source descendants are included because the zero-mode Kato packet is part
of the same contraction-complete source state.  Chronological V7 normalization
freezes each lower-loop zero-mode/source matching when first chosen, proving
projectivity in terminal loop order.
\end{proof}

\begin{corollary}[The index-zero theorem is recovered without a separate case]
\label{cor:v11-kzero-recovery}
For $k=0$, $K=C=0$, the saturation line is trivial and the minimal stratum is
the V10 balanced big cell.  Theorem~\ref{thm:v11-index-universality} reduces to
the V8--V10 regulator-universality theorem.
\end{corollary}

\begin{proof}
All zero-mode projectors and line factors vanish from the construction, while
$D^\dagger=D^{-1}$.  The positive- and negative-index graph descriptions both
reduce to the balanced graph cell of V10.
\end{proof}

\subsection{A local phase-lifting theorem on a contractible zero-mode chart}

The preceding theorem is conditional on a native admissible path.  The
finite-dimensional topological part of constructing such a path can now be
closed completely on a local chart.

\begin{theorem}[Local minimal-stratum connectedness after zero-mode matching]
\label{thm:v11-local-stratum-connected}
Let $X$ be a contractible compact parameter chart.  Let
$P_0,P_1:X\to\operatorname{Gr}(m+k,\mathcal H)$ be two smooth families in the
minimal index-$k$ stratum.  Choose a smooth unitary identification of their
kernel bundles when $k>0$, or of their cokernel bundles when $k<0$.  Then there
is a smooth homotopy $P_t$ through the same minimal stratum which realizes this
zero-mode identification and linearly interpolates the remaining affine graph
coordinate.  In particular, after the zero-mode bundle is matched, there is no
further local projector-level phase obstruction.
\end{theorem}

\begin{proof}
Assume first $k>0$.  By Theorem~\ref{thm:v11-positive-stratum}, each family is
specified by a kernel subbundle $K_i\subset X\times\mathcal H_-$ and a section
$A_i$ of
$\operatorname{Hom}(\underline{\mathcal H_+},K_i^\perp)$.  The chosen unitary
bundle identification and contractibility of $X$ allow a smooth block-unitary
homotopy in $\mathcal H_-$ carrying $K_0$ to $K_1$.  Apply it to the first
projector family; this preserves the minimal stratum.  We may then assume the
kernel bundle is the same $K$.  Both graph coordinates are sections of the
same vector bundle, so
$A_t=(1-t)A_0+tA_1$ remains in the affine fibre and therefore in the minimal
stratum.

For $k<0$, use Theorem~\ref{thm:v11-negative-stratum} and the same argument on
the cokernel bundle $C\subset X\times\mathcal H_+$.  The case $k=0$ is the V10
big-cell interpolation.
\end{proof}

\begin{firewall}[Projector connectedness is not yet a global native lift]
Theorem~\ref{thm:v11-local-stratum-connected} closes the finite-dimensional
local topology.  It does not by itself prove that an arbitrary projector
homotopy is generated by a cutoff-uniform finite-range or exponentially local,
gauge/spin-covariant microscopic Wilson kernel on an entire nontrivial gauge
background component.  V11 regulator universality therefore uses the native
admissibility hypothesis of Definition~\ref{def:v11-index-admissible}.  A
global locality-preserving lift is the next separate gate.
\end{firewall}

\subsection{What becomes of the determinant line at fixed nonzero index}

The geometry above also identifies precisely where a global obstruction can
remain.  On the minimal stratum the reduced graph fibre is affine, while the
zero-mode plane is not.  The determinant/Pfaffian phase therefore lives on the
zero-mode bundle after the reduced phase has been normalized.

\begin{proposition}[Residual line is the zero-mode line]
\label{prop:v11-residual-line}
On a protected minimal index-$k$ chart, choose the canonical positive graph
frame of Proposition~\ref{prop:v11-reduced-block} for the pure overlap kinetic
complement and the inherited finite normalization for the interacting reduced
block.  Then every residual topological line factor is carried by the kernel
and cokernel exterior powers.  In particular, on a contractible chart the
complex determinant line is locally trivial, while on a noncontractible
background component its possible holonomy is exactly a zero-mode/Fredholm-line
problem and is not removed by the reduced pseudodeterminant.
\end{proposition}

\begin{proof}
The pure kinetic reduced determinant is positive in the canonical graph frame
by \eqref{eq:v11-reduced-det}.  For the interacting block, the inherited
zero-mode factorization writes the full line as the zero-mode exterior factor
tensored with the reduced invertible determinant line.  The latter is locally
scalarized by the V7/V11 normalization and Schur cocycle.  Therefore any
remaining line topology is in the zero-mode exterior factor and its transition
functions.  Local triviality on a contractible chart is standard for complex
line bundles.  No statement about loop holonomy follows without examining the
global background component.
\end{proof}

\subsection{V11 anti-circularity audit}
\label{sec:v11-audit}

\begin{enumerate}[leftmargin=2.4em]
\item \textbf{Zero modes were not deleted.}
They are retained as the kernel/cokernel bundle and its exterior line.  Only
the reduced complement is inverted.
\item \textbf{The pseudodeterminant was not promoted to a vacuum amplitude.}
A scalar amplitude requires the explicit dual zero-mode saturation writer of
Definition~\ref{def:v11-saturation-writer}.
\item \textbf{Fixed numerical index was not assumed to classify a global
family.}
The minimal stratum deformation retracts to the zero-mode Grassmannian; global
kernel/cokernel bundle data and line holonomy remain visible.
\item \textbf{Accidental paired zero modes are excluded by a stated reduced
singular-value floor.}
A crossing of that floor leaves the protected minimal stratum and is not hidden
inside a divergent inverse.
\item \textbf{Hard shells never integrate the zero modes.}
Every hard fermion contraction uses the reduced inverse $D^\dagger$ on
$K^\perp\to C^\perp$.
\item \textbf{Source derivatives of the zero-mode projectors are retained.}
They are Riesz-projector/Kato derivatives and enter the same finite contact
packet as the determinant-line frame.
\item \textbf{The V8 two-scale estimate is applied only after the exact local
zero-mode factorization.}
Finite-rank zero-mode factors do not supply the ultraviolet scale gain; the
gain still comes from one irrelevant microscopic regulator mark on the reduced
hard forest.
\item \textbf{Projector-level connectedness was not silently called a native
locality theorem.}
The native lift remains an explicit admissibility gate outside the local
finite-dimensional topology.
\end{enumerate}

\section{V11 status ledger}
\label{sec:v11-status}

\subsection*{Proved in V11}
\begin{enumerate}[leftmargin=2.2em]
\item The finite overlap Weyl kinetic map is exactly
$D_P=2a^{-1}P_+|_{\operatorname{Ran}P}$, and its Fredholm index is
$\operatorname{rank}P-m$.
\item The minimal fixed-index projector stratum is an affine vector bundle over
the forced kernel Grassmannian for $k>0$ or the forced cokernel Grassmannian for
$k<0$.  Hence all nonzero-mode projector data are contractible once the
zero-mode plane is fixed.
\item In canonical graph frames the reduced kinetic block is explicitly
$2a^{-1}(I+A^*A)^{-1/2}$, with a positive reduced determinant and an exact
singular-value floor.
\item On a protected constant-rank full chiral/Yukawa chart, the kernel and
cokernel projectors have Riesz formulas, Kato frames and fixed-order source jets;
the Moore--Penrose inverse is uniformly bounded on the reduced complement.
\item The inherited finite zero-mode determinant-line factorization is stable
under exact shell Schur elimination: the zero-mode exterior factor is carried
unchanged while the reduced determinant obeys the ordinary hard-block/Schur
product.  All mixed source contacts are retained.
\item A nonzero local saturation writer in the dual zero-mode line converts the
line-valued factor into a scalar saturated amplitude, and this scalarization
commutes exactly with shell elimination.
\item Along every fixed-index V11-admissible native regulator homotopy, the V8
relative forest bounds survive with $D^{-1}$ replaced by the reduced
pseudoinverse.  The same $aM+\mu/M$ two-scale bound and
$O(\!\sqrt{a\mu}\!)$ optimized regulator difference follow.
\item Therefore endpoints of such a homotopy have the same matched
coefficientwise saturated continuum, including the complete source/BV/contact
and reduced determinant/Pfaffian packets, with strict perturbative-order
projectivity.
\item On a contractible protected parameter chart, after matching the
kernel/cokernel bundle there is no additional local projector-level phase
obstruction: the remaining graph coordinate interpolates linearly in an affine
fibre.
\end{enumerate}

\subsection*{Inherited}
\begin{enumerate}[leftmargin=2.2em]
\item The Grand Tensor finite zero-mode factorization, saturation identity,
Fredholm determinant line and the statement that zero modes carry the residual
massless line obstruction.
\item V1--V7: source-complete fixed-order continuum, common BV restoration,
arbitrary-fixed-loop forest theorem, ultraviolet Cauchy convergence,
scheme/reference covariance and exact loop-order projectivity.
\item V8: regulator-marked relative forests, line transgression and the
two-scale universality theorem along admissible gapped homotopies.
\item V9--V10: regulator-class phase lifting in the index-zero
Clifford and general balanced multiband single-cone sectors.
\end{enumerate}

\subsection*{Conditional}
\begin{enumerate}[leftmargin=2.2em]
\item V11 works on a minimal constant-index stratum with a positive reduced
singular-value floor.  An accidental zero-mode pair or index-changing crossing
requires a different stratified analysis.
\item Scalar normalized amplitudes require a declared nonzero zero-mode
saturation writer; a real Pfaffian sector additionally requires its physical
orientation.
\item The local projector connectedness theorem does not itself manufacture a
uniformly local native microscopic path on an arbitrary global topological
gauge component.  The regulator-universality theorem assumes the path satisfies
Definition~\ref{def:v11-index-admissible}.
\item All continuum claims remain coefficientwise at fixed finite
$(L,\Delta,r,N)$ and positive IR/reference scale.  No uniform large-loop or
massless-IR limit is asserted.
\end{enumerate}

\subsection*{Open beyond V11}
\begin{enumerate}[leftmargin=2.2em]
\item Prove a \emph{global native fixed-index lifting theorem}: starting from a
projector homotopy in the minimal stratum on a nontrivial gauge-background
component, construct a cutoff-uniform finite-range or exponentially local,
gauge/spin-covariant Wilson-kernel homotopy with the required source jets and
reduced chiral gap.
\item Compute and trivialize, when physically appropriate, the global
zero-mode determinant/Pfaffian line holonomy.  Equal index and local anomaly
cancellation alone do not supply this.
\item Analyze crossings between minimal strata: creation/annihilation of an
opposite-chirality zero-mode pair and genuine index change.  This requires a
spectral-flow/divisor theorem rather than a bounded pseudoinverse.
\item The independent stronger problems remain: positive-IR removal, analytic
finite-coupling RG trajectories, perturbative summability and nonperturbative
completion.
\end{enumerate}

\subsection*{Exact next load-bearing theorem}
\begin{center}
\fbox{\parbox{0.92\textwidth}{
\textbf{Global zero-mode line transport and native fixed-index lifting.}
On a connected nonzero-index background component, promote the local V11
kernel/cokernel Kato bundles to a globally patched zero-mode determinant or
Pfaffian line, construct a locality-preserving native lift of minimal-stratum
projector homotopies, and compute the resulting loop holonomy/spectral-flow
cocycle.  Prove regulator-class universality precisely on the subcomponent
where this global line datum and the reduced chiral gap agree.  This is the
remaining gate between the local fixed-index theorem of V11 and a global
nonzero-index regulator-class theorem.}}
\end{center}

\section*{Append-only continuation marker: V11}
\addcontentsline{toc}{section}{Append-only continuation marker: V11}
\textbf{End of V11.}  The complete V1--V10 body above is retained verbatim.
V11 replaces the full stratified-Fredholm picture by the minimal fixed-index
geometry actually needed by the shell theorem.  The overlap Weyl map is the
projection $2a^{-1}P_+|_{\operatorname{Ran}P}$; its minimal index-$k$ stratum is
an affine bundle over the forced kernel or cokernel Grassmannian.  The reduced
complement has an explicit positive inverse, while the zero-mode bundle and its
exterior line are carried as finite source data.  Riesz/Kato source calculus,
Moore--Penrose hard propagation, exact zero-mode-preserving Schur factorization
and an explicit saturation writer extend the V8 two-scale universality theorem
to every native homotopy which preserves the fixed-index reduced gap.  The
remaining obstruction is global: locality-preserving lifting and holonomy of
the zero-mode determinant/Pfaffian line on noncontractible background
components.


\clearpage
\section{V12 continuation: global zero-mode bundles, faithful-spin line transport, and the native locality obstruction}
\label{sec:v12-continuation}

V11 closed the fixed-index shell calculus on a contractible protected chart.
The remaining gate has three logically distinct layers which should not be
collapsed into one ``global phase'' condition:
\begin{equation}
 \boxed{
 \text{determinant/Pfaffian line datum}
 \quad\subsetneq\quad
 \text{full kernel/cokernel bundle datum}
 \quad\subsetneq\quad
 \text{native local-operator path datum}.}
 \label{eq:v12-three-layers}
\end{equation}
The first layer controls scalarization and phase/orientation.  The second
controls the homotopy class of a family inside the V11 minimal fixed-index
stratum.  The third asks whether such a family homotopy is realized by a
cutoff-uniform local microscopic Wilson kernel with the required source jets.
The inclusions in \eqref{eq:v12-three-layers} are strict in general.

The Grand Tensor manuscript already contains two load-bearing inputs for the
first layer.  First, its finite zero-mode theorem identifies the Fredholm line
with the kernel/cokernel exterior factor tensored with the reduced nonzero
line.  Second, on the ordinary-spin faithful Standard-Model branch with
\(G_6=(\mathrm{SU}(3)\times\mathrm{SU}(2)\times\mathrm{U}(1))/\mathbb Z_6\),
it proves vanishing of the complete local anomaly polynomial together with the
five-dimensional faithful-spin bordism/eta character, and hence a unit
parallel section of the complete complex determinant line on every smooth
connected family to which its geometric hypotheses apply.  These statements
are inherited and are not reproved below.

The new work is instead:
\begin{enumerate}[label=\textup{(V12.\arabic*)},leftmargin=3.8em]
\item to patch the V11 local Kato zero-mode frames globally and identify their
exact Cech/connection cocycle;
\item to combine the inherited faithful-spin parallel determinant section with
V11's reduced nonzero determinant to construct a global zero-mode saturation
writer for the complex Standard-Model determinant packet;
\item to identify the full global projector obstruction with the homotopy class
of the kernel or cokernel classifying map, not merely with its determinant
line;
\item to prove a finite-dimensional stable-range theorem under which an
isomorphism of zero-mode bundles already gives the required projector homotopy;
\item to show by an explicit volume-growing example that an abstract zero-mode
projector homotopy need not be a native quasilocal regulator homotopy;
\item to reduce global fixed-index regulator universality to path connectedness
of an open protected subset of the native exponentially local operator space,
and to prove that every such path may be replaced by a finite polygonal native
chain to which the V11--V8 estimates apply.
\end{enumerate}

Throughout, let \(B\) be a connected compact smooth parameter space (or a
finite CW complex with smooth charts where derivatives are used) carrying a
finite-cutoff chiral/Yukawa family
\begin{equation}
 D_b:E_b\longrightarrow F_b,
 \qquad b\in B,
 \label{eq:v12-family}
\end{equation}
which stays in the V11 minimal index-\(k\) stratum and has a uniform reduced
singular-value floor
\begin{equation}
 s_{\rm nonzero}(D_b)\ge \tau_*>0.
 \label{eq:v12-reduced-floor}
\end{equation}
Write
\begin{equation}
 K_b:=\ker D_b,
 \qquad
 C_b:=\ker D_b^*,
 \qquad
 k_+=\dim K_b,
 \qquad
 k_-=\dim C_b.
 \label{eq:v12-KC}
\end{equation}
On the minimal stratum \(k_+=\max(k,0)\) and
\(k_-=\max(-k,0)\).

\subsection{Global Cech--Kato patching of the zero-mode bundles}
\label{sec:v12-Cech-Kato}

The V11 Riesz formulas are global even when a single Kato frame is not.
Consequently the kernel and cokernel themselves patch without additional
choices.

\begin{theorem}[Global kernel and cokernel bundles]
\label{thm:v12-global-KC}
Under \eqref{eq:v12-reduced-floor}, the V11 Riesz projectors
\(\Pi_K(b)\) and \(\Pi_C(b)\) define smooth Hermitian vector bundles
\begin{equation}
 K\longrightarrow B,
 \qquad
 C\longrightarrow B
 \label{eq:v12-global-KC}
\end{equation}
of ranks \(k_+\) and \(k_-\).  On every finite good cover
\((U_\alpha)_\alpha\) of \(B\), choose V11 unitary Kato frames
\(e^K_\alpha,e^C_\alpha\).  On an overlap,
\begin{equation}
 e^K_\beta=e^K_\alpha u^K_{\alpha\beta},
 \qquad
 e^C_\beta=e^C_\alpha u^C_{\alpha\beta},
 \label{eq:v12-frame-transition}
\end{equation}
with
\(u^K_{\alpha\beta}:U_{\alpha\beta}\to U(k_+)\) and
\(u^C_{\alpha\beta}:U_{\alpha\beta}\to U(k_-)\), satisfying the exact
Cech cocycle identities on triple overlaps.  Every fixed source derivative of
the transition maps required by the coefficient packet is bounded on the
compact cover.
\end{theorem}

\begin{proof}
The Riesz-projector formulas of V11 use a common contour separated from the
nonzero singular spectrum by \(\tau_*^2\).  They are therefore globally smooth
functions of \(b\), and their constant ranks give smooth subbundles of the
ambient finite-cutoff bundles.  On each contractible \(U_\alpha\), V11 Kato
transport gives orthonormal frames.  Two orthonormal frames of the same fibre
differ by a unique unitary matrix, giving
\eqref{eq:v12-frame-transition}.  Uniqueness immediately gives
\[
 u^K_{\alpha\beta}u^K_{\beta\gamma}
 =u^K_{\alpha\gamma},
 \qquad
 u^C_{\alpha\beta}u^C_{\beta\gamma}
 =u^C_{\alpha\gamma}.
\]
The fixed source derivatives of the frames are bounded by the V11 Riesz/Kato
calculus; differentiating
\(u^K_{\alpha\beta}=(e^K_\alpha)^*e^K_\beta\) and similarly for \(C\)
gives the claimed transition bounds.
\end{proof}

We use the V11 determinant convention for the zero-mode factor,
\begin{equation}
 \boxed{
 \mathcal Z_D:=\det C\otimes(\det K)^*.}
 \label{eq:v12-zero-line}
\end{equation}
If one uses the opposite Fredholm-line convention, every formula below is
dualized and nothing substantive changes.

\begin{corollary}[Exact zero-mode line cocycle]
\label{cor:v12-zero-line-cocycle}
Let
\begin{equation}
 z_\alpha:=
 \det(e^C_\alpha)\otimes\det(e^K_\alpha)^*
 \label{eq:v12-local-z}
\end{equation}
be the local unit frame of \(\mathcal Z_D\).  Then
\begin{equation}
 z_\beta=g_{\alpha\beta}z_\alpha,
 \qquad
 \boxed{
 g_{\alpha\beta}
 =\frac{\det u^C_{\alpha\beta}}
        {\det u^K_{\alpha\beta}}
 \in U(1),}
 \label{eq:v12-zero-line-transition}
\end{equation}
and
\begin{equation}
 g_{\alpha\beta}g_{\beta\gamma}=g_{\alpha\gamma}.
 \label{eq:v12-zero-line-cocycle}
\end{equation}
Thus the V11 local saturation writers patch as sections of
\(\mathcal Z_D^*\), not as unrelated scalar functions.
\end{corollary}

\begin{proof}
Take top exterior powers in
\eqref{eq:v12-frame-transition}.  The determinant of the cokernel frame
contributes \(\det u^C_{\alpha\beta}\), while dualizing the kernel exterior
power contributes \((\det u^K_{\alpha\beta})^{-1}\).  The triple-overlap
identity follows by taking determinants of the vector-bundle cocycles.
\end{proof}

\subsection{The zero-mode Berry connection and its exact obstruction}
\label{sec:v12-zero-connection}

The local Kato frames also produce a canonical unitary connection on the
zero-mode line.  Define the imaginary-valued one-forms
\begin{equation}
 \mathcal A_\alpha^K
 :=\operatorname{Tr}\bigl((e_\alpha^K)^*\,d e_\alpha^K\bigr),
 \qquad
 \mathcal A_\alpha^C
 :=\operatorname{Tr}\bigl((e_\alpha^C)^*\,d e_\alpha^C\bigr),
 \label{eq:v12-Berry-forms}
\end{equation}
and
\begin{equation}
 \mathcal A_\alpha^Z:=\mathcal A_\alpha^C-\mathcal A_\alpha^K.
 \label{eq:v12-zero-connection}
\end{equation}

\begin{proposition}[Connection patching and curvature]
\label{prop:v12-zero-connection}
On overlaps,
\begin{equation}
 \boxed{
 \mathcal A_\beta^Z
 =\mathcal A_\alpha^Z+d\log g_{\alpha\beta}.}
 \label{eq:v12-connection-transform}
\end{equation}
Hence the local forms define a unitary connection \(\nabla^Z\) on
\(\mathcal Z_D\).  Its curvature is the globally defined two-form
\begin{equation}
 \boxed{
 \mathcal F_Z
 =\operatorname{Tr}(\Pi_C\,d\Pi_C\wedge d\Pi_C)
  -\operatorname{Tr}(\Pi_K\,d\Pi_K\wedge d\Pi_K),}
 \label{eq:v12-zero-curvature}
\end{equation}
with the usual harmless convention-dependent factor of \(i\) if one uses
real-valued connection forms.  In particular
\begin{equation}
 c_1(\mathcal Z_D)
 =\left[\frac{i}{2\pi}\mathcal F_Z\right]
 \in H^2(B;\mathbb Z).
 \label{eq:v12-c1-zero-line}
\end{equation}
\end{proposition}

\begin{proof}
For a unitary frame change \(e_\beta=e_\alpha u\),
\[
 \operatorname{Tr}(e_\beta^*de_\beta)
 =\operatorname{Tr}(e_\alpha^*de_\alpha)
  +\operatorname{Tr}(u^{-1}du)
 =\operatorname{Tr}(e_\alpha^*de_\alpha)+d\log\det u.
\]
Subtract the kernel formula from the cokernel formula and use
\eqref{eq:v12-zero-line-transition}.  The curvature formula is the standard
finite-rank projector identity, obtained by differentiating the orthogonal
projector \(ee^*\), using \(e^*e=I\), and taking the trace.  Chern--Weil theory
then gives \eqref{eq:v12-c1-zero-line}.
\end{proof}

\begin{theorem}[Global saturation criterion]
\label{thm:v12-global-saturation-criterion}
A nowhere-zero global complex saturation writer
\begin{equation}
 \omega_Z\in\Gamma(\mathcal Z_D^*)
 \label{eq:v12-global-writer}
\end{equation}
exists if and only if \(\mathcal Z_D\) is topologically trivial.  For a
paracompact background this is equivalent to
\begin{equation}
 c_1(\mathcal Z_D)=0.
 \label{eq:v12-line-trivial-c1}
\end{equation}
If, in addition, \(\nabla^Z\) is flat, a global nonzero \emph{parallel}
writer exists if and only if its holonomy character
\begin{equation}
 \operatorname{Hol}_Z:\pi_1(B)\longrightarrow U(1)
 \label{eq:v12-hol-character}
\end{equation}
is trivial.
\end{theorem}

\begin{proof}
A nonzero section of a complex line bundle trivializes it, and a trivial line
has an everywhere nonzero section.  Complex line bundles over a paracompact
base are classified by their first Chern class.  If the connection is flat,
parallel transport of one nonzero fibre vector is path independent exactly
when every loop has unit holonomy; this then produces a global parallel
section.  The converse is immediate.
\end{proof}

\begin{firewall}[Local anomaly cancellation need not flatten the zero-mode Berry connection by itself]
The complete chiral determinant line factors into the zero-mode line and the
reduced invertible line.  Vanishing of the curvature of the \emph{complete}
Standard-Model determinant connection therefore implies cancellation between
the two factor curvatures; it does not require the raw Kato connection
\(\nabla^Z\) on the zero-mode line to have zero curvature separately.  Only
the combined line-valued amplitude is physical before a factorized
normalization is chosen.
\end{firewall}

\subsection{The faithful-spin Standard Model supplies a global complex zero-mode writer}
\label{sec:v12-faithful-global-writer}

The inherited Grand Tensor result now closes the complex line-level problem on
the ordinary-spin faithful Standard-Model branch.  We formulate only the new
consequence for the V11 factorization.

Let \(\mathcal L_{\rm full}\to B\) denote the complete complex Standard-Model
chiral determinant line, and let \(\mathcal L_{\rm red}\to B\) be the reduced
nonzero-mode determinant line.  V11 gives a canonical line isomorphism
\begin{equation}
 \Phi:\mathcal L_{\rm full}
 \xrightarrow{\ \simeq\ }
 \mathcal Z_D\otimes\mathcal L_{\rm red}.
 \label{eq:v12-line-factorization}
\end{equation}
The reduced invertible map \(D^\perp\) has the nowhere-zero determinant section
\begin{equation}
 s_{\rm red}:=\det(D^\perp)
 \in\Gamma(\mathcal L_{\rm red}).
 \label{eq:v12-red-section}
\end{equation}

\begin{theorem}[Global complex zero-mode saturation on the faithful ordinary-spin branch]
\label{thm:v12-global-SM-writer}
Assume the inherited ordinary-spin \(G_6\) Standard-Model hypotheses under
which the complete determinant line has a unit parallel section
\begin{equation}
 s_{\rm full}\in\Gamma(\mathcal L_{\rm full}),
 \qquad
 \nabla^{\rm full}s_{\rm full}=0,
 \label{eq:v12-full-parallel}
\end{equation}
unique up to one constant phase on the connected component.  On any V11
minimal fixed-index branch with reduced floor \eqref{eq:v12-reduced-floor},
define
\begin{equation}
 \boxed{
 z_Z:=\Phi(s_{\rm full})\otimes s_{\rm red}^{-1}
 \in\Gamma(\mathcal Z_D).}
 \label{eq:v12-global-z}
\end{equation}
Then \(z_Z\) is smooth and nowhere zero.  Consequently
\(\mathcal Z_D\) is globally trivial and there is a canonical dual saturation
writer \(\omega_Z\), unique after fixing the constant phase of
\(s_{\rm full}\), such that
\begin{equation}
 \boxed{\langle\omega_Z,z_Z\rangle=1.}
 \label{eq:v12-global-pairing}
\end{equation}
All retained source derivatives patch globally.  Under exact shell elimination
and cutoff transport, \(z_Z\) and \(\omega_Z\) obey the same exact line cocycle
as the inherited complete determinant section; no additional complex
zero-mode holonomy remains.
\end{theorem}

\begin{proof}
The inherited faithful-spin theorem supplies the nowhere-zero parallel section
\(s_{\rm full}\).  The V11 reduced floor makes \(D^\perp\) invertible at every
point of \(B\), so its top exterior map is a nowhere-zero smooth section
\(s_{\rm red}\) of the reduced line.  The line factorization
\eqref{eq:v12-line-factorization} therefore makes
\eqref{eq:v12-global-z} a nowhere-zero smooth section of \(\mathcal Z_D\).
Its fibrewise dual normalized by \eqref{eq:v12-global-pairing} is smooth and
nowhere zero.

On a Kato cover, the local representatives of \(s_{\rm full}\),
\(s_{\rm red}\), and \(z_Z\) multiply by the corresponding transition
functions.  Differentiation therefore produces exactly the Cech/contact terms
already stored in V7 and V11; the transition factors cancel in the pairing
\eqref{eq:v12-global-pairing}.  The inherited complete line comparisons compose
exactly, while V11's reduced Schur determinants compose exactly.  Taking their
quotient shows that the zero-mode factors compose exactly as well.  Thus no
new loop character is introduced by the V11 splitting.
\end{proof}

\begin{corollary}[Factorized holonomy cancellation]
\label{cor:v12-factor-holonomy}
For any loop \(\gamma\subset B\), the factorized connections satisfy
\begin{equation}
 \operatorname{Hol}_{\rm full}(\gamma)
 =\operatorname{Hol}_{Z}(\gamma)
  \operatorname{Hol}_{\rm red}(\gamma).
 \label{eq:v12-hol-factor}
\end{equation}
On the inherited faithful ordinary-spin Standard-Model branch,
\(\operatorname{Hol}_{\rm full}(\gamma)=1\), and hence
\begin{equation}
 \boxed{
 \operatorname{Hol}_{Z}(\gamma)
 =\operatorname{Hol}_{\rm red}(\gamma)^{-1}.}
 \label{eq:v12-hol-cancel}
\end{equation}
The scalar saturated product is therefore single valued even when the two
factor connections separately have nontrivial holonomy.
\end{corollary}

\begin{proof}
Holonomy is multiplicative under tensor products of line bundles and inverse
under dualization.  Apply this to \eqref{eq:v12-line-factorization} and the
inherited unit holonomy of the complete determinant line.
\end{proof}

\subsection{Relative line transport between regulators is path independent}
\label{sec:v12-relative-line}

Now let \(t\in[0,1]\) parametrize a smooth native fixed-index regulator family
on \(B\).  Assume the V11 reduced floor and minimal nullities hold uniformly,
so the kernel/cokernel bundles extend over the cylinder
\begin{equation}
 \widetilde B:=B\times[0,1].
 \label{eq:v12-cylinder}
\end{equation}

\begin{theorem}[Cylinder transport of the zero-mode bundle and line]
\label{thm:v12-cylinder-transport}
The endpoint zero-mode bundles are restrictions of one bundle over
\(\widetilde B\).  Kato transport in the regulator parameter gives unitary
bundle isomorphisms
\begin{equation}
 J_K:K_0\xrightarrow{\simeq}K_1,
 \qquad
 J_C:C_0\xrightarrow{\simeq}C_1,
 \label{eq:v12-KC-transport}
\end{equation}
and therefore an induced unitary zero-line map
\begin{equation}
 J_Z:\mathcal Z_{D_0}\xrightarrow{\simeq}\mathcal Z_{D_1}.
 \label{eq:v12-Z-transport}
\end{equation}
On the faithful ordinary-spin Standard-Model branch, the corresponding
\emph{complete} complex determinant-line identification is independent of the
chosen fixed-index native regulator path with the same endpoints.  The induced
saturated scalar comparison is therefore path independent as well.
\end{theorem}

\begin{proof}
The Riesz projectors are smooth on \(\widetilde B\), hence define bundles there.
The Kato ODE in \(t\) gives \eqref{eq:v12-KC-transport}, and taking determinant
exterior powers gives \eqref{eq:v12-Z-transport}.  For two regulator paths with
the same endpoints, traverse one path and the reverse of the other.  This gives
a closed loop in the enlarged faithful-spin determinant family.  The inherited
eta/bordism theorem gives unit holonomy for the complete Standard-Model line on
every such paired loop.  Hence the two complete line identifications agree.
The reduced determinant section and the factorization theorem then force the
same equality for the factorized saturated amplitude.
\end{proof}

\subsection{The determinant line is not the complete fixed-index projector invariant}
\label{sec:v12-full-zero-bundle}

The previous theorem closes the complex amplitude line on the faithful branch,
but it does not classify the family inside the V11 minimal stratum.  The full
kernel or cokernel bundle is the relevant projector datum.

Let \(\mathcal M_k\) denote the finite-dimensional minimal index-\(k\) stratum
of V11.  Its bundle projection is
\begin{equation}
 \rho_k:\mathcal M_k\longrightarrow
 \begin{cases}
 \operatorname{Gr}(k,\mathcal H_-),&k>0,\\
 \operatorname{Gr}(-k,\mathcal H_+),&k<0,
 \end{cases}
 \label{eq:v12-rho-k}
\end{equation}
with affine vector-space fibres.

\begin{theorem}[Exact global projector classification inside the minimal stratum]
\label{thm:v12-minimal-classification}
For every compact CW parameter space \(B\), the projection \(\rho_k\) induces a
bijection of homotopy classes
\begin{equation}
 \boxed{
 [B,\mathcal M_k]
 \cong
 \begin{cases}
 [B,\operatorname{Gr}(k,\mathcal H_-)],&k>0,\\
 [B,\operatorname{Gr}(-k,\mathcal H_+)],&k<0.
 \end{cases}}
 \label{eq:v12-homotopy-classification}
\end{equation}
Thus two minimal fixed-index projector families are homotopic through the
minimal stratum if and only if their kernel-plane classifying maps
(respectively cokernel-plane maps) are homotopic.  Once such a zero-mode
homotopy is fixed, the affine nonzero-mode graph coordinates create no further
obstruction.
\end{theorem}

\begin{proof}
V11 proves that \(\mathcal M_k\) is the total space of a vector bundle over the
Grassmannian in \eqref{eq:v12-rho-k}.  Fibrewise multiplication by \(s\in[0,1]\)
is a strong deformation retraction onto the zero section.  Hence \(\rho_k\) is
a homotopy equivalence, which induces the claimed bijection on homotopy classes
from every CW source.  The final statement is the explicit fibre contraction.
\end{proof}

\begin{corollary}[The zero-mode line is only a shadow of the full bundle datum]
\label{cor:v12-line-shadow}
The determinant-line class
\begin{equation}
 [\mathcal Z_D]
 =\bigl[\det C\otimes(\det K)^*\bigr]
 \label{eq:v12-line-shadow}
\end{equation}
is a functor of the zero-mode bundle class, but it need not determine that
class.  Consequently equality or triviality of the global complex determinant
line is not, by itself, a projector-level connectedness theorem when
\(|k|>1\).
\end{corollary}

\begin{proof}
Taking the top exterior power forgets all characteristic data not seen by the
determinant.  For example, over \(S^4\) a rank-two complex bundle obtained from
a clutching map \(S^3\to SU(2)\) of degree one has
\(c_1=0\) and hence trivial determinant line, but has nonzero \(c_2\) and is
not isomorphic, hence not homotopic as a classifying map, to the trivial
rank-two bundle.  Theorem~\ref{thm:v12-minimal-classification} therefore
distinguishes the corresponding minimal-stratum families although their
determinant lines agree.
\end{proof}

This observation is load-bearing: the inherited faithful-spin eta theorem
trivializes the \emph{complex amplitude line}; it does not erase the full
zero-mode bundle from the microscopic phase-lifting problem.

\subsection{A finite-cutoff stable range for zero-mode classifying maps}
\label{sec:v12-stable-range}

The full Grassmannian obstruction simplifies when the ambient finite-cutoff
space is large compared with the dimension of the parameter panel.  This is
the regime relevant to a fixed finite background/source panel as the cutoff is
refined.

\begin{lemma}[Connectivity of the finite complex Grassmannian]
\label{lem:v12-Gr-connectivity}
Let
\begin{equation}
 i:\operatorname{Gr}(r,\mathbb C^m)
 \longrightarrow
 \operatorname{Gr}(r,\mathbb C^\infty)=BU(r)
 \label{eq:v12-Gr-inclusion}
\end{equation}
be the standard inclusion.  Then \(i\) is
\((2(m-r)+1)\)-connected.  Consequently, if \(B\) is a finite CW complex of
dimension
\begin{equation}
 \dim B\le 2(m-r),
 \label{eq:v12-stable-range-dim}
\end{equation}
then
\begin{equation}
 [B,\operatorname{Gr}(r,\mathbb C^m)]
 \longrightarrow [B,BU(r)]
 \label{eq:v12-Gr-map-bijection}
\end{equation}
is a bijection.
\end{lemma}

\begin{proof}
Use the Schubert CW decomposition.  The finite Grassmannian contains exactly
the cells indexed by partitions \(\lambda\) with at most \(r\) rows and
\(\lambda_1\le m-r\).  The first cell present in \(BU(r)\) but absent from the
finite Grassmannian corresponds to \(\lambda=(m-r+1)\), of complex dimension
\(m-r+1\) and hence real dimension \(2(m-r+1)\).  Therefore the relative CW
pair has no cells below that real dimension, so the inclusion is
\((2(m-r)+1)\)-connected.  Standard cellular obstruction theory then gives a
bijection on homotopy classes from any CW complex of dimension strictly below
the connectivity, which is exactly
\eqref{eq:v12-stable-range-dim}.
\end{proof}

\begin{theorem}[Bundle isomorphism implies projector homotopy in the stable ambient range]
\label{thm:v12-bundle-to-projector}
Suppose \(B\) has finite CW dimension \(d\).  For \(k>0\), let the negative
chiral ambient space have dimension \(m\) and assume
\begin{equation}
 d\le 2(m-k).
 \label{eq:v12-positive-stable-range}
\end{equation}
If two V11 minimal index-\(k\) projector families have isomorphic kernel bundles
\(K_0\simeq K_1\), then they are homotopic through the minimal stratum.  For
\(k<0\), the analogous statement holds for isomorphic cokernel bundles under
\(d\le2(m-|k|)\).
\end{theorem}

\begin{proof}
An isomorphism of rank-\(|k|\) complex bundles means that their classifying
maps become homotopic in \(BU(|k|)\).  Under the stable-range hypothesis,
Lemma~\ref{lem:v12-Gr-connectivity} says that the finite Grassmannian
classifying maps are already homotopic before stabilization.  Apply
Theorem~\ref{thm:v12-minimal-classification} to lift that zero-mode homotopy
through the affine minimal stratum.
\end{proof}

\begin{corollary}[Why the finite-cutoff ambient rank is not the main topological bottleneck]
\label{cor:v12-large-cutoff-stable}
For a fixed finite-dimensional parameter panel \(B\) and fixed index \(k\), the
stable-range inequality of Theorem~\ref{thm:v12-bundle-to-projector} holds at
all sufficiently large cutoffs whenever the ambient chiral dimension grows
with the regulator.  In that regime the projector-level obstruction is the
ordinary isomorphism class of the finite-rank zero-mode bundle.
\end{corollary}

\begin{proof}
The left-hand side \(d\) and \(|k|\) are fixed while \(m\to\infty\).  Hence
\(2(m-|k|)\ge d\) eventually.
\end{proof}

\subsection{Why an abstract zero-mode homotopy is not a native locality theorem}
\label{sec:v12-locality-obstruction}

The preceding topology still does not solve the microscopic problem.  A
zero-mode projector is typically a global operator in spacetime.  Rotating it
directly can destroy every cutoff-uniform locality moment even though its rank
is finite.

\begin{proposition}[Countertheorem: finite rank does not imply cutoff-uniform quasilocality]
\label{cth:v12-finite-rank-nonlocal}
Let
\(\Lambda_L=(\mathbb Z/L\mathbb Z)^d\) and
\(V=L^d\).  Let
\begin{equation}
 u_0(x)=V^{-1/2},
 \qquad
 (P_0f)(x)=u_0(x)\sum_y\overline{u_0(y)}f(y).
 \label{eq:v12-constant-projector}
\end{equation}
Then \(P_0\) has rank one and operator norm one, but for every \(s>0\) its
weighted row locality norm obeys
\begin{equation}
 \boxed{
 \sup_x\sum_{y\in\Lambda_L}
 (1+d_L(x,y))^s|P_0(x,y)|
 \ge c_{d,s}L^s.}
 \label{eq:v12-projector-nonlocal}
\end{equation}
Hence the family is not uniformly quasilocal as \(L\to\infty\).
Moreover, if
\(u_1(x)=V^{-1/2}e^{2\pi ix_1/L}\) and
\(u_t=\cos t\,u_0+\sin t\,u_1\), then the Kato generator of the rank-one
projector \(P_t=|u_t\rangle\langle u_t|\) has the same volume-growing weighted
moments.  Thus a smooth finite-rank zero-mode rotation does not by itself
produce a native local regulator homotopy.
\end{proposition}

\begin{proof}
The kernel is \(P_0(x,y)=V^{-1}\).  A fixed positive fraction of the lattice
sites lie at periodic distance at least \(cL\) from any chosen \(x\).  Therefore
\[
 V^{-1}\sum_y(1+d_L(x,y))^s
 \ge c' (1+cL)^s,
\]
which proves \eqref{eq:v12-projector-nonlocal}.  For \(P_t\),
\(\dot P_t=|\dot u_t\rangle\langle u_t|+|u_t\rangle\langle\dot u_t|\), and the
Kato generator \([\dot P_t,P_t]\) is a finite sum of rank-one kernels built
from \(u_0,u_1\).  Their pointwise magnitudes are \(O(V^{-1})\) on a positive
fraction of pairs, so the same weighted-moment lower bound follows (up to a
fixed constant for generic \(t\)).
\end{proof}

\begin{firewall}[Do not build the microscopic path by adding zero-mode projectors]
The V11 Riesz/Kato projectors are the correct variables for differentiating the
\emph{finite zero-mode source packet}.  Countertheorem
\ref{cth:v12-finite-rank-nonlocal} shows that they are generally the wrong
variables for constructing the microscopic ultraviolet regulator itself.
A native lift must be built inside the local Wilson-operator algebra; it cannot
be obtained by inserting a finite-rank projector correction and appealing only
to its operator norm.
\end{firewall}

\subsection{The correct native object: the protected local-operator component}
\label{sec:v12-native-component}

We therefore formulate the microscopic problem directly in a locality norm.
Fix an exponential weight \(\alpha>0\), the finite source-derivative order of
\(\mathfrak q\), and a compact background/source panel \(B\).  Let
\(\mathfrak A_{\alpha,\mathfrak q}\) be the real Banach affine space of
Hermitian, exactly gauge/spin-covariant Wilson-kernel families on \(B\) whose
weighted row and column norms and all required source derivatives are finite,
with the common physical principal Dirac--Yukawa jet frozen.  Finite-range
kernels are contained in this space.

Let \(\mathfrak P_{k,\mathfrak q}\subset\mathfrak A_{\alpha,\mathfrak q}\)
be the subset for which:
\begin{enumerate}[label=\textup{(P\arabic*)},leftmargin=3em]
\item the Wilson gap is positive uniformly on \(B\);
\item the full chiral/Yukawa family has minimal index \(k\) and a positive
reduced singular-value floor;
\item the declared complex zero-mode determinant data and, when present, real
Pfaffian orientation stay on the chosen protected component;
\item the V7 finite normalization Jacobian has the required constant rank.
\end{enumerate}

\begin{lemma}[The protected native set is open]
\label{lem:v12-native-open}
Inside the affine space with the physical principal jet frozen,
\(\mathfrak P_{k,\mathfrak q}\) is open in the
\(\mathfrak A_{\alpha,\mathfrak q}\) norm, after fixing one connected Wilson
rank/index component and one real Pfaffian orientation component when relevant.
\end{lemma}

\begin{proof}
A positive Wilson spectral gap is stable under an operator-norm perturbation
smaller than the gap.  On a fixed Wilson-rank component the chiral Fredholm
index is locally constant.  A positive smallest nonzero singular value is
stable under perturbations smaller than that value, so no accidental
kernel--cokernel pair is created.  The finite normalization Jacobian remains of
constant rank under a sufficiently small perturbation because its relevant
minor is nonzero.  Complex line data vary continuously; a real orientation is
locally constant away from a zero.  The weighted locality norm dominates the
operator and the finitely many source norms used in these tests, proving
openness.
\end{proof}

\begin{theorem}[Finite polygonal native lifting inside one protected operator component]
\label{thm:v12-polygonal-native}
Let \(A,B\in\mathfrak P_{k,\mathfrak q}\) lie in the same path-connected
component.  Then there exist finitely many kernels
\begin{equation}
 A=H_0,H_1,\ldots,H_M=B
 \label{eq:v12-polygonal-nodes}
\end{equation}
such that every straight segment
\begin{equation}
 H_{j,t}=(1-t)H_j+tH_{j+1},
 \qquad 0\le t\le1,
 \label{eq:v12-polygonal-segment}
\end{equation}
remains in \(\mathfrak P_{k,\mathfrak q}\).  Every segment is therefore a
locality-preserving V11 fixed-index native homotopy.  If the original kernels
are finite range and all nodes are chosen from one fixed finite-range affine
subspace, the polygonal chain is finite range; in the general case it is
uniformly exponentially local with the same weight \(\alpha\).
\end{theorem}

\begin{proof}
Choose a continuous path \(\gamma:[0,1]\to\mathfrak P_{k,\mathfrak q}\)
from \(A\) to \(B\).  By Lemma~\ref{lem:v12-native-open}, for every \(t\) there
is a norm ball \(B(\gamma(t),r_t)\) contained in the protected set.  Compactness
of \([0,1]\) and uniform continuity of \(\gamma\) give a finite subdivision
\(0=t_0<\cdots<t_M=1\) such that consecutive path points lie in one common
protected ball whose convex hull is still inside that ball.  Set
\(H_j=\gamma(t_j)\).  Norm balls in an affine Banach space are convex, so every
segment \eqref{eq:v12-polygonal-segment} remains protected.

Hermiticity, covariance and the frozen physical principal jet are affine linear
constraints and hence are preserved by the segment.  The exponential locality
norm is convex.  If the path lives in a fixed finite-range coefficient space,
linear interpolation does not enlarge that finite support.
\end{proof}

This theorem is the correct form of a global native lifting statement that can
be proved without classifying the path components of the local operator
algebra.  It is stronger than merely assuming a smooth projector path, because
its hypotheses are posed in the same locality topology used by the shell
estimates.

\subsection{Global fixed-index regulator universality on a native component}
\label{sec:v12-global-universality}

\begin{theorem}[Global complex fixed-index universality on the faithful-spin native component]
\label{thm:v12-global-index-universality}
Fix finite coefficient data \(\mathfrak q\), a compact connected ordinary-spin
faithful \(G_6\) Standard-Model background component \(B\), a fixed index
\(k\), and the positive IR/reference scale.  Let
\(A,B\in\mathfrak P_{k,\mathfrak q}\) lie in the same native path component.
Assume the inherited faithful-family hypotheses for the complete complex
Standard-Model determinant line.  Then:
\begin{enumerate}[label=\textup{(U\arabic*)},leftmargin=3em]
\item the zero-mode complex line admits the global saturation writer of
Theorem~\ref{thm:v12-global-SM-writer};
\item a finite polygonal native chain joining \(A\) to \(B\) exists by
Theorem~\ref{thm:v12-polygonal-native};
\item the complete complex determinant-line comparison along the chain is
path independent and has trivial paired loop holonomy;
\item the V11 relative hard-shell theorem applies on every segment and the
segment matching maps compose to an invertible continuum map
\begin{equation}
 F_{\infty,\mathfrak q,k}^{B\leftarrow A};
 \label{eq:v12-global-map}
\end{equation}
\item the globally saturated coefficientwise continua satisfy
\begin{equation}
 \boxed{
 \widehat{\mathfrak C}_{\infty,\lambda}^{B,(k),\rm sat}
 =F_{\infty,\mathfrak q,k}^{B\leftarrow A}
  \widehat{\mathfrak C}_{\infty,\lambda}^{A,(k),\rm sat}.}
 \label{eq:v12-global-universality}
\end{equation}
\end{enumerate}
The equivalence contains all retained connected, composite, stress/contact,
ghost/antifield, BV/BRST/master and reduced determinant-line coefficients and
is strictly projective in perturbative order.
\end{theorem}

\begin{proof}
Item (U1) is Theorem~\ref{thm:v12-global-SM-writer}; item (U2) is
Theorem~\ref{thm:v12-polygonal-native}.  Every polygonal segment lies in the
V11 admissible fixed-index region and hence carries the reduced
\(aM+\mu/M\) relative-forest estimate and the optimized
\(O(\sqrt{a\mu})\) regulator difference.  The inherited faithful-spin
parallel determinant section and Theorem~\ref{thm:v12-cylinder-transport}
make the complex line identifications independent of the chain and compatible
with the global saturation writer.  Apply V11 local universality to each
segment and compose the finitely many invertible matching maps.  The V7
chronological construction makes each segment map projective in terminal loop
order, hence so is their finite composition.
\end{proof}

\begin{corollary}[The complex line is no longer a load-bearing obstruction on this branch]
\label{cor:v12-complex-line-closed}
Within the hypotheses of Theorem~\ref{thm:v12-global-index-universality}, the
remaining classification problem for general fixed-index regulator
universality is not the complex determinant-line holonomy.  It is whether two
microscopic kernels lie in the same protected path component of the native
local-operator space.  Projector-level zero-mode bundle topology supplies a
necessary obstruction to such a path, but not by itself a locality-preserving
construction.
\end{corollary}

\subsection{The real Pfaffian branch retains a mod-two obstruction}
\label{sec:v12-real-Pfaffian}

The preceding faithful-spin theorem concerns the complete complex chiral
determinant line.  A genuinely real Majorana/Pfaffian sector has a different
global obstruction which must not be erased by complexification.

\begin{proposition}[Real Pfaffian orientation local system]
\label{prop:v12-Pfaffian-w1}
Let \(\mathcal P\to B\) be a real one-dimensional Pfaffian line on a
constant-rank branch.  Its local orientation frames have transition signs
\(\epsilon_{\alpha\beta}\in\{\pm1\}\), defining
\begin{equation}
 w_1(\mathcal P)\in H^1(B;\mathbb Z_2).
 \label{eq:v12-Pf-w1}
\end{equation}
A global real orientation exists if and only if
\(w_1(\mathcal P)=0\).  For a generic loop of real skew operators with only
simple crossings, the orientation holonomy is
\begin{equation}
 \boxed{
 \operatorname{Hol}_{\mathcal P}(\gamma)
 =(-1)^{\operatorname{SF}_2(\gamma)},}
 \label{eq:v12-Pf-sf2}
\end{equation}
where \(\operatorname{SF}_2\) is the parity of the real spectral flow.  The
formula extends by homotopy to arbitrary loops avoiding an endpoint rank
jump.
\end{proposition}

\begin{proof}
The transition signs of a real line form a Cech \(\mathbb Z_2\)-cocycle whose
class is the first Stiefel--Whitney class.  It is a coboundary exactly when the
local orientations can be rephased to agree globally.  At a simple real
crossing one Pfaffian eigenvalue changes sign, reversing the local orientation.
Along a generic loop the total sign is therefore the parity of the number of
crossings.  Pair creation/annihilation under a homotopy changes that number by
an even integer, so the parity extends to the homotopy-invariant mod-two
spectral flow.
\end{proof}

\begin{firewall}[The faithful complex eta theorem does not orient an independent real Pfaffian]
The inherited vanishing of the faithful-spin complex determinant holonomy
cannot be used to set \(w_1(\mathcal P)=0\) for an independently physical real
Majorana line.  Global scalar universality in such a sector requires either a
separate proof that the mod-two character vanishes or a declared physical
orientation local system which is transported unchanged between regulators.
\end{firewall}

\subsection{What is and is not closed by V12}
\label{sec:v12-boundary}

The global picture can now be stated sharply.

\begin{theorem}[Three-level global obstruction theorem]
\label{thm:v12-three-level}
On a V11 minimal fixed-index family the following implications hold.
\begin{enumerate}[label=\textup{(G\arabic*)},leftmargin=3em]
\item A native local-operator homotopy implies a minimal-stratum projector
homotopy and hence equality of the kernel/cokernel classifying-map class.
\item A minimal-stratum projector homotopy implies isomorphism of the
zero-mode determinant lines and equality of all their characteristic and
holonomy data transported by that homotopy.
\item Neither converse holds as a formal consequence of the preceding level:
the determinant line forgets higher zero-mode-bundle data, and an abstract
projector homotopy need not have cutoff-uniform locality.
\end{enumerate}
On the faithful ordinary-spin complex Standard-Model branch the line-level
obstruction is trivialized by the inherited eta/bordism theorem, but the two
stronger levels remain distinct.
\end{theorem}

\begin{proof}
Item (G1) follows by applying spectral functional calculus continuously along
a gapped native Wilson path.  Item (G2) follows by taking top exterior powers
of the kernel/cokernel bundle over the homotopy cylinder.  Failure of the first
converse is Corollary~\ref{cor:v12-line-shadow}; failure of the second as a
proof principle is Countertheorem~\ref{cth:v12-finite-rank-nonlocal}.  The last
sentence is Theorem~\ref{thm:v12-global-SM-writer}.
\end{proof}

\begin{center}
\fbox{\parbox{0.92\textwidth}{
\textbf{V12 closes global complex line transport, but it does not identify
projector equivalence with microscopic locality.}
For the faithful ordinary-spin Standard-Model determinant packet, the
complete complex line has trivial paired holonomy and V11's reduced nonzero
section therefore produces a global zero-mode saturation writer.  Global
fixed-index universality follows for every pair of regulators in the same
protected path component of the native exponentially local operator space.
The remaining load-bearing problem is to characterize those native path
components from intrinsic local data rather than from an already-given native
path.}}
\end{center}

\section{V12 anti-circularity audit}
\label{sec:v12-audit}

\begin{enumerate}[leftmargin=2.4em]
\item \textbf{The Grand Tensor eta theorem is inherited, not rederived.}
V12 uses its unit parallel complex determinant section and exact cutoff cocycle
only to obtain the new zero-mode saturation consequence after V11
factorization.
\item \textbf{The raw Kato zero-mode connection is not silently declared
flat.}
Its curvature is \eqref{eq:v12-zero-curvature}; only the complete determinant
connection has the inherited anomaly/eta trivialization.
\item \textbf{Determinant-line triviality is not promoted to full bundle
triviality.}
The rank-two \(S^4\) example in Corollary~\ref{cor:v12-line-shadow} keeps the
higher zero-mode-bundle obstruction explicit.
\item \textbf{Stable bundle classification is used only in its stated finite
ambient range.}
Outside \eqref{eq:v12-stable-range-dim}, V12 retains the actual finite
Grassmannian classifying map as the projector invariant.
\item \textbf{A finite-rank zero-mode projector is not called local because its
operator norm is bounded.}
Countertheorem~\ref{cth:v12-finite-rank-nonlocal} exhibits the missing
weighted-locality control explicitly.
\item \textbf{The native lift is constructed only inside an already identified
protected local-operator component.}
The polygonal theorem is a Banach-space statement in the native locality norm,
not a reconstruction from an abstract Grassmannian homotopy.
\item \textbf{The real Pfaffian is kept separate.}
Its \(\mathbb Z_2\) orientation character is not inferred from the complex
faithful-spin determinant theorem.
\end{enumerate}

\section{V12 status ledger}
\label{sec:v12-status}

\subsection*{Proved in V12}
\begin{enumerate}[leftmargin=2.2em]
\item The local V11 kernel/cokernel Kato frames patch to global smooth bundles
on every protected constant-rank component.  Their transition matrices have
all required fixed source derivatives.
\item The zero-mode determinant factor
\(\mathcal Z_D=\det C\otimes(\det K)^*\) has the exact unitary Cech cocycle
\(g_{\alpha\beta}=\det u^C_{\alpha\beta}/\det u^K_{\alpha\beta}\), a global
Kato/Berry connection and curvature
\(\operatorname{Tr}(\Pi_Cd\Pi_C\wedge d\Pi_C)-
  \operatorname{Tr}(\Pi_Kd\Pi_K\wedge d\Pi_K)\).
\item On the inherited faithful ordinary-spin \(G_6\) Standard-Model branch,
the complete parallel determinant section divided by the global nonzero
reduced determinant section gives a nowhere-zero global zero-mode section and
hence a global saturation writer.  The factorized holonomies cancel exactly.
\item Along a smooth fixed-index native regulator cylinder, Kato transport
gives global kernel/cokernel and zero-line endpoint isomorphisms.  On the
faithful complex branch the complete line identification and saturated scalar
comparison are path independent.
\item The homotopy class of a family inside the V11 minimal stratum is exactly
the homotopy class of its kernel-plane map for \(k>0\) or cokernel-plane map
for \(k<0\).  The determinant line alone is not complete when higher
zero-mode-bundle data are present.
\item The inclusion
\(\operatorname{Gr}(r,\mathbb C^m)\to BU(r)\) is
\((2(m-r)+1)\)-connected.  Hence on a finite parameter panel of dimension
\(d\le2(m-r)\), isomorphic zero-mode bundles already give homotopic finite
Grassmannian classifying maps and therefore a minimal-stratum projector
homotopy.
\item Finite rank does not imply native locality: the constant-mode projector
has weighted row moments growing like \(L^s\).  Thus an abstract zero-mode
Kato rotation is not a valid microscopic regulator lift without an independent
locality theorem.
\item The protected fixed-index native operator set is open in the weighted
exponentially local/source norm.  Any two regulators in one path component are
connected by a finite polygonal chain of native protected homotopies.
\item Combining that polygonal chain with V11 and the inherited faithful-spin
line theorem yields global complex fixed-index regulator universality on every
protected native path component, with exact loop-order projectivity.
\item A real Pfaffian line carries the independent obstruction
\(w_1\in H^1(B;\mathbb Z_2)\), with loop sign given by mod-two spectral flow.
\end{enumerate}

\subsection*{Inherited}
\begin{enumerate}[leftmargin=2.2em]
\item The Grand Tensor finite zero-mode factorization/saturation theorem,
Fredholm determinant line, faithful-spin \(G_6\) eta/bordism trivialization,
and exact smooth cutoff line cocycle.
\item V1--V7: the source-complete fixed-order continuum, common BV restoration,
forest theorem, UV Cauchy convergence, scheme/reference covariance and exact
loop-order projectivity.
\item V8--V10: regulator-marked relative forests, the two-scale universality
estimate, and native regulator-class lifting on the protected index-zero
single-cone component.
\item V11: minimal fixed-index geometry, reduced pseudoinverse/source calculus,
zero-mode-preserving shell factorization, saturation writers and local
fixed-index universality.
\end{enumerate}

\subsection*{Conditional}
\begin{enumerate}[leftmargin=2.2em]
\item The global complex saturation theorem uses the inherited ordinary-spin
faithful \(G_6\) family hypotheses.  Other spin--gauge structures require their
own eta/holonomy calculation.
\item The finite-Grassmannian stable-range theorem assumes a finite-dimensional
CW parameter panel.  It is not a classification theorem for an unrestricted
infinite-dimensional gauge configuration space.
\item Global regulator-class universality is proved for regulators already in
the same path component of the protected native local-operator space.  V12 does
not prove that every projector-level zero-mode-bundle equivalence lifts to that
native component.
\item A real Pfaffian sector still requires a vanishing mod-two orientation
character or an explicitly supplied physical orientation local system.
\item All continuum claims remain coefficientwise at prescribed finite
\((L,\Delta,r,N)\) and positive IR/reference scale.
\end{enumerate}

\subsection*{Open beyond V12}
\begin{enumerate}[leftmargin=2.2em]
\item Classify the path components of the protected exponentially local Wilson
operator space relative to the fixed physical Dirac--Yukawa principal jet.
Equivalently, determine when the zero-mode-bundle/projector homotopy produced by
V11--V12 admits a uniformly local native lift without inserting nonlocal
finite-rank projectors.
\item For a genuinely real Majorana/Pfaffian sector, compute the mod-two
spectral-flow character on the relevant global background component and prove
(or falsify) orientability.
\item Analyze crossings out of the minimal fixed-index stratum, including
creation of accidental zero-mode pairs and true index change, with divisor and
spectral-flow source packets.
\item The independent stronger problems remain: removal of the positive IR
scale, analytic finite-coupling RG trajectories, perturbative summability and
nonperturbative completion.
\end{enumerate}

\subsection*{Exact next load-bearing theorem}
\begin{center}
\fbox{\parbox{0.92\textwidth}{
\textbf{Local-operator phase classification / native lifting.}
Fix the physical Standard-Model principal Dirac--Yukawa jet, index \(k\), the
minimal reduced-gap condition and the global zero-mode bundle class.  Determine
the path components of the corresponding gapped finite-range or uniformly
exponentially local, gauge/spin-covariant Wilson-operator space.  Prove that
projector-equivalent regulators lie in one protected native component, or
exhibit an additional invariant of the local operator algebra.  The proof must
avoid finite-rank zero-mode projector corrections, whose weighted locality
fails by Countertheorem~\ref{cth:v12-finite-rank-nonlocal}.  Combined with
Theorem~\ref{thm:v12-global-index-universality}, such a result would upgrade
V12 from universality within a native component to a complete global
fixed-index regulator-class theorem.}}
\end{center}

\section*{Append-only continuation marker: V12}
\addcontentsline{toc}{section}{Append-only continuation marker: V12}
\textbf{End of V12.}  The complete V1--V11 body above is retained verbatim.
V12 patches the V11 zero-mode bundles globally, identifies the exact Cech and
Berry data of their determinant line, and uses the inherited faithful-spin
Standard-Model eta/bordism trivialization together with the global reduced
determinant section to construct a global complex zero-mode saturation writer.
It proves that the full kernel/cokernel classifying map, not merely its
determinant line, is the projector-level fixed-index invariant; in a finite
stable ambient range, bundle isomorphism already gives projector homotopy.  It
then isolates the remaining microscopic obstruction by showing that finite-rank
zero-mode rotations need not be uniformly quasilocal.  Finally, it proves a
finite polygonal native lifting theorem inside every protected path component
of the exponentially local operator space and obtains global complex
fixed-index regulator universality on that component.  The next gate is the
intrinsic classification of those native local-operator components.



\clearpage
\section{V13 continuation: local Riesz projectors, controlled zero-mode modules, and stable native lifting}
\label{sec:v13-continuation}

V12 isolated the remaining fixed-index problem as a locality problem rather
than a determinant-line problem.  Its counterexample showed correctly that an
\emph{arbitrary} finite-rank projector can be completely nonlocal even when its
operator norm is one.  For the actual V11 protected zero modes, however, there
is one additional hypothesis which was not used in that counterexample: the
zero eigenspace is the isolated spectral subspace of a uniformly local operator
whose remaining singular spectrum has a cutoff-uniform floor.  Going upstream
to that spectral gap changes the local classification problem.

The first result below is therefore a correction of variable, not a
retraction of V12: Riesz projectors of the protected chiral operator are
uniformly exponentially local (with a possibly smaller exponential rate).
Hence Kato transport along an already native fixed-index path is itself local.
The genuine obstruction is whether the endpoint zero-mode bundles admit a
\emph{controlled} local identification.  Ordinary bundle isomorphism and
complex determinant-line triviality forget this metric/locality datum.

V13 proves four statements.
\begin{enumerate}[label=\textup{(V13.\arabic*)},leftmargin=3.8em]
\item A gap plus exponential locality implies exponential locality of the
resolvent, sign, zero-mode Riesz projectors, Moore--Penrose inverse, and every
fixed source derivative, after a harmless loss of exponent.
\item The V11 minimal fixed-index geometry can be reconstructed entirely from
its local zero-mode projector and a local affine graph coordinate.  Thus the
full native phase problem reduces to a finite-rank \emph{controlled zero-mode
module} problem.
\item Ordinary zero-mode bundle data are not complete for this controlled
problem: a pair of macroscopically separated site zero modes has the same rank,
trivial determinant line and uniform reduced gap, but any intertwining partial
isometry has exponentially growing locality norm.
\item If the two zero-mode packets are controlled-isomorphic, then after
adjoining one inert massive spectator copy there is an explicit protected
local homotopy between the two regulators.  The spectator contributes only a
source-independent factor and disappears exactly from the normalized connected
functional.  Therefore controlled-isomorphic regulators are \emph{stably}
regulator-universal at every fixed coefficient order.
\end{enumerate}

The remaining issue after V13 is an unstable cancellation problem: whether the
spectator can always be removed in the exact Standard-Model carrier, or whether
there is a genuine controlled/coarse obstruction to such cancellation.

\subsection{The weighted local kernel algebra}
\label{sec:v13-local-algebra}

Let $\Lambda$ be one of the finite periodic lattices in the regulator sequence,
with periodic distance $d(\cdot,\cdot)$ and uniformly polynomial volume growth.
For a matrix kernel $A=(A_{xy})_{x,y\in\Lambda}$ put
\begin{equation}
 \|A\|_{\alpha}
 :=\max\left\{
 \sup_x\sum_y e^{\alpha d(x,y)}\|A_{xy}\|,
 \sup_y\sum_x e^{\alpha d(x,y)}\|A_{xy}\|
 \right\}.
 \label{eq:v13-local-norm}
\end{equation}
For a compact background/source panel $B$, the same notation means the
supremum over $b\in B$ together with the finitely many source derivatives
required by the fixed coefficient packet $\mathfrak q$.  Exact gauge/spin
covariance and the fixed representation grading are imposed as closed linear
constraints.

\begin{lemma}[Weighted Schur algebra]
\label{lem:v13-Schur-algebra}
For every $\alpha>0$, kernels with finite norm
\eqref{eq:v13-local-norm} form a Banach $*$-algebra and
\begin{equation}
 \boxed{
 \|AB\|_{\alpha}\le \|A\|_{\alpha}\|B\|_{\alpha},
 \qquad
 \|A^*\|_{\alpha}=\|A\|_{\alpha}.}
 \label{eq:v13-Schur-product}
\end{equation}
Moreover $\|A\|_{\ell^2\to\ell^2}\le \|A\|_{\alpha}$.
The same statements hold for the finite source-jet norm.
\end{lemma}

\begin{proof}
The triangle inequality gives
$d(x,z)\le d(x,y)+d(y,z)$ and hence
$e^{\alpha d(x,z)}\le e^{\alpha d(x,y)}e^{\alpha d(y,z)}$.
Insert the kernel product
$(AB)_{xz}=\sum_yA_{xy}B_{yz}$ and sum first in $z$ and then in $y$.
The column estimate is identical.  Taking adjoints exchanges the two Schur
rows.  The ordinary Schur test gives the $\ell^2$ operator bound.  A fixed
source derivative is a finite Leibniz sum, so the same argument applies after
summing the finitely many derivative seminorms.
\end{proof}

The fixed exponent is not preserved by every inverse with the same numerical
value.  What is preserved under a uniform spectral gap is \emph{some} positive
exponent.  This is precisely what the shell proof needs.

\begin{theorem}[Gap-to-resolvent locality with exponent loss]
\label{thm:v13-resolvent-locality}
Let $A=A^*$ obey
\begin{equation}
 \|A\|_{\alpha}\le M,
 \qquad
 \operatorname{dist}(z,\operatorname{spec}A)\ge\delta>0
 \label{eq:v13-resolvent-hyp}
\end{equation}
with constants uniform in cutoff, volume and the chosen compact source panel.
Then there exists
\begin{equation}
 0<\beta=\beta(\alpha,M,\delta)<\alpha
 \label{eq:v13-beta}
\end{equation}
for which
\begin{equation}
 \boxed{
 \|(z-A)^{-1}\|_{\beta}
 \le C_{\alpha,M,\delta}.}
 \label{eq:v13-resolvent-local}
\end{equation}
The same conclusion holds for every prescribed fixed source derivative of the
resolvent if the corresponding source jets of $A$ are bounded in
$\|\cdot\|_\alpha$.
\end{theorem}

\begin{proof}
Fix a lattice point $x_0$ and let
$\phi_{x_0}(x)=d(x,x_0)$.  For $\eta>0$ define the diagonal multiplication
$W_{\eta,x_0}f(x)=e^{\eta\phi_{x_0}(x)}f(x)$.  The conjugated kernel is
\[
 (W_{\eta,x_0}AW_{\eta,x_0}^{-1})_{xy}
 =e^{\eta(\phi_{x_0}(x)-\phi_{x_0}(y))}A_{xy}.
\]
Since $|\phi_{x_0}(x)-\phi_{x_0}(y)|\le d(x,y)$, the Schur norm of the
difference from $A$ is bounded by
\[
 \rho_A(\eta):=
 \max\left\{
 \sup_x\sum_y (e^{\eta d(x,y)}-1)\|A_{xy}\|,
 \sup_y\sum_x (e^{\eta d(x,y)}-1)\|A_{xy}\|
 \right\}.
\]
Uniform exponential summability at rate $\alpha$ implies
$\rho_A(\eta)\to0$ as $\eta\downarrow0$, uniformly on the declared family:
split the sum at a fixed distance $R$, use ordinary dominated convergence on
the finite head, and use the $e^{\alpha d}$ tail beyond $R$.  Choose
$0<\eta<\alpha$ so that $\rho_A(\eta)<\delta/2$.
Then
\[
 \operatorname{dist}\bigl(z,
 \operatorname{spec}(W_{\eta,x_0}AW_{\eta,x_0}^{-1})\bigr)
 \ge\delta/2
\]
in the resolvent sense supplied by the Neumann identity, and therefore
\begin{equation}
 \|W_{\eta,x_0}(z-A)^{-1}W_{\eta,x_0}^{-1}\|_{2\to2}
 \le 2/\delta.
 \label{eq:v13-CT-l2}
\end{equation}
Applying \eqref{eq:v13-CT-l2} to a coordinate vector at $x_0$ gives a uniform
exponentially weighted $\ell^2$ column bound for the resolvent.  Apply the same
argument to the adjoint for rows.  Choose $0<\beta<\eta$.  Polynomial lattice
volume growth gives
\[
 \sup_x\sum_y e^{-2(\eta-\beta)d(x,y)}<\infty
\]
uniformly in the finite periodic volumes.  Cauchy--Schwarz therefore converts
the weighted $\ell^2$ row and column bounds into
\eqref{eq:v13-resolvent-local}.

For source derivatives use
\[
 \partial(z-A)^{-1}
 =(z-A)^{-1}(\partial A)(z-A)^{-1}
\]
and iterate.  Every fixed derivative is a finite sum of resolvent words.
Apply Lemma~\ref{lem:v13-Schur-algebra}, decreasing the exponent once if
necessary to a common positive value for the finite family of words.
\end{proof}

\begin{corollary}[Local holomorphic/Riesz calculus on a protected spectral split]
\label{cor:v13-local-functional-calculus}
Let $A=A^*$ be exponentially local and suppose its spectrum is separated into
two closed sets by a uniform gap.  Every Riesz projector defined by a contour
inside that gap is exponentially local at some uniform positive rate, together
with every prescribed fixed source derivative.  The same is true of
$\operatorname{sgn}A$ when zero is uniformly excluded from the spectrum.
\end{corollary}

\begin{proof}
Use the contour representation
\[
 P_\Gamma={1\over2\pi i}\oint_\Gamma(z-A)^{-1}\,dz
\]
and Theorem~\ref{thm:v13-resolvent-locality}.  A fixed compact contour permits
one common exponent and constant.  The sign is a difference of the positive
and negative Riesz projectors.
\end{proof}

\subsection{Protected zero-mode projectors are local}
\label{sec:v13-zero-local}

We now apply the preceding theorem to the actual V11 zero modes.  This is the
load-bearing distinction from an arbitrary finite-rank projector.

\begin{theorem}[Uniform locality of V11 kernel/cokernel projectors]
\label{thm:v13-zero-projector-local}
Let $D_b$ be a V11 protected constant-rank chiral/Yukawa family.  Assume its
kernel is generated by a native finite-range or exponentially local operator
family and its nonzero singular values obey
\begin{equation}
 s_{\rm nonzero}(D_b)\ge\tau_*>0.
 \label{eq:v13-tau}
\end{equation}
Then there is a cutoff-uniform exponent $\beta>0$ such that
\begin{equation}
 \boxed{
 \|\Pi_K\|_\beta+\|\Pi_C\|_\beta
 \le C,}
 \label{eq:v13-zero-projector-local}
\end{equation}
and the same estimate holds for every fixed source derivative required by
$\mathfrak q$.  The Moore--Penrose inverse $D^\dagger$ is exponentially local
at a possibly smaller positive exponent.
\end{theorem}

\begin{proof}
The products $D^*D$ and $DD^*$ are exponentially local by
Lemma~\ref{lem:v13-Schur-algebra}.  Their spectra are contained in
$\{0\}\cup[\tau_*^2,\infty)$.  Choose a circle around zero of radius
$\tau_*^2/2$.  The V11 Riesz formulas for $\Pi_K$ and $\Pi_C$ and
Corollary~\ref{cor:v13-local-functional-calculus} prove
\eqref{eq:v13-zero-projector-local}, including source derivatives.

For the pseudoinverse use the V11 formula
\[
 D^\dagger=(D^*D+\Pi_K)^{-1}D^*.
\]
The positive operator $D^*D+\Pi_K$ has spectral floor
$\min\{1,\tau_*^2\}$ and is exponentially local.  Apply
Theorem~\ref{thm:v13-resolvent-locality} at $z=0$ and multiply by $D^*$.
\end{proof}

\begin{firewall}[The exact scope of the V12 nonlocal projector counterexample]
\label{assess:v13-v12-counterexample}
Countertheorem~\ref{cth:v12-finite-rank-nonlocal} remains correct: finite rank
alone does not imply locality.  The constant-mode projector used there is not
the isolated kernel projector of a uniformly local family with a
cutoff-uniform reduced spectral floor.  Theorem
\ref{thm:v13-zero-projector-local} shows that the V11 protected zero-mode
packet has additional analytic structure which rules out that particular
failure mode.  It does \emph{not} imply that two endpoint zero-mode projectors
are connected by a uniformly local path; that is the separate controlled
module question treated below.
\end{firewall}

\subsection{Kato transport along a native path is itself native-local}
\label{sec:v13-local-Kato}

\begin{theorem}[Local Kato transport]
\label{thm:v13-local-Kato}
Let $t\mapsto D_t$ be a $C^1$ V11 fixed-index native homotopy on a compact
parameter interval with one common locality exponent at the microscopic level
and one common reduced singular-value floor $\tau_*>0$.  Then for some
$\beta>0$ the zero-mode projectors and their $t$ derivatives satisfy
\begin{equation}
 \sup_t\bigl(
 \|\Pi_t\|_\beta+\|\dot\Pi_t\|_\beta
 \bigr)<\infty.
 \label{eq:v13-Pi-dot-local}
\end{equation}
The Kato equation
\begin{equation}
 \dot U_t=[\dot\Pi_t,\Pi_t]U_t,
 \qquad U_0=I,
 \label{eq:v13-local-Kato-ODE}
\end{equation}
has a unitary solution with
\begin{equation}
 \boxed{
 \sup_t\bigl(
 \|U_t\|_\beta+\|U_t^{-1}\|_\beta
 \bigr)<\infty,
 \qquad
 \Pi_t=U_t\Pi_0U_t^*.}
 \label{eq:v13-local-Kato-result}
\end{equation}
All prescribed fixed source derivatives obey the same statement after one
finite exponent loss.
\end{theorem}

\begin{proof}
Theorem~\ref{thm:v13-zero-projector-local} applied to the two-parameter family
$(t,b)$ gives \eqref{eq:v13-Pi-dot-local}; equivalently one differentiates the
Riesz formula in $t$.  The generator
$G_t=[\dot\Pi_t,\Pi_t]$ is anti-Hermitian and bounded in the Banach algebra by
Lemma~\ref{lem:v13-Schur-algebra}.  The Dyson series for
\eqref{eq:v13-local-Kato-ODE} therefore converges in the same Banach algebra,
with
\[
 \|U_t\|_\beta\le
 \exp\left(\int_0^t\|G_s\|_\beta\,ds\right).
\]
Anti-Hermiticity gives unitarity on $\ell^2$.  The standard Kato computation
shows $U_t\Pi_0U_t^*=\Pi_t$.  Source derivatives satisfy finite inhomogeneous
Dyson equations with coefficients already controlled in the same algebra.
\end{proof}

Thus a native homotopy produces more than an abstract bundle isomorphism: it
produces an exponentially local one.  This motivates the correct invariant.

\subsection{Controlled zero-mode modules}
\label{sec:v13-controlled-module}

Let $\mathcal A_{\exp,\mathfrak q}$ denote the union over positive exponents of
the preceding gauge/spin-covariant weighted algebras, with the fixed source
jet package understood.  The union is used only so that the harmless exponent
losses of Theorem~\ref{thm:v13-resolvent-locality} do not create artificial
phase boundaries.

\begin{definition}[Controlled local equivalence of zero-mode packets]
\label{def:v13-controlled-equivalence}
Two finite-rank zero-mode projector families $\Pi_0,\Pi_1$ on the same compact
background panel are \emph{controlled-equivalent}, written
\begin{equation}
 \Pi_0\sim_{\rm ctl}\Pi_1,
 \label{eq:v13-ctl-equivalence}
\end{equation}
if there is a partial isometry
$V\in\mathcal A_{\exp,\mathfrak q}$ such that
\begin{equation}
 \boxed{
 V^*V=\Pi_0,
 \qquad
 VV^*=\Pi_1,}
 \label{eq:v13-controlled-partial-isometry}
\end{equation}
with all retained background/source derivatives in one common positive
locality class.  The isomorphism class of
$\Pi\mathcal A_{\exp,\mathfrak q}$ is called the
\emph{controlled zero-mode module class}.
\end{definition}

\begin{proposition}[Native homotopy preserves the controlled module class]
\label{prop:v13-native-implies-controlled}
Endpoints of every V11 fixed-index native homotopy have
controlled-equivalent kernel bundles for $k>0$ and controlled-equivalent
cokernel bundles for $k<0$.
\end{proposition}

\begin{proof}
Use the local Kato unitary $U_1$ of
Theorem~\ref{thm:v13-local-Kato}.  For the kernel case set
$V=U_1\Pi_{K,0}$.  Then
$V^*V=\Pi_{K,0}$ and
$VV^*=\Pi_{K,1}$.  The cokernel case is identical.
\end{proof}

There are natural forgetful maps
\begin{equation}
 \boxed{
 \text{controlled module class}
 \longrightarrow
 \text{ordinary zero-mode bundle class}
 \longrightarrow
 \text{determinant/Pfaffian line class}.}
 \label{eq:v13-forgetful-hierarchy}
\end{equation}
V12 proved that the second arrow loses information.  The next proposition shows
that the first arrow can also lose information even when the reduced gap is
uniform.

\begin{proposition}[Macroscopic transport obstruction invisible to the ordinary bundle]
\label{prop:v13-macro-obstruction}
Let $\Lambda_L$ be an even one-dimensional periodic lattice and choose sites
$x_L=0$ and $y_L=L/2$.  Put
\begin{equation}
 \Pi_{0,L}=|\delta_{x_L}\rangle\langle\delta_{x_L}|,
 \qquad
 \Pi_{1,L}=|\delta_{y_L}\rangle\langle\delta_{y_L}|.
 \label{eq:v13-site-projectors}
\end{equation}
Both are range-zero local rank-one projections with trivial complex determinant
line.  They are also exact kernel projectors of the range-zero local operators
\begin{equation}
 D_{i,L}:=I-\Pi_{i,L},
 \label{eq:v13-site-D}
\end{equation}
whose nonzero singular spectrum is identically one.  Nevertheless, for every
fixed $\beta>0$ there is no family of partial isometries $V_L$ satisfying
\eqref{eq:v13-controlled-partial-isometry} with
$\sup_L\|V_L\|_\beta<\infty$.
Indeed any such partial isometry is
\begin{equation}
 V_L=e^{i\theta_L}
 |\delta_{y_L}\rangle\langle\delta_{x_L}|
 \label{eq:v13-site-partial-isometry}
\end{equation}
and hence
\begin{equation}
 \boxed{
 \|V_L\|_\beta=e^{\beta L/2}.}
 \label{eq:v13-site-cost}
\end{equation}
Thus ordinary rank/bundle data and a uniform reduced spectral gap do not by
themselves determine the controlled native phase.
\end{proposition}

\begin{proof}
A partial isometry with initial projection onto the one-dimensional line
$\mathbb C\delta_{x_L}$ and final projection onto
$\mathbb C\delta_{y_L}$ is uniquely of the form
\eqref{eq:v13-site-partial-isometry} up to phase.  Its only nonzero matrix entry
connects sites at periodic distance $L/2$, so its weighted row and column norms
are exactly $e^{\beta L/2}$.  The remaining statements are immediate from the
diagonal formula \eqref{eq:v13-site-D}.
\end{proof}

\begin{firewall}[What the macroscopic example does and does not show]
Proposition~\ref{prop:v13-macro-obstruction} is a native-locality obstruction,
not an additional perturbative gauge anomaly.  It may be embedded as a
spectator defect block next to a common physical Dirac block, so it shows that
``same principal physical jet + same numerical index + same ordinary zero-mode
bundle'' is not a theorem of controlled path connectedness in the unrestricted
local-operator category.  Whether the exact Standard-Model representation and
background grammar exclude such coarse sectors is a further physical
restriction and is not assumed here.
\end{firewall}

\subsection{The affine reduced graph coordinate is local}
\label{sec:v13-local-graph}

The controlled zero-mode class is the only new finite-rank datum needed by the
V11 minimal-stratum geometry.  The remaining graph coordinate can be recovered
by local algebraic operations.

\begin{proposition}[Local graph coordinate for positive index]
\label{prop:v13-positive-local-graph}
Assume $k>0$.  Let $P$ be the positive Wilson projector, let
$\Pi_K$ be the V11 kernel projector, and put
\begin{equation}
 R:=P-\Pi_K.
 \label{eq:v13-R-positive}
\end{equation}
On $\mathcal H_+$ define
\begin{equation}
 S_+:=P_+RP_+\big|_{\mathcal H_+}.
 \label{eq:v13-Splus}
\end{equation}
On the minimal stratum $S_+$ is positive and invertible, and the V11 graph
coordinate is exactly
\begin{equation}
 \boxed{
 A=QRP_+S_+^{-1}:\mathcal H_+\longrightarrow K^\perp.}
 \label{eq:v13-A-positive}
\end{equation}
If $P$ is exponentially local and the reduced chiral singular value has a
uniform positive floor, then $A$ and every fixed source derivative are
exponentially local at some uniform positive rate.
\end{proposition}

\begin{proof}
In the V11 decomposition
$W=K\oplus\operatorname{Graph}(A)$, the projection $R$ onto the graph has the
block matrix
\[
 R=
 \begin{pmatrix}I\\A\end{pmatrix}
 (I+A^*A)^{-1}
 \begin{pmatrix}I&A^*\end{pmatrix}
\]

after ordering $\mathcal H_+\oplus K^\perp$.  Hence
$S_+=(I+A^*A)^{-1}$ and
$QRP_+=A(I+A^*A)^{-1}$, proving
\eqref{eq:v13-A-positive}.  The V11 reduced singular floor is equivalent to a
positive lower bound on $S_+$.  Theorem~\ref{thm:v13-zero-projector-local}
makes $\Pi_K$ local, so $R$ and $S_+$ are local.  Apply
Theorem~\ref{thm:v13-resolvent-locality} to $S_+^{-1}$ and use the weighted
algebra product bound.  Source derivatives follow by the inverse rule.
\end{proof}

\begin{proposition}[Local graph coordinate for negative index]
\label{prop:v13-negative-local-graph}
Assume $k<0$.  Let $\Pi_C$ be the cokernel projection in $\mathcal H_+$ and set
\begin{equation}
 E:=P_+-\Pi_C,
 \qquad
 \widehat S:=\Pi_C+EPE
 \quad\text{on }\mathcal H_+.
 \label{eq:v13-Snegative}
\end{equation}
Then $\widehat S$ is positive and invertible on the protected minimal stratum,
and the V11 graph coordinate is
\begin{equation}
 \boxed{
 A=QPE\widehat S^{-1}E:C^\perp\longrightarrow\mathcal H_-.}
 \label{eq:v13-A-negative}
\end{equation}
It is exponentially local with all prescribed fixed source derivatives.
\end{proposition}

\begin{proof}
On $C^\perp$, the plus--plus block of the graph projector is
$(I+A^*A)^{-1}$; on $C$ the added term $\Pi_C$ is the identity.  Thus
$\widehat S$ has a uniform positive floor.  The block identity
$QPE=A(I+A^*A)^{-1}$ on $C^\perp$ gives
\eqref{eq:v13-A-negative}.  Locality follows exactly as in
Proposition~\ref{prop:v13-positive-local-graph}.
\end{proof}

\subsection{Reduction of the native phase problem to the controlled zero-mode module}
\label{sec:v13-reduction}

\begin{theorem}[Controlled zero-mode reduction of the fixed-index phase problem]
\label{thm:v13-controlled-reduction}
Fix the physical principal Dirac--Yukawa jet, the Standard-Model representation
grading, one positive IR/reference scale, a fixed index $k$, and one protected
minimal reduced-gap component.  Then:
\begin{enumerate}[label=\textup{(R\arabic*)},leftmargin=3em]
\item a native local-operator homotopy induces a controlled equivalence of the
kernel packets for $k>0$ or cokernel packets for $k<0$;
\item after the controlled zero-mode packet has been transported, the remaining
V11 affine graph coordinate can be connected by a straight-line homotopy in
the local algebra, with a uniform reduced singular-value floor;
\item hence no additional nonzero-mode local invariant survives inside the
minimal fixed-index overlap geometry.
\end{enumerate}
The only remaining microscopic datum beyond the already-fixed principal/line
information is the controlled zero-mode module class.
\end{theorem}

\begin{proof}
Item (R1) is Proposition~\ref{prop:v13-native-implies-controlled}.  For (R2),
assume first $k>0$ and let $U_t$ be a local Kato transport of the zero-mode
projection.  Extend it by the identity on $\mathcal H_+$.  Transport the first
graph coordinate by
$A_0(t)=U_tA_0$; this remains a local map into $K_t^\perp$ and therefore defines
a minimal-stratum local projector path by V11.  At the endpoint the kernel is
$K_1$.  Keeping $K_1$ fixed, linearly interpolate
\[
 A_s=(1-s)U_1A_0+sA_1.
\]
The graph fibre is a vector space, so the path stays in the minimal stratum.
Its norm is uniformly bounded, and the explicit V11 formula
$2a^{-1}(I+A_s^*A_s)^{-1/2}$ gives a positive reduced kinetic floor.  On the
compact interacting source chart one shrinks once, as in V9--V11, to preserve
the full chiral/Yukawa reduced floor.  Propositions
\ref{prop:v13-positive-local-graph} and
\ref{prop:v13-negative-local-graph} give the required locality.  The negative
index case is identical with the Kato transport performed in $\mathcal H_+$ on
the cokernel bundle.

All zero-mode transports lie in the local compact/soft ideal and hence do not
alter the frozen physical principal quotient.  The endpoints have the same
principal jet, so the linear graph interpolation may be performed in the affine
subspace with that quotient fixed.  This proves (R2), and (R3) follows from the
affineness of the V11 graph fibre.
\end{proof}

The theorem is deliberately phrased in terms of a controlled zero-mode
transport.  Definition~\ref{def:v13-controlled-equivalence} supplies an endpoint
partial isometry but not automatically an unstabilized path of such projectors.
The next subsection gives an explicit stable construction which is sufficient
for the coefficientwise continuum.

\subsection{Stable controlled lifting with an inert massive spectator}
\label{sec:v13-stable-lift}

Adjoin one auxiliary copy of the same finite chiral carrier, with no physical
sources, no gauge/Yukawa couplings and fixed Wilson sign
\begin{equation}
 \varepsilon_{\rm sp}=\gamma_5.
 \label{eq:v13-spectator-sign}
\end{equation}
Then
\begin{equation}
 D_{{\rm ov},{\rm sp}}
 =a^{-1}(I+\gamma_5^2)=2a^{-1}I,
 \label{eq:v13-spectator-overlap}
\end{equation}
so the spectator is massive, has index zero and has no zero mode.

\begin{lemma}[Exact decoupling of the inert spectator]
\label{lem:v13-spectator-decouple}
At either endpoint, the spectator of
\eqref{eq:v13-spectator-sign} contributes only a source-independent positive
finite factor to the fermion determinant.  It has no source vertices,
BV/BRST descendants, anomaly coefficients, stress/composite writers or line
holonomy.  Consequently it cancels exactly from
\begin{equation}
 W[J]=\log Z[J]-\log Z[0]
 \label{eq:v13-W-normalized}
\end{equation}
and from every retained normalized connected/source/master coefficient.
Adding or deleting finitely many such decoupled spectators at the endpoints
therefore does not change the coefficientwise continuum proved in V7.
\end{lemma}

\begin{proof}
The spectator action is a fixed nonzero quadratic form independent of every
physical source and background writer.  Its Grassmann integral is its constant
nonzero determinant.  The same constant multiplies $Z[J]$ and $Z[0]$ and hence
cancels in \eqref{eq:v13-W-normalized}.  With no source or gauge couplings it
has no differentiated vertices, no BV bracket entries and no anomaly trace.
Its determinant section is the fixed positive scalar section, so its line
transport is trivial.
\end{proof}

Let $J:\mathcal H_-\to\mathcal H_-^{\rm sp}$ be the fixed onsite copy
identification.  Suppose now that
$V$ is a controlled partial isometry from the initial zero-mode projector
$\Pi_0$ to the final one $\Pi_1$.  Put
\begin{equation}
 \widetilde V:=JV:
 \mathcal H_-\longrightarrow\mathcal H_-^{\rm sp}.
 \label{eq:v13-Vtilde}
\end{equation}

\begin{lemma}[Controlled rotation into the spectator copy]
\label{lem:v13-spectator-rotation}
For $0\le\theta\le\pi/2$ the block operator on
$\mathcal H_-\oplus\mathcal H_-^{\rm sp}$
\begin{equation}
 \Pi_\theta=
 \begin{pmatrix}
 \cos^2\theta\,\Pi_0 &
 \cos\theta\sin\theta\,\widetilde V^*\\
 \cos\theta\sin\theta\,\widetilde V &
 \sin^2\theta\,J\Pi_1J^*
 \end{pmatrix}
 \label{eq:v13-spectator-rotation}
\end{equation}
is an orthogonal rank-$|k|$ projection, exponentially local uniformly in
$\theta$, with
\begin{equation}
 \Pi_{0}=\Pi_0\oplus0,
 \qquad
 \Pi_{\pi/2}=0\oplus J\Pi_1J^*.
 \label{eq:v13-rotation-endpoints}
\end{equation}
A second identical onsite rotation using $J^*$ moves
$0\oplus J\Pi_1J^*$ to $\Pi_1\oplus0$.  Hence controlled-equivalent zero-mode
packets are connected after one inert-copy stabilization by an explicit local
projection path of the same rank.
\end{lemma}

\begin{proof}
Use
$\widetilde V^*\widetilde V=\Pi_0$ and
$\widetilde V\widetilde V^*=J\Pi_1J^*$.  Direct block multiplication shows
$\Pi_\theta^2=\Pi_\theta$ and self-adjointness is manifest.  Locality follows
from Definition~\ref{def:v13-controlled-equivalence}, the onsite nature of $J$,
and Lemma~\ref{lem:v13-Schur-algebra}.  The endpoint formulas are immediate.
The second rotation is the same construction with the roles of the two copies
reversed and partial isometry $J^*J\Pi_1=\Pi_1$ on the relevant range.
\end{proof}

\begin{theorem}[Stable native lifting from controlled zero-mode equivalence]
\label{thm:v13-stable-native-lift}
Let two V11 protected minimal fixed-index regulators $A,B$ have the same
physical principal jet, representation grading and global complex line datum
(and the same real orientation datum when a Pfaffian sector is present).
Assume their kernel packets for $k>0$, or cokernel packets for $k<0$, are
controlled-equivalent.  After adjoining one inert massive spectator copy to
both endpoints, there is a finite concatenation of uniformly exponentially
local native homotopies joining the stabilized regulators such that:
\begin{enumerate}[label=\textup{(S\arabic*)},leftmargin=3em]
\item the index and minimal zero-mode rank remain fixed;
\item the physical principal Dirac--Yukawa jet remains fixed;
\item the reduced kinetic singular value has a positive uniform floor, and the
full chiral/Yukawa reduced floor remains positive after one common finite
shrinking of the protected source chart;
\item all required fixed source jets remain in one positive locality class;
\item the zero-mode line/saturation packet is transported continuously.
\end{enumerate}
Thus every segment is V11-admissible.
\end{theorem}

\begin{proof}
Treat $k>0$; the cokernel proof is the chiral dual.  Add the inert balanced copy
of \eqref{eq:v13-spectator-sign}.  The endpoint kernel projectors become
$\Pi_0\oplus0$ and $\Pi_1\oplus0$ in the enlarged negative-chirality carrier.
Lemma~\ref{lem:v13-spectator-rotation} supplies a controlled local path between
them, of constant rank $k$.

Use Theorem~\ref{thm:v13-local-Kato} on this explicit local projector path to
obtain a local unitary transport of the negative-chirality complement.  Extend
it by the identity on the enlarged positive-chirality carrier and transport the
initial local graph coordinate supplied by
Proposition~\ref{prop:v13-positive-local-graph}.  This gives the first stage of
a minimal-stratum projector homotopy.  Once the target kernel projection is
reached, linearly interpolate the transported graph coordinate to the target
graph coordinate (with zero graph on the spectator block).  The V11 affine
geometry keeps the index minimal throughout and gives the explicit reduced
kinetic floor
\[
 {2\over a\sqrt{1+\|A_t\|^2}}.
\]
All operators used are finite products, sums, inverses across a positive gap,
or Dyson evolutions inside the weighted local algebra; hence they retain some
uniform positive exponential rate by the preceding theorems.

The zero-mode rotation is a compact/soft local deformation and the common
principal quotient is frozen.  The endpoint graph coordinates have the same
physical first jet, so the affine interpolation is taken inside that fixed
principal-jet fibre.  The full Yukawa/source pivot is an open condition on the
compact path, so the usual one-time chart shrinking of V9--V12 supplies a
common positive reduced floor.  Kato exterior powers transport the zero-mode
line and saturation writer.
\end{proof}

\begin{theorem}[Stable controlled regulator universality]
\label{thm:v13-stable-universality}
Under the hypotheses of Theorem~\ref{thm:v13-stable-native-lift}, the original
unstabilized regulators have the same matched coefficientwise continuum.  More
precisely, for every prescribed finite coefficient packet $\mathfrak q$ there
is an invertible finite local map
\begin{equation}
 F_{\infty,\mathfrak q,k}^{B\leftarrow A}
 \label{eq:v13-stable-map}
\end{equation}
such that
\begin{equation}
 \boxed{
 \widehat{\mathfrak C}^{B,(k),\rm sat}_{\infty,\lambda}
 =F_{\infty,\mathfrak q,k}^{B\leftarrow A}
  \widehat{\mathfrak C}^{A,(k),\rm sat}_{\infty,\lambda}.}
 \label{eq:v13-stable-universality}
\end{equation}
The equality includes the complete normalized connected/source packet,
operator mixing, stress/contact terms, ghosts/antifields, BV/BRST/master
relations and determinant/Pfaffian line insertions, and the maps can be chosen
projectively in perturbative order.
\end{theorem}

\begin{proof}
Apply V11 fixed-index universality to the stabilized local homotopy of
Theorem~\ref{thm:v13-stable-native-lift}.  At both endpoints remove the inert
spectator using Lemma~\ref{lem:v13-spectator-decouple}.  Since this removal is
an exact equality of normalized coefficient systems rather than a limiting
argument, the matching map descends to the original regulators.  The V7
chronological construction and the spectator's trivial source packet preserve
strict perturbative-order projectivity.
\end{proof}

\subsection{A positive criterion: uniformly tight zero modes}
\label{sec:v13-tight}

The controlled module condition is stronger than ordinary bundle isomorphism,
but it is automatic for many localized defect sectors.

\begin{definition}[Uniformly tight zero-mode packet]
\label{def:v13-tight-packet}
A finite-rank zero-mode projector family $\Pi$ is \emph{uniformly tight} if
there is a family of multiplication cutoffs $\chi_R$ supported in sets of
uniformly bounded diameter (independent of cutoff and volume) such that
\begin{equation}
 \|(I-\chi_R)\Pi\|+\|\Pi(I-\chi_R)\|
 \xrightarrow[R\to\infty]{}0
 \label{eq:v13-tightness}
\end{equation}
uniformly on the compact background panel, and the same holds for the retained
source derivatives.
\end{definition}

\begin{theorem}[Bundle isomorphism plus uniform tightness gives a controlled isomorphism]
\label{thm:v13-tight-controlled}
Let $\Pi_0,\Pi_1$ be uniformly tight zero-mode packets of the same rank over a
compact background panel.  Assume their ordinary complex zero-mode bundles are
isomorphic by a smooth unitary bundle map.  Then
\begin{equation}
 \boxed{\Pi_0\sim_{\rm ctl}\Pi_1.}
 \label{eq:v13-tight-ctl}
\end{equation}
Consequently Theorem~\ref{thm:v13-stable-universality} applies.
\end{theorem}

\begin{proof}
Let $W$ be a smooth unitary bundle isomorphism
$\operatorname{Ran}\Pi_0\to\operatorname{Ran}\Pi_1$, extended by zero on the
orthogonal complement.  Because the panel is compact and the rank is finite,
$W$ is uniformly bounded with all retained source derivatives.  Choose a common
localizing cutoff $\chi_R$ from Definition~\ref{def:v13-tight-packet} so that
\begin{equation}
 T_R:=\Pi_1\chi_RW\chi_R\Pi_0
 \label{eq:v13-TR}
\end{equation}
differs from $W$ by less than $1/4$ on the initial zero-mode range, uniformly on
the panel.  The support of $\chi_R$ has uniformly bounded diameter, so $T_R$ is
a controlled local finite-rank kernel.  Its restriction to
$\operatorname{Ran}\Pi_0$ is invertible onto
$\operatorname{Ran}\Pi_1$.  Put
\begin{equation}
 V=T_R\bigl(T_R^*T_R+(I-\Pi_0)\bigr)^{-1/2}\Pi_0.
 \label{eq:v13-polar-V}
\end{equation}
The operator in parentheses has a uniform positive floor.  Theorem
\ref{thm:v13-resolvent-locality} gives a local inverse square root, so $V$ is
controlled.  Polar decomposition gives
$V^*V=\Pi_0$ and $VV^*=\Pi_1$.  Source derivatives follow from the same
functional calculus.
\end{proof}

\begin{corollary}[Localized fixed-index sectors have no extra controlled obstruction]
\label{cor:v13-localized-sector}
On a uniformly tight zero-mode sector, the V12 ordinary kernel/cokernel bundle
class together with the already-declared complex/Pfaffian line data is enough
to imply stable native regulator universality.
\end{corollary}

\subsection{Controlled phase hierarchy and exact boundary of V13}
\label{sec:v13-boundary}

The preceding results sharpen the V12 three-level hierarchy to
\begin{equation}
 \boxed{
 \text{line class}
 \ \leftarrow\ 
 \text{ordinary zero-mode bundle class}
 \ \leftarrow\ 
 \text{controlled zero-mode module class}
 \ \leftarrow\ 
 \text{unstabilized native component}.}
 \label{eq:v13-four-level}
\end{equation}
The first two arrows can lose information by V12.  Proposition
\ref{prop:v13-macro-obstruction} shows that the second-to-third arrow can also
lose information.  Theorem~\ref{thm:v13-stable-native-lift} shows that the
third level is already sufficient after one physically inert stabilization.
What remains is cancellation of that stabilization in the microscopic path,
not in the normalized continuum where it has already disappeared exactly.

\begin{theorem}[V13 phase-reduction theorem]
\label{thm:v13-phase-reduction}
For the protected minimal fixed-index overlap problem with a positive reduced
singular-value floor:
\begin{enumerate}[label=\textup{(C\arabic*)},leftmargin=3em]
\item zero-mode Riesz projectors and Kato transports are automatically
quasilocal; arbitrary nonlocal finite-rank projectors are not the correct
spectral variables;
\item the nonzero-mode minimal-stratum fibre contributes no independent local
phase invariant;
\item a necessary invariant of an unstabilized native component is the
controlled zero-mode module class;
\item equality of that class is sufficient for native lifting after one inert
massive stabilization and hence is sufficient for equality of the original
normalized coefficientwise continua;
\item in uniformly tight zero-mode sectors, ordinary bundle isomorphism already
implies controlled equivalence.
\end{enumerate}
\end{theorem}

\begin{proof}
Items (C1)--(C5) are respectively Theorems
\ref{thm:v13-zero-projector-local},
\ref{thm:v13-controlled-reduction},
Proposition~\ref{prop:v13-native-implies-controlled},
Theorems~\ref{thm:v13-stable-native-lift}--\ref{thm:v13-stable-universality},
and Theorem~\ref{thm:v13-tight-controlled}.
\end{proof}

\subsection{V13 anti-circularity audit}
\label{sec:v13-audit}

\begin{enumerate}[leftmargin=2.4em]
\item \textbf{V12's counterexample was not discarded.}
V13 identifies the extra hypothesis it lacked: a local parent operator with a
uniform spectral separation of the zero eigenspace.  Riesz functional calculus
then forces locality of the actual V11 zero-mode projector.
\item \textbf{Locality of each endpoint projector was not promoted to locality
of an endpoint identification.}
Proposition~\ref{prop:v13-macro-obstruction} exhibits two perfectly local,
uniformly gapped rank-one zero modes whose unique intertwiner has locality cost
$e^{\beta L/2}$.
\item \textbf{The reduced graph fibre is not treated as a new topological
sector.}
Its coordinate is reconstructed by the local formulas
\eqref{eq:v13-A-positive} and \eqref{eq:v13-A-negative} and interpolated in an
affine space.
\item \textbf{The stable lift does not add physical massless degrees of
freedom.}
The auxiliary copy has overlap operator $2a^{-1}I$, no zero modes and no
sources.  Its endpoint contribution cancels exactly from normalized
coefficients.
\item \textbf{Stable continuum equivalence is distinguished from an
unstabilized microscopic path.}
V13 proves the former.  Removing the spectator from the path itself is the next
cancellation problem.
\item \textbf{The complex and real line data remain separate.}
All V13 controlled constructions carry the V12 complex saturation writer and,
when present, require the same real Pfaffian orientation component.
\end{enumerate}

\section{V13 status ledger}
\label{sec:v13-status}

\subsection*{Proved in V13}
\begin{enumerate}[leftmargin=2.2em]
\item Weighted exponentially local kernels form a Banach $*$-algebra, and a
uniform spectral gap gives a cutoff-uniform exponentially local resolvent after
a finite loss of exponent.  Riesz projectors, signs and fixed source derivatives
inherit the same property.
\item The V11 kernel/cokernel Riesz projectors and reduced Moore--Penrose inverse
are therefore uniformly exponentially local on every protected minimal
fixed-index chart with a positive reduced singular-value floor.
\item Kato transport along an already-native fixed-index path is itself
exponentially local, with all retained source derivatives.  Hence native paths
preserve a controlled zero-mode module class, stronger than the ordinary bundle
class.
\item Two macroscopically separated onsite zero modes give an explicit
counterexample showing that ordinary rank/bundle/determinant data plus a uniform
reduced gap do not imply controlled equivalence: the unique partial isometry
has locality norm $e^{\beta L/2}$.
\item The V11 affine graph coordinate is reconstructed from local projectors
and a gapped inverse by the exact formulas
\eqref{eq:v13-A-positive}--\eqref{eq:v13-A-negative}; it carries no independent
nonzero-mode local phase invariant.
\item Controlled-equivalent zero-mode packets admit an explicit stable native
lift after adjoining one inert massive spectator copy.  The path preserves
index, minimal nullity, the physical principal jet, a positive reduced gap,
source locality and the zero-mode line packet.
\item The inert spectator cancels exactly from normalized connected/source/BV
coefficients.  Therefore controlled equivalence implies equality of the
original unstabilized coefficientwise continua, despite the temporary
microscopic stabilization used to construct the path.
\item In uniformly tight localized zero-mode sectors, an ordinary smooth bundle
isomorphism can be localized and polar-corrected to a controlled partial
isometry.  Such sectors therefore have no extra controlled obstruction beyond
the V12 bundle/line data.
\end{enumerate}

\subsection*{Inherited}
\begin{enumerate}[leftmargin=2.2em]
\item V1--V7: source-complete fixed-order continuum, anomaly-free common BV
restoration, arbitrary-fixed-loop forests, ultraviolet Cauchy convergence,
scheme/reference covariance and exact loop-order projectivity.
\item V8--V10: regulator-marked two-scale comparison and regulator-class
universality on the protected index-zero single-cone sectors.
\item V11: minimal fixed-index affine geometry, reduced pseudoinverse calculus,
zero-mode Schur factorization, saturation writers and local fixed-index
universality.
\item V12: global zero-mode bundle/line patching, faithful-spin complex line
trivialization, global saturation, the stable finite-Grassmannian bundle
classification, and global universality within a protected native component.
\end{enumerate}

\subsection*{Conditional}
\begin{enumerate}[leftmargin=2.2em]
\item The local resolvent/Riesz theorem uses a cutoff-uniform exponential
locality bound and a positive reduced singular-value floor.  If the floor
collapses with volume, the V12 nonlocal projector phenomenon can reappear.
\item Stable native lifting requires a controlled partial isometry of the
kernel/cokernel packets.  Ordinary bundle isomorphism supplies one automatically
only under additional locality hypotheses such as the uniform-tightness theorem.
\item The spectator argument proves equality of normalized continuum
coefficients, not an unstabilized microscopic homotopy in the original carrier.
\item The V12 real Pfaffian mod-two orientation requirement remains in force.
\item All statements remain coefficientwise at fixed finite
$(L,\Delta,r,N)$ and positive IR/reference scale.
\end{enumerate}

\subsection*{Open beyond V13}
\begin{enumerate}[leftmargin=2.2em]
\item Prove \emph{controlled cancellation}: determine whether a controlled
zero-mode module isomorphism always yields an unstabilized native homotopy in
the exact Standard-Model carrier, or exhibit a genuine obstruction to removing
the inert spectator from the microscopic path.
\item Classify controlled zero-mode modules on physically relevant
noncontractible gauge-background components.  The macroscopic-transport example
shows that the answer can contain coarse/locality information absent from
ordinary $K$-theory.
\item Compute the V12 real Pfaffian mod-two spectral-flow character on any
independently physical Majorana sector.
\item Analyze exits from the minimal fixed-index stratum, including accidental
zero-mode-pair creation and true index change.
\item The independent stronger problems remain: positive-IR removal, analytic
finite-coupling RG trajectories, perturbative summability and nonperturbative
completion.
\end{enumerate}

\subsection*{Exact next load-bearing theorem}
\begin{center}
\fbox{\parbox{0.92\textwidth}{
\textbf{Controlled cancellation / unstabilized native lifting.}
Fix the Standard-Model representation, physical principal Dirac--Yukawa jet,
index $k$, reduced-gap component and global zero-mode line data.  Given a
controlled partial isometry between the two local kernel/cokernel packets,
prove that it extends inside the same native gauge/spin-covariant exponentially
local carrier to a path of minimal fixed-index Wilson projectors without the
auxiliary massive copy, or identify the exact obstruction (for example a
controlled complement-unitary/$K_1$ or coarse-transport class).  A positive
cancellation theorem, combined with V13 and
Theorem~\ref{thm:v12-global-index-universality}, would upgrade stable controlled
universality to a complete unstabilized global fixed-index regulator-class
theorem.}}
\end{center}

\section*{Append-only continuation marker: V13}
\addcontentsline{toc}{section}{Append-only continuation marker: V13}
\textbf{End of V13.}  The complete V1--V12 body above is retained verbatim.
V13 goes upstream from the abstract-projector locality obstruction and proves
that the actual V11 zero-mode Riesz projectors are uniformly exponentially
local whenever their native parent operator has the declared reduced spectral
floor.  It identifies the remaining microscopic datum as a controlled local
zero-mode module class, shows by a macroscopically separated onsite example
that this datum is finer than ordinary bundle/determinant information, and
proves that the V11 reduced graph coordinate carries no further local phase
invariant.  Controlled-equivalent zero-mode packets admit an explicit stable
native lift after one inert massive spectator copy; the spectator cancels
exactly from all normalized coefficientwise observables, yielding stable
global regulator universality.  The next gate is controlled cancellation: can
the spectator be removed from the microscopic path itself?


\clearpage
\section{V14 continuation: controlled cancellation, complement defects, and unstabilized native lifting}
\label{sec:v14-continuation}

V13 reduced the protected fixed-index phase problem to a controlled local
zero-mode module.  It also proved that controlled-equivalent packets become
connected after adjoining one inert massive copy.  The remaining question is
therefore a genuine \emph{cancellation} problem: can the auxiliary copy be
removed from the microscopic path itself?

The strongest possible answer is false in an arbitrary local coefficient
algebra.  Murray--von Neumann equivalence of two zero-mode projections gives a
controlled partial isometry between their ranges, but an unstabilized unitary
extension additionally requires the two complementary modules to cancel.  Even
when that complement cancellation holds, the extending unitary must lie in the
identity component of the controlled unitary group in order to generate a
projection homotopy.  These are logically distinct requirements.

The present continuation makes that distinction exact.  It proves:
\begin{enumerate}[label=\textup{(V14.\arabic*)},leftmargin=3.8em]
\item a necessary-and-sufficient complement criterion for extending a
controlled partial isometry to a controlled unitary on the original carrier;
\item a necessary-and-sufficient identity-component criterion for an
unstabilized controlled projection homotopy;
\item a propagation-zero finite-dimensional counterexample showing that
controlled equivalence alone does not imply cancellation before stabilization;
\item a quantitative small-angle cancellation theorem and an explicit
``internal buffer'' theorem which remove the spectator whenever the original
carrier already contains one suitable local massive reserve;
\item an unstabilized fixed-index native lifting/universality theorem whenever
the zero-mode projections lie in the same controlled $U_0$ orbit;
\item a stable-range corollary showing that on fixed finite-dimensional
parameter panels the purely finite-dimensional complement-bundle obstruction
disappears at sufficiently large ambient rank, leaving a genuinely controlled
locality/unitary-component problem as the remaining gate.
\end{enumerate}

Throughout V14 let $\mathfrak A_\beta$ denote one of the V13 unital Banach
$*$-algebras of exponentially local, gauge/spin-covariant kernel families with
the prescribed finite source derivatives, at some fixed exponent
$0<\beta<\alpha$.  All projections, partial isometries and unitaries below are
required to obey the same representation/source constraints.  Write
\begin{equation}
 \mathcal U(\mathfrak A_\beta)
 :=\{u\in\mathfrak A_\beta:u^*u=uu^*=I\},
 \qquad
 \mathcal U_0(\mathfrak A_\beta)
 \label{eq:v14-U0}
\end{equation}
for its unitary group and identity component.

\subsection{Exact complement cancellation for a controlled partial isometry}
\label{sec:v14-complement-cancellation}

Let $p,q\in\mathfrak A_\beta$ be projections and let
\begin{equation}
 v^*v=p,
 \qquad
 vv^*=q
 \label{eq:v14-v}
\end{equation}
be a controlled partial isometry, as in V13.

\begin{theorem}[Controlled unitary extension if and only if the complements cancel]
\label{thm:v14-unitary-extension}
The following are equivalent.
\begin{enumerate}[label=\textup{(E\arabic*)},leftmargin=3em]
\item There exists $u\in\mathcal U(\mathfrak A_\beta)$ such that
\begin{equation}
 up=v.
 \label{eq:v14-up-v}
\end{equation}
\item There exists a controlled partial isometry $w\in\mathfrak A_\beta$ with
\begin{equation}
 \boxed{
 w^*w=I-p,
 \qquad
 ww^*=I-q.}
 \label{eq:v14-complement-MvN}
\end{equation}
\end{enumerate}
When \textup{(E2)} holds,
\begin{equation}
 \boxed{u=v+w}
 \label{eq:v14-u-vw}
\end{equation}
is a unitary extension.  Conversely every extension has
$w=u(I-p)$.
\end{theorem}

\begin{proof}
Assume \textup{(E2)}.  Since $v=qvp$ and
$w=(I-q)w(I-p)$, the initial and final projections of $v$ and $w$ are
orthogonal.  Hence
\[
 v^*w=w^*v=vw^*=wv^*=0.
\]
Therefore
\[
 (v+w)^*(v+w)=p+(I-p)=I,
 \qquad
 (v+w)(v+w)^*=q+(I-q)=I.
\]
Also $(v+w)p=v$.  Thus \eqref{eq:v14-u-vw} is the desired unitary.

Conversely suppose $u$ is unitary and $up=v$.  Then
$upu^*=vv^*=q$.  Put $w=u(I-p)$.  One has
\[
 w^*w=(I-p)u^*u(I-p)=I-p
\]
and
\[
 ww^*=u(I-p)u^*=I-upu^*=I-q.
\]
Thus \textup{(E2)} holds.
\end{proof}

This theorem isolates the first unstable obstruction exactly.  The endpoint
partial isometry of V13 proves $p\sim q$, but unstabilized extension asks in
addition for
\begin{equation}
 \boxed{I-p\sim I-q}
 \label{eq:v14-complement-cancel}
\end{equation}
in the \emph{same controlled algebra}.  Equality in stabilized $K_0$ does not,
by itself, imply this unstabilized Murray--von Neumann cancellation.

\begin{definition}[Controlled complement defect]
\label{def:v14-complement-defect}
For controlled-equivalent $p,q$ define the \emph{unstable complement defect}
to vanish when the complements are Murray--von Neumann equivalent
\emph{inside the original controlled carrier}:
\begin{equation}
 \boxed{
 \mathfrak c_\beta(p,q)=0
 \quad\Longleftrightarrow\quad
 \exists\,w\in\mathfrak A_\beta:\ 
 w^*w=I-p,\ ww^*=I-q.}
 \label{eq:v14-complement-defect}
\end{equation}
No matrix amplification or added free summand is allowed in this definition.
Thus $\mathfrak c_\beta$ is deliberately an unstable relation, stronger than
equality of the corresponding group-completed $K_0$ classes.
\end{definition}

By Theorem~\ref{thm:v14-unitary-extension},
$\mathfrak c_\beta(p,q)=0$ is exactly the existence of a controlled unitary
extension of a chosen controlled equivalence after its complement has been
filled in inside the same carrier.

\subsection{The second obstruction: the controlled unitary component}
\label{sec:v14-unitary-component}

A unitary extension does not yet give a path from the identity.  For a
\emph{fixed} controlled partial isometry $v$, the residual component is also
exact.

Suppose \eqref{eq:v14-complement-MvN} holds and fix one complement partial
isometry $w_0$.  Put
\begin{equation}
 u_0:=v+w_0\in\mathcal U(\mathfrak A_\beta).
 \label{eq:v14-u0}
\end{equation}
If $w_1$ is another complement equivalence, then
\begin{equation}
 z:=w_0^*w_1
 \label{eq:v14-zcorner}
\end{equation}
is a unitary in the corner
$(I-p)\mathfrak A_\beta(I-p)$ and
\begin{equation}
 v+w_1=u_0(p+z).
 \label{eq:v14-extension-difference}
\end{equation}

\begin{definition}[Relative controlled unitary-component class]
\label{def:v14-kappa}
When $\mathfrak c_\beta(p,q)=0$, define
\begin{equation}
 \kappa_\beta(v)
 :=[u_0]\,
 \operatorname{im}\!\left(
 \pi_0\mathcal U((I-p)\mathfrak A_\beta(I-p))
 \longrightarrow
 \pi_0\mathcal U(\mathfrak A_\beta)
 \right)
 \label{eq:v14-kappa}
\end{equation}
as a right coset in the component group.  Equation
\eqref{eq:v14-extension-difference} makes this coset independent of the chosen
$w_0$.
\end{definition}

\begin{proposition}[Exact component criterion for an extension of a fixed $v$]
\label{prop:v14-kappa-criterion}
Under $\mathfrak c_\beta(p,q)=0$, there exists a controlled unitary extension
$u$ of the fixed $v$ with
\begin{equation}
 u\in\mathcal U_0(\mathfrak A_\beta)
 \label{eq:v14-extension-U0}
\end{equation}
if and only if the coset $\kappa_\beta(v)$ contains the identity component.
\end{proposition}

\begin{proof}
Every extension has the form $u_0(p+z)$ with $z$ a unitary in the complement
corner, by \eqref{eq:v14-extension-difference}.  Hence one can move $u_0$ into
the identity component while keeping its restriction to $p$ equal to $v$
exactly when the component class of $u_0$ can be cancelled by the image of a
complement-corner unitary.  This is precisely the stated coset condition.
\end{proof}

If the endpoint zero-mode identification is not fixed a priori, one may also
change $v$ by a unitary in the $p$ corner.  The corresponding obstruction is a
double-coset version of \eqref{eq:v14-kappa}; V14 will not need separate
notation for it because the projection-homotopy criterion below is cleaner.

\subsection{Projection paths are exactly controlled $U_0$ orbits}
\label{sec:v14-U0-orbits}

We first record the local direct rotation between two nearby projections.

\begin{lemma}[Controlled direct rotation for a small projection angle]
\label{lem:v14-direct-rotation}
Let $p,q\in\mathfrak A_\beta$ be projections with
\begin{equation}
 \delta:=\|p-q\|_{\rm op}<1.
 \label{eq:v14-small-angle}
\end{equation}
Put
\begin{equation}
 X:=qp+(I-q)(I-p).
 \label{eq:v14-X}
\end{equation}
Then
\begin{align}
 X-I&=(q-p)(2p-I),
 \label{eq:v14-X-near-I}\\
 X^*X&=I-(p-q)^2,
 \label{eq:v14-XstarX}
\end{align}
so $X$ is invertible and
\begin{equation}
 \boxed{
 u_{q,p}:=X(X^*X)^{-1/2}}
 \label{eq:v14-direct-unitary}
\end{equation}
is a unitary in $\mathfrak A_{\beta'}$ for every sufficiently small fixed
loss $0<\beta'<\beta$.  It satisfies
\begin{equation}
 u_{q,p}pu_{q,p}^*=q,
 \qquad
 u_{q,p}\in\mathcal U_0(\mathfrak A_{\beta'}).
 \label{eq:v14-direct-properties}
\end{equation}
All prescribed fixed source derivatives obey uniform bounds when
$\delta\le1-\eta$ for a fixed $\eta>0$.
\end{lemma}

\begin{proof}
The first identity is direct expansion.  For the second,
$Xp=qX$, so the cross terms in $X^*X$ vanish between $p$ and $I-p$; expanding
gives
$X^*X=I-p-q+pq+qp=I-(p-q)^2$.  Thus
$X^*X\succeq(1-\delta^2)I$.

The V13 weighted Banach-algebra/resolvent functional calculus gives
$(X^*X)^{-1/2}\in\mathfrak A_{\beta'}$ with all fixed source derivatives.
The polar factor \eqref{eq:v14-direct-unitary} is unitary.  Since
$Xp=qX$ and $X^*X$ commutes with $p$, one has $u_{q,p}p=qu_{q,p}$, proving the
projection identity.

Finally \eqref{eq:v14-X-near-I} gives $\|X-I\|<1$.  Hence
$X_s=I+s(X-I)$ is invertible for every $s\in[0,1]$.  Its polar factors
$u_s=X_s(X_s^*X_s)^{-1/2}$ form a continuous controlled unitary path from
$I$ to $u_{q,p}$.  The uniform source bounds follow from the fixed lower
spectral margin $1-\delta^2\ge\eta(2-\eta)$.
\end{proof}

\begin{theorem}[Exact controlled projection-homotopy criterion]
\label{thm:v14-projection-U0}
For two projections $p,q\in\mathfrak A_\beta$, the following are equivalent
(after one harmless fixed loss of exponential weight if needed for functional
calculus).
\begin{enumerate}[label=\textup{(H\arabic*)},leftmargin=3em]
\item There exists a norm-continuous path of controlled projections
$p_t\in\mathfrak A_{\beta'}$, $0\le t\le1$, with $p_0=p$ and $p_1=q$.
\item There exists a controlled unitary
\begin{equation}
 u\in\mathcal U_0(\mathfrak A_{\beta'})
 \label{eq:v14-U0-orbit}
\end{equation}
such that
\begin{equation}
 \boxed{q=upu^*.}
 \label{eq:v14-U0-conjugacy}
\end{equation}
\end{enumerate}
The same equivalence holds in the source-differentiable gauge/spin-equivariant
subalgebra.
\end{theorem}

\begin{proof}
If \textup{(H2)} holds, choose a continuous controlled unitary path
$u_t$ from $I$ to $u$ and set $p_t=u_tpu_t^*$.  This gives \textup{(H1)}.

Conversely suppose \textup{(H1)} holds.  Uniform continuity of $t\mapsto p_t$
on the compact interval gives a finite partition
$0=t_0<\cdots<t_M=1$ such that
\[
 \|p_{t_{j+1}}-p_{t_j}\|_{\rm op}<1/2.
\]
Apply Lemma~\ref{lem:v14-direct-rotation} to each adjacent pair.  It gives
$u_j\in\mathcal U_0(\mathfrak A_{\beta'})$ with
$u_jp_{t_j}u_j^*=p_{t_{j+1}}$.  The finite product
$u=u_{M-1}\cdots u_0$ lies in the identity component and satisfies
$q=upu^*$.  Fixed source and covariance constraints are preserved at each
step because the direct-rotation formula uses only algebra operations and
functional calculus inside the same constrained subalgebra.
\end{proof}

\begin{definition}[Unstable controlled cancellation obstruction]
\label{def:v14-unstable-obstruction}
For controlled-equivalent zero-mode projections $p,q$, define
\begin{equation}
 \boxed{
 \operatorname{Can}_\beta(p,q)=0
 \quad\Longleftrightarrow\quad
 q=upu^*\text{ for some }
 u\in\mathcal U_0(\mathfrak A_\beta).}
 \label{eq:v14-Can}
\end{equation}
Thus $\operatorname{Can}_\beta$ is the exact obstruction to an unstabilized
controlled zero-mode projection path.  It contains both complement cancellation
and unitary-component information.
\end{definition}

\subsection{Controlled equivalence alone does not cancel in general}
\label{sec:v14-counterexample}

The next counterexample is deliberately propagation zero.  It shows that no
proof of V14 cancellation can follow from V13's endpoint partial isometry alone;
one must use additional structure of the actual Standard-Model carrier or a
stable-range/local-buffer hypothesis.

\begin{theorem}[Propagation-zero unstable cancellation failure]
\label{cth:v14-unstable-cancellation}
There exist a compact parameter space $B$, a finite-dimensional onsite carrier,
and smooth propagation-zero projections $p,q$ such that:
\begin{enumerate}[label=\textup{(C\arabic*)},leftmargin=3em]
\item $p$ and $q$ have the same finite rank and are connected by a smooth
propagation-zero partial isometry $v$;
\item their range bundles are isomorphic and their determinant lines agree;
\item nevertheless no continuous onsite unitary family $u(b)$ satisfies
$upu^*=q$.
\end{enumerate}
Consequently controlled Murray--von Neumann equivalence does not imply
$\operatorname{Can}_\beta(p,q)=0$ in an arbitrary native coefficient algebra.
The same example embeds into the V11 minimal positive-index projector geometry
by taking zero graph coordinate.
\end{theorem}

\begin{proof}
Take $B=S^6$.  Complex rank-two bundles over $S^6$ are classified by the
clutching group
\[
 \pi_5(U(2))\cong\pi_5(S^3)\cong\mathbb Z_2.
\]
Let $F\to S^6$ be the nontrivial bundle represented by the nonzero element.
Under stabilization, the clutching class maps to
$\pi_5(U)\cong\mathbb Z$ by Bott periodicity.  Every homomorphism
$\mathbb Z_2\to\mathbb Z$ is zero, so $F$ is stably trivial.  Hence for some
$r\ge1$ there is a unitary bundle isomorphism
\begin{equation}
 F\oplus\underline{\mathbb C}^{\,r}
 \simeq
 \underline{\mathbb C}^{\,r+2}.
 \label{eq:v14-stably-trivial-F}
\end{equation}
Use such an isomorphism to regard the trivial rank-$r$ summand as a subbundle
$E_0$ of the trivial rank-$(r+2)$ bundle whose orthogonal complement is $F$.
Let $E_1$ be the standard trivial rank-$r$ coordinate subbundle, whose
orthogonal complement is the trivial rank-two bundle.  Let $p,q$ be the
orthogonal projections onto $E_0,E_1$.  Since both range bundles are trivial,
choose a smooth unitary bundle isomorphism
$v:E_0\to E_1$ and extend it by zero.  Then
$v^*v=p$ and $vv^*=q$.

If a continuous unitary family $u$ of the ambient trivial bundle satisfied
$upu^*=q$, it would restrict to a unitary isomorphism
\[
 E_0^\perp\simeq E_1^\perp,
\]
that is, $F\simeq\underline{\mathbb C}^2$, contradicting the nontrivial
clutching class.  All matrices act onsite in the physical lattice variable, so
the locality cost is exactly zero propagation.

For the V11 positive-index embedding, take
$\mathcal H_-\cong\mathbb C^{r+2}$, a fixed
$\mathcal H_+$ of the same ambient dimension, kernel plane $K_i=E_i$, and
zero affine graph coordinate $A_i=0$.  Then
$W_i=K_i\oplus\mathcal H_+$ lies in the minimal index-$r$ stratum, while a
minimal-stratum homotopy would in particular homotope the kernel projections,
which has just been excluded.
\end{proof}

\begin{firewall}[The counterexample is unstable by construction]
The bundle $F$ in Countertheorem~\ref{cth:v14-unstable-cancellation} is
stably trivial.  Adding enough trivial ambient reserve removes this particular
obstruction.  This is exactly why V13's one-copy stabilization can succeed even
when an unstabilized finite carrier fails.  On a fixed finite-dimensional
parameter panel, the V12 stable-range theorem similarly removes such purely
finite-Grassmannian cancellation defects once the complement rank is large
enough.  The counterexample therefore identifies the logical obstruction; it
does not claim that this specific $S^6$ defect survives in every large-cutoff
Standard-Model carrier.
\end{firewall}

\subsection{Two constructive regimes where cancellation is automatic}
\label{sec:v14-positive-regimes}

The exact obstruction theorem is useful because it reveals simple sufficient
conditions which are native and quantitative.

\begin{corollary}[Small-angle unstabilized cancellation]
\label{cor:v14-small-angle}
If two protected zero-mode projectors satisfy
$\|p-q\|_{\rm op}<1$, then
\begin{equation}
 \operatorname{Can}_{\beta'}(p,q)=0
 \label{eq:v14-small-angle-cancel}
\end{equation}
for every sufficiently small fixed loss $\beta'<\beta$.  The path and its fixed
source derivatives are given explicitly by the direct rotation of
Lemma~\ref{lem:v14-direct-rotation}.
\end{corollary}

\begin{proof}
The canonical unitary $u_{q,p}$ lies in
$\mathcal U_0(\mathfrak A_{\beta'})$ and conjugates $p$ to $q$.
Apply Definition~\ref{def:v14-unstable-obstruction}.
\end{proof}

The second criterion is the unstabilized analogue of the V13 spectator
rotation, but with the reserve already present in the original carrier.

\begin{definition}[Controlled internal buffer]
\label{def:v14-buffer}
A projection $r\in\mathfrak A_\beta$ is a \emph{controlled internal buffer}
for $p,q$ if
\begin{equation}
 rp=rq=0,
 \label{eq:v14-buffer-orthogonal}
\end{equation}
and there are controlled partial isometries $a,b$ such that
\begin{equation}
 a^*a=p,
 \quad aa^*=r,
 \qquad
 b^*b=q,
 \quad bb^*=r.
 \label{eq:v14-buffer-isometries}
\end{equation}
The buffer is required to carry the same representation/source grading as the
zero-mode packet but may be a gapped massive subpacket of the original
microscopic carrier.
\end{definition}

\begin{lemma}[Exact controlled rotation through an internal buffer]
\label{lem:v14-buffer-rotation}
Let $p,r$ be orthogonal projections and let
$a^*a=p$, $aa^*=r$.  Then
\begin{equation}
 U_a(\theta)
 :=I-p-r+\cos\theta\,(p+r)+\sin\theta\,(a-a^*),
 \qquad 0\le\theta\le\frac\pi2,
 \label{eq:v14-buffer-U}
\end{equation}
is a controlled unitary path with
\begin{equation}
 U_a(0)=I,
 \qquad
 U_a(\pi/2)pU_a(\pi/2)^*=r.
 \label{eq:v14-buffer-endpoint}
\end{equation}
\end{lemma}

\begin{proof}
On the orthogonal sum $p\mathcal H\oplus r\mathcal H$, set
$J=a-a^*$.  Since $a^2=(a^*)^2=0$ and
$a^*a=p$, $aa^*=r$,
\[
 J^*=-J,
 \qquad
 J^2=-(p+r).
\]
Hence the restriction of \eqref{eq:v14-buffer-U} to this sum is
$e^{\theta J}$, while the operator is the identity on the orthogonal
complement.  It is therefore unitary and continuous in the same controlled
algebra.  Moreover $Jp=a$, so at $\theta=\pi/2$ the $p$-range is mapped onto
the $r$-range.
\end{proof}

\begin{theorem}[Internal-buffer cancellation theorem]
\label{thm:v14-buffer-cancellation}
If $p,q$ admit a controlled internal buffer $r$, then
\begin{equation}
 \boxed{\operatorname{Can}_\beta(p,q)=0.}
 \label{eq:v14-buffer-cancel}
\end{equation}
More precisely, concatenate the path moving $p$ to $r$ by
$U_a(\theta)$ with the reverse of the path moving $q$ to $r$ by
$U_b(\theta)$.  This gives an explicit controlled projection path from $p$ to
$q$ entirely inside the original carrier.
\end{theorem}

\begin{proof}
Lemma~\ref{lem:v14-buffer-rotation} gives a controlled path
$p\rightsquigarrow r$.  Applying the same lemma to $b$ gives a path
$q\rightsquigarrow r$; reversing it gives $r\rightsquigarrow q$.  Concatenation
proves a controlled projection path, so
Theorem~\ref{thm:v14-projection-U0} gives
$\operatorname{Can}_\beta(p,q)=0$.
\end{proof}

\begin{corollary}[When the V13 spectator is already internal]
\label{cor:v14-internal-spectator}
Suppose the original Standard-Model carrier contains, on the entire protected
background/source panel, a source-smooth exponentially local massive subpacket
$r$ of rank $|k|$ which is orthogonal to both endpoint zero-mode packets and is
controlled-equivalent to each of them in the appropriate chiral sector.  Then
the V13 stable rotation can be carried out inside the original carrier; no
external spectator copy is required.
\end{corollary}

\subsection{Unstabilized native lifting from the controlled $U_0$ criterion}
\label{sec:v14-native-lift}

The V13 controlled-reduction theorem already proved that the nonzero-mode graph
coordinate carries no additional phase invariant.  We can therefore promote
the exact projection criterion to the full native regulator.

\begin{theorem}[Unstabilized fixed-index native lifting]
\label{thm:v14-unstable-native-lift}
Let $A,B$ be two V11 protected minimal fixed-index native regulators with the
same physical principal Dirac--Yukawa jet, Standard-Model representation,
positive IR/reference prescription and global complex/Pfaffian line datum.
Let $p_A,p_B$ denote their kernel projectors for $k>0$ or cokernel projectors
for $k<0$.  Assume that, in the corresponding equivariant source-differentiable
local algebra,
\begin{equation}
 \boxed{
 p_B=u p_Au^*
 \quad\text{for some}\quad
 u\in\mathcal U_0(\mathfrak A_\beta).}
 \label{eq:v14-native-U0-hyp}
\end{equation}
Then $A$ and $B$ are joined \emph{without stabilization} by a finite
concatenation of V11-admissible native fixed-index homotopies after the usual
one-time shrinking of the compact protected source chart.
\end{theorem}

\begin{proof}
Choose a controlled unitary path $u_t$ from $I$ to $u$ and put
$p_t=u_tp_Au_t^*$.  This is a controlled zero-mode projector path with all
prescribed source derivatives.  Apply the constructive part of
Theorem~\ref{thm:v13-controlled-reduction}: for $k>0$ transport the initial
local graph coordinate by $u_t$ in $\mathcal H_-$, thereby obtaining a
minimal-stratum projector path whose kernel is $p_t$; when the target kernel is
reached, linearly interpolate the transported affine graph coordinate to the
target graph coordinate.  For $k<0$ use the chiral-dual construction in
$\mathcal H_+$.  Propositions~\ref{prop:v13-positive-local-graph} and
\ref{prop:v13-negative-local-graph} keep every graph coordinate in the local
algebra and give the explicit positive reduced kinetic floor.

As in V13, the zero-mode unitary transport is a controlled compact/soft
deformation and the physical principal quotient is frozen.  The endpoint graph
coordinates have the same prescribed principal jet, so the final affine
interpolation is taken inside that fixed-jet fibre.  The compact path and the
open reduced/Yukawa pivot conditions permit one finite shrinking of the source
chart to keep the full reduced singular floor positive.  Kato exterior powers
carry the global line/saturation packet.  Every segment therefore satisfies
Definition~\ref{def:v11-index-admissible} without adding a new carrier copy.
\end{proof}

\begin{theorem}[Unstabilized global fixed-index regulator universality]
\label{thm:v14-unstable-universality}
Under the hypotheses of Theorem~\ref{thm:v14-unstable-native-lift}, for every
fixed coefficient packet $\mathfrak q$ there is an invertible finite local
continuum matching map
\begin{equation}
 F_{\infty,\mathfrak q,k}^{B\leftarrow A}
 \label{eq:v14-unstable-map}
\end{equation}
such that
\begin{equation}
 \boxed{
 \widehat{\mathfrak C}^{B,(k),\rm sat}_{\infty,\lambda}
 =F_{\infty,\mathfrak q,k}^{B\leftarrow A}
  \widehat{\mathfrak C}^{A,(k),\rm sat}_{\infty,\lambda}.}
 \label{eq:v14-unstable-universality}
\end{equation}
No spectator is introduced at any stage.  The equality contains the complete
normalized connected/source packet, operator mixing, stress/contact terms,
ghosts/antifields, BV/BRST/master relations and determinant/Pfaffian insertions,
and remains strictly projective in perturbative order.
\end{theorem}

\begin{proof}
Apply V11 fixed-index universality to the unstabilized native path of
Theorem~\ref{thm:v14-unstable-native-lift}.  The V12 global line transport
supplies the saturated complex comparison on the faithful branch, with the
same separate real-orientation hypothesis where applicable.  The V7
chronological construction gives exact projectivity in the terminal loop order.
\end{proof}

\begin{corollary}[Small-angle and internal-buffer universality]
\label{cor:v14-positive-universality}
The conclusion of Theorem~\ref{thm:v14-unstable-universality} holds, without
stabilization, whenever either:
\begin{enumerate}[label=\textup{(P\arabic*)},leftmargin=3em]
\item the two endpoint zero-mode projectors have operator-norm distance less
than one; or
\item they admit a controlled internal buffer in the sense of
Definition~\ref{def:v14-buffer}.
\end{enumerate}
\end{corollary}

\begin{proof}
Use Corollary~\ref{cor:v14-small-angle} or
Theorem~\ref{thm:v14-buffer-cancellation} to verify
\eqref{eq:v14-native-U0-hyp}, then apply
Theorem~\ref{thm:v14-unstable-universality}.
\end{proof}

\subsection{Stable range removes the purely finite-dimensional complement defect}
\label{sec:v14-stable-range}

V12 already proved that on a finite parameter panel and sufficiently large
ambient rank, ordinary bundle isomorphism gives a finite-Grassmannian projector
homotopy.  Combined with Theorem~\ref{thm:v14-projection-U0}, this has a useful
interpretation for V14.

\begin{proposition}[Onsite stable-range cancellation]
\label{prop:v14-onsite-stable-range}
Let $B$ be a finite CW complex of dimension $d$, and consider the propagation-zero
algebra $C(B,M_m)$ (or its smooth analogue on a compact manifold).  Let $p,q$
be rank-$r$ projections whose range bundles are isomorphic.  If
\begin{equation}
 d\le2(m-r),
 \label{eq:v14-onsite-stable-range}
\end{equation}
then
\begin{equation}
 q=upu^*
 \end{equation}
for some $u$ in the identity component of the onsite unitary group.
Consequently the purely finite-dimensional complement defect and unitary
component obstruction vanish in this stable range.
\end{proposition}

\begin{proof}
The V12 finite-Grassmannian stable-range theorem says that the two classifying
maps $B\to\operatorname{Gr}(r,\mathbb C^m)$ are homotopic when their vector
bundles are isomorphic under \eqref{eq:v14-onsite-stable-range}.  This is a
path of onsite projections.  Apply
Theorem~\ref{thm:v14-projection-U0} in the onsite algebra to obtain the desired
identity-component unitary.
\end{proof}

\begin{corollary}[What remains at large cutoff on a fixed finite panel]
\label{cor:v14-large-cutoff-boundary}
Fix a finite-dimensional background/source panel and a fixed zero-mode rank.
When the cutoff makes the ambient chiral complement sufficiently large for
\eqref{eq:v14-onsite-stable-range}, an obstruction to unstabilized native
lifting cannot be attributed merely to an unstable finite-dimensional bundle
complement.  Any remaining obstruction must use the genuinely controlled
spacetime-local algebra (or the independent real Pfaffian orientation data).
\end{corollary}

\subsection{The exact relationship between V13 stability and V14 cancellation}
\label{sec:v14-stable-vs-unstable}

The previous results identify precisely what the inert V13 copy accomplishes.

\begin{theorem}[Stable versus unstabilized controlled equivalence]
\label{thm:v14-stable-unstable}
For protected zero-mode projections $p,q$:
\begin{enumerate}[label=\textup{(S\arabic*)},leftmargin=3em]
\item an unstabilized native projection path implies controlled equivalence
$p\sim_{\rm ctl}q$;
\item controlled equivalence always gives a stabilized controlled projection
path by the V13 two-copy rotation;
\item removing the stabilization is equivalent to
$\operatorname{Can}_\beta(p,q)=0$;
\item the obstruction $\operatorname{Can}_\beta$ can be analyzed in two stages:
first the complement defect $\mathfrak c_\beta$, then, when it vanishes, the
relative controlled unitary-component class of
Definition~\ref{def:v14-kappa}.
\end{enumerate}
Thus V13 already proves continuum universality at the stable controlled level,
while V14 identifies the exact extra datum required for a microscopic path in
the original carrier.
\end{theorem}

\begin{proof}
Item (S1) is V13 Kato transport.  Item (S2) is
Lemma~\ref{lem:v13-spectator-rotation}.  Item (S3) is
Theorem~\ref{thm:v14-projection-U0}.  The two-stage statement is
Theorem~\ref{thm:v14-unitary-extension} and
Proposition~\ref{prop:v14-kappa-criterion}.
\end{proof}

\begin{center}
\fbox{\parbox{0.92\textwidth}{
\textbf{The universal controlled-cancellation conjecture is false without
additional carrier hypotheses.}
V13's controlled partial isometry is a stable datum.  An unstabilized native
path exists exactly when the two zero-mode projections belong to the same
identity-component orbit of the controlled local unitary group.  Small-angle
pairs and pairs sharing a controlled internal buffer satisfy this criterion
constructively.  A propagation-zero unstable bundle example shows that no
proof can follow from controlled Murray--von Neumann equivalence alone.}}
\end{center}

\section{V14 anti-circularity audit}
\label{sec:v14-audit}

\begin{enumerate}[leftmargin=2.4em]
\item \textbf{Stable equivalence was not silently called cancellation.}
Theorem~\ref{thm:v14-unitary-extension} exhibits the missing complement
partial isometry explicitly.
\item \textbf{A unitary extension was not silently called a path.}
Theorem~\ref{thm:v14-projection-U0} isolates the identity-component condition.
\item \textbf{The obstruction is not inferred only from abstract $K_1$.}
The primary object is the actual controlled $U_0$ orbit.  The complement
semigroup defect and the relative unitary-component coset are exact diagnostics
of that orbit problem.
\item \textbf{The counterexample is local.}
All matrices in Countertheorem~\ref{cth:v14-unstable-cancellation} have zero
spatial propagation.  Its failure is unstable complement topology, not a hidden
long-range operator.
\item \textbf{The counterexample is not overpromoted to the high-cutoff
physical branch.}
It is stably trivial, and Proposition~\ref{prop:v14-onsite-stable-range} shows
that the purely finite-dimensional obstruction disappears in the appropriate
large ambient stable range.
\item \textbf{Positive cancellation criteria are constructive.}
The small-angle unitary is an explicit polar direct rotation; the internal
buffer theorem gives an explicit exponential rotation in the original carrier.
\item \textbf{The V11 hard-shell theorem is invoked only after an actual native
path has been constructed.}
Theorem~\ref{thm:v14-unstable-native-lift} uses the $U_0$ path first and only
then applies the inherited fixed-index shell/universality estimates.
\end{enumerate}

\section{V14 status ledger}
\label{sec:v14-status}

\subsection*{Proved in V14}
\begin{enumerate}[leftmargin=2.2em]
\item A controlled partial isometry $v:p\to q$ extends to a controlled unitary
on the original carrier if and only if the complementary projections
$I-p$ and $I-q$ are themselves controlled Murray--von Neumann equivalent.  An
extension is exactly $u=v+w$.
\item For a fixed $v$, once complement cancellation holds, the remaining
extension obstruction is the relative component coset
$\kappa_\beta(v)$; it vanishes exactly when $v$ has an extension in the
identity component of the controlled unitary group.
\item Two controlled projections are joined by a controlled projection path if
and only if they are conjugate by a unitary in
$\mathcal U_0(\mathfrak A_\beta)$.  The proof uses a finite subdivision and an
explicit local direct-rotation formula for pairs at norm distance less than
one.
\item Controlled zero-mode equivalence alone does not imply unstabilized
cancellation in an arbitrary local coefficient algebra.  A propagation-zero
$S^6$ clutching example gives equivalent zero-mode ranges with nonisomorphic
complements and therefore no ambient unitary extension.
\item The unstable obstruction vanishes constructively for small-angle
zero-mode packets and for packets sharing a controlled internal buffer already
inside the original carrier.
\item The controlled $U_0$ criterion lifts through the V13 affine reduced graph
coordinate to a complete unstabilized V11-admissible native fixed-index path.
Hence regulators satisfying this criterion have the same matched global
coefficientwise continuum without any spectator stabilization.
\item On a finite-dimensional parameter panel in the V12 stable ambient range,
purely onsite bundle isomorphism already implies identity-component unitary
conjugacy.  Any remaining high-cutoff obstruction must therefore be genuinely
controlled/spatial (or the separate real Pfaffian orientation class).
\item V13 stable controlled equivalence and V14 unstabilized equivalence differ
precisely by the cancellation obstruction
$\operatorname{Can}_\beta$.
\end{enumerate}

\subsection*{Inherited}
\begin{enumerate}[leftmargin=2.2em]
\item V1--V7: source-complete coefficientwise continuum, common BV restoration,
arbitrary-fixed-loop forest theorem, ultraviolet Cauchy convergence,
scheme/reference covariance and exact perturbative-order projectivity.
\item V8--V10: relative regulator forests, the two-scale universality estimate,
and regulator-class lifting on the protected index-zero branches.
\item V11--V12: minimal fixed-index geometry, reduced pseudoinverse/source
calculus, global zero-mode bundle/line transport and global fixed-index
universality inside one protected native component.
\item V13: exponential locality of protected Riesz zero modes, the controlled
zero-mode module invariant, local reconstruction of the affine reduced graph,
and stable native lifting/universality after one inert massive copy.
\end{enumerate}

\subsection*{Conditional}
\begin{enumerate}[leftmargin=2.2em]
\item The exact cancellation criteria are statements in the chosen weighted
local algebra.  A loss of the V13 exponential-resolvent bound or reduced
singular-value floor can change the controlled component problem.
\item Countertheorem~\ref{cth:v14-unstable-cancellation} proves that universal
cancellation is false in general, but its unstable finite-dimensional defect
may disappear in the large-cutoff stable range of a fixed physical panel.
\item The internal-buffer theorem is load-bearing only when such a buffer is
actually present in the original Standard-Model representation/source carrier;
V14 does not assume one exists universally.
\item The V12 real Pfaffian mod-two orientation condition remains independent.
\item All continuum statements remain coefficientwise at prescribed finite
$(L,\Delta,r,N)$ and positive IR/reference scale.
\end{enumerate}

\subsection*{Open beyond V14}
\begin{enumerate}[leftmargin=2.2em]
\item Prove or falsify \emph{automatic controlled cancellation in the actual
Renewal Standard-Model carrier}: show that the protected reduced hard
complement supplies a uniformly local internal buffer (or an equivalent
stable-rank/cancellation property) for every fixed finite source/background
panel at sufficiently large cutoff.
\item If automatic cancellation fails, construct a genuine physical witness of
$\operatorname{Can}_\beta\ne0$ in the native spacetime-local algebra, not only
an unstable finite-dimensional coefficient example.
\item Classify the controlled unitary component/corner coset on noncontractible
gauge-background components and relate it, where possible, to a computable
local/coarse $K_1$ or transport index.
\item Compute the V12 real Pfaffian mod-two spectral-flow character on any
independently physical Majorana sector.
\item Analyze exits from the minimal fixed-index stratum and the independent
positive-IR/nonperturbative problems retained by V13.
\end{enumerate}

\subsection*{Exact next load-bearing theorem}
\begin{center}
\fbox{\parbox{0.92\textwidth}{
\textbf{Automatic internal cancellation in the Renewal ESM carrier.}
For each prescribed finite coefficient/source/background panel and fixed index
$k$, analyze the actual protected reduced Standard-Model chiral carrier at large
cutoff.  Prove that it contains a source-smooth exponentially local internal
buffer of rank $|k|$ (or prove an equivalent cancellation/stable-rank theorem
for the relevant local operator corner) so that every V13 controlled zero-mode
equivalence has $\operatorname{Can}_\beta=0$ without adding a spectator.
Otherwise produce an explicit native controlled complement/unitary-component
witness.  A positive theorem, combined with
Theorem~\ref{thm:v14-unstable-universality}, would close the distinction between
stable and unstabilized global fixed-index regulator universality.}}
\end{center}

\section*{Append-only continuation marker: V14}
\addcontentsline{toc}{section}{Append-only continuation marker: V14}
\textbf{End of V14.}  The complete V1--V13 body above is retained verbatim.
V14 proves that controlled zero-mode equivalence is a stable
Murray--von Neumann statement and isolates the exact extra data needed for an
unstabilized microscopic path: controlled complement cancellation followed by
the identity-component condition in the local unitary group.  It proves the
exact $U_0$-orbit criterion for controlled projection homotopy, gives a
propagation-zero unstable cancellation counterexample, and supplies explicit
small-angle and internal-buffer cancellation constructions.  Whenever the
$U_0$ criterion holds, the V13 affine graph reduction produces an unstabilized
V11-admissible native path and hence global fixed-index regulator universality
without a spectator.  The next gate is to decide whether the actual protected
Renewal ESM carrier automatically supplies the required internal cancellation
at large cutoff.



\clearpage
\section{V15 continuation: finite-rank localization, coefficientwise internal reserve, and automatic unstabilized lifting}
\label{sec:v15-continuation}

V14 identified the exact distinction between stable and unstabilized native
phase lifting.  A controlled equivalence of two zero-mode projections need not
cancel in an arbitrary local coefficient algebra; an unstabilized path requires
membership in the same controlled $\mathcal U_0$ orbit.  The V14 handoff asked
whether the actual Renewal ESM carrier supplies the missing reserve
automatically at sufficiently large cutoff.

There are two reasons not to attack this by a blanket stable-rank assertion.
First, the microscopic spacetime algebra is four-dimensional and contains far
more projections than the fixed-rank protected zero-mode packet.  Cancellation
of the entire local operator algebra is therefore a much stronger question than
the one needed here.  Second, the continuum theorem is coefficientwise: for a
prescribed finite source/contact order one needs the corresponding finite source
jet of the regulator path, not a single cutoff-independent finite-amplitude
source ball.  The same quantifier correction already used in V2 is load-bearing
again.

The present continuation proves a positive theorem at exactly that scope.  A
rank-$\kappa$ protected Riesz projector with a cutoff-uniform exponential
kernel bound has a uniformly localized orthonormal frame.  A finite bank of
bounded gauge/spin-covariant lattice translations of that frame supplies, at
large enough cutoff, a finite local candidate space of arbitrarily large fixed
multiplicity.  Because the retained source algebra is finite, one can choose a
rank-$\kappa$ subpacket which is orthogonal to \emph{every} source jet of both
endpoint zero-mode packets through the requested order.  In the truncated
source algebra this is an exact V14 internal buffer.  Its additional translation distance from the original localized zero-mode
centres is bounded in lattice units independently of the cutoff and therefore
has physical transport scale $O(a)$.

Thus the V13 spectator is unnecessary for every fixed coefficient/source
packet on a protected local chart.  This is deliberately not advertised as
cancellation of the full four-dimensional Roe-type algebra, nor as a global
finite-amplitude path over an arbitrary noncontractible gauge component.

Throughout V15 let $X_a$ be the periodic regulator lattice at spacing $a$,
with its graph distance $d_a$ measured in lattice units and a fixed finite
onsite fermion fibre $\mathbb C^{d_f}$ in the relevant chiral/representation
block.  Bounded geometry means that for every fixed integer $R$,
\begin{equation}
 N_R:=\sup_{a,x}|B_R(x)|<\infty.
 \label{eq:v15-bounded-geometry}
\end{equation}
The V13 weighted kernel algebra at exponent $\alpha>0$ is denoted
$\mathfrak A_{\alpha,a}$.

\subsection{A finite-rank exponentially local projector has a localized frame}
\label{sec:v15-localized-frame}

We begin with a quantitative statement which is independent of the overlap
specialization.

\begin{lemma}[Finite-propagation truncation]
\label{lem:v15-truncation}
Let $A\in\mathfrak A_{\alpha,a}$ and define
\begin{equation}
 A^{(R)}_{xy}:=\mathbf 1_{\{d_a(x,y)\le R\}}A_{xy}.
 \label{eq:v15-truncate}
\end{equation}
Then
\begin{equation}
 \boxed{
 \|A-A^{(R)}\|_{\rm op}
 \le e^{-\alpha R}\|A\|_\alpha .}
 \label{eq:v15-truncation-bound}
\end{equation}
The same estimate holds coefficientwise for every retained source derivative.
\end{lemma}

\begin{proof}
For every row,
\[
 \sum_{d_a(x,y)>R}\|A_{xy}\|
 \le e^{-\alpha R}
      \sum_y e^{\alpha d_a(x,y)}\|A_{xy}\|
 \le e^{-\alpha R}\|A\|_\alpha,
\]
and the same estimate holds for columns.  The Schur bound gives
\eqref{eq:v15-truncation-bound}.  Source derivatives are separate elements of
the same weighted algebra and obey the identical argument.
\end{proof}

\begin{lemma}[Uniform diagonal floor for a finite-rank local projection]
\label{lem:v15-diagonal-floor}
Let $p=p^*=p^2\in\mathfrak A_{\alpha,a}$ be nonzero and assume
\begin{equation}
 \|p\|_\alpha\le C_p
 \label{eq:v15-p-local-bound}
\end{equation}
uniformly in the cutoff.  Choose $R$ so that
$C_pe^{-\alpha R}\le1/4$.  Then
\begin{equation}
 \boxed{
 \max_x\|p_{xx}\|_{\rm op}
 \ge c_0:=\frac{3}{4N_R}.}
 \label{eq:v15-diagonal-floor}
\end{equation}
Consequently there are a site $x_0$ and a unit vector
$\xi\in\mathbb C^{d_f}$ such that
\begin{equation}
 \|p\iota_{x_0}\xi\|^2
 =\langle\xi,p_{x_0x_0}\xi\rangle
 \ge c_0,
 \label{eq:v15-column-floor}
\end{equation}
where $\iota_x$ is the canonical inclusion of the onsite fibre.
\end{lemma}

\begin{proof}
Lemma~\ref{lem:v15-truncation} gives
$\|p-p^{(R)}\|\le1/4$.  Since a nonzero projection has operator norm one,
\begin{equation}
 \|p^{(R)}\|\ge3/4.
 \label{eq:v15-pr-trunc-norm}
\end{equation}
Put $M:=\max_x\|p_{xx}\|$.  Positivity of the $2\times2$ block compression of
$p$ to the onsite fibres at $x,y$ gives the operator Cauchy--Schwarz estimate
\begin{equation}
 \|p_{xy}\|^2\le\|p_{xx}\|\,\|p_{yy}\|\le M^2.
 \label{eq:v15-block-CS}
\end{equation}
Every row and column of $p^{(R)}$ has at most $N_R$ nonzero blocks, hence its
Schur norm is at most $N_RM$.  Therefore
$3/4\le\|p^{(R)}\|\le N_RM$, proving
\eqref{eq:v15-diagonal-floor}.  Taking a top eigenvector of the positive block
$p_{x_0x_0}$ gives \eqref{eq:v15-column-floor}.
\end{proof}

\begin{theorem}[Uniformly localized orthonormal frame]
\label{thm:v15-localized-frame}
Fix a rank $\kappa<\infty$.  Let
$p_a\in\mathfrak A_{\alpha,a}$ be rank-$\kappa$ projections with
$\|p_a\|_\alpha\le C_p$ uniformly in the cutoff.  Then there exist an exponent
$0<\beta<\alpha$, constants $C_{\rm fr},c_{\rm fr}>0$ depending only on
$(\kappa,\alpha,C_p,d_f)$ and the bounded-geometry constants, sites
$x_{1,a},\ldots,x_{\kappa,a}$, and an orthonormal basis
$\psi_{1,a},\ldots,\psi_{\kappa,a}$ of $\Ran p_a$ such that
\begin{equation}
 \boxed{
 \sum_x e^{\beta d_a(x,x_{j,a})}\|\psi_{j,a}(x)\|
 \le C_{\rm fr}}
 \label{eq:v15-frame-l1}
\end{equation}
and therefore
\begin{equation}
 \|\psi_{j,a}(x)\|
 \le C_{\rm fr}e^{-\beta d_a(x,x_{j,a})}.
 \label{eq:v15-frame-pointwise}
\end{equation}
All constants are independent of the volume and cutoff.

If $p_a(\zeta)$ lies on one star-shaped V11 protected source/background chart
and its prescribed source derivatives are uniformly local, the reference frame
may be transported by the V13 local Kato equation.  For every fixed source
multi-index $I$ in the packet,
\begin{equation}
 \sum_x e^{\beta' d_a(x,x_{j,a})}
 \|\partial_I\psi_{j,a}(0;x)\|
 \le C_I
 \label{eq:v15-frame-source-jets}
\end{equation}
for one fixed $0<\beta'<\beta$.
\end{theorem}

\begin{proof}
Apply Lemma~\ref{lem:v15-diagonal-floor} to $p_a$.  Choose $x_{1,a},\xi_1$
as there and put
\begin{equation}
 \psi_{1,a}:=
 \frac{p_a\iota_{x_{1,a}}\xi_1}
      {\|p_a\iota_{x_{1,a}}\xi_1\|}.
 \label{eq:v15-first-frame-vector}
\end{equation}
The weighted column bound for $p_a$ and the denominator floor give
\eqref{eq:v15-frame-l1} for $j=1$ at exponent $\alpha$.
The rank-one projection
$e_{1,a}=|\psi_{1,a}\rangle\langle\psi_{1,a}|$ therefore belongs to
$\mathfrak A_{\beta_1,a}$, for any fixed $\beta_1<\alpha$, with a uniform
bound: use
$d_a(x,y)\le d_a(x,x_{1,a})+d_a(y,x_{1,a})$ and
\eqref{eq:v15-frame-l1}.

Because $\psi_{1,a}\in\Ran p_a$, one has $e_{1,a}\le p_a$ and
$p_a-e_{1,a}$ is a projection of rank $\kappa-1$.  Its weighted norm at
$\beta_1$ is bounded by a constant depending only on the preceding uniform
bounds.  Apply Lemma~\ref{lem:v15-diagonal-floor} again to this residual
projection.  Iterating $\kappa$ times and choosing in advance a decreasing
finite exponent chain
\[
 \alpha>\beta_1>\cdots>\beta_\kappa=:\beta>0
\]
produces the claimed orthonormal localized frame with uniform constants.  The
rank is fixed, so no constant depends on the cutoff.

For the source statement, use this frame at the reference point and solve the
V13 Kato transport along the radial path in the star-shaped source chart.  V13
proved that the Kato generator and all prescribed fixed derivatives belong to a
weighted algebra after an arbitrarily small exponent loss.  Acting with a
weighted local operator on a vector satisfying \eqref{eq:v15-frame-l1}
preserves an exponential moment about the same centre.  Finite differentiation
of the Kato ODE therefore proves \eqref{eq:v15-frame-source-jets}.
\end{proof}

\subsection{Bounded covariant translations provide arbitrarily many local copies}
\label{sec:v15-translation-bank}

The V2 native lattice already contains gauge/spin parallel shifts.  For a fixed
integer displacement $s\in\mathbb Z^4$, choose one bounded path from $x$ to
$x+s$ and denote the corresponding unitary parallel shift by
$\mathsf T_s^{\mathcal U}$.  Its propagation is at most a constant multiple of
$|s|_1$, it is exactly gauge/spin covariant, and every prescribed finite source
derivative is finite range.

\begin{lemma}[Overlap decay of translated localized vectors]
\label{lem:v15-shift-overlap}
Let $\psi,\phi$ satisfy exponential bounds of the form
\eqref{eq:v15-frame-pointwise} about centres $x_\psi,x_\phi$.  For every
$0<\beta_0<\beta$ there is $C_{\beta_0}$ such that
\begin{equation}
 \boxed{
 |\langle\mathsf T_s^{\mathcal U}\psi,
          \mathsf T_t^{\mathcal U}\phi\rangle|
 \le C_{\beta_0}
 e^{-\beta_0 d_a(x_\psi+s,x_\phi+t)}.}
 \label{eq:v15-shift-overlap}
\end{equation}
The bound is uniform in the cutoff and in the unitary parallel transports.
\end{lemma}

\begin{proof}
Parallel transport does not change pointwise norms.  Hence the absolute value
of the inner product is bounded by a convolution of two exponentials centred
at $x_\psi+s$ and $x_\phi+t$.  The triangle inequality gives an overall factor
$e^{-\beta_0d_a(x_\psi+s,x_\phi+t)}$ after reducing the exponent from $\beta$
to $\beta_0$; the remaining lattice convolution is uniformly summable by
bounded geometry.
\end{proof}

\begin{lemma}[Finite bounded-displacement separation]
\label{lem:v15-shift-selection}
Fix integers $\kappa,M\ge1$ and a separation radius $D$.  For every sufficiently
large periodic lattice and every set of centres
$X=\{x_1,\ldots,x_\kappa\}$ there exist displacements
\begin{equation}
 s_1,\ldots,s_M\in\mathbb Z^4,
 \qquad
 |s_j|_1\le S_*(\kappa,M,D),
 \label{eq:v15-shifts}
\end{equation}
with $S_*$ independent of the cutoff and of $X$, such that
\begin{equation}
 d_a(x_b+s_i,x_c+s_j)\ge D
 \quad
 (i\\ne j,\\ 1\\le b,c\\le\\kappa).
 \label{eq:v15-shift-separated}
\end{equation}
\end{lemma}

\begin{proof}
Choose an integer $L>2D$ and consider the candidate displacements
$nLe_1$, $1\le n\le N$.  For a fixed candidate $n$ and a fixed ordered pair
$(b,c)$, there is at most one other candidate $m$ for which
$x_b+nLe_1$ and $x_c+mLe_1$ are within distance $D$, provided the torus side
length exceeds $2NL+2D$.  Hence the conflict graph on the $N$ candidates has
maximum degree at most $2\kappa^2$.  Taking
$N\ge M(2\kappa^2+1)$, the greedy independent-set bound gives $M$ candidates
with no conflicts.  Their maximum displacement is
$S_*:=NL$.  All these numbers depend only on $(\kappa,M,D)$.
\end{proof}

\begin{corollary}[A uniformly conditioned translated-copy bank]
\label{cor:v15-copy-bank}
Fix $M$.  For all sufficiently large cutoffs, the displacements in
Lemma~\ref{lem:v15-shift-selection} may be chosen so that the $M\kappa$ vectors
\begin{equation}
 \mathsf T_{s_j}^{\mathcal U}\psi_b,
 \qquad
 1\le j\le M,\ 1\le b\le\kappa,
 \label{eq:v15-copy-vectors}
\end{equation}
have Gram matrix $G_0$ satisfying
\begin{equation}
 \boxed{\frac12I\preceq G_0\preceq\frac32I.}
 \label{eq:v15-copy-Gram}
\end{equation}
The raw synthesis operator of these vectors and its orthonormalization
$G_0^{-1/2}$ are uniformly exponentially local with propagation/locality
constants independent of the cutoff.
\end{corollary}

\begin{proof}
Choose $D$ in Lemma~\ref{lem:v15-shift-selection} so large that the right-hand
side of \eqref{eq:v15-shift-overlap} is less than
$[2M\kappa]^{-1}$ whenever the two copy labels differ.  Within one translated
copy the vectors remain exactly orthonormal because
$\mathsf T_s^{\mathcal U}$ is unitary.  The row-sum estimate for the off-diagonal
part of the Gram matrix is therefore less than $1/2$, proving
\eqref{eq:v15-copy-Gram}.  The synthesis operator is a finite sum of bounded
propagation shifts acting on the localized frame.  Functional calculus of the
finite Gram matrix preserves all uniform fixed-source bounds.
\end{proof}

\subsection{The finite source-jet algebra and the exact annihilator count}
\label{sec:v15-source-jet}

Let $\mathbf R_\varrho$ denote the complete finite graded
source/BV algebra retained by the prescribed packet through total source degree
$\varrho$; it includes the even central writers and the odd ghost/antifield
writers already used in V4--V7.  Only its finite vector-space dimension is
load bearing:
\begin{equation}
 \boxed{\nu_\varrho:=\dim_{\mathbb C}\mathbf R_\varrho<\infty.}
 \label{eq:v15-source-monomials}
\end{equation}
For $s$ commuting even coordinates alone this specializes to
$\mathbf R_\varrho=\mathbb C[J^1,\ldots,J^s]/(J)^{\varrho+1}$ and
$\nu_\varrho=\binom{s+\varrho}{\varrho}$.  With odd writers present, the
ordinary monomial basis is replaced by the corresponding finite graded basis;
the argument below uses only the number $\nu_\varrho$ and the fact that
multiplication is truncated at total degree $\varrho$.

Fix two rank-$\kappa$ protected zero-mode packets $p(J),q(J)$ in this truncated
algebra and a V13 controlled equivalence
\begin{equation}
 v(J)^*v(J)=p(J),
 \qquad
 v(J)v(J)^*=q(J).
 \label{eq:v15-controlled-v}
\end{equation}
Choose the V15 localized Kato frame synthesis map
\begin{equation}
 \Psi(J):\mathbb C^\kappa\longrightarrow\mathcal H_a,
 \qquad
 \Psi(J)^*\Psi(J)=I,
 \qquad
 \Psi(J)\Psi(J)^*=p(J),
 \label{eq:v15-Psi}
\end{equation}
and put
\begin{equation}
 \mathsf X(J):=v(J)\Psi(J),
 \qquad
 \mathsf X(J)^*\mathsf X(J)=I,
 \qquad
 \mathsf X(J)\mathsf X(J)^*=q(J).
 \label{eq:v15-Chi}
\end{equation}
All identities are in $\mathbf R_\varrho$.

Choose
\begin{equation}
 M\ge2\nu_\varrho+1
 \label{eq:v15-M}
\end{equation}
and the bounded shifts of
Corollary~\ref{cor:v15-copy-bank}.  Let
$\widetilde W(J):\mathbb C^{M\kappa}\to\mathcal H_a$ be the raw synthesis
operator whose $j$th block is
$\mathsf T_{s_j}^{\mathcal U}(J)\Psi(J)$.  Since the constant coefficient of
its Gram matrix has the floor \eqref{eq:v15-copy-Gram}, its inverse square root
exists uniquely in the finite nilpotent extension of the positive constant
Gram matrix by analytic matrix functional calculus (equivalently, by the
finite recursive Sylvester equations for its source coefficients).
Define
\begin{equation}
 W(J):=\widetilde W(J)
 \bigl(\widetilde W(J)^*\widetilde W(J)\bigr)^{-1/2}.
 \label{eq:v15-W}
\end{equation}
Then
\begin{equation}
 \boxed{W(J)^*W(J)=I_{M\kappa}}
 \label{eq:v15-W-isometry}
\end{equation}
exactly in $\mathbf R_\varrho$.

\begin{theorem}[Finite-jet common-complement reserve]
\label{thm:v15-jet-reserve}
For every fixed finite source order $\varrho$, and every sufficiently large
cutoff, there is an isometry
\begin{equation}
 Z:\mathbb C^\kappa\longrightarrow\mathbb C^{M\kappa}
 \label{eq:v15-Z}
\end{equation}
with constant scalar coefficients such that, with
\begin{equation}
 \mathsf E(J):=W(J)Z,
 \qquad
 r(J):=\mathsf E(J)\mathsf E(J)^*,
 \label{eq:v15-Eta-r}
\end{equation}
one has in the truncated source algebra
\begin{align}
 \mathsf E(J)^*\mathsf E(J)&=I_\kappa,
 \label{eq:v15-Eta-isometry}\\
 p(J)r(J)&=r(J)p(J)=0,
 \label{eq:v15-pr-zero}\\
 q(J)r(J)&=r(J)q(J)=0.
 \label{eq:v15-qr-zero}
\end{align}
The projection $r(J)$ and every coefficient of its retained source jet are
uniformly exponentially local.  Its propagation/locality radius depends on
$(\kappa,\nu_\varrho)$ and the fixed locality constants, but not on the
cutoff.  Its added displacement from the original zero-mode localization
centres is bounded in lattice units, hence has physical transport scale $O(a)$.
\end{theorem}

\begin{proof}
Orthogonality to $p$ and $q$ is equivalent, respectively, to annihilation by
$\Psi^*$ and $\mathsf X^*$.  Define
\begin{equation}
 F(J):=
 \begin{pmatrix}
  \Psi(J)^*W(J)\\
  \mathsf X(J)^*W(J)
 \end{pmatrix}
 :\mathbb C^{M\kappa}\longrightarrow\mathbb C^{2\kappa}.
 \label{eq:v15-F}
\end{equation}
Choose once and for all a homogeneous vector-space basis
$(m_A)_{A=1}^{\nu_\varrho}$ of $\mathbf R_\varrho$ and write the finite graded
source expansion
\begin{equation}
 F=\sum_{A=1}^{\nu_\varrho}m_A F_A.
 \label{eq:v15-F-expansion}
\end{equation}
(The Koszul signs are already contained in the multiplication table of
$\mathbf R_\varrho$ and do not affect the following linear dimension count.)
Stack all coefficient maps:
\begin{equation}
 \mathbf F:
 \mathbb C^{M\kappa}
 \longrightarrow
 \bigoplus_{A=1}^{\nu_\varrho}\mathbb C^{2\kappa},
 \qquad
 \mathbf Fz=(F_Az)_A.
 \label{eq:v15-stacked-F}
\end{equation}
Pure dimension counting gives
\begin{equation}
 \dim\ker\mathbf F
 \ge M\kappa-2\kappa\nu_\varrho
 \ge\kappa
 \label{eq:v15-kernel-dim}
\end{equation}
by \eqref{eq:v15-M}.  Choose any orthonormal $\kappa$-frame in this kernel and
let $Z$ be its synthesis isometry.  Then every coefficient of $FZ$ in the retained graded basis vanishes; equivalently,
\begin{equation}
 FZ=0\quad\text{in }\mathbf R_\varrho.
 \label{eq:v15-FZ-zero}
\end{equation}
Since $W^*W=I$ and $Z^*Z=I$,
\eqref{eq:v15-Eta-isometry} follows.  Equation
\eqref{eq:v15-FZ-zero} says
$\Psi^*\mathsf E=0=\mathsf X^*\mathsf E$.  Multiplying by the frame synthesis maps gives
\eqref{eq:v15-pr-zero}--\eqref{eq:v15-qr-zero}.

Locality follows from Corollary~\ref{cor:v15-copy-bank}: $W$ is obtained from a
finite bank of fixed bounded-displacement native shifts and a finite Gram
inverse whose constant coefficient has a uniform floor.  Multiplication by the
fixed finite matrix $Z$ changes no spatial support estimate.  Every source
coefficient is a finite algebraic combination of the same bounded local jets.
Finally the largest displacement $S_*$ is independent of $a$, so its physical
length is $aS_*=O(a)$.
\end{proof}

\begin{remark}[Frame and gauge/spin covariance]
\label{rem:v15-frame-covariance}
The proof used one V11 Kato frame only to write the finite-dimensional
annihilator problem.  Under a change of zero-mode frame, the synthesis maps
$\Psi,\mathsf X$ and the translated-copy bank are right multiplied by the
corresponding finite unitary matrices.  Transporting the coefficient isometry
$Z$ by the same finite change leaves the physical projection $r=\mathsf E\mathsf E^*$
and the partial isometries below conjugated by the exact gauge/spin action.
Thus the construction defines a covariant buffer on the framed protected
chart; different local frame choices give V7/V12-equivalent auxiliary paths,
not different endpoint regulators.
\end{remark}

\subsection{Automatic coefficientwise internal buffer}
\label{sec:v15-buffer}

Define
\begin{equation}
 a(J):=\mathsf E(J)\Psi(J)^*,
 \qquad
 b(J):=\mathsf E(J)\mathsf X(J)^*.
 \label{eq:v15-buffer-partials}
\end{equation}

\begin{theorem}[Automatic internal reserve at fixed source order]
\label{thm:v15-automatic-buffer}
Under the hypotheses of Theorem~\ref{thm:v15-jet-reserve},
\begin{align}
 a^*a&=p,& aa^*&=r,
 \label{eq:v15-a-buffer}\\
 b^*b&=q,& bb^*&=r,
 \label{eq:v15-b-buffer}
\end{align}
exactly in $\mathbf R_\varrho$, while $r$ is orthogonal to both endpoint
zero-mode packets.  Hence $r$ is a V14 internal buffer of rank $\kappa$ for the
complete retained source jet.

Moreover, because $r$ lies in the orthogonal complement of the protected zero
modes, the V11 reduced singular-value floor makes it a nonzero/massive packet at
both endpoints.  No new fermion species or carrier copy has been added.
\end{theorem}

\begin{proof}
Use $\mathsf E^*\mathsf E=I$ and
$\Psi^*\Psi=\mathsf X^*\mathsf X=I$:
\[
 a^*a=\Psi\mathsf E^*\mathsf E\Psi^*=\Psi\Psi^*=p,
 \qquad
 aa^*=\mathsf E\Psi^*\Psi\mathsf E^*=\mathsf E\mathsf E^*=r,
\]
and similarly for $b$.  Orthogonality is
Theorem~\ref{thm:v15-jet-reserve}.  On a V11 protected zero-mode chart,
$D^*D\succeq\tau_*^2$ on the orthogonal complement of the kernel (and likewise
on the cokernel side).  Thus the buffer is contained in the reduced nonzero
packet at each endpoint.  It is an internal direction of the existing carrier.
\end{proof}

The V14 exponential rotation can now be performed coefficientwise.  Put
\begin{equation}
 U_a(\theta)
 =I-p-r+\cos\theta\,(p+r)+\sin\theta\,(a-a^*),
 \label{eq:v15-Ua}
\end{equation}
with the analogous $U_b$.  Because
\eqref{eq:v15-a-buffer} and the orthogonality relations are exact identities in
$\mathbf R_\varrho$, the ordinary V14 proof gives
\begin{equation}
 U_a(\theta)^*U_a(\theta)=I,
 \qquad
 U_a(\pi/2)pU_a(\pi/2)^*=r
 \label{eq:v15-Ua-exact}
\end{equation}
inside that same truncated algebra.  Concatenating with the reversed $q$ to
$r$ rotation gives an exact source-jet projection path from $p$ to $q$.

\begin{corollary}[Coefficientwise cancellation obstruction vanishes at large cutoff]
\label{cor:v15-Can-zero}
For a prescribed finite source order $\varrho$, define
$\operatorname{Can}^{[\varrho]}_\beta(p,q)$ by the V14 $\mathcal U_0$ criterion
inside the truncated source algebra.  Then for every V13 controlled-equivalent
rank-$\kappa$ protected pair there is $a_0>0$ such that
\begin{equation}
 \boxed{
 \operatorname{Can}^{[\varrho]}_\beta(p_a,q_a)=0
 \qquad(0<a<a_0).}
 \label{eq:v15-Can-zero}
\end{equation}
The required path is uniformly local and uses only the original Renewal ESM
fermion carrier.
\end{corollary}

\begin{proof}
The only cutoff threshold is the geometric requirement in
Lemma~\ref{lem:v15-shift-selection}: the torus must be large enough to contain
the finite bounded-displacement copy bank without wraparound collisions.
Theorem~\ref{thm:v15-automatic-buffer} then verifies the V14 internal-buffer
hypothesis in $\mathbf R_\varrho$.  Apply
Theorem~\ref{thm:v14-buffer-cancellation} coefficientwise.
\end{proof}

\subsection{Why the coefficientwise quantifier is sufficient}
\label{sec:v15-quantifier}

The preceding theorem does not construct one analytic buffer on a
cutoff-independent finite-amplitude ball of all sources.  Such a theorem is not
needed for the continuum target.

\begin{theorem}[Finite-jet native lifting is sufficient for a fixed coefficient packet]
\label{thm:v15-jet-suffices}
Fix a perturbative/source packet
\begin{equation}
 \mathfrak q=(L,\Delta,\varrho,N,\mathcal P_{\rm ext})
 \label{eq:v15-fixed-q}
\end{equation}
and the positive IR/reference prescription.  Suppose two V11 protected native
regulators have V13 controlled-equivalent zero-mode packets.  At sufficiently
large cutoff, the coefficientwise internal-buffer path constructed above,
together with the V13 affine reduced-graph interpolation, supplies every
regulator-path source jet required by the V8--V11 relative forest theorem
through order $\varrho$.  Therefore no analytic finite-amplitude source radius
is required to compare the order-$\mathfrak q$ continuum coefficients.
\end{theorem}

\begin{proof}
Every operation entering the fixed-order shell proof---inverse differentiation,
Riesz/Kato differentiation, determinant/Pfaffian transgression, forest
subtraction, BV brackets and source/contact extraction---is a finite polynomial
or rational jet operation whose denominators are evaluated on the protected
reference pivots.  This is exactly the V2 reference-pivot coefficient theorem
and the V5 contraction-complete induction.  The V15 rotation identities are
exact in $\mathbf R_\varrho$, hence all of their coefficients through source
order $\varrho$ agree with those of an exact path.  Terms of degree
$\varrho+1$ or higher cannot feed a retained coefficient.  The reduced graph
coordinate is treated by V13 in the same finite source algebra.  Thus the V8
marked-forest derivative has all required regulator/source coefficients without
postulating a common nonzero source radius.
\end{proof}

\begin{firewall}[No analytic-source promotion]
Theorem~\ref{thm:v15-jet-suffices} is coefficientwise.  It does not assert that
one V15 buffer remains exactly orthogonal to both zero-mode packets on an open
finite-amplitude source ball independent of the cutoff.  Proving such an
analytic-family cancellation theorem would be a stronger problem, analogous
to the stronger finite-coupling RG trajectory explicitly excluded from E1--E6.
\end{firewall}

\subsection{Unstabilized coefficientwise universality without a spectator}
\label{sec:v15-universality}

\begin{theorem}[Automatic unstabilized universality for the fixed-rank protected packet]
\label{thm:v15-unstabilized-universality}
Fix finite $\mathfrak q$ as in \eqref{eq:v15-fixed-q}.  Let $A$ and $B$ be two
V11 minimal fixed-index native regulators on one protected local source chart,
with the same physical Standard-Model principal Dirac--Yukawa jet,
representation, positive IR/reference prescription and global complex line
data.  Assume their rank-$\kappa=|k|$ kernel packets ($k>0$) or cokernel packets
($k<0$) are V13 controlled-equivalent with the required source jets.

Then, for all sufficiently large cutoffs, the complete regulator comparison can
be performed in the original carrier with no spectator at every source order
retained by $\mathfrak q$.  Consequently the matched continuum states obey
\begin{equation}
 \boxed{
 \widehat{\mathfrak C}^{B,(k),\rm sat}_{\infty,\lambda}
 =F_{\infty,\mathfrak q,k}^{B\leftarrow A}
  \widehat{\mathfrak C}^{A,(k),\rm sat}_{\infty,\lambda}}
 \label{eq:v15-unstabilized-universality}
\end{equation}
without adjoining the V13 inert massive copy.
The equality contains the complete retained connected/source system,
operator mixing, stress/contact coefficients, ghosts/antifields,
BV/BRST/master descendants and determinant/Pfaffian insertions, and remains
strictly projective in perturbative order.
\end{theorem}

\begin{proof}
Corollary~\ref{cor:v15-Can-zero} supplies the V14 unstabilized zero-mode path in
the finite source algebra of the chosen packet.  Theorem
\ref{thm:v15-jet-suffices} promotes it to exactly the regulator/source jets
consumed by the V8--V11 relative forest theorem.  V13 supplies the local affine
reduced-graph interpolation after the zero-mode packet has been transported.
The V11 reduced singular floor and the openness argument permit the usual
one-time shrinking of the protected chart.  Thus the V11 two-scale estimate
applies to the complete coefficient packet without stabilization.  Passing to
the V6 ultraviolet limit gives \eqref{eq:v15-unstabilized-universality}.  V12
supplies the global complex saturation on the faithful branch; the independent
real Pfaffian orientation condition is unchanged.  Chronological V7
normalization gives strict loop-order projectivity.
\end{proof}

\begin{corollary}[Stable and unstabilized coefficientwise universality coincide]
\label{cor:v15-stable-equals-unstable}
At fixed finite coefficient/source order and fixed finite zero-mode rank, the
V13 spectator does not enlarge the class of matched coefficientwise continua on
a protected local Renewal ESM chart.  Every V13 controlled-equivalent pair is
already unstabilized-equivalent for all sufficiently large cutoffs.
\end{corollary}

\begin{proof}
V13 proved the stable statement.  Theorem
\ref{thm:v15-unstabilized-universality} proves the same comparison without the
spectator under exactly the fixed-rank/source-jet hypotheses of the present
packet.
\end{proof}

\subsection{Why this is not cancellation of the full four-dimensional local algebra}
\label{sec:v15-Roe-firewall}

The fixed-rank/source-jet hypotheses are essential.  They must not be replaced
by the statement that every pair of projections in a four-dimensional local
operator algebra cancels.

\begin{remark}[External comparison: cancellation is a low-dimensional property]
\label{rem:v15-Roe-comparison}
For a bounded-geometry metric space $X$, Li and Willett proved that cancellation
of projections in the uniform Roe $C^*$-algebra $C_u^*(X)$ is equivalent to
asymptotic dimension zero; see K.~Li and R.~Willett,
\emph{Low-dimensional properties of uniform Roe algebras}, J. London Math.
Soc. 97 (2018), 98--124, Theorem~2.2.  A four-dimensional lattice has positive
asymptotic dimension, so no unrestricted Roe-algebra cancellation principle is
available.  This external theorem is not used in the proof above; it explains
why V15 keeps the fixed-rank and finite-source-jet quantifiers explicit.
\end{remark}

The V15 proof escapes that obstruction for a concrete reason.  Its candidate
multiplicity $M$ is finite and depends only on the fixed number of source
monomials and the fixed zero-mode rank.  The extra copies are not new fields:
they are bounded covariant translations by at most $S_*$ lattice steps of the
same finite zero-mode packet.  The large cutoff is used only to make room for
this finite local construction.

\subsection{What remains globally}
\label{sec:v15-global-boundary}

A noncontractible global background component can still carry information not
visible in one truncated source chart.  V12 already separated the full
zero-mode bundle and determinant/Pfaffian line data.  V14 further identified
the controlled unitary component.  V15 removes the latter locally at every
fixed source order by constructing an internal reserve, but it does not assert
that the independently chosen local reserves glue to one finite-amplitude
microscopic path around every large background loop.

\begin{proposition}[Overlap of two V15 buffers]
\label{prop:v15-buffer-overlap}
On an overlap of two protected source charts, two V15 buffer choices produce
unstabilized regulator comparison paths with the same endpoints.  Their
concatenation is a closed controlled native loop.  The induced normalized
continuum maps therefore differ, if at all, by the V7 finite scheme/source
transition together with the V12 global line/orientation transport.  On the
faithful complex Standard-Model branch the determinant-line part is trivial;
a real Pfaffian mod-two character remains separate.
\end{proposition}

\begin{proof}
Each local construction is a V11-admissible coefficientwise path with the same
endpoint regulators.  Concatenating one with the reverse of the other gives a
closed sequence of native regulator/scheme transports.  V7 identifies the
residual finite local-coordinate action with its scheme groupoid, and V12
identifies the independent global line transport.  The inherited faithful-spin
complex determinant section has trivial paired holonomy, so only the declared
finite local scheme change and any independent real orientation character can
remain.
\end{proof}

\begin{firewall}[Local coefficientwise cancellation is not a global $\mathcal U_0$ theorem]
V15 closes the spectator distinction for every fixed finite coefficient packet
on a protected local chart.  It does not prove that the full source-dependent
family of zero-mode projections on a noncontractible background component lies
in one identity-component orbit of a global exponentially local unitary group.
That stronger statement requires descent of the local buffer choices, or an
independent computation of the resulting controlled loop class.
\end{firewall}

\section{V15 anti-circularity audit}
\label{sec:v15-audit}

\begin{enumerate}[leftmargin=2.4em]
\item \textbf{The buffer is not assumed to exist.}
It is built from a localized frame of the actual protected Riesz projector and
a finite bank of native bounded parallel shifts.
\item \textbf{Uniform localization is proved from the native spectral/locality
hypotheses.}
The diagonal-floor argument uses only exponential kernel control, finite rank
and bounded geometry; it does not assume prelocalized zero modes.
\item \textbf{The large cutoff is used only geometrically.}
All required displacements are bounded in lattice units by
$S_*(\kappa,M,D)$, independent of the cutoff.  Their physical radius shrinks
like $a$.
\item \textbf{Only finitely many source constraints are paid.}
The rank count in \eqref{eq:v15-kernel-dim} is finite because source order and
source panel are fixed before the cutoff is removed.  No estimate uniform as
$\varrho\to\infty$ is asserted.
\item \textbf{Finite-amplitude source invertibility is not smuggled in.}
The internal-buffer identities are exact in the truncated source algebra; V2's
reference-pivot theorem is what makes that the correct quantifier for the
coefficientwise continuum.
\item \textbf{No global Roe-algebra cancellation is claimed.}
The theorem is a fixed-rank finite-jet result, consistent with the known failure
of unrestricted cancellation on positive-dimensional coarse spaces.
\item \textbf{The independent real Pfaffian orientation is untouched.}
A complex internal buffer does not calculate the V12 mod-two orientation
character.
\end{enumerate}

\section{V15 status ledger}
\label{sec:v15-status}

\subsection*{Proved in V15}
\begin{enumerate}[leftmargin=2.2em]
\item A uniformly exponentially local projection of fixed finite rank admits a
uniformly exponentially localized orthonormal frame.  The proof is constructive:
a finite-propagation truncation gives a cutoff-independent diagonal floor, one
localized projection column is extracted, and the argument iterates on the
orthogonal residual projection.
\item V13 local Kato transport preserves the localized-frame estimates for every
prescribed fixed source derivative after one harmless exponent loss.
\item Finitely many bounded gauge/spin-covariant lattice translations of that
frame can be chosen with a uniformly conditioned Gram matrix; the number and
maximum displacement depend only on the fixed zero-mode rank and desired copy
multiplicity, not on the cutoff.
\item At source order $\varrho$, taking
$M\ge2\nu_\varrho+1$, where $\nu_\varrho$ is the dimension of the complete
finite graded source/BV algebra (equal to $\binom{s+\varrho}{\varrho}$ in the
purely even $s$-writer case), leaves a
$\kappa$-dimensional common annihilator of all $p$- and $q$-overlap coefficients
through that order.  Orthogonalizing the copy bank and selecting this
annihilator yields a rank-$\kappa$ exponentially local projection $r$ which is
orthogonal to both endpoint zero-mode packets in the exact truncated source
algebra.
\item The explicit partial isometries
$a=\mathsf E\Psi^*$ and $b=\mathsf E\mathsf X^*$ make $r$ a V14 internal buffer.  Hence the
V14 cancellation obstruction vanishes coefficientwise through every prescribed
finite source order once the cutoff is sufficiently large.
\item The V14 buffer rotation and V13 reduced-graph interpolation therefore give
an unstabilized fixed-index native regulator path at the level of every retained
source jet.  V8--V11 require no stronger source quantifier.
\item Consequently stable and unstabilized regulator universality coincide for
every prescribed finite coefficient/source packet on a protected local Renewal
ESM chart; the V13 massive spectator is not needed for the continuum
coefficients.
\item The required reserve has bounded lattice radius and hence shrinking
physical radius $O(a)$; it is a genuine ultraviolet internal reserve rather
than an added macroscopic degree of freedom.
\end{enumerate}

\subsection*{Inherited}
\begin{enumerate}[leftmargin=2.2em]
\item V1--V7: source-complete coefficientwise continuum, common BV restoration,
fixed-loop forests, ultraviolet Cauchy convergence, scheme/reference covariance
and exact loop-order projectivity.
\item V8--V12: relative regulator forests, fixed-index reduced shell calculus,
global zero-mode bundle/line transport and complex faithful-spin saturation.
\item V13: exponential locality of protected Riesz zero modes, controlled
zero-mode equivalence, local affine graph reconstruction and stable spectator
universality.
\item V14: the exact complement/$\mathcal U_0$ cancellation criterion and the
explicit internal-buffer rotation giving unstabilized native lifting whenever a
buffer is present.
\end{enumerate}

\subsection*{Conditional}
\begin{enumerate}[leftmargin=2.2em]
\item The automatic reserve theorem is coefficientwise on a protected framed
source/background chart.  It does not construct one analytic finite-amplitude
buffer over an arbitrary large source ball.
\item The zero-mode rank $\kappa$, source order $\varrho$, source-panel dimension
$s$ and representation fibre are fixed before the cutoff is enlarged.  The
buffer radius may depend badly on these finite data.
\item The global noncontractible-background problem is reduced to descent of the
local buffer paths and the controlled loop/unitary component; V15 does not claim
a single global microscopic $\mathcal U_0$ path on every component.
\item The V12 real Pfaffian mod-two orientation condition remains independent.
\item All continuum conclusions retain the positive IR/reference scale and no
uniform $L\to\infty$ estimate.
\end{enumerate}

\subsection*{Open beyond V15}
\begin{enumerate}[leftmargin=2.2em]
\item Globalize the local V15 reserve over a noncontractible gauge/background
component.  The overlap of local reserves defines a controlled unitary-loop
Cech cocycle; compute whether this class vanishes on the faithful Standard-Model
native component or identify a genuine coarse/$K_1$ obstruction.
\item Compute the V12 real Pfaffian mod-two spectral-flow character on any
independently physical Majorana sector.
\item Analyze crossings out of the V11 minimal fixed-index stratum, including
accidental zero-mode pair creation and genuine index change, with the required
divisor/spectral-flow source packet.
\item The independent stronger problems remain: positive-IR removal, analytic
finite-coupling RG trajectories, perturbative summability and nonperturbative
completion.
\end{enumerate}

\subsection*{Exact next load-bearing theorem}
\begin{center}
\fbox{\parbox{0.92\textwidth}{
\textbf{Global descent of the coefficientwise internal reserve.}
Cover a connected protected fixed-index gauge/background component by V15
source charts and construct the overlap controlled-unitary cocycle relating
their internal buffers.  Prove that this cocycle is a coboundary in the native
exponentially local unitary group (hence the local unstabilized paths glue to one
global path), or identify its computable controlled $K_1$/coarse-transport
class.  On the faithful complex Standard-Model branch the determinant-line
holonomy is already trivial by V12, so any surviving obstruction is genuinely
a microscopic local-unitary descent obstruction rather than an anomaly phase.}}
\end{center}

\section*{Append-only continuation marker: V15}
\addcontentsline{toc}{section}{Append-only continuation marker: V15}
\textbf{End of V15.}  The complete V1--V14 body above is retained verbatim.
V15 proves automatic coefficientwise internal cancellation for the actual
fixed-rank protected zero-mode packet at sufficiently large cutoff.  It first
extracts a uniformly localized orthonormal frame from the exponentially local
Riesz projector, then uses finitely many bounded native parallel shifts to
build a uniformly conditioned copy bank.  Since only finitely many source jets
are retained, a dimension count produces a rank-$|k|$ common annihilator of all
endpoint overlap jets, yielding an exact V14 internal buffer in the truncated
source algebra.  The buffer lives in the existing carrier, uses only a bounded lattice
displacement from the original localized zero-mode centres, and therefore has
shrinking added physical transport scale; it removes the V13 spectator from all
fixed-order continuum coefficients.  The remaining stronger gate is global
descent of these local buffer choices around noncontractible background loops.


\clearpage
\section{V16 continuation: phenomenology pivot --- physical calibration, predictive quotients, and the native RG generator}
\label{sec:v16-continuation}

V1--V7 established the source-complete coefficientwise perturbative continuum,
and V8--V15 strengthened its regulator and fixed-index realization.  Those
results remain inherited.  The present continuation deliberately changes the
main research priority.  The global descent problem isolated at the end of V15
is retained as a secondary mathematical branch; it is no longer the next
load-bearing problem for phenomenology.

The phenomenological line begins from a different question.  Once a continuum
coefficient system exists, which of its local coordinates are physical inputs,
which are conventional or redundant, which genuinely run, and which relations
among the remaining renormalized coefficients are predictions rather than
fits?  The distinction is essential in gravitational EFT, because existence of
a finite local Wilson bank does not determine the numerical point in that bank.

Fix throughout a finite coefficient packet
\begin{equation}
 \mathfrak q=(L,\Delta,r,N,\mathcal P_{\rm ext})
 \label{eq:v16-q}
\end{equation}
and a positive reference scale \(\mu_0>0\).  The V7 projective continuum gives
a finite coefficient state at this order.  We write its local part abstractly
as a point of a finite-dimensional real manifold
\begin{equation}
 \mathcal C_{\mathfrak q}(\mu_0).
 \label{eq:v16-C}
\end{equation}
The coordinates include the Standard-Model marginal couplings and masses,
Newton/gravitational EFT coordinates, mixed gravity--matter Wilson
coefficients, local composite/contact coordinates, field/source normalization
coordinates and BV-exact/redundant directions admitted by the fixed bank.  No
claim is made that all these coordinates are independently measurable.

\subsection{The new P1--P9 phenomenology programme}
\label{sec:v16-programme}

The main line from V16 onward is:
\begin{enumerate}[label=\textup{(P\arabic*)},leftmargin=3.2em]
\item \textbf{Physical calibration.}  Choose a reference scale \(\mu_0\), a
finite physical input panel and a common field/source normalization, and solve
for the corresponding native local coefficients.
\item \textbf{Native RG vector field.}  Populate the actual beta functions and
operator-mixing/anomalous-dimension matrices of the chosen finite bank.
\item \textbf{Standard-scheme matching.}  Compute the finite conversion from the
native calibrated coordinates to a standard phenomenological scheme such as a
chosen \(\overline{\rm MS}\)/on-shell/EFT basis.
\item \textbf{Predictive quotient.}  Separate input directions, freely
matchable Wilson directions, redundant/scheme directions and genuinely
Renewal-selected relations.
\item \textbf{Threshold matching.}  Cross every retained heavy threshold with a
single common matching action/source prescription.
\item \textbf{Numerical RG trajectory.}  Integrate the calibrated finite-order
RG system with uncertainty transport.
\item \textbf{Benchmark observable panel.}  Select a small set of observables by
expected signal-to-uncertainty rather than by formal elegance.
\item \textbf{Observable map and RG cancellation.}  Compute the matrix elements
and verify cancellation of the reference-scale dependence to the retained
perturbative order.
\item \textbf{Prediction audit.}  For every claimed effect record whether each
coefficient is an experimental input, an unconstrained Wilson coordinate, a
finite scheme convention, or a genuine scheme-invariant consequence of the
Renewal matching law.
\end{enumerate}

The present V16 attacks P1 and P4 rigorously and constructs the exact formal
object needed for P2 and P3.  It does \emph{not} assign uncomputed numerical
beta coefficients.

\subsection{A finite physical calibration panel}
\label{sec:v16-calibration-panel}

Let \(\mathcal A_{\rm cal}=\{\mathcal N_a\}_{a=1}^{n_{\rm in}}\) be a finite
set of physical normalization functionals evaluated at \(\mu_0\).  A
normalization functional is a fixed coefficient of the normalized connected or
1PI source family on a declared kinematic panel.  Typical entries may include
field residues, gauge vertices, a Higgs mass/vev convention, Yukawa vertices,
a gravitational two-point normalization defining Newton's coupling, and
selected higher-derivative anchors.  The theorem below uses only that the
chosen functionals are actual coefficients of the source-complete continuum
state and that their calibration Jacobian is nonsingular.

Choose local coordinates
\begin{equation}
 c=(x,w,z)
 \label{eq:v16-coordinate-split}
\end{equation}
on \(\mathcal C_{\mathfrak q}(\mu_0)\), where:
\begin{enumerate}[label=\textup{(C\arabic*)},leftmargin=3em]
\item \(x\in\mathbb R^{n_{\rm in}}\) are the coordinates to be fixed by the
physical normalization panel;
\item \(w\in\mathbb R^{n_W}\) are independent Wilson/EFT coordinates not fixed
by that panel;
\item \(z\in\mathbb R^{n_R}\) are local redundant, BV-exact, basis or contact
coordinates retained by the source-complete state but not counted as independent
phenomenological inputs.
\end{enumerate}
The split is a local coordinate choice, not an assertion that the three blocks
are canonically defined before a phenomenological quotient is chosen.

Write the normalization map as
\begin{equation}
 \mathcal N_{\mu_0}(x,w,z;\hbar)
 =\sum_{\ell=0}^{L}\hbar^\ell
   \mathcal N^{(\ell)}_{\mu_0}(x,w,z)
 \quad \bmod \hbar^{L+1},
 \label{eq:v16-N-map}
\end{equation}
and let
\begin{equation}
 p=(p_1,\ldots,p_{n_{\rm in}})
 \label{eq:v16-p}
\end{equation}
be the declared numerical physical input vector at \(\mu_0\).

\begin{definition}[Calibrated physical chart]
\label{def:v16-calibrated-chart}
A calibration chart at \((p_0,w_0,z_0)\) is admissible if the tree/reference
Jacobian
\begin{equation}
 J_0:=
 \left.\frac{\partial\mathcal N^{(0)}_{\mu_0}}
 {\partial x}\right|_{(x_0,w_0,z_0)}
 \label{eq:v16-J0}
\end{equation}
is invertible, where
\(\mathcal N^{(0)}_{\mu_0}(x_0,w_0,z_0)=p_0\).
The matrix \(J_0\) is the finite physical calibration analogue of the
constant-rank normalization gates already used in V6--V7.
\end{definition}

\subsection{P1: coefficientwise physical calibration theorem}
\label{sec:v16-P1}

The first point is that no finite-amplitude analytic RG trajectory is needed to
calibrate a fixed perturbative order.

\begin{theorem}[Formal physical calibration at a reference scale]
\label{thm:v16-calibration}
Assume Definition~\ref{def:v16-calibrated-chart}.  Fix \((p,w,z)\) in the local
finite coefficient chart.  There exists a unique coefficientwise series
\begin{equation}
 \boxed{
 x=\sigma_{\mu_0}^{[L]}(p,w,z)
 =x_0(p,w,z)+\sum_{\ell=1}^{L}\hbar^\ell x_\ell(p,w,z)}
 \quad\bmod\hbar^{L+1}
 \label{eq:v16-sigma}
\end{equation}
satisfying
\begin{equation}
 \boxed{
 \mathcal N_{\mu_0}
 \bigl(\sigma_{\mu_0}^{[L]}(p,w,z),w,z;\hbar\bigr)
 =p\quad\bmod\hbar^{L+1}.}
 \label{eq:v16-calibration-eq}
\end{equation}
At each loop order \(\ell\ge1\), \(x_\ell\) is obtained from one linear system
with the same invertible leading Jacobian \(J_0\).  No common nonzero radius in
\(\hbar\) or in the source amplitudes is required.
\end{theorem}

\begin{proof}
At order zero, the ordinary finite-dimensional implicit-function theorem gives
a unique local solution \(x_0=x_0(p,w,z)\) because \(J_0\) is invertible.
Suppose \(x_1,\ldots,x_{\ell-1}\) have been chosen so that
\eqref{eq:v16-calibration-eq} holds through order \(\hbar^{\ell-1}\).  Insert
\[
 x=x_0+\hbar x_1+\cdots+\hbar^\ell x_\ell
\]
into \eqref{eq:v16-N-map}.  The coefficient of \(\hbar^\ell\) has the form
\begin{equation}
 J_0x_\ell+R_\ell
 \label{eq:v16-recursive-cal}
\end{equation}
where \(R_\ell\) is a finite polynomial in the previously fixed
\(x_1,\ldots,x_{\ell-1}\), the derivatives of the lower-loop normalization
functionals, and \(\mathcal N^{(j)}\) with \(j\le\ell\).  The target vector
\(p\) has no positive powers of \(\hbar\), so the order-\(\ell\) condition is
\(J_0x_\ell=-R_\ell\).  Invertibility of \(J_0\) gives the unique solution.
Induction proves existence and uniqueness through \(L\).  The construction is
purely coefficientwise and therefore requires no convergence of the resulting
formal series.
\end{proof}

\begin{corollary}[Exact perturbative-order compatibility of calibration]
\label{cor:v16-calibration-projective}
If the normalization panel and coordinate ordering are chosen chronologically
as in V7, then
\begin{equation}
 \boxed{
 \pi_L^{L+1}\sigma_{\mu_0}^{[L+1]}
 =\sigma_{\mu_0}^{[L]}.}
 \label{eq:v16-calibration-projective}
\end{equation}
Thus increasing the terminal loop order never refits a previously calibrated
lower-order physical parameter.
\end{corollary}

\begin{proof}
The recursion in the proof of Theorem~\ref{thm:v16-calibration} determines
\(x_\ell\) from data of order at most \(\ell\).  Appending the
\((L+1)\)-st equation does not alter any earlier right-hand side.  Therefore
the coefficients through order \(L\) are identical.
\end{proof}

\begin{definition}[Native-to-physical calibration map]
\label{def:v16-native-physical-map}
On the admissible chart define
\begin{equation}
 \boxed{
 \Phi_{\mu_0}^{\rm phys}:
 c=(x,w,z)
 \longmapsto
 \bigl(p=\mathcal N_{\mu_0}(c),w,z\bigr).}
 \label{eq:v16-native-physical-map}
\end{equation}
Its inverse is
\begin{equation}
 (p,w,z)\longmapsto
 \bigl(\sigma_{\mu_0}^{[L]}(p,w,z),w,z\bigr)
 \label{eq:v16-native-physical-inverse}
\end{equation}
coefficientwise through order \(L\).
\end{definition}

This is the precise version of the first phenomenological task: the native
Renewal coefficients become recognizable physical inputs only after a declared
normalization panel has been supplied and its finite Jacobian has been checked.

\subsection{Field/source normalization is part of the same map}
\label{sec:v16-wavefunction}

Wave-function and source normalizations must not be performed after the fact as
independent graphwise choices.  Let \(Z\) denote the finite block of allowed
linear field/source transformations on the chosen physical quotient, including
flavour matrices where required.  Let \(\mathcal R_a(c,Z)\) be the residue and
vertex normalization conditions used to fix them.

\begin{proposition}[Simultaneous coupling and field/source calibration]
\label{prop:v16-Z-calibration}
Suppose the combined leading Jacobian
\begin{equation}
 \frac{\partial(\mathcal N^{(0)},\mathcal R^{(0)})}
      {\partial(x,Z)}
 \label{eq:v16-combined-J}
\end{equation}
is invertible at the reference point after quotienting exact gauge/BV
redundancies.  Then Theorem~\ref{thm:v16-calibration} applies verbatim to the
combined variable \((x,Z)\).  Consequently all calibrated couplings, residues,
source normalizations and their contact descendants arise from one common
finite coefficient map.
\end{proposition}

\begin{proof}
Replace \(x\) by the finite coordinate vector containing both the coupling and
field/source-normalization entries and replace \(\mathcal N\) by
\((\mathcal N,\mathcal R)\).  The recursive proof uses only invertibility of
the common leading Jacobian, so it is unchanged.
\end{proof}

\begin{firewall}[Calibration is not prediction]
Theorem~\ref{thm:v16-calibration} tells us how to express native coefficients
in terms of a declared physical input vector and whatever Wilson coordinates
remain free.  It does not determine those unfixed Wilson coordinates and does
not turn a fitted EFT coefficient into a prediction.
\end{firewall}

\subsection{P4 first: the phenomenological quotient}
\label{sec:v16-predictive-quotient}

Before computing a large beta table, it is useful to identify the space on
which a genuinely predictive relation could live.

Let \(\mathcal G_{\rm red}\) denote the finite groupoid generated, at the fixed
coefficient order, by admissible local field redefinitions, BV-exact
coordinate changes, composite/source basis changes and finite scheme
transformations which preserve the matched physical observable algebra.  It
acts locally on \(\mathcal C_{\mathfrak q}\).  We write
\begin{equation}
 \mathcal M_{\mathfrak q}^{\rm phys}
 :=\mathcal C_{\mathfrak q}/\mathcal G_{\rm red}
 \label{eq:v16-physical-moduli}
\end{equation}
for the local phenomenological quotient.  This quotient is used only on a
constant-rank chart; no global orbifold statement is required.

The calibrated physical input map descends to
\begin{equation}
 I_{\mu_0}:\mathcal M_{\mathfrak q}^{\rm phys}
 \longrightarrow\mathcal P_{\rm in},
 \label{eq:v16-input-map}
\end{equation}
and the calibrated fibre at the measured input \(p\) is
\begin{equation}
 \mathcal F_p:=I_{\mu_0}^{-1}(p).
 \label{eq:v16-calibrated-fibre}
\end{equation}
Its coordinates are precisely the physical Wilson/EFT information left after
calibration and redundancy removal.

\begin{definition}[Renewal matching locus and predictive relation]
\label{def:v16-prediction}
Let \(\Theta\) denote the set of genuinely free microscopic matching parameters
which remain after the experimental inputs \(p\) are fixed.  A
\emph{Renewal matching law} is a map
\begin{equation}
 \mathfrak M_{\mu_0}:\Theta\longrightarrow\mathcal F_p.
 \label{eq:v16-Renewal-matching}
\end{equation}
A smooth scalar relation \(R\) on \(\mathcal F_p\) is a
\emph{local predictive relation} if
\begin{equation}
 R\circ\mathfrak M_{\mu_0}=0
 \label{eq:v16-predictive-relation}
\end{equation}
and \(R\) is a scalar under all admissible finite scheme/basis transitions.
It is nontrivial if its differential is nonzero on the matching locus.
\end{definition}

\begin{theorem}[Local count of independent predictions]
\label{thm:v16-predictive-codimension}
Assume \(\mathfrak M_{\mu_0}\) has constant rank \(r_M\) near
\(\theta_0\), and put
\begin{equation}
 d_F:=\dim\mathcal F_p.
 \label{eq:v16-dF}
\end{equation}
Then the space of independent first-order predictive covectors at
\(c_0=\mathfrak M_{\mu_0}(\theta_0)\) has dimension
\begin{equation}
 \boxed{d_F-r_M.}
 \label{eq:v16-predictive-count}
\end{equation}
More precisely it is the conormal space
\begin{equation}
 \boxed{
 \mathcal P_{c_0}^*
 :=\{\ell\in T_{c_0}^*\mathcal F_p:
       \ell\circ D\mathfrak M_{\mu_0}|_{\theta_0}=0\}.}
 \label{eq:v16-predictive-conormal}
\end{equation}
This dimension and conormal subspace transform covariantly under every
invertible finite scheme change.
\end{theorem}

\begin{proof}
The constant-rank theorem makes the local matching locus an immersed
\(r_M\)-dimensional submanifold of the \(d_F\)-dimensional calibrated fibre.
Its conormal space is the annihilator of the tangent image of
\(D\mathfrak M\), hence has dimension \(d_F-r_M\).  If
\(F:\mathcal F_p\to\mathcal F'_p\) is an invertible scheme transition, the new
matching map is \(F\circ\mathfrak M\).  Its tangent image is
\(DF(T\operatorname{im}\mathfrak M)\); covectors transform by
\((DF^{-1})^*\).  Therefore the conormal dimension and the vanishing of a
covariant relation are invariant.
\end{proof}

\begin{corollary}[Matrix test for infinitesimal predictivity]
\label{cor:v16-prediction-matrix}
Choose local Wilson coordinates \(W^I\) on \(\mathcal F_p\) and microscopic
parameters \(\theta^a\).  Let
\begin{equation}
 S^I{}_a
 :=\frac{\partial W^I}{\partial\theta^a}
 \label{eq:v16-selection-matrix}
\end{equation}
be the matching sensitivity matrix.  A linearized coefficient combination
\begin{equation}
 \ell_I\,\delta W^I
 \label{eq:v16-linear-prediction}
\end{equation}
is predictive precisely when
\begin{equation}
 \boxed{\ell_I S^I{}_a=0\quad\text{for every }a.}
 \label{eq:v16-left-null}
\end{equation}
Thus the left nullspace of the populated matching sensitivity matrix is the
first object to inspect for candidate phenomenological relations.
\end{corollary}

\begin{proof}
Equation \eqref{eq:v16-left-null} is exactly the coordinate form of
\eqref{eq:v16-predictive-conormal}.
\end{proof}

\subsection{A no-prediction theorem for continuum existence alone}
\label{sec:v16-no-prediction}

The preceding definition makes it possible to state sharply what has and has
not yet been obtained from the continuum theorem itself.

\begin{theorem}[Continuum existence alone does not numerically select free Wilson coordinates]
\label{thm:v16-no-prediction}
Suppose that after physical calibration the proved continuum theorem retains an
open set \(U\subset\mathbb R^{n_W}\) of independent Wilson coordinates \(w\)
as admissible coordinates of the same fixed-order state, with no additional
microscopic matching equation selecting a proper subset of \(U\).  Then no
nonzero smooth relation
\begin{equation}
 R(w)=0
 \label{eq:v16-fake-relation}
\end{equation}
can be inferred from continuum existence alone on that open set.  In
particular, if the current theorem allows \(w\) to vary independently in \(U\),
its predictive conormal space in the \(w\)-directions is zero.
\end{theorem}

\begin{proof}
A relation logically implied by the theorem must hold for every admissible
point in \(U\).  Hence \(R\) vanishes identically on a nonempty open set.
Therefore every derivative of \(R\) on that set vanishes, so it supplies no
nontrivial local constraint.  Equivalently the matching locus is the whole
open Wilson fibre and its tangent map has full rank \(n_W\); Theorem
\ref{thm:v16-predictive-codimension} gives zero conormal dimension.
\end{proof}

\begin{corollary}[What must be added to obtain new phenomenology]
\label{cor:v16-need-selection}
To obtain a new numerical relation among independent EFT/Wilson coefficients,
one must add a microscopic or dynamical selection law
\(\mathfrak M_{\mu_0}\) whose image has positive codimension in the calibrated
Wilson fibre, or compute an observable whose value is independent of the
remaining free directions.  Renormalizability/continuum existence by itself is
not such a selection law.
\end{corollary}

This corollary is the operational reason P4 is moved ahead of expensive P2
computations: a beta function transports whatever boundary data have been
chosen, but it does not by itself reduce the number of independent boundary
coordinates.

\subsection{P2: the RG vector field from reference-scale transport}
\label{sec:v16-RG}

V7 already supplied finite reference-scale covariance on compact positive
scale intervals.  We now reinterpret that family as the physical RG transport.
Let
\begin{equation}
 \mathcal U(\mu_2,\mu_1):
 \mathcal M_{\mathfrak q}^{\rm phys}(\mu_1)
 \longrightarrow
 \mathcal M_{\mathfrak q}^{\rm phys}(\mu_2)
 \label{eq:v16-RG-transport}
\end{equation}
be the continuum finite transition obtained by holding the underlying
coefficientwise continuum state fixed while recalibrating it at the new
positive reference scale.  It obeys
\begin{equation}
 \mathcal U(\mu_3,\mu_2)\mathcal U(\mu_2,\mu_1)
 =\mathcal U(\mu_3,\mu_1),
 \qquad
 \mathcal U(\mu,\mu)=I.
 \label{eq:v16-RG-cocycle}
\end{equation}

Choose dimensionless physical coordinates \(g^A\) on the calibrated quotient.
For an operator \(\mathcal O_I\) of canonical dimension \(d_I\), one may for
example replace a dimensionful coefficient \(c_I\) by
\begin{equation}
 \bar c_I(\mu):=\mu^{d_I-4}c_I(\mu),
 \label{eq:v16-dimensionless-Wilson}
\end{equation}
with analogous conventional definitions for masses and Newton's coupling.  The
choice of dimensionless coordinates is part of the phenomenological scheme.

\begin{definition}[Native beta vector field]
\label{def:v16-beta}
Assume \(\mathcal U\) is \(C^1\) in \(t=\log\mu\) on the chosen positive
reference interval.  Define
\begin{equation}
 \boxed{
 \beta^A_{\rm nat}(g,\mu)
 :=\left.
 \frac{d}{ds}
 \Bigl[\mathcal U(e^s\mu,\mu)g\Bigr]^A
 \right|_{s=0}.}
 \label{eq:v16-beta}
\end{equation}
This is a definition of the actual native RG vector field; it is not a
numerical value for its coefficients.
\end{definition}

\begin{theorem}[Cocycle implies the RG differential equation]
\label{thm:v16-RG-ODE}
Let
\begin{equation}
 g(t):=\mathcal U(e^t\mu_0,\mu_0)g_0.
 \label{eq:v16-running-g}
\end{equation}
Then
\begin{equation}
 \boxed{
 \frac{dg^A}{dt}
 =\beta^A_{\rm nat}(g(t),e^t\mu_0).}
 \label{eq:v16-RG-ODE}
\end{equation}
If the chosen dimensionless scheme is mass independent on the interval and
there is no explicit reference-scale dependence apart from the running
coordinates, \(\beta_{\rm nat}\) is autonomous.
\end{theorem}

\begin{proof}
For small \(s\), the cocycle gives
\[
 g(t+s)
 =\mathcal U(e^{t+s}\mu_0,e^t\mu_0)g(t).
\]
Subtract \(g(t)\), divide by \(s\), and let \(s\to0\).  Definition
\ref{def:v16-beta} at the reference scale \(e^t\mu_0\) gives
\eqref{eq:v16-RG-ODE}.  In a mass-independent dimensionless scheme, scale
covariance identifies transports at equal scale ratios, removing explicit
\(\mu\)-dependence from the generator.
\end{proof}

\begin{proposition}[Canonical scaling plus loop contribution]
\label{prop:v16-canonical-beta}
For a dimensionless Wilson coordinate
\(\bar c_I=\mu^{d_I-4}c_I\), the beta function decomposes as
\begin{equation}
 \boxed{
 \beta_{\bar c_I}
 =(d_I-4)\bar c_I+\beta_{I}^{\rm loop},}
 \label{eq:v16-canonical-beta}
\end{equation}
where \(\beta_I^{\rm loop}\) is the logarithmic scale derivative of the
renormalized dimensionful coefficient at fixed canonical units.  The first
term is kinematic and the second requires an actual loop calculation.
\end{proposition}

\begin{proof}
Differentiate
\(\bar c_I=\mu^{d_I-4}c_I\) with respect to \(\log\mu\).  The derivative of
the explicit power gives \((d_I-4)\bar c_I\); the remaining term is
\(\mu^{d_I-4}dc_I/d\log\mu\), which is the loop/matching contribution by
definition.
\end{proof}

\subsection{The beta function is a vector, not a table of scheme-independent numbers}
\label{sec:v16-beta-scheme}

Let \(g'=F(g,\mu)\) be an invertible finite phenomenological scheme change,
including a finite operator-basis transformation where appropriate.

\begin{theorem}[Finite-scheme transformation law for the beta vector]
\label{thm:v16-beta-transform}
The beta vector fields in the two schemes obey
\begin{equation}
 \boxed{
 \beta'(g',\mu)
 =D_gF(g,\mu)\,\beta(g,\mu)
  +\mu\,\partial_\mu F(g,\mu),
 \qquad g=F^{-1}(g',\mu).}
 \label{eq:v16-beta-transform}
\end{equation}
For a scale-independent finite redefinition,
\begin{equation}
 \beta'=F_*\beta.
 \label{eq:v16-beta-pushforward}
\end{equation}
Hence the RG flow is coordinate covariant even though individual higher-loop
beta coefficients need not be scheme invariant.
\end{theorem}

\begin{proof}
Along an RG trajectory,
\[
 \frac{dg'}{d\log\mu}
 =D_gF\frac{dg}{d\log\mu}
  +\mu\partial_\mu F.
\]
Insert \(dg/d\log\mu=\beta(g,\mu)\) and express the result in the primed
coordinates.
\end{proof}

\subsection{Composite/source mixing as an RG connection}
\label{sec:v16-operator-mixing}

Let \(\mathbf O(\mu)\) be a finite column of renormalized composite operators in
the retained phenomenological bank and suppose scale transport is
\begin{equation}
 \mathbf O(\mu_2)=
 R(\mu_2,\mu_1;g)\,\mathbf O(\mu_1).
 \label{eq:v16-operator-transport}
\end{equation}
The source vector transforms contragrediently so that
\(J^T\mathbf O\) is unchanged.

\begin{definition}[Anomalous-dimension/mixing matrix]
\label{def:v16-gamma}
Define
\begin{equation}
 \boxed{
 \Gamma(g,\mu)
 :=\left.
 \frac{d}{ds}R(e^s\mu,\mu;g)
 \right|_{s=0}.}
 \label{eq:v16-gamma}
\end{equation}
Then
\begin{equation}
 \frac{d\mathbf O}{d\log\mu}=\Gamma\mathbf O,
 \qquad
 \frac{dJ^T}{d\log\mu}=-J^T\Gamma.
 \label{eq:v16-operator-source-RG}
\end{equation}
\end{definition}

\begin{proposition}[Operator-basis covariance of the mixing matrix]
\label{prop:v16-gamma-transform}
Under a finite scale-dependent basis change
\begin{equation}
 \mathbf O'=B(g,\mu)\mathbf O,
 \label{eq:v16-O-basis}
\end{equation}
one has along the RG vector field
\begin{equation}
 \boxed{
 \Gamma'
 =B\Gamma B^{-1}
  +\bigl[(\mu\partial_\mu+\beta^A\partial_A)B\bigr]B^{-1}.}
 \label{eq:v16-gamma-transform}
\end{equation}
Thus \(\Gamma\) is a connection on the finite operator/source bundle, not a
basis-independent matrix of numbers.
\end{proposition}

\begin{proof}
Differentiate \(\mathbf O'=B\mathbf O\) along an RG trajectory:
\[
 \frac{d\mathbf O'}{d\log\mu}
 =\bigl[(\mu\partial_\mu+\beta\cdot\partial)B\bigr]\mathbf O
   +B\Gamma\mathbf O.
\]
Replace \(\mathbf O=B^{-1}\mathbf O'\) to obtain
\eqref{eq:v16-gamma-transform}.
\end{proof}

\subsection{Predictive relations must themselves be RG transported}
\label{sec:v16-prediction-RG}

A relation imposed only at \(\mu_0\) is phenomenologically meaningful at other
scales only if its matching locus is transported by the RG flow.

\begin{theorem}[RG transport of a Renewal prediction]
\label{thm:v16-prediction-RG}
Suppose the matching locus satisfies
\begin{equation}
 \mathfrak M_{\mu}
 =\mathcal U(\mu,\mu_0)\circ\mathfrak M_{\mu_0}
 \label{eq:v16-M-RG}
\end{equation}
and let \(R(g,\mu)\) be a smooth scalar relation satisfying
\(R=0\) on \(\operatorname{im}\mathfrak M_\mu\) for every \(\mu\).  Then on the
matching locus
\begin{equation}
 \boxed{
 (\mu\partial_\mu+\beta^A\partial_A)R=0.}
 \label{eq:v16-CS-relation}
\end{equation}
Conversely, a relation satisfying \eqref{eq:v16-CS-relation} and vanishing on
the matching locus at \(\mu_0\) vanishes on its RG transport for every scale in
the common interval.
\end{theorem}

\begin{proof}
For a fixed microscopic parameter \(\theta\), define
\(g(\mu)=\mathfrak M_\mu(\theta)\).  By
\eqref{eq:v16-M-RG}, \(g\) obeys Theorem~\ref{thm:v16-RG-ODE}.  Since
\(R(g(\mu),\mu)=0\), differentiating in \(\log\mu\) gives
\eqref{eq:v16-CS-relation}.  Conversely the same differential equation says
that \(R\) is constant along the transported trajectory; zero initial value
therefore remains zero.
\end{proof}

This theorem is the formal bridge between P4 and P2.  A candidate predictive
coefficient relation is not merely a numerical identity at the calibration
scale; it defines an RG-invariant submanifold or else it must be supplemented
by its scale-dependent completion.

\subsection{P3: coefficientwise matching to a standard phenomenological scheme}
\label{sec:v16-standard-matching}

The native coordinates are not themselves the final comparison language.  Let
\(c^{\rm nat}\) and \(c^{\rm std}\) be coefficient vectors in the native and a
chosen standard phenomenological scheme at the same \(\mu_0\).  Pick a finite
matching panel \(\mathcal M\) of renormalized amplitudes/source coefficients
large enough to fix the desired conversion coordinates and impose
\begin{equation}
 \mathcal M^{\rm nat}(c^{\rm nat})
 =\mathcal M^{\rm std}(c^{\rm std}).
 \label{eq:v16-standard-match-eq}
\end{equation}

\begin{theorem}[Formal finite native-to-standard matching map]
\label{thm:v16-standard-match}
Assume that at tree/reference order the standard matching Jacobian
\begin{equation}
 J_{\rm std}:=
 \frac{\partial\mathcal M^{\rm std,(0)}}
      {\partial c^{\rm std}}
 \label{eq:v16-Jstd}
\end{equation}
is invertible on the finite coefficient block being matched.  Then there is a
unique coefficientwise finite map
\begin{equation}
 \boxed{
 c^{\rm std}=F_{\rm std\leftarrow nat}^{[L]}
 (c^{\rm nat})}
 \quad\bmod\hbar^{L+1}
 \label{eq:v16-standard-map}
\end{equation}
solving \eqref{eq:v16-standard-match-eq}.  If the tree coordinates have been
identified, the one-loop shift is
\begin{equation}
 \boxed{
 \delta c^{\rm std}_{(1)}
 =J_{\rm std}^{-1}
  \left(
   \mathcal M^{\rm nat,(1)}
   -\mathcal M^{\rm std,(1)}
  \right),}
 \label{eq:v16-one-loop-match}
\end{equation}
with all quantities evaluated at the common tree point.  Higher orders follow
recursively from the same leading Jacobian.
\end{theorem}

\begin{proof}
The proof is the same coefficientwise implicit recursion as
Theorem~\ref{thm:v16-calibration}.  For the one-loop formula, write
\(c^{\rm std}=c^{(0)}+\hbar\delta c_{(1)}^{\rm std}\).  Expanding the standard
side of \eqref{eq:v16-standard-match-eq} gives
\[
 \mathcal M^{\rm std,(0)}(c^{(0)})
 +\hbar\left(
   J_{\rm std}\delta c_{(1)}^{\rm std}
   +\mathcal M^{\rm std,(1)}(c^{(0)})
 \right).
\]
Equating it with the native expansion and using the common tree value gives
\eqref{eq:v16-one-loop-match}.  At higher order the new coefficient appears
linearly through the same \(J_{\rm std}\), while all nonlinear terms are known
from lower orders.
\end{proof}

\begin{firewall}[The conversion theorem is not the conversion table]
Theorem~\ref{thm:v16-standard-match} proves the existence and recursive form of
the finite matching map once a concrete common amplitude panel is populated.
V16 has not evaluated the native loop amplitudes needed for an actual
\(\overline{\rm MS}\), on-shell or SMEFT conversion table.
\end{firewall}

\subsection{Phenomenological prediction ledger}
\label{sec:v16-ledger-definition}

For every candidate observable or coefficient relation, define four Boolean
labels:
\begin{equation}
 \bigl(
 \mathrm{Input},\mathrm{FreeWilson},\mathrm{RenewalSelected},
 \mathrm{SchemeInvariant}
 \bigr).
 \label{eq:v16-ledger-labels}
\end{equation}
A quantity is called a new prediction only if it is not an input, is not an
independent free Wilson coordinate, lies on a populated Renewal matching locus
of positive codimension, and is expressed as a scheme-covariant observable or
relation.

\begin{proposition}[No double counting between calibration and prediction]
\label{prop:v16-no-double-count}
Let \(Q\) be one of the coordinates used as an independent component of the
input vector \(p\).  Then no relation whose only nonzero differential is in the
\(Q\)-direction is an independent prediction of the calibrated theory.
Likewise, a Wilson coefficient varied freely in \(\Theta\) cannot be counted as
a prediction merely because it is subsequently run by the RG equation.
\end{proposition}

\begin{proof}
Fixing \(Q\) as an input removes its tangent direction from
\(T\mathcal F_p\), so its dual does not belong to the predictive conormal
\eqref{eq:v16-predictive-conormal}.  A freely varied Wilson coefficient belongs
to the image of \(D\mathfrak M\); every predictive covector annihilates that
image by definition.  RG transport is a local diffeomorphism and therefore
cannot change this dimension count.
\end{proof}

\subsection{What V16 has and has not populated}
\label{sec:v16-population-boundary}

The structural work above permits the next numerical programme to be specified
without ambiguity.

\begin{definition}[Minimal P2 coefficient block]
\label{def:v16-P2-block}
The first populated RG packet should contain:
\begin{enumerate}[label=\textup{(R\arabic*)},leftmargin=3em]
\item the calibrated Standard-Model marginal block
\(g_1,g_2,g_3,Y_u,Y_d,Y_e,\lambda_H\), together with the field/source
anomalous dimensions needed to define them;
\item the chosen Higgs mass/vev convention and its mass-dependent running;
\item Newton's coupling in the declared dimensionless coordinate and the
independent gravitational/mixed EFT coordinates through the first derivative
order relevant to the selected observable panel;
\item the finite composite/source mixing matrix for the operators entering that
panel.
\end{enumerate}
Every coordinate must be defined by the P1 calibration map before its beta
coefficient is quoted.
\end{definition}

The old V80 line already carried physical gauge couplings and independent
Wilson coordinates as state variables while explicitly leaving their detailed
beta functions uncomputed.  V16 does not retrospectively assign values to
those missing rows.  They are now the next primary calculation.

\begin{theorem}[Logarithmic-shell extraction principle]
\label{thm:v16-log-shell}
Suppose a calibrated native shell calculation for one dimensionless coordinate
\(g^A\), with scale ratio \(b>1\), has local increment
\begin{equation}
 g^A(b\mu)-g^A(\mu)
 =\sum_{\ell=1}^{L}\hbar^\ell
  \left[
   B^{A}_{\ell}(g)\log b
   +R^{A}_{\ell}(g,b)
  \right]
 \label{eq:v16-shell-increment}
\end{equation}
where
\begin{equation}
 R^{A}_{\ell}(g,b)=o(\log b)
 \quad\text{as }b\downarrow1
 \label{eq:v16-shell-remainder}
\end{equation}
uniformly on the calibrated compact chart.  Then
\begin{equation}
 \boxed{
 \beta^A_{\rm nat}(g)
 =\sum_{\ell=1}^{L}\hbar^\ell B^A_\ell(g)
 \quad\bmod\hbar^{L+1}}
 \label{eq:v16-beta-log}
\end{equation}
plus the explicit canonical-scaling contribution when \(g^A\) was made
dimensionless by a power of \(\mu\).
\end{theorem}

\begin{proof}
Set \(b=e^s\).  Divide \eqref{eq:v16-shell-increment} by \(s=\log b\) and let
\(s\downarrow0\).  The remainder vanishes by
\eqref{eq:v16-shell-remainder}, while the coefficient of \(\log b\) survives.
Definition~\ref{def:v16-beta} gives \eqref{eq:v16-beta-log}.  Proposition
\ref{prop:v16-canonical-beta} supplies any explicit canonical term.
\end{proof}

This theorem identifies exactly what the next loop calculations must return:
the local logarithmic coefficient in the \emph{already calibrated common
action/source scheme}, not a graphwise divergence before wave-function and
operator mixing have been assembled.

\subsection{V16 anti-circularity audit}
\label{sec:v16-audit}

\begin{enumerate}[leftmargin=2.4em]
\item \textbf{The phenomenology pivot does not revoke the continuum theorem.}
V1--V15 remain inherited; their stronger regulator questions are simply moved
off the shortest path to observables.
\item \textbf{A physical coupling is not identified with the formal loop
parameter.}
The calibrated coordinate is defined by an actual normalization functional at
\(\mu_0\).
\item \textbf{Wave-function, source and coupling normalization are solved
simultaneously.}
No graph receives its own post-hoc finite rescaling.
\item \textbf{Continuum existence is not advertised as a numerical
prediction.}
Theorem~\ref{thm:v16-no-prediction} proves that an independently retained open
Wilson fibre carries no nontrivial numerical relation until a matching law
selects a proper submanifold.
\item \textbf{The beta vector is not treated as a scheme-invariant table.}
Its exact finite-scheme transformation law is
\eqref{eq:v16-beta-transform}; operator mixing is a connection with
\eqref{eq:v16-gamma-transform}.
\item \textbf{No uncomputed beta coefficient is filled from memory or from a
continuum reference theory.}
V16 proves the extraction theorem and identifies the packet which must be
calculated natively.
\item \textbf{A prediction at \(\mu_0\) must survive RG transport.}
The matching locus and its relations obey
Theorem~\ref{thm:v16-prediction-RG}.
\item \textbf{Native-to-standard matching is not equated with equality of raw
coefficient arrays.}
It is fixed by equality of a common physical amplitude/source panel and solved
coefficientwise by Theorem~\ref{thm:v16-standard-match}.
\end{enumerate}

\section{V16 status ledger}
\label{sec:v16-status}

\subsection*{Proved in V16}
\begin{enumerate}[leftmargin=2.2em]
\item At every prescribed finite loop/EFT/source order, any declared physical
normalization panel with nonsingular leading calibration Jacobian determines a
unique coefficientwise native-to-physical calibration map at the reference
scale \(\mu_0\).  Couplings, field residues and source normalizations can be
solved in one common finite system, and the construction is exactly projective
in perturbative order.
\item The calibrated phenomenological coefficient space is the quotient by
admissible finite scheme, field-redefinition, BV-exact and operator/source-basis
coordinates.  After fixing the experimental inputs, new predictions live on
the conormal to the populated Renewal matching locus inside the calibrated
Wilson fibre.
\item The number of independent first-order predictive relations is the
codimension of that matching locus, and the left nullspace of its sensitivity
matrix is a practical coordinate representation of the predictive conormal.
This object is covariant under finite scheme changes.
\item A continuum theorem which retains an open set of independent Wilson
coordinates but supplies no additional microscopic selection law implies no
nontrivial numerical relation among those coordinates.  RG evolution alone
does not change this conclusion.
\item The V7 reference-scale transport defines a native beta vector field.  Its
cocycle gives the RG ODE, dimensionless Wilson coefficients split into canonical
and loop running, and finite scheme changes transform the beta vector by
\eqref{eq:v16-beta-transform}.
\item Composite/source operator mixing is naturally an RG connection.  Its
basis transformation law is \eqref{eq:v16-gamma-transform}, with sources
running contragrediently.
\item Scheme-invariant predictive relations are transported by the RG flow and
obey the coefficientwise Callan--Symanzik equation
\eqref{eq:v16-CS-relation} on the matching locus.
\item A native-to-standard phenomenological conversion exists coefficientwise
once a common matching panel with nonsingular standard Jacobian is populated;
the explicit one-loop finite shift is \eqref{eq:v16-one-loop-match}.
\item The local logarithmic coefficient of a calibrated infinitesimal shell is
exactly the corresponding native beta coefficient, as stated in
Theorem~\ref{thm:v16-log-shell}.
\end{enumerate}

\subsection*{Inherited}
\begin{enumerate}[leftmargin=2.2em]
\item V1--V7: the source-complete coefficientwise Einstein--Standard Model
continuum, common BV/master restoration, arbitrary-fixed-loop forest theorem,
UV Cauchy limit, scheme/reference covariance and exact loop-order projectivity.
\item V8--V15: stronger regulator/fixed-index universality results, including
local unstabilized cancellation of the protected finite-rank zero-mode packet
at every fixed coefficient/source order.
\item The earlier bosonic two-loop construction's explicit distinction between
physical normalization conditions, independently retained EFT/Wilson
coordinates and as-yet-uncomputed detailed beta functions.
\end{enumerate}

\subsection*{Conditional}
\begin{enumerate}[leftmargin=2.2em]
\item P1 becomes a numerical calibration only after the actual physical input
panel, kinematic normalization points and measured values are declared.  V16
proves the coefficientwise inversion theorem but does not choose experimental
numbers.
\item The predictive quotient becomes nontrivial only after a microscopic
Renewal matching law \(\mathfrak M_{\mu_0}\) is populated.  At present the
continuum theorem alone leaves the independent Wilson coordinates free.
\item Definition~\ref{def:v16-beta} assumes differentiable positive-scale
reference transport.  The coefficientwise finite maps exist from V7; the
actual beta coefficients still require the logarithmic native shell
calculation of Theorem~\ref{thm:v16-log-shell}.
\item The standard-scheme conversion theorem requires an explicit common
matching panel and loop amplitudes in both schemes.  No numerical
\(\overline{\rm MS}\)/on-shell conversion table is claimed here.
\item Gravity remains an EFT: higher-derivative Wilson coefficients are not
removed merely because a physical calibration of the lower-weight coordinates
has been imposed.
\end{enumerate}

\subsection*{Open --- primary phenomenology line}
\begin{enumerate}[leftmargin=2.2em]
\item Populate the P1 input panel explicitly: choose \(\mu_0\), the precise
renormalized gauge/Higgs/Yukawa/gravity input observables and the independent
covariant gravitational/mixed Wilson basis after topological, field-redefinition
and equation-of-motion identifications appropriate to the intended observable
panel.
\item Compute the P2 native logarithmic shell coefficients for the calibrated
Standard-Model marginal block and the selected gravitational/mixed EFT bank,
including the complete field and composite/source mixing matrices.  This is the
next major calculation; no graphwise beta coefficient is accepted before the
common calibration and mixing rows are assembled.
\item Populate the P3 native-to-standard finite matching map on one common
amplitude/source panel.
\item Construct the P4 microscopic matching sensitivity matrix
\(S^I{}_a\) and identify its left-nullspace.  Only nonzero conormal directions
are candidates for new Renewal predictions.
\item After the predictive directions are known, perform P5 threshold matching,
P6 numerical RG integration with uncertainty transport, and P7--P9 observable
selection, matrix-element evaluation and prediction audit.
\end{enumerate}

\subsection*{Secondary mathematical branch}
The V15 global descent problem, real Pfaffian mod-two transport, index-changing
crossings, positive-IR removal, perturbative summability and nonperturbative
completion remain mathematically open but are no longer the next gates for the
phenomenology programme.

\subsection*{Exact next load-bearing theorem}
\begin{center}
\fbox{\parbox{0.92\textwidth}{
\textbf{P2a: populated native one-loop RG packet in calibrated physical
coordinates.}
Fix an explicit P1 reference scale and normalization panel.  In one common
covariant EFT/operator basis, compute the logarithmic shell coefficient for the
calibrated Standard-Model marginal coordinates
\(g_1,g_2,g_3,Y_u,Y_d,Y_e,\lambda_H\), the required field/source anomalous
dimensions, and the first phenomenologically selected gravity/mixed Wilson
block.  Assemble the result only after all wave-function, BV/source/contact and
operator-mixing contributions are included.  Prove that the resulting vector
and mixing matrix satisfy the native Ward/BV identities and the scheme
transformation laws of V16.  The output is the first machine-readable
\(\{\beta^A,\Gamma^I{}_J\}\) table; numerical RG integration should not begin
before this common packet is populated.}}
\end{center}

\section*{Append-only continuation marker: V16}
\addcontentsline{toc}{section}{Append-only continuation marker: V16}
\textbf{End of V16.}  The complete V1--V15 body above is retained verbatim.
V16 deliberately pivots the main line from stronger regulator-universality
questions to phenomenology.  It proves the coefficientwise physical
calibration theorem at a reference scale, including simultaneous field/source
normalization; constructs the calibrated phenomenological quotient and a
scheme-covariant predictive conormal; proves that continuum existence alone
cannot select independently retained Wilson coefficients; defines the native
RG vector field from positive reference-scale transport and derives its scheme
and operator-mixing covariance; proves RG transport of predictive relations;
and gives the recursive native-to-standard matching theorem.  The next primary
gate is no longer global buffer descent.  It is the explicit one-loop
logarithmic shell calculation which populates the first calibrated native beta
and anomalous-dimension table.


\clearpage
\section{V17 continuation: the first calibrated one-loop RG packet --- universal Standard-Model residues and the Higgs--curvature block}
\label{sec:v17-continuation}

V16 changed the main programme from stronger continuum/regulator theorems to
physical calibration, RG population, matching, predictive quotients and
observables.  The first numerical-looking task is therefore to populate a beta
vector.  There is, however, one upstream distinction that must be fixed before
quoting familiar Standard-Model coefficients.  One-loop beta coefficients are
unchanged by a finite \emph{scale-independent} one-loop coordinate change, but
a mass-dependent physical scheme can contain explicit ratios such as
$m/\mu$ and may shift its one-loop beta vector by an explicit
$\mu\partial_\mu$ term.  Thus a standard $\overline{\rm MS}$ table can be
identified with the native one-loop table only after the calibration window is
chosen so that the finite native-to-standard matching has no explicit scale.

V17 makes precisely that choice for the first packet.  We work on a flat
reference geometry, to zeroth order in the dimensionless quantum-gravity
coupling $G_N\mu^2$, in the unbroken Standard Model at nonexceptional Euclidean
external momenta of the form
\begin{equation}
 p_i=\mu\,\widehat p_i,
 \qquad
 \widehat p_i\cdot\widehat p_j\ \hbox{fixed},
 \label{eq:v17-nonexceptional-panel}
\end{equation}
with no mass or infrared scale in the defining dimensionless marginal
calibration amplitudes.  This is a UV calibration panel, not a removal of the
positive IR/reference scale from the complete theory.  Thresholds, physical
masses and on-shell input conversion are postponed to P5.

The outcome has two parts.  First, the one-loop Standard-Model marginal beta
vector is fixed in the native calibrated coordinates by ordinary one-loop
scheme universality and can therefore be populated explicitly.  Second, the
same argument populates the first mixed gravity--matter coordinate, the
nonminimal Higgs-curvature coupling $\xi$ multiplying $R H^\dagger H$, at
order $G_N^0$.  Gauge-dependent field anomalous dimensions are deliberately
not promoted to phenomenological numbers: V17 instead specifies the exact
calibrated residue combinations which must be evaluated in the native BV gauge
when a field/source table is required.

\subsection{Conventions for the first phenomenological RG packet}
\label{sec:v17-conventions}

We use the $SU(5)$-normalized hypercharge coupling
\begin{equation}
 g_1=\sqrt{\frac53}\,g_Y,
 \label{eq:v17-g1}
\end{equation}
and the Higgs potential convention
\begin{equation}
 V(H)=-m_H^2 H^\dagger H+\lambda_H(H^\dagger H)^2.
 \label{eq:v17-Higgs-potential}
\end{equation}
The Yukawa matrices are $Y_u,Y_d,Y_e$ in family space.  Define
\begin{align}
 T&:=\Tr\!\left(
  3Y_u^\dagger Y_u+3Y_d^\dagger Y_d+Y_e^\dagger Y_e
 \right),
 \label{eq:v17-T}\\
 Y_4&:=\Tr\!\left(
  3(Y_u^\dagger Y_u)^2
 +3(Y_d^\dagger Y_d)^2
 +(Y_e^\dagger Y_e)^2
 \right).
 \label{eq:v17-Y4}
\end{align}
For the mixed curvature block we write the local action term as
\begin{equation}
 S_{\xi}=\int d^4x\sqrt g\;\xi\,R H^\dagger H,
 \label{eq:v17-xi-convention}
\end{equation}
so that the conformal scalar value in this sign convention is $\xi=1/6$.
Changing the overall sign convention for the curvature term requires the
corresponding translation of the formula below.

Put
\begin{equation}
 \epsilon_{16}:=\frac{1}{16\pi^2}.
 \label{eq:v17-epsilon16}
\end{equation}
The V17 displayed beta functions are one-loop equations of the form
\begin{equation}
 \frac{dg^A}{d\log\mu}
 =\epsilon_{16}\,\beta_{(1)}^A+O(\epsilon_{16}^2)
 \label{eq:v17-beta-convention}
\end{equation}
unless stated otherwise.

\begin{definition}[V17 massless nonexceptional UV calibration window]
\label{def:v17-UV-window}
The first RG packet is evaluated on a calibration panel satisfying:
\begin{enumerate}[label=\textup{(U\arabic*)},leftmargin=3em]
\item the reference background is flat and the marginal Standard-Model block
is extracted at order $(G_N\mu^2)^0$;
\item every defining 1PI momentum configuration is nonexceptional and has the
fixed dimensionless shape \eqref{eq:v17-nonexceptional-panel};
\item masses are omitted in the logarithmic UV residue; their conversion and
threshold effects are retained for the later P5 matching problem;
\item the native and standard schemes use the same tree-level coupling and
operator conventions \eqref{eq:v17-g1}--\eqref{eq:v17-xi-convention};
\item all wave-function, source/contact and BV restoring terms required by the
V16 calibration Jacobian are included before the logarithmic coefficient is
read.
\end{enumerate}
\end{definition}

\begin{firewall}[This is not the final low-energy input scheme]
The V17 UV window is designed to populate the universal one-loop logarithmic
vector.  Experimental calibration will normally be performed with massive,
possibly on-shell or mixed inputs near the electroweak scale.  Converting those
inputs to the UV running coordinates requires finite threshold/matching terms
and can generate explicit mass-ratio dependence.  Those terms belong to P3--P5
and are not silently absorbed into the V17 beta table.
\end{firewall}

\subsection{One-loop scheme universality and its exact boundary}
\label{sec:v17-scheme-universality}

The V16 transformation law immediately determines what may be imported from a
standard one-loop computation.

\begin{theorem}[One-loop transformation under a finite coordinate change]
\label{thm:v17-one-loop-scheme}
Let $g$ denote any finite collection of dimensionless scalar or matrix
couplings and suppose
\begin{equation}
 \beta(g,\mu)=\epsilon_{16}\beta_{(1)}(g)+O(\epsilon_{16}^2).
 \label{eq:v17-beta-expand}
\end{equation}
Let another calibrated coordinate system be
\begin{equation}
 g'=g+\epsilon_{16}f(g,\mu)+O(\epsilon_{16}^2).
 \label{eq:v17-finite-change}
\end{equation}
Then
\begin{equation}
 \boxed{
 \beta'_{(1)}(g)
 =\beta_{(1)}(g)+\mu\,\partial_\mu f(g,\mu).}
 \label{eq:v17-one-loop-transform}
\end{equation}
In particular, if $f$ has no explicit $\mu$ dependence,
\begin{equation}
 \boxed{\beta'_{(1)}=\beta_{(1)}.}
 \label{eq:v17-one-loop-universal}
\end{equation}
\end{theorem}

\begin{proof}
Use the V16 exact transformation law
\[
 \beta'=D_gF\,\beta+\mu\partial_\mu F
\]
with $F(g,\mu)=g+\epsilon_{16}f(g,\mu)+O(\epsilon_{16}^2)$.  Since
$D_gF=I+\epsilon_{16}D_gf+O(\epsilon_{16}^2)$ and
$\beta=O(\epsilon_{16})$, the product
$\epsilon_{16}D_gf\,\beta$ starts at two loops.  The only additional one-loop
term is therefore $\epsilon_{16}\mu\partial_\mu f$, which proves
\eqref{eq:v17-one-loop-transform}.
\end{proof}

\begin{proposition}[Scale independence of the finite matching map in the V17 window]
\label{prop:v17-scale-independent-match}
In the massless nonexceptional window of Definition~\ref{def:v17-UV-window},
the continuum one-loop native-to-$\overline{\rm MS}$ matching of the
\emph{dimensionless marginal block}, after the common tree conventions are
identified, has the form
\begin{equation}
 g^{\overline{\rm MS}}
 =g^{\rm nat}
 +\epsilon_{16}f(g^{\rm nat};\widehat p_i\!\cdot\!\widehat p_j)
 +O(\epsilon_{16}^2),
 \label{eq:v17-massless-match}
\end{equation}
with no explicit $\mu$ dependence.  Consequently the one-loop native beta
vector equals the one-loop $\overline{\rm MS}$ vector on this packet.
\end{proposition}

\begin{proof}
After the native continuum limit has been taken, the difference of two finite
renormalization prescriptions for a dimensionless marginal coefficient is a
finite local dimensionless function.  In the declared window there is no mass,
IR regulator or gravitational length entering the defining nonexceptional 1PI
coefficient.  All external momenta are $\mu$ times a fixed dimensionless
configuration.  Dimensional analysis therefore leaves only the couplings and
the fixed dimensionless momentum angles as arguments of the finite one-loop
matching term.  In particular there is no independent dimensionless ratio
containing $\mu$.  Equation \eqref{eq:v17-one-loop-universal} applies.
\end{proof}

\begin{firewall}[Mass-dependent schemes are different]
If a later physical scheme contains a mass ratio $m/\mu$, a threshold ratio
$M/\mu$, a finite curvature ratio $R/\mu^2$, or another explicit scale, then
$\mu\partial_\mu f$ in \eqref{eq:v17-one-loop-transform} need not vanish.
One must transform the beta vector together with the finite matching map.  V17
therefore does not claim that the one-loop table below is numerically identical
in every on-shell or threshold scheme.
\end{firewall}

\subsection{The calibrated Standard-Model one-loop marginal vector}
\label{sec:v17-SM-vector}

The standard one-loop coefficients in the conventions above are now promoted
to the native calibrated V17 UV packet by
Proposition~\ref{prop:v17-scale-independent-match}.  These coefficients are
not new Renewal predictions; they are a nontrivial validation target for the
native shell implementation.

\begin{theorem}[Populated one-loop Standard-Model marginal packet at $O(G_N^0)$]
\label{thm:v17-SM-beta}
On the V17 massless nonexceptional UV calibration panel,
\begin{align}
 16\pi^2\beta_{g_1}
 &=\frac{41}{10}g_1^3,
 \label{eq:v17-beta-g1}\\
 16\pi^2\beta_{g_2}
 &=-\frac{19}{6}g_2^3,
 \label{eq:v17-beta-g2}\\
 16\pi^2\beta_{g_3}
 &=-7g_3^3.
 \label{eq:v17-beta-g3}
\end{align}
The matrix Yukawa beta functions are
\begin{align}
 16\pi^2\beta_{Y_u}
 &=Y_u\left[
 \frac32(Y_u^\dagger Y_u-Y_d^\dagger Y_d)
 +T
 -\frac{17}{20}g_1^2
 -\frac94g_2^2
 -8g_3^2
 \right],
 \label{eq:v17-beta-Yu}\\
 16\pi^2\beta_{Y_d}
 &=Y_d\left[
 \frac32(Y_d^\dagger Y_d-Y_u^\dagger Y_u)
 +T
 -\frac14g_1^2
 -\frac94g_2^2
 -8g_3^2
 \right],
 \label{eq:v17-beta-Yd}\\
 16\pi^2\beta_{Y_e}
 &=Y_e\left[
 \frac32Y_e^\dagger Y_e
 +T
 -\frac94g_1^2
 -\frac94g_2^2
 \right].
 \label{eq:v17-beta-Ye}
\end{align}
For the quartic coupling,
\begin{align}
 16\pi^2\beta_{\lambda_H}
 ={}&24\lambda_H^2
 -\left(\frac95g_1^2+9g_2^2\right)\lambda_H
 +\frac{27}{200}g_1^4
 +\frac9{20}g_1^2g_2^2
 +\frac98g_2^4
 \notag\\
 &+4\lambda_HT-2Y_4.
 \label{eq:v17-beta-lambda}
\end{align}
Finally, for the Higgs mass parameter in
\eqref{eq:v17-Higgs-potential},
\begin{equation}
 \boxed{
 16\pi^2\beta_{m_H^2}
 =m_H^2\left[
 12\lambda_H+2T-\frac9{10}g_1^2-\frac92g_2^2
 \right].}
 \label{eq:v17-beta-mass}
\end{equation}
All equations are at one loop and at order $(G_N\mu^2)^0$.
\end{theorem}

\begin{proof}
The displayed coefficients are the standard one-loop Standard-Model beta
functions in the unbroken massless theory with $SU(5)$-normalized $g_1$ and
potential convention \eqref{eq:v17-Higgs-potential}.  The gauge coefficients
come from the gauge, fermion and Higgs contributions to the three gauge
vacuum-polarization/vertex calibration rows.  The Yukawa expressions are the
matrix one-loop equations, and the quartic/mass equations are written in terms
of the invariant traces \eqref{eq:v17-T}--\eqref{eq:v17-Y4}.
Proposition~\ref{prop:v17-scale-independent-match} identifies these one-loop
logarithmic residues with the native calibrated residues on the V17 UV panel.
The theorem does not assert equality of finite parts or of higher-loop
coefficients.
\end{proof}

\begin{corollary}[Third-generation/top-dominant check]
\label{cor:v17-top-check}
If $Y_d,Y_e$ and all up-type Yukawas except $y_t$ are neglected, then
\begin{equation}
 \boxed{
 16\pi^2\beta_{y_t}
 =y_t\left(
 \frac92y_t^2
 -\frac{17}{20}g_1^2
 -\frac94g_2^2
 -8g_3^2
 \right),}
 \label{eq:v17-beta-yt}
\end{equation}
while \eqref{eq:v17-beta-lambda} reduces to
\begin{align}
 16\pi^2\beta_{\lambda_H}
 ={}&24\lambda_H^2
 +\lambda_H\left(
 12y_t^2-\frac95g_1^2-9g_2^2
 \right)
 -6y_t^4
 \notag\\
 &+\frac{27}{200}g_1^4
 +\frac9{20}g_1^2g_2^2
 +\frac98g_2^4.
 \label{eq:v17-beta-lambda-top}
\end{align}
\end{corollary}

\begin{proof}
Set $T=3y_t^2$ and $Y_4=3y_t^4$ in
Theorem~\ref{thm:v17-SM-beta}.
\end{proof}

\subsection{First mixed gravity--matter row: the nonminimal Higgs coupling}
\label{sec:v17-xi}

A phenomenology-oriented gravity packet should begin with a coefficient whose
operator and normalization are unambiguous before tackling the entire
curvature-squared basis.  The natural first row is the nonminimal coupling
\eqref{eq:v17-xi-convention}.  At order $G_N^0$, it is generated already by
Standard-Model matter loops on a classical curved background.

\begin{theorem}[One-loop Higgs--curvature running at $O(G_N^0)$]
\label{thm:v17-beta-xi}
In the convention \eqref{eq:v17-xi-convention}, the one-loop matter-induced
running of the nonminimal Higgs-curvature coupling is
\begin{equation}
 \boxed{
 16\pi^2\beta_\xi\big|_{G_N^0}
 =\left(\xi-\frac16\right)
 \left[
  12\lambda_H+2T
  -\frac9{10}g_1^2
  -\frac92g_2^2
 \right].}
 \label{eq:v17-beta-xi}
\end{equation}
Thus the conformal value $\xi=1/6$ is a one-loop fixed surface for the
matter-induced flow in this convention.
\end{theorem}

\begin{proof}
The curved-space one-loop divergence multiplying $RH^\dagger H$ is
proportional to $\xi-1/6$.  In the hypercharge convention using $g_Y$, the
gauge term is $-(3/2)g_Y^2-(9/2)g_2^2$.  Using
$g_Y^2=(3/5)g_1^2$ gives $-(9/10)g_1^2-(9/2)g_2^2$.  The Yukawa trace is
$2T$ and the quartic term is $12\lambda_H$.  As for the marginal flat-space
packet, the one-loop coefficient is invariant under a finite scale-independent
change of the dimensionless $\xi$ coordinate.  Quantum-gravity corrections
proportional to powers of $G_N\mu^2$ are not included.
\end{proof}

\begin{corollary}[Common multiplicative bracket for $m_H^2$ and $\xi-1/6$]
\label{cor:v17-common-bracket}
At order $G_N^0$ define
\begin{equation}
 \mathcal B_H
 :=12\lambda_H+2T-\frac9{10}g_1^2-\frac92g_2^2.
 \label{eq:v17-BH}
\end{equation}
Then
\begin{equation}
 16\pi^2\frac{d\log m_H^2}{d\log\mu}=\mathcal B_H,
 \qquad
 16\pi^2\frac{d\log|\xi-1/6|}{d\log\mu}=\mathcal B_H
 \label{eq:v17-common-bracket-flow}
\end{equation}
where the logarithms are interpreted away from the corresponding zero surfaces.
Hence
\begin{equation}
 \boxed{
 \frac{\xi(\mu)-1/6}{m_H^2(\mu)}
 =\text{constant at one loop and }O(G_N^0)}
 \label{eq:v17-ratio-invariant}
\end{equation}
within a mass-independent interval in which both equations apply.
\end{corollary}

\begin{proof}
Divide \eqref{eq:v17-beta-mass} and \eqref{eq:v17-beta-xi} by the corresponding
coordinates.  Their logarithmic derivatives are identical, so their difference
vanishes at the retained order.
\end{proof}

\begin{firewall}[The ratio is not yet a Renewal prediction]
Equation \eqref{eq:v17-ratio-invariant} is an RG invariant of the one-loop
matter system.  It does not fix its numerical value.  Unless the Renewal
matching law selects either $m_H^2$, $\xi$, or their ratio beyond the
experimental input set, the constant is a boundary datum rather than a new
prediction.  This is exactly the V16 distinction between RG structure and a
positive-codimension predictive matching locus.
\end{firewall}

\subsection{Native extraction: the calibrated logarithmic-residue equation}
\label{sec:v17-native-extraction}

The explicit table above is useful only if the native calculation has a
well-defined row-by-row target.  V17 now gives that target independently of
any gauge choice for the intermediate wave-function factors.

Let $g^A$ denote the calibrated coordinates
\begin{equation}
 g^A\in
 \{g_1,g_2,g_3,Y_u,Y_d,Y_e,\lambda_H,m_H^2,\xi\}
 \label{eq:v17-coordinate-list}
\end{equation}
and let
\begin{equation}
 \mathcal R_a(g,\mu)=0,
 \qquad a=1,\ldots,n_{\rm cal},
 \label{eq:v17-calibration-residual}
\end{equation}
be the complete finite set of calibration equations defining these coordinates
from normalized 1PI/source coefficients.  For example, a residual may be a
normalized transverse gauge vertex minus the declared coupling, a normalized
Yukawa vertex minus $Y_f$, or a four-Higgs coefficient minus $\lambda_H$.
Field residues and source normalizations are included in the same residual
system rather than treated as external graphwise constants.

\begin{theorem}[Calibrated beta vector from native logarithmic residues]
\label{thm:v17-residue-equation}
Let
\begin{equation}
 J_{aA}:=\frac{\partial\mathcal R_a^{(0)}}{\partial g^A}
 \label{eq:v17-residue-J}
\end{equation}
be the invertible tree/reference calibration Jacobian on the chosen independent
coordinate block.  Let
\begin{equation}
 \mathcal L_a(g)
 :=\left.
 \mu\frac{\partial}{\partial\mu}
 \mathcal R_a^{(1)}(g,\mu)
 \right|_{g\,\mathrm{fixed}}
 \label{eq:v17-log-residue-vector}
\end{equation}
be the one-loop explicit logarithmic derivative of the \emph{complete}
calibration residual, including the calibrated wave-function, contact, BV and
operator-basis terms.  Then
\begin{equation}
 \boxed{
 J_{aA}\,\beta_{(1)}^A+\mathcal L_a=0,}
 \label{eq:v17-residue-linear-system}
\end{equation}
and therefore
\begin{equation}
 \boxed{
 \beta_{(1)}^A=-(J^{-1})^{Aa}\mathcal L_a.}
 \label{eq:v17-residue-beta}
\end{equation}
This is the native machine-extraction rule for the V17 table.
\end{theorem}

\begin{proof}
Along the same physical continuum state, the defining calibration equations
remain zero when the reference scale is changed and the calibrated coordinates
are transported by their RG vector:
\[
 0=\frac{d}{d\log\mu}\mathcal R_a(g(\mu),\mu)
 =\beta^A\partial_A\mathcal R_a
  +\mu\partial_\mu\mathcal R_a.
\]
At one loop, replace $\partial_A\mathcal R_a$ by its tree value $J_{aA}$ and
retain the explicit one-loop scale derivative
\eqref{eq:v17-log-residue-vector}.  This gives
\eqref{eq:v17-residue-linear-system}; invertibility of the calibration
Jacobian gives \eqref{eq:v17-residue-beta}.
\end{proof}

\begin{corollary}[Why graphwise logarithms are not beta functions]
\label{cor:v17-no-graphwise-beta}
A logarithmic divergence or shell logarithm in a single unnormalized graph is
not, by itself, an entry of $\beta_{(1)}$.  Only its contribution to the common
residue vector $\mathcal L_a$ after the full calibration/source/BV combination
enters \eqref{eq:v17-residue-beta}.  Two different vertex choices for the same
gauge coupling must give the same calibrated beta coefficient once the native
Ward/BV identities are satisfied.
\end{corollary}

\begin{proof}
Equation \eqref{eq:v17-residue-beta} depends only on the complete calibration
residual.  A decomposition of that residual into individual graphs is not
coordinate invariant.  Equality between different admissible gauge-coupling
vertices follows because the native Ward/BV identities make their normalized
calibration residuals equivalent physical coordinates.
\end{proof}

\subsection{Field anomalous dimensions: required internal rows, not yet a physical table}
\label{sec:v17-field-gamma}

A numerical beta table for physical couplings must not be confused with a
numerical table of gauge-fixed elementary-field anomalous dimensions.
Individual wave-function factors depend on gauge and field parameterization,
while the calibrated coupling beta functions above are gauge independent.

\begin{definition}[V17 internal wave-function residue packet]
\label{def:v17-wave-packet}
For a chosen native BV gauge and field basis, let
\begin{equation}
 \Gamma_{\rm wf}^{(1)}
 =\{\gamma_H,\gamma_{Q_L},\gamma_{u_R},\gamma_{d_R},
   \gamma_{L_L},\gamma_{e_R},
   \gamma_{B},\gamma_W,\gamma_G,\ldots\}
 \label{eq:v17-wave-packet}
\end{equation}
be the one-loop logarithmic residue rows needed to assemble the normalization
residuals \eqref{eq:v17-calibration-residual}.  Their gauge-fixed numerical
values are \emph{internal data}.  The phenomenological output is the
calibrated combination \eqref{eq:v17-residue-beta}, together with operator
mixing matrices for gauge-invariant composites.
\end{definition}

\begin{proposition}[Gauge/basis changes cancel in the calibrated marginal vector]
\label{prop:v17-gauge-cancel}
Suppose two native BV gauges or field bases are related on the protected one-loop
chart by a finite local change preserving the same physical calibration
functionals.  Then their individual wave-function/vertex residue rows may
differ, but the beta vector obtained from
\eqref{eq:v17-residue-beta} is the same on the V17 UV packet.
\end{proposition}

\begin{proof}
The two descriptions are finite local coordinate/basis changes of the complete
calibration residual system.  On the massless nonexceptional panel their
one-loop finite matching is scale independent, so
Theorem~\ref{thm:v17-one-loop-scheme} leaves the one-loop physical beta vector
unchanged.  The cancellation is therefore a statement about the complete
calibration system, not equality of individual $Z$ factors.
\end{proof}

\begin{firewall}[Why V17 does not print elementary-field anomalous dimensions]
The V16 handoff asked for the field/source anomalous dimensions required to
define the calibrated couplings.  V17 identifies the exact rows and proves how
they enter the common residue system, but does not assign gauge-dependent
numbers before a native BV gauge and field normalization convention are frozen.
Doing so would create convention-dependent output while adding no new
phenomenological invariant.  Gauge-invariant composite/source mixing remains a
separate P2 calculation and should be tabulated once the benchmark observable
panel is chosen.
\end{firewall}

\subsection{Machine-readable first RG packet}
\label{sec:v17-machine-table}

The companion file
\begin{center}
\texttt{Renewal\_P2a\_OneLoop\_RG\_Packet\_V17.json}
\end{center}
contains the conventions, equations and population status of the V17 packet.
Its entries are deliberately typed as either
\texttt{POPULATED\_UNIVERSAL\_1L},
\texttt{POPULATED\_MATTER\_CURVED\_1L}, or
\texttt{OPEN\_NATIVE\_GAUGE\_DEPENDENT}.  This prevents a benchmark standard
coefficient from being silently confused with a native-specific finite part.

\begin{table}[ht]
\centering
\small
\begin{tabular}{@{}p{0.24\linewidth}p{0.28\linewidth}p{0.38\linewidth}@{}}
\toprule
Block & V17 status & Meaning\\
\midrule
$g_1,g_2,g_3$ & populated universal 1L & Native calibrated UV residues fixed by one-loop scheme universality.\\
$Y_u,Y_d,Y_e$ & populated universal 1L & Full matrix equations populated; finite native parts remain scheme dependent.\\
$\lambda_H$ & populated universal 1L & Full flavour-trace expression populated.\\
$m_H^2$ & populated mass parameter 1L & Mass-independent one-loop anomalous running; threshold conversion still open.\\
$\xi RH^\dagger H$ & populated matter-curved 1L & $O(G_N^0)$ mixed gravity--matter row; quantum-gravity corrections open.\\
Elementary-field $\gamma$ rows & internal/open native gauge convention & Needed to reconstruct native residues graphically, not a scheme-invariant phenomenological table.\\
Curvature-squared / higher mixed Wilson bank & open & Basis, field-redefinition quotient and native logarithmic coefficients still to be populated.\\
$O(G_N\mu^2)$ corrections & open & Genuine quantum-gravity contribution, distinct from the matter-curved $\xi$ row.\\
\bottomrule
\end{tabular}
\caption{Population status of the first calibrated phenomenological RG packet.}
\label{tab:v17-status-packet}
\end{table}

\subsection{Prediction audit: standard running is a validation row, not new phenomenology}
\label{sec:v17-prediction-audit}

The purpose of P2 is partly to verify that the native continuum reproduces the
known Standard-Model logarithmic flow, but reproducing it is not itself a novel
Renewal prediction.

\begin{theorem}[No new prediction from a universally matched beta row]
\label{thm:v17-beta-not-prediction}
Assume a calibrated coordinate $g^A(\mu_0)$ is an independent phenomenological
input and its one-loop beta coefficient is fixed solely by the universal
Standard-Model logarithmic residue of
Theorem~\ref{thm:v17-SM-beta}.  Then the resulting one-loop trajectory
$g^A(\mu)$ is a consistency consequence of that input and the Standard-Model
field content, not an independent Renewal prediction.  A new prediction
requires additional positive-codimension information in the V16 matching locus,
such as a relation among Wilson coefficients or between a Wilson coefficient
and the calibrated input vector.
\end{theorem}

\begin{proof}
The initial value $g^A(\mu_0)$ is free by hypothesis.  Solving an ODE with a
universal vector field transports that freely chosen boundary datum but does
not reduce the dimension of the allowed initial-condition manifold.  By V16,
a predictive covector must annihilate the tangent image of the matching map and
therefore requires positive codimension.  Universal running changes the
location of the allowed set at another scale but not its dimension.
\end{proof}

\begin{corollary}[What would count as a first Renewal RG prediction]
\label{cor:v17-first-RG-prediction}
A relation such as
\begin{equation}
 R\bigl(g_i,Y_f,\lambda_H,\xi,C_I^{\rm grav},C_J^{\rm mixed}\bigr)=0
 \label{eq:v17-candidate-relation}
\end{equation}
becomes a genuine RG prediction only if:
\begin{enumerate}[label=\textup{(P\arabic*)},leftmargin=3em]
\item it is selected by the populated microscopic Renewal matching law rather
than imposed as an input;
\item its differential is nonzero on the calibrated Wilson fibre;
\item it survives finite scheme/basis/redundant transformations; and
\item it obeys the V16 RG transport equation
$(\mu\partial_\mu+\beta\cdot\partial)R=0$ on the matching locus.
\end{enumerate}
\end{corollary}

\subsection{External calibration references used for the V17 benchmark}
\label{sec:v17-external-references}

The explicit one-loop Standard-Model coefficients above are benchmarked
against the conventional perturbative literature rather than inferred from an
uncomputed native graph table.  Useful references include:
\begin{enumerate}[leftmargin=2.4em]
\item L.~N.~Mihaila, J.~Salomon and M.~Steinhauser,
\emph{Gauge Coupling Beta Functions in the Standard Model to Three Loops},
\href{https://arxiv.org/abs/1201.5868}{arXiv:1201.5868};
\item A.~V.~Bednyakov, A.~F.~Pikelner and V.~N.~Velizhanin,
\emph{Yukawa coupling beta-functions in the Standard Model at three loops},
\href{https://arxiv.org/abs/1212.6829}{arXiv:1212.6829}, and
\emph{Three-loop Higgs self-coupling beta-function in the Standard Model with
complex Yukawa matrices},
\href{https://arxiv.org/abs/1310.3806}{arXiv:1310.3806};
\item T.~Markkanen, S.~Nurmi, A.~Rajantie and S.~Stopyra,
\emph{The 1-loop effective potential for the Standard Model in curved
spacetime},
\href{https://arxiv.org/abs/1804.02020}{arXiv:1804.02020};
\item A.~Yamada,
\emph{Lattice perturbation theory in the overlap formulation for the Yukawa
and gauge interactions},
\href{https://arxiv.org/abs/hep-lat/9802013}{hep-lat/9802013}, as a useful
methodological comparison for one-loop overlap-regulator self-energy and
vertex calculations;
\item A.~V.~Bednyakov, A.~S.~Fedoruk and D.~I.~Kazakov,
\emph{Renormalization-group analysis of the SM: Loops, uncertainties, and
vacuum stability},
\href{https://arxiv.org/abs/2509.03369}{arXiv:2509.03369}, for a recent review
of the Standard-Model RG hierarchy.
\end{enumerate}

\subsection{V17 anti-circularity audit}
\label{sec:v17-audit}

\begin{enumerate}[leftmargin=2.4em]
\item \textbf{Standard one-loop coefficients are imported only after proving
the relevant scheme condition.}  Theorem~\ref{thm:v17-one-loop-scheme} keeps
the explicit $\mu\partial_\mu f$ term; Proposition
\ref{prop:v17-scale-independent-match} removes it only on the massless
nonexceptional panel.
\item \textbf{Quantum gravity is not silently set equal to zero everywhere.}
The explicit Standard-Model table is the $O(G_N^0)$ block.  Terms proportional
to $G_N\mu^2$ are a separate EFT calculation.
\item \textbf{Thresholds are not hidden inside the beta table.}  The table is
an unbroken high-energy logarithmic packet.  Physical electroweak input
conversion and decoupling belong to P3--P5.
\item \textbf{Gauge-dependent wave-function anomalous dimensions are not called
observables.}  They are retained as internal rows of the native calibration
system and may be populated after the native BV gauge is frozen.
\item \textbf{The mixed curvature row is convention tagged.}  The
$\xi-1/6$ factor refers specifically to
\eqref{eq:v17-xi-convention}; another sign convention must be translated before
comparison.
\item \textbf{The RG invariant \eqref{eq:v17-ratio-invariant} is not called a
new prediction.}  Its integration constant remains free until the Renewal
matching law fixes it.
\item \textbf{No graphwise logarithm is promoted directly to a beta
coefficient.}  Native extraction uses the complete calibrated linear system
\eqref{eq:v17-residue-beta}.
\end{enumerate}

\section{V17 status ledger}
\label{sec:v17-status}

\subsection*{Proved in V17}
\begin{enumerate}[leftmargin=2.2em]
\item A finite one-loop coordinate change modifies the one-loop beta vector
only by its explicit scale derivative; scale-independent finite changes leave
the complete one-loop vector unchanged.
\item In the massless nonexceptional UV calibration window, the finite
native-to-$\overline{\rm MS}$ matching of the dimensionless Standard-Model
marginal block has no explicit scale, so its one-loop native beta vector equals
the conventional one-loop Standard-Model vector.
\item The full one-loop matrix beta functions for
$g_1,g_2,g_3,Y_u,Y_d,Y_e,\lambda_H$ are populated in the stated conventions,
together with the multiplicative running of $m_H^2$.
\item The first mixed gravity--matter row is populated at order $G_N^0$:
$\beta_\xi=(\xi-1/6)\mathcal B_H/(16\pi^2)$.  The same bracket governs
$m_H^2$, giving the one-loop matter RG invariant
$(\xi-1/6)/m_H^2$ within a mass-independent interval.
\item A common native residue equation
$J\beta_{(1)}+\mathcal L=0$ is proved.  It identifies the exact finite set of
normalized logarithmic shell coefficients which the native implementation must
compute to reproduce the populated table.
\item Gauge-dependent elementary-field anomalous dimensions are separated from
the gauge-independent calibrated marginal beta vector.  They remain internal
normalization rows until a native BV gauge and field basis are fixed.
\item A machine-readable companion packet records the conventions, equations
and population status of every first-block row.
\item The familiar Standard-Model running is classified as a validation and
transport result, not a new Renewal prediction; genuine new phenomenology still
requires a positive-codimension populated matching relation in the V16
predictive quotient.
\end{enumerate}

\subsection*{Inherited}
\begin{enumerate}[leftmargin=2.2em]
\item V1--V7: source-complete coefficientwise continuum, common BV restoration,
fixed-loop forests, ultraviolet Cauchy convergence, scheme/reference covariance
and exact loop-order projectivity.
\item V8--V15: regulator/fixed-index strengthening, culminating in automatic
coefficientwise unstabilized cancellation on a protected local chart.
\item V16: the phenomenology pivot, coefficientwise physical calibration,
predictive quotient, native RG generator, scheme/basis covariance and formal
native-to-standard matching theorem.
\end{enumerate}

\subsection*{Conditional}
\begin{enumerate}[leftmargin=2.2em]
\item Equality of the native and conventional one-loop marginal tables is
proved on the massless nonexceptional UV calibration panel.  A mass-dependent
on-shell or threshold scheme requires the explicit finite matching map and its
$\mu$ derivative.
\item The $\xi$ row is the matter-induced curved-background coefficient at
$O(G_N^0)$.  Quantum-gravity corrections and the complete curvature-squared
mixing block are not populated.
\item Elementary-field/source anomalous dimensions remain gauge/basis
convention dependent and have not been assigned a native numerical table.
\item All statements remain coefficientwise at the prescribed finite loop,
source and EFT order.  V17 does not start a numerical RG integration and does
not fit experimental initial conditions.
\end{enumerate}

\subsection*{Open beyond V17 --- main phenomenology line}
\begin{enumerate}[leftmargin=2.2em]
\item Freeze and populate the first \emph{predictive} gravity/mixed Wilson
quotient: determine whether the Renewal microscopic matching law fixes any
scheme-invariant combination of $\xi$, curvature-squared, Higgs-curvature and
other selected EFT coefficients after physical inputs are removed.
\item Populate the gauge-invariant composite/source anomalous-dimension block
for the first benchmark observable panel.
\item Compute the explicit native-to-standard finite one-loop matching constants
for a practical electroweak input scheme, then perform the $m_t,m_H,m_W,m_Z$
threshold handoff.
\item Add the first $O(G_N\mu^2)$ quantum-gravity corrections only after the
physical gravitational EFT basis and field-redefinition quotient are frozen.
\item Once the previous rows are populated, integrate the numerical RG
trajectory with input and theory-uncertainty propagation.
\end{enumerate}

\subsection*{Secondary mathematical branch}
\begin{enumerate}[leftmargin=2.2em]
\item Global descent of the V15 coefficientwise internal reserve, real Pfaffian
orientation, minimal-stratum crossings, positive-IR removal and stronger
nonperturbative questions remain available but are deliberately deprioritized
relative to the phenomenology line.
\end{enumerate}

\subsection*{Exact next load-bearing theorem}
\begin{center}
\fbox{\parbox{0.92\textwidth}{
\textbf{P4a: populate the first predictive gravity/mixed Wilson quotient.}
Choose a concrete covariant EFT basis at the calibration scale containing at
least the nonminimal Higgs-curvature coordinate and the first
curvature-squared/mixed operators relevant to a small observable panel.  Quotient
by field redefinitions, equations of motion, topological densities and the V7
finite scheme groupoid; then evaluate the actual Renewal microscopic matching
map into this Wilson fibre.  Compute the rank of its Jacobian and the left-null
predictive conormal.  Either exhibit the first nontrivial scheme-invariant
Renewal coefficient relation and transport it with the V17 beta vector, or
prove that this entire selected block remains freely matchable and hence yields
no new prediction.  Only after this rank test should threshold matching and
numerical RG integration be promoted to the main computational task.}}
\end{center}

\section*{Append-only continuation marker: V17}
\addcontentsline{toc}{section}{Append-only continuation marker: V17}
\textbf{End of V17.}  The complete V1--V16 body above is retained verbatim.
V17 populates the first phenomenological one-loop RG packet.  It first proves
the exact one-loop finite-scheme transformation law and isolates a massless
nonexceptional UV calibration panel on which the finite native-to-standard map
is scale independent.  This promotes the conventional one-loop
Standard-Model marginal beta vector to the native calibrated $O(G_N^0)$
packet, including full matrix Yukawa equations, the Higgs quartic and mass
running.  It also populates the first mixed gravity--matter row, the
matter-induced running of $\xi RH^\dagger H$, proves the common native
logarithmic-residue extraction equation, and separates gauge-dependent
wave-function rows from the phenomenological beta vector.  The populated
Standard-Model flow is explicitly classified as a validation row rather than a
new Renewal prediction.  The next primary gate is P4a: compute the actual
microscopic matching rank in a selected gravity/mixed Wilson fibre and either
extract the first scheme-invariant Renewal relation or prove that the block is
still entirely free.



\clearpage
\section{V18 continuation: P4a predictivity audit of the first gravity--matter Wilson block}
\label{sec:v18-continuation}

V16 defined the calibrated phenomenological quotient and the predictive
conormal, while V17 populated the first one-loop Standard-Model RG packet and
one mixed curved-space row.  The V17 handoff proposed that the first genuine
phenomenology test should now be a rank calculation in a selected
gravity--matter Wilson fibre.  The present continuation performs that test.

The result is negative in a mathematically useful sense.  The continuum
construction inherited from V80 and closed in V1--V7 deliberately retains the
independent ghost-zero EFT/Wilson quotient as matching data.  After physical
calibration, an arbitrary small finite addition in any retained independent
Wilson direction is still allowed.  Therefore the current matching map has
identity columns in those directions, its Jacobian has full row rank on the
selected Wilson block, and the V16 predictive conormal is zero.  Running those
coefficients, or choosing a scheme in which some of them vanish at one scale,
does not convert them into predictions.

This is not a failure of the continuum theorem.  It identifies the next
phenomenological datum which is genuinely missing: a microscopic Renewal
\emph{selection law} which removes some of the independent finite Wilson
coordinates.  Without such a law, computing more local beta functions only
transports freely chosen boundary data.

There is, however, a second benchmark sector which is already invariant under
finite local counterterms: nonanalytic momentum dependence of renormalized
amplitudes.  V18 therefore separates
\begin{center}
\fbox{\parbox{0.88\textwidth}{\centering
\textbf{local Wilson relations requiring microscopic selection}\\[0.25em]
$\oplus$\\[0.25em]
\textbf{nonanalytic loop tails fixed by calibrated low-energy data}}}
\end{center}
The second sector is useful as a stringent native validation target, although
it is not a new Renewal-specific prediction when it coincides with ordinary
low-energy Einstein--Standard Model EFT.

\subsection{A closed pilot operator bank}
\label{sec:v18-pilot-bank}

Work at the V16 calibration scale \(\mu_0>0\), at one fixed finite coefficient
packet \(\mathfrak q\), and on a background/observable panel with no physical
boundary contribution.  Define
\begin{align}
 \mathcal O_\xi&:=R\,H^\dagger H,
 \label{eq:v18-Oxi}\\
 \mathcal O_C&:=C_{\mu\nu\rho\sigma}C^{\mu\nu\rho\sigma},
 \label{eq:v18-OC}\\
 \mathcal O_R&:=R^2,
 \label{eq:v18-OR}\\
 E_4&:=R_{\mu\nu\rho\sigma}R^{\mu\nu\rho\sigma}
      -4R_{\mu\nu}R^{\mu\nu}+R^2.
 \label{eq:v18-E4}
\end{align}
The raw pilot bank is not just the three displayed representatives.  Local
metric and matter field redefinitions can move curvature-squared coefficients
into matter operators through the classical equations of motion.  We therefore
close the bank before quotienting:
\begin{equation}
 \mathcal B_{\rm pilot}^{\rm raw}
 :=\operatorname{span}\{\mathcal O_\xi,\mathcal O_C,\mathcal O_R,
                          E_4,\Box R\}
   +\mathcal B_{\rm ind},
 \label{eq:v18-raw-pilot}
\end{equation}
where \(\mathcal B_{\rm ind}\) is the finite set of matter/mixed local operators
of the same prescribed EFT order generated by the allowed field-redefinition
and equation-of-motion closure on the selected panel.  Its explicit size is
finite because \(\mathfrak q\) is fixed.

The V16 physical quotient removes the redundant/BV-exact, field-redefinition,
operator-basis and finite-scheme directions.  Let
\begin{equation}
 \mathcal W_{\rm pilot}
 \subset T\mathcal F_p
 \label{eq:v18-Wpilot}
\end{equation}
be the resulting finite-dimensional physical Wilson subspace generated by the
image of \eqref{eq:v18-raw-pilot}; write
\begin{equation}
 d_{\rm pilot}:=\dim\mathcal W_{\rm pilot}.
 \label{eq:v18-dpilot}
\end{equation}
Choose any local physical section
\begin{equation}
 (W^1,\ldots,W^{d_{\rm pilot}})
 \label{eq:v18-Wcoords}
\end{equation}
with dual functionals \(\ell_I\).  The class of
\(R H^\dagger H\) is included whenever it is nonredundant on the chosen panel.
The curvature-squared part is represented by any independent quotient section;
the notation \(\mathcal O_C,\mathcal O_R\) is not used to claim that these two
raw representatives remain separately invariant under every field
redefinition.

\begin{firewall}[Topological and boundary terms]
\label{fire:v18-topological}
On the boundaryless fixed-topology pilot panel, the integrated Euler density
\(E_4\) is topological and \(\Box R\) is a total derivative, so neither is an
independent bulk Wilson coordinate.  If a future observable panel contains a
physical boundary, the associated boundary operators must be restored before
quotienting; V18 does not discard them in that different problem.
\end{firewall}

\subsection{Independent finite Wilson additions survive physical calibration}
\label{sec:v18-independent-additions}

The load-bearing input is already present in the inherited continuum theorem:
independent ghost-zero EFT coordinates remain in the Wilson quotient after BV
restoration, while lower-weight physical normalization conditions are imposed
separately.  V16 then proved that the calibrated lower-weight coordinates are
smooth functions of the retained Wilson variables.

\begin{proposition}[Local Wilson-translation chart after calibration]
\label{prop:v18-Wilson-translation}
Fix a calibrated point \(c_0\in\mathcal F_p\) and the pilot quotient
\(\mathcal W_{\rm pilot}\).  After shrinking to a constant-rank chart, there
is a neighborhood \(U\subset\mathbb R^{d_{\rm pilot}}\) and a smooth map
\begin{equation}
 \Phi_{\mu_0}:U\longrightarrow\mathcal F_p,
 \qquad
 \Phi_{\mu_0}(0)=c_0,
 \label{eq:v18-Phi}
\end{equation}
such that
\begin{equation}
 \boxed{
 W^I\bigl(\Phi_{\mu_0}(u)\bigr)
 =W^I(c_0)+u^I,
 \qquad 1\le I\le d_{\rm pilot}.}
 \label{eq:v18-translation-coords}
\end{equation}
The physical input coordinates remain exactly equal to the calibrated values
\(p\); any lower-weight compensating shift is supplied by the V16 calibration
section rather than by changing the Wilson displacement \(u\).
\end{proposition}

\begin{proof}
Choose local representatives \(\mathcal O_I^{\rm phys}\) of the quotient
coordinates \eqref{eq:v18-Wcoords}.  The inherited common-action construction
allows a finite local addition
\begin{equation}
 \delta S_{\rm fin}(u)=\sum_Iu^I\mathcal O_I^{\rm phys}
 \label{eq:v18-finite-addition}
\end{equation}
inside the retained ghost-zero EFT bank.  By construction of the quotient
section, its differential in the Wilson coordinates is the identity.

This finite addition can shift lower-weight calibration observables through
loops or field redefinitions.  V16's coefficientwise calibration theorem
solves the lower-weight normalization equations uniquely for the calibrated
coordinates as smooth functions of the retained Wilson data.  Substitute that
calibration section after \eqref{eq:v18-finite-addition}.  It restores the input
vector \(p\) without altering the declared quotient coordinate \(u\).  This
defines \(\Phi_{\mu_0}\) and proves
\eqref{eq:v18-translation-coords}.
\end{proof}

\begin{corollary}[Finite schemes do not remove the local Wilson freedom]
\label{cor:v18-scheme-Wilson-freedom}
Let \(F\) be any admissible V7/V16 finite scheme/basis change on the calibrated
fibre.  The image of \(\Phi_{\mu_0}\) under \(F\) still has tangent rank
\(d_{\rm pilot}\).  Thus a finite scheme may move the numerical origin and mix
the pilot coordinates, but cannot create a predictive codimension which was
absent before the transformation.
\end{corollary}

\begin{proof}
The transformed tangent map is
\(DF\circ D\Phi_{\mu_0}\).  Since \(DF\) is invertible and
\(D\Phi_{\mu_0}|_0\) has rank \(d_{\rm pilot}\), their composition has the
same rank.
\end{proof}

\subsection{P4a rank theorem: the current pilot block has zero predictive conormal}
\label{sec:v18-rank-test}

The current proof state treats the independent Wilson coordinates themselves
as legitimate microscopic matching data.  Make this explicit by writing the
current matching parameter space locally as
\begin{equation}
 \Theta_{\rm cur}=\Theta_{\rm rest}\times U,
 \label{eq:v18-current-theta}
\end{equation}
where \(u\in U\subset\mathbb R^{d_{\rm pilot}}\) are the independent finite
pilot Wilson additions of Proposition~\ref{prop:v18-Wilson-translation}.
Let
\begin{equation}
 \mathfrak M_{\mu_0}^{\rm cur}:\Theta_{\rm cur}\to\mathcal F_p
 \label{eq:v18-current-match}
\end{equation}
be the matching map currently supported by the continuum theorem.

\begin{theorem}[P4a negative rank theorem for freely retained Wilson data]
\label{thm:v18-P4a-negative}
At every regular point of \(\mathfrak M_{\mu_0}^{\rm cur}\), the differential
of the pilot Wilson projection has full row rank:
\begin{equation}
 \boxed{
 \operatorname{rank}
 D\bigl(W\circ\mathfrak M_{\mu_0}^{\rm cur}\bigr)
 =d_{\rm pilot}.}
 \label{eq:v18-full-rank}
\end{equation}
Indeed, after the coordinate choice of
Proposition~\ref{prop:v18-Wilson-translation}, the columns corresponding to the
independent pilot matching variables obey
\begin{equation}
 \boxed{
 \frac{\partial W^I}{\partial u^J}=\delta^I{}_J.}
 \label{eq:v18-identity-columns}
\end{equation}
Consequently the V16 predictive conormal restricted to the pilot block is zero:
\begin{equation}
 \boxed{
 \mathcal P^*_{c_0}\cap\mathcal W_{\rm pilot}^*=\{0\}.}
 \label{eq:v18-zero-conormal}
\end{equation}
No nontrivial scheme-covariant local relation among the pilot Wilson
coordinates follows from the present continuum/RG theorem alone.
\end{theorem}

\begin{proof}
Hold \(\theta\in\Theta_{\rm rest}\) fixed and vary only the coordinate
\(u\).  Proposition~\ref{prop:v18-Wilson-translation} gives
\eqref{eq:v18-identity-columns}.  Therefore the pilot rows of the matching
sensitivity matrix contain an identity block, proving
\eqref{eq:v18-full-rank}.  A covector \(\ell_I\) in the V16 predictive
conormal must annihilate every column of the sensitivity matrix, including the
identity columns.  Hence \(\ell_I=0\) for all pilot directions, proving
\eqref{eq:v18-zero-conormal}.  Corollary
\ref{cor:v18-scheme-Wilson-freedom} makes the statement scheme covariant.
\end{proof}

\begin{corollary}[No current prediction of the conformal Higgs value]
\label{cor:v18-xi-not-predicted}
The current theorem does not predict
\begin{equation}
 \xi(\mu_0)=\frac16
 \label{eq:v18-xi-one-sixth}
\end{equation}
or any other numerical value of the nonminimal Higgs--curvature coefficient.
The fact that \(\xi=1/6\) is a fixed locus of the V17 matter-induced one-loop
flow does not select that locus as a boundary condition.
\end{corollary}

\begin{proof}
Whenever the \(R H^\dagger H\) class is an independent pilot coordinate,
Theorem~\ref{thm:v18-P4a-negative} permits its value to be varied locally while
the physical input panel is recalibrated.  An RG fixed locus is a property of
the vector field, not a selection of an initial condition.
\end{proof}

\begin{corollary}[The V17 ratio invariant still carries a free constant]
\label{cor:v18-ratio-free}
At the V17 one-loop \(O(G_N^0)\) level,
\begin{equation}
 \frac{\xi(\mu)-1/6}{m_H^2(\mu)}
 =K_\xi
 \label{eq:v18-Kxi}
\end{equation}
is RG invariant on a mass-independent interval, but \(K_\xi\) is not fixed by
the current Renewal continuum theorem.  It is therefore an invariant label of
RG trajectories, not yet a prediction.
\end{corollary}

\begin{proof}
V17 proves the scale invariance of the ratio.  Theorem
\ref{thm:v18-P4a-negative} leaves \(\xi(\mu_0)\) free after calibration, so the
integration constant \(K_\xi\) varies across the admitted matching fibre.
\end{proof}

\begin{corollary}[Zero bare curvature coefficients are a convention unless selected]
\label{cor:v18-zero-not-prediction}
A prescription such as
\begin{equation}
 W^I(\mu_0)=0
 \label{eq:v18-zero-boundary}
\end{equation}
for one or more independent pilot Wilson coordinates is a boundary condition
or scheme choice unless a separate microscopic Renewal selection theorem makes
\eqref{eq:v18-zero-boundary} invariantly true on the matching locus.  Its
subsequent radiative running does not retroactively make it a prediction.
\end{corollary}

\subsection{A matching-incidence compiler for genuine Renewal predictions}
\label{sec:v18-incidence-compiler}

The negative theorem identifies the correct next object.  One must replace the
formal Wilson matching variables \(u^I\) by the \emph{actual} microscopic
parameters of the finite Renewal action, if that action supplies fewer
independent directions.

Let \(\vartheta=(\vartheta^a)\) denote those genuinely microscopic parameters
after the V16 experimental inputs have been fixed and after any independent
Wilson coordinates which are to be claimed as predictions have been removed
from the parameter list.  Let
\begin{equation}
 S_{\rm loc}^{\rm match}(\vartheta)
 \label{eq:v18-native-local-action}
\end{equation}
be the finite local coefficient packet obtained from the common native action,
shell matching and physical calibration.  Denote by
\begin{equation}
 \Pi_{\rm phys}:\mathcal B_{\rm pilot}^{\rm raw}\to\mathcal W_{\rm pilot}
 \label{eq:v18-Piphys}
\end{equation}
the declared finite quotient projection and by \(\ell_I\) the dual Wilson
coordinates.

\begin{definition}[Microscopic matching-incidence matrix]
\label{def:v18-incidence}
The P4 matching-incidence matrix is
\begin{equation}
 \boxed{
 \mathsf S^I{}_a
 :=\ell_I\!\left(
       \Pi_{\rm phys}
       \frac{\partial S_{\rm loc}^{\rm match}}
            {\partial\vartheta^a}
       \right)_{\mu_0}.}
 \label{eq:v18-incidence}
\end{equation}
Every derivative is taken after the V16 physical calibration map is imposed,
so changes of lower-weight input coordinates are not counted as Wilson
predictions.
\end{definition}

\begin{theorem}[Exact rank criterion for a populated microscopic matching law]
\label{thm:v18-incidence-rank}
Assume \(S_{\rm loc}^{\rm match}(\vartheta)\) is populated on a constant-rank
chart.  Then the number of independent first-order Renewal predictions in the
pilot block is
\begin{equation}
 \boxed{
 d_{\rm pilot}-\operatorname{rank}\mathsf S.}
 \label{eq:v18-incidence-count}
\end{equation}
Candidate linearized predictive combinations are exactly the left-null vectors
\begin{equation}
 \boxed{
 \ell_I\mathsf S^I{}_a=0.}
 \label{eq:v18-incidence-leftnull}
\end{equation}
Under an invertible Wilson-basis/scheme change and an invertible microscopic
parameter reparametrization,
\begin{equation}
 \mathsf S\longmapsto A\,\mathsf S\,B
 \label{eq:v18-incidence-transform}
\end{equation}
with \(A,B\) invertible.  Therefore the rank, codimension and existence of a
nonzero predictive conormal are invariant.
\end{theorem}

\begin{proof}
Equation \eqref{eq:v18-incidence} is the coordinate differential of the V16
matching map restricted to the pilot quotient.  The prediction count and left
nullspace therefore follow directly from Theorem
\ref{thm:v16-predictive-codimension} and Corollary
\ref{cor:v16-prediction-matrix}.  An invertible change of Wilson coordinates
acts on rows by an invertible Jacobian \(A\), while an invertible
reparametrization of \(\vartheta\) acts on columns by an invertible Jacobian
\(B\).  Matrix rank and nullity are unchanged.
\end{proof}

\begin{corollary}[Identity columns are the exact predictivity killer]
\label{cor:v18-identity-killer}
If even one claimed predictive Wilson coordinate \(W^I\) is also retained as
an independent microscopic matching parameter \(\vartheta^I=W^I\), then the
corresponding row of \(\mathsf S\) contains an identity column.  If this holds
for every pilot coordinate, \(\operatorname{rank}\mathsf S=d_{\rm pilot}\)
and the predictive conormal is zero.
\end{corollary}

This is the precise algebraic form of the phenomenological bottleneck: the next
Renewal calculation must derive the pilot Wilson coefficients from a smaller
microscopic parameter set, not merely rename those Wilson coefficients as
``native inputs.''

\subsection{Local logarithmic running cannot fix the missing boundary data}
\label{sec:v18-running-not-selection}

Suppose the V17/V16 RG system for the pilot coordinates has the form
\begin{equation}
 \frac{dW^I}{d\log\mu}=\beta_W^I(g,W).
 \label{eq:v18-W-running}
\end{equation}
Then for any initial point \(W_0\) at \(\mu_0\), local existence gives a
trajectory
\begin{equation}
 W(\mu)=\mathcal U_W(\mu,\mu_0)W_0.
 \label{eq:v18-W-trajectory}
\end{equation}

\begin{theorem}[RG transport does not reduce a full-dimensional matching fibre]
\label{thm:v18-RG-no-codim}
If the matching image at \(\mu_0\) contains an open pilot Wilson set, then its
RG transport at every scale where the finite-order flow is locally invertible
also contains an open set of the same dimension.  In particular, local beta
functions cannot generate a predictive codimension absent from the matching
boundary data.
\end{theorem}

\begin{proof}
At fixed finite perturbative order the RG transport is a local diffeomorphism
on the compact calibrated chart for a sufficiently short scale interval; V7
and V16 give its inverse by reversing the reference-scale transport.  A local
diffeomorphism maps open sets to open sets and preserves their dimension.  The
statement extends across finitely many threshold-free intervals by composition.
\end{proof}

\begin{corollary}[Computing curvature beta tables is not yet the highest-value step]
\label{cor:v18-beta-deprioritize}
Until a microscopic matching law of positive codimension is populated, a more
complete local curvature/mixed beta table transports freely matchable Wilson
coefficients but does not create a new relation among them.  Such a table
remains useful for later running and consistency, but it is downstream of the
P4 incidence calculation if the immediate objective is a new prediction.
\end{corollary}

\subsection{The immediately scheme-invariant benchmark sector: nonanalytic amplitudes}
\label{sec:v18-nonanalytic}

The absence of a local Wilson prediction does not imply that every loop effect
is scheme dependent.  Local finite counterterms are analytic polynomials in
external momenta around a nonexceptional low-energy point.  They therefore
cannot change genuinely nonanalytic momentum dependence generated by massless
propagation.

Let a renormalized amplitude coefficient be decomposed locally as
\begin{equation}
 \mathcal A(q)
 =\mathcal P(q;W,\mu)+\mathcal N(q;g,\mu),
 \label{eq:v18-A-PN}
\end{equation}
where \(\mathcal P\) is analytic in the external invariants in the declared
window and \(\mathcal N\) contains nonanalytic terms such as logarithms or
branch functions.

\begin{theorem}[Finite local counterterms cannot change a nonanalytic coefficient]
\label{thm:v18-nonanalytic-invariant}
Assume two admissible renormalization schemes differ by a finite local action
through the EFT order retained by \(\mathfrak q\).  Then their amplitude
difference is analytic in the nonexceptional momentum window.  Consequently
the coefficient of any linearly independent nonanalytic structure, for example
\begin{equation}
 q^4\log\!\frac{-q^2-i0}{\mu^2},
 \label{eq:v18-qlogq}
\end{equation}
is invariant under that finite local scheme change.
\end{theorem}

\begin{proof}
A finite local operator with finitely many derivatives contributes a polynomial
in the external momenta, and field redefinitions or local source-basis changes
modify amplitudes by the same analytic local class together with terms
proportional to the equations of motion.  Therefore the difference of two
matched finite local schemes is analytic at the chosen nonexceptional point.
A nonzero multiple of \eqref{eq:v18-qlogq}, or any other branch/nonanalytic
basis element, cannot equal an analytic polynomial on a punctured neighborhood.
Its coefficient is therefore unchanged.
\end{proof}

\begin{corollary}[Universal long-distance loop tails are immediate validation observables]
\label{cor:v18-nonanalytic-validation}
Once the physical low-energy inputs and field content are fixed, nonanalytic
long-distance loop coefficients provide scheme-independent tests of the native
continuum calculation which do not require fixing the independent local
curvature Wilson coefficients.  Agreement with ordinary Einstein--Standard
Model EFT validates the native matching/continuum machinery but is not, by
itself, a new Renewal-specific prediction.
\end{corollary}

\begin{firewall}[Nonanalytic does not mean experimentally large]
\label{fire:v18-small-nonanalytic}
Scheme invariance is a statement about predictivity, not observability.
Gravitational loop tails can be parametrically tiny at accessible energies.
V18 therefore promotes them as clean validation benchmarks, not automatically
as the highest-significance experimental targets.
\end{firewall}

\subsection{External EFT comparison}
\label{sec:v18-external-EFT}

The separation displayed above is the standard logic of
gravity as an effective field theory: unknown short-distance physics is encoded
in local Wilson coefficients, whereas sufficiently long-distance nonanalytic
loop effects are fixed by the low-energy propagating degrees of freedom.  Useful
references are J.~F.~Donoghue,
\emph{General relativity as an effective field theory: The leading quantum
corrections}, \href{https://arxiv.org/abs/gr-qc/9405057}{gr-qc/9405057}, and
C.~P.~Burgess, \emph{Quantum Gravity in Everyday Life: General Relativity as an
Effective Field Theory},
\href{https://arxiv.org/abs/gr-qc/0311082}{gr-qc/0311082}.  These references are
used only for EFT comparison and motivation; the rank theorem above follows
from the inherited Renewal Wilson quotient and V16 calibration/predictivity
geometry.

\subsection{Machine-readable P4a audit}
\label{sec:v18-machine-audit}

A companion JSON artifact records the first predictivity audit with the
following logical status:
\begin{center}
\begin{tabular}{@{}p{0.25\linewidth}p{0.62\linewidth}@{}}
\toprule
Item & V18 result\\
\midrule
Pilot block & Higgs--curvature plus curvature-squared quotient, closed under the induced matter/EOM directions before quotienting.\\
Current Wilson status & Independent finite matching coordinates retained by the continuum theorem.\\
Sensitivity matrix & Contains a \(d_{\rm pilot}\times d_{\rm pilot}\) identity block.\\
Current predictive conormal & Zero in the pilot Wilson directions.\\
\(\xi=1/6\) & RG fixed locus, not selected boundary value.\\
Curvature coefficients set to zero & Convention/boundary condition unless a microscopic selection law is added.\\
Immediate invariant benchmark & Coefficients of nonanalytic momentum structures.\\
Next required object & Populated microscopic matching-incidence matrix \(\mathsf S^I{}_a\) with the independent Wilson identity columns removed from the microscopic parameter list.\\
\bottomrule
\end{tabular}
\end{center}

\subsection{V18 anti-circularity audit}
\label{sec:v18-audit}

\begin{enumerate}[leftmargin=2.4em]
\item \textbf{A field-redefinition-sensitive raw operator basis is not called
physical.}  The pilot bank is closed under induced matter/EOM directions before
the V16 quotient is taken.
\item \textbf{Topological and boundary rows are declared before quotienting.}
The first pilot panel is boundaryless; a future boundary observable must restore
its boundary operators.
\item \textbf{A chosen zero boundary coefficient is not called a prediction.}
Independent finite Wilson translations survive calibration and produce identity
columns in the sensitivity matrix.
\item \textbf{RG running is not used to manufacture codimension.}  Theorem
\ref{thm:v18-RG-no-codim} separates transport from boundary selection.
\item \textbf{The V17 \(\xi\) invariant is not overpromoted.}  Its integration
constant remains free on the current matching fibre.
\item \textbf{Nonanalytic loop coefficients are separated from local Wilson
coefficients.}  Their scheme invariance follows because finite local
counterterms are analytic; this does not make the effect numerically large or
Renewal specific.
\item \textbf{The negative result is constructive.}  It identifies the exact
missing matrix \(\mathsf S^I{}_a\) whose left nullspace would constitute new
Renewal phenomenology.
\end{enumerate}

\section{V18 status ledger}
\label{sec:v18-status}

\subsection*{Proved in V18}
\begin{enumerate}[leftmargin=2.2em]
\item A first curved-Higgs/gravity pilot bank is defined only after closure
under the matter operators induced by local field redefinitions and equations
of motion; topological/boundary directions are separated at the declared
panel.
\item Every independent physical pilot Wilson direction admitted by the current
continuum theorem has a local finite-translation chart after V16 physical
calibration.  Lower-weight observables are recalibrated without constraining
that Wilson displacement.
\item The current matching sensitivity matrix therefore contains identity
columns for all freely retained pilot Wilson coordinates.  Its pilot row rank
is maximal and the V16 predictive conormal is zero.  This statement is invariant
under admissible finite scheme/basis changes.
\item The current theory does not predict \(\xi(\mu_0)=1/6\), vanishing
curvature-squared coefficients, or the V17 ratio constant
\((\xi-1/6)/m_H^2\).  These remain boundary/matching data until a microscopic
selection theorem removes their independent Wilson freedom.
\item A basis-independent microscopic matching-incidence matrix
\(\mathsf S^I{}_a\) is constructed.  Its left nullspace is exactly the first
linearized predictive conormal, and its rank is invariant under invertible
Wilson-scheme and microscopic-parameter coordinate changes.
\item A finite-order RG flow cannot generate predictive codimension absent at
the matching scale; it transports the full-dimensional Wilson fibre by a local
diffeomorphism.
\item Coefficients of linearly independent nonanalytic momentum structures are
invariant under finite local counterterm/scheme changes.  They define immediate
scheme-independent native validation observables even when the local Wilson
boundary data are free.
\end{enumerate}

\subsection*{Inherited}
\begin{enumerate}[leftmargin=2.2em]
\item V1--V7: source-complete coefficientwise continuum, common BV restoration,
fixed-loop forests, ultraviolet Cauchy convergence, finite-scheme covariance
and exact loop-order projectivity.
\item V8--V15: regulator/fixed-index strengthening, retained as the secondary
mathematical branch.
\item V16: physical calibration, the phenomenological quotient, predictive
conormal, native RG generator, scheme covariance and finite native-to-standard
matching theorem.
\item V17: the calibrated one-loop Standard-Model marginal beta vector,
\(O(G_N^0)\) running of \(\xi\), and the common native logarithmic-residue
extraction equation.
\end{enumerate}

\subsection*{Conditional}
\begin{enumerate}[leftmargin=2.2em]
\item The exact numerical dimension \(d_{\rm pilot}\) depends on the chosen
observable panel and on completing the induced matter operator closure before
taking the physical quotient.  V18's rank theorem does not require choosing a
basis-dependent number prematurely.
\item The no-prediction theorem describes the \emph{current} Renewal matching
state, in which the independent Wilson quotient is retained as matching data.
A new microscopic source-selection principle could reduce the rank and create
predictive relations.
\item Nonanalytic benchmark coefficients are universal only after the physical
low-energy field content and calibrated couplings are fixed; threshold changes
must be matched before comparing across different effective theories.
\item The real Pfaffian/global reserve issues of V12--V15 remain independent and
deprioritized relative to the phenomenology line.
\end{enumerate}

\subsection*{Open beyond V18 --- main phenomenology line}
\begin{enumerate}[leftmargin=2.2em]
\item Populate the \emph{actual microscopic} matching law
\(S_{\rm loc}^{\rm match}(\vartheta)\) from the finite Renewal dynamics after
removing any Wilson coordinate which is merely being renamed as a microscopic
input.  Compute the incidence matrix \(\mathsf S^I{}_a\) on a small pilot
operator/observable panel and inspect its left nullspace.
\item If that incidence matrix has positive codimension, integrate the V17 RG
vector on the resulting constrained locus and convert the relation to a
standard phenomenological scheme.  If it remains full row rank, enlarge or
change the observable/operator panel rather than calling a fit a prediction.
\item In parallel, compute one native nonanalytic amplitude coefficient as a
scheme-invariant validation target; the natural candidates are low-momentum
stress/gravity or mixed stress--Higgs form factors for which the corresponding
ordinary EFT tail is known.
\item Only after a nonzero predictive conormal is found should the full
curvature/mixed anomalous-dimension bank, threshold matching and numerical RG
integration regain priority.
\end{enumerate}

\subsection*{Secondary mathematical branch}
The V15 global reserve descent, real Pfaffian orientation, index-changing
crossings, positive-IR removal and stronger nonperturbative questions remain
open and deliberately secondary to the matching-incidence calculation.

\subsection*{Exact next load-bearing theorem}
\begin{center}
\fbox{\parbox{0.92\textwidth}{
\textbf{P4b: populated microscopic matching-incidence matrix.}
Choose the smallest finite Renewal microscopic parameter packet which actually
controls a phenomenologically useful gravity--matter operator panel after the
experimental P1 inputs are fixed.  Compute the common local matched action
\(S_{\rm loc}^{\rm match}(\vartheta)\), project it to the V16 physical Wilson
quotient, and evaluate
\(\mathsf S^I{}_a=\ell_I\Pi_{\rm phys}\partial_{\vartheta^a}
S_{\rm loc}^{\rm match}\).  Independent Wilson coefficients are not permitted
to reappear as tautological microscopic parameters.  Determine the rank and
left nullspace.  A nonzero left-null covector is the first candidate
Renewal-specific coefficient relation; full row rank is a rigorous negative
phenomenology result and requires changing the panel or supplying additional
microscopic structure.  In either case, do not spend the next main block on
numerical RG integration until this rank test has been performed.}}
\end{center}

\section*{Append-only continuation marker: V18}
\addcontentsline{toc}{section}{Append-only continuation marker: V18}
\textbf{End of V18.}  The complete V1--V17 body above is retained verbatim.
V18 performs the first P4 rank audit.  It closes a curved-Higgs/curvature pilot
bank under field-redefinition/EOM directions before taking the physical
quotient, proves that every currently independent Wilson direction survives
physical calibration as a local finite-translation coordinate, and shows that
the present matching sensitivity matrix therefore contains an identity block.
The pilot predictive conormal is zero: neither \(\xi=1/6\), vanishing local
curvature coefficients, nor the V17 ratio invariant is selected by the current
continuum/RG theorem.  V18 constructs the exact microscopic matching-incidence
matrix whose left nullspace would encode a genuine Renewal prediction and
proves that RG running cannot manufacture codimension absent from the matching
boundary data.  It also separates the immediately scheme-invariant nonanalytic
loop sector as a clean native validation target.  The next primary gate is P4b:
populate the actual microscopic incidence matrix without reintroducing the
Wilson coordinates themselves as free inputs.



\clearpage
\section{V19 continuation: populated microscopic calibration incidence, residual race selectors, and the missing Wilson-selection bridge}
\label{sec:v19-continuation}

V18 reduced new phenomenology to a concrete rank problem: remove every Wilson
coordinate which is merely being renamed as a microscopic input, populate an
actual microscopic map into the physical Wilson quotient, and inspect its left
nullspace.  The first upstream audit shows that one must distinguish two
questions before taking that rank:
\begin{equation}
 \boxed{
 \text{is a microscopic incidence matrix populated?}
 \qquad\text{and only then}\qquad
 \text{what is its rank?}}
 \label{eq:v19-two-gates}
\end{equation}
A missing matching map is not the zero map.  Assigning zero sensitivity where
no source-to-action theorem has been supplied would manufacture predictions by
absence of data.

The inherited Grand Tensor construction does, however, contain one explicit
finite microscopic incidence calculation which can be used as a control.  Its
locally constant race law has six positive symmetric rates, six drift rates,
and one independently calibrated density coordinate.  The decoded geometric
packet consists of the six components of a positive spatial tensor, three shift
coordinates, and one density/scale coordinate.  V19 computes this incidence
matrix exactly, identifies the microscopic directions which remain after P1
geometric calibration, and then audits whether the existing corpus supplies a
map from those residual microscopic directions into the V18 gravity--matter
Wilson quotient.

The answer is sharp.  The geometric incidence has full row rank ten, while a
three-dimensional drift packet remains invisible to the decoded geometry but
is still visible to the underlying race law.  Those three directions are the
smallest currently explicit candidate microscopic selector packet after the
classical geometry is fixed.  What is \emph{not} present in the current corpus
is a theorem assigning the renormalized gravity--matter Wilson coefficients as
functions of those three variables.  The signed action/source row remains an
independent physical object.  Therefore the P4b Wilson incidence matrix cannot
yet be evaluated; the next load-bearing object is a source-selection bridge,
not another beta table.

\subsection{The explicit microscopic race packet}
\label{sec:v19-race-packet}

Let
\begin{equation}
 \vartheta_{\rm race}
 :=(s_1,\ldots,s_6,d_1,\ldots,d_6,a_\rho)
 \in\mathbb R^{13},
 \qquad a_\rho:=\log\rho,
 \label{eq:v19-race-theta}
\end{equation}
with $s_a>0$.  Let $R=(r_1\ \cdots\ r_6)$ be the fixed $3\times6$ matrix of
race directions inherited from the calibrated geometry branch.  The standing
hypotheses are
\begin{equation}
 \operatorname{rank}R=3,
 \qquad
 \{r_ar_a^T\}_{a=1}^6
 \text{ are linearly independent in }\operatorname{Sym}_3.
 \label{eq:v19-R-hyp}
\end{equation}
Write
\begin{equation}
 S:=\operatorname{diag}(s_1,\ldots,s_6),
 \qquad
 \mathcal B:=RSR^T,
 \qquad
 \beta:=Rd.
 \label{eq:v19-race-decoder}
\end{equation}
The ten-dimensional decoded calibration packet is
\begin{equation}
 \mathcal G_{\rm race}(\vartheta_{\rm race})
 :=\bigl(\operatorname{vec}_{\rm sym}\mathcal B,\,\beta,\,a_\rho\bigr)
 \in\mathbb R^{6}\oplus\mathbb R^3\oplus\mathbb R.
 \label{eq:v19-race-G}
\end{equation}
Here $\operatorname{vec}_{\rm sym}$ is any fixed linear identification
$\operatorname{Sym}_3\simeq\mathbb R^6$.

Define the $6\times6$ matrix
\begin{equation}
 \mathsf V_R
 :=\begin{pmatrix}
 |&&|\\
 \operatorname{vec}_{\rm sym}(r_1r_1^T)&\cdots&
 \operatorname{vec}_{\rm sym}(r_6r_6^T)\\
 |&&|
 \end{pmatrix}.
 \label{eq:v19-VR}
\end{equation}
By \eqref{eq:v19-R-hyp}, $\mathsf V_R$ is invertible.

\begin{theorem}[Fully populated race-to-geometry incidence matrix]
\label{thm:v19-race-incidence}
The differential of the actual race calibration map is the constant block
matrix
\begin{equation}
 \boxed{
 D\mathcal G_{\rm race}
 =\mathsf J_{\rm race}
 =\begin{pmatrix}
   \mathsf V_R&0&0\\
   0&R&0\\
   0&0&1
  \end{pmatrix}.}
 \label{eq:v19-race-J}
\end{equation}
Consequently
\begin{equation}
 \boxed{
 \operatorname{rank}\mathsf J_{\rm race}=10,
 \qquad
 \dim\ker\mathsf J_{\rm race}=3,}
 \label{eq:v19-race-rank}
\end{equation}
and
\begin{equation}
 \boxed{
 \ker\mathsf J_{\rm race}
 =\{(0,\delta d,0):R\delta d=0\}.}
 \label{eq:v19-race-kernel}
\end{equation}
In particular the left nullspace is zero: the explicit race law imposes no
first-order relation among the ten decoded geometry/calibration coordinates.
\end{theorem}

\begin{proof}
Differentiate \eqref{eq:v19-race-decoder}.  The derivative with respect to
$s_a$ is $r_ar_a^T$, so in the chosen symmetric-vector coordinates the first
block is exactly $\mathsf V_R$.  The derivative of $\beta=Rd$ with respect to
$d$ is $R$, and the density coordinate differentiates to one.  No cross block
occurs in \eqref{eq:v19-race-decoder}.  Hence
\eqref{eq:v19-race-J} follows.

The first block has rank six by linear independence of the dyadics, the second
has rank three by \eqref{eq:v19-R-hyp}, and the last has rank one.  Thus the
row rank is ten.  The $s$ and $a_\rho$ components of a kernel vector must vanish;
the remaining condition is exactly $R\delta d=0$.  Rank--nullity in the
six-dimensional $d$ space gives a three-dimensional kernel.  Since the matrix
has full row rank, its left nullspace is zero.
\end{proof}

\begin{corollary}[The first actual microscopic rank audit is calibration, not prediction]
\label{cor:v19-race-no-prediction}
Treat the ten coordinates in \eqref{eq:v19-race-G} as P1 physical geometry and
scale inputs.  Then the race law has enough microscopic freedom to calibrate
all ten independently at first order.  Its incidence matrix therefore supplies
a consistency/calibration result, not a positive-codimension phenomenological
prediction.
\end{corollary}

\begin{proof}
The V16 predictive conormal of a populated map is the annihilator of its tangent
image.  Theorem~\ref{thm:v19-race-incidence} says that tangent image is the
entire ten-dimensional output space, so the annihilator is zero.
\end{proof}

\subsection{The three residual race directions are genuine microscopic directions}
\label{sec:v19-residual-race}

The kernel in \eqref{eq:v19-race-kernel} should not be confused with a null
direction of the microscopic probability law.  The inherited race Fisher form
is positive on the complete $(s,d)$ packet under the protected rate
inequalities.  Hence the three directions in $\ker R$ change the law even
though they do not change the decoded shift $\beta$.

Choose a $6\times3$ matrix $N$ with orthonormal columns spanning $\ker R$ and a
fixed right inverse $R^+$ of $R$.  Every drift vector may be decomposed as
\begin{equation}
 d=R^+\beta+N\eta,
 \qquad
 \eta\in\mathbb R^3,
 \label{eq:v19-d-decomp}
\end{equation}
with $R N=0$.  After the geometry variables $(s,\beta,a_\rho)$ have been fixed
by P1, the coordinate
\begin{equation}
 \boxed{\eta:=N^Td}
 \label{eq:v19-eta}
\end{equation}
is the smallest explicitly displayed microscopic race packet which remains
unconsumed by that geometry calibration.

\begin{proposition}[Residual race packet is not a calibration gauge]
\label{prop:v19-eta-physical}
Under the protected race-rate inequalities, every nonzero variation
$\delta\eta$ induces a nonzero variation of the finite race likelihood and has
strictly positive Fisher cost, even though
\begin{equation}
 \delta\mathcal B=0,
 \qquad
 \delta\beta=0,
 \qquad
 \delta a_\rho=0.
 \label{eq:v19-eta-geometry-invisible}
\end{equation}
Therefore the three residual directions are legitimate candidate microscopic
selectors rather than redundant reparameterizations of the decoded geometry.
\end{proposition}

\begin{proof}
A residual variation has $\delta s=0$ and
$\delta d=N\delta\eta\ne0$.  The inherited race information form contains the
positive quadratic block
\[
 \lambda^{-1}h^2\,\delta d^T
 (S-h^2DS^{-1}D)^{-1}\delta d,
\]
when $\delta s=0$.  The protected rate inequalities make the matrix inside the
inverse positive definite.  Hence the cost is strictly positive for
$\delta d\ne0$.  Equation \eqref{eq:v19-eta-geometry-invisible} follows from
$RN=0$ and the decoder formulas.
\end{proof}

\begin{firewall}[Residual does not mean predictive]
The existence of three microscopic directions which survive geometry
calibration does not imply that any gravity--matter Wilson coefficient depends
on them.  That dependence is a separate source-to-action statement.  V19 uses
$\eta$ only as a candidate domain for the next matching map.
\end{firewall}

\subsection{Why the P4b Wilson incidence is not yet a matrix}
\label{sec:v19-missing-map}

V18 defined the desired incidence matrix only after a populated local matching
action
\begin{equation}
 S_{\rm loc}^{\rm match}(\vartheta)
 \label{eq:v19-required-local-action}
\end{equation}
exists as a function of genuine microscopic variables.  The inherited Renewal
architecture keeps several logically different objects separate: the normalized
predictive/event law, absolute physical calibration, the signed action/source
writer, and the independent physical variations used to test that action.  In
particular, the constructive vacuum branch explicitly does not promote an
arbitrary renewal law into a selector of the Einstein action.

The distinction can be formalized as follows.

\begin{definition}[Populated microscopic Wilson bridge]
\label{def:v19-selection-bridge}
Fix the P1 inputs $p$ and the V18 physical pilot quotient
$\mathcal W_{\rm pilot}$.  A \emph{populated microscopic Wilson bridge} on a
microscopic packet $\eta$ is a coefficientwise finite map
\begin{equation}
 \boxed{
 \mathcal Q_{\rm sel}:\eta
 \longmapsto S_{\rm bare/finite}^{\rm ESM}(\eta)}
 \label{eq:v19-Qsel}
\end{equation}
with the following properties.
\begin{enumerate}[label=\textup{(B\arabic*)},leftmargin=3em]
\item The input coordinates are genuine microscopic law/action coordinates;
      no claimed predictive Wilson coefficient is copied into the input list.
\item The output is one common local ESM action/source packet through the
      derivative/EFT order of the pilot block, including all contact and BV
      descendants required by the V5 contraction-complete state.
\item V16 physical calibration removes the declared experimental inputs but
      does not independently refit the pilot Wilson coordinates after
      \(\mathcal Q_{\rm sel}\) is evaluated.
\item The map has the source derivatives needed to form the V18 incidence
      matrix on the protected coefficientwise chart.
\end{enumerate}
After the common renormalization/calibration map $\mathcal R_{\mu_0}$, the
physical microscopic matching map is
\begin{equation}
 \mathfrak M_{\rm pilot}
 :=\Pi_{\rm phys}\circ\mathcal R_{\mu_0}\circ\mathcal Q_{\rm sel}.
 \label{eq:v19-Mpilot}
\end{equation}
\end{definition}

\begin{theorem}[Incidence is defined only after the selection bridge is populated]
\label{thm:v19-incidence-readiness}
Let $\eta$ be a genuine microscopic packet.  The V18 matrix
\begin{equation}
 \mathsf S^I{}_A
 =\ell_I\,D\mathfrak M_{\rm pilot}(\partial_{\eta^A})
 \label{eq:v19-ready-S}
\end{equation}
is mathematically defined only when a map
$\mathcal Q_{\rm sel}$ satisfying Definition~\ref{def:v19-selection-bridge}
has been supplied.  If the corpus supplies no such map, then:
\begin{enumerate}[label=\textup{(R\arabic*)},leftmargin=3em]
\item assigning $\mathsf S=0$ is unjustified and would create a spurious
      maximal predictive conormal;
\item reintroducing the pilot Wilson coefficients as coordinates of $\eta$
      restores the V18 identity columns and gives zero predictive conormal;
\item the correct status is \emph{unpopulated incidence}, not rank zero.
\end{enumerate}
\end{theorem}

\begin{proof}
Equation \eqref{eq:v19-ready-S} is the derivative of the composite map
\eqref{eq:v19-Mpilot}; without its first arrow there is no function to
differentiate.  Replacing absence of a function by the zero function adds the
extra assertion that the Wilson output is independent of every microscopic
coordinate and fixed to a known base value, which is not contained in the
premises.  It would therefore turn missing data into equations.

Conversely, if a Wilson coefficient is added directly to the microscopic
coordinate list, differentiation in that direction produces the identity
column proved in V18.  The three alternatives are exhaustive at the level of
the present rank compiler.
\end{proof}

\subsection{Corpus audit: which arrows are actually populated?}
\label{sec:v19-corpus-audit}

The present source corpus supports the following typed inventory.

\begin{table}[ht]
\centering
\small
\begin{tabular}{@{}p{0.22\linewidth}p{0.30\linewidth}p{0.38\linewidth}@{}}
\toprule
Native packet & Populated output & P4 status\\
\midrule
Race variables $(s,d,a_\rho)$ &
$\mathcal B=R\operatorname{diag}(s)R^T$, $\beta=Rd$, density/scale &
Fully populated calibration incidence; rank $10$, with three residual $\ker R$ directions.\\
Absolute calibration packet &
Clock, length, volume, action/pressure/measure anchors &
Physical calibration data; not by itself a Wilson-selection law.\\
Classical signed action/source writer &
Euler/action covector and same-record consistency tests &
Independent physical row; the current corpus does not derive its finite quantum Wilson coordinates from $\eta\in\ker R$.\\
Gauge Wilson clock/parameter &
Gauge source normalization/running benchmark on its declared branch &
Primarily a calibrated coupling/clock direction; not a loaded selector of the gravity--matter pilot quotient.\\
V18 pilot Wilson coordinates &
$\mathcal W_{\rm pilot}$ &
Populated only as independent finite matching data; retaining them as inputs gives the V18 identity block.\\
\bottomrule
\end{tabular}
\caption{V19 audit of actual microscopic/calibration arrows already present in the inherited Renewal corpus.}
\label{tab:v19-corpus-audit}
\end{table}

\begin{proposition}[Current P4b Wilson bridge is unpopulated]
\label{prop:v19-unpopulated-bridge}
Within the inherited V1--V18 plus Grand Tensor corpus used by this lineage, no
populated coefficientwise map of the form
\begin{equation}
 \eta\in\ker R
 \longmapsto
 \Pi_{\rm phys}S_{\rm loc}^{\rm match}(\eta)
 \in\mathcal W_{\rm pilot}
 \label{eq:v19-missing-eta-map}
\end{equation}
is supplied.  The available results instead either:
\begin{enumerate}[label=\textup{(C\arabic*)},leftmargin=3em]
\item decode/calibrate physical geometry from the race law;
\item test or transport a separately occurring signed action/source row; or
\item retain the pilot Wilson quotient as independent finite matching data.
\end{enumerate}
Therefore the P4b Wilson incidence matrix requested at the end of V18 remains
unpopulated rather than zero or full rank after the Wilson identity columns are
removed.
\end{proposition}

\begin{proof}
Theorem~\ref{thm:v19-race-incidence} explicitly supplies item (C1).  The
inherited classical closure architecture treats the complete action covector
and source/action synthesis error as separate same-record rows, which is item
(C2).  V18 proved item (C3) for the renormalized pilot Wilson quotient.  None
of these constructions defines the composite arrow
\eqref{eq:v19-missing-eta-map}.  The conclusion follows from
Theorem~\ref{thm:v19-incidence-readiness}.
\end{proof}

\begin{firewall}[A corpus audit is not a theorem that no deeper Renewal law can exist]
Proposition~\ref{prop:v19-unpopulated-bridge} is relative to the mathematical
objects actually loaded in the present manuscripts.  It does not prove that no
future microscopic Renewal principle can select Wilson coefficients.  It says
precisely that such a principle has not yet been supplied in a form on which
P4's rank calculation can operate.
\end{firewall}

\subsection{Conditional dimension promise of the residual race sector}
\label{sec:v19-dimension-promise}

Although the Wilson bridge is absent, the race audit identifies a sharply
bounded candidate domain.  Suppose, as a future additional hypothesis, that
all microscopic coordinates relevant to the chosen pilot Wilson block after P1
calibration factor through the residual race packet $\eta\in\mathbb R^3$.
Then
\begin{equation}
 \mathfrak M_{\rm race}:\mathbb R^3\longrightarrow\mathcal W_{\rm pilot}
 \label{eq:v19-race-M}
\end{equation}
is the complete pilot matching map.

\begin{theorem}[Conditional rank ceiling from the race sector]
\label{thm:v19-race-rank-ceiling}
Under the additional completeness hypothesis above,
\begin{equation}
 \boxed{
 \operatorname{rank}D\mathfrak M_{\rm race}\le3.}
 \label{eq:v19-race-rank-ceiling}
\end{equation}
Hence at a regular point the number of independent first-order pilot relations
obeys
\begin{equation}
 \boxed{
 \dim\mathcal P^*_{\rm pilot}
 \ge d_{\rm pilot}-3.}
 \label{eq:v19-race-promise}
\end{equation}
whenever $d_{\rm pilot}>3$.
\end{theorem}

\begin{proof}
The tangent domain of \eqref{eq:v19-race-M} is three-dimensional, so the rank
of its differential is at most three.  Apply the V16/V18 predictive-codimension
formula.
\end{proof}

\begin{firewall}[This is a promise, not a prediction]
Theorem~\ref{thm:v19-race-rank-ceiling} is conditional on the residual race
packet being complete for the pilot Wilson sector.  The present corpus does not
prove that hypothesis.  Additional microscopic action, matter, boundary or
source parameters could enlarge the domain and raise the rank.  No numerical
Wilson relation is asserted in V19 from the dimension count alone.
\end{firewall}

\subsection{The exact next calculational object: the source-selection derivative}
\label{sec:v19-source-selection-derivative}

Once a bridge is supplied, its incidence factors by the chain rule.  Let
$\iota_{\rm res}:\mathbb R^3\to T\Theta_{\rm race}$ denote the embedding
$\delta\eta\mapsto(0,N\delta\eta,0)$.  Then
\begin{equation}
 \boxed{
 \mathsf S
 =D_W\Pi_{\rm phys}\,
  D\mathcal R_{\mu_0}\,
  D\mathcal Q_{\rm sel}\,
  \iota_{\rm res},}
 \label{eq:v19-S-factorization}
\end{equation}
where the first factor means evaluation in the chosen Wilson coordinate basis.
Everything except
\begin{equation}
 \boxed{
 D\mathcal Q_{\rm sel}\,\iota_{\rm res}}
 \label{eq:v19-missing-derivative}
\end{equation}
is already structurally defined by V1--V18.  This derivative is therefore the
minimal upstream object whose population can create the first genuine
Renewal-specific Wilson prediction.

\begin{proposition}[Source-selection derivative is the unique new P4b payload]
\label{prop:v19-unique-payload}
Fix the V16 calibration scheme, the V18 pilot quotient, and the residual race
coordinates \eqref{eq:v19-eta}.  Any two microscopic models with the same
coefficientwise derivative
\eqref{eq:v19-missing-derivative} at the calibrated point give the same
linearized P4b incidence matrix and hence the same first-order predictive
conormal.  Conversely, changing that derivative can change the rank even when
all V1--V18 continuum/RG data are held fixed.
\end{proposition}

\begin{proof}
The first statement is immediate from the factorization
\eqref{eq:v19-S-factorization}.  For the converse, the remaining factors are
fixed finite linear maps; a change in the image of
$D\mathcal Q_{\rm sel}\iota_{\rm res}$ which is not annihilated by their
composition changes the column space of $\mathsf S$ and can therefore change
its rank or left nullspace.
\end{proof}

\subsection{A concrete loading protocol for P4c}
\label{sec:v19-loading-protocol}

The next calculation should not start from an arbitrary counterterm ansatz.  A
valid microscopic loading should proceed in the following order.

\begin{definition}[P4c same-record action-selection packet]
\label{def:v19-P4c-packet}
A P4c packet at $\mu_0$ consists of:
\begin{enumerate}[label=\textup{(P\arabic*)},leftmargin=3em]
\item the actual finite event/race law with the P1 geometry coordinates fixed
      and residual coordinates $\eta$ varied;
\item one signed local action/source writer evaluated on the \emph{same}
      microscopic records, including the matter and geometric source rows
      required by the pilot EFT order;
\item a declared finite subtraction/reference prescription shared across all
      $\eta$, with no independent refit of the pilot Wilson coefficients;
\item the V16 calibration functionals for the lower-weight experimental inputs;
\item the V18 quotient projection and Wilson coordinate readout.
\end{enumerate}
The packet is complete only if these data determine
$\mathcal Q_{\rm sel}(\eta)$ coefficientwise through the requested source and
loop order.
\end{definition}

\begin{theorem}[P4c packet turns the rank question into a finite derivative calculation]
\label{thm:v19-P4c-finite}
For every fixed coefficient packet
$\mathfrak q=(L,\Delta,r,N,\mathcal P_{\rm ext})$, a populated
Definition~\ref{def:v19-P4c-packet} reduces the first P4b rank test to finitely
many derivatives of finitely many local/source coefficients at the calibrated
point.  No finite-amplitude source ball, all-loop trajectory, or further
regulator-universality theorem is required.
\end{theorem}

\begin{proof}
At fixed $\mathfrak q$, the V5 contraction-complete bank, V16 calibration
functionals and V18 pilot quotient are finite.  The V2 reference-pivot theorem
reduces every inverse/source insertion to finite jets at the reference point.
Thus $D\mathcal Q_{\rm sel}\iota_{\rm res}$ requires only the finite coefficient
array through the declared loop/EFT/source orders.  Composition with the
finite maps in \eqref{eq:v19-S-factorization} produces a finite matrix.
\end{proof}

\subsection{Parallel validation track: use what is already computable}
\label{sec:v19-validation-track}

The absence of the local Wilson-selection bridge does not block all
phenomenological work.  V18 separated the nonanalytic momentum sector because
finite local counterterms cannot change it.  The earlier native matter--metric
continuum calculation already supplies a well-defined renormalized nonlocal
two-metric coefficient with a controlled lattice-to-continuum limit.  Its
nonanalytic continuum part is therefore a natural P7/P8 validation target while
P4c is being populated.

\begin{proposition}[Validation and prediction can proceed on different rows]
\label{prop:v19-validation-vs-prediction}
Computing a scheme-invariant nonanalytic continuum coefficient can validate the
native regulator, source normalization and RG/matching conventions even when
all local pilot Wilson coefficients remain free.  Such a calculation does not
replace the P4c rank test and should be labelled \emph{validation} unless the
Renewal microscopic model changes the low-energy field content or fixes an
otherwise free boundary datum.
\end{proposition}

\begin{proof}
V18 proved that local finite counterterms affect only analytic momentum
polynomials, so the coefficient of a linearly independent nonanalytic
structure is scheme invariant.  The same theorem also proved that the local
Wilson fibre is currently full-dimensional.  Therefore success on the
nonanalytic row tests the continuum implementation but does not reduce the
local matching dimension.
\end{proof}

\subsection{Machine-readable microscopic-incidence inventory}
\label{sec:v19-machine}

The companion file
\begin{center}
\texttt{Renewal\_P4b\_Microscopic\_Incidence\_Inventory\_V19.json}
\end{center}
records the populated race Jacobian, its rank/kernel, the residual candidate
packet, and the readiness status of the pilot Wilson bridge.  In particular,
it uses the status label
\begin{center}
\texttt{UNPOPULATED\_MAP\_NOT\_ZERO\_MATRIX}
\end{center}
for the current gravity--matter Wilson incidence.

\subsection{V19 anti-circularity audit}
\label{sec:v19-audit}

\begin{enumerate}[leftmargin=2.4em]
\item \textbf{A missing microscopic map is not replaced by the zero map.}
Theorem~\ref{thm:v19-incidence-readiness} keeps absence of data distinct from a
prediction of independence.
\item \textbf{Wilson coordinates are not smuggled back into the microscopic
packet.}  Doing so would recreate the V18 identity block and zero predictive
conormal.
\item \textbf{The first populated incidence is an actual native one.}
Theorem~\ref{thm:v19-race-incidence} differentiates the explicit race decoder,
not a hypothetical EFT ansatz.
\item \textbf{The residual three-dimensional packet is not declared redundant.}
Its Fisher cost is positive even though it is invisible to the decoded
geometry.
\item \textbf{The three-dimensional rank ceiling is explicitly conditional.}
Other genuine microscopic action/matter/source coordinates may enlarge the
selector domain.
\item \textbf{Same-record action incidence is required.}
The next bridge must evaluate the signed action/source packet on the actual
microscopic law; a separately fitted counterterm table does not populate P4b.
\item \textbf{RG integration remains downstream.}
Until \eqref{eq:v19-missing-derivative} is known, additional beta coefficients
transport free boundary data but do not answer the selection question.
\end{enumerate}

\section{V19 status ledger}
\label{sec:v19-status}

\subsection*{Proved in V19}
\begin{enumerate}[leftmargin=2.2em]
\item The explicit 13-parameter race law has a fully populated ten-output
geometry/calibration incidence matrix with rank ten.  Its kernel is exactly the
three-dimensional drift subspace $\ker R$.
\item The race-to-geometry map has zero predictive conormal: it is a calibration
map, not a source of relations among the ten decoded physical coordinates.
\item The three residual drift directions are genuine microscopic law
directions with positive Fisher cost, despite being invisible to the decoded
geometry.  They form the smallest currently explicit candidate selector packet
after P1 geometry calibration.
\item A V18 Wilson incidence matrix is defined only after a genuine microscopic
source-to-action bridge has been populated.  Missing bridge data are neither a
zero sensitivity matrix nor a prediction.
\item The current source corpus does not supply a coefficientwise map from the
residual race packet into the V18 gravity--matter Wilson quotient.  It instead
keeps geometry calibration, the signed action/source writer, and independent
Wilson matching data logically separate.
\item Conditional on the residual race packet being complete for a future
pilot Wilson bridge, its three-dimensional tangent gives
$\operatorname{rank}\mathsf S\le3$ and hence at least
$d_{\rm pilot}-3$ first-order predictive covectors when $d_{\rm pilot}>3$.
This is a rank promise, not a presently populated prediction.
\item The only new object needed to populate the linearized P4b rank is the
same-record source-selection derivative
$D\mathcal Q_{\rm sel}\iota_{\rm res}$; all downstream quotient/calibration
maps are already structurally available.
\item A fixed-order P4c packet reduces that source-selection problem to a finite
coefficient/source derivative calculation.  No further foundational continuum
or regulator-universality result is needed for the rank test.
\end{enumerate}

\subsection*{Inherited}
\begin{enumerate}[leftmargin=2.2em]
\item V1--V15: the coefficientwise continuum and its stronger regulator/fixed-
index extensions, now secondary to phenomenology.
\item V16: physical calibration, predictive quotient/conormal, native RG vector
and finite scheme matching.
\item V17: the first calibrated one-loop Standard-Model RG packet and the native
logarithmic-residue extraction equation.
\item V18: the negative pilot-Wilson predictivity theorem, the microscopic
incidence compiler, the no-codimension-from-RG theorem and the nonanalytic
validation sector.
\item The Grand Tensor race decoder and efficient-information packet, including
rank-three $R$, six independent dyadics $r_ar_a^T$, and independent density
calibration.
\end{enumerate}

\subsection*{Conditional}
\begin{enumerate}[leftmargin=2.2em]
\item The corpus-audit statement is relative to the mathematical objects
actually supplied in the current manuscripts; a future Renewal microscopic
action-selection theorem can change the status immediately.
\item The residual race packet is only a candidate complete selector.  V19 does
not prove that all higher-action/matter/source microscopic freedom factors
through $\eta\in\ker R$.
\item The numerical dimension $d_{\rm pilot}$ remains basis independent but
panel dependent, exactly as in V18.
\item The parallel nonanalytic calculation is a validation target unless an
additional Renewal-specific matching condition changes its low-energy content.
\item All statements remain coefficientwise at fixed finite perturbative,
source and EFT order and at positive reference/IR scale.
\end{enumerate}

\subsection*{Open beyond V19 --- primary phenomenology line}
\begin{enumerate}[leftmargin=2.2em]
\item Populate the P4c same-record action-selection bridge
$\mathcal Q_{\rm sel}$, beginning with the derivative along the three residual
race coordinates $\eta$.  This requires an actual common gravity--matter
signed action/source writer whose finite Wilson output is not independently
refitted after the microscopic law is varied.
\item Evaluate the finite matrix \eqref{eq:v19-S-factorization}.  If its left
nullspace is nonzero, convert the resulting relation to a standard
phenomenological basis and transport it with the V17 RG.  If the map remains
full row rank after all genuine microscopic coordinates are included, record a
negative phenomenology result and change the panel.
\item In parallel, extract one closed-form nonanalytic native continuum
coefficient from the already proved matter--metric amplitude and compare it to
the ordinary EFT benchmark as a scheme-invariant validation row.
\item Only after a populated nonzero predictive conormal is obtained should the
large anomalous-dimension bank, threshold matching and numerical trajectory
return to the top of the queue.
\end{enumerate}

\subsection*{Secondary mathematical branch}
The V15 global reserve descent, real Pfaffian orientation, index-changing
crossings, full IR removal, summability and nonperturbative completion remain
secondary and are not load-bearing for P4c.

\subsection*{Exact next load-bearing theorem}
\begin{center}
\fbox{\parbox{0.92\textwidth}{
\textbf{P4c: same-record microscopic action-selection derivative.}
Keep the P1 geometry and Standard-Model inputs fixed.  Vary the explicitly
identified residual race coordinates $\eta\in\ker R$ on the same native finite
records and evaluate one common signed gravity--matter action/source packet
through the V18 pilot EFT order, with no independent refit of the pilot Wilson
coefficients.  Prove coefficientwise existence of
$D\mathcal Q_{\rm sel}\iota_{\rm res}$, compose it with the already constructed
renormalization/calibration and physical quotient maps as in
\eqref{eq:v19-S-factorization}, and compute the rank/left nullspace.  If the
current microscopic law does not determine this derivative, the correct next
result is an explicit source-selection obstruction identifying which action
writer or incidence record is missing; do not replace it by a fitted Wilson
boundary condition.}}
\end{center}

\section*{Append-only continuation marker: V19}
\addcontentsline{toc}{section}{Append-only continuation marker: V19}
\textbf{End of V19.}  The complete V1--V18 body above is retained verbatim.
V19 goes upstream from the formal P4b rank compiler and audits which microscopic
arrows are actually populated.  The inherited 13-parameter race law provides a
fully explicit rank-ten incidence into ten geometry/calibration coordinates,
leaving a three-dimensional physically distinguishable drift packet in
$\ker R$ after P1 geometry calibration.  This is the smallest currently
explicit candidate microscopic selector packet.  However, the present corpus
keeps the signed action/source writer independent from the renewal-law decoder
and supplies no coefficientwise map from those residual race variables into
the V18 gravity--matter Wilson quotient.  V19 therefore proves that the P4b
Wilson matrix is presently unpopulated, not zero, constructs the exact missing
source-selection derivative, and reduces the next phenomenology gate to a
finite same-record action-incidence calculation.  A conditional three-column
rank ceiling shows why this bridge could become predictive, but no Wilson
relation is claimed before it is populated.



\clearpage
\section{V20 continuation: residual race skewness, decoded-action silence, and the exact nuisance-incidence compiler}
\label{sec:v20-continuation}

V19 isolated the missing P4c object as the same-record source-selection
derivative
\(D\mathcal Q_{\rm sel}\,\iota_{\rm res}\).  The first upstream question is
whether the three residual race coordinates really enter the already identified
native Einstein--Standard-Model action, or whether they belong to a finer
microscopic sector which the physical decoder has discarded.  The distinction
is decisive: a derivative of an absent map must not be set to zero, but a
populated map which factors through the decoded physical coordinates may have an
\emph{actually proved} zero derivative.

The inherited race law and first-jet action theorem are sufficient to answer
this at the classical native level.  The residual directions are not arbitrary
hidden parameters.  They are exactly the tetrahedral circulation directions
which preserve both the drift and diffusion tensors.  Their first nonzero
scheduler effect occurs at cubic derivative order and carries an explicit
factor \(h^2\).  By contrast, the identified first-jet Einstein--SM action
factors through the decoded fields
\((\mathcal B,\beta,\rho,\text{matter heads})\), so its coefficient array is
constant along the residual race fibre once P1 calibration is held fixed.

This continuation proves those statements exactly, identifies the three
residual coordinates with an injective cubic-moment tensor, and then uses the
inherited information-dual decomposition to show where any nonzero P4c
sensitivity could still reside.  It lies entirely in the source-orthogonal
``nuisance action'' covector.  Finally, V20 turns that abstract remainder into a
finite same-record covariance compiler: for any explicitly occurring action or
Wilson writer, the residual selector matrix is its cross covariance with the
three residual race scores, plus any direct writer derivative.  This is the
smallest data object which must now be populated.

\subsection{Exact residual race score and Fisher metric}
\label{sec:v20-residual-score}

Retain the V19 variables.  For \(\sigma\in\{+1,-1\}\), write the twelve
directed race rates as
\begin{equation}
 k_a^\sigma
 =\frac{s_a+\sigma h d_a}{2h^2},
 \qquad
 \mathcal K_h:=\sum_{a,\sigma}k_a^\sigma
 =\frac{\lambda}{h^2},
 \qquad
 \lambda:=\sum_as_a.
 \label{eq:v20-directed-rates}
\end{equation}
The winner/waiting-time density is
\begin{equation}
 p(a,\sigma,t)=k_a^\sigma e^{-\mathcal K_ht}.
 \label{eq:v20-race-density}
\end{equation}
Choose the V19 matrix \(N\in\mathbb R^{6\times3}\) with orthonormal columns
spanning \(\ker R\), and write
\begin{equation}
 d=R^+\beta+N\eta.
 \label{eq:v20-d-eta}
\end{equation}
For the \(\alpha\)-th residual coordinate set
\(n_{a\alpha}:=N_{a\alpha}\).

\begin{theorem}[Exact residual score bank]
\label{thm:v20-residual-score}
At fixed \((s,\beta,a_\rho)\), the score of the complete winner/waiting-time
experiment in the residual coordinate \(\eta^\alpha\) is
\begin{equation}
 \boxed{
 \mathfrak s_\alpha(a,\sigma,t)
 =\frac{\sigma h\,n_{a\alpha}}
        {s_a+\sigma h d_a}.}
 \label{eq:v20-residual-score}
\end{equation}
It is independent of the waiting time because the total race intensity is
independent of \(d\), and it has mean zero.  Its Fisher matrix is
\begin{equation}
 \boxed{
 G_\eta
 =\frac{h^2}{\lambda}\,
 N^T\bigl(S-h^2DS^{-1}D\bigr)^{-1}N\succ0.}
 \label{eq:v20-residual-Fisher}
\end{equation}
Thus the residual circulation packet remains a faithful three-dimensional
microscopic experiment after P1 geometry calibration.
\end{theorem}

\begin{proof}
Differentiate \(\log p\).  Since
\(\partial_{\eta^\alpha}\mathcal K_h=0\), only the winning-rate term remains:
\[
 \partial_{\eta^\alpha}\log k_a^\sigma
 =\frac{\sigma h\,n_{a\alpha}}{s_a+\sigma h d_a}.
\]
For fixed \(a\), summing the two signs against their winner probabilities
\((s_a+\sigma hd_a)/(2\lambda)\) cancels, proving the zero mean.

For the Fisher matrix,
\begin{align*}
 \mathbb E(\mathfrak s_\alpha\mathfrak s_\beta)
 &=\sum_{a,\sigma}
 \frac{s_a+\sigma hd_a}{2\lambda}
 \frac{h^2n_{a\alpha}n_{a\beta}}
 {(s_a+\sigma hd_a)^2}\\
 &=\frac{h^2}{\lambda}
 \sum_a\frac{n_{a\alpha}n_{a\beta}}
 {s_a-h^2d_a^2/s_a},
\end{align*}
which is exactly \eqref{eq:v20-residual-Fisher}.  The protected rate
inequalities make the diagonal matrix in the inverse positive definite, and
\(N\) has full column rank.
\end{proof}

\begin{corollary}[Whitened residual scores]
\label{cor:v20-whitened-scores}
The score synthesis
\begin{equation}
 \widehat{\mathfrak s}
 :=\mathfrak s\,G_\eta^{-1/2}
 \label{eq:v20-whitened-score}
\end{equation}
has identity Fisher Gram.  Any invertible reparameterization of \(\eta\)
changes the unwhitened selector matrix only by an invertible right factor and
therefore leaves its left nullspace and rank invariant.
\end{corollary}

\subsection{Residual race directions are exactly cubic scheduler skewness}
\label{sec:v20-cubic-skewness}

The residual packet has a more intrinsic interpretation than ``three unused
drift coordinates''.  Let the scheduler act on a smooth test function by
\begin{equation}
 (\mathcal L_hf)(x)
 :=\sum_{a=1}^6\sum_{\sigma=\pm1}
 k_a^\sigma\bigl[f(x+\sigma hr_a)-f(x)\bigr].
 \label{eq:v20-race-generator}
\end{equation}
The first two Kramers--Moyal tensors are exactly the already calibrated
\(\beta=Rd\) and \(\mathcal B=RSR^T\).

\begin{theorem}[Residual generator begins two lattice powers beyond the drift]
\label{thm:v20-residual-generator}
For \(v\in\mathbb R^3\) put \(\delta d=Nv\).  Then
\begin{equation}
 \boxed{
 D_\eta\mathcal L_h[v]f
 =\frac1{2h}\sum_a\delta d_a
 \bigl[f(x+hr_a)-f(x-hr_a)\bigr].}
 \label{eq:v20-exact-residual-generator}
\end{equation}
Because \(R\delta d=0\), the first-derivative term cancels.  For
\(f\in C^5\) on a common compact chart,
\begin{equation}
 \boxed{
 D_\eta\mathcal L_h[v]f
 =\frac{h^2}{6}\sum_a\delta d_a(r_a\!\cdot\!\nabla)^3f
 +\mathcal R_{h,v}f,}
 \label{eq:v20-cubic-generator}
\end{equation}
with
\begin{equation}
 \|\mathcal R_{h,v}f\|_\infty
 \le C h^4\|v\|\,\|f\|_{C^5}.
 \label{eq:v20-cubic-remainder}
\end{equation}
Equivalently, its flat Fourier symbol is
\begin{equation}
 \boxed{
 \widehat{D_\eta\mathcal L_h[v]}(p)
 =\frac{i}{h}\sum_a\delta d_a\sin(h r_a\!\cdot p)
 =-\frac{i h^2}{6}\sum_a\delta d_a(r_a\!\cdot p)^3
 +O(h^4|p|^5\|v\|).}
 \label{eq:v20-cubic-symbol}
\end{equation}
Thus the residual race fibre is invisible to the physical drift and diffusion
but survives as a cutoff-suppressed cubic scheduler tensor.
\end{theorem}

\begin{proof}
Differentiate \eqref{eq:v20-race-generator} using
\(\partial_d k_a^\sigma=\sigma/(2h)\), proving
\eqref{eq:v20-exact-residual-generator}.  The centered Taylor formula gives
\[
 \frac{f(x+hr)-f(x-hr)}{2h}
 =(r\cdot\nabla)f
 +\frac{h^2}{6}(r\cdot\nabla)^3f
 +O(h^4\|f\|_{C^5}).
\]
The linear term summed against \(\delta d_a\) is
\((R\delta d)\cdot\nabla f=0\).  The root set is fixed, so the finitely many
fifth-directional-derivative constants are absorbed into \(C\).  The Fourier
formula is the exact multiplier of
\eqref{eq:v20-exact-residual-generator}; expand the sine and use the same
cancellation.
\end{proof}

For the explicit \(A_3\) roots of the inherited scheduler, the cubic map is
itself faithful on the circulation space.

\begin{theorem}[The cubic-moment map is injective on tetrahedral circulation]
\label{thm:v20-cubic-injective}
Let
\begin{align*}
 r_1&=e_1,&r_2&=e_2,&r_3&=e_3,\\
 r_4&=e_2-e_1,&r_5&=e_3-e_1,&r_6&=e_3-e_2.
\end{align*}
Every \(\delta d\in\ker R\) is uniquely
\begin{equation}
 \delta d
 =(a+b,-a+c,-b-c,a,b,c).
 \label{eq:v20-cycle-param}
\end{equation}
Define the cubic polynomial encoding its symmetric third moment by
\begin{equation}
 \mathfrak T_3(\delta d;x)
 :=\sum_{j=1}^6\delta d_j(r_j\cdot x)^3.
 \label{eq:v20-T3}
\end{equation}
Then
\begin{equation}
 \boxed{
 \mathfrak T_3(\delta d;x)
 =3a\,x_1x_2(x_1-x_2)
 +3b\,x_1x_3(x_1-x_3)
 +3c\,x_2x_3(x_2-x_3).}
 \label{eq:v20-T3-explicit}
\end{equation}
The three displayed cubic polynomials are linearly independent.  Hence
\begin{equation}
 \boxed{
 \delta d\in\ker R,
 \quad \mathfrak T_3(\delta d)=0
 \Longrightarrow \delta d=0.}
 \label{eq:v20-T3-injective}
\end{equation}
Thus the three V19 residual coordinates are equivalently the three independent
cubic circulation/skewness coordinates of the scheduler.
\end{theorem}

\begin{proof}
The equation \(R\delta d=0\) gives
\[
 \delta d_1=\delta d_4+\delta d_5,
 \quad
 \delta d_2=-\delta d_4+\delta d_6,
 \quad
 \delta d_3=-\delta d_5-\delta d_6,
\]
which is \eqref{eq:v20-cycle-param}.  Substitute into
\eqref{eq:v20-T3}.  The \(a\), \(b\), and \(c\) parts reduce respectively to
\[
 x_1^3-x_2^3+(x_2-x_1)^3,
 \quad
 x_1^3-x_3^3+(x_3-x_1)^3,
 \quad
 x_2^3-x_3^3+(x_3-x_2)^3,
\]
which are the three terms in \eqref{eq:v20-T3-explicit}.  Their monomial
supports contain respectively \(x_1^2x_2\), \(x_1^2x_3\), and
\(x_2^2x_3\) with nonzero coefficients not shared in the same pattern by the
other two, so no nontrivial linear combination vanishes.
\end{proof}

\begin{remark}[Correct microscopic interpretation of the V19 residual packet]
The P1-invisible race directions are not generic hidden coordinates.  They are
the scheduler's tetrahedral circulation sector, equivalently its first
unfixed odd Kramers--Moyal tensor.  Any microscopic selection principle which
uses only the physical first two scheduler moments necessarily annihilates
this packet.  A nonzero dependence must see finer same-record information,
such as the cubic circulation tensor or an independently occurring action
writer correlated with its score.
\end{remark}

\subsection{The inherited first-jet Einstein--SM action is exactly silent along the residual fibre}
\label{sec:v20-decoded-action-silence}

The Grand Tensor first-jet action theorem supplies an independently identified
native action of the form
\begin{equation}
 S_h[y]
 =h^4\sum_nF_h^{(1)}
 \bigl((T_\sigma y,T_\sigma\delta_\mu^+y)_{(\sigma,\mu)\in I_1}\bigr)
 +\mathcal B_h[y],
 \label{eq:v20-first-jet-action}
\end{equation}
with the minimal gravitational, gauge, Higgs and bilinear Dirac--Yukawa
sectors and fixed calibrated coefficients.  Its native decoder depends on the
race packet only through
\begin{equation}
 \mathcal B=RSR^T,
 \qquad
 \beta=Rd,
 \qquad
 \rho,
 \label{eq:v20-decoder-moments}
\end{equation}
plus separately retained matter heads.

\begin{definition}[Decoded minimal-action bridge]
\label{def:v20-decoded-bridge}
Hold the P1 geometry, Standard-Model inputs and matter heads fixed and define
\begin{equation}
 \mathcal Q_{\rm dec}(\eta)
 :=S_h\!\left[\mathcal D_{\rm dec}
 (s,R^+\beta+N\eta,\rho;\text{matter})\right].
 \label{eq:v20-Qdec}
\end{equation}
This is a genuinely populated microscopic-to-action map, but only for the
identified minimal first-jet action itself; it does not include an independently
chosen higher-Wilson boundary functional.
\end{definition}

\begin{theorem}[Exact decoded-action silence]
\label{thm:v20-decoded-action-silence}
On the V19 residual fibre,
\begin{equation}
 \boxed{
 D_\eta\mathcal D_{\rm dec}=0,
 \qquad
 D_\eta\mathcal Q_{\rm dec}=0.}
 \label{eq:v20-decoded-action-zero}
\end{equation}
The same statement holds after every fixed source derivative already present in
the decoded first-jet action packet.  Thus the literal currently identified
minimal Einstein--SM action does not use the three residual circulation
coordinates as classical coupling selectors.
\end{theorem}

\begin{proof}
Along \(d=R^+\beta+N\eta\), one has
\(R d=\beta\) because \(RN=0\).  The symmetric rates \(s\) and density
\(\rho\) are fixed, so \(\mathcal B\), \(\beta\), \(\rho\), and the declared
matter heads are independent of \(\eta\).  Hence the decoded field record is
constant along the residual fibre.  Equation \eqref{eq:v20-Qdec} is therefore
constant as well.  Fixed external source differentiation commutes with this
parameter derivative because the source variables are independent coordinates
in the finite action packet.
\end{proof}

\begin{firewall}[A proved zero derivative here is not the missing quantum Wilson matrix]
Theorem~\ref{thm:v20-decoded-action-silence} concerns the populated minimal
first-jet action map.  The continuum construction additionally carries
cutoff-dependent local restoration and independent ghost-zero Wilson matching
coordinates.  An irrelevant microscopic circulation insertion can, in
principle, feed finite \emph{local} Wilson matching data through ultraviolet
loops.  V20 therefore does not set the full quantum P4c matrix to zero from
\eqref{eq:v20-decoded-action-zero}.  It proves only that any such nonzero
selection is not contained in the already decoded classical action.
\end{firewall}

\subsection{Information-dual factorization: only the nuisance action covector can see the residual fibre}
\label{sec:v20-information-dual}

The inherited race-information theorem already contains the correct dual
language.  Let \(I\succ0\) be the complete microscopic information matrix and
let
\begin{equation}
 J:=D\mathcal G_{\rm race}
 \label{eq:v20-J}
\end{equation}
be the surjective physical calibration tangent map.  Put
\begin{equation}
 Q:=(JI^{-1}J^*)^{-1},
 \qquad
 L:=I^{-1}J^*Q,
 \qquad
 JL=I.
 \label{eq:v20-information-lift}
\end{equation}
For any scalar action/source covector \(e\) the inherited orthogonal
decomposition is
\begin{equation}
 e=J^*E+e_\perp,
 \qquad
 E=L^*e,
 \qquad
 L^*e_\perp=0.
 \label{eq:v20-action-dual-decomp}
\end{equation}

\begin{theorem}[Residual action derivative equals the nuisance covector]
\label{thm:v20-nuisance-only}
Let \(\iota_{\rm res}v=(0,Nv,0)\) be the V19 embedding of the residual race
tangent.  Then every action/source covector satisfies
\begin{equation}
 \boxed{
 e\,\iota_{\rm res}
 =e_\perp\,\iota_{\rm res}.}
 \label{eq:v20-nuisance-only}
\end{equation}
In particular, every covector which factors entirely through the calibrated
physical variables, \(e=J^*E\), annihilates the residual race packet.
For a finite bank of action/Wilson writers this holds row by row; hence the
entire residual selector matrix is the restriction of the bank's nuisance
covector block to \(\ker J\).
\end{theorem}

\begin{proof}
By definition \(J\iota_{\rm res}=0\).  Therefore
\[
 (J^*E)\iota_{\rm res}
 =E(J\iota_{\rm res})=0.
\]
Apply \eqref{eq:v20-action-dual-decomp}.  The bank-valued statement follows by
applying the scalar identity to every output covector.
\end{proof}

\begin{corollary}[The decoded first-jet action has zero nuisance block]
\label{cor:v20-decoded-no-nuisance}
For the bridge \(\mathcal Q_{\rm dec}\) of
Definition~\ref{def:v20-decoded-bridge}, the corresponding differential
factors through \(J\).  Hence its nuisance covector block vanishes and
Theorem~\ref{thm:v20-decoded-action-silence} is recovered from
Theorem~\ref{thm:v20-nuisance-only}.
\end{corollary}

\subsection{A finite same-record compiler for the missing nuisance incidence}
\label{sec:v20-same-record-compiler}

The missing P4c data can now be written in an experimentally and numerically
loadable form.  Let \(\omega\) denote one complete finite native record and
let \(p_\eta(\omega)\) be a normalized same-support family carrying the
residual race deformation.  Let
\begin{equation}
 \mathcal A_I(\omega;\eta),
 \qquad I=1,\ldots,d_{\rm pilot},
 \label{eq:v20-writer-bank}
\end{equation}
be an actually occurring finite bank of signed action/source or Wilson readout
writers on that same record.  The writers may include direct parameter
dependence; no positivity is assumed.

\begin{theorem}[Exact score--action incidence formula]
\label{thm:v20-score-action-incidence}
At the calibrated base point,
\begin{equation}
 \boxed{
 \partial_{\eta^\alpha}
 \mathbb E_\eta[\mathcal A_I]
 =\mathbb E[\partial_{\eta^\alpha}\mathcal A_I]
 +\operatorname{Cov}(\mathcal A_I,\mathfrak s_\alpha).}
 \label{eq:v20-score-action-incidence}
\end{equation}
Define the finite matrices
\begin{align}
 \mathsf D_{I\alpha}
 &:=\mathbb E[\partial_{\eta^\alpha}\mathcal A_I],
 \label{eq:v20-direct-action-matrix}\\
 \mathsf C_{I\alpha}
 &:=\operatorname{Cov}(\mathcal A_I,\mathfrak s_\alpha),
 \label{eq:v20-cov-action-matrix}\\
 \mathsf A_{\rm res}
 &:=\mathsf D+\mathsf C.
 \label{eq:v20-residual-action-matrix}
\end{align}
Then \(\mathsf A_{\rm res}\) is the exact linearized same-record action
incidence of the pilot writer bank with the three residual race directions.
Its Fisher-whitened form
\begin{equation}
 \boxed{
 \widehat{\mathsf A}_{\rm res}
 :=\mathsf A_{\rm res}G_\eta^{-1/2}}
 \label{eq:v20-whitened-incidence}
\end{equation}
has the same rank and left nullspace as \(\mathsf A_{\rm res}\) and is
invariant under invertible residual-coordinate changes up to a right orthogonal
factor after whitening.
\end{theorem}

\begin{proof}
Differentiate the finite normalized expectation:
\[
 \partial_\alpha\sum_\omega p_\eta(\omega)\mathcal A_I(\omega;\eta)
 =\sum_\omega p\,\partial_\alpha\mathcal A_I
  +\sum_\omega p\,\mathcal A_I\,\partial_\alpha\log p.
\]
The second term is \(\mathbb E[\mathcal A_I\mathfrak s_\alpha]\).
The score has zero mean by Theorem~\ref{thm:v20-residual-score}, so this is the
covariance.  The remaining statements are finite-dimensional linear algebra;
right multiplication by an invertible matrix does not change the image or left
nullspace, and Fisher whitening chooses a canonical Euclidean metric on the
residual score space.
\end{proof}

\begin{corollary}[Pure occurrence response is a cross Gram]
\label{cor:v20-cross-Gram}
If the writer bank itself has no direct residual dependence after P1
calibration, \(\mathsf D=0\), then
\begin{equation}
 \boxed{
 \mathsf A_{\rm res}=\mathsf C
 =\bigl[\langle\mathcal A_I-\mathbb E\mathcal A_I,
                 \mathfrak s_\alpha\rangle_{L^2(p)}\bigr]_{I,\alpha}.}
 \label{eq:v20-cross-Gram}
\end{equation}
Hence P4c reduces to one signed writer--score cross Gram on the same native
records.
\end{corollary}

\begin{proposition}[Action-work drift is one concrete loader of the cross Gram]
\label{prop:v20-action-work-loader}
For a normalized transition row \(P_\eta(x,y)\) and a same-record action value
\(a_h\) independent of \(\eta\), define
\begin{equation}
 w(x,y):=a_h(y)-a_h(x),
 \qquad
 d_\eta(x):=\sum_yP_\eta(x,y)w(x,y).
 \label{eq:v20-action-work}
\end{equation}
If
\(\mathfrak s_\alpha(x,y)=\partial_{\eta^\alpha}\log P_\eta(x,y)|_0\), then
\begin{equation}
 \boxed{
 \partial_{\eta^\alpha}d_\eta(x)|_0
 =\mathbb E\!\left[w(X,Y)\mathfrak s_\alpha(X,Y)\mid X=x\right].}
 \label{eq:v20-action-work-score}
\end{equation}
Thus the inherited same-record action-work diagnostic is precisely a conditional
version of the incidence compiler above once the residual score is included in
the transition law.
\end{proposition}

\begin{proof}
Differentiate the finite transition sum.  Normalization gives zero conditional
mean score, so no separate centering term is required.
\end{proof}

\subsection{P4c factorization after the upstream reduction}
\label{sec:v20-P4c-factorization}

Let
\begin{equation}
 \mathsf K_{\rm phys}
 :=D_W\Pi_{\rm phys}\,D\mathcal R_{\mu_0}
 \label{eq:v20-Kphys}
\end{equation}
be the already constructed downstream renormalization/calibration/Wilson
readout differential.  Whenever the same-record writer bank in
\eqref{eq:v20-writer-bank} is chosen to be the local matching-action packet fed
into that downstream map, the V19 factorization becomes
\begin{equation}
 \boxed{
 \mathsf S
 =\mathsf K_{\rm phys}\,\mathsf A_{\rm res}.}
 \label{eq:v20-S-factorization}
\end{equation}
This identifies the missing P4c object more sharply than
\(D\mathcal Q_{\rm sel}\iota_{\rm res}\): it is a finite bank of direct
writer derivatives plus same-record score--writer cross covariances.

\begin{theorem}[Decoded-action branch gives a genuine zero selector but no prediction]
\label{thm:v20-zero-selector-no-prediction}
If the only microscopic action bridge retained is the populated decoded
minimal bridge \(\mathcal Q_{\rm dec}\), then
\begin{equation}
 \boxed{
 \mathsf A_{\rm res}=0,
 \qquad
 \mathsf S=0.}
 \label{eq:v20-zero-selector}
\end{equation}
This zero is now a proved property of that bridge, not a placeholder for missing
data.  Nevertheless the current continuum theory still contains the
independent V18 physical Wilson matching coordinates.  Including those genuine
coordinates in the complete microscopic/matching domain restores the V18
identity columns and full row rank.  Therefore
\begin{equation}
 \boxed{
 \text{the decoded-action zero selector does not create a Renewal prediction.}}
 \label{eq:v20-zero-no-prediction}
\end{equation}
\end{theorem}

\begin{proof}
The first statement is Theorem~\ref{thm:v20-decoded-action-silence} followed by
\eqref{eq:v20-S-factorization}.  For the second, V18 proved that differentiating
the physical Wilson outputs with respect to their independently retained finite
matching coordinates gives an identity block.  Appending zero columns from
\(\eta\) to a matrix which already contains that identity block cannot reduce
its row rank or create a left-null covector.
\end{proof}

\begin{firewall}[Do not confuse scheduler irrelevance with quantum matching silence]
The cubic residual scheduler insertion carries an explicit \(h^2\) factor at
fixed soft momentum, but it is not uniformly small at loop momentum
\(|p|\asymp h^{-1}\).  Ultraviolet graphs can therefore convert this
microscopic difference into finite \emph{local} matching shifts.  Such shifts
are exactly the object measured by \(\mathsf A_{\rm res}\) after the complete
same-record action/coarse-graining packet is specified.  V20 proves no theorem
that all quantum Wilson coefficients are independent of \(\eta\).  It proves
that this dependence is not present in the already decoded classical action and
must enter through an additional nuisance/action-incidence or ultraviolet
matching mechanism.
\end{firewall}

\subsection{What V20 changes in the phenomenology programme}
\label{sec:v20-phenomenology-handoff}

V19 left the residual race packet as a generic three-dimensional candidate
selector.  V20 resolves its physical meaning and narrows the missing payload:
\begin{equation}
 \boxed{
 \eta
 \longleftrightarrow
 \text{tetrahedral cubic circulation tensor}
 \longrightarrow
 \mathsf A_{\rm res}
 \longrightarrow
 \mathsf S.}
 \label{eq:v20-selector-chain}
\end{equation}
The first arrow is now explicit and injective.  The second arrow is zero for the
literal decoded minimal action and is otherwise exactly the direct-plus-cross-
Gram matrix of Theorem~\ref{thm:v20-score-action-incidence}.  The third arrow is
already supplied by V16--V19.

This gives a terminating decision rule for the next loading:
\begin{enumerate}[label=\textup{(D\arabic*)},leftmargin=3em]
\item if an actual Wilson/source writer bank is loaded on the same residual-race
      records, compute \(\mathsf D\), \(\mathsf C\), and the rank of
      \(\mathsf K_{\rm phys}(\mathsf D+\mathsf C)\);
\item if the resulting left nullspace is nonzero after all genuine microscopic
      action coordinates are included, transport those relations with the V17
      RG and convert them to a standard phenomenological basis;
\item if the bank is action-orthogonal to the residual scores and has no direct
      dependence, record the residual race sector as a regulator/occurrence
      direction rather than a selector and stop spending P4 effort on it;
\item if no same-record writer bank exists, the remaining obstruction is now a
      finite acquisition problem: populate the three score--action cross rows,
      not another continuum-existence theorem.
\end{enumerate}

\section{V20 anti-circularity audit}
\label{sec:v20-audit}

\begin{enumerate}[leftmargin=2.4em]
\item \textbf{The V19 missing map was not silently set to zero.}
V20 first constructs the actually populated decoded-action bridge and proves
its derivative is zero.  It separately retains the unpopulated nuisance/action
incidence which could generate quantum Wilson matching.
\item \textbf{The residual race packet is identified intrinsically.}
Its score/Fisher metric is explicit and its first unfixed scheduler moment is
an injective cubic tensor; it is not merely ``three parameters left over by
rank--nullity''.
\item \textbf{Classical action silence is proved from factorization through the
physical decoder.}
No statement about the quantum Wilson quotient is inferred from that fact.
\item \textbf{Same-record selection is a measurable finite object.}
The exact missing linear response is a direct action derivative plus a signed
score--writer cross covariance, not an arbitrary fitted counterterm matrix.
\item \textbf{An \(h^2\) soft expansion is not summed into the ultraviolet.}
The firewall above keeps cutoff-scale matching separate, where the exact sine
symbol rather than the soft Taylor expansion must be used.
\item \textbf{Independent Wilson coordinates remain independent until a
microscopic selection law removes them.}
Zero residual sensitivity does not reduce the V18 identity block and therefore
does not create a prediction.
\end{enumerate}

\section{V20 status ledger}
\label{sec:v20-status}

\subsection*{Proved in V20}
\begin{enumerate}[leftmargin=2.2em]
\item The complete residual race score is
\(\mathfrak s_\alpha=\sigma h n_{a\alpha}/(s_a+\sigma hd_a)\), with positive
three-dimensional Fisher matrix
\(G_\eta=(h^2/\lambda)N^T(S-h^2DS^{-1}D)^{-1}N\).
\item The residual scheduler derivative has the exact centered-difference
form \eqref{eq:v20-exact-residual-generator}.  Since \(RN=0\), its drift term
cancels and the first soft contribution is the \(h^2\) cubic tensor in
\eqref{eq:v20-cubic-generator}.
\item On the explicit \(A_3\) root system the cubic-moment map is injective on
\(\ker R\).  In cycle coordinates its three basis polynomials are
\(3x_1x_2(x_1-x_2)\), \(3x_1x_3(x_1-x_3)\), and
\(3x_2x_3(x_2-x_3)\).
\item The populated native first-jet Einstein--SM action and decoder are
exactly constant along the residual race fibre after P1 geometry calibration.
Thus the residual race packet is not a classical coupling selector in the
literal decoded-action branch.
\item For every action/source covector, its derivative on the residual race
fibre is exactly the restriction of the source-orthogonal nuisance covector.
All physical-decoder covectors annihilate the residual packet.
\item For any finite same-record signed writer bank, the exact residual
selector matrix is
\(\mathsf A_{\rm res}=\mathsf D+\mathsf C\), where \(\mathsf D\) is the
direct writer derivative and \(\mathsf C\) is the score--writer covariance.
Its Fisher-whitened form has the same rank and predictive left nullspace.
\item The V19 P4c matrix factors as
\(\mathsf S=\mathsf K_{\rm phys}\mathsf A_{\rm res}\).  On the decoded
minimal-action branch this matrix is genuinely zero, but the independent V18
Wilson matching coordinates retain full row rank and therefore no prediction is
created.
\end{enumerate}

\subsection*{Inherited}
\begin{enumerate}[leftmargin=2.2em]
\item V1--V15: coefficientwise continuum existence and its stronger regulator
extensions, now secondary to the phenomenology line.
\item V16: physical calibration, predictive quotient/conormal and finite-scheme
RG geometry.
\item V17: the first calibrated one-loop SM/mixed RG packet.
\item V18: the negative free-Wilson predictivity theorem and physical Wilson
quotient.
\item V19: the fully populated rank-ten race-to-geometry incidence, the
three-dimensional residual microscopic packet and the P4c factorization.
\item The Grand Tensor first-jet minimal action, race-information dual
Pythagoras and same-record action-work diagnostic.
\end{enumerate}

\subsection*{Conditional}
\begin{enumerate}[leftmargin=2.2em]
\item The same-record covariance compiler becomes a Wilson-selection theorem
only when the writers \(\mathcal A_I\) are the actual local matching-action
coordinates feeding the V18 physical quotient.  An arbitrary diagnostic writer
is not automatically a Wilson coefficient.
\item The residual race packet need not be the complete microscopic selector
for a future action theory.  Additional microscopic action or matter variables
can enlarge the domain and the rank.
\item The \(h^2\) cubic expansion is a soft-momentum statement.  Full-zone
ultraviolet matching must use the exact native symbol.
\item The V12 real Pfaffian condition and all positive-IR/fixed-loop firewalls
remain unchanged.
\end{enumerate}

\subsection*{Open beyond V20 --- primary phenomenology line}
\begin{enumerate}[leftmargin=2.2em]
\item Populate the finite same-record matrix
\(\mathsf A_{\rm res}=\mathsf D+\mathsf C\) for an actual gravity--matter
matching writer bank.  Only three residual scores are required.  This is now a
finite direct-action/cross-covariance measurement or exact finite calculation.
\item If the first-order matrix vanishes on a symmetry-preserving reference
branch, compute the second residual Hessian before abandoning the sector; the
quadratic response can contain information invisible to the linear conormal.
\item In parallel, extract one scheme-invariant nonanalytic native
matter--metric coefficient as the first observable validation benchmark.
\item Once a nonzero predictive conormal is actually populated, resume P3/P5:
standard-scheme conversion, threshold matching and numerical RG transport of
that relation.
\end{enumerate}

\subsection*{Secondary mathematical branch}
Global reserve descent, real Pfaffian orientation, index-changing crossings,
full IR removal, perturbative summability and nonperturbative completion remain
secondary to the P4/P7 phenomenology line.

\subsection*{Exact next load-bearing theorem}
\begin{center}
\fbox{\parbox{0.92\textwidth}{
\textbf{P4d: populate the residual score--action cross Gram.}
Choose a smallest actual gravity--matter matching writer bank on the same
finite records as the residual race experiment.  For the three explicit scores
\(\mathfrak s_\alpha\), compute the direct derivatives
\(\mathsf D_{I\alpha}\) and same-record covariances
\(\mathsf C_{I\alpha}=\operatorname{Cov}(\mathcal A_I,\mathfrak s_\alpha)\).
Feed \(\mathsf A_{\rm res}=\mathsf D+\mathsf C\) through the already defined
renormalization/calibration quotient map \(\mathsf K_{\rm phys}\), then compute
the rank and left nullspace of
\(\mathsf S=\mathsf K_{\rm phys}\mathsf A_{\rm res}\).
If the required writer bank does not occur on the same records, return that
specific occurrence/incidence failure.  If the first-order matrix vanishes by a
proved symmetry, promote the next calculation to the finite second residual
Hessian rather than fitting a Wilson boundary condition.}}
\end{center}

\section*{Append-only continuation marker: V20}
\addcontentsline{toc}{section}{Append-only continuation marker: V20}
\textbf{End of V20.}  The complete V1--V19 body above is retained verbatim.
V20 resolves the physical meaning of the three P1-invisible race coordinates:
they are exactly the tetrahedral circulation sector and are faithfully encoded
by the scheduler's first unfixed cubic Kramers--Moyal tensor.  The populated
first-jet Einstein--SM action factors only through the decoded drift/diffusion,
so its residual derivative is exactly zero.  Information-dual decomposition
then shows that any nonzero P4c sensitivity must live in a source-orthogonal
nuisance action covector.  V20 turns that remainder into an exact finite
same-record compiler: direct writer derivatives plus score--writer cross
covariances.  The decoded-action zero does not create a prediction because the
independent V18 Wilson matching coordinates remain.  The next phenomenology
gate is therefore to populate the three-column score--action incidence matrix
for an actual gravity--matter matching writer bank.


\clearpage
\section{V21 continuation: first-jump selector compression, cycle incidence, and the same-record common-action obstruction}
\label{sec:v21-continuation}

V20 identified the first residual microscopic Renewal coordinates after physical
calibration: the three-dimensional circulation packet
\(\eta\in\ker R\), with a positive score Fisher metric and a first soft
appearance in the cubic Kramers--Moyal tensor.  It also reduced the desired
phenomenological selector to a finite same-record score--action incidence
matrix.  The present continuation attacks that matrix before introducing a
larger observable bank.

There is an upstream simplification and an upstream obstruction.  On one
resolved first-jump race, the residual score has an exact cancellation against
the event probability.  Consequently its covariance with an arbitrary
same-event writer depends only on six reversal contrasts and, after projection
to \(\ker R\), on three cycle coordinates per writer.  Thus the first-order
P4d acquisition problem is much smaller than a generic twelve-cell joint table.
On the other hand, the inherited physical-population programme requires the
full-field common-action score to occur on the \emph{same} first-jump cylinder.
The presently supplied corpus provides the compiler for this row but does not
supply the complete nonlinear Store--SM--gravity action writer on that selected
race cylinder.  That row is therefore upstream unloaded; it must not be
replaced by zero.

The second correction concerns V20's suggested Hessian fallback.  If a
first-jump writer has no explicit \(\eta\)-dependence, its expectation is
\emph{exactly affine} in \(\eta\), because the total hazard is fixed and the
winner probabilities are affine.  Hence its entire residual Hessian vanishes.
A zero first derivative on such a writer cannot be rescued by differentiating
the same one-jump statistic twice.  One must instead acquire an explicitly
\(\eta\)-dependent writer, pass to a multi-jump/Feynman--Kac path observable, or
use a genuinely nonlinear downstream matching map.

Throughout V21 retain the V20 notation
\begin{equation}
 k_a^\sigma(\eta)
 =\frac{s_a+\sigma h\bigl(d_a+(N\eta)_a\bigr)}{2h^2},
 \qquad
 \lambda:=\sum_{a=1}^6s_a,
 \qquad
 RN=0,
 \qquad N^TN=I_3,
 \label{eq:v21-rates}
\end{equation}
with \(\sigma\in\{+1,-1\}\), and restrict to the open positivity chart on
which every numerator in \eqref{eq:v21-rates} is positive.  The total hazard is
\begin{equation}
 \mathcal K_h=\sum_{a,\sigma}k_a^\sigma=\frac{\lambda}{h^2},
 \label{eq:v21-total-hazard}
\end{equation}
independent of \(\eta\).

\subsection{The residual first-jump law and its exact score cancellation}
\label{sec:v21-first-jump-law}

Integrating the V20 winner--time density over the holding time gives the
winner probabilities
\begin{equation}
 \boxed{
 q_{a\sigma}(\eta)
 =\frac{s_a+\sigma h\bigl(d_a+(N\eta)_a\bigr)}{2\lambda}.}
 \label{eq:v21-winner-prob}
\end{equation}
Because the total hazard \eqref{eq:v21-total-hazard} is fixed, the holding time
is independent of the winner and has the common density
\begin{equation}
 \varpi_h(t)=\frac{\lambda}{h^2}e^{-\lambda t/h^2},
 \qquad t\ge0,
 \label{eq:v21-holding-density}
\end{equation}
for every \(\eta\).

At the reference point \(\eta=0\), write
\begin{equation}
 \mathfrak s_\alpha(a,\sigma)
 =\partial_{\eta^\alpha}\log q_{a\sigma}(\eta)\big|_{\eta=0}
 =\frac{\sigma h n_{a\alpha}}{s_a+\sigma h d_a},
 \label{eq:v21-score}
\end{equation}
which is the V20 residual score.  The next identity is the elementary algebra
which makes the first-jump selector compressible.

\begin{lemma}[Probability--score cancellation]
\label{lem:v21-score-cancellation}
For every \(a,\sigma,\alpha\),
\begin{equation}
 \boxed{
 q_{a\sigma}(0)\,\mathfrak s_\alpha(a,\sigma)
 =\frac{\sigma h}{2\lambda}n_{a\alpha}.}
 \label{eq:v21-score-cancel}
\end{equation}
In particular the right-hand side is independent of the calibrated drift
\(d\), and
\begin{equation}
 \sum_{a,\sigma}q_{a\sigma}(0)\mathfrak s_\alpha(a,\sigma)=0.
 \label{eq:v21-score-centered}
\end{equation}
\end{lemma}

\begin{proof}
Substitute \eqref{eq:v21-winner-prob} and \eqref{eq:v21-score}; the factors
\(s_a+\sigma hd_a\) cancel.  Summing over \(\sigma=\pm1\) then gives zero for
each \(a\), proving centeredness.
\end{proof}

\subsection{Exact compression of an arbitrary same-first-jump writer}
\label{sec:v21-writer-compression}

Let \(\mathcal A_I(a,\sigma,t;\eta)\), \(1\le I\le m\), be any finite bank of
integrable physical writers occurring on the same first-jump cylinder.  The
index \(I\) may label a local action coefficient, stress/contact writer,
common-action score coordinate, or any later physical Wilson readout before the
V18 physical quotient is applied.  Define its conditional holding-time mean
\begin{equation}
 \overline{\mathcal A}_{I,a\sigma}(\eta)
 :=\int_0^\infty
 \mathcal A_I(a,\sigma,t;\eta)\,\varpi_h(t)\,dt.
 \label{eq:v21-conditional-mean}
\end{equation}
The common holding density is independent of \(\eta\), so explicit writer
variation is the only direct derivative inside this conditional mean.

\begin{theorem}[Exact residual derivative of a same-first-jump writer]
\label{thm:v21-exact-selector}
At \(\eta=0\),
\begin{equation}
 \boxed{
 \partial_{\eta^\alpha}\mathbb E_\eta[\mathcal A_I]\big|_0
 =\mathsf D_{I\alpha}
 +\frac{h}{2\lambda}
   \sum_{a=1}^6 n_{a\alpha}
   \Delta\mathcal A_{Ia},}
 \label{eq:v21-exact-selector}
\end{equation}
where
\begin{align}
 \mathsf D_{I\alpha}
 &:={\sum_{a,\sigma}}q_{a\sigma}(0)
   \partial_{\eta^\alpha}
   \overline{\mathcal A}_{I,a\sigma}(\eta)\big|_0,
 \label{eq:v21-direct-array}\\
 \Delta\mathcal A_{Ia}
 &:=\overline{\mathcal A}_{I,a+}(0)
     -\overline{\mathcal A}_{I,a-}(0).
 \label{eq:v21-reversal-contrast}
\end{align}
Equivalently, with
\(\Delta\mathcal A\in\mathbb R^{m\times6}\),
\begin{equation}
 \boxed{
 \mathsf A_{\rm res}
 =\mathsf D+\frac{h}{2\lambda}\Delta\mathcal A\,N.}
 \label{eq:v21-Ares}
\end{equation}
The second term is exactly the V20 score--writer covariance matrix.
\end{theorem}

\begin{proof}
Differentiate the finite expectation
\[
 \mathbb E_\eta[\mathcal A_I]
 =\sum_{a,\sigma}q_{a\sigma}(\eta)
   \overline{\mathcal A}_{I,a\sigma}(\eta).
\]
The derivative of the second factor gives \(\mathsf D\).  For the probability
factor, the score identity gives
\[
 \partial_{\eta^\alpha}q_{a\sigma}(0)
 =q_{a\sigma}(0)\mathfrak s_\alpha(a,\sigma)
 =\frac{\sigma h}{2\lambda}n_{a\alpha}
\]
by Lemma~\ref{lem:v21-score-cancellation}.  Summing the \(\sigma=\pm1\)
contributions gives the reversal contrast
\eqref{eq:v21-reversal-contrast}.  The covariance interpretation follows from
\(\mathbb E\mathfrak s_\alpha=0\):
\[
 \operatorname{Cov}(\mathcal A_I,\mathfrak s_\alpha)
 =\mathbb E[\mathcal A_I\mathfrak s_\alpha]
 =\frac{h}{2\lambda}\sum_a n_{a\alpha}\Delta\mathcal A_{Ia}.
\]
This is \eqref{eq:v21-Ares}.
\end{proof}

\begin{corollary}[Only the reversal-odd writer sector is visible at first order]
\label{cor:v21-odd-only}
If
\begin{equation}
 \overline{\mathcal A}_{I,a+}(0)
 =\overline{\mathcal A}_{I,a-}(0)
 \qquad\text{for every }a,
 \label{eq:v21-reversal-even}
\end{equation}
and the writer has no direct \(\eta\)-dependence, then
\begin{equation}
 \partial_\eta\mathbb E_\eta[\mathcal A_I]\big|_0=0.
 \label{eq:v21-even-zero}
\end{equation}
Thus a purely reversal-even local writer cannot linearly detect the residual
circulation packet on one first-jump cylinder.
\end{corollary}

\begin{proof}
Under the hypotheses both terms on the right of
\eqref{eq:v21-exact-selector} vanish.
\end{proof}

\subsection{Cycle projection is the complete first-order stochastic selector}
\label{sec:v21-cycle-projection}

The six reversal contrasts contain three drift-visible and three
circulation-visible components.  The residual score sees only the latter.
Define
\begin{equation}
 P_{\rm cyc}:=NN^T,
 \qquad
 P_{\rm grad}:=I_6-P_{\rm cyc}.
 \label{eq:v21-cycle-projector}
\end{equation}
Because \(N\) is an orthonormal basis of \(\ker R\),
\begin{equation}
 \Ran P_{\rm cyc}=\ker R,
 \qquad
 \Ran P_{\rm grad}=\Ran R^T.
 \label{eq:v21-cycle-grad}
\end{equation}

\begin{theorem}[Cycle-incidence criterion]
\label{thm:v21-cycle-incidence}
For a writer bank with no direct \(\eta\)-dependence,
\begin{equation}
 \mathsf A_{\rm res}
 =\frac{h}{2\lambda}\Delta\mathcal A N.
 \label{eq:v21-cycle-Ares}
\end{equation}
For each writer row \(I\), the following are equivalent:
\begin{enumerate}[label=\textup{(C\arabic*)},leftmargin=3em]
\item its residual first-order selector vanishes;
\item \(\Delta\mathcal A_I N=0\);
\item \(\Delta\mathcal A_I P_{\rm cyc}=0\);
\item \(\Delta\mathcal A_I\in\Ran R^T\).
\end{enumerate}
Hence the complete stochastic first-order P4d content of the six directional
writer contrasts is their projection onto the three-dimensional cycle space
\(\ker R\).
\end{theorem}

\begin{proof}
The first two statements are identical by
\eqref{eq:v21-cycle-Ares}.  Since \(NN^T\) is the orthogonal projector onto
\(\Ran N\), \(\Delta\mathcal A_I N=0\) if and only if
\(\Delta\mathcal A_I NN^T=0\).  Finally
\[
 (\ker R)^\perp=\Ran R^T
\]
by finite-dimensional linear algebra, proving the last equivalence.
\end{proof}

\begin{corollary}[Target-minimal first-order acquisition count]
\label{cor:v21-minimal-acquisition}
For \(m\) writer outputs with no explicit \(\eta\)-dependence, the first-order
residual selector requires only the \(3m\) cycle-projected numbers
\begin{equation}
 \boxed{
 \mathcal C_{I\alpha}
 :=\sum_{a=1}^6\Delta\mathcal A_{Ia}n_{a\alpha},
 \qquad 1\le I\le m,\quad1\le\alpha\le3.}
 \label{eq:v21-cycle-table}
\end{equation}
The full twelve-event mean table is sufficient but not target minimal.
However, the \(3m\) numbers must be produced from same-record labelled writer
data; they are not determined by unlabelled writer marginals.
\end{corollary}

\begin{proof}
Equation \eqref{eq:v21-cycle-Ares} depends on the writer table only through
\(\Delta\mathcal A N\), which has \(3m\) entries.  The necessity of same-record
labelling is proved in Countertheorem~\ref{cth:v21-marginals-no-cross} below.
\end{proof}

\subsection{Separate score and writer marginals do not determine the cross Gram}
\label{sec:v21-identifiability}

The preceding compression does not remove the occurrence requirement.  A
cross Gram is a property of a joint same-record law, not of two separate
marginal distributions.

\begin{proposition}[Counterexample: fixed marginals, arbitrary cross incidence]
\label{cth:v21-marginals-no-cross}
Let \(X,Y\in\{-1,+1\}\) have joint law
\begin{equation}
 \mathbb P_\rho(X=x,Y=y)
 =\frac{1+\rho xy}{4},
 \qquad -1\le\rho\le1.
 \label{eq:v21-rho-coupling}
\end{equation}
Then the marginals of \(X\) and \(Y\) are uniform for every \(\rho\), while
\begin{equation}
 \boxed{\operatorname{Cov}_\rho(X,Y)=\rho.}
 \label{eq:v21-rho-cov}
\end{equation}
Therefore separate knowledge of a residual score distribution and a writer
distribution does not determine their same-record covariance or the matrix
\(\mathsf A_{\rm res}\).
\end{proposition}

\begin{proof}
Summing \eqref{eq:v21-rho-coupling} over either variable gives probability
\(1/2\) for each sign, independently of \(\rho\).  Both means are zero and
\[
 \mathbb E_\rho[XY]
 =\sum_{x,y}xy\frac{1+\rho xy}{4}=\rho.
\]
\end{proof}

\begin{corollary}[The same-record common-action row is irreducible]
\label{cor:v21-same-record-irreducible}
A separately reconstructed race score packet and a separately reconstructed
full-field common-action writer bank cannot be substituted for the joint
first-jump action-score incidence demanded by P4d.  The event labels or an
exact equivalent joint coupling are physical input to the cross Gram.
\end{corollary}

\begin{proof}
Any such substitution would determine the covariance from its two marginals,
contradicting Countertheorem~\ref{cth:v21-marginals-no-cross}.
\end{proof}

\subsection{Upstream correction: the one-jump Hessian cannot rescue a silent fixed writer}
\label{sec:v21-hessian-correction}

The winner probabilities \eqref{eq:v21-winner-prob} are exactly affine in
\(\eta\).  This removes a tempting but incorrect fallback strategy.

\begin{theorem}[Affine one-jump no-Hessian theorem]
\label{thm:v21-no-hessian}
Suppose the conditional writer means are independent of \(\eta\):
\begin{equation}
 \overline{\mathcal A}_{I,a\sigma}(\eta)
 \equiv\overline{\mathcal A}_{I,a\sigma}.
 \label{eq:v21-fixed-writer}
\end{equation}
Then throughout the positivity chart
\begin{equation}
 \boxed{
 \mathbb E_\eta[\mathcal A_I]
 =\mathbb E_0[\mathcal A_I]
 +\frac{h}{2\lambda}
   \sum_{a=1}^6(N\eta)_a\Delta\mathcal A_{Ia},}
 \label{eq:v21-affine-expectation}
\end{equation}
and consequently
\begin{equation}
 \boxed{
 \partial_{\eta^\alpha}\partial_{\eta^\beta}
 \mathbb E_\eta[\mathcal A_I]=0}
 \label{eq:v21-hessian-zero}
\end{equation}
identically.

Therefore, if the first-order selector vanishes for such a one-jump writer, a
second derivative of the same statistic contains no hidden P4 information.
Second-order work becomes meaningful only after at least one of the following
changes occurs: the writer itself depends on \(\eta\), a multi-jump/path
observable is used, or the downstream physical matching map is nonlinear in
the selected expectation packet.
\end{theorem}

\begin{proof}
Substitute \eqref{eq:v21-winner-prob} and the fixed means into
\(\sum_{a,\sigma}q_{a\sigma}(\eta)\overline{\mathcal A}_{I,a\sigma}\).
The constant terms give the reference expectation; the coefficient of
\((N\eta)_a\) is
\(h\Delta\mathcal A_{Ia}/(2\lambda)\).  There are no quadratic or higher
terms because \(q_{a\sigma}\) is affine.  Differentiating twice proves
\eqref{eq:v21-hessian-zero}.
\end{proof}

\begin{firewall}[Do not confuse a zero one-jump Hessian with microscopic silence]
Theorem~\ref{thm:v21-no-hessian} concerns a fixed conditional writer evaluated
on one jump.  A path action contains products, occupations, repeated jumps,
Perron/Doob normalization and possibly explicit deformation dependence.  Those
objects can be nonlinear in \(\eta\) even when each elementary winner
probability is affine.  The theorem says only that one must move to that richer
object before charging a second-order acquisition.
\end{firewall}

\subsection{A populated control: entropy-production incidence}
\label{sec:v21-entropy-control}

Although the gravity--matter common-action cross Gram is not yet populated on
the selected residual race cylinder, the race law itself supplies a useful
same-record action control: the stochastic reversal/entropy increment.  Define
\begin{equation}
 \ell_a(d)
 :=\log\frac{s_a+hd_a}{s_a-hd_a},
 \qquad
 \mathcal E(a,\sigma;d):=\sigma\ell_a(d).
 \label{eq:v21-entropy-increment}
\end{equation}
Its mean per jump is
\begin{equation}
 \boxed{
 \Sigma_{\rm jump}(d)
 =\frac{h}{\lambda}
  \sum_{a=1}^6d_a\ell_a(d).}
 \label{eq:v21-entropy-mean}
\end{equation}
Every summand is nonnegative because \(x\log((1+x)/(1-x))\ge0\) for
\(|x|<1\).

\begin{proposition}[Exact residual gradient of the entropy control]
\label{prop:v21-entropy-gradient}
Along \(d(\eta)=d+N\eta\),
\begin{equation}
 \boxed{
 \partial_{\eta^\alpha}\Sigma_{\rm jump}
 =\frac{h}{\lambda}\sum_{a=1}^6n_{a\alpha}
 \left[
  \ell_a(d)
  +\frac{2hs_ad_a}{s_a^2-h^2d_a^2}
 \right].}
 \label{eq:v21-entropy-gradient}
\end{equation}
At the detailed-balance reference \(d=0\), the gradient vanishes.
\end{proposition}

\begin{proof}
Differentiate \(d_a\ell_a(d_a)\).  One has
\[
 \ell_a'(d_a)
 =\frac{h}{s_a+hd_a}+\frac{h}{s_a-hd_a}
 =\frac{2hs_a}{s_a^2-h^2d_a^2}.
\]
The chain rule with \(\partial_{\eta^\alpha}d_a=n_{a\alpha}\) gives the
formula.  At \(d=0\), both \(\ell_a\) and the product \(d_a\ell'_a\) vanish.
\end{proof}

\begin{theorem}[Entropy Hessian equals four times the residual Fisher metric at equilibrium]
\label{thm:v21-entropy-fisher}
At \(d=0\),
\begin{equation}
 \boxed{
 \nabla_\eta^2\Sigma_{\rm jump}\big|_{d=0}
 =\frac{4h^2}{\lambda}N^TS^{-1}N
 =4G_\eta\big|_{d=0}.}
 \label{eq:v21-entropy-fisher}
\end{equation}
In particular the entropy-production Hessian is strictly positive on every
nonzero residual race direction.
\end{theorem}

\begin{proof}
Near zero,
\[
 \ell_a(d_a)
 =\frac{2hd_a}{s_a}+O(d_a^3),
\]
so
\[
 \Sigma_{\rm jump}(d)
 =\frac{2h^2}{\lambda}
   d^TS^{-1}d+O(\|d\|^4).
\]
Restricting to \(d=N\eta\) gives the first Hessian in
\eqref{eq:v21-entropy-fisher}.  V20 gives
\(G_\eta|_{d=0}=h^2\lambda^{-1}N^TS^{-1}N\), proving the second equality.
Positivity follows from \(S\succ0\) and \(N^TN=I\).
\end{proof}

\begin{remark}[What the entropy control establishes]
\label{assess:v21-entropy-control}
The identity \eqref{eq:v21-entropy-fisher} is a populated same-record
score--action consistency test.  It verifies that the residual race sector
carries a real nonlinear stochastic action despite the vanishing first-order
response at equilibrium.  It does \emph{not} identify that stochastic entropy
writer with the gravity--matter common action or with a Wilson matching
coefficient.  Substituting it for the missing full-field action row would
violate the P4 provenance requirement.
\end{remark}

\subsection{The inherited common-action compiler and the current physical stop}
\label{sec:v21-common-action-stop}

The physical population programme already provides the correct abstract
compiler for the row P4d now needs.  In its notation, the path-likelihood score
synthesis \(S_{\rm jump}\) and the independently declared full-field common
action score synthesis \(S_{\rm com}\) are compared only after shorting the
certified BRST-exact, boundary, endpoint-coboundary and coordinate-gauge score
bank \(N_{\rm nui}\).  With
\begin{equation}
 P_{\rm nui}
 :=N_{\rm nui}(N_{\rm nui}^*N_{\rm nui})^\dagger N_{\rm nui}^*,
 \label{eq:v21-Pnui}
\end{equation}
the inherited physical incidence Gram is
\begin{equation}
 \boxed{
 \mathbb C_{\rm jump\leftrightarrow com}^{\rm act}
 =(S_{\rm com}-S_{\rm jump})^*
  (I-P_{\rm nui})
  (S_{\rm com}-S_{\rm jump})\succeq0.}
 \label{eq:v21-inherited-action-Gram}
\end{equation}
Zero of this Gram is the exact criterion for one certified common action
representative on the same path; its polar range is otherwise the unique
source-minimal missing action writer.

The load-bearing fact for the present phenomenology line is not a new theorem
about \eqref{eq:v21-inherited-action-Gram}.  It is the population status of its
inputs.  The inherited foundational and physical-population programmes state
that the compilers exist, but the complete packet
\begin{equation}
 \boxed{
 \text{resolved accepted path law}
 +\text{ executable full-action writer}
 +\text{ Perron/control rows}
 +\text{ complex writer}
 +\text{ adjacent cross Gram}}
 \label{eq:v21-required-packet}
\end{equation}
is not displayed on one selected nonlinear Store--SM--gravity bridge.  In
particular, the full-field common-action score on the residual V20 race cylinder
is not currently a populated table from which the reversal contrasts
\(\Delta\mathcal A\) could be evaluated.

\begin{theorem}[P4d identifiability alternative at the current corpus boundary]
\label{thm:v21-P4d-alternative}
Let \(\mathcal W_{\rm pilot}\) be the V18 physical Wilson quotient and let
\(\mathsf K_{\rm phys}\) be the populated downstream calibration/renormalization
map of V20.  Then exactly one of the following applies to the current residual
race packet:
\begin{enumerate}[label=\textup{(P\arabic*)},leftmargin=3em]
\item \textbf{Same-record action row populated.}  The conditional means and
      direct derivatives of the declared common-action writer determine
      \(\mathsf A_{\rm res}\) by \eqref{eq:v21-Ares}, and the physical
      predictivity matrix is
      \begin{equation}
       \boxed{
       \mathsf S
       =\mathsf K_{\rm phys}
        \left(
         \mathsf D+\frac{h}{2\lambda}\Delta\mathcal A N
        \right).}
       \label{eq:v21-physical-S}
      \end{equation}
      Its left nullspace is then a genuine candidate predictive conormal.
\item \textbf{Same-record action row unpopulated.}  The matrix
      \(\mathsf A_{\rm res}\), and hence \(\mathsf S\), is
      \emph{unidentified}.  Neither the zero matrix nor a fitted Wilson boundary
      condition is licensed.  The minimum next acquisition is the labelled
      reversal/cycle action table of Corollary~\ref{cor:v21-minimal-acquisition},
      together with the direct derivative array \(\mathsf D\) whenever the
      writer explicitly depends on \(\eta\).
\end{enumerate}
For the currently supplied corpus the second branch is the justified status of
the selected full-field gravity--matter action bridge.
\end{theorem}

\begin{proof}
If the joint writer exists, Theorem~\ref{thm:v21-exact-selector} computes its
microscopic derivative and V20 gives the composition with
\(\mathsf K_{\rm phys}\).  The V16 predictive-conormal theorem then identifies
the left nullspace as the first-order relation space.  If the joint writer does
not exist, Countertheorem~\ref{cth:v21-marginals-no-cross} forbids recovering
the covariance from separate marginals, while V19 forbids treating a missing
map as zero and V18 forbids reinstating the Wilson outputs as free microscopic
inputs.  The inherited population programme explicitly classifies this packet
as an upstream acquisition row.  Therefore the second branch is an
identifiability statement, not a numerical zero.
\end{proof}

\begin{center}
\fbox{\parbox{0.92\textwidth}{
\textbf{P4d has been compressed to a three-column physical action-incidence
row, but that row is not yet populated for the selected nonlinear
Store--SM--gravity bridge.}
At first-jump level the residual race score sees only the cycle projection of
the reversal-odd common-action contrast.  A fixed one-jump writer is affine in
the residual race coordinates, so a vanishing linear selector cannot be
rescued by its second derivative.  The next valid move is same-record
common-action acquisition, or, if that first-jump cycle contrast is proved
zero, promotion to a multi-jump/Feynman--Kac common-action observable.}}
\end{center}

\subsection{Rank consequences once the action row lands}
\label{sec:v21-rank-consequences}

Although the physical writer is not yet loaded, it is useful to record the
exact rank test that will terminate immediately when it is.  Let the selected
pilot quotient have dimension \(d_{\rm pilot}\), and let
\begin{equation}
 \mathsf S\in\mathbb R^{d_{\rm pilot}\times3}
 \label{eq:v21-S-shape}
\end{equation}
be \eqref{eq:v21-physical-S}.  Then
\begin{equation}
 \boxed{
 \dim\ker\mathsf S^T
 =d_{\rm pilot}-\operatorname{rank}\mathsf S
 \ge d_{\rm pilot}-3.}
 \label{eq:v21-nullity-bound}
\end{equation}
This becomes a phenomenological prediction count only under the V19
completeness hypothesis that the three residual race directions constitute the
entire unfitted microscopic selector packet for this pilot quotient.

\begin{corollary}[Immediate P4 decision after acquisition]
\label{cor:v21-P4-decision}
Once the physical three-column matrix \(\mathsf S\) is populated:
\begin{enumerate}[label=\textup{(D\arabic*)},leftmargin=3em]
\item if \(\operatorname{rank}\mathsf S=3\), then every pilot output direction transverse to
      its three-dimensional image gives a candidate first-order Renewal
      relation;
\item if \(\operatorname{rank}\mathsf S<3\), the relation space is larger and the microscopic
      residual packet itself contains inactive combinations;
\item if additional independently occurring microscopic selector coordinates
      are subsequently found, their columns must be appended before any left
      null vector is called a prediction.
\end{enumerate}
\end{corollary}

\begin{proof}
The nullity identity is elementary rank--nullity.  The final qualification is
V16--V19's no-free-prediction rule: the predictive conormal is the annihilator
of the complete microscopic matching tangent, not of a prematurely truncated
column set.
\end{proof}

\subsection{V21 anti-circularity audit}
\label{sec:v21-audit}

\begin{enumerate}[leftmargin=2.4em]
\item \textbf{A missing common-action table is not represented by zero.}
The physical selector remains unidentified until the writer is attached to the
same labelled race records.
\item \textbf{Separate writer and score marginals are not substituted for a
cross Gram.}
Countertheorem~\ref{cth:v21-marginals-no-cross} gives an exact coupling
counterexample.
\item \textbf{The first-jump acquisition is compressed only after provenance is
fixed.}
The \(3m\) cycle table is target sufficient, but it must be computed from
same-record reversal-labelled writer values.
\item \textbf{A zero first-order one-jump selector is not followed blindly by a
Hessian.}
For an \(\eta\)-independent one-jump writer the Hessian is identically zero by
Theorem~\ref{thm:v21-no-hessian}; the correct escalation is to a richer path
observable or an explicitly deforming writer.
\item \textbf{The stochastic entropy row is a control, not the ESM action.}
Its Hessian/Fisher identity validates the incidence machinery but supplies no
gravity--matter Wilson matching condition.
\item \textbf{The three-column rank bound is conditional on selector
completeness.}
Any newly identified microscopic action/source coordinate must be appended to
the incidence matrix before prediction counting.
\end{enumerate}

\section{V21 status ledger}
\label{sec:v21-status}

\subsection*{Proved in V21}
\begin{enumerate}[leftmargin=2.2em]
\item The residual first-jump winner probabilities are affine in \(\eta\), the
holding-time law is \(\eta\)-independent, and the exact probability--score
product is \(\sigma h n_{a\alpha}/(2\lambda)\).
\item For every same-first-jump writer bank, the exact residual derivative is
\(\mathsf A_{\rm res}=\mathsf D+(h/2\lambda)\Delta\mathcal A N\).  The second
term is exactly the score--writer covariance.
\item With no explicit writer deformation, only the cycle projection of the
reversal-odd six-direction contrast contributes.  The selector vanishes
exactly when the contrast lies in \(\Ran R^T\).
\item The target-minimal stochastic acquisition for \(m\) first-jump writers is
the \(m\times3\) cycle contrast table, after same-record event labels are
available.
\item Separate score and writer marginals do not determine their covariance;
a one-parameter binary coupling gives fixed marginals and arbitrary cross
incidence.
\item The expectation of an \(\eta\)-independent one-jump writer is exactly
affine in \(\eta\), so its complete residual Hessian vanishes.  A zero linear
selector must therefore be escalated to a richer physical object rather than
the same writer's second derivative.
\item The race entropy-production control is fully populated.  At detailed
balance its residual gradient vanishes and its Hessian equals four times the
V20 residual Fisher metric:
\(\nabla_\eta^2\Sigma_{\rm jump}=4G_\eta\succ0\).
\item Once a physical common-action writer is loaded, the complete P4d matrix is
\(\mathsf S=\mathsf K_{\rm phys}[\mathsf D+(h/2\lambda)\Delta\mathcal A N]\),
so its predictive conormal is immediately computable.
\end{enumerate}

\subsection*{Inherited}
\begin{enumerate}[leftmargin=2.2em]
\item V16--V20: physical calibration, native RG covariance, the no-free-Wilson
prediction theorem, the physical Wilson quotient, the populated race
calibration matrix, the three-dimensional residual circulation packet, and its
positive score/Fisher geometry.
\item The physical quantum-packet population programme: the path-score/common-
action incidence Gram and its nuisance-short Pythagoras, with the polar residual
as the unique source-minimal missing action writer.
\item The foundational closure programme: an executable full-action writer and
its same-history Perron/control rows are upstream physical inputs, and the
complete selected nonlinear Store--SM--gravity packet is not supplied merely by
having the abstract compiler.
\end{enumerate}

\subsection*{Conditional}
\begin{enumerate}[leftmargin=2.2em]
\item The matrix \(\mathsf S\) becomes a numerical/analytic P4 prediction
matrix only after the selected full-field common-action writer is populated on
the same residual race cylinder.
\item The lower bound \(d_{\rm pilot}-3\) on candidate relations assumes the
three residual race coordinates exhaust the unfitted microscopic selector
packet for the pilot quotient.  V21 does not prove that completeness.
\item The entropy control is an exact stochastic path-action diagnostic, not a
substitute for the Einstein--Standard-Model action writer.
\item All phenomenological statements retain the V16--V20 fixed finite
loop/EFT/source panel and positive reference-scale quantifiers.
\end{enumerate}

\subsection*{Open beyond V21}
\begin{enumerate}[leftmargin=2.2em]
\item Populate the full-field common-action score on the same twelve directed
residual race events (or an exactly equivalent joint cylinder), and compute its
six reversal contrasts, three cycle projections, and any direct
\(\eta\)-derivative rows.
\item Apply the V21 matrix immediately to the smallest physical gravity--matter
Wilson quotient and compute its rank and left nullspace.
\item If the proved first-jump common-action cycle contrast vanishes and the
writer has no explicit \(\eta\)-dependence, move to the minimum multi-jump or
Feynman--Kac common-action observable; do not compute the identically zero
one-jump Hessian.
\item Audit whether additional independently occurring microscopic
full-action/source coordinates supplement \(\eta\); append their columns before
promoting any null relation to a prediction.
\item The parallel P2/P3 programme remains: populate higher-loop native beta
rows, threshold matching, standard-scheme conversion, and eventually numerical
RG evolution after a predictive boundary relation has actually been found.
\end{enumerate}

\subsection*{Exact next load-bearing theorem}
\begin{center}
\fbox{\parbox{0.92\textwidth}{
\textbf{P4e: same-history full-field action-score landing on the residual race
cylinder.}
Attach one executable Einstein--Standard-Model common-action writer (or its
source-minimal physical score bank) to the same twelve directed first-jump
records whose residual scores are already populated.  After the independently
certified nuisance short, evaluate the six reversal contrasts and their three
cycle projections, add any direct residual derivative, and form
\(\mathsf S=\mathsf K_{\rm phys}\mathsf A_{\rm res}\).  Compute its rank and
left nullspace before fitting any pilot Wilson coefficient.  If the first-jump
cycle row is proved zero and has no direct deformation, promote the acquisition
to the minimum multi-jump/Feynman--Kac action writer; the fixed one-jump Hessian
is already proved identically zero.}}
\end{center}

\section*{Append-only continuation marker: V21}
\addcontentsline{toc}{section}{Append-only continuation marker: V21}
\textbf{End of V21.}  The complete V1--V20 body above is retained verbatim.
V21 compresses the residual score--action selector to its exact first-jump
content.  The total race hazard is independent of the three residual
circulation coordinates, so the score covariance with a same-event writer is
exactly \((h/2\lambda)\Delta\mathcal A N\): only the reversal-odd cycle
projection matters.  Separate marginals do not determine this cross Gram, and
a fixed one-jump writer is affine in the residual coordinates, so a vanishing
linear selector cannot be rescued by its Hessian.  The populated stochastic
entropy control satisfies the exact equilibrium identity
\(\nabla_\eta^2\Sigma_{\rm jump}=4G_\eta\), validating the incidence compiler
without substituting stochastic entropy for the ESM common action.  The
physical P4d stop is therefore precise: the selected full-field common-action
score has not yet been populated on the same residual race cylinder.  Once it
lands, only a three-column cycle-incidence matrix remains before the first
predictivity rank test.



\clearpage
\section{V22 continuation: common-action coboundary no-go, reversal parity, and the first nontrivial path selector}
\label{sec:v22-continuation}

V21 reduced the residual race/action problem to a three-column first-jump
incidence matrix and stated the correct escalation rule: if the physical
first-jump common-action row is proved zero, do not compute the same writer's
identically vanishing Hessian; promote the acquisition to the minimum richer
path observable.  The present continuation carries out that upstream test.

Two distinct facts must be separated.  First, the inherited protected
Palatini--Cartan--Holst and first-jet Einstein--Standard-Model actions are
\emph{scalar functions of the determining field tuple}.  The fast residual
race winner and waiting time are separate scheduler marks.  On the scheduler-
only first-jump cylinder the determining field is held fixed, so that decoded
scalar action writer has no reversal contrast.  Second, even when a scalar
common-action increment is evaluated between two physical endpoint fields, it
is an exact coboundary.  Summing it along a path telescopes to an endpoint
term and therefore has zero extensive Feynman--Kac pressure.  Merely replacing
one jump by many jumps does not create phenomenology.

The first extensive selector must consequently use a \emph{non-coboundary}
physical action observable: for example an independently acquired accepted
Gibbs/occurrence cost or the reflected action-response/Dirichlet writer of the
Grand Tensor transfer packet.  On a reflected base such action-response
writers are reversal even, whereas the V20 residual circulation score is
reversal odd.  Their linear cross-incidence therefore vanishes by symmetry;
the first possible reflected signal is a second-order circulation
susceptibility of a path statistic.  V22 derives an explicit finite Green-
operator compiler for that susceptibility and routes it through the inherited
three-cylinder action-response closure.

Throughout, the physical calibration and finite coefficient packet
\(\mathfrak q\) are those of V16--V21.  No Wilson coefficient is fitted in this
section.

\subsection{The decoded common action really gives a zero first-jump row}
\label{sec:v22-first-jump-zero}

Let \(\Theta\) denote the complete protected determining-field tuple on which
the inherited scalar common-action writer
\begin{equation}
 \mathcal S_{\rm com}(\Theta)
 \label{eq:v22-common-action}
\end{equation}
is defined.  It may be taken to be the protected
Palatini--Cartan--Holst/common first-jet Einstein--SM action on the declared
branch.  Let
\begin{equation}
 \Gamma=(B,T),
 \qquad B=(a,\sigma),\quad a=1,\ldots,6,\quad \sigma=\pm1,
 \label{eq:v22-race-mark}
\end{equation}
be the fast root winner--waiting mark of the V20 residual race cylinder.  By
construction this mark is not itself the accepted determining-field update;
the initial field, terminal field and update/reverse label belong to the
separate full-field mark of the same-history packet.

\begin{theorem}[Scheduler-only scalar-action silence]
\label{thm:v22-scheduler-silence}
On the scheduler-only first-jump cylinder, hold \(\Theta\) fixed while varying
\(\eta\in\ker R\).  Then every scalar writer which factors through \(\Theta\)
has zero V21 residual incidence.  In particular,
\begin{equation}
 \boxed{
 \partial_{\eta^\alpha}\mathcal S_{\rm com}(\Theta)=0,
 \qquad
 \Delta\mathcal A_{a}=0,
 \qquad
 \mathsf A_{\rm res}=0}
 \label{eq:v22-scheduler-zero}
\end{equation}
for the decoded scalar common action on this cylinder, through every retained
source derivative.
\end{theorem}

\begin{proof}
The scheduler mark changes only the winner/waiting variables of the fast root
race.  The determining-field tuple is held fixed, so any writer of the form
\(F(\Theta)\) is independent of \(B,T\) and of the residual coordinate
\(\eta\).  Its forward and reverse conditional means are equal, hence the V21
reversal contrast \(\Delta\mathcal A_a\) is zero.  There is also no direct
\(\eta\)-derivative.  Substitute in
\(\mathsf A_{\rm res}=\mathsf D+(h/2\lambda)\Delta\mathcal A N\).
The same argument commutes with any fixed external/source differentiation
because those sources are independent coordinates of the declared action
packet.
\end{proof}

\begin{firewall}[This zero is not the accepted common-action incidence]
Theorem~\ref{thm:v22-scheduler-silence} closes only the temptation to identify
the fast scheduler mark itself with a physical field/action update.  The full
same-history event additionally carries initial and terminal determining
fields, an update/reverse label and an independently evaluated action
increment.  Its score--action cross incidence is a different joint record and
is not set to zero by \eqref{eq:v22-scheduler-zero}.
\end{firewall}

\subsection{Exact common-action increments are coboundaries, not extensive rewards}
\label{sec:v22-coboundary}

The next possible shortcut is to keep the same scalar action but sum its signed
increments over many jumps.  This also fails for a structural reason.

Let \(L\) be a finite continuous-time Markov generator, acting on column
functions, with off-diagonal entries \(L(x,y)\ge0\) and \(L\mathbf 1=0\).
For a scalar state function \(S\), define the directed increment
\begin{equation}
 g_S(x,y):=S(y)-S(x).
 \label{eq:v22-coboundary-reward}
\end{equation}
The Feynman--Kac jump tilt is
\begin{equation}
 [B_k](x,y)=L(x,y)e^{k g_S(x,y)}\quad(x\ne y),
 \qquad
 [B_k](x,x)=L(x,x).
 \label{eq:v22-coboundary-tilt}
\end{equation}

\begin{theorem}[Coboundary Feynman--Kac no-go]
\label{thm:v22-coboundary-no-go}
Put \(D_k=\operatorname{diag}(e^{kS})\).  Then
\begin{equation}
 \boxed{B_k=D_k^{-1}LD_k.}
 \label{eq:v22-coboundary-similarity}
\end{equation}
Hence \(B_k\) is similar to \(L\), its Perron/spectral bound is identically
zero,
\begin{equation}
 \boxed{\psi_S(k)=s(B_k)=0,}
 \label{eq:v22-coboundary-scgf-zero}
\end{equation}
and along every path
\begin{equation}
 \boxed{
 \sum_{0<t\le T}g_S(X_{t-},X_t)
 =S(X_T)-S(X_0).}
 \label{eq:v22-coboundary-telescope}
\end{equation}
If \(L=L_\eta\) varies in any fixed-support regulator family while the same
scalar \(S\) is used, then
\begin{equation}
 \boxed{
 \partial_\eta^I\partial_k^m\psi_S(\eta,k)=0
 \quad\text{for all }m\ge1}
 \label{eq:v22-coboundary-all-zero}
\end{equation}
at every point where the family is defined.
\end{theorem}

\begin{proof}
For \(x\ne y\),
\[
 [D_k^{-1}LD_k](x,y)
 =e^{-kS(x)}L(x,y)e^{kS(y)}
 =L(x,y)e^{k(S(y)-S(x))}.
\]
The diagonal is unchanged, proving
\eqref{eq:v22-coboundary-similarity}.  Similar matrices have the same spectrum,
while a Markov generator has spectral bound zero.  The path identity is the
ordinary telescoping sum.  The parameter-dependent statement follows because
the same similarity identity holds for every \(\eta\); therefore the entire
function \(\psi_S(\eta,k)\) is zero.
\end{proof}

\begin{corollary}[Do not promote the same signed action increment to more jumps]
\label{cor:v22-no-multijump-coboundary}
If the full-field common-action increment is literally the difference of one
scalar common action at its two endpoints, then using two, three or arbitrarily
many such signed increments in a Feynman--Kac pressure cannot create an
extensive P4 selector.  A nonzero extensive selector must use a physical
quantity which is not an exact state-function coboundary, such as a local
stationarity/occurrence cost, a positive action-response writer, or an
explicitly source-dependent action density.
\end{corollary}

\subsection{Reversal parity routes the residual circulation}
\label{sec:v22-parity}

The V20 residual deformation is a circulation: reversing the directed root
label sends \(\eta\) to \(-\eta\).  We record the resulting parity rule in a
form independent of the detailed action model.

\begin{definition}[Reflected residual family]
\label{def:v22-reflected-family}
Let \(\vartheta\) be an involution on finite path records.  A one-parameter
slice \(\mathbb P_t\) of the residual family is \emph{reflected} if
\begin{equation}
 \mathbb P_t(\vartheta\omega)=\mathbb P_{-t}(\omega).
 \label{eq:v22-reflected-law}
\end{equation}
A path writer \(A\) is even or odd if
\(A\circ\vartheta=+A\) or \(-A\), respectively.
\end{definition}

\begin{theorem}[Parity router for residual score incidence]
\label{thm:v22-parity-router}
At \(t=0\), the residual score
\begin{equation}
 \mathfrak s(\omega)=\left.\partial_t\log\mathbb P_t(\omega)\right|_{t=0}
 \label{eq:v22-score-odd}
\end{equation}
is reversal odd.  Consequently every integrable reversal-even writer satisfies
\begin{equation}
 \boxed{
 \mathbb E_0[A\,\mathfrak s]=0.}
 \label{eq:v22-even-linear-zero}
\end{equation}
More generally, if a tilted path exponent \(\psi(t,k)\) uses a reversal-even
reward, then
\begin{equation}
 \boxed{\psi(t,k)=\psi(-t,k)}
 \label{eq:v22-psi-even}
\end{equation}
and every odd residual derivative of \(\psi\) vanishes at \(t=0\).
\end{theorem}

\begin{proof}
Differentiate \eqref{eq:v22-reflected-law} at zero and divide by the positive
base law to obtain
\(\mathfrak s(\vartheta\omega)=-\mathfrak s(\omega)\).  The base law is
\(\vartheta\)-invariant, so the integral of an odd function is zero, proving
\eqref{eq:v22-even-linear-zero}.  For an even reward the tilted path weight is
also invariant under simultaneous \(t\mapsto-t\) and
\(\omega\mapsto\vartheta\omega\), yielding
\eqref{eq:v22-psi-even}.
\end{proof}

The signed endpoint increment of a scalar action is reversal odd, but
Theorem~\ref{thm:v22-coboundary-no-go} removes it from extensive pressure
statistics.  By contrast, the inherited reflected action-response writer
\begin{equation}
 \mathcal A_F(x,y)
 =\frac12\bigl(F(y)-F(x)\bigr)^*\bigl(F(y)-F(x)\bigr)
 \label{eq:v22-pair-action}
\end{equation}
is reversal even and can be extensive, but
Theorem~\ref{thm:v22-parity-router} kills its \emph{linear} circulation
selector on a reflected base.  Thus the first possible reflected coupling is
quadratic in the residual circulation.

\subsection{The first nontrivial reflected path compiler}
\label{sec:v22-path-susceptibility}

We now derive the finite matrix required for that quadratic response.  The
result is useful independently of whether the currently loaded full-field
packet already supplies all its entries.

Let \(L\) be an irreducible finite generator with stationary row vector
\(\pi\), \(\pi L=0\), \(L\mathbf1=0\).  Let
\begin{equation}
 Q:=I-\mathbf1\pi,
 \qquad
 G:=Q(-L)^{-1}Q
 \label{eq:v22-Green}
\end{equation}
be the centered Green operator.  Consider an affine circulation deformation
\begin{equation}
 L_\eta=L+\eta^\alpha K_\alpha,
 \qquad
 K_\alpha\mathbf1=0,
 \qquad
 \pi K_\alpha=0,
 \label{eq:v22-Leta}
\end{equation}
so \(\pi\) remains stationary.  Let \(g(x,y)\) be an
\(\eta\)-independent directed jump reward and define the first two tilt
matrices at \(k=0\),
\begin{align}
 J_\eta(x,y)&=L_\eta(x,y)g(x,y),\qquad x\ne y,
 \label{eq:v22-Jeta}\\
 H_\eta(x,y)&=L_\eta(x,y)g(x,y)^2,
 \qquad x\ne y,
 \label{eq:v22-Heta}
\end{align}
with zero diagonal.  Write
\(J=J_0\), \(H=H_0\), and
\(J_\alpha=K_\alpha\circ g\),
\(H_\alpha=K_\alpha\circ g^2\) on the off-diagonal support.

\begin{theorem}[Asymptotic action variance and its circulation Hessian]
\label{thm:v22-variance-susceptibility}
Let \(\psi(\eta,k)\) be the Perron exponent of the jump-reward tilted
generator.  At \(k=0\),
\begin{equation}
 \boxed{
 \mathcal V(\eta)
 :=\partial_k^2\psi(\eta,0)
 =\pi H_\eta\mathbf1
  +2\pi J_\eta G_\eta J_\eta\mathbf1,}
 \label{eq:v22-variance-rate}
\end{equation}
where
\(G_\eta=Q(-L_\eta)^{-1}Q\).
Moreover,
\begin{equation}
 \partial_{\eta^\alpha}G_\eta\big|_0
 =G K_\alpha G,
 \label{eq:v22-Ga}
\end{equation}
\begin{equation}
 \partial_{\eta^\alpha}\partial_{\eta^\beta}G_\eta\big|_0
 =G K_\alpha G K_\beta G
  +G K_\beta G K_\alpha G.
 \label{eq:v22-Gab}
\end{equation}
For the affine deformation \eqref{eq:v22-Leta}, the quadratic circulation
susceptibility of the asymptotic reward variance is
\begin{equation}
 \boxed{
 \begin{aligned}
 \frac12\partial_{\alpha\beta}\mathcal V(0)
 =\pi\Big[&J_\alpha GJ_\beta
          +J_\beta GJ_\alpha\\
 &+J_\alpha GK_\beta GJ
  +J_\beta GK_\alpha GJ\\
 &+JGK_\alpha GJ_\beta
  +JGK_\beta GJ_\alpha\\
 &+\frac12JG(K_\alpha GK_\beta+K_\beta GK_\alpha)GJ
 \Big]\mathbf1 .
 \end{aligned}}
 \label{eq:v22-susceptibility}
\end{equation}
The linear term from \(\pi H_\eta\mathbf1\) has zero second
\(\eta\)-derivative.  If the reflected symmetry of
Theorem~\ref{thm:v22-parity-router} holds and \(g\) is reversal even, then
\begin{equation}
 \partial_{\eta^\alpha}\mathcal V(0)=0,
 \label{eq:v22-var-linear-zero}
\end{equation}
so \eqref{eq:v22-susceptibility} is the first allowed reflected circulation
response of this path statistic.
\end{theorem}

\begin{proof}
For a simple Perron eigenvalue, second-order perturbation in the tilt gives
\[
 \partial_k^2\psi(\eta,0)
 =\pi H_\eta\mathbf1
 +2\pi J_\eta Q(-L_\eta)^{-1}QJ_\eta\mathbf1,
\]
which is \eqref{eq:v22-variance-rate}.  On the centered subspace,
\(G_\eta=(-L_\eta)^{-1}\); differentiating the inverse gives
\eqref{eq:v22-Ga}--\eqref{eq:v22-Gab}.  Differentiate
\(\pi J_\eta G_\eta J_\eta\mathbf1\) twice.  Since \(J_\eta\) is affine,
there are no second derivatives of \(J_\eta\); the seven product-rule terms
are exactly those displayed in \eqref{eq:v22-susceptibility}.  The
\(H_\eta\) term is affine and therefore contributes no Hessian.  The parity
statement follows from \eqref{eq:v22-psi-even} after two derivatives in the
tilt variable.
\end{proof}

\begin{firewall}[The susceptibility is not a prediction until the physical reward lands]
Equation~\eqref{eq:v22-susceptibility} is a finite compiler.  It becomes a
Renewal phenomenology matrix only when \(L\), the three residual generator
marks \(K_\alpha\), and a physical non-coboundary gravity--matter reward
\(g\) are populated on the same accepted path law.  Substituting stochastic
entropy or an arbitrary positive pair writer for the common-action packet is
not licensed.
\end{firewall}

\subsection{Routing through the inherited three-cylinder action-response packet}
\label{sec:v22-three-cylinder}

The Grand Tensor already provides a finite stopping rule for a reflected
physical action-response bank.  Let \(F\) be the complete bounded physical
writer bank, let
\(\Gamma_P(F)\) be its conditional pair-action writer, and put
\begin{equation}
 M_n^\Gamma
 =\mathbb E[\gamma_F(X_n)\gamma_F(X_0)^*],
 \qquad n=0,1,2,
 \label{eq:v22-M012}
\end{equation}
with the inherited centered \(\gamma_F\).  Its supported compression and first
memory innovation are
\begin{align}
 A_\Gamma
 &=(M_0^\Gamma)^{\dagger/2}M_1^\Gamma
   (M_0^\Gamma)^{\dagger/2},
 \label{eq:v22-Agamma}\\
 \mathbb I_\Gamma^{(1)}
 &=(M_0^\Gamma)^{\dagger/2}
 \bigl[M_2^\Gamma
  -M_1^\Gamma(M_0^\Gamma)^\dagger M_1^\Gamma\bigr]
 (M_0^\Gamma)^{\dagger/2}.
 \label{eq:v22-Igamma}
\end{align}

\begin{theorem}[Three-cylinder routing of the first nontrivial selector]
\label{thm:v22-three-cylinder-router}
Assume the common-action incidence residual for the selected reflected writer
bank has been independently closed.
\begin{enumerate}[label=\textup{(R\arabic*)},leftmargin=3em]
\item If \(\mathbb I_\Gamma^{(1)}\ne0\), the current action-response bank is
      not temporally closed.  The inherited polar writer
      \(J_\Gamma^{\rm new}\) must be appended before any Green/path
      susceptibility such as \eqref{eq:v22-susceptibility} is interpreted as
      complete.
\item If \(\mathbb I_\Gamma^{(1)}=0\), the supported writer range reduces the
      accepted transfer and all higher correlation cylinders are generated by
      powers of \(A_\Gamma\).  On the strictly centered contractive part,
      the integrated correlation/Green series is therefore the finite matrix
      function
      \begin{equation}
       \boxed{
       \sum_{n\ge0}A_\Gamma^n
       =(I-A_\Gamma)^{-1}.}
       \label{eq:v22-head-Green}
      \end{equation}
      Hence the path susceptibility is computable from the finite reducing
      head together with the residual generator marks.
\end{enumerate}
\end{theorem}

\begin{proof}
The first item is exactly the meaning of the positive innovation
\eqref{eq:v22-Igamma}: it is the Gram of the part of one-step evolution which
leaves the current writer range.  For the second item, vanishing of the
innovation makes that range reducing, so the inherited three-cylinder theorem
gives
\(J_\Gamma^*P^nJ_\Gamma=A_\Gamma^n\) for every \(n\ge0\).  On the centered
subspace with spectral radius below one the geometric series converges to
\((I-A_\Gamma)^{-1}\).  This is precisely the finite correlation resolvent
needed by a Green/path susceptibility calculation.
\end{proof}

\begin{remark}[Why the positive pair action is useful but cannot replace the signed common-action row]
The pair action \(\mathcal A_F\) is a physically meaningful extensive
Dirichlet/response observable once its writer bank and transfer are acquired.
Its reversal evenness makes it an appropriate second-order circulation probe,
not a replacement for the missing signed first-order common-action incidence.
The two roles are complementary: the signed row tests direction-sensitive
action provenance, whereas the pair-action packet tests reflected dissipation
and temporal closure.
\end{remark}

\subsection{P4e decision after the upstream correction}
\label{sec:v22-P4e-decision}

V21's first-jump stop can now be refined into an exact decision tree.

\begin{theorem}[P4e action-selector alternative]
\label{thm:v22-P4e-alternative}
For the three residual race coordinates of V19--V21, exactly one of the
following routes is appropriate for a declared physical action packet.
\begin{enumerate}[label=\textup{(A\arabic*)},leftmargin=3em]
\item \textbf{Directed non-coboundary row.}  A same-history physical action
      writer has a reversal-odd component which is not the increment of one
      scalar state function.  Use the V21 linear three-column incidence matrix
      immediately.
\item \textbf{Exact scalar-action increment.}  The writer is
      \(S(y)-S(x)\).  Its extensive Feynman--Kac selector is identically zero by
      Theorem~\ref{thm:v22-coboundary-no-go}; do not promote it to longer paths.
\item \textbf{Reflected even action-response row.}  The physical reward is
      reversal even.  Its linear circulation selector vanishes by
      Theorem~\ref{thm:v22-parity-router}.  Close the three-cylinder writer
      packet and use the second-order path susceptibility
      \eqref{eq:v22-susceptibility} as the first admissible circulation
      diagnostic.
\item \textbf{Unloaded physical row.}  If neither a directed non-coboundary
      common-action reward nor a reflected physical action-response bank has
      been acquired on the same accepted path law, the P4 matrix remains
      unpopulated.  Missing data are not represented by zero.
\end{enumerate}
\end{theorem}

\begin{proof}
The four cases are the exhaustive combination of (i) whether the action row is
loaded, (ii) whether it is an exact coboundary, and (iii) its reversal parity.
The corresponding conclusions are V21, Theorem
\ref{thm:v22-coboundary-no-go}, Theorems
\ref{thm:v22-parity-router}--\ref{thm:v22-three-cylinder-router}, and the
V19--V21 no-missing-map rule.
\end{proof}

For the currently populated decoded/PCH scalar action on the scheduler-only
race cylinder, route \textup{(A2)} applies after the trivial first-jump zero of
Theorem~\ref{thm:v22-scheduler-silence}.  For the full accepted/reflected
Einstein--SM action-response packet, route \textup{(A3)} becomes available only
after the same-history accepted transfer and common-action incidence residual
are populated.  The historical locked global-cylinder export remains a
separate acquisition; the explicit active RPESM construction provides a model
carrier but does not numerically identify that historical row.

\subsection{V22 anti-circularity audit}
\label{sec:v22-audit}

\begin{enumerate}[leftmargin=2.4em]
\item \textbf{The deterministic PCH writer is not confused with a scheduler
jump.}  Its field tuple is held fixed on the fast race mark, giving a proved
zero first-jump row rather than an assumed zero physical incidence.
\item \textbf{Longer paths do not rescue an exact scalar action increment.}
Its Feynman--Kac tilt is a similarity transform and the SCGF is identically
zero.
\item \textbf{Reversal-even action response is not forced into a linear
selector.}  Circulation parity kills the first derivative; V22 moves to the
first allowed second-order path statistic.
\item \textbf{The second-order compiler is finite and explicit.}  It uses only
the accepted generator, its three residual marks, the physical reward matrices,
and one centered Green operator.
\item \textbf{Temporal incompleteness is tested before a Green coefficient is
promoted.}  A nonzero three-cylinder innovation returns the exact missing
writer to append.
\item \textbf{No stochastic control row is substituted for the physical
Einstein--SM action.}  The susceptibility remains conditional on the physical
writer/transfer incidence actually landing.
\end{enumerate}

\section{V22 status ledger}
\label{sec:v22-status}

\subsection*{Proved in V22}
\begin{enumerate}[leftmargin=2.2em]
\item The populated decoded/PCH scalar common action is independent of the fast
scheduler-only first-jump residual mark; its V21 reversal contrast, direct
residual derivative and three-column first-jump selector are exactly zero on
that cylinder.
\item A signed increment of one scalar common action is an exact Markov
coboundary.  Its Feynman--Kac tilted generator is similar to the untilted
generator, so its scaled cumulant generator and every mixed regulator/tilt
derivative vanish identically.  Promoting the same increment to more jumps
cannot create an extensive selector.
\item The residual circulation score is reversal odd.  Every reversal-even
writer has zero linear residual incidence on a reflected base, and its tilted
pressure is even in the circulation coordinate.
\item The first two tilt matrices and centered Green operator give the exact
asymptotic jump-reward variance
\(\mathcal V=\pi H\mathbf1+2\pi JGJ\mathbf1\).  For an affine circulation
deformation with fixed stationary law, V22 derives the explicit quadratic
circulation susceptibility \eqref{eq:v22-susceptibility}.
\item The inherited three-cylinder action-response innovation is the correct
finite gate before interpreting that susceptibility: a nonzero innovation
returns a new physical writer, while a zero innovation makes the correlation
Green function a finite rational function of the reducing compression.
\item P4e therefore has an exact router between a directed non-coboundary
linear selector, a coboundary no-go, a reflected second-order action-response
selector, and an unloaded-row stop.
\end{enumerate}

\subsection*{Inherited}
\begin{enumerate}[leftmargin=2.2em]
\item V16--V21: physical calibration and RG covariance, the no-free-Wilson
prediction theorem, the pilot physical quotient, the populated race incidence,
the three residual circulation coordinates and their exact first-jump score
compiler.
\item The Grand Tensor protected scalar PCH/common first-jet action writer,
finite common-action exactness, accepted/common-action occurrence compiler,
finite Feynman--Kac theorem, and reflected three-cylinder action-response
closure.
\item The active RPESM same-history event architecture carrying fast root,
initial/terminal fields, update/reverse labels, action and line/locality marks;
the historical locked cylinder remains a separate export/acquisition problem.
\end{enumerate}

\subsection*{Conditional}
\begin{enumerate}[leftmargin=2.2em]
\item Formula \eqref{eq:v22-susceptibility} assumes an affine residual
generator deformation with a common stationary law, as holds for the idealized
circulation slice used to derive the compiler.  A general accepted kernel with
\(\eta\)-dependent stationary law requires the ordinary full Perron-eigenvector
perturbation terms.
\item The reflected second-order selector becomes physical only after the
accepted transfer, complete action-response bank and common-action incidence
residual are all populated on one same-history path law.
\item A real Pfaffian orientation and all V16 fixed finite loop/EFT/source and
positive-reference-scale qualifications remain in force.
\end{enumerate}

\subsection*{Open beyond V22}
\begin{enumerate}[leftmargin=2.2em]
\item Populate one non-coboundary full-field gravity--matter action observable
on the active same-history accepted path law.  The preferred reflected choice
is the complete physical action-response bank after common-action incidence is
verified.
\item Evaluate its three-cylinder packet
\((M_0^\Gamma,M_1^\Gamma,M_2^\Gamma)\).  If the innovation is nonzero, append
its polar missing writer; if zero, compute the finite reducing Green head.
\item Contract the three residual generator marks with that physical path
packet using \eqref{eq:v22-susceptibility}, then push the resulting finite
selector through \(\mathsf K_{\rm phys}\) and inspect the rank/left nullspace
before fitting Wilson coefficients.
\item If a genuinely directed non-coboundary signed action row is acquired,
return to the cheaper V21 linear selector instead of the quadratic reflected
compiler.
\item The parallel P2/P3 line remains available after a predictive boundary
relation is found: higher-loop native beta rows, threshold matching,
standard-scheme conversion and numerical RG evolution.
\end{enumerate}

\subsection*{Exact next load-bearing theorem}
\begin{center}
\fbox{\parbox{0.92\textwidth}{
\textbf{P4f: physical non-coboundary action-response landing.}
On one explicit same-history accepted RPESM path law, populate the complete
bounded gravity--matter writer bank and verify its incidence with the
independently reconstructed common action.  Compute the three stationary
action-response cylinders.  A nonzero first innovation returns the exact
missing writer to append; a zero innovation closes the reducing action head.
Then evaluate the residual-circulation path susceptibility of
\eqref{eq:v22-susceptibility} (or the V21 linear selector if a directed
non-coboundary signed row is present), push it through the V18 physical Wilson
quotient, and compute the first actual predictivity rank/nullspace.  Exact
scalar endpoint-action increments are excluded from this target by the V22
coboundary no-go.}}
\end{center}

\section*{Append-only continuation marker: V22}
\addcontentsline{toc}{section}{Append-only continuation marker: V22}
\textbf{End of V22.}  The complete V1--V21 body above is retained verbatim.
V22 performs the upstream escalation demanded by V21.  The populated decoded
PCH/common first-jet action is silent on the scheduler-only residual race mark,
and a signed increment of any scalar common action is an exact coboundary whose
Feynman--Kac pressure vanishes identically; increasing the path length with the
same increment would therefore be circular.  On a reflected branch the
physical pair action-response writer is reversal even, so the residual
circulation has no linear incidence.  V22 derives the first admissible
second-order path compiler: an explicit Green-operator formula for the
quadratic circulation susceptibility of the asymptotic action-reward variance,
and routes it through the inherited three-cylinder temporal-closure test.
The next phenomenology gate is thus not another abstract action increment but a
physically acquired non-coboundary gravity--matter action-response bank on one
same-history accepted path law.


\clearpage
\section{V23 continuation: action--dynamics incidence, refresh calibration, and the marked circulation gate}
\label{sec:v23-continuation}

V22 identified the first nontrivial reflected circulation diagnostic, but it
left one hypothesis too compressed.  The projective RPESM regulator supplies,
on one literal same-history event law, the full field record, a reflection
positive stationary path law, action/stress/contact writers, and the finite
action-response cylinders required by the universal compilers.  What it does
\emph{not} do automatically is identify the three residual race coordinates
\(\eta\) of V19--V22 with a deformation of the accepted full-field generator.
The displayed RPESM generator
\begin{equation}
 L=\rho(R-I)+L^{\rm loc}
 \label{eq:v23-RP-generator}
\end{equation}
is chosen after the stationary cylinder has been fixed.  Its local
conditional-expectation part is not, merely by the population theorem, the
push-forward of the residual circulation in the fast scheduler.

This distinction is load bearing.  A complete physical writer bank and a
closed three-cylinder head can coexist with
\begin{equation}
 \partial_{\eta^\alpha}L=0.
 \label{eq:v23-zero-operation-incidence}
\end{equation}
In that case every V22 path susceptibility with respect to \(\eta\) is zero,
not because the writer is absent but because the microscopic circulation has
not been made incident with the accepted operation.  V23 therefore separates
three finite questions:
\begin{equation}
 \boxed{
 \text{writer population}
 \quad\longrightarrow\quad
 \text{action--transfer incidence}
 \quad\longrightarrow\quad
 \text{circulation--operation incidence}.}
 \label{eq:v23-three-incidences}
\end{equation}
The first is existentially populated by the inherited RPESM construction.  The
second is a finite common-action calibration residual.  The third is the new
three-column dynamical derivative
\begin{equation}
 \boxed{K_\alpha:=\partial_{\eta^\alpha}L_{\rm acc}\big|_{\eta=0}.}
 \label{eq:v23-Kalpha}
\end{equation}
Only after all three levels are populated may the V22 residual-circulation
susceptibility be interpreted as phenomenological matching data.

\subsection{A completely explicit same-history control branch: pure refresh}
\label{sec:v23-pure-refresh}

We first isolate an exact branch of the populated RPESM regulator on which the
action-response cylinders can be computed for an arbitrary bounded physical
writer bank.  Let \((\Omega,\mu)\) be one finite-cutoff RPESM field space with
its interacting stationary law and let
\begin{equation}
 Rf:=\int f\,\dd\mu,
 \qquad
 Q:=I-R.
 \label{eq:v23-RQ}
\end{equation}
Fix \(\rho>0\) and take the allowed refresh-only member
\begin{equation}
 \boxed{L_\rho=\rho(R-I)=-\rho Q.}
 \label{eq:v23-refresh-generator}
\end{equation}
At a physical sampling duration \(\tau>0\),
\begin{equation}
 \boxed{
 P_{\rho,\tau}:=e^{\tau L_\rho}
 =R+rQ,
 \qquad
 r:=e^{-\rho\tau}\in(0,1).}
 \label{eq:v23-refresh-transfer}
\end{equation}
This branch retains the full interacting stationary cylinder \(\mu\); it
changes only the accepted temporal transport.

Let \(F=(F_1,\ldots,F_d)\) be any bounded physical writer bank and let the
centered action-response synthesis \(S_\Gamma\), covariance
\(M_0^\Gamma\), supported space \(\mathcal E_\Gamma\), and isometry
\(J_\Gamma\) be those of V22.  In particular the columns of
\(S_\Gamma\) are centered.

\begin{theorem}[Exact refresh three-cylinder packet]
\label{thm:v23-refresh-three-cylinder}
For the transfer \eqref{eq:v23-refresh-transfer}, every bounded physical
action-response bank satisfies
\begin{equation}
 \boxed{
 M_n^\Gamma=r^nM_0^\Gamma
 \qquad(n\ge0).}
 \label{eq:v23-refresh-Mn}
\end{equation}
Consequently
\begin{equation}
 \boxed{A_\Gamma=rI_{\mathcal E_\Gamma}.}
 \label{eq:v23-refresh-A}
\end{equation}
\begin{equation}
 \boxed{\mathbb I_\Gamma^{(1)}=0.}
 \label{eq:v23-refresh-I}
\end{equation}
\begin{equation}
 \boxed{J_\Gamma^*P_{\rho,\tau}^nJ_\Gamma=r^nI_{\mathcal E_\Gamma}.}
 \label{eq:v23-refresh-reducing}
\end{equation}
Thus the complete action-response writer head is temporally reducing on the
pure-refresh branch, independently of the detailed interacting stationary
measure and independently of which bounded writer bank is chosen.
\end{theorem}

\begin{proof}
Because \(S_\Gamma\) is centered,
\(RS_\Gamma=0\), hence
\(P_{\rho,\tau}S_\Gamma=rS_\Gamma\).  Iteration gives
\(P_{\rho,\tau}^nS_\Gamma=r^nS_\Gamma\), and therefore
\[
 M_n^\Gamma
 =S_\Gamma^*P_{\rho,\tau}^nS_\Gamma
 =r^nS_\Gamma^*S_\Gamma
 =r^nM_0^\Gamma.
\]
Substitution into the V22 three-cylinder formulas gives
\eqref{eq:v23-refresh-A} and
\[
 (M_0^\Gamma)^{\dagger/2}
 [r^2M_0^\Gamma-r^2M_0^\Gamma(M_0^\Gamma)^\dagger M_0^\Gamma]
 (M_0^\Gamma)^{\dagger/2}=0
\]
on the supported range.  The final identity follows in the same way.
\end{proof}

\begin{corollary}[Refresh closure does not imply physical action incidence]
\label{cor:v23-refresh-not-action}
Temporal completeness of the writer bank, even in the strongest form
\eqref{eq:v23-refresh-reducing}, does not identify the transfer Dirichlet
matrix with the independently reconstructed Einstein--SM common-action matrix.
That equality is a separate calibration equation.
\end{corollary}

\begin{proof}
Theorem~\ref{thm:v23-refresh-three-cylinder} used only the stationary law and
the chosen refresh clock.  It did not use any Hessian or coefficient of the
independently reconstructed common action.  Therefore the action-incidence row
cannot follow from the temporal closure identity alone.
\end{proof}

\subsection{Exact common-action calibration of the refresh clock}
\label{sec:v23-refresh-calibration}

The preceding control branch permits a sharp action-incidence test.  Let
\begin{equation}
 C_F:=Y_F^*QY_F\succeq0
 \label{eq:v23-CF}
\end{equation}
be the centered static writer covariance, and let
\(G_{\rm act}\succeq0\) denote the independently reconstructed common-action
matrix on the same writer coefficients.  For the pure-refresh transfer,
\begin{equation}
 \boxed{
 Y_F^*(I-P_{\rho,\tau})Y_F
 =(1-r)C_F.}
 \label{eq:v23-refresh-Dirichlet}
\end{equation}
Thus one scalar clock can match the action only along the ray spanned by
\(C_F\).

\begin{theorem}[Refresh-clock action-incidence criterion]
\label{thm:v23-refresh-action-incidence}
Assume \(C_F\ne0\).  There exists a finite refresh rate \(\rho>0\) satisfying
\begin{equation}
 G_{\rm act}=Y_F^*(I-P_{\rho,\tau})Y_F
 \label{eq:v23-refresh-action-match}
\end{equation}
if and only if
\begin{equation}
 \boxed{
 G_{\rm act}=aC_F
 \quad\text{for one }a\in(0,1).}
 \label{eq:v23-refresh-ray}
\end{equation}
Then
\begin{equation}
 \boxed{
 \rho=-\tau^{-1}\log(1-a).}
 \label{eq:v23-refresh-rho}
\end{equation}
More generally, the least-squares scalar action mismatch over
\(a\in[0,1]\) is
\begin{equation}
 \boxed{
 \Delta_{\rm ref}^{\min}
 =\|G_{\rm act}-a_*C_F\|_{\rm HS}^2,}
 \label{eq:v23-refresh-minres}
\end{equation}
where
\begin{equation}
 \boxed{
 a_*
 =\Pi_{[0,1]}
 \frac{\langle C_F,G_{\rm act}\rangle_{\rm HS}}
      {\|C_F\|_{\rm HS}^2}.}
 \label{eq:v23-refresh-astar}
\end{equation}
A positive residual proves that a scalar global refresh clock cannot realize
the independently reconstructed action-response matrix on that writer bank;
one must use additional local operation rows rather than re-fit the action.
\end{theorem}

\begin{proof}
Equation \eqref{eq:v23-refresh-Dirichlet} gives the first equivalence with
\(a=1-e^{-\rho\tau}\in(0,1)\).  Solving for \(\rho\) gives
\eqref{eq:v23-refresh-rho}.  The minimizer of the quadratic function
\(\|G_{\rm act}-aC_F\|_{\rm HS}^2\) on the closed interval \([0,1]\) is the
orthogonal scalar projection clipped to that interval, which is
\eqref{eq:v23-refresh-astar}.  The last statement follows because
\(G_{\rm act}\) is independently fixed and therefore cannot be changed to
force the residual to vanish.
\end{proof}

\begin{firewall}[Clock calibration is not Wilson matching]
The scalar \(\rho\) is a physical temporal normalization of the accepted
operation.  Equation \eqref{eq:v23-refresh-action-match} calibrates that clock
against an independently reconstructed action-response matrix.  It does not
select a gravity--matter EFT Wilson coefficient, and failure of the ray test
cannot be repaired by changing \(G_{\rm act}\) after the transfer has been
observed.
\end{firewall}

\subsection{De-refreshing does not change the temporal-closure decision}
\label{sec:v23-derefresh}

The pure-refresh calculation extends to the full RPESM form
\eqref{eq:v23-RP-generator}.  Assume the local part is self-adjoint,
\(L^{\rm loc}\mathbf1=0\), and \(\mu L^{\rm loc}=0\), as in the reflected
slab construction.  Then \(R L^{\rm loc}=L^{\rm loc}R=0\), and on the centered
subspace
\begin{equation}
 L=-\rho I+L^{\rm loc}.
 \label{eq:v23-centered-L}
\end{equation}
Hence for sampling duration \(\tau\),
\begin{equation}
 \boxed{
 P_\tau
 =R+e^{-\rho\tau}e^{\tau L^{\rm loc}}Q.}
 \label{eq:v23-derefresh-transfer}
\end{equation}

\begin{theorem}[Exact de-refreshing of the three-cylinder innovation]
\label{thm:v23-derefresh-innovation}
Let
\begin{equation}
 \widehat M_n^\Gamma
 :=S_\Gamma^*e^{n\tau L^{\rm loc}}S_\Gamma.
 \label{eq:v23-Mhat}
\end{equation}
Then
\begin{equation}
 \boxed{M_n^\Gamma=e^{-n\rho\tau}\widehat M_n^\Gamma.}
 \label{eq:v23-M-scale}
\end{equation}
Consequently
\begin{align}
 A_\Gamma
 &=e^{-\rho\tau}\widehat A_\Gamma,
 \label{eq:v23-A-scale}\\
 \boxed{
 \mathbb I_\Gamma^{(1)}
 =e^{-2\rho\tau}\widehat{\mathbb I}_\Gamma^{(1)}.}
 \label{eq:v23-I-scale}
\end{align}
In particular,
\begin{equation}
 \boxed{
 \operatorname{rank}\mathbb I_\Gamma^{(1)}
 =\operatorname{rank}\widehat{\mathbb I}_\Gamma^{(1)},}
 \label{eq:v23-I-rank}
\end{equation}
so the auxiliary global refresh can neither create nor erase an action-response
memory writer.  The temporal-closure decision belongs entirely to the local
accepted operation after exact de-refreshing.
\end{theorem}

\begin{proof}
Equation \eqref{eq:v23-derefresh-transfer} and centeredness of
\(S_\Gamma\) give \eqref{eq:v23-M-scale}.  Substitution in the definitions of
\(A_\Gamma\) and \(\mathbb I_\Gamma^{(1)}\) yields one factor
\(e^{-\rho\tau}\) and \(e^{-2\rho\tau}\), respectively.  The latter scalar is
strictly positive, so the innovation rank is unchanged.
\end{proof}

\subsection{The canonical residual race deformation does not yet deform the accepted field operation}
\label{sec:v23-operation-incidence}

We now return to the three residual race coordinates \(\eta\) of V19--V22.
They alter the fast root winner/waiting law while leaving the calibrated drift,
diffusion and density coordinates fixed.  On the canonical RPESM population
construction, the stationary field cylinder and the accepted generator
\eqref{eq:v23-RP-generator} are subsequently chosen from the common action and
conditional-expectation packet.  No theorem in the inherited population step
identifies the residual race parameter with a coefficient of
\(L^{\rm loc}\).

\begin{theorem}[Canonical RPESM operation-incidence zero]
\label{thm:v23-canonical-Kzero}
Consider the canonical residual deformation in which the V20 race rates vary
with \(\eta\), while the P1-calibrated decoded geometry, the common stationary
field cylinder, the refresh clock and the coefficients of
\(L^{\rm loc}\) are held fixed.  Then
\begin{equation}
 \boxed{
 K_\alpha
 :=\partial_{\eta^\alpha}L_{\rm acc}\big|_0
 =0
 \qquad(\alpha=1,2,3).}
 \label{eq:v23-Kzero}
\end{equation}
Every residual-circulation derivative of an action-response path statistic
which depends on \(\eta\) only through the accepted field generator therefore
vanishes on this branch.  In particular, the V22 quadratic susceptibility is
zero.
\end{theorem}

\begin{proof}
By construction, the residual race parameter is varied only in the fast
scheduler law.  The decoded variables used to build the field cylinder are
held fixed by P1 and, along the V19 residual subspace, are unchanged by
\(\eta\).  The displayed accepted generator is held fixed as part of the
canonical population construction.  Therefore its derivative with respect to
\(\eta\) is zero.  The V22 susceptibility is a polynomial in the marked
generator derivatives \(K_\alpha\), so it vanishes when all three are zero.
\end{proof}

\begin{corollary}[Population complete is not selector complete]
\label{cor:v23-population-not-selector}
The RPESM population theorem can simultaneously supply a complete same-history
action writer bank, a reflected transfer, and closed three-cylinder moments
without supplying any nonzero P4 coupling of the V19 residual race packet to
the accepted physical operation.  Temporal writer completeness is therefore
logically weaker than microscopic selector incidence.
\end{corollary}

\begin{proof}
Use the populated writer/transfer packet together with
Theorem~\ref{thm:v23-canonical-Kzero}.  The first concerns existence of the
physical rows; the second concerns the derivative of the accepted operation
with respect to a particular residual microscopic parameter.  They are
different data.
\end{proof}

\subsection{The law--material short identifies what a nonzero coupling must contain}
\label{sec:v23-law-material}

V13--V22 repeatedly distinguish motion of a probability law from motion of a
physical writer in a material frame.  That distinction gives a direct
selector criterion here.

Let \(t\mapsto\mu_t\) be the residual scheduler-law path, let
\(\psi_t=\sqrt{\dd\mu_t/\dd\mathrm{vol}}\), and let \(f_t\) be a physical
writer.  Write
\begin{equation}
 \Phi_t=\psi_tf_t,
 \qquad
 \dot\Phi_t=L_f+M_f,
 \label{eq:v23-law-material-split}
\end{equation}
where \(L_f\) is the law-motion column and \(M_f\) the material derivative
column in the inherited fixed-frame decomposition.  If \(P_L\) is the
orthogonal projection onto the already retained law-motion range, then
\begin{equation}
 \boxed{
 (L_f+M_f)^*(I-P_L)(L_f+M_f)
 =M_f^*(I-P_L)M_f.}
 \label{eq:v23-law-material-short}
\end{equation}

\begin{proposition}[Decoded writers have zero residual material incidence]
\label{prop:v23-decoded-material-zero}
Suppose \(f_t\) factors only through P1-calibrated decoded variables which are
constant along the residual race direction, and is constant along the
corresponding scheduler material flow once those decoded variables are fixed.
Then
\begin{equation}
 \boxed{M_f=0}
 \label{eq:v23-Mf-zero}
\end{equation}
and the complete residual derivative of \(\Phi_t\) is already contained in the
law-motion source range.  After shorting the scheduler law score, the writer
contributes no new microscopic action selector.
\end{proposition}

\begin{proof}
The material column is the sum of the explicit writer derivative and its
material transport derivative.  Both vanish under the stated hypotheses.
Equation \eqref{eq:v23-law-material-short} then gives zero residual after the
law-motion range has been removed.
\end{proof}

\begin{firewall}[A same-record covariance can still be nuisance ancestry]
A nonzero raw covariance between a physical writer and the residual scheduler
score does not by itself prove a new action selector.  If that covariance lies
in the already retained law-motion range, it is source ancestry of the changing
scheduler law.  P4 requires a nonzero material/operation residual after the
law source has been shorted.
\end{firewall}

\subsection{The missing P4 datum is a three-column operation derivative}
\label{sec:v23-K-compiler}

The correct replacement for the now-closed ``missing writer'' diagnosis is to
load the physical operation derivative itself.  Let a candidate same-history
accepted family have
\begin{equation}
 L_\eta=L+\eta^\alpha K_\alpha+O(|\eta|^2),
 \qquad
 K_\alpha\mathbf1=0,
 \label{eq:v23-Leta-physical}
\end{equation}
with the stationary-law derivative treated in the full Perron/Kato packet when
it is nonzero.  The three \(K_\alpha\) are not free Wilson coefficients; they
are microscopic derivatives of the accepted transition rule with respect to
the already identified residual race coordinates.

Even after the base action-response bank reduces the accepted transfer, the
marked generators can leave that head.  This has its own exact finite Gram.
Let \(J_\Gamma\) and \(\Pi_\Gamma\) be the V22 action-response isometry and
range projection, and define
\begin{equation}
 B_\alpha:=(I-\Pi_\Gamma)K_\alpha J_\Gamma.
 \label{eq:v23-Balpha}
\end{equation}
Stack these as
\begin{equation}
 B_K:\mathcal E_\Gamma\otimes\mathbb C^3\to\mathcal H,
 \qquad
 B_K(c_1,c_2,c_3)=\sum_{\alpha=1}^3B_\alpha c_\alpha.
 \label{eq:v23-BK}
\end{equation}

\begin{theorem}[Marked circulation-to-operation closure]
\label{thm:v23-marked-closure}
The block operator
\begin{equation}
 \boxed{
 \mathbb I_{\Gamma,K}
 :=B_K^*B_K
 =\bigl[
 J_\Gamma^*K_\alpha^*(I-\Pi_\Gamma)K_\beta J_\Gamma
 \bigr]_{\alpha,\beta=1}^3
 \succeq0}
 \label{eq:v23-marked-Gram}
\end{equation}
is the exact first operation-incidence innovation of the residual circulation
relative to the closed action-response head.  Moreover,
\begin{equation}
 \boxed{
 \mathbb I_{\Gamma,K}=0
 \iff
 K_\alpha\Ran J_\Gamma\subseteq\Ran J_\Gamma
 \quad\text{for all }\alpha.}
 \label{eq:v23-marked-zero}
\end{equation}
If the rank is positive, the source-minimal missing marked writer is the polar
range
\begin{equation}
 \boxed{
 J_{\Gamma,K}^{\rm new}
 =B_K(\mathbb I_{\Gamma,K})^{\dagger/2}}
 \label{eq:v23-marked-new}
\end{equation}
on \(\supp\mathbb I_{\Gamma,K}\).  It must be appended before a finite-head
V22 susceptibility is declared complete.
\end{theorem}

\begin{proof}
The displayed block matrix is a Gram by construction and is therefore
positive.  It vanishes exactly when \(B_K=0\), equivalently when every
\((I-\Pi_\Gamma)K_\alpha J_\Gamma\) vanishes.  This is precisely the marked
invariance condition in \eqref{eq:v23-marked-zero}.  Polar normalization of
the nonzero range gives \eqref{eq:v23-marked-new}; Moore--Penrose inversion
restricts it to the supported coefficient space, exactly as in the ordinary
three-cylinder innovation.
\end{proof}

\begin{corollary}[Finite compressed residual generator]
\label{cor:v23-compressed-K}
On the zero branch of \eqref{eq:v23-marked-Gram}, define
\begin{equation}
 \widehat K_\alpha
 :=J_\Gamma^*K_\alpha J_\Gamma.
 \label{eq:v23-Khat}
\end{equation}
Then
\begin{equation}
 \boxed{K_\alpha J_\Gamma=J_\Gamma\widehat K_\alpha.}
 \label{eq:v23-K-intertwine}
\end{equation}
If the base action-response head also reduces the accepted generator, every
word formed from the reduced Green operator and the marked generators which
starts in the action head can be evaluated using only the finite compressed
matrices \(\widehat L\) and \(\widehat K_\alpha\), provided the corresponding
reward-source vectors have first been verified to lie in the same head.
\end{corollary}

\begin{proof}
The marked-zero condition gives
\(K_\alpha J_\Gamma=\Pi_\Gamma K_\alpha J_\Gamma\).  Since
\(\Pi_\Gamma=J_\Gamma J_\Gamma^*\), this is
\eqref{eq:v23-K-intertwine}.  The final statement follows by repeated
intertwining and the same reduction of the centered Green operator used in
V22.
\end{proof}

\subsection{P4f status after V23}
\label{sec:v23-P4f-status}

The phenomenology bottleneck can now be stated without ambiguity.

\begin{theorem}[P4f population/incidence alternative]
\label{thm:v23-P4f-alternative}
For the explicit RPESM regulator and the V19 residual race packet:
\begin{enumerate}[label=\textup{(P\arabic*)},leftmargin=3em]
\item the same-history finite field/action/source writer bank and a positive
      reflected stationary path law are populated;
\item temporal closure of a selected action-response bank is decided by the
      ordinary three-cylinder innovation and is independent of the auxiliary
      global refresh after de-refreshing;
\item common-action incidence is a separate finite matrix calibration, with
      the pure-refresh branch obeying the exact ray test
      \eqref{eq:v23-refresh-ray};
\item the canonical residual race deformation supplies no nonzero derivative
      of the accepted field operation, so \(K_\alpha=0\) and the V22
      circulation susceptibility vanishes on that branch;
\item a nonzero P4 selector therefore requires a physically populated
      circulation-to-operation derivative \(K_\alpha\), or equivalently a
      nonzero material/operation residual after the scheduler law source is
      shorted;
\item once \(K_\alpha\) lands, the marked Gram
      \eqref{eq:v23-marked-Gram} decides whether the existing action-response
      bank is sufficient for the residual perturbation or returns the exact
      missing marked writer to append.
\end{enumerate}
\end{theorem}

\begin{proof}
Item (P1) is the inherited population construction; (P2) is
Theorem~\ref{thm:v23-derefresh-innovation}; (P3) is
Theorem~\ref{thm:v23-refresh-action-incidence}; (P4) is
Theorem~\ref{thm:v23-canonical-Kzero}; (P5) is
Proposition~\ref{prop:v23-decoded-material-zero} together with the V19--V22
selector definition; and (P6) is Theorem~\ref{thm:v23-marked-closure}.
\end{proof}

The correct current status is therefore
\begin{equation}
 \boxed{
 \texttt{WRITER\_BANK\_POPULATED;\ 
 ACTION\_INCIDENCE\_FINITE;\ 
 ETA\_TO\_OPERATION\_INCIDENCE\_UNLOADED}.}
 \label{eq:v23-status-code}
\end{equation}
This is strictly more informative than the V22 phrase ``physical reward not
yet landed'': the reward and path carrier exist on the explicit RPESM
regulator, but the particular microscopic residual parameter chosen for P4 has
not yet been made incident with the accepted operation.

\subsection{V23 anti-circularity audit}
\label{sec:v23-audit}

\begin{enumerate}[leftmargin=2.4em]
\item \textbf{Finite packet population is not promoted to physical incidence.}
The RPESM construction supplies action writers and a path law, but the common
action matrix must still match the transfer Dirichlet matrix.
\item \textbf{Temporal closure is not promoted to microscopic selection.}
The pure-refresh branch closes every three-cylinder head while carrying no
residual race operation derivative.
\item \textbf{The refresh clock is calibrated rather than fitted into the
action.}  The independently reconstructed action matrix is held fixed and the
scalar clock must pass the ray residual test.
\item \textbf{De-refreshing is exact.}  The auxiliary refresh rescales the
innovation by a positive scalar and cannot change its rank.
\item \textbf{Raw scheduler covariance is not called a new action coupling.}
The law--material short removes source ancestry before a material selector is
recognized.
\item \textbf{A nonzero residual susceptibility requires an actual
\(K_\alpha\).}  Missing circulation-to-operation incidence is not represented
by a guessed differential of the accepted generator.
\item \textbf{Marked temporal closure is tested separately.}  Even after the
base action head reduces the transfer, a residual generator mark can leave that
head and must return its own polar missing writer.
\end{enumerate}

\section{V23 status ledger}
\label{sec:v23-status}

\subsection*{Proved in V23}
\begin{enumerate}[leftmargin=2.2em]
\item The refresh-only member of the explicit RPESM slab has
\(M_n^\Gamma=e^{-n\rho\tau}M_0^\Gamma\) for every bounded physical writer
bank.  Its action-response head is exactly reducing and its three-cylinder
innovation vanishes.
\item Common-action incidence remains independent of that temporal closure.
For pure refresh the transfer action is \((1-e^{-\rho\tau})C_F\), so one
scalar clock matches an independently reconstructed action matrix exactly if
and only if that matrix lies on the positive covariance ray.  The optimal
scalar mismatch and clock are explicit.
\item For the general RPESM decomposition
\(L=\rho(R-I)+L^{\rm loc}\), exact de-refreshing gives
\(M_n=e^{-n\rho\tau}\widehat M_n\) and rescales the three-cylinder innovation
by \(e^{-2\rho\tau}\); its rank is invariant.
\item On the canonical P4 residual deformation, varying only the V20 fast race
law while holding the calibrated field cylinder and accepted-operation
coefficients fixed gives \(K_\alpha=0\).  Hence the V22 residual-circulation
path susceptibility is exactly zero on that populated branch.
\item The inherited law--material decomposition shows why decoded writers do
not evade this result: if they factor only through residual-invariant P1
coordinates, their post-short material incidence vanishes.
\item The first genuine P4 microscopic dynamical datum is therefore the
three-column accepted-operation derivative
\(K_\alpha=\partial_{\eta^\alpha}L_{\rm acc}\).
\item Given those marks, V23 constructs the positive marked innovation
\(\mathbb I_{\Gamma,K}\).  Its zero branch is exactly invariance of the closed
action-response head under all three residual operation marks; its positive
branch returns the source-minimal polar marked writer which must be appended.
\end{enumerate}

\subsection*{Inherited}
\begin{enumerate}[leftmargin=2.2em]
\item V16--V22: physical calibration/RG covariance, the no-free-Wilson theorem,
the microscopic race calibration and three residual circulation coordinates,
the first-jump action-incidence compiler, the scalar-action coboundary no-go,
and the quadratic reflected path susceptibility.
\item The Grand Tensor explicit RPESM cylinder and reflected slab, complete
same-history finite packet population, accepted/common-action occurrence
compiler, action-response three-cylinder theorem, calibrated Fisher--Green
bridge and law--material derivative decomposition.
\end{enumerate}

\subsection*{Conditional}
\begin{enumerate}[leftmargin=2.2em]
\item The pure-refresh and de-refreshing formulas are exact for the explicit
RPESM generator decomposition.  They are a control branch, not a claim that
the physical accepted dynamics is pure refresh.
\item The canonical \(K_\alpha=0\) theorem holds when the residual race law is
varied while the accepted field-operation coefficients are held fixed, exactly
as in the present population construction.  A future same-record coupling can
make \(K_\alpha\ne0\); that coupling is the desired new physical input.
\item The V22 susceptibility still requires the stationary-law derivative when
a loaded \(K_\alpha\) changes the stationary measure; V23's marked closure is
independent of that additional Perron/Kato term.
\item All fixed finite loop/EFT/source, positive-reference-scale, determinant
line and real-Pfaffian qualifications of V16--V22 remain in force.
\end{enumerate}

\subsection*{Open beyond V23}
\begin{enumerate}[leftmargin=2.2em]
\item Populate the same-record circulation-to-operation incidence
\(K_\alpha=\partial_{\eta^\alpha}L_{\rm acc}\) for one physical accepted
Einstein--SM update.  The most direct implementation is a joint score/operation
packet in which the residual root circulation changes the actual accepted
field update rather than only the fast scheduler record.
\item Evaluate the marked Gram \(\mathbb I_{\Gamma,K}\).  If nonzero, append
its polar marked writer and repeat; if zero, compress the three operation marks
to the closed finite action head.
\item Calibrate the common-action/transfer incidence on that same head.  A
scalar refresh clock is sufficient only on the covariance-ray branch;
otherwise populate the required local operation coefficients rather than
altering the action matrix.
\item Once the marked head closes, evaluate the V22 path susceptibility (with
full stationary-law derivatives when needed), push it through the V18 physical
Wilson quotient, and compute the first actual P4 rank/left-nullspace.
\item The parallel P2/P3 programme remains available after a predictive
boundary relation is found: higher-loop beta rows, threshold matching,
standard-scheme conversion and numerical RG evolution.
\end{enumerate}

\subsection*{Exact next load-bearing theorem}
\begin{center}
\fbox{\parbox{0.92\textwidth}{
\textbf{P4g: same-record circulation-to-operation incidence.}
On one explicit accepted Einstein--Standard-Model update, couple the three V20
residual root-circulation coordinates to the physical update kernel and compute
\(K_\alpha=\partial_{\eta^\alpha}L_{\rm acc}\) (or the discrete-transfer
analogue) from the same event record.  Short the already known scheduler-law
source before declaring a material operation derivative.  Then evaluate the
marked action-head Gram \(\mathbb I_{\Gamma,K}\): append its polar range if
positive, or compress the three marks if zero.  Only on the resulting closed
marked head should the V22 quadratic action-response susceptibility be pushed
through the physical Wilson quotient and tested for a nontrivial predictive
left nullspace.}}
\end{center}

\section*{Append-only continuation marker: V23}
\addcontentsline{toc}{section}{Append-only continuation marker: V23}
\textbf{End of V23.}  The complete V1--V22 body above is retained verbatim.
V23 goes upstream from the V22 ``physical writer landing'' diagnosis and
separates writer population from two stronger incidence requirements.  The
explicit RPESM regulator already carries a same-history physical writer bank
and reflected path law.  On its pure-refresh control branch every
three-cylinder action-response head is exactly reducing, while the
independently reconstructed common action passes a separate scalar clock/ray
calibration test.  Exact de-refreshing shows that the auxiliary reset cannot
change temporal-memory rank.  More importantly, the canonical residual race
deformation currently changes only the fast scheduler law and leaves the
accepted full-field generator fixed, so the V22 circulation susceptibility is
exactly zero on that branch.  The first missing phenomenological datum is
therefore not another writer but the same-record derivative of the accepted
operation with respect to residual circulation.  V23 packages those three
marks in a positive marked-closure Gram which either closes on the existing
action head or returns the exact polar marked writer to append.




\clearpage
\section{V24 continuation: terminal operation germs, commutator inversion, and microscopic coefficient incidence}
\label{sec:v24-continuation}

V23 isolated the first missing phenomenological datum as a three-column
circulation-to-operation derivative.  That statement must now be split more
carefully.  The explicit RPESM construction contains a reflected Markov
operation on the commutative field algebra, whereas the loaded Euclidean
operation packet also carries a represented noncommutative algebra with a
static Dirac/coefficient operator.  The Grand Tensor terminal-germ theorem
proves that the former does not determine the latter: a self-adjoint reflected
Markov generator supplies symmetric Euclidean time translation, but not the
antisymmetric same-carrier derivation
\begin{equation}
 \mathcal L^{\rm germ}=-i[D^{\rm germ},\,\cdot\,]
 \label{eq:v24-germ}
\end{equation}
on the represented operation algebra.

This distinction is load bearing for P4.  On the canonical residual-race
branch V23 proved that the accepted \emph{field} generator is held fixed when
\(\eta\) is varied.  Thus its derivative is exactly zero.  The terminal
operation germ, by contrast, has not yet been acquired as a same-history
channel derivative even at the reference point.  Its \(\eta\)-derivative is
therefore not a zero matrix: it is an unpopulated row.  The aim of V24 is to
turn that missing row into a finite compiler which, once the differentiated
CPTP panel is supplied, returns either a precise obstruction or the desired
microscopic coefficient sensitivity.

The new chain is
\begin{equation}
 \boxed{
 \eta
 \longrightarrow
 \partial_\eta\mathcal L^{\rm germ}
 \longrightarrow
 \partial_\eta D^{\rm germ}\!\!\pmod{\text{commutant}}
 \longrightarrow
 \partial_\eta c^{\rm phys}
 \longrightarrow
 \text{P4 matching rank}.}
 \label{eq:v24-chain}
\end{equation}
No arrow in \eqref{eq:v24-chain} is replaced by a static coefficient, a
symmetric Markov generator, or a fitted Wilson boundary value.

\subsection{Two distinct residual operation derivatives}
\label{sec:v24-two-derivatives}

Let \(\eta=(\eta^1,\eta^2,\eta^3)\) denote the V19--V20 residual circulation
coordinates.  We write
\begin{align}
 K_\alpha^{\rm fld}
 &:=\left.\partial_{\eta^\alpha}L_{\rm acc,\eta}\right|_{\eta=0},
 \label{eq:v24-K-field}\\
 \mathcal K_\alpha^{\rm op}
 &:=\left.\partial_{\eta^\alpha}
     \mathcal L_{\eta}^{\rm germ}\right|_{\eta=0}.
 \label{eq:v24-K-op}
\end{align}
The first acts on the accepted reflected field algebra.  The second is a
linear map on the represented operation core and, on the desired inner branch,
is a commutator derivation.

\begin{theorem}[Current population status of the two residual derivatives]
\label{thm:v24-two-derivative-status}
On the presently populated canonical RPESM branch:
\begin{enumerate}[label=\textup{(D\arabic*)},leftmargin=3em]
\item \(K_\alpha^{\rm fld}=0\) for \(\alpha=1,2,3\);
\item the reference operation germ \(\mathcal L_0^{\rm germ}\) is not yet
identified with the static coefficient operator by a same-carrier
differentiated channel panel;
\item consequently \(\mathcal K_\alpha^{\rm op}\) is presently
\emph{unpopulated}, not proved to vanish.
\end{enumerate}
In particular, the zero in item \textup{(D1)} cannot be substituted for the
missing operation-germ columns of item \textup{(D3)}.
\end{theorem}

\begin{proof}
Item \textup{(D1)} is the V23 canonical-branch theorem: the residual race law
is varied while the accepted field-operation coefficients are held fixed.
Items \textup{(D2)--(D3)} are the inherited terminal differential-incidence
boundary.  The loaded packet contains the static coefficient operator and the
self-adjoint reflected field generator, but the Grand Tensor terminal audit
returns a positive rank-one witness until a same-carrier forward/reflected
CPTP germ and its mixed incidence rows are supplied.  A missing derivative row
cannot be identified with the zero derivative of a different operator.
\end{proof}

\begin{firewall}[Symmetric Euclidean time is not the antisymmetric operation germ]
The field generator \(L_{\rm acc}\), the OS Hamiltonian, and
\(D^{\rm germ}\) are distinct objects.  Even exact knowledge of the complete
field Markov semigroup does not imply
\(\mathcal L^{\rm germ}=-i[D^{\rm coeff},\cdot]\) on the represented
noncommutative operation algebra.  P4g therefore has to acquire a differentiated
same-carrier operation panel rather than reinterpret the field clock.
\end{firewall}

\subsection{Finite commutator synthesis on a represented operation core}
\label{sec:v24-commutator-synthesis}

Fix one cutoff and one finite represented operation core
\begin{equation}
 \mathcal A_{\rm op}^{\rm fin}
 =\operatorname{alg}\{a_1,\ldots,a_m\}
 \subseteq \mathcal B(\mathcal H_{\rm op}),
 \label{eq:v24-op-core}
\end{equation}
where the chosen generators include every operation needed by the finite P4
panel.  Let \(\mathfrak h\) be the real Hilbert space of allowed Hermitian
coefficient variations on \(\mathcal H_{\rm op}\), with Hilbert--Schmidt inner
product, after imposing the fixed grading, representation and reality linear
constraints.  Define the compatible commutant
\begin{equation}
 \mathfrak c
 :=\{H\in\mathfrak h:[H,\pi(a_j)]=0\ \text{for all }j\},
 \qquad
 \mathfrak h_\perp:=\mathfrak c^\perp.
 \label{eq:v24-commutant}
\end{equation}
Put
\begin{equation}
 \mathcal Y
 :=\bigoplus_{j=1}^m \mathcal B_2(\mathcal H_{\rm op}),
 \qquad
 \mathcal C(H)
 :=\bigl(-i[H,\pi(a_j)]\bigr)_{j=1}^m.
 \label{eq:v24-C-map}
\end{equation}
Here \(\mathcal B_2\) denotes the finite Hilbert--Schmidt space.

\begin{lemma}[Kernel of the finite commutator synthesis]
\label{lem:v24-C-kernel}
The synthesis map \(\mathcal C:\mathfrak h\to\mathcal Y\) satisfies
\begin{equation}
 \boxed{\ker\mathcal C=\mathfrak c.}
 \label{eq:v24-C-kernel}
\end{equation}
Consequently its restriction
\(\mathcal C_\perp:=\mathcal C|_{\mathfrak h_\perp}\) is injective and the
finite Gram operator
\begin{equation}
 G_C:=\mathcal C_\perp^*\mathcal C_\perp
 \label{eq:v24-GC}
\end{equation}
is strictly positive on \(\mathfrak h_\perp\).
\end{lemma}

\begin{proof}
By definition, \(\mathcal C(H)=0\) exactly when \(H\) commutes with every
selected generator.  Since the generators produce the declared finite
operation core, this is precisely \(H\in\mathfrak c\).  Injectivity on the
orthogonal complement is immediate.  In finite dimension, an injective linear
map has positive-definite normal Gram.
\end{proof}

Define the commutator reconstruction operator and its range projector by
\begin{equation}
 \mathcal C^\sharp
 :=G_C^{-1}\mathcal C_\perp^*,
 \qquad
 P_C:=\mathcal C_\perp\mathcal C^\sharp.
 \label{eq:v24-C-sharp}
\end{equation}

\begin{theorem}[Exact inner-derivation incidence compiler]
\label{thm:v24-inner-compiler}
For any measured finite operation derivative vector \(Y\in\mathcal Y\), put
\begin{equation}
 H_Y:=\mathcal C^\sharp Y,
 \qquad
 R_Y:=(I-P_C)Y.
 \label{eq:v24-HY-RY}
\end{equation}
Then:
\begin{enumerate}[label=\textup{(I\arabic*)},leftmargin=3em]
\item \(H_Y\in\mathfrak h_\perp\) is the unique minimum-norm Hermitian
coefficient whose commutator image is the orthogonal projection of \(Y\) onto
\(\operatorname{Ran}\mathcal C\);
\item
\begin{equation}
 \boxed{R_Y=0
 \quad\Longleftrightarrow\quad
 Y=-i[H_Y,\cdot]\ \text{on the selected core};}
 \label{eq:v24-inner-zero}
\end{equation}
\item if \(R_Y\ne0\), its positive Gram
\begin{equation}
 \mathfrak r_Y:=R_Y^*R_Y
 \label{eq:v24-inner-residual-Gram}
\end{equation}
returns an explicit operation-germ component not represented by the current
static coefficient/commutator bank.
\end{enumerate}
The ambiguity of the reconstructed Hermitian coefficient is exactly the
compatible commutant \(\mathfrak c\), already quotiented by the choice
\(H_Y\in\mathfrak h_\perp\).
\end{theorem}

\begin{proof}
The formula \eqref{eq:v24-C-sharp} is the ordinary left Moore--Penrose inverse
of the injective finite map \(\mathcal C_\perp\).  Hence
\(P_C\) is the orthogonal projector onto its range and
\(Y=P_CY+(I-P_C)Y\) is the unique orthogonal decomposition.  The first two
claims follow.  The third is simply the squared norm on the unexplained
orthogonal component; it vanishes exactly on the represented commutator range.
\end{proof}

\subsection{Differentiating a same-history forward/reflected channel family}
\label{sec:v24-channel-family}

Let \(\Phi_{\eta,+}(t)\) and \(\Phi_{\eta,-}(t)\) be same-carrier CPTP curves
on the represented operation algebra, measured with one common duration and
source gauge, and satisfying
\begin{equation}
 \Phi_{\eta,\pm}(0)=\operatorname{id}.
 \label{eq:v24-channel-id}
\end{equation}
Define the antisymmetric difference quotient
\begin{equation}
 A_\eta(t)
 :=\frac{\Phi_{\eta,+}(t)-\Phi_{\eta,-}(t)}{2t}.
 \label{eq:v24-Aeta}
\end{equation}

\begin{theorem}[Residual differentiation of the terminal germ]
\label{thm:v24-germ-diff}
Assume there is a neighborhood of \((\eta,t)=(0,0)\) on which
\(A_\eta(t)\) is continuously differentiable in \(\eta\), and for each
\(\alpha\)
\begin{equation}
 \left\|
 \partial_{\eta^\alpha}A_\eta(t)
 -Y_\alpha(\eta)
 \right\|_{\mathcal Y}
 \le Ct^2
 \label{eq:v24-uniform-germ-rate}
\end{equation}
uniformly on that neighborhood.  Then
\begin{equation}
 \mathcal L_\eta^{\rm germ}
 :=\lim_{t\downarrow0}A_\eta(t)
 \label{eq:v24-Lgerm-limit}
\end{equation}
exists in the selected finite norm, is differentiable in \(\eta\), and
\begin{equation}
 \boxed{
 \partial_{\eta^\alpha}\mathcal L_\eta^{\rm germ}
 =Y_\alpha(\eta).}
 \label{eq:v24-germ-derivative}
\end{equation}
If in addition
\begin{equation}
 \mathcal L_\eta^{\rm germ}
 =-i[D_\eta^{\rm germ},\cdot]
 \label{eq:v24-inner-family}
\end{equation}
with \(D_\eta^{\rm germ}\) differentiable modulo the compatible commutant,
then
\begin{equation}
 \boxed{
 Y_\alpha(0)
 =-i[\dot D_\alpha^{\rm germ},\cdot],
 \qquad
 \dot D_\alpha^{\rm germ}
 =\mathcal C^\sharp Y_\alpha(0)}
 \label{eq:v24-Ddot}
\end{equation}
provided the residual \(R_{Y_\alpha}\) of
Theorem~\ref{thm:v24-inner-compiler} vanishes.
\end{theorem}

\begin{proof}
The estimate \eqref{eq:v24-uniform-germ-rate} gives uniform convergence of the
\(\eta\)-derivatives as \(t\downarrow0\).  The standard finite-dimensional
interchange theorem for a uniformly convergent derivative family therefore
permits differentiation through the limit.  Differentiating
\eqref{eq:v24-inner-family} gives the commutator formula.  The minimum-norm
representative modulo the commutant is exactly the reconstruction of
Theorem~\ref{thm:v24-inner-compiler}.
\end{proof}

\subsection{A nested-time finite acquisition formula}
\label{sec:v24-nested-time}

The inherited terminal handoff asks for forward and reflected rows at two
nested physical times plus a smaller-time held-out check.  The next lemma
explains why that packet is sufficient for a high-order germ estimator on a
reversal-symmetric branch.

\begin{lemma}[Two-time Richardson extraction]
\label{lem:v24-Richardson}
Assume the same-history reversal convention makes
\begin{equation}
 A_\eta(t)
 =\mathcal L_\eta^{\rm germ}
  +t^2B_\eta+t^4C_\eta(t),
 \qquad
 \sup\|C_\eta(t)\|\le C_4.
 \label{eq:v24-A-expansion}
\end{equation}
Define
\begin{equation}
 \widehat{\mathcal L}_{\eta,t}
 :=\frac{4A_\eta(t/2)-A_\eta(t)}{3}.
 \label{eq:v24-Richardson}
\end{equation}
Then
\begin{equation}
 \boxed{
 \|\widehat{\mathcal L}_{\eta,t}
   -\mathcal L_\eta^{\rm germ}\|
 \le C_Rt^4.}
 \label{eq:v24-Richardson-error}
\end{equation}
If \(\mathcal L_\eta^{\rm germ}\) is three times differentiable in \(\eta\)
and its third derivatives are bounded, then the symmetric residual estimator
\begin{equation}
 \widehat{\mathcal K}_{\alpha;t,\epsilon}
 :=\frac{
 \widehat{\mathcal L}_{\epsilon e_\alpha,t}
 -\widehat{\mathcal L}_{-\epsilon e_\alpha,t}
 }{2\epsilon}
 \label{eq:v24-K-estimator}
\end{equation}
satisfies
\begin{equation}
 \boxed{
 \|\widehat{\mathcal K}_{\alpha;t,\epsilon}
 -\mathcal K_\alpha^{\rm op}\|
 \le C_1t^4+C_2\epsilon^2.}
 \label{eq:v24-K-estimator-error}
\end{equation}
\end{lemma}

\begin{proof}
Insert \eqref{eq:v24-A-expansion} into
\eqref{eq:v24-Richardson}.  The \(t^2B_\eta\) terms cancel exactly and the
remainder is bounded by a constant times \(t^4\).  The second statement is the
ordinary centered finite-difference remainder in \(\eta\), combined with the
uniform Richardson error.  Since the panel is finite, the constants may be
taken uniformly over the selected operation bank.
\end{proof}

Thus one concrete raw acquisition panel for the three residual directions is
\begin{equation}
 \eta\in\{0,\pm\epsilon e_1,\pm\epsilon e_2,\pm\epsilon e_3\},
 \qquad
 t\in\{\tau,\tau/2\},
 \qquad
 \pm\ \text{forward/reflected branch},
 \label{eq:v24-raw-panel}
\end{equation}
namely twenty-eight channel rows plus the common identity row, duration
calibration and the inherited mixed incidence Grams.  A smaller held-out time
can then test the \(t^4\) Richardson law without changing the fitted germ.

\subsection{Terminal incidence as a microscopic matching law}
\label{sec:v24-terminal-matching}

The reconstructed operation germ becomes phenomenologically useful only after
it is matched to the static physical coefficient operator rather than treated
as an unrelated Hamiltonian.  Let
\begin{equation}
 c=(c^1,\ldots,c^d)
 \label{eq:v24-c}
\end{equation}
be a finite physical coefficient block at \(\mu_0\), after the V16 calibration
and V18 quotient by redundant directions.  Let
\begin{equation}
 D^{\rm coeff}(c)
 \label{eq:v24-Dcoeff}
\end{equation}
be its represented Hermitian synthesis on the same operation carrier.  At the
reference point \(c_0\), define
\begin{equation}
 \mathcal J_D:\mathbb R^d\longrightarrow\mathfrak h_\perp,
 \qquad
 \mathcal J_D\,\delta c
 :=P_{\mathfrak h_\perp}
   \left.
   \frac{d}{ds}D^{\rm coeff}(c_0+s\delta c)
   \right|_{s=0}.
 \label{eq:v24-JD}
\end{equation}

\begin{definition}[Static coefficient-incidence residual]
\label{def:v24-coeff-residual}
Assume \(\mathcal J_D\) has full column rank.  Put
\begin{equation}
 \mathcal J_D^\sharp
 :=(\mathcal J_D^*\mathcal J_D)^{-1}\mathcal J_D^*,
 \qquad
 P_D:=\mathcal J_D\mathcal J_D^\sharp.
 \label{eq:v24-JD-sharp}
\end{equation}
For the three reconstructed germ columns
\(\dot D_\alpha^{\rm germ}\), define
\begin{align}
 \mathsf M_{D,\alpha}
 &:=\mathcal J_D^\sharp\dot D_\alpha^{\rm germ}
 \in\mathbb R^d,
 \label{eq:v24-MD}\\
 \mathcal R_{D,\alpha}
 &:=(I-P_D)\dot D_\alpha^{\rm germ}.
 \label{eq:v24-RD}
\end{align}
The matrix \(\mathsf M_D\in\mathbb R^{d\times3}\) has these vectors as its
columns.
\end{definition}

\begin{theorem}[Differentiated terminal-incidence matching law]
\label{thm:v24-terminal-matching}
Suppose a same-history family satisfies the terminal incidence
\begin{equation}
 D_\eta^{\rm germ}
 =D^{\rm coeff}(c(\eta))+C_\eta,
 \qquad
 C_\eta\in\mathfrak c,
 \label{eq:v24-terminal-incidence}
\end{equation}
with \(c(0)=c_0\).  Then necessarily
\begin{equation}
 \boxed{
 \mathcal R_{D,\alpha}=0,
 \qquad
 \partial_{\eta^\alpha}c(0)
 =\mathsf M_{D,\alpha}.}
 \label{eq:v24-terminal-match-result}
\end{equation}
Conversely, if every \(\mathcal R_{D,\alpha}\) vanishes, then
\(\mathsf M_D\) is the unique first-order physical coefficient response which
matches the measured germ derivative modulo the compatible commutant.

Therefore the missing V19 selector derivative for this terminal coefficient
block is no longer abstract: once the same-history germ rows are acquired,
\begin{equation}
 \boxed{
 D\mathfrak M_{\rm terminal}(0)=\mathsf M_D.}
 \label{eq:v24-terminal-selector}
\end{equation}
\end{theorem}

\begin{proof}
Differentiate \eqref{eq:v24-terminal-incidence} at \(\eta=0\) and project to
\(\mathfrak h_\perp\).  The commutant derivative drops out and one obtains
\[
 \dot D_\alpha^{\rm germ}
 =\mathcal J_D\,\partial_{\eta^\alpha}c(0).
\]
Full column rank of \(\mathcal J_D\) gives the unique solution by
\(\mathcal J_D^\sharp\), while the orthogonal complement must vanish.  The
converse is the same finite linear-algebra statement read backwards.
\end{proof}

\begin{corollary}[Conditional predictivity rank from a complete three-parameter selector]
\label{cor:v24-rank-bound}
If the residual circulation packet \(\eta\in\mathbb R^3\) is proved to be the
complete microscopic selector for this terminal-matched coefficient block and
no coefficient in \(c\) is independently refitted afterward, then
\begin{equation}
 \boxed{
 \operatorname{rank}\mathsf M_D\le3,
 \qquad
 \dim\ker(\mathsf M_D^T)\ge d-3.}
 \label{eq:v24-rank-bound}
\end{equation}
Thus any block with \(d>3\) has a nontrivial first-order predictive conormal.
This conclusion is conditional on completeness of the microscopic selector;
it is not inferred merely from the existence of three presently named
coordinates.
\end{corollary}

\begin{proof}
The first inequality is the rank bound for a \(d\times3\) matrix.  Rank-nullity
for its transpose gives the second.  If additional microscopic variables act
on the block, their columns must be appended before the left nullspace is
interpreted as a prediction.
\end{proof}

\subsection{The mixed operation/action incidence remains a separate marked gate}
\label{sec:v24-operation-action}

The terminal germ lives on the represented operation algebra, whereas the V23
marked Gram was formulated on a selected action-response carrier.  A physical
comparison therefore requires the mixed incidence rows demanded by the
terminal handoff.  Let
\begin{equation}
 \mathcal T_{\Gamma}:
 \mathcal Y\longrightarrow
 \operatorname{End}(\mathcal H_\Gamma)
 \label{eq:v24-TGamma}
\end{equation}
be the finite operation-to-action incidence map reconstructed from that mixed
Gram once it is populated.  Define
\begin{equation}
 K_{\Gamma,\alpha}
 :=\mathcal T_\Gamma(\mathcal K_\alpha^{\rm op}).
 \label{eq:v24-KGamma}
\end{equation}
Then the V23 marked action-head innovation becomes
\begin{equation}
 \boxed{
 \mathbb I_{\Gamma,K}^{\rm op}
 =\left[
 J_\Gamma^*K_{\Gamma,\alpha}^*
 (I-\Pi_\Gamma)
 K_{\Gamma,\beta}J_\Gamma
 \right]_{\alpha,\beta=1}^3.}
 \label{eq:v24-marked-Gram}
\end{equation}

\begin{proposition}[No operation-to-action substitution without the mixed Gram]
\label{prop:v24-no-substitution}
Before \(\mathcal T_\Gamma\) is populated from same-history mixed incidence,
neither the static coefficient operator nor the symmetric field generator
licenses setting
\(K_{\Gamma,\alpha}=\mathcal K_\alpha^{\rm op}\) or
\(K_{\Gamma,\alpha}=K_\alpha^{\rm fld}\).  After the mixed incidence is
populated, \eqref{eq:v24-marked-Gram} is an ordinary finite positive Gram and
can be routed by the V23 marked-closure theorem.
\end{proposition}

\begin{proof}
The three maps act on different carriers.  Equality would require precisely
the missing same-history incidence identifying their source ranges.  Once the
mixed Gram supplies \(\mathcal T_\Gamma\), the displayed formula is a finite
composition on one carrier and the inherited marked-Gram theorem applies.
\end{proof}

\subsection{A terminating finite P4g compiler}
\label{sec:v24-router}

The previous results can be organized as a finite decision procedure.

\begin{theorem}[Terminal-germ P4g router]
\label{thm:v24-router}
For one finite physical operation/coefficient/action panel, perform the
following ordered tests.
\begin{enumerate}[label=\textup{(R\arabic*)},leftmargin=3em]
\item Acquire the common-identity forward/reflected channel panel
\eqref{eq:v24-raw-panel}, its duration calibration and held-out smaller-time
row.  Form \(\widehat{\mathcal L}_{0,t}\) and the three
\(\widehat{\mathcal K}_{\alpha;t,\epsilon}\).
\item Apply Theorem~\ref{thm:v24-inner-compiler} to the reference germ and the
three residual germ columns.  A nonzero inner residual is returned as the first
operation-germ obstruction.
\item Match the reconstructed Hermitian germ to
\(D^{\rm coeff}(c)\) modulo the compatible commutant.  A nonzero
\(\mathcal R_{D,\alpha}\) is returned as the first static coefficient-incidence
obstruction; otherwise output \(\mathsf M_D\).
\item Populate the mixed operation/action incidence map \(\mathcal T_\Gamma\)
and evaluate \(\mathbb I_{\Gamma,K}^{\rm op}\).  If positive, append the polar
marked writer and rerun the action-head closure; if zero, compress the marks to
the closed action head.
\item Only after the preceding rows close, push \(\mathsf M_D\) and the V22
path susceptibility through the V18 physical quotient and compute the actual
P4 matching rank/left nullspace.
\end{enumerate}
Every failure returns a typed finite residual.  No downstream step is licensed
by replacing an unpopulated row with zero.
\end{theorem}

\begin{proof}
Items \textup{(R1)--(R2)} are Lemma~\ref{lem:v24-Richardson} and
Theorem~\ref{thm:v24-inner-compiler}.  Item \textup{(R3)} is
Theorem~\ref{thm:v24-terminal-matching}.  Item \textup{(R4)} is
Proposition~\ref{prop:v24-no-substitution} followed by the V23 marked-closure
alternative.  Item \textup{(R5)} is the V18 definition of physical
predictivity applied only after the actual microscopic matching derivative has
been populated.  The ordering prevents a static coefficient, field clock or
missing mixed incidence from being reused as a substitute for the row it is
supposed to test.
\end{proof}

The current branch stops at \textup{(R1)}: the same-history differentiated
operation channel family has not yet been acquired.  The existing Grand Tensor
terminal audit already returns a positive rank-one witness at the reference
point, so this stop is independently certified and not merely an absence of
notation.

\subsection{V24 anti-circularity audit}
\label{sec:v24-audit}

\begin{enumerate}[leftmargin=2.4em]
\item \textbf{The proved field-generator zero is not reused as an operation-germ zero.}
The two derivatives act on distinct algebras and have different symmetry.
\item \textbf{A static Dirac/coefficient operator is not called a measured
generator.}  Terminal incidence is tested only after the forward/reflected
channel germ has landed.
\item \textbf{Commutant ambiguity is removed explicitly.}  The Hermitian germ
is reconstructed in \(\mathfrak c^\perp\); the residual tests innerness rather
than choosing an arbitrary anchor.
\item \textbf{The nested-time estimator has a declared error law.}  The
\(O(t^4+\epsilon^2)\) rate is conditional on the stated reversal-symmetric
Taylor hypotheses and is held out at a smaller time.
\item \textbf{Coefficient matching is separated from operation reconstruction.}
A measured inner germ can still fail to lie in the tangent range of the chosen
physical coefficient block; \(\mathcal R_D\) records that failure.
\item \textbf{A three-column rank bound is not automatically a prediction.}
It becomes predictive only if \(\eta\) is proved complete for the selected
microscopic block and no Wilson coordinate is independently refitted.
\item \textbf{Operation/action incidence is not assumed.}  The mixed Gram must
land before an operation derivation is promoted to a marked action-head
generator.
\end{enumerate}

\section{V24 status ledger}
\label{sec:v24-status}

\subsection*{Proved in V24}
\begin{enumerate}[leftmargin=2.2em]
\item P4 residual operation incidence splits into a symmetric accepted-field
derivative \(K_\alpha^{\rm fld}\) and an antisymmetric operation-germ derivative
\(\mathcal K_\alpha^{\rm op}\).  The former is zero on the current canonical
branch; the latter is unpopulated rather than zero.
\item On a finite represented operation core, the commutator synthesis map has
kernel exactly the compatible commutant.  Its restriction to the orthogonal
complement is injective and has an explicit positive normal Gram.
\item Every measured operation derivative therefore has a unique minimum-norm
Hermitian commutator reconstruction plus an exact orthogonal residual.  The
residual vanishes if and only if the measured derivative is an inner derivation
on the selected core.
\item Under a uniform differentiated small-time estimate, the residual
circulation derivative commutes with the terminal-germ limit.  On an inner
family it reconstructs \(\partial_\eta D^{\rm germ}\) modulo the commutant.
\item A two-nested-time Richardson estimator combined with symmetric residual
perturbations has error \(O(t^4+\epsilon^2)\) under the declared reversal-symmetric
Taylor hypotheses.  The corresponding raw three-direction panel has twenty-eight
forward/reflected channel rows plus one common identity row and calibration data.
\item Differentiated terminal incidence with a static physical coefficient
operator yields an explicit microscopic matching matrix
\(\mathsf M_D\).  A separate residual decides whether the measured germ tangent
lies in the physical coefficient tangent range.
\item If the three residual coordinates are independently proved to be the
complete microscopic selector for a \(d\)-dimensional terminal coefficient
block, then the first-order predictive conormal has dimension at least
\(d-3\).
\item The operation-to-action mixed incidence remains an independent finite
gate.  Once populated, its marked Gram is routed by the V23 closure theorem.
\end{enumerate}

\subsection*{Inherited}
\begin{enumerate}[leftmargin=2.2em]
\item V16--V23: physical calibration/RG covariance, the no-free-Wilson theorem,
three residual race coordinates, first-jump and path selector compilers, the
coboundary and reversal-parity no-go theorems, and the distinction between
writer population, common-action incidence and circulation-to-operation
incidence.
\item The Grand Tensor loaded Euclidean/OS packet, terminal same-history
differential demand, the theorem that a static coefficient and symmetric
reflected Markov generator do not determine the operation germ, and the
positive rank-one first terminal witness with its explicit germ handoff panel.
\end{enumerate}

\subsection*{Conditional}
\begin{enumerate}[leftmargin=2.2em]
\item The Richardson order \(O(t^4)\) assumes the reversal-symmetric expansion
\eqref{eq:v24-A-expansion}; without it, the raw channel rows still define the
germ in the small-time limit but require a lower-order extrapolation theorem.
\item Full column rank of \(\mathcal J_D\) is a calibration hypothesis for the
selected physical coefficient block.  If it fails, the block itself has
redundant/unidentified coordinates which must be quotiented before P4.
\item The rank promise \eqref{eq:v24-rank-bound} requires proof that no
additional microscopic selector coordinates enter the block.
\item Operation-to-action incidence still has to be physically populated on
the same history before the V23 marked Gram is evaluable.
\item All finite-loop/EFT/source, positive-reference-scale, determinant-line
and real-Pfaffian qualifications from V16--V23 remain in force.
\end{enumerate}

\subsection*{Open beyond V24}
\begin{enumerate}[leftmargin=2.2em]
\item Acquire the same-carrier forward/reflected CPTP operation panel jointly at
\(\eta=0\) and the six residual perturbations, with two nested physical times,
common duration calibration, a held-out smaller-time row, and the complete
mixed incidence Gram demanded by the terminal-germ handoff.
\item Evaluate the inner-derivation residuals and recover the three
\(\partial_{\eta^\alpha}D^{\rm germ}\) columns modulo the commutant.
\item Match those columns to a declared finite physical coefficient block and
compute \(\mathsf M_D\) and its coefficient-incidence residual.
\item Populate the operation/action mixed incidence, evaluate the V23 marked
Gram, close the marked action head, and only then compute the V22 path
susceptibility and V18 predictive left nullspace.
\item The parallel P2/P3 tasks remain available after the first nontrivial
matching relation lands: threshold matching, higher-loop RG rows, standard
scheme conversion and numerical trajectory evolution.
\end{enumerate}

\subsection*{Exact next load-bearing theorem}
\begin{center}
\fbox{\parbox{0.92\textwidth}{
\textbf{P4h: joint residual terminal-germ acquisition and coefficient matching.}
Populate one same-history represented-operation panel at
\(\eta\in\{0,\pm\epsilon e_1,\pm\epsilon e_2,\pm\epsilon e_3\}\) and two nested
physical durations, with forward/reflected branches, common source gauge,
duration calibration and the complete mixed incidence Gram.  Prove the
small-time/reversal extrapolation on that native panel, test the inner-derivation
residual, reconstruct the three Hermitian germ derivatives modulo the compatible
commutant, and match them to a declared physical coefficient block through
\(\mathcal J_D\).  The output is either a typed residual or the first actually
populated microscopic coefficient-sensitivity matrix \(\mathsf M_D\); only the
latter may be used for a P4 predictivity rank calculation.}}
\end{center}

\section*{Append-only continuation marker: V24}
\addcontentsline{toc}{section}{Append-only continuation marker: V24}
\textbf{End of V24.}  The complete V1--V23 body above is retained verbatim.
V24 goes upstream from the V23 circulation-to-operation handoff and separates
the proved zero derivative of the symmetric accepted field generator from the
still-unpopulated antisymmetric operation-algebra germ.  It constructs an exact
finite commutator-incidence compiler which reconstructs a measured inner germ
modulo the compatible commutant or returns its orthogonal residual, gives a
nested-time and residual-parameter acquisition scheme for the three microscopic
columns, and differentiates the terminal static-coefficient incidence to obtain
the first explicit formula for a microscopic physical coefficient-sensitivity
matrix.  The remaining gate is empirical/native population of the same-history
differentiated CPTP panel and its mixed operation/action incidence; no static
coefficient or symmetric field clock is substituted for those missing rows.



\clearpage
\section{V25 continuation: finite terminal-germ tomography, error budgets, and certified coefficient sensitivity}
\label{sec:v25-continuation}

V24 reduced the phenomenology selector problem to a finite same-history
terminal-germ acquisition.  It also identified the exact stopping rule: no
static coefficient, symmetric reflected Markov generator, determinant line, or
continuum argument may be substituted for the missing differentiated CPTP
packet.  The present continuation does not pretend that those channel rows have
already been measured.  Instead it proves that the declared acquisition is
finite, tomographically complete on the frozen operation core, quantitatively
stable under row error, and sufficient to produce a certified microscopic
coefficient-sensitivity matrix once the rows are populated.

This is the correct next mathematical task because P4 now has two logically
separate questions:
\begin{equation}
 \boxed{
 \begin{array}{c}
 \text{is the finite acquisition sufficient to identify the terminal germ?}
 \\
 \text{and, only after that, what numerical germ does the acquired packet carry?}
 \end{array}}
 \label{eq:v25-two-questions}
\end{equation}
V25 closes the first question at fixed finite operation/source panel.  The
second remains a population task and is reported as such in the ledger.

Throughout, let the represented source-minimal operation core be
$\mathcal A_{\rm op}\cong M_d(\mathbb C)$, let
$C:\mathbb C^d\to\mathcal H_F$ be the inherited source-fixing isometry, and
write
\begin{equation}
 \iota(X)=CXC^*,\qquad P=CC^*.
 \label{eq:v25-encoder}
\end{equation}
Let $(F_\mu)_{\mu=1}^{d^2}$ be a Hermitian Hilbert--Schmidt orthonormal basis of
$M_d(\mathbb C)$ and let $(G_A)_{A=1}^{D^2}$ be a Hermitian
Hilbert--Schmidt orthonormal basis of the finite full-cell matrix algebra
containing the landed carrier, chosen so that the first $d^2$ elements span the
encoded core.  All norms below are finite-dimensional Hilbert--Schmidt or the
induced operator norm unless stated otherwise.

\subsection{A scalar tomography realization of one channel row}
\label{sec:v25-tomography}

V24 counted one acquired CPTP map as one ``channel row''.  A physical
implementation of such a row must of course be finite scalar data.  The next
theorem makes that finite realization explicit.

\begin{definition}[Core-resolved channel tomography table]
\label{def:v25-channel-table}
For a linear Hermiticity-preserving map
$\Phi$ on the full-cell algebra define
\begin{equation}
 T_{A\mu}(\Phi)
 :=\left\langle G_A,\,\Phi(\iota(F_\mu))\right\rangle_{\rm HS},
 \qquad
 1\le A\le D^2,\quad 1\le\mu\le d^2.
 \label{eq:v25-channel-table}
\end{equation}
The subtable with $A\le d^2$ is the compressed core action; the rows
$A>d^2$ are the source-core leakage data required by the terminal audit.
\end{definition}

\begin{theorem}[Finite tomography of the channel restriction]
\label{thm:v25-channel-tomography}
The table $T(\Phi)$ uniquely determines the restricted superoperator
\begin{equation}
 \Phi\circ\iota:M_d(\mathbb C)\longrightarrow M_D(\mathbb C).
 \label{eq:v25-restricted-superoperator}
\end{equation}
More precisely,
\begin{equation}
 \boxed{
 \Phi(\iota(F_\mu))
 =\sum_{A=1}^{D^2}T_{A\mu}(\Phi)G_A.}
 \label{eq:v25-tomography-reconstruction}
\end{equation}
Hence the same table determines the corresponding core-input Choi block and
all finite residuals used in the V24 terminal-germ compiler.

If a measured table $\widehat T$ satisfies
\begin{equation}
 \sum_{A,\mu}
 |\widehat T_{A\mu}-T_{A\mu}|^2\le\delta_T^2,
 \label{eq:v25-table-error}
\end{equation}
then the reconstructed restricted superoperator obeys
\begin{equation}
 \boxed{
 \|\widehat\Phi\circ\iota-\Phi\circ\iota\|_{\rm HS\to HS}
 \le\delta_T.}
 \label{eq:v25-superoperator-error}
\end{equation}
\end{theorem}

\begin{proof}
Equation \eqref{eq:v25-tomography-reconstruction} is expansion in the complete
orthonormal output basis $(G_A)$.  Since $(F_\mu)$ is a basis of the input
core, these $d^2$ outputs determine the restriction linearly.  The Choi block
is a fixed linear reshuffling of the same matrix elements, so no further data
are needed.  For the error estimate, the Hilbert--Schmidt norm of the
superoperator matrix in the orthonormal input/output bases is exactly the
Euclidean norm of the coefficient table difference; the induced norm is bounded
by that Hilbert--Schmidt norm.
\end{proof}

\begin{corollary}[Finite scalar cost of the V24 channel panel]
\label{cor:v25-finite-cost}
For a frozen finite operation core, every one of the V24 channel rows is a
finite table of $D^2d^2$ real Hermitian-basis coefficients, or any equivalent
informationally complete process-tomography packet.  No infinite-dimensional
channel reconstruction is required for the coefficientwise P4 audit.
\end{corollary}

\subsection{The twenty-eight-row design and what it actually identifies}
\label{sec:v25-28-design}

Fix one reference time $t>0$ and one residual step $\epsilon>0$.  Put
\begin{equation}
 \mathcal E_\epsilon
 =\{0,\pm\epsilon e_1,\pm\epsilon e_2,\pm\epsilon e_3\}
 \subset\mathbb R^3,
 \qquad
 \mathcal T_t=\{t,t/2\}.
 \label{eq:v25-design-sets}
\end{equation}
For every $(\eta,\tau,\sigma)$ with
$\eta\in\mathcal E_\epsilon$, $\tau\in\mathcal T_t$ and
$\sigma\in\{+,-\}$, acquire the restricted CPTP map
$\Phi_{\eta,\sigma}(\tau)\circ\iota$.  There are
\begin{equation}
 \boxed{7\times2\times2=28}
 \label{eq:v25-28-count}
\end{equation}
channel rows.

\begin{proposition}[Role of the twenty-eight rows]
\label{prop:v25-design-role}
Under the V24 reversal-symmetric small-time and residual-parameter Taylor
hypotheses, the design \eqref{eq:v25-design-sets} determines:
\begin{enumerate}[label=\textup{(D\arabic*)},leftmargin=3em]
\item the reference antisymmetric terminal germ at $\eta=0$ through a two-time
Richardson cancellation of the leading $t^2$ bias;
\item the three first residual derivatives through centered differences in the
three coordinates, with the leading even-in-$\eta$ bias cancelled;
\item the source-core leakage and inner-Hamiltonian residual of each recovered
column because the full-cell output table is retained;
\item the static coefficient-incidence residual after applying the V24
coefficient synthesis map $\mathcal J_D$.
\end{enumerate}
The baseline point $\eta=0$ is required to reconstruct and audit the reference
terminal coefficient itself; the six nonzero residual settings are the
centrally symmetric first-derivative design.
\end{proposition}

\begin{proof}
Items (D1)--(D2) are the two extrapolations quantified in the next subsection.
Items (D3)--(D4) follow because Theorem~\ref{thm:v25-channel-tomography}
reconstructs the complete restricted channel outputs, not only their core
compression.
\end{proof}

\subsection{A complete acquisition error budget}
\label{sec:v25-error-budget}

For each residual coordinate $\eta$, define the exact antisymmetric channel
contrast
\begin{equation}
 A_\eta(\tau)
 :=\frac{\Phi_{\eta,+}(\tau)-\Phi_{\eta,-}(\tau)}{2\tau}.
 \label{eq:v25-Aeta}
\end{equation}
Assume the uniform expansions on the frozen finite core
\begin{align}
 A_\eta(\tau)
 &=\mathcal L_\eta+c_2(\eta)\tau^2+c_4(\eta)\tau^4
   +O(\tau^6),
 \label{eq:v25-time-expansion}\\
 \mathcal L_\eta
 &=\mathcal L_0+\eta^\alpha\mathcal K_\alpha
   +\frac12\eta^\alpha\eta^\beta\mathcal K_{\alpha\beta}
   +O(|\eta|^3).
 \label{eq:v25-eta-expansion}
\end{align}
These are finite regularity hypotheses for the acquisition panel, not claims
about a nonperturbative continuum channel family.

Define the exact Richardson germ estimator
\begin{equation}
 \mathcal R_t(\eta)
 :=\frac{4A_\eta(t/2)-A_\eta(t)}{3}
 \label{eq:v25-Richardson}
\end{equation}
and the exact centered residual estimator
\begin{equation}
 \mathcal K_{\alpha;t,\epsilon}^{\rm est}
 :=\frac{\mathcal R_t(+\epsilon e_\alpha)
          -\mathcal R_t(-\epsilon e_\alpha)}{2\epsilon}.
 \label{eq:v25-K-est}
\end{equation}

\begin{lemma}[Deterministic truncation error]
\label{lem:v25-truncation-error}
There are finite panel constants $C_t,C_\eta$ such that
\begin{equation}
 \boxed{
 \|\mathcal K_{\alpha;t,\epsilon}^{\rm est}-\mathcal K_\alpha\|
 \le C_t t^4+C_\eta\epsilon^2.}
 \label{eq:v25-det-error}
\end{equation}
\end{lemma}

\begin{proof}
The combination \eqref{eq:v25-Richardson} cancels the $c_2\tau^2$ term and
leaves a uniform $O(t^4)$ remainder.  The centered residual difference cancels
all even powers in $\eta$ and leaves an $O(\epsilon^2)$ derivative remainder.
The panel is finite, so the relevant Taylor suprema give finite constants.
\end{proof}

Now suppose every acquired channel restriction has a tomographic error bounded
by
\begin{equation}
 \|\widehat\Phi_{\eta,\sigma}(\tau)\circ\iota
       -\Phi_{\eta,\sigma}(\tau)\circ\iota\|
 \le\delta
 \label{eq:v25-row-error}
\end{equation}
for the two declared times and seven residual settings.

\begin{theorem}[Noise amplification of the terminal-germ acquisition]
\label{thm:v25-noise-amplification}
Let $\widehat{\mathcal K}_{\alpha;t,\epsilon}$ be obtained by applying
\eqref{eq:v25-Aeta}--\eqref{eq:v25-K-est} to the measured channel rows.  Then
\begin{equation}
 \boxed{
 \|\widehat{\mathcal K}_{\alpha;t,\epsilon}-\mathcal K_\alpha\|
 \le
 C_t t^4+C_\eta\epsilon^2
 +\frac{3\delta}{\epsilon t}.}
 \label{eq:v25-total-error}
\end{equation}
The same inequality holds for the full-cell restricted superoperator and hence
for its source-core leakage part.
\end{theorem}

\begin{proof}
At time $t$, the two branch errors contribute at most
$\delta/t$ to $A_\eta(t)$.  At time $t/2$, they contribute at most
$2\delta/t$.  Therefore the Richardson combination has noise at most
\[
 \frac{4(2\delta/t)+\delta/t}{3}=\frac{3\delta}{t}.
\]
Taking a centered difference of the $+\epsilon e_\alpha$ and
$-\epsilon e_\alpha$ Richardson rows divides the sum of two such errors by
$2\epsilon$, giving $3\delta/(\epsilon t)$.  Add
Lemma~\ref{lem:v25-truncation-error}.
\end{proof}

\begin{corollary}[Balanced acquisition scaling]
\label{cor:v25-balanced-scaling}
If the regularity constants remain bounded and the tomographic error
$\delta\downarrow0$, the choices
\begin{equation}
 t\asymp\delta^{1/7},
 \qquad
 \epsilon\asymp\delta^{2/7}
 \label{eq:v25-balanced-scales}
\end{equation}
balance the three terms in \eqref{eq:v25-total-error} and give
\begin{equation}
 \boxed{
 \|\widehat{\mathcal K}_{\alpha;t,\epsilon}-\mathcal K_\alpha\|
 =O(\delta^{4/7}).}
 \label{eq:v25-balanced-rate}
\end{equation}
This is an acquisition-design statement, not an assertion that a particular
physical experiment has already reached this asymptotic regime.
\end{corollary}

\begin{proof}
With $t=\delta^{1/7}$ and $\epsilon=\delta^{2/7}$, the three terms scale as
$t^4=\delta^{4/7}$, $\epsilon^2=\delta^{4/7}$ and
$\delta/(\epsilon t)=\delta^{4/7}$.
\end{proof}

\subsection{A held-out temporal falsification row}
\label{sec:v25-heldout}

The V24 handoff required a smaller-time row not used in the primary fit.  It
has a precise purpose: falsify the assumed two-time asymptotic model before the
reconstructed germ is used downstream.

At any selected residual setting define
\begin{equation}
 \mathcal R_{t/2}(\eta)
 :=\frac{4A_\eta(t/4)-A_\eta(t/2)}{3}.
 \label{eq:v25-heldout-R}
\end{equation}

\begin{proposition}[Held-out Richardson consistency]
\label{prop:v25-heldout}
Under \eqref{eq:v25-time-expansion},
\begin{equation}
 \boxed{
 \|\mathcal R_t(\eta)-\mathcal R_{t/2}(\eta)\|
 \le C_{\rm ho}t^4.}
 \label{eq:v25-heldout-bound}
\end{equation}
With noisy channel data the right side acquires the explicit sum of the
corresponding Richardson noise bounds.  A held-out discrepancy larger than the
declared truncation-plus-noise budget invalidates the extrapolation hypothesis
for that row and stops the terminal-germ compiler before any physical
coefficient is reconstructed.
\end{proposition}

\begin{proof}
Both Richardson estimators equal $\mathcal L_\eta+O(t^4)$, with the second
having the smaller $O((t/2)^4)$ remainder.  Subtraction gives the result.
\end{proof}

\subsection{Robust inner-derivation reconstruction}
\label{sec:v25-inner-robust}

Retain the V24 commutator synthesis
\begin{equation}
 \mathcal C_\perp:\mathfrak c^\perp\longrightarrow\mathcal Y,
 \label{eq:v25-Cperp}
\end{equation}
where $\mathfrak c$ is the compatible commutant and $\mathcal Y$ is the finite
operation-derivative data space.  Put
\begin{equation}
 \sigma_C:=s_{\min}(\mathcal C_\perp)>0,
 \qquad
 \mathcal C_\perp^\sharp
 =(\mathcal C_\perp^*\mathcal C_\perp)^{-1}\mathcal C_\perp^*.
 \label{eq:v25-Csharp}
\end{equation}
Let $P_C=\mathcal C_\perp\mathcal C_\perp^\sharp$.

\begin{theorem}[Certified Hamiltonian tangent reconstruction]
\label{thm:v25-H-cert}
Suppose the exact residual column is $Y_\alpha$ and the measured estimator
$\widehat Y_\alpha$ satisfies
\begin{equation}
 \|\widehat Y_\alpha-Y_\alpha\|\le\varepsilon_K,
 \label{eq:v25-eK}
\end{equation}
where $\varepsilon_K$ may be taken from
\eqref{eq:v25-total-error}.  Define
\begin{equation}
 \widehat H_\alpha
 :=\mathcal C_\perp^\sharp\widehat Y_\alpha,
 \qquad
 \widehat R_\alpha
 :=(I-P_C)\widehat Y_\alpha.
 \label{eq:v25-H-R-hat}
\end{equation}
Then
\begin{align}
 \|\widehat H_\alpha-H_\alpha^{\rm proj}\|
 &\le\sigma_C^{-1}\varepsilon_K,
 \label{eq:v25-H-error}\\
 \|\widehat R_\alpha-R_\alpha\|
 &\le\varepsilon_K,
 \label{eq:v25-R-error}
\end{align}
where
$H_\alpha^{\rm proj}=\mathcal C_\perp^\sharp Y_\alpha$ and
$R_\alpha=(I-P_C)Y_\alpha$.
In particular,
\begin{equation}
 \boxed{
 \|\widehat R_\alpha\|>\varepsilon_K
 \quad\Longrightarrow\quad
 R_\alpha\ne0.}
 \label{eq:v25-residual-reject}
\end{equation}
Thus a residual exceeding the acquisition budget is a certified rejection of
the inner-derivation branch on the selected operation core.
\end{theorem}

\begin{proof}
The pseudoinverse norm is
$\|\mathcal C_\perp^\sharp\|=\sigma_C^{-1}$.  The projector $I-P_C$ has norm
one.  Apply these operators to
$\widehat Y_\alpha-Y_\alpha$.  The final implication is the reverse triangle
inequality.
\end{proof}

\begin{firewall}[Small measured residual is not exact innerness]
A finite noisy packet satisfying
$\|\widehat R_\alpha\|\le\varepsilon_K$ is compatible with the inner branch but
does not prove $R_\alpha=0$.  Exact innerness is proved only by an exact
identity, by a convergent acquisition sequence whose residual tends to zero, or
by a separate structural theorem.  V25 therefore distinguishes
\emph{certified rejection} from \emph{failure to reject}.
\end{firewall}

\subsection{Certified physical coefficient sensitivity}
\label{sec:v25-coefficient-cert}

Let
\begin{equation}
 \mathcal J_D:\mathbb R^m\longrightarrow\mathfrak c^\perp
 \label{eq:v25-JD}
\end{equation}
be the V24 physical coefficient tangent synthesis for a declared block
$c=(c^1,\ldots,c^m)$, after redundant/commutant directions have been removed.
Assume full column rank and put
\begin{equation}
 \sigma_D:=s_{\min}(\mathcal J_D)>0,
 \qquad
 \mathcal J_D^\sharp
 =(\mathcal J_D^*\mathcal J_D)^{-1}\mathcal J_D^*.
 \label{eq:v25-JDsharp}
\end{equation}

\begin{theorem}[Coefficient-incidence certificate]
\label{thm:v25-coefficient-certificate}
For the reconstructed Hamiltonian columns define
\begin{equation}
 \widehat m_\alpha
 :=\mathcal J_D^\sharp\widehat H_\alpha,
 \qquad
 \widehat Q_\alpha
 :=(I-\mathcal J_D\mathcal J_D^\sharp)\widehat H_\alpha.
 \label{eq:v25-m-Q-hat}
\end{equation}
If the exact projected Hamiltonian error obeys
$\|\widehat H_\alpha-H_\alpha^{\rm proj}\|\le\varepsilon_H$, then
\begin{align}
 \|\widehat m_\alpha-m_\alpha^{\rm proj}\|
 &\le\sigma_D^{-1}\varepsilon_H,
 \label{eq:v25-m-error}\\
 \|\widehat Q_\alpha-Q_\alpha\|
 &\le\varepsilon_H.
 \label{eq:v25-Q-error}
\end{align}
Therefore
\begin{equation}
 \boxed{
 \|\widehat Q_\alpha\|>\varepsilon_H
 \quad\Longrightarrow\quad
 Q_\alpha\ne0,}
 \label{eq:v25-Q-reject}
\end{equation}
which certifies that the declared physical coefficient block is incomplete or
incorrect for that germ direction.
\end{theorem}

\begin{proof}
Identical to Theorem~\ref{thm:v25-H-cert}, with
$\mathcal C_\perp$ replaced by $\mathcal J_D$.
\end{proof}

Collect the three reconstructed coefficient columns into
\begin{equation}
 \widehat{\mathsf M}_D
 =\begin{pmatrix}\widehat m_1&\widehat m_2&\widehat m_3\end{pmatrix}
 \in\mathbb R^{m\times3}.
 \label{eq:v25-MD-hat}
\end{equation}
The propagated operator-norm uncertainty is denoted
$\varepsilon_M$.

\subsection{When can a predictive left nullspace be certified?}
\label{sec:v25-rank-cert}

Let $\mathsf K_{\rm phys}$ denote the already calibrated linear map from the
selected terminal coefficient block to the V18 physical Wilson/observable
quotient, and define
\begin{equation}
 \mathsf S
 :=\mathsf K_{\rm phys}\mathsf M_D,
 \qquad
 \widehat{\mathsf S}
 :=\widehat{\mathsf K}_{\rm phys}\widehat{\mathsf M}_D.
 \label{eq:v25-S}
\end{equation}
Assume all calibration and acquisition errors have been propagated to a single
bound
\begin{equation}
 \|\widehat{\mathsf S}-\mathsf S\|_{\rm op}
 \le\varepsilon_S.
 \label{eq:v25-eS}
\end{equation}

\begin{proposition}[Stable numerical rank and predictive subspace]
\label{prop:v25-rank-stability}
Suppose the exact sensitivity matrix $\mathsf S$ has rank $r$ and smallest
nonzero singular value
\begin{equation}
 g:=\sigma_r(\mathsf S)>2\varepsilon_S.
 \label{eq:v25-gap}
\end{equation}
Then the $r$-dimensional column space of $\mathsf S$ and the corresponding
column space extracted from $\widehat{\mathsf S}$ are separated by an angle of
order at most
\begin{equation}
 \boxed{
 \frac{\varepsilon_S}{g-\varepsilon_S}.}
 \label{eq:v25-subspace-angle}
\end{equation}
The same bound controls the orthogonal predictive conormal spaces
$\ker\mathsf S^T$ and the estimated left-null subspace.
\end{proposition}

\begin{proof}
This is the finite-dimensional singular-subspace perturbation estimate obtained
by applying the resolvent/projector bound to the separated singular spectrum,
or equivalently the standard Wedin bound.  The singular gap between the
nonzero sector and zero is $g$; perturbation by at most $\varepsilon_S<g/2$
preserves that separation and yields \eqref{eq:v25-subspace-angle}.
\end{proof}

\begin{firewall}[Predictivity still requires selector completeness]
Even an exactly measured nonzero left nullspace of the three-column matrix
$\mathsf S$ is a genuine Renewal prediction only if the residual coordinates
$\eta^1,\eta^2,\eta^3$ have independently been proved to exhaust the
microscopic selector directions for the chosen physical coefficient block after
P1 calibration.  Without that completeness theorem, V25 produces a
\emph{conditional selector conormal}: a relation insensitive to the three
measured residual directions, not a relation insensitive to every allowed
microscopic matching parameter.
\end{firewall}

\begin{corollary}[Structural codimension if the three-selector theorem is proved]
\label{cor:v25-structural-codim}
If the three residual coordinates are proved to be the complete microscopic
selector packet for an $m$-dimensional physical coefficient block, then
\begin{equation}
 \boxed{
 \dim\ker\mathsf S^T\ge m-3.}
 \label{eq:v25-m-minus-3}
\end{equation}
This lower bound is exact linear algebra and does not depend on the measured
rank being three.  What remains data dependent is the orientation of the
predictive conormal and whether the physical block itself has been chosen
correctly.
\end{corollary}

\begin{proof}
A matrix with three columns has rank at most three.  Apply rank-nullity to its
transpose.
\end{proof}

\subsection{An exact unitary calibration control}
\label{sec:v25-control}

The compiler admits a simple exact control family which tests the algebra
without being confused with a physical prediction.  Let
\begin{equation}
 D_\eta=D_0+\eta^\alpha H_\alpha,
 \qquad
 H_\alpha=H_\alpha^*,
 \label{eq:v25-control-D}
\end{equation}
and define
\begin{equation}
 \Phi_{\eta,+}(t)=\operatorname{Ad}_{e^{-itD_\eta}},
 \qquad
 \Phi_{\eta,-}(t)=\operatorname{Ad}_{e^{+itD_\eta}}.
 \label{eq:v25-control-channel}
\end{equation}

\begin{proposition}[Exact control recovery]
\label{prop:v25-control-recovery}
For the unitary control family,
\begin{equation}
 \mathcal L_\eta=-i[D_\eta,\cdot],
 \qquad
 \mathcal K_\alpha=-i[H_\alpha,\cdot].
 \label{eq:v25-control-germ}
\end{equation}
The inner-derivation residual is identically zero, and the V24/V25
reconstruction returns $H_\alpha$ modulo the compatible commutant.  If the
chosen coefficient tangent synthesis contains the $H_\alpha$, its
coefficient-incidence residual is also zero.  Failure of an implementation to
recover this control therefore diagnoses the acquisition/compiler pipeline,
not new phenomenology.
\end{proposition}

\begin{proof}
Differentiate \eqref{eq:v25-control-channel} at $t=0$ and then in $\eta$.
The remaining statements are the defining orthogonal projections of the V24
compiler.
\end{proof}

\subsection{The P4h acquisition certificate}
\label{sec:v25-certificate}

The preceding results can be summarized as a finite typed certificate.  A
populated terminal-germ acquisition is accepted only if the following ordered
checks pass:
\begin{enumerate}[label=\textup{(A\arabic*)},leftmargin=3em]
\item every one of the twenty-eight declared channel settings is reconstructed
on the frozen source core with its leakage table and row uncertainty;
\item duration/source gauge is common across all rows;
\item the held-out smaller-time residual is within its declared
truncation-plus-noise budget;
\item each of the three residual germ columns has an inner-derivation residual
consistent with zero at the claimed precision, or else the branch is rejected;
\item each reconstructed Hamiltonian tangent lies in the declared physical
coefficient tangent range, or else that coefficient block is rejected;
\item the mixed operation/action incidence required by V23 is populated and its
marked Gram is closed before an action-response susceptibility is evaluated;
\item selector completeness is separately proved before a left nullspace is
called a physical prediction.
\end{enumerate}

\begin{theorem}[Finite terminal-germ sufficiency]
\label{thm:v25-finite-sufficiency}
At fixed finite operation/source panel, assumptions
\textup{(A1)--(A7)} are sufficient to turn the V24 raw acquisition into either:
\begin{enumerate}[label=\textup{(S\arabic*)},leftmargin=3em]
\item a typed falsification residual identifying the first failed incidence
row; or
\item a finite coefficient-sensitivity matrix $\widehat{\mathsf M}_D$ with an
explicit propagated error budget and, after the V23 marked action gate, a
finite physical sensitivity matrix $\widehat{\mathsf S}$ with a certified
stable predictive subspace whenever the selector-completeness theorem is also
available.
\end{enumerate}
No continuum, RG, or EFT theorem downstream can replace a failed acquisition
condition.
\end{theorem}

\begin{proof}
Theorems~\ref{thm:v25-channel-tomography} and
\ref{thm:v25-noise-amplification} reconstruct the germ columns with controlled
error.  Theorem~\ref{thm:v25-H-cert} gives the inner-derivation decision and
Hamiltonian tangent, Theorem~\ref{thm:v25-coefficient-certificate} gives the
coefficient tangent, and Proposition~\ref{prop:v25-rank-stability} gives the
stable conormal after the V23 action-incidence gate.  Each operation is finite
linear algebra on the frozen panel.  The firewalls above explain why no failed
row can be supplied by a later theorem.
\end{proof}

\begin{firewall}[V25 is a sufficiency theorem, not a fabricated acquisition]
The present version proves that the P4h panel is finite and mathematically
sufficient.  It does not assign numerical values to any of the twenty-eight
forward/reflected channel maps, does not set their residuals to zero, and does
not report a numerical $\mathsf M_D$ or predictive left nullspace.  Those
objects remain unpopulated until a literal same-history differentiated CPTP
packet is supplied.
\end{firewall}

\section{V25 anti-circularity audit}
\label{sec:v25-audit}

\begin{enumerate}[leftmargin=2.4em]
\item \textbf{A channel row was reduced to finite tomography, not treated as an
oracle.}  The core-input full-cell coefficient table explicitly includes the
leakage data required by the terminal audit.
\item \textbf{Derivative amplification is paid explicitly.}  The
$3\delta/(\epsilon t)$ term prevents arbitrarily small $t$ or $\epsilon$ from
being called a free accuracy gain.
\item \textbf{Small residual is not exact incidence.}  V25 permits certified
rejection but distinguishes it from exact innerness or exact coefficient
membership.
\item \textbf{A left nullspace is not yet a prediction without selector
completeness.}  The three-column matrix can only constrain the microscopic
three-direction image actually measured.
\item \textbf{The held-out row is a falsification row.}  It is not used to tune
the same Richardson estimator it is supposed to test.
\item \textbf{No downstream substitution is allowed.}  RG running, continuum
existence, determinant-line triviality and reflected field dynamics cannot
manufacture an unmeasured terminal-germ row.
\end{enumerate}

\section{V25 status ledger}
\label{sec:v25-status}

\subsection*{Proved in V25}
\begin{enumerate}[leftmargin=2.2em]
\item One finite forward/reflected channel row on the frozen source core is
fully determined by a finite Hermitian process-tomography table, including the
source-core leakage needed by the terminal audit.
\item The twenty-eight-row V24 design is sufficient for one reference germ and
three centrally differenced residual germ columns under the declared two-time
and reversal-symmetric expansion.
\item The complete deterministic-plus-tomographic error budget is
$C_tt^4+C_\eta\epsilon^2+3\delta/(\epsilon t)$; the balanced asymptotic design
has $t\asymp\delta^{1/7}$, $\epsilon\asymp\delta^{2/7}$ and total error
$O(\delta^{4/7})$.
\item A held-out smaller-time Richardson row gives an explicit falsification
test for the assumed temporal expansion.
\item The V24 inner-derivation inversion and physical coefficient matching have
explicit stability constants given by the smallest singular values of their
finite design Grams.  Residuals larger than the propagated acquisition budget
certify rejection of the corresponding branch.
\item A finite error bound on the physical sensitivity matrix gives a stable
predictive conormal whenever its exact nonzero singular sector is separated.
\item If the three residual microscopic coordinates are independently proved to
be the complete selector packet for an $m$-dimensional physical block, the
predictive conormal has dimension at least $m-3$.
\item A unitary calibration-control family exactly exercises the compiler and
separates implementation validation from physical prediction.
\end{enumerate}

\subsection*{Inherited}
\begin{enumerate}[leftmargin=2.2em]
\item V16--V24: physical calibration/RG covariance, no-free-Wilson theorem,
residual race coordinates, action/path-selector no-go results, operation/action
incidence separation, the terminal same-history differential demand, exact
commutator projection and finite physical coefficient tangent synthesis.
\item The Grand Tensor forward/reverse CPTP tangent theorem, source-core leakage
Pythagoras, canonical Choi-source reconstruction, and the positive rank-one
first terminal witness.
\end{enumerate}

\subsection*{Conditional}
\begin{enumerate}[leftmargin=2.2em]
\item The $O(t^4)$ Richardson bias uses the V24 reversal-symmetric expansion;
a weaker time regularity changes the design order and error budget.
\item Tomographic uncertainty must be calibrated in a norm controlling the
finite superoperator residuals used by the compiler.
\item Exact innerness and exact coefficient incidence cannot be inferred from a
single finite noisy packet merely because its measured residual lies below the
error bar.
\item Predictivity requires a separate theorem that the three residual
coordinates exhaust all microscopic selector directions for the declared
physical coefficient block.
\item The operation-to-action mixed incidence and V23 marked action-head closure
remain physically unpopulated.
\end{enumerate}

\subsection*{Open beyond V25}
\begin{enumerate}[leftmargin=2.2em]
\item Populate the twenty-eight forward/reflected channel maps and the held-out
smaller-time row on one literal same-history represented ESM carrier, with
common duration/source calibration and full-cell leakage tomography.
\item Evaluate the three inner-derivation residuals and, on the accepted branch,
recover the Hermitian germ derivatives and their physical coefficient columns.
\item Populate the mixed operation/action incidence, close the V23 marked Gram,
and push the resulting finite sensitivity through the V18 physical quotient.
\item Prove or falsify completeness of the three residual race coordinates for
the first selected physical coefficient block.  Only then compute and interpret
the physical predictive left nullspace.
\item The parallel phenomenology programme remains available after the first
nontrivial relation lands: threshold matching, standard-scheme conversion,
higher-loop RG and numerical trajectory/observable propagation.
\end{enumerate}

\subsection*{Exact next load-bearing theorem}
\begin{center}
\fbox{\parbox{0.92\textwidth}{
\textbf{P4i: populated terminal-germ certificate and selector-completeness test.}
Acquire the V24/V25 twenty-eight-row same-history forward/reflected CPTP panel,
its full-cell leakage tables, common duration/source calibration and held-out
smaller-time data.  Apply the V25 error budget, test the inner-derivation and
physical-coefficient residuals, and obtain the first numerically populated
three-column matrix $\mathsf M_D$ or a typed falsification witness.  In parallel,
prove whether the three residual race coordinates exhaust the microscopic
selector directions of the chosen physical coefficient block.  Only if both
the acquisition and completeness gates pass may the resulting left nullspace be
promoted to the first Renewal-specific phenomenological relation.}}
\end{center}

\section*{Append-only continuation marker: V25}
\addcontentsline{toc}{section}{Append-only continuation marker: V25}
\textbf{End of V25.}  The complete V1--V24 body above is retained verbatim.
V25 does not fabricate the missing terminal channel rows.  It proves instead
that the V24 acquisition is finite and tomographically complete on the frozen
operation core, derives its full truncation/noise amplification budget and a
balanced acquisition rate, supplies a held-out temporal falsification test,
and propagates the acquisition uncertainty through the inner-derivation and
physical-coefficient inversions to a stable predictivity subspace.  The remaining
primary gate is literal population of the same-history channel panel together
with an independent selector-completeness theorem; without those data no
numerical microscopic sensitivity or Renewal-specific coefficient relation is
claimed.



\clearpage
\section{V26 continuation: exact race-fibre selector completeness and the no-extra-selector certificate}
\label{sec:v26-continuation}

V25 closed the finite acquisition compiler conditional on two logically
separate inputs: literal terminal-germ data and a theorem that the measured
microscopic selector columns exhaust the selector directions relevant to the
chosen physical coefficient block.  The second requirement should be attacked
before interpreting any numerical left nullspace.  The present continuation
goes upstream to the calibrated race manifold and separates two notions which
were deliberately left adjacent in V19--V25:
\begin{equation}
 \boxed{
 \begin{gathered}
 \text{completeness inside the declared race family}\\
 \neq\\
 \text{completeness for the full Renewal ESM coefficient block}
 \end{gathered}}
 \label{eq:v26-two-completenesses}
\end{equation}

The first statement can now be proved exactly.  The V19 race decoder is a
rank-ten map from thirteen microscopic coordinates, and its three-dimensional
fibre is not merely one convenient candidate slice: after P1 calibration it is
the complete residual manifold of the independently exponential race model.
The second statement requires auditing every additional unfixed microscopic
coordinate which can move the selected Wilson quotient.  V26 gives an exact
finite residual for that audit.  A left nullspace of the three-column race
matrix is a full Renewal prediction only when the residual vanishes.

Throughout V26 retain the V19 variables
\begin{equation}
 \vartheta_{\rm race}=(s,d,a_\rho)
 \in\mathbb R^6\times\mathbb R^6\times\mathbb R,
 \qquad
 S=\operatorname{diag}(s),
 \qquad
 D=\operatorname{diag}(d),
 \label{eq:v26-race-vars}
\end{equation}
and the calibrated packet
\begin{equation}
 \mathcal G_{\rm race}(s,d,a_\rho)
 =\bigl(\operatorname{vec}_{\rm sym}(RSR^T),Rd,a_\rho\bigr).
 \label{eq:v26-race-calibration}
\end{equation}
The inherited assumptions are
\(\operatorname{rank}R=3\) and linear independence of the six dyadics
\(r_ar_a^T\).  Choose once and for all a full-column-rank matrix
\begin{equation}
 N:\mathbb R^3\longrightarrow\mathbb R^6,
 \qquad
 \operatorname{im}N=\ker R.
 \label{eq:v26-N}
\end{equation}

\subsection{The calibrated race fibre is exactly three-dimensional}
\label{sec:v26-race-fibre}

Fix a protected P1 calibration point
\begin{equation}
 p_{\rm race}=(\mathcal B_0,\beta_0,a_{\rho,0})
 \label{eq:v26-prace}
\end{equation}
inside the open rate inequalities of V19--V20.  Let
\begin{equation}
 \mathcal F_{\rm race}(p_{\rm race})
 :=\{(s,d,a_\rho):
       \mathcal G_{\rm race}(s,d,a_\rho)=p_{\rm race}\}
 \label{eq:v26-race-fibre-def}
\end{equation}
intersected with that protected rate domain.

\begin{theorem}[Exact selector completeness of the calibrated race family]
\label{thm:v26-race-complete}
On the protected component containing the V19 base point, there are unique
\(s_0\) and \(a_{\rho,0}\) fixed by
\(\mathcal B_0\) and \(a_{\rho,0}\), and every point of the calibrated race
fibre has the form
\begin{equation}
 \boxed{
 s=s_0,
 \qquad
 a_\rho=a_{\rho,0},
 \qquad
 d=d_*+N\eta,
 \quad \eta\in U_\eta\subset\mathbb R^3,}
 \label{eq:v26-race-fibre-param}
\end{equation}
where \(Rd_*=\beta_0\) and \(U_\eta\) is the open set cut out by the protected
rate inequalities.  Consequently
\begin{equation}
 \boxed{
 T\mathcal F_{\rm race}
 =\{(0,N\delta\eta,0):\delta\eta\in\mathbb R^3\}.}
 \label{eq:v26-race-fibre-tangent}
\end{equation}
Thus the V19 residual coordinates \(\eta\) exhaust every microscopic
deformation of the declared independently exponential race law which leaves
all ten P1 race-derived calibration coordinates fixed.
\end{theorem}

\begin{proof}
The six dyadics \(r_ar_a^T\) are linearly independent in the six-dimensional
space \(\operatorname{Sym}_3\).  Hence the linear map
\[
 s\longmapsto R\operatorname{diag}(s)R^T
\]
is an isomorphism onto its image, which is all of
\(\operatorname{Sym}_3\) on the declared branch.  Therefore
\(\mathcal B=\mathcal B_0\) fixes \(s=s_0\) uniquely.  The last calibration
coordinate fixes \(a_\rho=a_{\rho,0}\).  The remaining equation is
\(Rd=\beta_0\).  Its solution set is the affine space
\(d_*+\ker R=d_*+\operatorname{im}N\).  Intersecting with the strict protected
rate inequalities leaves an open subset \(U_\eta\) of this affine
three-plane.  Differentiating gives \eqref{eq:v26-race-fibre-tangent}.
\end{proof}

\begin{corollary}[Nonlinear, not merely infinitesimal, completeness]
\label{cor:v26-nonlinear-complete}
Let \(F\) be any coefficientwise smooth microscopic functional defined only on
the protected race family.  After P1 fixes \(p_{\rm race}\), its restriction to
the calibrated race fibre is a unique function
\begin{equation}
 F\big|_{\mathcal F_{\rm race}}
 =\widetilde F(\eta).
 \label{eq:v26-F-eta}
\end{equation}
Therefore no additional independent race-law selector coordinate exists on
that calibrated component, at finite amplitude or at tangent level.
\end{corollary}

\begin{proof}
Theorem~\ref{thm:v26-race-complete} is a global affine parametrization of the
protected fibre by \(\eta\), so restriction and composition with the inverse
chart give \(\widetilde F\).
\end{proof}

\subsection{The Fisher geometry singles out the same residual fibre}
\label{sec:v26-Fisher-fibre}

The completeness theorem is algebraic.  The V19--V20 information metric gives
a second, independent characterization of the same three directions.  Put
\begin{equation}
 W:=\bigl(S-h^2DS^{-1}D\bigr)^{-1}\succ0.
 \label{eq:v26-W}
\end{equation}
At fixed \(s\), the race Fisher form on drift variations is
\begin{equation}
 \langle u,v\rangle_{\rm race}
 :=\frac{h^2}{\lambda}u^TWv.
 \label{eq:v26-race-inner}
\end{equation}
The V20 residual Fisher matrix is
\begin{equation}
 G_\eta=\frac{h^2}{\lambda}N^TWN.
 \label{eq:v26-Geta}
\end{equation}

\begin{proposition}[Canonical information-orthogonal residual projector]
\label{prop:v26-residual-projector}
Define
\begin{equation}
 \boxed{
 P_{\rm res}^{W}
 :=N(N^TWN)^{-1}N^TW.}
 \label{eq:v26-Pres}
\end{equation}
Then \(P_{\rm res}^{W}\) is the \(W\)-orthogonal projection onto
\(\ker R\).  Its complement gives the unique information-minimal lift of a
fixed physical drift variation \(\delta\beta\).  In particular, every drift
tangent decomposes uniquely as
\begin{equation}
 \boxed{
 \delta d=\delta d_{\rm phys}+N\delta\eta,
 \qquad
 \langle\delta d_{\rm phys},N\delta\eta\rangle_{\rm race}=0,}
 \label{eq:v26-info-split}
\end{equation}
and the Fisher cost obeys the exact Pythagoras identity
\begin{equation}
 \|\delta d\|_{\rm race}^2
 =\|\delta d_{\rm phys}\|_{\rm race}^2
 +\delta\eta^TG_\eta\delta\eta.
 \label{eq:v26-info-Pythagoras}
\end{equation}
\end{proposition}

\begin{proof}
The formula \eqref{eq:v26-Pres} is the standard weighted projection onto the
full-column-rank range of \(N\): it is idempotent, has range
\(\operatorname{im}N=\ker R\), and satisfies
\(W P_{\rm res}^{W}=(P_{\rm res}^{W})^TW\).  The inherited efficient race
lift is precisely the minimizer in the affine constraint \(R\delta d=\delta
\beta\); its Euler equation is orthogonality to \(\ker R\) in the
\(W\)-metric.  This gives \eqref{eq:v26-info-split}.  Orthogonality and
\eqref{eq:v26-Geta} give \eqref{eq:v26-info-Pythagoras}.
\end{proof}

\begin{corollary}[Whitening does not change selector completeness]
\label{cor:v26-whiten-complete}
Replacing \(\eta\) by the V20 whitened coordinate
\(\widehat\eta=G_\eta^{1/2}\eta\) multiplies every race-selector sensitivity
matrix on the right by an invertible \(3\times3\) matrix.  Its image, rank and
left nullspace are unchanged.
\end{corollary}

\subsection{Race-family completeness is not full ESM selector completeness}
\label{sec:v26-full-selector}

Theorem~\ref{thm:v26-race-complete} closes the selector problem only inside
the declared race model.  The complete Renewal ESM construction also carries
other microscopic or finite matching coordinates: matter/action heads,
source-incidence variables, boundary data, and--at the present V18 stage--the
independent finite Wilson coordinates themselves.  Only those directions
which remain unfixed after P1 and actually move the selected physical
coefficient block belong in the selector domain for P4.

Let \(\mathcal W\) be a finite physical coefficient quotient of dimension
\(m\).  Suppose the calibrated microscopic tangent space relevant to this block
has been split as
\begin{equation}
 T\Theta_{\rm sel}
 =V_\eta\oplus V_\chi,
 \qquad
 \dim V_\eta=3,
 \label{eq:v26-selector-split}
\end{equation}
where \(V_\eta\) is the race fibre and \(V_\chi\) collects every additional
unfixed microscopic selector direction.  Write the physical sensitivity as
\begin{equation}
 \boxed{
 \mathsf S_{\rm full}
 =\bigl[\mathsf S_\eta\ \ \mathsf S_\chi\bigr]:
 V_\eta\oplus V_\chi\longrightarrow\mathcal W.}
 \label{eq:v26-Sfull}
\end{equation}
Here \(\mathsf S_\eta\) is the V24--V25 three-column terminal-germ sensitivity
once populated.

\begin{definition}[Block-specific selector-completeness residual]
\label{def:v26-selector-residual}
Let
\begin{equation}
 P_\eta:=\mathsf S_\eta\mathsf S_\eta^\dagger
 \label{eq:v26-Peta}
\end{equation}
be the orthogonal projector onto the race-selector image in the chosen physical
coefficient metric.  Define
\begin{equation}
 \boxed{
 \mathcal C_{\rm sel}
 :=(I-P_\eta)\mathsf S_\chi.}
 \label{eq:v26-Csel}
\end{equation}
We say that the race packet is \emph{complete for the coefficient block} when
\(\mathcal C_{\rm sel}=0\).
\end{definition}

\begin{theorem}[Exact no-extra-selector criterion]
\label{thm:v26-no-extra-selector}
The following are equivalent:
\begin{enumerate}[label=\textup{(C\arabic*)},leftmargin=3em]
\item \(\mathcal C_{\rm sel}=0\);
\item every additional microscopic selector sensitivity lies in the image of
the race-selector sensitivity,
\begin{equation}
 \operatorname{im}\mathsf S_\chi
 \subseteq
 \operatorname{im}\mathsf S_\eta;
 \label{eq:v26-image-inclusion}
\end{equation}
\item the full and race-only physical tangent images coincide,
\begin{equation}
 \operatorname{im}\mathsf S_{\rm full}
 =\operatorname{im}\mathsf S_\eta;
 \label{eq:v26-image-equality}
\end{equation}
\item the race-only and full predictive conormals coincide,
\begin{equation}
 \boxed{
 \ker\mathsf S_\eta^T
 =\ker\mathsf S_{\rm full}^T.}
 \label{eq:v26-conormal-equality}
\end{equation}
\end{enumerate}
Thus \(\mathcal C_{\rm sel}\) is the exact finite obstruction to promoting a
three-column left nullspace to the complete first-order predictive conormal of
the chosen coefficient block.
\end{theorem}

\begin{proof}
Because \(P_\eta\) is the orthogonal projector onto
\(\operatorname{im}\mathsf S_\eta\), condition
\((I-P_\eta)\mathsf S_\chi=0\) is equivalent to
\eqref{eq:v26-image-inclusion}.  The image of the block matrix
\eqref{eq:v26-Sfull} is the linear span of the two column spaces, so the
inclusion is equivalent to \eqref{eq:v26-image-equality}.  Finally, in a finite
inner-product space, the left nullspace of a map is the orthogonal complement
of its image.  Equal images are therefore equivalent to equal left nullspaces.
\end{proof}

\begin{corollary}[Full predictive conormal is an intersection]
\label{cor:v26-conormal-intersection}
Without selector completeness one has exactly
\begin{equation}
 \boxed{
 \ker\mathsf S_{\rm full}^T
 =\ker\mathsf S_\eta^T
  \cap
  \ker\mathsf S_\chi^T.}
 \label{eq:v26-conormal-intersection}
\end{equation}
Hence additional microscopic directions can only remove, never create,
predictive relations inferred from the three-column race matrix.
\end{corollary}

\begin{proof}
A covector annihilates the block matrix if and only if it annihilates each
column block separately.
\end{proof}

\subsection{The V18 identity block is an automatic completeness failure}
\label{sec:v26-identity-block}

The preceding criterion immediately explains why V18 had to keep the pilot
Wilson quotient nonpredictive.  Suppose a coefficient block of dimension \(m\)
remains among the independent finite microscopic matching inputs.  Those
coordinates contribute an identity sensitivity
\begin{equation}
 \mathsf S_{\rm Wilson}=I_m.
 \label{eq:v26-Wilson-identity}
\end{equation}

\begin{proposition}[Independent Wilson inputs destroy race-selector completeness]
\label{prop:v26-identity-kills}
If \(I_m\) occurs as a column block of \(\mathsf S_\chi\), then
\begin{equation}
 \operatorname{im}\mathsf S_{\rm full}=\mathcal W,
 \qquad
 \ker\mathsf S_{\rm full}^T=\{0\}.
 \label{eq:v26-identity-no-prediction}
\end{equation}
Moreover
\begin{equation}
 \mathcal C_{\rm sel}=I_m-P_\eta.
 \label{eq:v26-identity-Csel}
\end{equation}
Thus the race packet can be complete only if
\(\operatorname{rank}\mathsf S_\eta=m\), which is impossible for \(m>3\).
\end{proposition}

\begin{proof}
The identity columns span all of \(\mathcal W\), so the full image is
\(\mathcal W\) and its annihilator is zero.  Substitution into
\eqref{eq:v26-Csel} gives \eqref{eq:v26-identity-Csel}.  Vanishing requires
\(P_\eta=I_m\), hence full row rank of \(\mathsf S_\eta\), which is bounded by
three.
\end{proof}

\begin{firewall}[No prediction while the predicted coefficient is still an input]
Proposition~\ref{prop:v26-identity-kills} is the finite-dimensional form of the
V18 no-free-prediction theorem.  A future Renewal selection principle must
remove the claimed predictive Wilson coordinates from the independent input
list before the three residual race variables can constrain them.  Measuring a
three-column sensitivity matrix while retaining the identity block does not
change this conclusion.
\end{firewall}

\subsection{Conditional promotion theorem for a locked race-selector model}
\label{sec:v26-conditional-promotion}

A useful positive theorem is nevertheless available once the microscopic model
is actually locked.  By a \emph{locked race-selector model} for a coefficient
block we mean:
\begin{enumerate}[label=\textup{(L\arabic*)},leftmargin=3em]
\item all P1 physical inputs are fixed;
\item all additional microscopic coordinates affecting the block are either
fixed by independently measured inputs or are explicit functions of the race
fibre \(\eta\);
\item no claimed predictive Wilson coefficient is retained as an independent
matching coordinate;
\item the complete local matching map is a populated function
\begin{equation}
 \mathfrak M_{\rm lock}:U_\eta\longrightarrow\mathcal W.
 \label{eq:v26-Mlock}
\end{equation}
\end{enumerate}

\begin{theorem}[Race-fibre completeness promotes the three-column conormal]
\label{thm:v26-locked-promotion}
In a locked race-selector model, the V25 three-column terminal-germ sensitivity
is the complete microscopic sensitivity of the selected coefficient block:
\begin{equation}
 \boxed{
 D\mathfrak M_{\rm lock}=\mathsf S_\eta.}
 \label{eq:v26-DMlock}
\end{equation}
Consequently
\begin{equation}
 \boxed{
 \mathcal P^*_{\rm block}
 =\ker\mathsf S_\eta^T,}
 \label{eq:v26-block-prediction}
\end{equation}
and at a regular point
\begin{equation}
 \dim\mathcal P^*_{\rm block}
 =m-\operatorname{rank}\mathsf S_\eta
 \ge m-3.
 \label{eq:v26-block-dimension}
\end{equation}
Every such relation is invariant under an invertible reparametrization or
Fisher whitening of \(\eta\).
\end{theorem}

\begin{proof}
Theorem~\ref{thm:v26-race-complete} says that, once the P1 calibration is fixed,
there are no further independent coordinates inside the declared race family.
The locked-model hypotheses remove every extra selector direction outside that
family from the chosen coefficient problem.  Therefore the differential of the
full matching map has exactly the three race columns.  The V16 predictive
conormal is the annihilator of their image, which is
\eqref{eq:v26-block-prediction}; rank--nullity gives
\eqref{eq:v26-block-dimension}.  Right multiplication by an invertible
coordinate change does not alter the image.
\end{proof}

\begin{firewall}[Locked-model prediction is conditional on the model lock]
Theorem~\ref{thm:v26-locked-promotion} does not assert that the current full
Renewal ESM corpus already satisfies items \textup{(L1)--(L4)}.  It states the
precise additional theorem needed to turn the V25 terminal-germ matrix into a
complete block prediction.  At present the independent action/source rows and
finite Wilson matching coordinates prevent that promotion for the unrestricted
ESM coefficient bank.
\end{firewall}

\subsection{Finite selector-completeness certificate with uncertainty}
\label{sec:v26-noisy-certificate}

The V25 terminal-germ acquisition is noisy, and any future extra-selector
packet will be as well.  The completeness residual has a stable finite-data
version.

Let \(\widehat{\mathsf S}_\eta\) and \(\widehat{\mathsf S}_\chi\) be measured
matrices with
\begin{equation}
 \|\widehat{\mathsf S}_\eta-\mathsf S_\eta\|\le\varepsilon_\eta,
 \qquad
 \|\widehat{\mathsf S}_\chi-\mathsf S_\chi\|\le\varepsilon_\chi.
 \label{eq:v26-matrix-errors}
\end{equation}
Assume the smallest nonzero singular value of \(\mathsf S_\eta\) is
\(g_\eta>2\varepsilon_\eta\).  Let
\(\widehat P_\eta\) be the spectral projector of
\(\widehat{\mathsf S}_\eta\widehat{\mathsf S}_\eta^T\) onto its resolved
nonzero singular sector and define
\begin{equation}
 \widehat{\mathcal C}_{\rm sel}
 :=(I-\widehat P_\eta)\widehat{\mathsf S}_\chi.
 \label{eq:v26-Csel-hat}
\end{equation}

\begin{proposition}[Certified rejection of selector completeness]
\label{prop:v26-completeness-rejection}
There is a constant \(C\), depending only on the resolved singular-sector
conditioning, such that
\begin{equation}
 \|\widehat{\mathcal C}_{\rm sel}-\mathcal C_{\rm sel}\|
 \le
 \varepsilon_\chi
 +C\frac{\varepsilon_\eta}{g_\eta-\varepsilon_\eta}
   \bigl(\|\mathsf S_\chi\|+\varepsilon_\chi\bigr)
 =:\varepsilon_{\rm sel}.
 \label{eq:v26-Csel-error}
\end{equation}
Hence
\begin{equation}
 \boxed{
 \|\widehat{\mathcal C}_{\rm sel}\|>\varepsilon_{\rm sel}
 \quad\Longrightarrow\quad
 \mathcal C_{\rm sel}\ne0.}
 \label{eq:v26-Csel-reject}
\end{equation}
A measured residual below the error budget is only failure to reject
incompleteness; it is not promoted to exact selector completeness.
\end{proposition}

\begin{proof}
The resolved left singular subspace is separated from zero by
\(g_\eta\).  Standard finite-dimensional spectral-projector perturbation gives
\[
 \|\widehat P_\eta-P_\eta\|
 \le C\frac{\varepsilon_\eta}{g_\eta-\varepsilon_\eta}.
\]
Then
\[
 \widehat{\mathcal C}_{\rm sel}-\mathcal C_{\rm sel}
 =(I-\widehat P_\eta)(\widehat{\mathsf S}_\chi-\mathsf S_\chi)
 +(P_\eta-\widehat P_\eta)\mathsf S_\chi,
\]
and the triangle inequality gives \eqref{eq:v26-Csel-error}.  The rejection
criterion follows immediately.
\end{proof}

\subsection{Current selector-completeness status of the phenomenology programme}
\label{sec:v26-current-status}

The preceding results make the present status precise.

\begin{table}[ht]
\centering
\small
\begin{tabular}{@{}p{0.25\linewidth}p{0.23\linewidth}p{0.42\linewidth}@{}}
\toprule
Selector question & Status & Meaning\\
\midrule
Race family after P1 calibration & \textbf{proved complete} & The calibrated fibre is exactly \(d_*+\ker R\simeq\mathbb R^3\).\\
Residual race Fisher packet & \textbf{proved faithful} & \(G_\eta\succ0\); whitening/reparametrization does not change rank or conormal.\\
V25 terminal-germ sensitivity & \textbf{unpopulated} & The finite 28-row acquisition is sufficient, but literal channel data have not been supplied.\\
Extra microscopic selector block \(\chi\) & \textbf{not locked} & Independent action/source and finite matching coordinates remain outside the race fibre.\\
Current unrestricted Wilson bank & \textbf{not predictive} & Independent Wilson coordinates supply an identity block, as in V18 and Proposition~\ref{prop:v26-identity-kills}.\\
Locked race-selector model & \textbf{conditional theorem available} & Once extra directions are removed or proved image-redundant, \(\ker\mathsf S_\eta^T\) is the complete block conormal.\\
\bottomrule
\end{tabular}
\caption{V26 selector-completeness audit for the phenomenology line.}
\label{tab:v26-selector-status}
\end{table}

\begin{center}
\fbox{\parbox{0.92\textwidth}{
\textbf{V26 closes selector completeness inside the calibrated race model, but
not yet for the unrestricted Renewal ESM Wilson bank.}
The three residual coordinates are exactly all race-law deformations invisible
to P1 geometry.  For any chosen physical coefficient block, the only remaining
completeness question is finite: do the sensitivities of every additional
unfixed microscopic coordinate lie in the three-column image?  The exact
obstruction is
\((I-P_\eta)\mathsf S_\chi\).  Independent Wilson inputs make this obstruction
nonzero automatically unless the block has dimension at most three and the
race matrix is full row rank.}}
\end{center}

\section{V26 anti-circularity audit}
\label{sec:v26-audit}

\begin{enumerate}[leftmargin=2.4em]
\item \textbf{The three residual directions are no longer merely a candidate
slice.}  Their completeness is proved for the declared race family by the exact
calibration fibre, not inferred from a rank ceiling alone.
\item \textbf{Race-family completeness is not promoted to full ESM
completeness.}  Additional action, matter, source and finite matching directions
remain explicit in \(V_\chi\).
\item \textbf{A block-specific criterion replaces a global microscopic census.}
P4 needs only to know whether extra selectors enlarge the image in the chosen
coefficient quotient; this is exactly what \(\mathcal C_{\rm sel}\) tests.
\item \textbf{Independent Wilson inputs are not hidden.}  Their identity block
forces zero predictive conormal and therefore must be removed by a genuine
microscopic matching law before a Renewal relation can be claimed.
\item \textbf{No numerical terminal-germ matrix is fabricated.}  V25's 28-row
packet remains unpopulated; V26 only closes the domain/completeness geometry
which must accompany that acquisition.
\item \textbf{Small measured completeness residual is not exact completeness.}
The noisy certificate permits one-sided rejection and retains the acquisition
error budget explicitly.
\end{enumerate}

\section{V26 status ledger}
\label{sec:v26-status}

\subsection*{Proved in V26}
\begin{enumerate}[leftmargin=2.2em]
\item The P1-calibrated independently exponential race family is exactly a
three-dimensional affine fibre \(d_*+\ker R\) on its protected component.  The
V19 coordinates \(\eta\) exhaust all race-law selector directions which leave
\((\mathcal B,\beta,a_\rho)\) fixed.
\item The inherited race Fisher metric gives an exact orthogonal decomposition
between the information-minimal physical lift and the residual selector fibre;
the residual cost is \(\delta\eta^TG_\eta\delta\eta\).
\item For a general physical coefficient block with additional microscopic
selectors \(\chi\), the exact obstruction to three-selector completeness is
\(\mathcal C_{\rm sel}=(I-P_\eta)\mathsf S_\chi\).  It vanishes if and only if
the full and race-only tangent images, and hence predictive conormals, coincide.
\item The full predictive conormal is exactly
\(\ker\mathsf S_\eta^T\cap\ker\mathsf S_\chi^T\).
\item Retaining the physical Wilson coefficients as independent matching inputs
inserts an identity block, making the unrestricted block nonpredictive and
preventing race-selector completeness for any block of dimension greater than
three.
\item In a genuinely locked race-selector model, the V25 three-column
sensitivity is complete and its left nullspace is the full first-order block
prediction, with codimension at least \(m-3\).
\item A finite noisy selector-completeness residual gives a certified one-sided
rejection test once both race and extra-selector sensitivity matrices are
populated.
\end{enumerate}

\subsection*{Inherited}
\begin{enumerate}[leftmargin=2.2em]
\item V16--V25: physical calibration and RG covariance, no-free-Wilson theorem,
explicit race incidence, residual score/Fisher packet, action/operation
incidence compilers, the finite terminal-germ inversion and the complete
28-row acquisition/error certificate.
\item The Grand Tensor calibrated race-information theorem and its information-
minimal physical lift, including the strict positive residual Fisher metric.
\end{enumerate}

\subsection*{Conditional}
\begin{enumerate}[leftmargin=2.2em]
\item Full Renewal-specific predictivity still requires a microscopic model in
which the claimed predictive Wilson coordinates are not independently free.
\item The extra-selector block \(\chi\) must be explicitly enumerated for each
chosen physical coefficient quotient; V26 does not assume that the race family
is the entire Renewal ESM microscopic theory.
\item The V25 terminal channel packet and the operation-to-action mixed
incidence remain physically unpopulated.
\item The noisy completeness certificate assumes a resolved nonzero singular
sector for the measured race-selector matrix.
\end{enumerate}

\subsection*{Open beyond V26}
\begin{enumerate}[leftmargin=2.2em]
\item Freeze the first concrete physical coefficient block and enumerate every
microscopic coordinate which remains unfixed after P1 calibration.  Partition
those coordinates into measured inputs, redundant directions, the complete
race fibre and genuine extra selectors \(\chi\).
\item For every genuine extra selector, populate or bound its physical
sensitivity and evaluate the no-extra-selector residual
\((I-P_\eta)\mathsf S_\chi\).
\item Populate the V25 twenty-eight-row same-history terminal-germ acquisition
and obtain the first numerical \(\mathsf S_\eta\) or a typed falsification
witness.
\item Only after both the acquisition and no-extra-selector gates pass should a
left nullspace be promoted to a Renewal-specific phenomenological relation.
\item Threshold matching, standard-scheme conversion, numerical RG trajectories
and observable propagation remain downstream of the first genuine relation.
\end{enumerate}

\subsection*{Exact next load-bearing theorem}
\begin{center}
\fbox{\parbox{0.92\textwidth}{
\textbf{P4j: selector-domain locking and no-extra-selector incidence.}
Choose the first finite physical gravity--matter coefficient block.  After P1
calibration, list every microscopic coordinate that can still move that block;
remove measured inputs and redundant/BV-exact directions, retain the exact
three-dimensional race fibre, and collect the remainder as \(\chi\).  Populate
or rigorously bound \(\mathsf S_\chi\), then evaluate
\((I-P_\eta)\mathsf S_\chi\).  If the residual vanishes, V26 promotes the V25
three-column left nullspace to the complete first-order block prediction.  If it
is nonzero, append the detected extra selector columns and recompute the rank.
In parallel, populate the V25 terminal-germ panel needed to obtain the actual
three race-selector columns.}}
\end{center}

\section*{Append-only continuation marker: V26}
\addcontentsline{toc}{section}{Append-only continuation marker: V26}
\textbf{End of V26.}  The complete V1--V25 body above is retained verbatim.
V26 proves that the residual race variables are exactly, and nonlinearly, the
complete P1-calibrated fibre of the declared independently exponential race
family.  It then separates this race-family theorem from full ESM selector
completeness and introduces the exact block-specific obstruction
\((I-P_\eta)\mathsf S_\chi\): the three-selector packet is complete for a
physical coefficient block if and only if every additional microscopic
sensitivity lies in its image.  Independent Wilson matching inputs fail this
test automatically and keep the unrestricted EFT nonpredictive.  A locked
race-selector model, by contrast, promotes the V25 three-column conormal to the
complete first-order prediction.  The next gate is to lock one concrete
selector domain and populate the extra-selector incidence together with the
already-designed terminal-germ data.



\clearpage
\section{V27 continuation: P4j upstream visibility gate and selector-domain locking for the curved-Higgs pilot}
\label{sec:v27-continuation}

V26 closed an important domain question: inside the declared independently
exponential race family, the P1-calibrated microscopic fibre is exactly the
three-dimensional residual packet $\eta\in\ker R$.  The next handoff asked for
selector-domain locking on a concrete gravity--matter coefficient block.  There
is, however, one further gate which must be crossed before a selector rank can
be interpreted physically: the chosen coefficient block must be \emph{visible
on the acquisition carrier}.  V24--V25 encoded this as the full-column-rank
hypothesis on a static coefficient synthesis $\mathcal J_D$, but did not prove
that the terminal operation germ actually sees the V18 curvature/Higgs Wilson
block.

The present continuation resolves that issue by separating two tasks:
\begin{center}
\fbox{\parbox{0.90\textwidth}{\centering
\textbf{coefficient visibility on a declared physical source/amplitude panel}
\\[0.25em]$\oplus$\\[0.25em]
\textbf{microscopic selector completeness for the now-visible coefficient block}.}}
\end{center}
For the curved-Higgs anchor, the natural readout is not the unsourced terminal
Dirac/operation germ.  It is a small source-lifted action/amplitude panel: one
mixed metric--Higgs vertex and two independent quadratic metric projectors.
This panel separates the three anchor coefficients at tree level.  The
28-channel terminal-germ compiler remains valid for coefficient blocks which
actually enter that operation carrier, but it is no longer used as a surrogate
for curvature-Wilson visibility.

\subsection{A concrete three-row curved-Higgs anchor}
\label{sec:v27-anchor}

Work on the V18 boundaryless fixed-topology pilot panel and retain the physical
quotient by BV-exact, equation-of-motion, field-redefinition, total-derivative
and finite-scheme directions.  Define the anchor representatives
\begin{equation}
 \mathcal O_\xi=R H^\dagger H,\qquad
 \mathcal O_C=C_{\mu\nu\rho\sigma}C^{\mu\nu\rho\sigma},\qquad
 \mathcal O_R=R^2.
 \label{eq:v27-anchor-ops}
\end{equation}
Whenever their physical classes are nonzero on the selected source panel, let
\begin{equation}
 \mathcal W_0
 :=\operatorname{span}
 \{[\mathcal O_\xi],[\mathcal O_C],[\mathcal O_R]\}
 \subset\mathcal W_{\rm pilot}.
 \label{eq:v27-W0}
\end{equation}
The coefficient vector is
\begin{equation}
 W_0=(\xi,c_C,c_R)^T.
 \label{eq:v27-W0-coords}
\end{equation}
Topological $E_4$ and the total derivative $\Box R$ remain excluded exactly as
in V18.  If a future panel makes one of the three displayed classes redundant,
that class is removed before applying the statements below; no rank is assigned
to a vanished quotient class.

\begin{definition}[Source-lifted anchor readout]
\label{def:v27-anchor-readout}
Choose nonexceptional Euclidean momenta and three normalized physical
functionals:
\begin{enumerate}[label=\textup{(A\arabic*)},leftmargin=3em]
\item $\ell_\xi$: one metric-curvature insertion together with one $H$ and one
$H^\dagger$ insertion;
\item $\ell_C$: the quadratic metric vertex evaluated on a transverse-traceless
polarization;
\item $\ell_R$: the quadratic metric vertex evaluated on a conformal/scalar
polarization.
\end{enumerate}
The normalizations are chosen so that the resulting anchor readout
\begin{equation}
 \mathcal A_0(W_0)
 :=\bigl(\ell_\xi S,\ell_C S,\ell_R S\bigr)^T
 \label{eq:v27-A0}
\end{equation}
has unit tree Jacobian in the coordinates \eqref{eq:v27-W0-coords}.
\end{definition}

\begin{theorem}[The three anchor classes are separated by finite source data]
\label{thm:v27-anchor-visible}
On a flat reference background with vanishing background Higgs field, and with
nonexceptional source momenta, the three functionals of
Definition~\ref{def:v27-anchor-readout} can be chosen so that
\begin{equation}
 \boxed{
 D_{W_0}\mathcal A_0\big|_{\rm tree}=I_3.}
 \label{eq:v27-anchor-identity}
\end{equation}
Consequently the three-row anchor is fully visible on this source-lifted panel.
\end{theorem}

\begin{proof}
The mixed operator $RH^\dagger H$ is the only displayed anchor representative
with two Higgs legs.  A functional derivative with one metric insertion and two
Higgs insertions therefore annihilates $\mathcal O_C$ and $\mathcal O_R$ and is
nonzero on $\mathcal O_\xi$ for a curvature-carrying metric source.  This gives
$\ell_\xi$ after normalization.

For the pure-gravity pair, linearize around flat space.  On a
transverse-traceless metric perturbation, the linearized scalar curvature
vanishes while the linearized Weyl tensor is generically nonzero at
nonexceptional momentum.  Hence the quadratic $R^2$ vertex vanishes on that
polarization whereas the $C^2$ vertex does not; normalize this functional to
obtain $\ell_C$.  On a conformal metric perturbation, the metric is
infinitesimally conformally flat, so the linearized Weyl tensor vanishes while
the linearized scalar curvature is generically nonzero.  The corresponding
quadratic functional therefore isolates $R^2$ and yields $\ell_R$ after
normalization.  The resulting tree matrix is diagonal with nonzero entries and
can be normalized to the identity.
\end{proof}

\begin{corollary}[Finite extension to the complete V18 pilot quotient]
\label{cor:v27-pilot-visible}
Let $\mathcal W_{\rm pilot}$ be any finite physical quotient obtained after the
V18 closure under induced matter/EOM directions.  If its additional independent
classes are represented by local source polynomials on the declared finite
panel, then there exists a finite enlargement of the source functional bank
for which the tree-level visibility map
\begin{equation}
 \mathcal J_{\rm vis}:\mathcal W_{\rm pilot}\longrightarrow\mathcal Y_{\rm vis}
 \label{eq:v27-Jvis}
\end{equation}
has full column rank.
\end{corollary}

\begin{proof}
The physical quotient is finite dimensional.  Independent local classes define
independent polynomial source vertices modulo the already-quotiented null
relations.  Finite-dimensional duality therefore supplies finitely many source
evaluation functionals which separate a chosen basis.  Appending those
functionals to the three explicit anchor rows gives a full-column-rank
visibility matrix.
\end{proof}

\subsection{Why the unsourced terminal germ is not the curvature-Wilson readout}
\label{sec:v27-terminal-visibility}

The V24--V25 terminal compiler is a valid finite reconstruction theorem, but it
contains a hypothesis which is block dependent: the static coefficient map
must be visible on the selected operation core.  For the curvature anchor this
cannot be inferred merely from the existence of the terminal germ.

\begin{proposition}[Flat unsourced reference invisibility of the anchor at tree level]
\label{prop:v27-unsourced-invisible}
Consider an unsourced terminal operation panel evaluated at flat metric and
vanishing Higgs background, with no metric or Higgs source derivatives.  Then
the first tree-level variation of the terminal data with respect to the three
anchor coefficients vanishes:
\begin{equation}
 \boxed{
 D_{W_0}\mathcal O_{\rm term}\big|_{\rm tree}=0.}
 \label{eq:v27-terminal-zero}
\end{equation}
Accordingly the V24 full-column-rank coefficient-incidence hypothesis cannot be
certified for $W_0$ from that unsourced panel.
\end{proposition}

\begin{proof}
At flat metric and zero Higgs background, each anchor interaction begins at
higher field order than the unsourced linearized datum: $C^2$ and $R^2$ begin at
quadratic order in metric perturbations, while $RH^\dagger H$ requires one
curvature-carrying metric perturbation and two Higgs fields.  Their first
variation on the zero-source reference configuration therefore vanishes.
The required nonzero coefficient derivatives appear only after the source
insertions used in Theorem~\ref{thm:v27-anchor-visible}.  Loop-induced terminal
dependence is a separate higher-order effect and cannot be substituted for the
missing tree visibility hypothesis.
\end{proof}

\begin{firewall}[Operation-germ acquisition remains useful]
Proposition~\ref{prop:v27-unsourced-invisible} does not invalidate V24--V25.
Those compilers remain the appropriate route for coefficient blocks which enter
the represented terminal operator directly, or for a source-differentiated
operation channel whose visibility matrix is independently proved full rank.
The correction is only that the gravity--matter pilot must not be inferred from
an unsourced operation carrier which does not see it.
\end{firewall}

\subsection{Visible quotient before predictive quotient}
\label{sec:v27-visible-quotient}

The preceding distinction has a general finite-dimensional form.

\begin{definition}[Visibility quotient]
\label{def:v27-visible-quotient}
For a physical coefficient space $\mathcal W$ and declared observable/acquisition
map $\mathcal O$, define its linearized visibility map
\begin{equation}
 \mathcal J_{\mathcal O}:T\mathcal W\longrightarrow\mathcal Y,
 \qquad
 \mathcal J_{\mathcal O}=D\mathcal O.
 \label{eq:v27-JO}
\end{equation}
The observable identifies only the quotient
\begin{equation}
 \boxed{
 \mathcal W_{\rm vis}:=
 T\mathcal W/\ker\mathcal J_{\mathcal O}.}
 \label{eq:v27-Wvis}
\end{equation}
\end{definition}

\begin{proposition}[Measurement invisibility is not predictivity]
\label{prop:v27-invisible-not-predictive}
Let $\mathsf S:T\Theta_{\rm sel}\to T\mathcal W$ be the true microscopic
matching differential and let
\begin{equation}
 \mathsf S_{\rm obs}=\mathcal J_{\mathcal O}\mathsf S.
 \label{eq:v27-Sobs}
\end{equation}
If $\mathcal J_{\mathcal O}$ is not injective, a null direction of the observed
matrix need not define a relation on $\mathcal W$.  Predictive rank may be read
from observed data without further qualification only after either
\begin{enumerate}[label=\textup{(V\arabic*)},leftmargin=3em]
\item $\mathcal J_{\mathcal O}$ is proved injective on the chosen coefficient
block; or
\item the prediction is explicitly formulated on the visible quotient
$\mathcal W_{\rm vis}$.
\end{enumerate}
\end{proposition}

\begin{proof}
If $0\ne w\in\ker\mathcal J_{\mathcal O}$, then changing the physical
coefficient by $w$ leaves the declared data unchanged to first order.  No
linear-algebra operation on $\mathsf S_{\rm obs}$ can decide whether that
unseen physical direction is selected or freely matchable.  Quotienting the
kernel, or proving it absent, is therefore necessary before an observed
left-null relation is interpreted as a coefficient relation.
\end{proof}

For the anchor readout of Theorem~\ref{thm:v27-anchor-visible},
$\mathcal J_{\mathcal O}=I_3$ at tree order, so this ambiguity is absent.

\subsection{Concrete selector census for the anchor block}
\label{sec:v27-selector-census}

We may now perform P4j on a block whose physical visibility is certified.
After P1 calibration and the V16/V18 quotient, the microscopic directions
known in the present corpus divide as follows:
\begin{enumerate}[label=\textup{(C\arabic*)},leftmargin=3em]
\item measured/calibrated inputs, including the decoded geometry and the
chosen Standard-Model input coordinates --- fixed, hence not selectors;
\item BV-exact, field-redefinition, basis/scheme and compatible-commutant
directions --- quotiented, hence not physical selectors;
\item the exact three-dimensional race fibre $\eta$ of V26;
\item the independent finite Wilson coordinates of the anchor itself, retained
by the current continuum theorem;
\item any further microscopic common-action/source coordinates which a future
same-record selection law allows to move the anchor after P1.  Collect these,
without setting them to zero, into $\chi_*$. 
\end{enumerate}
Thus a current local selector chart has the form
\begin{equation}
 (\eta,W_0,\chi_*).
 \label{eq:v27-selector-chart}
\end{equation}
In the normalized anchor coordinates of
Theorem~\ref{thm:v27-anchor-visible}, the matching differential has the block
form
\begin{equation}
 \boxed{
 \mathsf S_{\rm cur}
 =\bigl[\,\mathsf S_\eta\ \ I_3\ \ \mathsf S_{\chi_*}\,\bigr].}
 \label{eq:v27-current-S}
\end{equation}

\begin{theorem}[Current unrestricted anchor has zero predictive conormal]
\label{thm:v27-anchor-no-prediction}
For the current Renewal matching state, in which the three anchor Wilson
coordinates remain independently matchable,
\begin{equation}
 \boxed{
 \operatorname{rank}\mathsf S_{\rm cur}=3,
 \qquad
 \ker\mathsf S_{\rm cur}^T=\{0\}.}
 \label{eq:v27-anchor-no-prediction}
\end{equation}
This conclusion is independent of the still-unpopulated values of
$\mathsf S_\eta$ and $\mathsf S_{\chi_*}$.
\end{theorem}

\begin{proof}
The middle block in \eqref{eq:v27-current-S} is $I_3$, so the row rank is
three.  The transpose therefore has trivial kernel in the three-dimensional
anchor coefficient space.
\end{proof}

\subsection{Why an arbitrary ``locking section'' would manufacture a result}
\label{sec:v27-section-arbitrariness}

It is tempting to remove the identity block by declaring the anchor Wilson
coordinates to be functions of $\eta$.  That is not a deduction unless the
function is supplied by an independent microscopic law.

\begin{theorem}[Arbitrary-section theorem]
\label{thm:v27-arbitrary-section}
Suppose the current local matching chart contains an open product
$U_\eta\times U_W$ with independent Wilson coordinate $W\in\mathcal W_0$.
For every linear map
\begin{equation}
 A:\mathbb R^3\longrightarrow\mathcal W_0
 \label{eq:v27-A-arbitrary}
\end{equation}
and every base point $(\eta_0,W_0)$, there exists a smooth local section
\begin{equation}
 W=f(\eta)
 \label{eq:v27-section}
\end{equation}
through the base point with
\begin{equation}
 Df(\eta_0)=A.
 \label{eq:v27-Df-A}
\end{equation}
Therefore a Wilson ``locking rule'' chosen only as a section of the already
free matching fibre can manufacture an arbitrary three-column sensitivity
matrix and has no predictive content by itself.
\end{theorem}

\begin{proof}
In local coordinates choose
\[
 f(\eta)=W_0+A(\eta-\eta_0)
\]
inside a sufficiently small neighborhood contained in $U_W$.  Its derivative
at $\eta_0$ is $A$.  Since the unrestricted theory already permits independent
variation of $W$, selecting this graph is an additional model restriction, not
a consequence of the preceding continuum/calibration theorems.
\end{proof}

\begin{corollary}[Zero-Wilson and conformal-$\xi$ boundary values are not locks]
\label{cor:v27-zero-not-lock}
Conditions such as $c_C=c_R=0$ or $\xi=1/6$ imposed directly at $\mu_0$ are
particular choices of section in the current free Wilson fibre.  They become
physical selections only if an independent microscopic incidence law implies
them.
\end{corollary}

\subsection{A genuine lock is an implicit microscopic matching equation}
\label{sec:v27-genuine-lock}

The correct replacement for an arbitrary section is a same-record microscopic
constraint whose Wilson Jacobian is nonsingular.

\begin{definition}[Physical lock equation]
\label{def:v27-lock-equation}
Let $\theta=(\eta,\chi_*)$ denote all non-Wilson microscopic selectors retained
for the anchor.  A \emph{physical lock} is a finite same-record equation
\begin{equation}
 \mathcal F(W_0,\theta)=0,
 \qquad
 \mathcal F:\mathcal W_0\times\Theta\to\mathbb R^3,
 \label{eq:v27-lock-F}
\end{equation}
obtained from an independently populated action/amplitude/operation incidence,
with
\begin{equation}
 \det D_{W_0}\mathcal F\ne0
 \label{eq:v27-lock-nonsing}
\end{equation}
at the calibrated point.
\end{definition}

\begin{theorem}[Implicit selector law]
\label{thm:v27-implicit-lock}
Under Definition~\ref{def:v27-lock-equation}, there is a unique local matching
law $W_0=W_0(\theta)$ and
\begin{equation}
 \boxed{
 D_\theta W_0
 =-\bigl(D_{W_0}\mathcal F\bigr)^{-1}D_\theta\mathcal F.}
 \label{eq:v27-lock-sensitivity}
\end{equation}
In particular, if $\theta=\eta$ is proved complete for the block, the complete
first-order anchor sensitivity is the $3\times3$ matrix
\begin{equation}
 \mathsf S_\eta
 =-\bigl(D_{W_0}\mathcal F\bigr)^{-1}D_\eta\mathcal F.
 \label{eq:v27-Seta-lock}
\end{equation}
\end{theorem}

\begin{proof}
This is the finite-dimensional implicit-function theorem applied at the
calibrated point.  Differentiating $\mathcal F(W_0(\theta),\theta)=0$ gives
\[
 D_{W_0}\mathcal F\,D_\theta W_0+D_\theta\mathcal F=0,
\]
and invertibility of the Wilson block gives
\eqref{eq:v27-lock-sensitivity}.
\end{proof}

The source-lifted anchor of Theorem~\ref{thm:v27-anchor-visible} makes the lock
especially transparent.  Let
\begin{equation}
 \mathcal A_{\rm ren}^{\rm anchor}(W_0)
 \label{eq:v27-Aren-anchor}
\end{equation}
be the three renormalized source amplitudes and let
\begin{equation}
 \mathcal A_{\rm mic}^{\rm anchor}(\theta)
 \label{eq:v27-Amic-anchor}
\end{equation}
be the same three amplitudes computed from an independently populated
microscopic same-record law with no subsequent Wilson refit.  Set
\begin{equation}
 \mathcal F
 :=\mathcal A_{\rm ren}^{\rm anchor}(W_0)
   -\mathcal A_{\rm mic}^{\rm anchor}(\theta).
 \label{eq:v27-F-anchor}
\end{equation}
At tree/reference order,
$D_{W_0}\mathcal A_{\rm ren}^{\rm anchor}=I_3$, so
\begin{equation}
 \boxed{
 D_\theta W_0
 =D_\theta\mathcal A_{\rm mic}^{\rm anchor}
 }
 \label{eq:v27-direct-anchor-incidence}
\end{equation}
at that normalized order.  Higher-loop corrections replace the identity by the
already-defined V16 calibration/matching Jacobian but do not alter the logic.

\subsection{Exact P4j no-extra-selector test after locking}
\label{sec:v27-no-extra-test}

Suppose a genuine lock removes the independent $I_3$ Wilson columns.  The
remaining sensitivity is
\begin{equation}
 \mathsf S_{\rm lock}
 =\bigl[\,\mathsf S_\eta\ \ \mathsf S_{\chi_*}\,\bigr].
 \label{eq:v27-Slock}
\end{equation}
Let
\begin{equation}
 P_\eta=\mathsf S_\eta\mathsf S_\eta^\dagger.
 \label{eq:v27-Peta}
\end{equation}
Then the V26 block-specific completeness obstruction becomes
\begin{equation}
 \boxed{
 \mathcal C_{\rm lock}
 =(I-P_\eta)\mathsf S_{\chi_*}.}
 \label{eq:v27-Clock}
\end{equation}

\begin{theorem}[Locked-anchor selector criterion]
\label{thm:v27-locked-anchor-criterion}
After a genuine physical lock has removed the independent Wilson coordinates,
the three race selectors are complete for the anchor coefficient image if and
only if
\begin{equation}
 \mathcal C_{\rm lock}=0.
 \label{eq:v27-Clock-zero}
\end{equation}
If, in addition, $\chi_*$ is absent, fixed, or image-redundant, then
\begin{equation}
 \boxed{
 \mathcal P_0^*=\ker\mathsf S_\eta^T.}
 \label{eq:v27-anchor-conormal}
\end{equation}
For the three-dimensional anchor specifically:
\begin{equation}
 \dim\mathcal P_0^*=3-\operatorname{rank}\mathsf S_\eta.
 \label{eq:v27-anchor-pred-dim}
\end{equation}
Thus no relation is dimensionally forced: a full-rank $3\times3$ race map gives
zero first-order anchor relation, while rank two, one, or zero gives respectively
one, two, or three independent first-order relations.
\end{theorem}

\begin{proof}
The completeness equivalence is V26 applied after the independent Wilson block
has genuinely been removed from the microscopic domain.  With no transverse
extra-selector image, the full matching image equals
$\operatorname{im}\mathsf S_\eta$.  Rank-nullity on the three-dimensional
coefficient space gives \eqref{eq:v27-anchor-pred-dim}.
\end{proof}

\begin{firewall}[The anchor is a rank test, not a guaranteed prediction]
The three anchor outputs and three complete race selectors have equal
dimension.  V27 therefore does not promise a nontrivial relation in this
minimal block.  Its purpose is to produce the first correctly visible,
correctly locked rank test.  A dimensionally guaranteed conormal requires an
extended visible block of dimension greater than three with the same complete
selector image, exactly as in V24/V26.
\end{firewall}

\subsection{Re-prioritizing the acquisition programme}
\label{sec:v27-acquisition-priority}

For the curvature/Higgs pilot, the preceding theorems change the order of
operations.

\begin{enumerate}[label=\textup{(P\arabic*)},leftmargin=3em]
\item Use the source-lifted anchor amplitudes of
Theorem~\ref{thm:v27-anchor-visible} as the primary Wilson readout.
\item Populate their same-record microscopic dependence on
$\eta$ (and on any declared $\chi_*$) with no independent Wilson refit.
\item Form the lock residual \eqref{eq:v27-F-anchor}; reject the proposed
microscopic model if the three amplitude rows cannot be matched by the anchor
coefficient block plus the already-certified induced closure.
\item Once locked, compute $\mathsf S_\eta$, the no-extra-selector residual
\eqref{eq:v27-Clock}, and the first actual left nullspace.
\item Retain the V25 28-channel terminal-germ acquisition as the primary route
for operation/Dirac coefficient blocks or as an additional cross-check after a
source-differentiated operation carrier is shown to see the curvature block.
\end{enumerate}

This prevents an expensive terminal-channel tomography from being treated as
the unique path to a coefficient block which is already separated directly by
three low-order source amplitudes.

\subsection{Minimal source-lifted acquisition error law}
\label{sec:v27-source-error}

The direct anchor route also has a simpler finite-difference error budget.  Let
$\widehat{\mathcal A}_{\rm mic}(\eta)$ be the three measured/microscopically
computed anchor amplitudes with uniform row error at most $\delta_A$.  For the
central estimator
\begin{equation}
 \widehat{\mathsf S}_{\eta,\alpha}
 :=\frac{
 \widehat{\mathcal A}_{\rm mic}(+\epsilon e_\alpha)
 -\widehat{\mathcal A}_{\rm mic}(-\epsilon e_\alpha)
 }{2\epsilon},
 \label{eq:v27-central-anchor}
\end{equation}
assume a bounded third $\eta$ derivative on the protected chart.  Then
\begin{equation}
 \boxed{
 \|\widehat{\mathsf S}_{\eta,\alpha}-\mathsf S_{\eta,\alpha}\|
 \le C_\eta\epsilon^2+\frac{\delta_A}{\epsilon}.}
 \label{eq:v27-anchor-error}
\end{equation}
Balancing the two terms gives
\begin{equation}
 \epsilon\asymp\delta_A^{1/3},
 \qquad
 \boxed{
 \|\widehat{\mathsf S}_{\eta,\alpha}-\mathsf S_{\eta,\alpha}\|
 =O(\delta_A^{2/3}).}
 \label{eq:v27-anchor-balance}
\end{equation}
This is not a statement about currently available data.  It is the finite
acquisition law for the direct P4 anchor once those same-record amplitudes are
populated.

\subsection{V27 anti-circularity audit}
\label{sec:v27-audit}

\begin{enumerate}[leftmargin=2.4em]
\item \textbf{A coefficient must be visible before its selector rank is read.}
The unsourced terminal operation germ is not assumed to see curvature/Higgs
Wilson coefficients.
\item \textbf{The pilot readout is now explicit.}  The $RH^\dagger H$,
$C^2$ and $R^2$ classes are separated by one mixed Higgs--metric source row and
two independent metric-polarization rows.
\item \textbf{Measurement invisibility is not called a prediction.}  Any
noninjective observable map is quotiented before a left nullspace is interpreted.
\item \textbf{Independent Wilson inputs remain visible.}  Their identity block
proves the current unrestricted anchor is nonpredictive without knowing the
race sensitivity numerically.
\item \textbf{A hand-chosen section is not a microscopic law.}  The
arbitrary-section theorem shows that an arbitrary first derivative can be
manufactured by choosing a graph inside the already free Wilson fibre.
\item \textbf{A genuine lock has a falsifiable equation.}  The Wilson Jacobian
of the same-record lock must be nonsingular and its residual can fail.
\item \textbf{The three-selector anchor does not guarantee a relation.}  A
nontrivial conormal appears only if the measured locked $3\times3$ sensitivity
is rank deficient.
\item \textbf{The 28-channel terminal design is retained but re-scoped.}  It is
not discarded; it is no longer the primary visibility instrument for the
curvature/Higgs anchor unless source-differentiated terminal visibility is
proved.
\end{enumerate}

\section{V27 status ledger}
\label{sec:v27-status}

\subsection*{Proved in V27}
\begin{enumerate}[leftmargin=2.2em]
\item The concrete curved-Higgs anchor
$([RH^\dagger H],[C^2],[R^2])$ is separated on a finite source-lifted panel by
one mixed Higgs--metric functional and transverse-traceless/scalar metric
projectors; the normalized tree visibility Jacobian is $I_3$.
\item On a flat zero-Higgs unsourced terminal panel, the anchor has zero
reference/tree visibility.  Therefore V24's full-column-rank static incidence
cannot be assumed for this block without source lifting or a separate higher-
order proof.
\item For any finite physical coefficient block, predictivity must be evaluated
only after quotienting the kernel of the declared observable/visibility map, or
after proving that map injective.
\item After P1 and the physical quotient, the currently known anchor selector
chart contains the exact race fibre $\eta$, the independent anchor Wilson
coordinates themselves, and any further unresolved same-record selectors
$\chi_*$.  The Wilson identity columns make the current unrestricted anchor
full row rank and nonpredictive.
\item Any smooth local race-to-Wilson derivative can be manufactured by an
arbitrary section of the free Wilson fibre.  Such a section has no predictive
status without an independent microscopic lock equation.
\item A genuine lock with invertible Wilson Jacobian yields the exact
sensitivity formula
$D_\theta W=-(D_W\mathcal F)^{-1}D_\theta\mathcal F$.
For the normalized source-lifted anchor this reduces at tree/reference order to
the derivative of the microscopic amplitude packet.
\item After genuine locking, the V26 no-extra-selector obstruction becomes
$(I-P_\eta)\mathsf S_{\chi_*}$.  If it vanishes, the complete first-order
anchor conormal is $\ker\mathsf S_\eta^T$.
\item The direct source-lifted central-difference acquisition has error
$O(\epsilon^2+\delta_A/\epsilon)$ and balances at
$\epsilon\asymp\delta_A^{1/3}$, giving $O(\delta_A^{2/3})$ sensitivity error.
\end{enumerate}

\subsection*{Inherited}
\begin{enumerate}[leftmargin=2.2em]
\item V16--V18: physical calibration, predictive quotient, RG covariance and
the original curved-Higgs/curvature pilot quotient with independent Wilson
matching freedom.
\item V19--V23: the exact residual race packet, same-record action/operation
incidence compilers and the distinction between scheduler, field-operation and
action-response rows.
\item V24--V25: finite terminal operation-germ inversion and its complete
28-channel data/error certificate for coefficient blocks visible on that
carrier.
\item V26: exact three-dimensional race-fibre completeness and the general
no-extra-selector obstruction $(I-P_\eta)\mathsf S_\chi$.
\end{enumerate}

\subsection*{Conditional}
\begin{enumerate}[leftmargin=2.2em]
\item The source-lifted anchor assumes the three displayed physical classes are
nonredundant on the chosen boundaryless panel; any vanishing class is removed
before rank analysis.
\item Full Renewal-specific selection still requires a populated microscopic
same-record amplitude/action law.  V27 supplies the dual readout and lock
compiler but does not fabricate its values.
\item Additional selectors $\chi_*$ remain unenumerated until a concrete
microscopic action/source model is frozen.  Their absence is not assumed.
\item Loop-level visibility and finite matching corrections must be propagated
through the V16 Jacobian when the anchor is used beyond tree/reference order.
\end{enumerate}

\subsection*{Open beyond V27}
\begin{enumerate}[leftmargin=2.2em]
\item Populate the three source-lifted microscopic anchor amplitudes on the
seven selector settings $\eta=0,\pm\epsilon e_1,\pm\epsilon e_2,\pm\epsilon e_3$
with one common P1 calibration and no independent anchor-Wilson refit.
\item Evaluate the lock residual and its Wilson Jacobian.  Reject or enlarge the
coefficient bank if the microscopic amplitude vector does not lie in the
visible anchor image after the already-certified induced closure.
\item Enumerate every remaining non-race selector $\chi_*$ of the chosen
same-record microscopic model and compute the transverse completeness residual
$(I-P_\eta)\mathsf S_{\chi_*}$.
\item If the lock and completeness gates pass, compute the rank/determinant of
the resulting $3\times3$ race sensitivity.  Only a rank deficiency constitutes
a first-order relation in the minimal anchor.
\item Extend the visible source bank to the full V18 pilot quotient.  If its
physical dimension exceeds three while the selector image remains the same,
the dimension bound of V24/V26 yields guaranteed first-order relations.
\item Keep the V25 terminal-germ acquisition as a parallel operation/Dirac
coefficient programme rather than using it as a surrogate for curvature-Wilson
visibility.
\end{enumerate}

\subsection*{Exact next load-bearing theorem}
\begin{center}
\fbox{\parbox{0.92\textwidth}{
\textbf{P4k: direct source-lifted Wilson locking on the curved-Higgs anchor.}
On one same-history microscopic Renewal ESM model, evaluate the three normalized
source amplitudes dual to $RH^\dagger H$, $C^2$ and $R^2$ at
$\eta=0,\pm\epsilon e_1,\pm\epsilon e_2,\pm\epsilon e_3$, using one common P1
calibration and no independent refit of the three anchor Wilson coefficients.
Form the finite lock equation
$\mathcal F=\mathcal A_{\rm ren}^{\rm anchor}-\mathcal A_{\rm mic}^{\rm anchor}$,
verify the Wilson Jacobian is nonsingular and the held-out/source-consistency
rows close, then compute the three-column race sensitivity with its finite
error budget.  Enumerate any additional selectors $\chi_*$ and evaluate
$(I-P_\eta)\mathsf S_{\chi_*}$.  If that residual vanishes, the left nullspace
of the locked $3\times3$ matrix is the complete first-order prediction of the
minimal anchor; if the matrix is full rank, record the rigorous negative result
and enlarge the visible pilot block before doing further RG phenomenology.}}
\end{center}

\section*{Append-only continuation marker: V27}
\addcontentsline{toc}{section}{Append-only continuation marker: V27}
\textbf{End of V27.}  The complete V1--V26 body above is retained verbatim.
V27 goes upstream of selector rank and closes the missing coefficient-visibility
gate.  The minimal curved-Higgs anchor $([RH^\dagger H],[C^2],[R^2])$ is
separated by an explicit finite source-lifted amplitude panel, while the
unsourced flat terminal germ is proved not to see that anchor at tree/reference
order.  The current unrestricted selector chart therefore contains an explicit
Wilson identity block and remains nonpredictive.  An arbitrary graph inside
that free Wilson fibre can manufacture any desired race sensitivity, so a
physical lock must instead arise from an independently populated same-record
matching equation with nonsingular Wilson Jacobian.  Once such a lock lands,
V26's no-extra-selector obstruction applies directly and the minimal anchor
becomes a genuine $3\times3$ rank test.  The next gate is to populate those
three source-lifted microscopic amplitudes and their extra-selector incidence;
the V25 terminal-germ packet remains a parallel operation/Dirac programme.


% =====================================================================
% Version V28: upstream even-cumulant selector split and populated gravity lock
% =====================================================================
\clearpage
\section{Version V28: upstream even-cumulant selector split, populated curvature lock, and the first branch-specific gravity relation}
\label{sec:v28}

Version V27 deliberately stopped before assigning microscopic values to the
source-lifted curvature/Higgs anchor.  Its next proposed acquisition varied only
the three residual race coordinates $\eta\in\ker R$.  An upstream audit of the
already-developed renewal Einstein-dynamics sector shows that this would target
the wrong microscopic layer for the first curvature-squared lock.  The finite-
window gravity construction contains an independently populated \emph{even
quartic} selector, namely the normalized fourth displacement cumulant.  On a
specified isotropic compound-Poisson channel that selector maps explicitly into
the curvature-squared Wilson coefficients.

This continuation therefore performs two corrections at once.  First, it makes
precise the distinction between the V26 three-dimensional \emph{calibrated race
fibre} and the broader higher-cumulant microscopic selector domain.  Second, it
imports the already-populated fourth-cumulant matching map and projects it onto
the V27 visible $(C^2,R^2)$ quotient.  The outcome is the first nontrivial
branch-specific coefficient relation in the phenomenology lineage.  It is not
promoted to a universal Renewal--ESM prediction: the relation is tied to the
specified isotropic compound-Poisson finite-window branch, to a declared local
matching scheme, and to the absence of additional curvature-squared selectors
within that branch.

\subsection{The Gaussian calibration does not fix the quartic renewal layer}
\label{sec:v28-quartic-selector}

Let a centered isotropic compound-Poisson displacement stream have event rate
$r>0$ and one-step covariance $s^2 I$.  Its second and fourth connected scales
are
\begin{equation}
 D=rs^2,
 \qquad
 Q=rs^4.
 \label{eq:v28-DQ}
\end{equation}
The leading metric/Gaussian calibration depends on $D$, whereas the first
isotropic finite-window correction depends on the normalized quartic quantity
$Q/D^2$.  We keep the renewal reset scale entering the finite-window gravity
matching distinct and denote it by $\lambda_{\rm r}$.

\begin{proposition}[An explicit selector transverse to the P1 Gaussian data]
\label{prop:v28-quartic-transverse}
Fix $D>0$.  For every $r>0$, choose
\begin{equation}
 s^2=\frac{D}{r}.
 \label{eq:v28-s-fixed-D}
\end{equation}
Then all channels in this one-parameter family have the same Gaussian diffusion
scale $D$, but
\begin{equation}
 \frac{Q}{D^2}
 =\frac{rs^4}{D^2}
 =\frac1r.
 \label{eq:v28-q4-variable}
\end{equation}
Consequently the fourth-cumulant selector is not determined by the P1 Gaussian
metric calibration.
\end{proposition}

\begin{proof}
Substitute \eqref{eq:v28-s-fixed-D} into \eqref{eq:v28-DQ}.  The second cumulant
remains $D$, while $Q=r(D/r)^2=D^2/r$.
\end{proof}

\begin{definition}[Even Burnett selector]
\label{def:v28-a4}
On the calibrated finite-window gravity branch define
\begin{equation}
 a_4
 :=G_{\rm ren}\lambda_{\rm r}^{2}\frac{Q}{D^2}.
 \label{eq:v28-a4}
\end{equation}
Once $G_{\rm ren}$, $\lambda_{\rm r}$, and the Gaussian metric scale are fixed,
$a_4$ is equivalent to the remaining isotropic fourth-cumulant coordinate.
\end{definition}

\begin{theorem}[Correction to the scope of V26 selector completeness]
\label{thm:v28-v26-scope}
The V26 statement that $\eta\in\ker R\simeq\mathbb R^3$ exhausts the microscopic
fibre is exact inside the declared fixed-support independently exponential race
family.  It does not imply that $\eta$ exhausts the selector domain of a broader
renewal process whose higher connected displacement cumulants may vary at fixed
first two moments.  The family of Proposition~\ref{prop:v28-quartic-transverse}
provides an explicit transverse microscopic direction $a_4$ invisible to the
Gaussian calibration.
\end{theorem}

\begin{proof}
V26 fixes the race variables $(s,d,\log\rho)$ through the calibration map of that
specific exponential race family and identifies its residual tangent fibre with
$\ker R$.  Proposition~\ref{prop:v28-quartic-transverse} changes the microscopic
jump-rate/step-scale realization while preserving the calibrated second moment,
but changes $Q/D^2$.  This deformation is therefore not a fourth vector inside
the V26 race fibre; it belongs to a larger microscopic model class.
\end{proof}

\begin{firewall}[Race-fibre completeness is not higher-cumulant completeness]
The conclusion of Theorem~\ref{thm:v28-v26-scope} is a scope correction, not a
contradiction.  V26 correctly solved the fibre problem it posed.  V28 shows that
curvature-squared phenomenology lives one cumulant order higher and therefore
requires an enlarged selector census whenever the microscopic model is enlarged
beyond that fixed-support race family.
\end{firewall}

\subsection{Imported microscopic curvature matching on the worked channel}
\label{sec:v28-worked-lock}

The upstream renewal Einstein-dynamics construction has already evaluated the
volume contribution of the specified isotropic compound-Poisson channel at the
first finite-window order.  In its fixed covariant heat-kernel matching
convention, the curvature-squared action is written
\begin{equation}
 S_{4\partial}^{\rm grav}
 =\int d^4x\sqrt{-g}\,
 \left(
 c_1R^2+c_2R_{\mu\nu}R^{\mu\nu}
 +c_3C_{\mu\nu\rho\sigma}C^{\mu\nu\rho\sigma}
 \right),
 \label{eq:v28-stelle-basis}
\end{equation}
and the worked microscopic branch yields
\begin{equation}
 \boxed{
 (c_1,c_2,c_3)
 =a_4\left(\frac54,\frac53,-\frac14\right).}
 \label{eq:v28-worked-triple}
\end{equation}
Equivalently, the Wilson ray is proportional to $(15,20,-3)$.  V28 treats
\eqref{eq:v28-worked-triple} as an inherited populated microscopic matching law,
not as a consequence of V1--V27 alone.

\begin{theorem}[The worked channel supplies a genuine microscopic lock]
\label{thm:v28-worked-lock}
For the specified isotropic compound-Poisson finite-window branch, with all
lower-order calibration data held fixed, the pure curvature-squared matching map
at the displayed order factors through the single even selector $a_4$ as
\begin{equation}
 \mathfrak M_{4\partial}^{\rm grav}(a_4)
 =a_4\left(\frac54,\frac53,-\frac14\right).
 \label{eq:v28-micro-lock}
\end{equation}
Hence its microscopic image in the three-coordinate Stelle basis has rank one.
This lock is physically different from the arbitrary graph construction excluded
by V27: the coefficient ratios are inherited from an independently evaluated
fourth-cumulant response rather than chosen as a section of a free Wilson fibre.
\end{theorem}

\begin{proof}
Equation~\eqref{eq:v28-worked-triple} is the explicit upstream finite-window
calculation.  At fixed calibration, $a_4$ is the only scalar amplitude multiplying
the isotropic fourth-cumulant tensor, so differentiation with respect to $a_4$
gives the displayed rank-one image.  No free $c_i$ coordinate is differentiated
in constructing this map.
\end{proof}

\begin{remark}[What has and has not been locked]
\label{rem:v28-lock-scope}
The theorem locks the \emph{worked-channel contribution} to the pure
curvature-squared block at the declared order.  It does not assert that every
Renewal microscopic model has only one fourth-cumulant selector, nor that matter,
mixed, anisotropic, parity-odd, or independently added local Wilson sectors are
absent in the unrestricted continuum theory.
\end{remark}

\subsection{Projection to the V27 physical curvature quotient}
\label{sec:v28-GB-projection}

V27 uses the visible physical anchor coordinates $c_C$ and $c_R$ multiplying
$C^2$ and $R^2$.  In four dimensions, on a boundaryless fixed-topology panel,
the Euler density
\begin{equation}
 E_4
 =R_{\mu\nu\rho\sigma}R^{\mu\nu\rho\sigma}
 -4R_{\mu\nu}R^{\mu\nu}+R^2
 \label{eq:v28-E4}
\end{equation}
has no local bulk variation.  Combining it with the Weyl decomposition gives
\begin{equation}
 R_{\mu\nu}R^{\mu\nu}
 =\frac12 C_{\mu\nu\rho\sigma}C^{\mu\nu\rho\sigma}
 +\frac13 R^2-\frac12E_4.
 \label{eq:v28-Ricci-GB}
\end{equation}
Therefore
\begin{equation}
 c_1R^2+c_2R_{\mu\nu}R^{\mu\nu}+c_3C^2
 \equiv
 c_RR^2+c_CC^2
 \quad(\mathrm{mod}\ E_4),
 \label{eq:v28-physical-basis}
\end{equation}
with
\begin{equation}
 c_R=c_1+\frac13c_2,
 \qquad
 c_C=c_3+\frac12c_2.
 \label{eq:v28-projection}
\end{equation}

\begin{theorem}[Populated pure-gravity anchor ray]
\label{thm:v28-anchor-ray}
On the worked branch of Theorem~\ref{thm:v28-worked-lock}, the V27 pure-gravity
anchor coordinates are
\begin{equation}
 \boxed{
 c_C=\frac7{12}a_4,
 \qquad
 c_R=\frac{65}{36}a_4.}
 \label{eq:v28-anchor-coeffs}
\end{equation}
Hence the physical matching derivative with respect to $a_4$ is
\begin{equation}
 \mathsf S_{a_4}^{\rm grav}
 =
 \begin{pmatrix}
  7/12\\[2pt]
  65/36
 \end{pmatrix},
 \qquad
 \operatorname{rank}\mathsf S_{a_4}^{\rm grav}=1.
 \label{eq:v28-anchor-sensitivity}
\end{equation}
Its one-dimensional left conormal is generated by
\begin{equation}
 \boxed{
 \ell_{\rm grav}=(65,-21),
 \qquad
 65c_C-21c_R=0.}
 \label{eq:v28-relation}
\end{equation}
Equivalently,
\begin{equation}
 \boxed{\frac{c_C}{c_R}=\frac{21}{65}}
 \label{eq:v28-ratio}
\end{equation}
whenever $a_4\ne0$.
\end{theorem}

\begin{proof}
Substitute \eqref{eq:v28-worked-triple} into
\eqref{eq:v28-projection}:
\begin{equation}
 c_R
 =a_4\left(\frac54+\frac59\right)
 =\frac{65}{36}a_4,
 \qquad
 c_C
 =a_4\left(-\frac14+\frac56\right)
 =\frac7{12}a_4.
 \end{equation}
The sensitivity vector is nonzero, hence rank one.  Finally
\begin{equation}
 65\frac7{12}-21\frac{65}{36}=0,
 \end{equation}
which proves the conormal relation.
\end{proof}

\begin{corollary}[Full V27 anchor with free $\xi$]
\label{cor:v28-full-anchor-free-xi}
Suppose the mixed nonminimal curvature--Higgs coordinate $\xi$ remains an
independent calibrated/matching coordinate while the two pure-gravity
curvature-squared coefficients are locked by the worked channel.  In the V27
ordering
\begin{equation}
 W_0=(\xi,c_C,c_R)^T,
 \end{equation}
the selector sensitivity with coordinates $(\xi,a_4)$ is
\begin{equation}
 \mathsf S_{\xi,a_4}
 =
 \begin{pmatrix}
 1&0\\
 0&7/12\\
 0&65/36
 \end{pmatrix}.
 \label{eq:v28-full-anchor-sens}
\end{equation}
It has rank two and left conormal generated by
\begin{equation}
 \boxed{(0,65,-21).}
 \label{eq:v28-full-conormal}
\end{equation}
Thus leaving $\xi$ free does not by itself destroy the pure-gravity relation.
Any additional selector that moves $(c_C,c_R)$ transversely to
\eqref{eq:v28-anchor-sensitivity}, however, removes it.
\end{corollary}

\begin{firewall}[Branch-specific relation, not universal full-ESM prediction]
Equation~\eqref{eq:v28-relation} is proved for the specified isotropic
compound-Poisson finite-window matching branch, after quotienting the Euler
term in the declared local basis.  The unrestricted V18 continuum theory still
allows independent physical Wilson matching data.  The relation becomes a
prediction only in a microscopic model which \emph{replaces} those independent
$c_C,c_R$ inputs by the populated lock of
Theorem~\ref{thm:v28-worked-lock} and which excludes additional transverse
curvature-squared selectors at the same order.
\end{firewall}

\subsection{Selector completeness on the worked isotropic curvature branch}
\label{sec:v28-branch-completeness}

The upstream finite-window calculation contains a useful special feature: for a
centered isotropic Gaussian jump law inside a compound-Poisson stream, the fully
symmetric $O(3)$-invariant fourth-order tensor is one-dimensional.  Thus the
volume contribution to the pure curvature-squared sector has one scalar quartic
amplitude $Q$, not three independent microscopic fourth-order tensors.

\begin{theorem}[Branch-level selector completeness for the imported volume term]
\label{thm:v28-branch-complete}
Restrict to the following explicitly declared branch:
\begin{enumerate}[label=\textup{(B\arabic*)},leftmargin=3em]
\item the centered isotropic compound-Poisson displacement stream with Gaussian
one-jump law;
\item the fixed covariant finite-window/heat-kernel matching convention of
Theorem~\ref{thm:v28-worked-lock};
\item the pure gravitational volume contribution at the first four-derivative
order;
\item no independently appended curvature-squared Wilson coordinate and no
additional anisotropic, odd, mixed, or matter-induced microscopic selector in
this subblock.
\end{enumerate}
Then the selector image of the visible $(c_C,c_R)$ quotient is exactly the line
spanned by \eqref{eq:v28-anchor-sensitivity}.  Consequently
\begin{equation}
 \boxed{
 \mathcal P_{\rm grav,worked}^*
 =\operatorname{span}\{(65,-21)\}.}
 \label{eq:v28-worked-conormal}
\end{equation}
Within this branch, \eqref{eq:v28-relation} is therefore a complete first-order
and finite-amplitude coefficient relation at the matching order, not merely a
left nullspace of an incompletely sampled selector matrix.
\end{theorem}

\begin{proof}
Under (B1), rotational invariance makes the fully symmetric fourth cumulant a
single scalar multiple of the isotropic tensor.  The imported microscopic
calculation maps that scalar to the Wilson ray
\eqref{eq:v28-worked-triple}.  Assumptions (B2)--(B4) exclude additional
curvature-squared columns in this restricted branch.  The Gauss--Bonnet quotient
then maps the one-dimensional microscopic image to the nonzero vector
\eqref{eq:v28-anchor-sensitivity}.  The annihilator of a nonzero line in a
two-dimensional space is one-dimensional and was computed explicitly in
Theorem~\ref{thm:v28-anchor-ray}.
\end{proof}

\begin{remark}[The V26 completeness obstruction detects the new selector]
\label{rem:v28-v26-obstruction}
If one attempted to analyze the broader worked channel using only the V26 race
packet, the quartic column \eqref{eq:v28-anchor-sensitivity} belongs to the
extra-selector block $\mathsf S_\chi$.  Unless that column already lies in the
race-selector image, the V26 residual
$(I-P_\eta)\mathsf S_\chi$ is nonzero.  Thus the V26 compiler behaves correctly:
it detects precisely the higher-cumulant direction which V28 has now identified.
\end{remark}

\begin{corollary}[Parity-separated product extension]
\label{cor:v28-parity-product}
Consider a product extension in which $\eta$ is the reversal-odd circulation
packet of V20 and $a_4$ is a reversal-even fourth-cumulant coordinate.  If the
pure curvature-squared matching map satisfies
\begin{equation}
 W_{\rm grav}(-\eta,a_4)=W_{\rm grav}(\eta,a_4),
 \label{eq:v28-even-eta}
\end{equation}
then at the reflected base point
\begin{equation}
 D_\eta W_{\rm grav}(0,a_4)=0.
 \label{eq:v28-eta-zero-curv}
\end{equation}
Hence the first curvature-squared selector is the even quartic direction rather
than the linear circulation direction on that symmetric extension.
\end{corollary}

\begin{proof}
Differentiate \eqref{eq:v28-even-eta} with respect to $\eta$ at $\eta=0$.
\end{proof}

\subsection{Scheme and RG transport of the worked relation}
\label{sec:v28-scheme-rg}

The numerical linear form in \eqref{eq:v28-relation} belongs to the fixed
matching coordinates in which \eqref{eq:v28-worked-triple} was evaluated.  V16
already established that predictive relations are transported by finite scheme
changes rather than being discarded.

Let $F_{{\rm phys}\leftarrow{\rm HK}}$ be the finite local change from the
worked heat-kernel matching coordinates to a chosen phenomenological coordinate
system.  Define
\begin{equation}
 R_{\rm HK}(c_C,c_R):=65c_C-21c_R,
 \label{eq:v28-RHK}
\end{equation}
and
\begin{equation}
 \boxed{
 R_{\rm phys}(W_{\rm phys})
 :=R_{\rm HK}\!\left(
 F_{{\rm HK}\leftarrow{\rm phys}}(W_{\rm phys})
 \right).}
 \label{eq:v28-Rphys}
\end{equation}

\begin{proposition}[Scheme-covariant zero set]
\label{prop:v28-scheme-cov}
If the finite scheme map is invertible on the selected coefficient block, then
\begin{equation}
 R_{\rm HK}=0
 \quad\Longleftrightarrow\quad
 R_{\rm phys}=0.
 \label{eq:v28-scheme-zero}
\end{equation}
Thus the branch relation is scheme-covariant as a geometric submanifold even
when its coordinate expression ceases to be the literal linear combination
$65c_C-21c_R$.
\end{proposition}

\begin{proof}
This is composition with an invertible coordinate map.  No claim is made that a
nonlinear or scale-dependent finite scheme change preserves the same numerical
linear equation.
\end{proof}

Let $U(\mu,\mu_0)$ denote the physical RG transport of V16.  At the matching
scale $\mu_0$ define the transported relation
\begin{equation}
 R_\mu(W)
 :=R_{\mu_0}\!\left(U(\mu_0,\mu)W\right).
 \label{eq:v28-Rmu}
\end{equation}

\begin{corollary}[RG transport rather than re-imposition]
\label{cor:v28-rg-relation}
Along the RG flow,
\begin{equation}
 \boxed{
 (\mu\partial_\mu+\beta^A\partial_A)R_\mu=0,}
 \label{eq:v28-CS-relation}
\end{equation}
and $R_\mu=0$ is the RG image of the microscopic matching relation.  The
matching condition is imposed once at $\mu_0$; it is not separately refitted at
each scale.
\end{corollary}

\begin{proof}
Differentiate \eqref{eq:v28-Rmu} using the cocycle property of the V16 RG
transport.
\end{proof}

\subsection{Observable form of the first worked-channel relation}
\label{sec:v28-observable-form}

The V27 source panel was normalized so that its leading visibility matrix on the
pure-gravity anchor is the identity.  Beyond leading order, write the finite
renormalized amplitude map at $\mu$ as
\begin{equation}
 \mathcal A_{\rm grav}(W;\mu)
 =\mathcal A_{\rm known}(\mu)
 +J_A(\mu)
 \begin{pmatrix}c_C\\ c_R\end{pmatrix}
 +O(W^2_{4\partial}),
 \label{eq:v28-amplitude-map}
\end{equation}
where $J_A(\mu_0)=I_2$ at the normalized tree/reference point and
$\mathcal A_{\rm known}$ contains the calibrated lower-order and explicitly
computed loop pieces.

\begin{proposition}[Amplitude-level conormal]
\label{prop:v28-amplitude-conormal}
Whenever $J_A(\mu)$ is invertible on the two visible curvature rows, the worked
matching relation is equivalent, to the retained order, to
\begin{equation}
 \boxed{
 (65,-21)\,J_A(\mu)^{-1}
 \left[
 \mathcal A_{\rm grav}^{\rm obs}(\mu)
 -\mathcal A_{\rm known}(\mu)
 \right]
 =0,}
 \label{eq:v28-observable-relation}
\end{equation}
with the corresponding transported conormal when a nontrivial finite scheme or
RG map is used.  Thus the first microscopic coefficient relation already has a
finite observable projection target; no direct measurement of bare Wilson
coordinates is required.
\end{proposition}

\begin{proof}
Invert \eqref{eq:v28-amplitude-map} on the visible two-dimensional block and
substitute the result into \eqref{eq:v28-relation}.  Scheme/RG transport follows
from Proposition~\ref{prop:v28-scheme-cov} and
Corollary~\ref{cor:v28-rg-relation}.
\end{proof}

\subsection{V28 anti-circularity audit}
\label{sec:v28-audit}

\begin{enumerate}[leftmargin=2.4em]
\item \textbf{The V27 seven-setting $\eta$ experiment is not used to manufacture
curvature-squared coefficients.}  An independently populated even fourth-
cumulant matching map already exists upstream.
\item \textbf{V26 is not contradicted.}  Its three variables are complete for
its declared calibrated race family; $a_4$ appears only after enlarging the
microscopic model to a higher-cumulant family.
\item \textbf{The new relation is not an arbitrary Wilson section.}  The vector
$(5/4,5/3,-1/4)$ is inherited from the evaluated microscopic finite-window
response.
\item \textbf{The Euler direction is quotiented explicitly.}  The relation is
read in the two-dimensional boundaryless local bulk quotient
$(c_C,c_R)$, not in a redundant three-coordinate basis.
\item \textbf{The relation is branch-specific.}  Additional anisotropic,
mixed, matter-induced, or independently appended curvature-squared selectors
may enlarge the matching image and remove the conormal.
\item \textbf{$\xi$ is not silently predicted.}  The current source audit has
not supplied a same-record microscopic lock for $RH^\dagger H$; it therefore
remains an open/free anchor coordinate in V28.
\item \textbf{Scheme covariance is distinguished from literal coefficient
ratios.}  The zero locus transports under finite scheme changes; the simple
numbers $65$ and $21$ belong to the declared worked matching coordinates.
\item \textbf{RG running does not create the relation.}  The matching relation
is fixed at $\mu_0$ and transported by the V16 RG flow.
\end{enumerate}

\section{V28 status ledger}
\label{sec:v28-status}

\subsection*{Proved in V28}
\begin{enumerate}[leftmargin=2.2em]
\item The Gaussian/P1 calibration does not fix the isotropic quartic renewal
selector: at fixed $D=rs^2$, one has $Q/D^2=1/r$.
\item The three-dimensional V26 $\eta$ packet is complete only for its declared
fixed-support race family; a broader higher-cumulant renewal family contains
explicit transverse selector directions.
\item Importing the already-evaluated isotropic compound-Poisson finite-window
branch gives the populated microscopic curvature lock
$(c_1,c_2,c_3)=a_4(5/4,5/3,-1/4)$.
\item After the four-dimensional Euler/Gauss--Bonnet quotient, the V27 visible
pure-gravity anchor is
$(c_C,c_R)=a_4(7/12,65/36)$.
\item The worked pure-gravity matching image therefore has rank one and exact
conormal
$65c_C-21c_R=0$ in the declared matching coordinates.
\item Under the explicit worked-channel branch assumptions excluding additional
curvature-squared selectors, that conormal is complete for the imported volume
contribution at the retained matching order.
\item Leaving the mixed Higgs-curvature coordinate $\xi$ independently free
adds one selector column but does not remove the pure-gravity conormal.
\item The relation transports covariantly under the V16 finite scheme map and
RG flow, and it admits the finite observable projection
\eqref{eq:v28-observable-relation} on the V27 source-lifted amplitude panel.
\end{enumerate}

\subsection*{Inherited}
\begin{enumerate}[leftmargin=2.2em]
\item V16--V18: physical calibration, predictive quotient, finite scheme/RG
transport and the no-prediction theorem for independently free Wilson data.
\item V19--V26: exact race-fibre coordinates, their information geometry,
operation/action incidence compilers, finite acquisition errors, and selector-
completeness criterion.
\item V27: the finite source-lifted visibility panel separating
$RH^\dagger H$, $C^2$, and $R^2$ and the arbitrary-section no-go.
\item Upstream renewal Einstein-dynamics finite-window calculation: for the
specified isotropic compound-Poisson channel, the fourth-cumulant Wilson ray is
$(15,20,-3)$ in the Stelle basis, with a one-dimensional isotropic quartic
selector.
\end{enumerate}

\subsection*{Conditional}
\begin{enumerate}[leftmargin=2.2em]
\item The complete conormal statement
\eqref{eq:v28-worked-conormal} is restricted to the explicitly declared worked
isotropic volume branch and its fixed matching convention.
\item A broader microscopic model may contain additional independent
curvature-squared selector columns.  V26's
$(I-P)\mathsf S_\chi$ test must be rerun before extending the relation beyond
the worked branch.
\item The literal linear numerical relation $65c_C-21c_R=0$ is tied to the
worked matching coordinates.  In a different finite scheme its zero set is
transported by the inverse scheme map.
\item Matter and mixed sectors, including a microscopic lock for $\xi$, are not
closed by the imported pure-gravity fourth-cumulant calculation.
\item Loop and threshold corrections must be included in
$\mathcal A_{\rm known}$ and in the finite scheme/RG transport before confronting
the amplitude-level relation with data.
\end{enumerate}

\subsection*{Open beyond V28}
\begin{enumerate}[leftmargin=2.2em]
\item Convert the worked heat-kernel matching ray explicitly into the chosen
V16 phenomenological scheme at a fixed reference scale $\mu_0$, including the
finite one-loop matching matrix for the two visible curvature operators.
\item Compute the operator-mixing/RG transport of the two-dimensional pure-
gravity conormal so that the relation can be quoted at a standard comparison
scale without re-fitting the microscopic selector.
\item Populate the two V27 source amplitudes (TT and conformal/scalar metric
rows) including the known Einstein/SM loop pieces and form the observable
combination \eqref{eq:v28-observable-relation}.
\item Audit additional fourth-cumulant and mixed microscopic sectors and compute
the V26 transverse-selector residual relative to the worked ray.  This is the
required universality test for extending the relation beyond the isotropic
compound-Poisson branch.
\item Search separately for a genuine microscopic lock of the mixed
$RH^\dagger H$ coordinate $\xi$; do not infer it from the pure-gravity cumulant
ray.
\item Retain the V25 terminal-germ programme for operation/Dirac coefficients;
it is orthogonal to the now-populated pure-gravity finite-window lock.
\end{enumerate}

\subsection*{Exact next load-bearing theorem}
\begin{center}
\fbox{\parbox{0.92\textwidth}{
\textbf{P4l: phenomenological transport of the worked curvature ray.}
Fix one reference scale $\mu_0$ and the V16 physical/standard coefficient
scheme.  Starting from the populated microscopic heat-kernel matching law
$(c_C,c_R)=a_4(7/12,65/36)$, compute the finite native-to-standard matching
matrix on the two visible curvature rows, transport the conormal
$R_{\rm HK}=65c_C-21c_R$ to the physical quotient, and derive its RG-evolved
form $R_\mu=0$.  On the V27 TT and conformal/scalar source amplitudes, subtract
all calibrated lower-order and explicitly computed loop/threshold pieces and
prove the corresponding observable relation
$\ell_\mu^T[\mathcal A_{\rm grav}^{\rm obs}-\mathcal A_{\rm known}]=0$ with a
finite uncertainty budget.  In parallel, enumerate every additional
fourth-cumulant/mixed microscopic selector which can move the same two physical
rows and evaluate the V26 transverse-completeness residual.  If a transverse
column survives, report the worked relation as branch-specific only; if none
survives in a locked enlarged model, promote the transported conormal to the
first complete Renewal gravity phenomenology relation.}}
\end{center}

\section*{Append-only continuation marker: V28}
\addcontentsline{toc}{section}{Append-only continuation marker: V28}
\textbf{End of V28.}  The complete V1--V27 body above is retained verbatim.
V28 goes upstream of the planned seven-setting circulation acquisition and
identifies the microscopic layer which already controls the first
curvature-squared gravity correction.  The calibrated race fibre remains exactly
three-dimensional inside its declared model, but a broader renewal process has
an independent even fourth-cumulant/Burnett selector invisible to the Gaussian
P1 data.  Importing the previously evaluated isotropic compound-Poisson
finite-window branch supplies a genuine microscopic Wilson lock rather than an
arbitrary section: in the Stelle basis
$(c_1,c_2,c_3)=a_4(5/4,5/3,-1/4)$.  After the four-dimensional Euler quotient,
the visible pure-gravity anchor is
$(c_C,c_R)=a_4(7/12,65/36)$ and therefore obeys the exact worked-branch relation
$65c_C-21c_R=0$.  This is the first nonzero branch-specific predictive
conormal in the phenomenology lineage.  It is not yet promoted to a universal
Renewal--ESM prediction: the next gate is to transport it into the physical
scheme/RG frame and to audit every additional microscopic selector capable of
moving the same two curvature rows.



\section{V29: physical-quotient transport of the worked Burnett curvature ray}
\label{sec:v29-main}

V28 produced the first nonzero branch-specific microscopic conormal in a
convenient off-shell curvature basis,
\begin{equation}
 (c_C,c_R)
 =a_4\left(\frac7{12},\frac{65}{36}\right),
 \qquad
 65c_C-21c_R=0.
 \label{eq:v29-v28-ray}
\end{equation}
Before transporting these two coordinates to a standard phenomenological
scheme, however, one must apply the physical quotient already built into V16.
That quotient removes admissible local field redefinitions and equation-of-
motion directions.  In four-dimensional boundaryless Einstein gravity this is
load-bearing at curvature-squared order: the pair $(c_C,c_R)$ is a useful
matching coordinate, but it is not by itself a pair of pure-gravity physical
moduli.  This section performs that quotient explicitly and identifies the
matter-coupled observable descendant of the V28 ray.

\subsection{Upstream correction: the four-derivative pure-gravity pair is
field-redefinition redundant}
\label{sec:v29-field-redef}

Work at the P1 flat reference branch after the cosmological/vacuum term has
been calibrated, and write the lower-weight action as
\begin{equation}
 S_0[g,\Psi]
 =\int d^4x\sqrt g\,
 \left\{\frac{M_g^2}{2}R+\mathcal L_{\rm m}(g,\Psi)\right\},
 \qquad M_g^2=(8\pi G_N)^{-1},
 \label{eq:v29-S0}
\end{equation}
with fixed lower-weight matter coordinates.  The V28 four-derivative local
piece is
\begin{equation}
 S_4[g]
 =\int d^4x\sqrt g\,
 \left(c_C C_{\mu\nu\rho\sigma}C^{\mu\nu\rho\sigma}
       +c_R R^2\right).
 \label{eq:v29-S4}
\end{equation}
In four dimensions,
\begin{equation}
 C^2=E_4+2R_{\mu\nu}R^{\mu\nu}-\frac23R^2,
 \label{eq:v29-C2-Euler}
\end{equation}
where $E_4$ is the Euler density.

\begin{proposition}[Explicit curvature-squared field redefinition]
\label{prop:v29-field-redefinition}
To first order in the four-derivative Wilson coefficients, the local metric
redefinition
\begin{equation}
 \boxed{
 \delta g_{\mu\nu}
 =\frac{2}{M_g^2}
 \left[
   2c_C R_{\mu\nu}
   -\left(c_R+\frac13c_C\right)g_{\mu\nu}R
 \right]}
 \label{eq:v29-dg}
\end{equation}
removes every non-topological pure-gravity term in \eqref{eq:v29-S4}.  More
precisely,
\begin{equation}
 S_0[g+\delta g,\Psi]+S_4[g+\delta g]
 =S_0[g,\Psi]
 +c_C\int d^4x\sqrt g\,E_4
 +\Delta S_{\rm m}
 +O(c^2),
 \label{eq:v29-field-redef-result}
\end{equation}
where $\Delta S_{\rm m}$ is the induced matter-sector term computed below.
Thus, on a boundaryless four-dimensional pure-gravity on-shell panel, the
local curvature-squared pair carries no independent bulk observable at this
order.
\end{proposition}

\begin{proof}
For a covariant metric variation,
\begin{equation}
 \delta S_{\rm EH}
 =-\frac{M_g^2}{2}\int d^4x\sqrt g\,
 G^{\mu\nu}\delta g_{\mu\nu}
 \label{eq:v29-dSEH}
\end{equation}
up to the usual boundary term.  For the more general variation
\begin{equation}
 \delta g_{\mu\nu}
 =\frac{2}{M_g^2}
 \left(aR_{\mu\nu}+bg_{\mu\nu}R\right),
 \label{eq:v29-dg-ab}
\end{equation}
one obtains in $d=4$
\begin{equation}
 \delta\mathcal L_{\rm EH}
 =-aR_{\mu\nu}R^{\mu\nu}
  +\left(\frac a2+b\right)R^2.
 \label{eq:v29-EH-shift}
\end{equation}
Using \eqref{eq:v29-C2-Euler}, the non-topological part of
\eqref{eq:v29-S4} is
\begin{equation}
 2c_C R_{\mu\nu}R^{\mu\nu}
 +\left(c_R-\frac23c_C\right)R^2.
 \label{eq:v29-nontop}
\end{equation}
The choice
\begin{equation}
 a=2c_C,
 \qquad
 b=-c_R-\frac13c_C
 \label{eq:v29-ab-choice}
\end{equation}
cancels \eqref{eq:v29-nontop} exactly to first order in the EFT
coefficients, leaving only $c_C E_4$ in the pure metric sector.
\end{proof}

\begin{firewall}[What V28 did and did not predict]
\label{fire:v29-v28-scope}
Equation~\eqref{eq:v29-v28-ray} remains an exact statement about the worked
microscopic matching branch in its declared off-shell curvature coordinates.
Proposition~\ref{prop:v29-field-redefinition} shows that it is \emph{not}, by
itself, a nontrivial pure-gravity element of the V16 physical quotient.  The
V27 TT/scalar metric rows therefore remain valuable matching and regulator
checks, but they cannot be advertised as a field-redefinition-invariant
four-dimensional vacuum-gravity observable.  The physical descendant is
obtained only after the induced matter/contact terms are retained.
\end{firewall}

\subsection{Landing the V28 ray in the coupled Einstein--SM physical quotient}
\label{sec:v29-matter-contact}

Define the Hilbert stress tensor of the fixed lower-weight matter action by
\begin{equation}
 \delta S_{\rm m}
 =\frac12\int d^4x\sqrt g\,
 T^{\mu\nu}\delta g_{\mu\nu}.
 \label{eq:v29-Tdef}
\end{equation}
The field redefinition \eqref{eq:v29-dg} then gives
\begin{equation}
 \boxed{
 \Delta\mathcal L_{\rm m}
 =\frac1{M_g^2}
 \left[
 2c_C T^{\mu\nu}R_{\mu\nu}
 -\left(c_R+\frac13c_C\right)TR
 \right].}
 \label{eq:v29-matter-shift}
\end{equation}
At the same EFT order, the lower-weight Einstein equation is
\begin{equation}
 M_g^2G_{\mu\nu}=T_{\mu\nu},
 \label{eq:v29-Einstein-EOM}
\end{equation}
so that, modulo the V16 equation-of-motion quotient,
\begin{equation}
 R_{\mu\nu}
 \simeq\frac1{M_g^2}
 \left(T_{\mu\nu}-\frac12g_{\mu\nu}T\right),
 \qquad
 R\simeq-\frac{T}{M_g^2}.
 \label{eq:v29-R-vs-T}
\end{equation}
Here and below $\simeq$ denotes equality in the fixed-order physical quotient,
up to higher-order terms, total derivatives and the already-declared local
redundancies.

\begin{theorem}[Stress-contact representative of the curvature-squared class]
\label{thm:v29-stress-contact}
At fixed lower-weight matter coordinates, the V28 curvature pair has the
physical-quotient representative
\begin{equation}
 \boxed{
 \Delta\mathcal L_{\rm phys}
 \simeq
 \kappa_1\,T_{\mu\nu}T^{\mu\nu}
 +\kappa_2\,T^2,}
 \label{eq:v29-contact-action}
\end{equation}
with
\begin{equation}
 \boxed{
 \kappa_1=\frac{2c_C}{M_g^4},
 \qquad
 \kappa_2=\frac{c_R-\frac23c_C}{M_g^4}.}
 \label{eq:v29-kappa-map}
\end{equation}
For the populated V28 Burnett ray,
\begin{equation}
 \boxed{
 (\kappa_1,\kappa_2)
 =\frac{a_4}{12M_g^4}(14,17).}
 \label{eq:v29-contact-ray}
\end{equation}
Consequently its branch-specific physical conormal is
\begin{equation}
 \boxed{
 17\kappa_1-14\kappa_2=0.}
 \label{eq:v29-contact-relation}
\end{equation}
\end{theorem}

\begin{proof}
Insert \eqref{eq:v29-R-vs-T} into \eqref{eq:v29-matter-shift}.  One finds
\begin{align}
 \Delta\mathcal L_{\rm m}
 &\simeq\frac1{M_g^4}
 \left[
 2c_C\left(T_{\mu\nu}T^{\mu\nu}-\frac12T^2\right)
 +\left(c_R+\frac13c_C\right)T^2
 \right]
 \notag\\
 &=\frac1{M_g^4}
 \left[
 2c_C T_{\mu\nu}T^{\mu\nu}
 +\left(c_R-\frac23c_C\right)T^2
 \right],
 \label{eq:v29-contact-derivation}
\end{align}
which proves \eqref{eq:v29-kappa-map}.  Substituting
$c_C=7a_4/12$ and $c_R=65a_4/36$ gives
\begin{equation}
 \kappa_1=\frac{7a_4}{6M_g^4},
 \qquad
 \kappa_2=\frac{17a_4}{12M_g^4},
 \label{eq:v29-kappa-values}
\end{equation}
namely \eqref{eq:v29-contact-ray}; its left null covector is $(17,-14)$.
\end{proof}

\begin{firewall}[Lower-weight matter conventions]
\label{fire:v29-stress-convention}
The stress tensor in Theorem~\ref{thm:v29-stress-contact} belongs to the fixed
P1 lower-weight Einstein--SM point.  In particular, if the nonminimal Higgs
coordinate $\xi$ is allowed to vary independently, its improvement of the
stress tensor and its own mixed curvature operators must be included as extra
selector/basis rows.  Equation~\eqref{eq:v29-contact-relation} is therefore not
used to silently lock $\xi$.
\end{firewall}

\subsection{Finite native-to-standard transport: tangent first, coordinates second}
\label{sec:v29-standard-transport}

The current corpus proves existence of a finite native-to-standard matching map
but does not yet populate its numerical gravity/matter contact matrix.  The
right object to transport is therefore the one-dimensional microscopic tangent,
not the literal pair of numbers appearing in one off-shell basis.

Let $\mathbf C_{\rm nat}$ be a finite physical contact/operator bank containing
\eqref{eq:v29-contact-action}, and let
\begin{equation}
 \mathbf C_{\rm std}
 =F_{\mu_0}(\mathbf C_{\rm nat})
 \label{eq:v29-standard-map}
\end{equation}
be the V16 finite standard-scheme map at the reference scale.  Write the worked
matching locus locally as
\begin{equation}
 \mathbf C_{\rm nat}(a_4)
 =\mathbf C_{\rm nat}^{(0)}
 +a_4 v_{\rm nat}+O(a_4^2),
 \label{eq:v29-native-tangent}
\end{equation}
where the contact part of $v_{\rm nat}$ is proportional to $(14,17)$.

\begin{proposition}[Scheme transport of the microscopic tangent and conormal]
\label{prop:v29-scheme-tangent}
Let
\begin{equation}
 J_F:=D F_{\mu_0}|_{\mathbf C_{\rm nat}^{(0)}}.
 \label{eq:v29-JF}
\end{equation}
Then
\begin{equation}
 \boxed{v_{\rm std}=J_Fv_{\rm nat}.}
 \label{eq:v29-vstd}
\end{equation}
Every native predictive covector $\ell_{\rm nat}$ satisfying
$\ell_{\rm nat}^Tv_{\rm nat}=0$ transports as
\begin{equation}
 \boxed{
 \ell_{\rm std}=J_F^{-T}\ell_{\rm nat}}
 \label{eq:v29-ellstd}
\end{equation}
whenever $J_F$ is invertible on the retained physical bank.  In a closed
2-by-2 contact block this is the complete standard-scheme conversion of the
relation.  If the standard scheme mixes the two contacts into additional
operators, the full mixing-closed bank must be used instead.
\end{proposition}

\begin{proof}
Differentiate \eqref{eq:v29-standard-map} with respect to $a_4$ at the
reference point.  The conormal transformation is the inverse-transpose tangent
map already proved abstractly in V16.
\end{proof}

\begin{corollary}[No numerical \texorpdfstring{$\overline{\rm MS}$}{MSbar}
ratio is presently licensed]
\label{cor:v29-no-msbar-number}
The current source-complete continuum corpus does not populate the numerical
matrix $J_F$ for the stress-contact/curvature sector.  Therefore V29 does not
replace \eqref{eq:v29-contact-relation} by an invented numerical
$\overline{\rm MS}$ relation.  The finite conversion problem has been reduced
to the explicit Jacobian \eqref{eq:v29-JF} on a common physical amplitude
panel.
\end{corollary}

\subsection{RG transport of a one-selector microscopic ray}
\label{sec:v29-RG-ray}

Let $\mathbf C(t,a_4)$ denote the chosen mixing-closed physical Wilson vector,
where $t=\log(\mu/\mu_0)$, and let
\begin{equation}
 \dot{\mathbf C}=\beta_C(\mathbf C,g(t)).
 \label{eq:v29-C-RG}
\end{equation}
Separate the $a_4$-independent reference trajectory and microscopic tangent by
\begin{equation}
 \mathbf C(t,a_4)
 =\mathbf C^{(0)}(t)+a_4v(t)+O(a_4^2).
 \label{eq:v29-RG-linearization}
\end{equation}

\begin{theorem}[Variational RG equation for the worked microscopic tangent]
\label{thm:v29-tangent-RG}
The microscopic tangent obeys
\begin{equation}
 \boxed{
 \dot v(t)=B(t)v(t),
 \qquad
 B(t)=\left.\partial_{\mathbf C}\beta_C\right|_{\mathbf C^{(0)}(t)}.}
 \label{eq:v29-v-RG}
\end{equation}
A transported conormal covector may be evolved by the adjoint equation
\begin{equation}
 \boxed{
 \dot\ell(t)=-B(t)^T\ell(t).}
 \label{eq:v29-ell-RG}
\end{equation}
If $\ell(0)^Tv(0)=0$, then
\begin{equation}
 \ell(t)^Tv(t)=0
 \label{eq:v29-orth-preserve}
\end{equation}
for every scale on the common RG interval.
\end{theorem}

\begin{proof}
Differentiate \eqref{eq:v29-C-RG} with respect to $a_4$ at $a_4=0$ to obtain
\eqref{eq:v29-v-RG}.  Then
\[
 \frac{d}{dt}(\ell^Tv)
 =(-B^T\ell)^Tv+\ell^TBv=0.
\]
\end{proof}

\begin{firewall}[Why a bare 2-by-2 curvature beta table is not the next
physical calculation]
\label{fire:v29-no-2x2}
After the V16 field-redefinition quotient the worked branch lands in
matter/contact operators, not in an autonomous pure-gravity $(c_C,c_R)$
physical block.  Standard-Model and gravitational loops may mix those contacts
with a larger dimension-eight bank.  A 2-by-2 RG table is sufficient only if
that two-dimensional contact subspace is independently proved invariant under
$B(t)$.  Otherwise the correct P2/P3 object is the smallest finite
mixing-closed bank containing $v(0)$.
\end{firewall}

\subsection{Scheme-free observable transport}
\label{sec:v29-observable}

The cleanest phenomenological statement avoids assigning physical meaning to
intermediate Wilson coordinates.  Let a finite on-shell matter amplitude/source
panel obey, to first order in the worked selector,
\begin{equation}
 \mathcal A(\mu,a_4)
 =\mathcal A^{(0)}(\mu)
 +a_4 h(\mu)+O(a_4^2),
 \qquad
 h(\mu)=J_{\rm phys}(\mu)v(\mu),
 \label{eq:v29-amplitude-ray}
\end{equation}
where $\mathcal A^{(0)}$ contains every calibrated lower-order, loop and
threshold contribution independent of the microscopic Burnett selector.

\begin{proposition}[Observable conormal of a one-selector amplitude ray]
\label{prop:v29-observable-conormal}
For any amplitude covector $q(\mu)$ satisfying
\begin{equation}
 q(\mu)^Th(\mu)=0,
 \label{eq:v29-q-ann}
\end{equation}
the worked branch obeys
\begin{equation}
 \boxed{
 q(\mu)^T
 \left[\mathcal A^{\rm obs}(\mu)-\mathcal A^{(0)}(\mu)\right]
 =O(a_4^2).}
 \label{eq:v29-observable-relation}
\end{equation}
For a two-amplitude panel with
$h=(h_1,h_2)^T\neq0$, one may choose
\begin{equation}
 \boxed{q=(h_2,-h_1)^T.}
 \label{eq:v29-q-two}
\end{equation}
At the reference point in a contact-normalized panel with
$J_{\rm phys}=I_2$, the covector is proportional to $(17,-14)$.
\end{proposition}

\begin{proof}
Equation~\eqref{eq:v29-amplitude-ray} gives
$\mathcal A-\mathcal A^{(0)}=a_4h+O(a_4^2)$.  Contract with any covector
annihilating $h$.
\end{proof}

\begin{corollary}[Observable relation is independent of the intermediate
finite scheme]
\label{cor:v29-observable-scheme}
If two coefficient schemes are related by an invertible finite matching map and
the same physical amplitude panel is used, the vector $h(\mu)$ in
\eqref{eq:v29-amplitude-ray} is unchanged.  Thus
\eqref{eq:v29-observable-relation} is the preferred final representation of the
worked phenomenological relation.
\end{corollary}

\subsection{Selector-completeness audit on the physical contact bank}
\label{sec:v29-selector-audit}

Let $v$ be the worked Burnett tangent in the physical mixing-closed bank and
let $S_\chi$ collect every additional microscopic fourth-cumulant, anisotropic,
mixed, matter or independently retained selector which can move the same bank.
Define
\begin{equation}
 P_v:=\frac{vv^T}{v^Tv}
 \label{eq:v29-Pv}
\end{equation}
for one nonzero selector column and
\begin{equation}
 \boxed{
 \mathcal C_{\rm Burnett}
 :=(I-P_v)S_\chi.}
 \label{eq:v29-Burnett-completeness}
\end{equation}

\begin{proposition}[Physical completeness criterion for the worked ray]
\label{prop:v29-completeness}
The one-selector worked relation is complete on the chosen physical bank if and
only if
\begin{equation}
 \boxed{\mathcal C_{\rm Burnett}=0.}
 \label{eq:v29-completeness-zero}
\end{equation}
Any nonzero column of \eqref{eq:v29-Burnett-completeness} is a direct
certificate that the branch-specific conormal must not be promoted to a
universal Renewal--ESM prediction.
\end{proposition}

\begin{proof}
This is the rank-one specialization of the V26 selector-completeness theorem:
all additional selector images must lie in the same one-dimensional tangent
image if they are not to enlarge the matching locus.
\end{proof}

\subsection{V29 anti-circularity audit}
\label{sec:v29-audit}

\begin{enumerate}[leftmargin=2.4em]
\item \textbf{V28 is retained, not erased.}  Its ratio is an exact worked-branch
statement in the declared heat-kernel/off-shell curvature coordinates.
\item \textbf{The V16 physical quotient is now applied before phenomenological
promotion.}  In four-dimensional boundaryless pure gravity, the non-topological
curvature-squared pair is removed by \eqref{eq:v29-dg}.
\item \textbf{The physical descendant is matter-coupled.}  The worked ray lands
in the stress-contact direction $(\kappa_1,\kappa_2)\propto(14,17)$ and obeys
$17\kappa_1-14\kappa_2=0$ at the reference point.
\item \textbf{No numerical standard-scheme matrix is fabricated.}  The current
corpus proves its existence but does not provide the required finite contact
matching amplitudes.
\item \textbf{No closed 2-by-2 RG block is assumed.}  Operator mixing is followed
in the smallest physical bank proved closed under the anomalous-dimension
matrix.
\item \textbf{The final relation is formulated at amplitude level.}  Known
selector-independent loops, thresholds and lower-weight pieces are subtracted
before taking the observable conormal.
\item \textbf{The relation remains branch-specific.}  Additional microscopic
selectors are tested by \eqref{eq:v29-Burnett-completeness} on the physical
contact bank, not in the redundant off-shell curvature coordinates.
\end{enumerate}

\section{V29 status ledger}
\label{sec:v29-status}

\subsection*{Proved in V29}
\begin{enumerate}[leftmargin=2.2em]
\item The V28 curvature coordinates $(c_C,c_R)$ are not independent
four-dimensional boundaryless pure-gravity physical moduli at the retained
curvature-squared order.  The explicit local redefinition
\eqref{eq:v29-dg} removes their non-topological pure-metric part.
\item The same redefinition transfers the worked microscopic information into
coupled matter curvature/stress operators.  After the leading Einstein EOM,
the physical representative is
$\kappa_1T_{\mu\nu}T^{\mu\nu}+\kappa_2T^2$ with
\eqref{eq:v29-kappa-map}.
\item The populated V28 Burnett ray maps to
$(\kappa_1,\kappa_2)=a_4(14,17)/(12M_g^4)$, giving the branch-specific physical
contact relation $17\kappa_1-14\kappa_2=0$ at the reference point.
\item Finite scheme conversion acts on the microscopic tangent by the Jacobian
$J_F$ and on conormals by $J_F^{-T}$.  No literal coefficient ratio is scheme
invariant, but the matching locus is.
\item The RG evolution of the one-selector tangent is the variational equation
$\dot v=Bv$; conormals evolve by the adjoint equation $\dot\ell=-B^T\ell$.
\item A closed 2-by-2 curvature/contact RG table is valid only after the
corresponding two-dimensional physical operator span is proved invariant under
mixing.  Otherwise the mixing-closed bank must be enlarged.
\item The scheme-free phenomenological target is the amplitude ray
\eqref{eq:v29-amplitude-ray} and its observable conormal
\eqref{eq:v29-observable-relation}.
\item The V26 selector-completeness test specializes to the worked physical
Burnett ray as \eqref{eq:v29-Burnett-completeness}.
\end{enumerate}

\subsection*{Inherited}
\begin{enumerate}[leftmargin=2.2em]
\item V16: the phenomenological quotient explicitly divides by admissible local
field redefinitions, BV-exact directions, finite scheme transformations and
operator/source basis changes.
\item V17--V27: calibrated RG/matching infrastructure, selector-incidence and
completeness compilers, source visibility tests and the separation of physical
input directions from genuine microscopic selectors.
\item V28: the populated isotropic compound-Poisson Burnett matching ray
$(c_C,c_R)=a_4(7/12,65/36)$ in the declared heat-kernel/off-shell curvature
basis.
\end{enumerate}

\subsection*{Conditional}
\begin{enumerate}[leftmargin=2.2em]
\item The stress-contact representative assumes the flat/P1 lower-weight
Einstein equation with calibrated vacuum energy and fixed lower-weight matter
coordinates.  Nonzero background curvature or a varying Higgs improvement
coordinate requires the corresponding enlarged mixed bank.
\item The relation $17\kappa_1-14\kappa_2=0$ is complete only inside the worked
microscopic branch until all additional fourth-cumulant/mixed selectors have
passed \eqref{eq:v29-completeness-zero}.
\item The standard-scheme conormal is numerical only after the common physical
matching amplitudes needed to populate $J_F$ are computed.
\item The RG-evolved two-amplitude relation is numerical only after the
mixing-closed anomalous-dimension bank and physical amplitude Jacobian are
populated.
\end{enumerate}

\subsection*{Open beyond V29}
\begin{enumerate}[leftmargin=2.2em]
\item Choose the smallest field-redefinition-invariant Einstein--SM contact
bank containing $T_{\mu\nu}T^{\mu\nu}$ and $T^2$ after matter EOM and source
identities, and prove that it is closed under the retained one-loop mixing.
\item Populate the finite native-to-standard matching Jacobian $J_F$ on a
common on-shell/source amplitude panel.  This replaces the obsolete target of
a bare pure-gravity 2-by-2 curvature conversion table.
\item Compute the variational RG matrix $B(t)$ on that mixing-closed physical
bank and transport the microscopic tangent $v(t)$ and its conormal.
\item Choose a two-or-more amplitude panel with nonzero visibility of the
stress-contact ray, subtract calibrated lower-order/loop/threshold pieces, and
evaluate the observable relation \eqref{eq:v29-observable-relation} with an
uncertainty budget.
\item Enumerate anisotropic fourth cumulants, higher even moments, mixed
matter-renewal selectors and any remaining independent contact inputs; evaluate
\eqref{eq:v29-Burnett-completeness} before any universal promotion.
\item Continue the separate search for a genuine microscopic lock of the
$RH^\dagger H$ coefficient $\xi$.
\end{enumerate}

\subsection*{Exact next load-bearing theorem}
\begin{center}
\fbox{\parbox{0.92\textwidth}{
\textbf{P4m: physical contact-bank closure and phenomenological transport.}
Starting from the populated V29 stress-contact ray
$(\kappa_1,\kappa_2)\propto(14,17)$, construct the minimal
field-redefinition/EOM-reduced Einstein--SM operator bank which contains its
on-shell matter amplitudes and is closed under the retained one-loop mixing.
Populate one common native-to-standard physical matching Jacobian on that
bank; solve the variational RG equation for the microscopic tangent; and
transport the conormal by the adjoint flow.  On an explicit on-shell amplitude
panel, subtract all selector-independent lower-order, loop and threshold pieces
and compute the scheme-free residual relation
$q(\mu)^T[\mathcal A^{\rm obs}-\mathcal A^{(0)}]=0$.  In parallel, enumerate
all additional microscopic selectors that can move the same physical bank and
evaluate the rank-one V26 completeness residual
$(I-P_v)S_\chi$.  A nonzero transverse column keeps the result branch-specific;
only a vanishing locked residual permits promotion to a complete Renewal--ESM
phenomenological relation.}}
\end{center}

\section*{Append-only continuation marker: V29}
\addcontentsline{toc}{section}{Append-only continuation marker: V29}
\textbf{End of V29.}  The complete V1--V28 body above is retained verbatim.
V29 goes upstream of the planned 2-by-2 curvature scheme conversion and applies
the V16 physical quotient first.  In four-dimensional boundaryless Einstein
gravity the V28 pair $(c_C,c_R)$ is an off-shell matching coordinate rather
than an independent pure-gravity physical modulus: its non-topological part is
removed by an explicit local metric redefinition.  In the coupled Einstein--SM
theory the same microscopic information lands instead in the physical
stress-contact ray
$(\kappa_1,\kappa_2)=a_4(14,17)/(12M_g^4)$ and therefore obeys the
branch-specific contact conormal $17\kappa_1-14\kappa_2=0$ at the reference
point.  Finite standard-scheme matching transports the tangent by its Jacobian,
RG evolution transports it by the variational beta matrix, and the final
scheme-free statement is an on-shell amplitude conormal after subtraction of
all selector-independent known pieces.  The next gate is therefore a
mixing-closed physical Einstein--SM contact bank and its populated standard
matching/RG/amplitude transport, together with the rank-one selector-
completeness audit on that physical bank.



\section{V30: reference-scale on-shell visibility of the physical stress-contact ray}
\label{sec:v30}

V29 moved the populated Burnett curvature ray through the four-dimensional
field-redefinition/EOM quotient and obtained the physical matter-contact
representative
\begin{equation}
 \Delta\mathcal L_{\rm cont}
 =\kappa_1 T_{\mu\nu}T^{\mu\nu}+\kappa_2 T^2,
 \qquad
 (\kappa_1,\kappa_2)
 =\kappa_0(14,17),
 \qquad
 \kappa_0:=\frac{a_4}{12M_g^4}.
 \label{eq:v30-contact-ray}
\end{equation}
The V29 next theorem proposed constructing a one-loop mixing-closed
Einstein--SM contact bank before writing an observable.  V30 goes one step
upstream.  A full anomalous-dimension bank is required to transport the ray
between separated scales, but it is \emph{not} required to formulate the first
scheme-independent benchmark at the reference scale $\mu_0$.  What is required
first is an on-shell amplitude panel whose physical Jacobian actually separates
the two stress-contact coordinates.

The section therefore does four things.  First, it proves that a panel made
only from classically traceless external stress tensors is rank deficient and
cannot test the V29 contact relation.  Second, it constructs an explicit
massive-scalar/Higgs four-point panel and derives its complete tree-level
contact amplitude.  Third, it converts the V29 ray into two finite
same-amplitude conormals at three physical kinematic points.  Fourth, it
separates this reference-scale observable theorem from the still-open one-loop
mixing and microscopic selector-completeness problems.

\subsection{Reference conventions and the physical contact bank}
\label{sec:v30-conventions}

Fix the P1 reference scale $\mu_0$, the lower-weight matter masses and residues,
and one stress-tensor improvement convention.  All amplitudes below are
amputated on-shell connected amplitudes in that fixed physical normalization.
Selector-independent Einstein--SM loops, thresholds and lower-weight pieces are
collected into $\mathcal M^{(0)}$ and subtracted before the Burnett relation is
tested.

For the explicit scalar derivation, use the canonical minimally coupled real
scalar representative
\begin{equation}
 \mathcal L_h
 =\frac12\partial_\mu h\,\partial^\mu h-\frac12m_h^2h^2,
 \label{eq:v30-scalar-L}
\end{equation}
with bilinear stress-tensor matrix element
\begin{equation}
 B_{\mu\nu}(p_i,p_j)
 :=p_{i\mu}p_{j\nu}+p_{j\mu}p_{i\nu}
   -\eta_{\mu\nu}\bigl(p_i\!\cdot\!p_j-m_h^2\bigr),
 \label{eq:v30-Bij}
\end{equation}
and trace
\begin{equation}
 b(p_i,p_j)
 :=\eta^{\mu\nu}B_{\mu\nu}(p_i,p_j)
 =4m_h^2-2p_i\!\cdot\!p_j.
 \label{eq:v30-bij}
\end{equation}
An overall common Feynman-rule normalization is irrelevant to every rank and
conormal statement below; it is stripped from the definition of the normalized
contact amplitude $\widehat{\mathcal M}_{4h}$.

\begin{firewall}[Improvement coordinate]
\label{fw:v30-improvement}
The explicit numbers in the scalar amplitude below use the canonical stress
tensor in \eqref{eq:v30-Bij}.  If a nonminimal Higgs improvement coordinate
$\xi$ is retained, it must be fixed as P1 data or included in an enlarged
source/contact bank before using the numerical scalar projector.  One must not
vary $\xi$ while simultaneously treating the stress tensor in
\eqref{eq:v30-contact-ray} as fixed.  The scheme-free amplitude-ray statement
survives a consistent change of stress-tensor convention, but the intermediate
Jacobian must then be recomputed.
\end{firewall}

\subsection{A traceless-panel no-go}
\label{sec:v30-traceless}

\begin{proposition}[Purely traceless external panels cannot separate the
V29 contact pair]
\label{prop:v30-traceless-no-go}
Let an on-shell amplitude panel be such that every external-state stress-tensor
matrix element entering the retained contact insertion is traceless.  Then the
linearized physical amplitude map for
$(\kappa_1,\kappa_2)$ has rank at most one:
\begin{equation}
 D_{(\kappa_1,\kappa_2)}\mathcal A_{\rm tr}
 =\bigl(J_1\ \ 0\bigr).
 \label{eq:v30-traceless-J}
\end{equation}
In particular, a panel made solely from classically traceless massless gauge or
conformal matter states cannot by itself test
$17\kappa_1-14\kappa_2=0$.
\end{proposition}

\begin{proof}
The second operator in \eqref{eq:v30-contact-ray} is $T^2$.  If every relevant
on-shell matrix element of $T$ vanishes, its first contact insertion vanishes
on the whole panel.  Only the $T_{\mu\nu}T^{\mu\nu}$ column remains.
\end{proof}

This is the first visibility correction to the V29 phenomenology plan: a
high-energy massless panel may be convenient for RG calculations but is
insufficient as the sole physical readout of the two contact coordinates.

\subsection{Exact four-scalar contact amplitude}
\label{sec:v30-scalar-amplitude}

Take four incoming on-shell scalar momenta
\begin{equation}
 p_i^2=m_h^2,
 \qquad
 p_1+p_2+p_3+p_4=0,
 \label{eq:v30-onshell}
\end{equation}
and define
\begin{equation}
 s=(p_1+p_2)^2,
 \qquad
 t=(p_1+p_3)^2,
 \qquad
 u=(p_1+p_4)^2,
 \qquad
 s+t+u=4m_h^2.
 \label{eq:v30-stu}
\end{equation}
Let $\mathfrak P$ denote the three unordered pairings
$(12)(34)$, $(13)(24)$ and $(14)(23)$.

\begin{lemma}[Pairing contractions]
\label{lem:v30-pairing}
For the bilinear tensor \eqref{eq:v30-Bij},
\begin{align}
 \sum_{(ij)(kl)\in\mathfrak P}
 B_{\mu\nu}(p_i,p_j)B^{\mu\nu}(p_k,p_l)
 &=s^2+t^2+u^2+12m_h^4,
 \label{eq:v30-Bsum}\\
 \sum_{(ij)(kl)\in\mathfrak P}
 b(p_i,p_j)b(p_k,p_l)
 &=s^2+t^2+u^2+60m_h^4.
 \label{eq:v30-bsum}
\end{align}
\end{lemma}

\begin{proof}
For generic vectors $a,b,c,d$ with equal scalar mass $m_h$,
\begin{equation}
 B(a,b):B(c,d)
 =2\bigl[(a\!\cdot\!c)(b\!\cdot\!d)
        +(a\!\cdot\!d)(b\!\cdot\!c)\bigr]
  -2m_h^2(a\!\cdot\!b+c\!\cdot\!d)+4m_h^4.
 \label{eq:v30-generic-contraction}
\end{equation}
Also
$b(a,b)=4m_h^2-2a\!\cdot\!b$.
Insert
$p_1\!\cdot\!p_2=(s-2m_h^2)/2$ and its $t,u$ analogues,
sum the three pairings, and use $s+t+u=4m_h^2$.
\end{proof}

\begin{theorem}[Two-dimensional on-shell visibility of the contact pair]
\label{thm:v30-scalar-visibility}
Up to one common nonzero convention factor, the normalized four-scalar contact
amplitude generated by \eqref{eq:v30-contact-ray} is
\begin{equation}
 \boxed{
 \widehat{\mathcal M}_{4h}(s,t,u)
 =\Sigma\,X+\Pi\,m_h^4,}
 \qquad
 X:=s^2+t^2+u^2,
 \label{eq:v30-M-affine}
\end{equation}
where
\begin{equation}
 \boxed{
 \Sigma:=\kappa_1+\kappa_2,
 \qquad
 \Pi:=12\kappa_1+60\kappa_2.}
 \label{eq:v30-SigmaPi}
\end{equation}
For two physical kinematic points with $m_h\neq0$ and $X_1\neq X_2$,
\begin{equation}
 \begin{pmatrix}
  \widehat{\mathcal M}_1\\
  \widehat{\mathcal M}_2
 \end{pmatrix}
 =
 \begin{pmatrix}
  X_1&m_h^4\\
  X_2&m_h^4
 \end{pmatrix}
 \begin{pmatrix}\Sigma\\\Pi\end{pmatrix},
 \label{eq:v30-two-point-J}
\end{equation}
and
\begin{equation}
 \det D_{(\Sigma,\Pi)}
 (\widehat{\mathcal M}_1,\widehat{\mathcal M}_2)
 =m_h^4(X_1-X_2)\neq0.
 \label{eq:v30-two-point-det}
\end{equation}
Hence the massive scalar channel separates the two physical contact
coordinates.
\end{theorem}

\begin{proof}
The quartic contact insertion is the sum over the three pairings of
$\kappa_1B:B+\kappa_2bb$.  Equations~\eqref{eq:v30-Bsum} and
\eqref{eq:v30-bsum} give \eqref{eq:v30-M-affine}.  The determinant in
\eqref{eq:v30-two-point-det} is immediate.
\end{proof}

\begin{corollary}[Explicit dual projectors]
\label{cor:v30-dual-projectors}
For $X_1\neq X_2$,
\begin{align}
 \Sigma
 &=\frac{\widehat{\mathcal M}_1-\widehat{\mathcal M}_2}{X_1-X_2},
 \label{eq:v30-Sigma-extract}\\
 \Pi
 &=\frac{X_1\widehat{\mathcal M}_2-X_2\widehat{\mathcal M}_1}
 {m_h^4(X_1-X_2)}.
 \label{eq:v30-Pi-extract}
\end{align}
Thus the two contact combinations are reconstructible from two on-shell
kinematic rows without referring to an off-shell curvature basis.
\end{corollary}

\subsection{The V29 Burnett ray becomes a direct amplitude relation}
\label{sec:v30-Burnett-amplitude}

Insert the populated physical ray
$(\kappa_1,\kappa_2)=\kappa_0(14,17)$ from
\eqref{eq:v30-contact-ray}.  Then
\begin{equation}
 \boxed{
 \Sigma=31\kappa_0,
 \qquad
 \Pi=1188\kappa_0.}
 \label{eq:v30-ray-SigmaPi}
\end{equation}
Consequently
\begin{equation}
 \boxed{31\Pi-1188\Sigma=0.}
 \label{eq:v30-PiSigma-relation}
\end{equation}

\begin{theorem}[Reference-scale two-point observable conormal]
\label{thm:v30-two-point-conormal}
Let
\begin{equation}
 \Delta\widehat{\mathcal M}_j
 :=\widehat{\mathcal M}^{\rm obs}_j
   -\widehat{\mathcal M}^{(0)}_j
 \label{eq:v30-residual-M}
\end{equation}
be the physical amplitude after subtraction of all selector-independent terms
at the common reference scale $\mu_0$.  To first order in the worked Burnett
selector,
\begin{equation}
 \Delta\widehat{\mathcal M}_j
 =\kappa_0\bigl(31X_j+1188m_h^4\bigr)+O(a_4^2).
 \label{eq:v30-ray-amplitude}
\end{equation}
Hence any two distinct kinematic points obey
\begin{equation}
 \boxed{
 \bigl(31X_1+1188m_h^4\bigr)
 \Delta\widehat{\mathcal M}_2
 -
 \bigl(31X_2+1188m_h^4\bigr)
 \Delta\widehat{\mathcal M}_1
 =O(a_4^2).}
 \label{eq:v30-two-point-observable}
\end{equation}
Equivalently, combining the two extracted contact coordinates gives
\eqref{eq:v30-PiSigma-relation}.
\end{theorem}

\begin{proof}
Equation~\eqref{eq:v30-ray-amplitude} follows by substituting
\eqref{eq:v30-ray-SigmaPi} into \eqref{eq:v30-M-affine}.  The determinant of
the two vectors
$(\Delta\widehat{\mathcal M}_1,
\Delta\widehat{\mathcal M}_2)$ and
$(31X_1+1188m_h^4,31X_2+1188m_h^4)$ then vanishes to the retained selector
order.
\end{proof}

\subsection{A three-row falsification panel}
\label{sec:v30-three-row}

Two kinematic rows are enough to reconstruct $(\kappa_1,\kappa_2)$, but they
cannot by themselves falsify the assumption that the residual amplitude lies
in the two-dimensional stress-contact bank.  Add a third point with distinct
$X_3$.

\begin{proposition}[Contact-bank closure row]
\label{prop:v30-three-point-closure}
Every amplitude generated solely by the two operators in
\eqref{eq:v30-contact-ray} satisfies
\begin{equation}
 \boxed{
 \det
 \begin{pmatrix}
  1&X_1&\Delta\widehat{\mathcal M}_1\\
  1&X_2&\Delta\widehat{\mathcal M}_2\\
  1&X_3&\Delta\widehat{\mathcal M}_3
 \end{pmatrix}
 =0}
 \label{eq:v30-bank-closure}
\end{equation}
through the retained local contact order.  Violation of
\eqref{eq:v30-bank-closure} after the known subtraction is a direct certificate
that the physical contact bank must be enlarged before testing the Burnett
conormal.
\end{proposition}

\begin{proof}
Equation~\eqref{eq:v30-M-affine} is affine in the single kinematic scalar $X$.
Three points of an affine function are linearly dependent in the displayed
$3\times3$ determinant.
\end{proof}

If \eqref{eq:v30-bank-closure} passes, the worked Burnett branch further
requires the one-dimensional ray
\begin{equation}
 h_j:=31X_j+1188m_h^4,
 \qquad
 \Delta\widehat{\mathcal M}_j=\kappa_0h_j+O(a_4^2).
 \label{eq:v30-hj}
\end{equation}
Two independent conormals may be chosen as
\begin{equation}
 q^{(12)}=(h_2,-h_1,0),
 \qquad
 q^{(13)}=(h_3,0,-h_1).
 \label{eq:v30-three-conormals}
\end{equation}
Thus
\begin{equation}
 \boxed{
 q^{(12)}\!\cdot\!\Delta\widehat{\mathbf M}=O(a_4^2),
 \qquad
 q^{(13)}\!\cdot\!\Delta\widehat{\mathbf M}=O(a_4^2).}
 \label{eq:v30-three-relations}
\end{equation}

\subsection{Concrete equal-mass benchmark kinematics}
\label{sec:v30-benchmark}

For reproducibility, take equal-mass $2\to2$ scattering in the centre-of-mass
frame at scattering angle $\theta=\pi/2$.  Then
\begin{equation}
 t=u=-\frac{s-4m_h^2}{2},
 \qquad
 X(s)=s^2+\frac12(s-4m_h^2)^2.
 \label{eq:v30-X90}
\end{equation}
Choose
\begin{equation}
 s_1=6m_h^2,
 \qquad
 s_2=8m_h^2,
 \qquad
 s_3=10m_h^2.
 \label{eq:v30-benchmark-s}
\end{equation}
Then
\begin{equation}
 (X_1,X_2,X_3)=m_h^4(38,72,118)
 \label{eq:v30-benchmark-X}
\end{equation}
and the Burnett amplitude ray is proportional to
\begin{equation}
 \boxed{h=m_h^4(2366,3420,4846).}
 \label{eq:v30-benchmark-h}
\end{equation}
One convenient integer-normalized pair of observable conormals is therefore
\begin{equation}
 \boxed{
 1710\,\Delta\widehat{\mathcal M}_1
 -1183\,\Delta\widehat{\mathcal M}_2
 =O(a_4^2),}
 \label{eq:v30-benchmark-r12}
\end{equation}
\begin{equation}
 \boxed{
 2423\,\Delta\widehat{\mathcal M}_1
 -1183\,\Delta\widehat{\mathcal M}_3
 =O(a_4^2).}
 \label{eq:v30-benchmark-r13}
\end{equation}
The common overall normalization of the contact Feynman rule cancels from both
relations.

\subsection{Reference-scale bypass of the unpopulated standard mixing table}
\label{sec:v30-reference-bypass}

V29 already proved that under an invertible finite coefficient change
$C'=F(C)$, the microscopic tangent and amplitude Jacobian transform as
\begin{equation}
 v'=J_Fv,
 \qquad
 J_{\mathcal A}'=J_{\mathcal A}J_F^{-1}.
 \label{eq:v30-scheme-transform}
\end{equation}
Therefore the physical amplitude tangent
\begin{equation}
 h=J_{\mathcal A}v
 \label{eq:v30-physical-h}
\end{equation}
is invariant:
\begin{equation}
 \boxed{J_{\mathcal A}'v'=J_{\mathcal A}v.}
 \label{eq:v30-h-invariant}
\end{equation}

\begin{corollary}[Reference-scale observable theorem does not require a
numerical native-to-standard Wilson map]
\label{cor:v30-reference-bypass}
At the common physical reference scale $\mu_0$, the amplitude relations
\eqref{eq:v30-two-point-observable} and \eqref{eq:v30-three-relations} may be
tested directly on physical amplitudes without first populating the numerical
finite coefficient map $J_F$ or a one-loop anomalous-dimension matrix.
Those ingredients become load-bearing only when the relation is transported to
a separated scale or compared in an intermediate Wilson basis.
\end{corollary}

This is the phenomenologically useful upstream simplification: the first
benchmark can be formulated in observables now, while the substantially larger
one-loop dimension-eight mixing problem remains a subsequent transport task.

\subsection{What the scalar panel does not prove}
\label{sec:v30-firewalls}

\begin{firewall}[Amplitude visibility is not microscopic selector completeness]
\label{fw:v30-selector-completeness}
The three scalar rows prove visibility of the two stress-contact coordinates and
provide a falsification test for amplitude components outside the affine
$X$-bank.  They do not prove that the Burnett selector is the only microscopic
source capable of producing the same one-dimensional amplitude ray.  Any
additional selector whose physical amplitude derivative is collinear with
$h$ survives this panel undetected.  The V29/V26 completeness residual must
still be evaluated on a larger multi-species/helicity bank before universal
promotion.
\end{firewall}

\begin{firewall}[One-loop closure is still open]
\label{fw:v30-loop-closure}
The two stress-contact operators are not assumed to form a closed one-loop
Einstein--SM operator subspace.  At a separated scale, or beyond the direct
reference-scale amplitude statement, one must enlarge to the smallest
field-redefinition/EOM-reduced operator bank closed under the retained mixing
matrix.  V30 makes no numerical anomalous-dimension claim.
\end{firewall}

\begin{firewall}[The scalar formula is a calibrated anchor, not all of the SM]
\label{fw:v30-scalar-scope}
Equation~\eqref{eq:v30-M-affine} is the canonical massive-scalar/Higgs anchor
in the declared stress-tensor convention.  Gauge, fermion, broken-phase mixed
and improved-Higgs channels should be added as orthogonal visibility rows for
the eventual full Einstein--SM audit; they are not inferred from the scalar
formula.
\end{firewall}

\section{V30 status ledger}
\label{sec:v30-status}

\subsection*{Proved in V30}
\begin{enumerate}[leftmargin=2.2em]
\item A physical panel made solely from traceless external stress-tensor
matrix elements has rank at most one on $(\kappa_1,\kappa_2)$ and cannot test
the V29 relation by itself.
\item For the canonical massive scalar/Higgs stress tensor, the complete
four-point contact amplitude is
$\widehat{\mathcal M}=\Sigma X+\Pi m_h^4$ with
$\Sigma=\kappa_1+\kappa_2$ and
$\Pi=12\kappa_1+60\kappa_2$.
\item Two distinct physical kinematic points with $m_h\neq0$ determine both
contact coordinates.  A third point provides the independent affine-bank
closure test \eqref{eq:v30-bank-closure}.
\item The populated V29 ray $(\kappa_1,\kappa_2)\propto(14,17)$ becomes the
scheme-free reference-scale relation $31\Pi-1188\Sigma=0$ and the direct
amplitude conormals \eqref{eq:v30-two-point-observable} and
\eqref{eq:v30-three-relations}.
\item The concrete $90^\circ$ benchmark at
$s/m_h^2=(6,8,10)$ has amplitude ray
$h/m_h^4=(2366,3420,4846)$ and integer conormals
\eqref{eq:v30-benchmark-r12}--\eqref{eq:v30-benchmark-r13}.
\item At the common reference scale, the physical amplitude tangent is
invariant under finite intermediate coefficient-scheme changes.  Therefore the
first observable benchmark does not require a populated numerical
native-to-standard contact map or anomalous-dimension matrix.
\end{enumerate}

\subsection*{Inherited}
\begin{enumerate}[leftmargin=2.2em]
\item V28: the worked isotropic compound-Poisson Burnett selector populates the
curvature matching ray.
\item V29: after the four-dimensional field-redefinition/EOM quotient, that ray
lands in the physical stress-contact direction
$(\kappa_1,\kappa_2)=a_4(14,17)/(12M_g^4)$.
\item V16--V29: selector-independent lower-weight/loop/threshold contributions
must be calibrated and subtracted on one common source/amplitude panel; finite
scheme maps transport tangents and conormals but do not alter the final
physical amplitude ray.
\end{enumerate}

\subsection*{Conditional}
\begin{enumerate}[leftmargin=2.2em]
\item The explicit scalar coefficients $12$ and $60$ use the canonical
minimally coupled stress tensor.  A varying Higgs improvement coordinate
requires an enlarged bank; a fixed alternative improvement convention requires
recomputing the physical Jacobian consistently.
\item The amplitude relations are stated to first order in the worked Burnett
selector, after subtraction of all selector-independent contributions through
the same perturbative/EFT order.
\item The branch-specific relation becomes a complete Renewal--ESM prediction
only after every additional microscopic selector has been shown to have no
transverse image on a sufficiently rich multi-species/helicity amplitude bank.
\end{enumerate}

\subsection*{Open beyond V30}
\begin{enumerate}[leftmargin=2.2em]
\item Add orthogonal SM external-state rows---at minimum massive fermion and
gauge/Higgs mixed channels---to test microscopic selector completeness beyond
the scalar amplitude ray.
\item Construct the smallest physical dimension-eight Einstein--SM bank needed
for those amplitudes and prove closure under the retained one-loop mixing.
\item Populate the native-to-standard physical matching Jacobian and the
variational RG matrix only after that bank is fixed; transport the V30
reference-scale conormal to separated scales.
\item Populate numerical/experimental or same-record microscopic values of the
three benchmark residual amplitudes and propagate their uncertainties through
the two conormal tests.
\item Continue the separate microscopic lock of the $RH^\dagger H$ coefficient
and its stress-tensor improvement dependence.
\end{enumerate}

\subsection*{Exact next load-bearing theorem}
\begin{center}
\fbox{\parbox{0.92\textwidth}{
\textbf{P4n: multi-species physical-amplitude closure.}
Starting from the explicit V30 massive-scalar reference-scale ray, add the
smallest set of independent on-shell Standard-Model external-state/helicity
channels which (i) keeps full visibility of the two physical stress-contact
coordinates, (ii) detects extra dimension-eight contact directions that are
collinear on four-Higgs scattering, and (iii) supports a finite
field-redefinition/EOM-reduced operator bank.  Compute the physical amplitude
Jacobian of that bank and the transverse selector-completeness residual.  Only
a vanishing residual promotes the V30 one-dimensional Burnett amplitude ray
beyond the worked branch.  Once the bank is fixed, populate its one-loop
standard-scheme mixing and transport the observable conormal away from the
reference scale.}}
\end{center}

\section*{Append-only continuation marker: V30}
\addcontentsline{toc}{section}{Append-only continuation marker: V30}
\textbf{End of V30.}  The complete V1--V29 body above is retained verbatim.
V30 goes upstream of the unpopulated full dimension-eight mixing table and
constructs the first explicit on-shell visibility map for the physical V29
stress-contact ray.  A traceless massless panel is proved insufficient, while
a canonical massive scalar/Higgs four-point panel separates the two contact
coordinates with
$\widehat{\mathcal M}=\Sigma(s^2+t^2+u^2)+\Pi m_h^4$.
The worked Burnett selector gives
$(\Sigma,\Pi)=\kappa_0(31,1188)$ and therefore a direct scheme-free
reference-scale amplitude relation.  Three physical kinematic rows provide one
contact-bank closure test plus two independent Burnett-ray conormals; the
explicit $90^\circ$ benchmark at $s/m_h^2=(6,8,10)$ yields the integer
relations in \eqref{eq:v30-benchmark-r12}--\eqref{eq:v30-benchmark-r13}.
Full one-loop transport and universal promotion remain downstream of a richer
multi-species/helicity selector-completeness audit.



% =============================================================================
% V31 append-only continuation: P4n multi-species amplitude closure
% =============================================================================
\section{P4n: multi-species amplitude closure of the physical stress-contact bank}
\label{sec:v31-multispecies-closure}

\subsection{Why the scalar anchor must be enlarged before one-loop transport}
\label{sec:v31-upstream}

V30 supplied a complete on-shell visibility map for the two universal stress
contacts on a canonical massive scalar/Higgs channel.  That calculation is
necessary but not sufficient for selector completeness.  A dimension-eight
operator which acts only in the scalar sector can be exactly collinear with the
V30 Burnett ray on every four-Higgs row while failing to possess the universal
cross-species structure forced by a contact built from the \emph{total}
stress tensor.  The next observable gate must therefore test species
universality before it tests scale transport.

Fix throughout this section one P1/reference convention for the improved Higgs
stress tensor.  Split the calibrated Standard-Model stress tensor into three
representative sectors
\begin{equation}
 T_{\mu\nu}=T^H_{\mu\nu}+T^F_{\mu\nu}+T^V_{\mu\nu},
 \label{eq:v31-stress-split}
\end{equation}
where $H$ is a massive scalar/Higgs sector, $F$ is a massive fermion sector,
and $V$ is a massless gauge sector.  At the tree/reference level used for the
visibility calculation,
\begin{equation}
 T^V:=g^{\mu\nu}T^V_{\mu\nu}=0,
 \label{eq:v31-gauge-trace-zero}
\end{equation}
while $T^H$ and $T^F$ are nonzero on generic massive external states.  Quantum
trace-anomaly contributions belong to the loop/reference subtraction or to an
enlarged loop-level bank; they are not folded into the tree visibility
Jacobian below.

\subsection{The nine-dimensional species-resolved stress bank}
\label{sec:v31-nine-bank}

Let $\mathfrak S=\{H,F,V\}$.  For $a\le b$ define the normalized symmetric
pair operators
\begin{equation}
 \mathcal O^{(2)}_{ab}
 :=(2-\delta_{ab})T_a^{\mu\nu}T^b_{\mu\nu},
 \label{eq:v31-O2}
\end{equation}
and, when $a,b\in\{H,F\}$,
\begin{equation}
 \mathcal O^{(0)}_{ab}
 :=(2-\delta_{ab})T_aT_b.
 \label{eq:v31-O0}
\end{equation}
The pair normalization is chosen so that
\begin{equation}
 T_{\mu\nu}T^{\mu\nu}
 =\sum_{a\le b}\mathcal O^{(2)}_{ab},
 \qquad
 T^2
 =\sum_{\substack{a\le b\\ a,b\in\{H,F\}}}\mathcal O^{(0)}_{ab}
 \label{eq:v31-total-expand}
\end{equation}
on the tree panel with $T^V=0$.

The smallest bank which can test species universality of the V29 stress
contacts is therefore
\begin{equation}
 \mathcal B_{\rm stress}
 =\operatorname{span}\left\{
 \mathcal O^{(2)}_{HH},\mathcal O^{(2)}_{HF},\mathcal O^{(2)}_{HV},
 \mathcal O^{(2)}_{FF},\mathcal O^{(2)}_{FV},\mathcal O^{(2)}_{VV},
 \mathcal O^{(0)}_{HH},\mathcal O^{(0)}_{HF},\mathcal O^{(0)}_{FF}
 \right\}.
 \label{eq:v31-bank}
\end{equation}
Write its coefficient vector as
\begin{equation}
 y=(b_2,b_0)\in\mathbb R^6\oplus\mathbb R^3,
 \label{eq:v31-y}
\end{equation}
with ordering
\begin{align}
 b_2&=(b_{HH},b_{HF},b_{HV},b_{FF},b_{FV},b_{VV}),
 \label{eq:v31-b2}\\
 b_0&=(c_{HH},c_{HF},c_{FF}).
 \label{eq:v31-b0}
\end{align}
The universal contact of V29 is the two-dimensional subspace
\begin{equation}
 \boxed{
 b_2=\kappa_1\mathbf 1_6,
 \qquad
 b_0=\kappa_2\mathbf 1_3.}
 \label{eq:v31-universal-subspace}
\end{equation}
The worked Burnett branch is the one-dimensional ray
\begin{equation}
 \boxed{
 y_{\rm B}(\kappa_0)
 =\kappa_0 v_{\rm B},
 \qquad
 v_{\rm B}:=(14\mathbf 1_6,17\mathbf 1_3).}
 \label{eq:v31-burnett-vector}
\end{equation}
Its Euclidean norm is
\begin{equation}
 \boxed{\|v_{\rm B}\|^2=6\cdot14^2+3\cdot17^2=2043.}
 \label{eq:v31-vnorm}
\end{equation}

\begin{theorem}[Exact species-universality decomposition]
\label{thm:v31-universality-decomposition}
For any $y=(b_2,b_0)\in\mathbb R^9$, define
\begin{equation}
 \bar b_2:=\frac16\mathbf 1_6^Tb_2,
 \qquad
 \bar b_0:=\frac13\mathbf 1_3^Tb_0,
 \label{eq:v31-means}
\end{equation}
\begin{equation}
 r_2:=\left(I_6-\frac16\mathbf 1_6\mathbf 1_6^T\right)b_2,
 \qquad
 r_0:=\left(I_3-\frac13\mathbf 1_3\mathbf 1_3^T\right)b_0,
 \label{eq:v31-species-residuals}
\end{equation}
and the Burnett-ratio residual
\begin{equation}
 \Delta_{\rm B}:=17\bar b_2-14\bar b_0.
 \label{eq:v31-ratio-residual}
\end{equation}
Then:
\begin{enumerate}[label=\textup{(\roman*)},leftmargin=2.8em]
\item $r_2=0$ iff all six tensor-contact pair coefficients are equal;
\item $r_0=0$ iff all three trace-contact pair coefficients are equal;
\item after \textup{(i)--(ii)}, $\Delta_{\rm B}=0$ iff the common coefficients
satisfy $(\kappa_1,\kappa_2)\propto(14,17)$;
\item the orthogonal distance to the Burnett ray decomposes exactly as
\begin{equation}
 \boxed{
 \left\|y-\frac{v_{\rm B}v_{\rm B}^T}{2043}y\right\|^2
 =\|r_2\|^2+\|r_0\|^2+\frac{2}{227}\Delta_{\rm B}^2.}
 \label{eq:v31-distance-decomp}
\end{equation}
\end{enumerate}
Consequently the nine-dimensional stress bank contains exactly eight
independent linear directions transverse to the worked Burnett ray: five
species-nonuniversal tensor directions, two species-nonuniversal trace
directions and one universal Burnett-ratio direction.
\end{theorem}

\begin{proof}
The first two assertions are the orthogonal decompositions of $\mathbb R^6$
and $\mathbb R^3$ into the all-ones line and its orthogonal complement.  The
universal two-plane is therefore spanned by
$e_2=(\mathbf 1_6,0)$ and $e_0=(0,\mathbf 1_3)$.  The Burnett ray in this plane
is $v_{\rm B}=14e_2+17e_0$.  An orthogonal vector in the same plane is
\begin{equation}
 w_{\rm B}=(17\mathbf 1_6,-28\mathbf 1_3),
 \qquad
 v_{\rm B}^Tw_{\rm B}=0,
 \qquad
 \|w_{\rm B}\|^2=4086.
 \label{eq:v31-wB}
\end{equation}
For the universal projection
$u=(\bar b_2\mathbf 1_6,\bar b_0\mathbf 1_3)$,
\begin{equation}
 u^Tw_{\rm B}=6(17\bar b_2-14\bar b_0)=6\Delta_{\rm B}.
 \label{eq:v31-uw}
\end{equation}
Hence the squared residual inside the universal plane is
$36\Delta_{\rm B}^2/4086=2\Delta_{\rm B}^2/227$.  The two species-residual
subspaces are mutually orthogonal and orthogonal to the universal plane, which
proves \eqref{eq:v31-distance-decomp} and the dimension count.
\end{proof}

A convenient least-squares estimate of the single worked-branch scale is
\begin{equation}
 \boxed{
 \widehat\kappa_0
 =\frac{v_{\rm B}^Ty}{2043}
 =\frac{14\sum_{i=1}^6(b_2)_i+17\sum_{j=1}^3(b_0)_j}{2043}.}
 \label{eq:v31-kappa-est}
\end{equation}
For data with a nontrivial covariance matrix one replaces the Euclidean
projection by the corresponding inverse-covariance weighted projection; the
algebraic nullspace is unchanged when the covariance is nonsingular.

\subsection{Why mixed-species rows are load bearing}
\label{sec:v31-mixed-essential}

\begin{proposition}[Self channels cannot certify total-stress universality]
\label{prop:v31-self-insufficient}
Measurements confined to $HH$, $FF$ and $VV$ self-scattering determine only
the diagonal entries of the tensor pair matrix and the $HH,FF$ diagonal trace
entries.  They leave the mixed coefficients
\begin{equation}
 b_{HF},\ b_{HV},\ b_{FV},\ c_{HF}
 \label{eq:v31-unseen-mixed}
\end{equation}
undetermined.  Therefore even exact equality of all visible self-channel
coefficients cannot imply \eqref{eq:v31-universal-subspace}.
\end{proposition}

\begin{proof}
Changing any of the four coordinates in \eqref{eq:v31-unseen-mixed} leaves all
three self-channel amplitudes unchanged at contact order.  In particular one
may keep all measured diagonal coefficients equal while varying an off-diagonal
coefficient away from the universal value.  Thus the self-channel restriction
has a nontrivial kernel on the species-universality quotient.
\end{proof}

This is the on-shell analogue of the same-history mixed-source requirement in
the Grand Tensor closure programme: cross-sector incidence is not reconstructible
from two separate marginal measurements.

\subsection{Finite physical amplitude projectors}
\label{sec:v31-projectors}

For each species pair $ab$, let $\mathcal A_{ab}$ denote the finite-dimensional
space of retained on-shell contact amplitudes as functions of the chosen
nonexceptional kinematic/helicity variables.  Let
\begin{equation}
 f^{(2)}_{ab}\in\mathcal A_{ab},
 \qquad
 f^{(0)}_{ab}\in\mathcal A_{ab}
 \label{eq:v31-amplitude-functions}
\end{equation}
be the amplitude functions generated by $\mathcal O^{(2)}_{ab}$ and
$\mathcal O^{(0)}_{ab}$ whenever the latter exists.

\begin{lemma}[Finite dual-row lemma]
\label{lem:v31-dual-row}
If $f_1,\ldots,f_r$ are linearly independent functions on a physical
kinematic/helicity domain, then there exist $r$ physical evaluation rows
$\varepsilon_1,\ldots,\varepsilon_r$ such that the matrix
\begin{equation}
 J_{ij}=\varepsilon_i(f_j)
 \label{eq:v31-eval-matrix}
\end{equation}
is invertible.
\end{lemma}

\begin{proof}
If every $r\times r$ evaluation matrix were singular, the evaluation vectors
$(f_1(x),\ldots,f_r(x))$ would span a proper subspace of $\mathbb R^r$ (or
$\mathbb C^r$).  A nonzero covector annihilating that span would give a
nontrivial linear combination of the $f_j$ vanishing at every physical point,
contradicting independence.
\end{proof}

For the traceful pairs
\begin{equation}
 HH,\quad HF,\quad FF,
 \label{eq:v31-traceful-pairs}
\end{equation}
choose two fit rows for which
$f^{(2)}_{ab},f^{(0)}_{ab}$ are independent.  For the gauge-containing pairs
\begin{equation}
 HV,\quad FV,\quad VV,
 \label{eq:v31-gauge-pairs}
\end{equation}
the tree trace column vanishes and one nonzero tensor row is sufficient for
coefficient extraction.  If the asserted independence fails for a proposed
channel, that channel is not a two-coordinate visibility panel and must be
replaced or its quotient dimension reduced; no pseudoinverse is promoted to a
physical inverse silently.

\subsection{A fifteen-row pairwise closure panel}
\label{sec:v31-fifteen-panel}

To separate coefficient extraction from bank validation, add one held-out row
to each species pair.  This gives the following finite panel:
\begin{center}
\begin{tabular}{c|c|c|c}
Pair & visible stress structures & fit rows & held-out rows\\ \hline
$HH$ & $\mathcal O^{(2)}_{HH},\mathcal O^{(0)}_{HH}$ & 2 & 1\\
$HF$ & $\mathcal O^{(2)}_{HF},\mathcal O^{(0)}_{HF}$ & 2 & 1\\
$FF$ & $\mathcal O^{(2)}_{FF},\mathcal O^{(0)}_{FF}$ & 2 & 1\\
$HV$ & $\mathcal O^{(2)}_{HV}$ & 1 & 1\\
$FV$ & $\mathcal O^{(2)}_{FV}$ & 1 & 1\\
$VV$ & $\mathcal O^{(2)}_{VV}$ & 1 & 1
\end{tabular}
\end{center}
Thus the complete pairwise panel uses
\begin{equation}
 \boxed{3+3+3+2+2+2=15\ \text{physical amplitude rows}.}
 \label{eq:v31-fifteen}
\end{equation}
The three $HH$ rows are already supplied by V30, so only twelve additional
rows are required to realize this declared panel.

For each pair let $J_{ab}^{\rm fit}$ be its fit Jacobian and
$A_{ab}^{\rm fit}$ the selector-subtracted measured amplitude vector.  Define
\begin{equation}
 \widehat c_{ab}:=(J_{ab}^{\rm fit})^{-1}A_{ab}^{\rm fit}
 \label{eq:v31-pair-fit}
\end{equation}
for the traceful two-column pairs, with the obvious one-column division on a
gauge-containing pair.  If $j_{ab}^{\rm ho}$ is the held-out Jacobian row,
define
\begin{equation}
 \boxed{
 \Delta_{ab}^{\rm ho}
 :=A_{ab}^{\rm ho}-j_{ab}^{\rm ho}\widehat c_{ab}.}
 \label{eq:v31-heldout-residual}
\end{equation}

\begin{theorem}[Pairwise stress-bank closure certificate]
\label{thm:v31-bank-closure}
If any $\Delta_{ab}^{\rm ho}$ is nonzero beyond its propagated uncertainty, the
nine-dimensional stress-bilinear bank \eqref{eq:v31-bank} is insufficient on
the corresponding physical amplitude panel.  The species-universality and
Burnett-ratio tests of Theorem~\ref{thm:v31-universality-decomposition} must then be
postponed until the missing amplitude direction is added.

If all held-out residuals vanish, the fifteen rows are consistent with the
nine-dimensional stress bank and determine one coefficient vector
$y=(b_2,b_0)$ on that panel.  The eight transverse diagnostics are then exactly
$r_2$, $r_0$ and $\Delta_{\rm B}$.
\end{theorem}

\begin{proof}
On each pair, every amplitude generated by the declared local bank lies in the
column space of its fit Jacobian and therefore predicts the held-out row.  A
failed held-out prediction cannot be repaired by changing the already fitted
coefficients without violating the fit rows.  Conversely, if all held-out rows
agree, the extracted pair coefficients consistently define $y$ on the selected
panel.  The final statement is Theorem~\ref{thm:v31-universality-decomposition}.
\end{proof}

\begin{firewall}[Finite panel closure is not a complete SMEFT basis theorem]
\label{fw:v31-panel-not-global}
An operator outside $\mathcal B_{\rm stress}$ whose amplitudes happen to be
collinear with the stress-bank columns on all fifteen selected rows is not
excluded by Theorem~\ref{thm:v31-bank-closure}.  Universal promotion therefore still
requires either a basis theorem for the retained on-shell amplitude sector or
additional rows that separate every allowed nuisance direction.  V31 proves a
finite falsification/visibility theorem, not completeness of the entire
dimension-eight Standard-Model EFT.
\end{firewall}

\subsection{Relation to the V30 scalar benchmark}
\label{sec:v31-v30-embed}

The V30 four-Higgs formula is the $HH$ restriction of the V31 bank.  Replacing
$(\kappa_1,\kappa_2)$ there by $(b_{HH},c_{HH})$ gives
\begin{equation}
 \widehat{\mathcal M}_{HH}
 =(b_{HH}+c_{HH})X
 +(12b_{HH}+60c_{HH})m_h^4.
 \label{eq:v31-HH-restriction}
\end{equation}
Hence the V30 two-coordinate inversion exactly determines the first tensor and
trace entries of the nine-component vector $y$.  What V31 adds are the five
remaining tensor pair coefficients and the two remaining trace pair
coefficients.  Only after those rows are populated can the statement
``the contact is built from the total stress tensor'' be tested rather than
assumed.

\subsection{A one-loop closure gate after the physical bank is fixed}
\label{sec:v31-loop-gate}

Let $\mathcal B_{\rm cand}$ be any field-redefinition/EOM-reduced physical
operator bank which contains $\mathcal B_{\rm stress}$ and is sufficient for
the selected multi-species amplitudes.  Let $P_{\rm cand}$ denote its projector
inside a larger finite operator bank at the retained EFT order, and let
$\Gamma^{(8)}$ be the one-loop mixing operator in a fixed standard scheme.
Define the mixing innovation
\begin{equation}
 \boxed{
 \mathcal I_{\rm mix}
 :=(I-P_{\rm cand})\Gamma^{(8)}P_{\rm cand}.}
 \label{eq:v31-mixing-innovation}
\end{equation}
Then
\begin{equation}
 \boxed{
 \mathcal I_{\rm mix}=0
 \iff
 \Gamma^{(8)}\operatorname{Ran}P_{\rm cand}
 \subseteq\operatorname{Ran}P_{\rm cand}.}
 \label{eq:v31-mixing-closure}
\end{equation}
Thus the full one-loop transport calculation should begin only after the tree
amplitude bank passes its visibility/held-out tests.  If
$\mathcal I_{\rm mix}\neq0$, the operator bank must be enlarged by the finite
range of $\mathcal I_{\rm mix}$ before transporting the Burnett conormal.

\subsection{Current phenomenological status after the multi-species split}
\label{sec:v31-phenom-status}

The worked microscopic branch makes the sharp nine-coordinate prediction
\begin{equation}
 \boxed{
 b_2=14\kappa_0\mathbf 1_6,
 \qquad
 b_0=17\kappa_0\mathbf 1_3.}
 \label{eq:v31-nine-prediction}
\end{equation}
Inside the declared stress-bilinear bank this is a codimension-eight ray.  It
contains three logically distinct tests:
\begin{enumerate}[leftmargin=2.2em]
\item tensor species universality: $r_2=0$ (five conditions);
\item trace species universality: $r_0=0$ (two conditions);
\item the renewal Burnett ratio: $\Delta_{\rm B}=0$ (one condition).
\end{enumerate}
This decomposition is preferable to testing the single four-Higgs ratio alone:
a species-specific scalar operator may satisfy the V30 relation accidentally,
but it will generally fail one or more mixed-species universality rows.

\begin{firewall}[Improvement and trace conventions]
\label{fw:v31-improvement}
The Higgs improvement/nonminimal coordinate must be fixed as P1/reference data
for the nine-coordinate formulas above.  If it is varied, its induced changes
in $T^H_{\mu\nu}$ and $T^H$ are an additional physical/nuisance direction and
must be appended to the amplitude Jacobian before evaluating $r_2$, $r_0$ or
$\Delta_{\rm B}$.  Likewise a loop-level gauge trace anomaly is not set to zero;
it belongs to the loop/reference subtraction or to the enlarged mixing-closed
bank.
\end{firewall}

\section{V31 status ledger}
\label{sec:v31-status}

\subsection*{Proved in V31}
\begin{enumerate}[leftmargin=2.2em]
\item The smallest species-resolved stress-bilinear bank for one massive
scalar/Higgs sector, one massive fermion sector and one massless gauge sector
has nine physical pair coefficients: six tensor contacts and three trace
contacts.
\item A universal total-stress contact is exactly the two-plane
$b_2=\kappa_1\mathbf 1_6$, $b_0=\kappa_2\mathbf 1_3$; the worked V29 Burnett
branch is the ray $(14\mathbf 1_6,17\mathbf 1_3)$ of squared norm $2043$.
\item The orthogonal residual to that ray decomposes exactly into five tensor
species-universality directions, two trace species-universality directions and
one ratio residual as in \eqref{eq:v31-distance-decomp}.
\item Same-species amplitudes alone cannot certify total-stress universality;
the mixed $HF$, $HV$ and $FV$ rows are load bearing.
\item A fifteen-row pairwise panel uses nine fit rows and six held-out rows.
It gives a finite closure and falsification compiler for the nine-dimensional
stress bank;
V30 already supplies the three $HH$ rows.
\item Once a candidate physical operator bank is fixed, its one-loop closure is
the finite condition $(I-P_{\rm cand})\Gamma^{(8)}P_{\rm cand}=0$.
\end{enumerate}

\subsection*{Inherited}
\begin{enumerate}[leftmargin=2.2em]
\item V29: after the four-dimensional field-redefinition/EOM quotient, the
worked microscopic curvature ray lands in
$(\kappa_1,\kappa_2)=\kappa_0(14,17)$.
\item V30: the $HH$ amplitude separates the two contact coordinates, supplies
a third-row closure test, and gives an explicit scheme-independent
reference-scale observable relation.
\item V16--V30: all selector-independent SM/gravity/threshold contributions
must be subtracted in one common physical calibration; finite coefficient
scheme changes do not alter the final on-shell amplitude tangent.
\end{enumerate}

\subsection*{Conditional}
\begin{enumerate}[leftmargin=2.2em]
\item For each traceful pair, the two declared stress amplitude functions must
be linearly independent on the selected physical kinematic/helicity domain;
otherwise that pair has a smaller visible quotient and the panel must be
changed.
\item The Higgs improvement convention is fixed.  A varying improvement
coordinate enlarges the bank.
\item The fifteen-row panel certifies closure only on the selected amplitude
rows.  Full dimension-eight EFT completeness requires a basis theorem or
additional separating rows.
\end{enumerate}

\subsection*{Open beyond V31}
\begin{enumerate}[leftmargin=2.2em]
\item Populate explicit nonexceptional $HF$, $HV$, $FF$, $FV$ and $VV$
kinematic/helicity rows and their held-out partners; compute the numerical
nine-coordinate amplitude vector and covariance.
\item Enumerate the smallest additional on-shell dimension-eight nuisance bank
which contributes to those channels and test whether the fifteen-row stress
bank is complete against it.
\item Only after that tree-level bank is fixed, compute its one-loop
standard-scheme mixing innovation \eqref{eq:v31-mixing-innovation}, append any
new range, and transport the Burnett conormal to separated scales.
\item Audit every remaining microscopic selector independently.  Require zero
transverse image on the enlarged multi-species amplitude bank before promoting
the branch-specific ray.
\end{enumerate}

\subsection*{Exact next load-bearing theorem}
\begin{center}
\fbox{\parbox{0.92\textwidth}{
\textbf{P4o: explicit mixed-species amplitude tomography.}
Populate one concrete nonexceptional on-shell row set for the five new species
pairs $HF,HV,FF,FV,VV$, with two independent fit rows plus one held-out row for
each traceful pair and one fit plus one held-out row for each gauge-containing
pair.  Invert the physical pair Jacobians, assemble the nine-component vector
$y=(b_2,b_0)$, and evaluate the exact residual decomposition
$(r_2,r_0,\Delta_{\rm B})$.  If any held-out row fails, append the source-minimal
missing on-shell operator direction before computing RG mixing.  If the bank
closes, enumerate the smallest nuisance operator bank that shares these
external states and evaluate its transverse image.  Only then start the
one-loop mixing-closure calculation.}}
\end{center}

\section*{Append-only continuation marker: V31}
\addcontentsline{toc}{section}{Append-only continuation marker: V31}
\textbf{End of V31.}  The complete V1--V30 body above is retained verbatim.
V31 goes upstream of the large one-loop dimension-eight mixing problem and
constructs the minimal species-resolved physical stress-contact bank needed to
test whether the V29/V30 relation is genuinely universal across Standard-Model
sectors.  The bank has six tensor and three trace pair coefficients.  The
worked Burnett branch is the one-dimensional ray
$(14\mathbf 1_6,17\mathbf 1_3)$, with an exact eight-direction residual
splitting into five tensor species-nonuniversality directions, two trace
species-nonuniversality directions and one Burnett-ratio direction.  A
fifteen-row on-shell panel, reusing the three V30 Higgs rows, provides a finite
pairwise closure compiler before any one-loop transport is attempted.



% =============================================================================
% V32 append-only continuation: explicit multi-species amplitude tomography
% =============================================================================
\section{P4o continued: explicit multi-species stress tomography}
\label{sec:v32-explicit-tomography}

\subsection{Upstream correction: use pole-subtracted forward rows for fitting and nonforward rows for closure}
\label{sec:v32-upstream}

V31 asked for explicit nonexceptional rows for the five new species pairs.
There is a sharper organization.  The conserved stress tensor has a universal
forward one-particle normalization, so the coefficient fit can be carried out
with pole-subtracted forward functionals, where spin and polarization
kinematics collapse.  A genuinely nonforward row should then be reserved for
the held-out closure test.  This is important: a bank tested only at $t=0$ can
miss derivative operators which vanish or become collinear in the forward
limit.

Let
\begin{equation}
 \Delta\mathcal M_{ab}
 :=\mathcal M_{ab}^{\rm obs}-\mathcal M_{ab}^{\rm known}
 \label{eq:v32-subtracted-amplitude}
\end{equation}
be the amplitude after subtracting, in the same P1 convention, all calibrated
lower-weight Standard-Model and Einstein pieces, massless exchange poles,
threshold terms already assigned to the reference action, and the declared
loop/reference contribution.  Let $\mathcal P_{ab}^{\rm dir}$ denote the
direct internal-state projector.  For identical fields it is implemented by
orthogonal colour/flavour labels or, equivalently, by subtracting the crossed
contraction before the contact projection.

\begin{firewall}[The forward fit is a contact functional, not an unsubtracted
forward cross section]
\label{fw:v32-forward-subtracted}
Whenever a massless exchange channel is present, the symbol $t=0$ below is
understood only after the known nonlocal pole has been removed in
\eqref{eq:v32-subtracted-amplitude}.  No singular forward observable is being
identified with a local Wilson coefficient.
\end{firewall}

\subsection{Universal forward stress theorem}
\label{sec:v32-forward-theorem}

Use relativistically normalized one-particle states.  Conservation and the
translation Ward identity give, for any physical one-particle species $a$,
\begin{equation}
 \boxed{
 \langle p,\sigma'|T_a^{\mu\nu}(0)|p,\sigma\rangle
 =2p^\mu p^\nu\,\delta_{\sigma'\sigma}.}
 \label{eq:v32-forward-stress}
\end{equation}
Consequently
\begin{equation}
 \langle p,\sigma|T_a(0)|p,\sigma\rangle=2m_a^2,
 \label{eq:v32-forward-trace}
\end{equation}
with $m_a=0$ for a massless gauge state.  The statement is independent of the
spin of the state.  A Higgs improvement term is proportional to momentum
transfer and therefore does not change \eqref{eq:v32-forward-stress}; it does,
however, affect generic nonforward rows and remains fixed as in V31.

For an elastic pair $ab$ define
\begin{equation}
 z_{ab}:=p_a\!\cdot p_b
 =\frac{s-m_a^2-m_b^2}{2}.
 \label{eq:v32-zab}
\end{equation}
Normalize the direct pole-subtracted forward contact functional by
\begin{equation}
 \mathfrak A_{ab}^{\rm fwd}(s)
 :=\frac18\,
 \mathcal P_{ab}^{\rm dir}
 \Delta\mathcal M_{ab}(s,t=0).
 \label{eq:v32-forward-functional}
\end{equation}

\begin{theorem}[Universal forward pair tomography]
\label{thm:v32-forward-pair}
On the V31 stress bank, every traceful pair $ab\in\{HH,HF,FF\}$ obeys
\begin{equation}
 \boxed{
 \mathfrak A_{ab}^{\rm fwd}(s)
 =b_{ab}z_{ab}^2+c_{ab}m_a^2m_b^2,}
 \label{eq:v32-forward-traceful}
\end{equation}
while every pair containing a tree-level traceless gauge leg obeys
\begin{equation}
 \boxed{
 \mathfrak A_{aV}^{\rm fwd}(s)
 =b_{aV}z_{aV}^2,\qquad
 \mathfrak A_{VV}^{\rm fwd}(s)=b_{VV}\frac{s^2}{4}.}
 \label{eq:v32-forward-gauge}
\end{equation}
The normalization is the same for self and mixed pairs because the V31
operator normalization and the two direct Wick contractions of a self pair
give the same connected factor.
\end{theorem}

\begin{proof}
Insert \eqref{eq:v32-forward-stress} in the matrix elements of
\eqref{eq:v31-O2}--\eqref{eq:v31-O0}.  The tensor contraction contributes
$8(p_a\!\cdot p_b)^2$ before the normalization in
\eqref{eq:v32-forward-functional}, while the trace contraction contributes
$8m_a^2m_b^2$.  The trace term is absent for a massless gauge leg by
\eqref{eq:v31-gauge-trace-zero}.  The direct projector removes exchange
contamination without changing these contact factors.
\end{proof}

This theorem is the reason to fit the pair coefficients in the forward
channel: neither explicit spinor bilinears nor gauge helicity algebra enter the
coefficient inversion.

\subsection{Explicit $HF$ and $FF$ inversions}
\label{sec:v32-traceful-inversion}

For a traceful pair $ab$ choose the two forward fit points
\begin{equation}
 z_{ab}^{(2)}=2m_am_b,
 \qquad
 z_{ab}^{(3)}=3m_am_b,
 \label{eq:v32-traceful-z}
\end{equation}
and define dimensionless residual rows
\begin{equation}
 u_{ab}^{(n)}
 :=\frac{\mathfrak A_{ab}^{\rm fwd}(s^{(n)})}{m_a^2m_b^2},
 \qquad n=2,3.
 \label{eq:v32-uab}
\end{equation}
Then
\begin{equation}
 \begin{pmatrix}u_{ab}^{(2)}\\[2pt]u_{ab}^{(3)}\end{pmatrix}
 =
 \begin{pmatrix}4&1\\9&1\end{pmatrix}
 \begin{pmatrix}b_{ab}\\c_{ab}\end{pmatrix},
 \qquad
 \det=-5,
 \label{eq:v32-traceful-J}
\end{equation}
so
\begin{equation}
 \boxed{
 b_{ab}=\frac{u_{ab}^{(3)}-u_{ab}^{(2)}}5,
 \qquad
 c_{ab}=\frac{9u_{ab}^{(2)}-4u_{ab}^{(3)}}5.}
 \label{eq:v32-traceful-inverse}
\end{equation}
This gives the requested explicit $HF$ inversion with
$(m_a,m_b)=(m_h,m_F)$ and the $FF$ inversion with
$(m_a,m_b)=(m_F,m_F)$.  In the $FF$ row one may use two orthogonal internal
labels of the selected massive fermion sector to isolate the direct channel.

A third forward point $z_{ab}^{(4)}=4m_am_b$ provides the internal polynomial
check
\begin{equation}
 \boxed{
 5u_{ab}^{(4)}-12u_{ab}^{(3)}+7u_{ab}^{(2)}=0.}
 \label{eq:v32-forward-third-check}
\end{equation}
It is useful as a calibration check but is \emph{not} the V31 bank-closure
row; the latter is chosen nonforward below.

On the worked Burnett ray,
\begin{equation}
 \boxed{
 u_{ab}^{(2)}=73\kappa_0,
 \qquad
 u_{ab}^{(3)}=143\kappa_0,
 \qquad
 u_{ab}^{(4)}=241\kappa_0.}
 \label{eq:v32-burnett-traceful-rows}
\end{equation}
Thus every traceful pair separately obeys the direct observable relation
\begin{equation}
 \boxed{
 73u_{ab}^{(3)}-143u_{ab}^{(2)}=0,}
 \label{eq:v32-burnett-traceful-relation}
\end{equation}
and, more stringently, the values of $\kappa_0$ extracted from $HH$, $HF$ and
$FF$ must agree.

\subsection{Explicit $HV$, $FV$ and $VV$ forward rows}
\label{sec:v32-gauge-inversion}

For $aV$ with $a=H$ or $F$, choose
\begin{equation}
 z_{aV}=m_a^2,
 \qquad s=3m_a^2.
 \label{eq:v32-aV-fit}
\end{equation}
Then
\begin{equation}
 \boxed{
 b_{aV}
 =\frac{\mathfrak A_{aV}^{\rm fwd}(3m_a^2)}{m_a^4}.}
 \label{eq:v32-aV-inverse}
\end{equation}
For $VV$, use the P1/reference scale to choose
\begin{equation}
 s=2\mu_0^2,
 \qquad z_{VV}=\mu_0^2,
 \label{eq:v32-VV-fit}
\end{equation}
and obtain
\begin{equation}
 \boxed{
 b_{VV}
 =\frac{\mathfrak A_{VV}^{\rm fwd}(2\mu_0^2)}{\mu_0^4}.}
 \label{eq:v32-VV-inverse}
\end{equation}
A second forward scale with $z\mapsto2z$ must scale the residual by four;
again this is a useful contact-scaling check rather than the nonforward closure
row.

The worked Burnett branch predicts the common gauge-containing normalization
\begin{equation}
 \boxed{
 \frac{\mathfrak A_{HV}^{\rm fwd}}{z_{HV}^2}
 =\frac{\mathfrak A_{FV}^{\rm fwd}}{z_{FV}^2}
 =\frac{\mathfrak A_{VV}^{\rm fwd}}{z_{VV}^2}
 =14\kappa_0.}
 \label{eq:v32-gauge-universality}
\end{equation}
Together with \eqref{eq:v32-traceful-inverse}, this populates all five new
species-pair coefficients of V31 from explicit forward contact rows.

\subsection{A concrete broken-phase benchmark panel}
\label{sec:v32-benchmark-panel}

For definiteness one may take $H$ to be the physical Higgs, $F$ a selected
massive quark (for example the top on a reference panel above threshold), and
$V$ a gluon.  Use distinct colour labels in the $FF$ and $VV$ direct rows.
The fit rows are
\begin{center}
\begin{tabular}{p{0.13\textwidth}p{0.28\textwidth}p{0.45\textwidth}}
Pair & forward fit kinematics & extracted coordinates\\ \hline
$HF$ & $z=2m_hm_F,\ 3m_hm_F$ & $b_{HF},c_{HF}$ from \eqref{eq:v32-traceful-inverse}\\
$FF$ & $z=2m_F^2,\ 3m_F^2$ & $b_{FF},c_{FF}$ from \eqref{eq:v32-traceful-inverse}\\
$HV$ & $s=3m_h^2$ & $b_{HV}$ from \eqref{eq:v32-aV-inverse}\\
$FV$ & $s=3m_F^2$ & $b_{FV}$ from \eqref{eq:v32-aV-inverse}\\
$VV$ & $s=2\mu_0^2$ & $b_{VV}$ from \eqref{eq:v32-VV-inverse}
\end{tabular}
\end{center}
All rows mean the pole-subtracted contact functional
\eqref{eq:v32-forward-functional}.  Another massive fermion and another
massless gauge species can be substituted without changing the compiler.

Combining these five extractions with the V30 $HH$ inversion produces the
complete V31 coefficient vector
\begin{equation}
 \boxed{
 y=(b_{HH},b_{HF},b_{HV},b_{FF},b_{FV},b_{VV};
 c_{HH},c_{HF},c_{FF})\in\mathbb R^9.}
 \label{eq:v32-populated-y}
\end{equation}
The eight diagnostics $r_2,r_0,\Delta_{\rm B}$ are then evaluated exactly by
\eqref{eq:v31-species-residuals}--\eqref{eq:v31-ratio-residual}.

\subsection{Nonforward held-out rows from stress form factors}
\label{sec:v32-nonforward}

Forward fitting alone cannot certify the local bank.  For a generic on-shell
row $r$ define the renormalized one-particle stress form factor
\begin{equation}
 F_a^{\mu\nu}[r]
 :=\langle a(p_a',\sigma_a')|T_a^{\mu\nu}(0)|a(p_a,\sigma_a)\rangle,
 \qquad
 F_a[r]:=g_{\mu\nu}F_a^{\mu\nu}[r].
 \label{eq:v32-stress-form-factor}
\end{equation}
The V31 contact operators generate the exactly computable row functions
\begin{align}
 J_{ab}^{(2)}[r]
 &:=\frac18\,
 \mathcal P_{ab}^{\rm dir}
 \langle ab|\mathcal O_{ab}^{(2)}|ab\rangle_r,
 \label{eq:v32-J2}\\
 J_{ab}^{(0)}[r]
 &:=\frac18\,
 \mathcal P_{ab}^{\rm dir}
 \langle ab|\mathcal O_{ab}^{(0)}|ab\rangle_r.
 \label{eq:v32-J0}
\end{align}
At forward kinematics these reduce to $z_{ab}^2$ and $m_a^2m_b^2$ by
Theorem~\ref{thm:v32-forward-pair}.  Away from the forward limit they retain
the spin, helicity and improvement information needed for a genuine bank test.

A concrete held-out choice is obtained at the larger fit energy for each pair,
with centre-of-mass scattering angle
\begin{equation}
 \boxed{\theta=\frac\pi2.}
 \label{eq:v32-heldout-angle}
\end{equation}
For $HF$ and $FF$ use the $z=3m_am_b$ energy; for $HV$ and $FV$ use the
energy in \eqref{eq:v32-aV-fit}; for $VV$ use \eqref{eq:v32-VV-fit}.  Fix one
helicity-preserving fermion row and one transverse gauge polarization row,
and keep the same row definition at every selector setting.  The momenta are
fully fixed by $s$, the masses and \eqref{eq:v32-heldout-angle}; for unequal
masses the centre-of-mass momentum is
\begin{equation}
 p_*^2=\frac{\lambda(s,m_a^2,m_b^2)}{4s},
 \qquad
 t=-2p_*^2,
 \label{eq:v32-heldout-t}
\end{equation}
with $\lambda$ the K\"all\'en polynomial.

For a traceful pair define
\begin{equation}
 \boxed{
 \Delta_{ab}^{\perp}
 :=\mathfrak A_{ab}^{\rm nf}
 -J_{ab}^{(2)}[r_\perp]\widehat b_{ab}
 -J_{ab}^{(0)}[r_\perp]\widehat c_{ab},}
 \label{eq:v32-nf-res-traceful}
\end{equation}
and for a gauge-containing pair
\begin{equation}
 \boxed{
 \Delta_{aV}^{\perp}
 :=\mathfrak A_{aV}^{\rm nf}
 -J_{aV}^{(2)}[r_\perp]\widehat b_{aV}.}
 \label{eq:v32-nf-res-gauge}
\end{equation}
These are the six V31 held-out rows, with the $HH$ row already supplied by V30.
A nonzero value beyond propagated uncertainty falsifies the stress-only bank
on that physical channel.

\begin{firewall}[Why the held-out row is nonforward]
\label{fw:v32-nonforward-needed}
A local derivative operator can vanish at $t=0$ or have the same forward
polynomial as a stress contact while becoming independent at $t\neq0$.
Therefore the optional forward scaling checks above do not replace
\eqref{eq:v32-nf-res-traceful}--\eqref{eq:v32-nf-res-gauge}.  At least one
nonforward row per pair is load bearing for the V31 closure claim.
\end{firewall}

\subsection{Source-minimal amplitude innovation if a row fails}
\label{sec:v32-amplitude-innovation}

There is no need to guess a complete dimension-eight SMEFT basis before using a
failed held-out row.  On any finite enlarged row panel let $A_{ab}$ be the
selector-subtracted amplitude vector and let $J_{ab}$ be the matrix whose
columns are the currently retained stress-bank amplitude functions.  Define
\begin{equation}
 P_{ab}:=J_{ab}J_{ab}^{\dagger},
 \qquad
 r_{ab}:=(I-P_{ab})A_{ab}.
 \label{eq:v32-pair-innovation}
\end{equation}

\begin{theorem}[Finite source-minimal amplitude innovation]
\label{thm:v32-amplitude-innovation}
If $r_{ab}\neq0$, then the normalized vector
\begin{equation}
 \boxed{
 f_{ab}^{\rm new}:=\frac{r_{ab}}{\|r_{ab}\|}}
 \label{eq:v32-new-amplitude-direction}
\end{equation}
is orthogonal to every retained stress amplitude column and spans the unique
rank-one enlargement required by the measured residual on the current row
panel.  Appending this amplitude direction makes the measured residual vanish
on that panel.  Identifying a local EFT operator whose on-shell amplitude maps
to $f_{ab}^{\rm new}$ is a subsequent operator-basis problem and is not assumed
here.
\end{theorem}

\begin{proof}
By construction $r_{ab}\in\operatorname{Ran}(I-P_{ab})$ and is orthogonal to
$\operatorname{Ran}J_{ab}$.  The orthogonal projection of $A_{ab}$ onto
$\operatorname{Ran}J_{ab}\oplus\operatorname{span}\{r_{ab}\}$ is therefore
exactly $A_{ab}$.  No smaller nonzero enlargement can absorb a nonzero vector
orthogonal to the previous range.
\end{proof}

This is deliberately an amplitude-space statement.  It prevents a failed
stress-bank row from being ``explained'' by arbitrarily choosing one operator
from a large dimension-eight catalogue before the data specify which on-shell
direction is actually missing.

\subsection{Explicit Burnett diagnostics on the populated row panel}
\label{sec:v32-burnett-diagnostics}

When all six nonforward held-out rows pass, the nine coefficients extracted by
V30 and \eqref{eq:v32-traceful-inverse}--\eqref{eq:v32-VV-inverse} can be
inserted directly into Theorem~\ref{thm:v31-universality-decomposition}.  The
worked microscopic branch requires
\begin{equation}
 \boxed{
 r_2=0,\qquad r_0=0,\qquad \Delta_{\rm B}=0.}
 \label{eq:v32-eight-tests}
\end{equation}
Equivalently, define the pairwise scale estimators
\begin{align}
 \widehat\kappa_{ab}^{(2)}&:=\frac{\widehat b_{ab}}{14},
 &&ab\in\{HH,HF,HV,FF,FV,VV\},
 \label{eq:v32-kappa-tensor}\\
 \widehat\kappa_{ab}^{(0)}&:=\frac{\widehat c_{ab}}{17},
 &&ab\in\{HH,HF,FF\}.
 \label{eq:v32-kappa-trace}
\end{align}
The Burnett ray is exactly the statement that all nine estimates coincide:
\begin{equation}
 \boxed{
 \widehat\kappa_{HH}^{(2)}=\widehat\kappa_{HF}^{(2)}
 =\widehat\kappa_{HV}^{(2)}=\widehat\kappa_{FF}^{(2)}
 =\widehat\kappa_{FV}^{(2)}=\widehat\kappa_{VV}^{(2)}
 =\widehat\kappa_{HH}^{(0)}=\widehat\kappa_{HF}^{(0)}
 =\widehat\kappa_{FF}^{(0)}.}
 \label{eq:v32-nine-kappa-equality}
\end{equation}
This representation is often more transparent for data analysis than the
orthogonal residual coordinates, while being algebraically equivalent to
$r_2=r_0=\Delta_{\rm B}=0$.

\subsection{What V32 does and does not populate}
\label{sec:v32-scope}

The row locations, physical projections and inversion formulas are now fully
specified.  Actual numerical values of the nine coefficients still require the
corresponding selector-subtracted amplitudes; no such numerical multi-species
amplitude data are present in the current corpus, and none are fabricated in
this continuation.  V32 therefore closes the \emph{tomographic design} and the
finite algebra needed to turn those amplitudes into $y$, but not the empirical
or microscopic population of the twelve new rows.

The translation-Ward forward fit has two practical advantages.  First, the fit
Jacobian is insensitive to external spin and to the Higgs improvement term.
Second, all spin/helicity dependence is moved into the held-out nonforward
rows, where it acts as a bank-completeness test rather than contaminating the
coefficient inversion.

\section{V32 status ledger}
\label{sec:v32-status}

\subsection*{Proved in V32}
\begin{enumerate}[leftmargin=2.2em]
\item After subtraction of known nonlocal poles and lower-weight pieces, the
translation Ward identity gives the universal forward pair formulas
\eqref{eq:v32-forward-traceful}--\eqref{eq:v32-forward-gauge} for every species
in the V31 stress bank.
\item The $HF$ and $FF$ tensor/trace coefficients are extracted explicitly
from the two rows $z/(m_am_b)=2,3$ by
\eqref{eq:v32-traceful-inverse}; the fit determinant is $-5$.
\item The $HV$, $FV$ and $VV$ tensor coefficients are extracted from one
explicit forward row each by \eqref{eq:v32-aV-inverse} and
\eqref{eq:v32-VV-inverse}.
\item Together with V30, these formulas populate the full nine-coordinate
coefficient vector $y$ once the twelve new physical amplitude rows are supplied.
\item The six held-out rows can be taken at $90^\circ$ nonforward kinematics
and are computed from the renormalized stress form factors.  A failed row is a
finite falsification of the stress-only bank on that channel.
\item A failed finite panel has a canonical source-minimal amplitude innovation
$f_{ab}^{\rm new}$, so no arbitrary dimension-eight operator need be selected
before the missing on-shell direction is identified.
\item On the worked Burnett branch, every traceful pair has normalized forward
rows $(73,143,241)\kappa_0$ at $z/(m_am_b)=(2,3,4)$, every gauge-containing pair
has normalized tensor coefficient $14\kappa_0$, and the complete nine-coordinate
test is the equality \eqref{eq:v32-nine-kappa-equality}.
\end{enumerate}

\subsection*{Inherited}
\begin{enumerate}[leftmargin=2.2em]
\item V29: the field-redefinition/EOM-reduced worked branch gives
$(\kappa_1,\kappa_2)=\kappa_0(14,17)$.
\item V30: three Higgs rows determine $b_{HH},c_{HH}$ and test the scalar
stress bank.
\item V31: the multi-species stress bank is nine-dimensional and its Burnett
ray has eight transverse diagnostics $(r_2,r_0,\Delta_{\rm B})$.
\end{enumerate}

\subsection*{Conditional}
\begin{enumerate}[leftmargin=2.2em]
\item The forward limits are taken only after the declared pole/reference
subtraction.  Without that subtraction they are not local Wilson projectors.
\item The direct $FF$ and $VV$ rows use orthogonal internal labels or an
explicit crossed-contraction subtraction.
\item The Higgs improvement convention is fixed for nonforward rows.  It drops
out of the forward fit but not from the held-out form factor.
\item The finite row panel is not yet a complete dimension-eight SMEFT basis
theorem; an operator invisible on all selected rows remains outside its reach.
\end{enumerate}

\subsection*{Open beyond V32}
\begin{enumerate}[leftmargin=2.2em]
\item Supply the twelve selector-subtracted physical amplitudes specified in
\S\ref{sec:v32-benchmark-panel} and \S\ref{sec:v32-nonforward}; assemble the
numerical nine-vector $y$ and covariance.
\item If any nonforward residual is nonzero, sample enough additional rows to
stabilize the amplitude innovation rank and then identify local operator
preimages for the measured innovation range.
\item If the finite stress bank closes, enumerate additional microscopic
selectors and require zero transverse image on the nine-coordinate physical
bank before promoting the branch relation.
\item Only after the tree amplitude bank is closed should the one-loop mixing
innovation of V31 be computed and the Burnett conormal transported across
scales.
\end{enumerate}

\subsection*{Exact next load-bearing theorem}
\begin{center}
\fbox{\parbox{0.92\textwidth}{
\textbf{P4p: amplitude-innovation closure before RG.}
Populate the twelve explicit V32 rows.  Use the five forward fit blocks to
construct $y=(b_2,b_0)$ and the six nonforward held-out rows to test the
stress-bank image.  If a held-out residual is nonzero, enlarge the row panel
and compute the stable finite innovation space
$(I-P_{\rm stress})\mathcal A$ before assigning local dimension-eight operator
names.  If the innovation space vanishes, evaluate
$(r_2,r_0,\Delta_{\rm B})$ and the nine pairwise $\widehat\kappa_0$ estimates.
Only after those tests pass should one construct the smallest local operator
bank whose on-shell image contains the closed amplitude range and begin the
one-loop mixing-closure calculation.}}
\end{center}

\section*{Append-only continuation marker: V32}
\addcontentsline{toc}{section}{Append-only continuation marker: V32}
\textbf{End of V32.}  The complete V1--V31 body above is retained verbatim.
V32 resolves the explicit multi-species tomography requested by V31 at the
level supported by the current corpus.  Pole-subtracted forward stress Ward
identities give universal, spin-independent fit rows and closed-form inversions
for $HF,HV,FF,FV,VV$.  Genuine bank closure is moved to nonforward held-out
rows so derivative nuisance operators cannot hide in the forward limit.  A
failed row generates a canonical source-minimal amplitude innovation rather
than an arbitrarily chosen EFT operator.  Numerical population of the twelve
new amplitude rows remains an acquisition problem and is not fabricated.



% =====================================================================
% Version V33: upstream completion audit for the finite-window gravity ray
% =====================================================================
\section{Version V33: upstream completion audit, conormal-first van Vleck gate, and completion-stable amplitude tests}
\label{sec:v33}

\subsection{Why V32 must pause before numerical row population}
\label{sec:v33-upstream-pause}

V32 left the twelve mixed-species amplitudes as an acquisition problem.  Before
using the worked isotropic finite-window ray to generate those rows, one further
upstream distinction is load bearing.  The detailed finite-window calculation
which produced the rational direction
\begin{equation}
 (c_1,c_2,c_3)_{\rm vol}
 =a_4\left(\frac54,\frac53,-\frac14\right)
 \label{eq:v33-volume-ray}
\end{equation}
was explicitly the \emph{Riemann-normal-coordinate volume-element} contribution.
Its own scope statement keeps the van Vleck/Gaussian-prefactor contribution as a
separate additive curvature-squared sector.  Thus the V28--V32 ratios are exact
for the imported volume sector, but the current corpus does not yet justify
identifying that sector with the complete four-derivative coefficient vector of
the worked channel.

A second source audit closes a possible basis ambiguity.  The detailed worked
calculation writes \eqref{eq:v33-volume-ray} in the Stelle-type basis
\begin{equation}
 \boxed{
 \mathcal L_{4\partial}
 =c_1R^2+c_2R_{\mu\nu}R^{\mu\nu}+c_3C_{\mu\nu\rho\sigma}C^{\mu\nu\rho\sigma}.}
 \label{eq:v33-basis}
\end{equation}
Therefore the V28 Gauss--Bonnet projection itself is retained.  What V33
corrects is the \emph{completion status} of the downstream phenomenological
ray, not the V28 basis algebra.

Write the complete pure-gravity coefficient vector at this order as
\begin{equation}
 c^{\rm full}=c^{\rm vol}+\delta c,
 \qquad
 \delta c=(\delta c_1,\delta c_2,\delta c_3)^T,
 \label{eq:v33-full-c}
\end{equation}
where $\delta c$ contains the uncomputed van Vleck/propagator contribution and,
if the geometric construction is later changed, any additional same-order
pure-gravity completion.  On the fixed geodesic-ball branch the immediate
upstream target is the van Vleck part of $\delta c$.

\begin{firewall}[Sector relation versus complete worked-channel relation]
\label{fw:v33-sector-completion}
The volume ray \eqref{eq:v33-volume-ray} is populated microscopic data.  It is
legitimate to derive its field-redefinition and amplitude descendants.  It is
not legitimate to set $\delta c=0$ merely because the additive sector has not
been evaluated.  Accordingly every V29--V32 numerical ratio inherited from
\eqref{eq:v33-volume-ray} is henceforth read as a \emph{volume-sector null}
until the completion obstruction below is shown to vanish.
\end{firewall}

\subsection{A single scalar controls survival of the V28 curvature conormal}
\label{sec:v33-curvature-obstruction}

Modulo the four-dimensional Euler density, V28 used
\begin{equation}
 c_R=c_1+\frac13c_2,
 \qquad
 c_C=c_3+\frac12c_2.
 \label{eq:v33-GB-map}
\end{equation}
The volume part obeys
\begin{equation}
 65c_C^{\rm vol}-21c_R^{\rm vol}=0.
 \label{eq:v33-volume-curvature-null}
\end{equation}
Define the additive completion obstruction
\begin{equation}
 \boxed{
 \Xi_{\rm add}
 :=65\,\delta c_C-21\,\delta c_R
 =-21\,\delta c_1+\frac{51}{2}\,\delta c_2+65\,\delta c_3.}
 \label{eq:v33-Xi-def}
\end{equation}

\begin{theorem}[Conormal-first completion theorem]
\label{thm:v33-conormal-first}
For the complete pure-gravity coefficient vector \eqref{eq:v33-full-c},
\begin{equation}
 \boxed{
 65c_C^{\rm full}-21c_R^{\rm full}=\Xi_{\rm add}.}
 \label{eq:v33-full-curvature-null}
\end{equation}
Consequently the V28 volume-sector zero locus survives the additive completion
if and only if
\begin{equation}
 \boxed{\Xi_{\rm add}=0.}
 \label{eq:v33-Xi-zero}
\end{equation}
Determining the entire vector $\delta c$ is therefore unnecessary if the first
question is only whether the already-populated conormal survives.
\end{theorem}

\begin{proof}
Insert \eqref{eq:v33-full-c} into \eqref{eq:v33-GB-map}.  The volume part
annihilates the V28 conormal by \eqref{eq:v33-volume-curvature-null}; the
remainder is precisely \eqref{eq:v33-Xi-def}.
\end{proof}

The gate \eqref{eq:v33-Xi-zero} is weaker than requiring the additive vector to
vanish.  Any additive correction tangent to the same physical curvature ray is
harmless for the relation, and more generally any $\delta c$ in the kernel of
the displayed conormal preserves it.

\subsection{The same scalar controls the physical stress-contact descendant}
\label{sec:v33-contact-obstruction}

V29 maps the pure-gravity curvature block to the matter-contact quotient
\begin{equation}
 \Delta\mathcal L_{\rm phys}
 \simeq
 \kappa_1T_{\mu\nu}T^{\mu\nu}+\kappa_2T^2.
 \label{eq:v33-contact-L}
\end{equation}
Using \eqref{eq:v33-GB-map}, the exact coefficient map is
\begin{equation}
 \boxed{
 \kappa_1=\frac{c_2+2c_3}{M_g^4},
 \qquad
 \kappa_2=\frac{c_1-\frac23c_3}{M_g^4}.}
 \label{eq:v33-contact-map}
\end{equation}
For the volume ray this reproduces V29,
\begin{equation}
 (\kappa_1,\kappa_2)_{\rm vol}
 =\frac{a_4}{12M_g^4}(14,17).
 \label{eq:v33-volume-contact-ray}
\end{equation}

\begin{theorem}[Physical completion obstruction]
\label{thm:v33-contact-obstruction}
For the complete pure-gravity sector,
\begin{equation}
 \boxed{
 17\kappa_1^{\rm full}-14\kappa_2^{\rm full}
 =\frac{2}{3M_g^4}\,\Xi_{\rm add}.}
 \label{eq:v33-contact-Xi}
\end{equation}
Hence the V29 $17{:}14$ contact relation is a complete worked-channel statement
at this order if and only if the additive conormal obstruction vanishes.
\end{theorem}

\begin{proof}
By \eqref{eq:v33-contact-map},
\begin{align}
 17\delta\kappa_1-14\delta\kappa_2
 &=\frac1{M_g^4}
 \left(
 -14\delta c_1+17\delta c_2+\frac{130}{3}\delta c_3
 \right)\\
 &=\frac{2}{3M_g^4}
 \left(
 -21\delta c_1+\frac{51}{2}\delta c_2+65\delta c_3
 \right),
\end{align}
which is \eqref{eq:v33-contact-Xi}.  The volume contribution itself annihilates
$17\kappa_1-14\kappa_2$.
\end{proof}

This identity is the useful upstream compression: the open propagator sector may
contain three curvature coordinates, but only one scalar combination can spoil
the already-derived volume-sector phenomenological ratio.

\subsection{Propagation to the V30 scalar amplitude panel}
\label{sec:v33-scalar-amplitude}

Retain the V30 definitions
\begin{equation}
 \Sigma:=\kappa_1+\kappa_2,
 \qquad
 \Pi:=12\kappa_1+60\kappa_2,
 \label{eq:v33-Sigma-Pi}
\end{equation}
for the pole-subtracted four-Higgs contact amplitude
\begin{equation}
 \widehat{\mathcal M}_{4h}=\Sigma X+\Pi m_h^4.
 \label{eq:v33-higgs-amp}
\end{equation}
A direct elimination gives
\begin{equation}
 31\Pi-1188\Sigma
 =-48(17\kappa_1-14\kappa_2).
 \label{eq:v33-Higgs-conormal-factor}
\end{equation}
Combining with Theorem~\ref{thm:v33-contact-obstruction} yields
\begin{equation}
 \boxed{
 31\Pi-1188\Sigma
 =-\frac{32}{M_g^4}\,\Xi_{\rm add}.}
 \label{eq:v33-Higgs-Xi}
\end{equation}
Thus the V30 integer amplitude relations are exact tests of the \emph{volume
sector}; for the complete worked channel they become null tests precisely after
\eqref{eq:v33-Xi-zero} is closed.

\begin{remark}[Why the V30 visibility calculation remains useful]
The massive-scalar panel still has full rank on $(\kappa_1,\kappa_2)$.  V33 does
not undo that result.  It changes only the microscopic image being tested: the
volume sector supplies one known line, while the unevaluated additive sector may
move the complete pure-gravity point within the same two-dimensional visible
contact plane.
\end{remark}

\subsection{Seven completion-stable species tests and one completion-sensitive ratio}
\label{sec:v33-multispecies-split}

The strongest salvage from V31--V32 is that a \emph{pure-gravity} additive
curvature-squared completion still enters matter through the total stress tensor.
Therefore it cannot distinguish the $H,F,V$ species labels introduced in V31.
For arbitrary complete pure-gravity coefficients $\kappa_1,\kappa_2$ one has
\begin{equation}
 \boxed{
 b_2=\kappa_1\mathbf 1_6,
 \qquad
 b_0=\kappa_2\mathbf 1_3.}
 \label{eq:v33-universal-completion}
\end{equation}
Consequently the V31 species-nonuniversality residuals obey
\begin{equation}
 \boxed{r_2=0,\qquad r_0=0}
 \label{eq:v33-stable-r20}
\end{equation}
for \emph{every} pure-gravity completion of the declared type, not just for the
volume ray.

By contrast, the V31 Burnett-ratio residual becomes
\begin{equation}
 \boxed{
 \Delta_{\rm B}
 :=17\bar b_2-14\bar b_0
 =\frac{2}{3M_g^4}\,\Xi_{\rm add}.}
 \label{eq:v33-DeltaB-Xi}
\end{equation}
Thus the eight V31 transverse diagnostics split canonically into
\begin{enumerate}[label=\textup{(\roman*)},leftmargin=2.4em]
\item five tensor-sector species-universality constraints, stable under any
pure-gravity additive completion;
\item two trace-sector species-universality constraints, likewise stable;
\item one $17{:}14$ volume-ray ratio constraint, which is completion-sensitive
and equivalent to $\Xi_{\rm add}=0$.
\end{enumerate}

\begin{firewall}[What can break the seven stable tests]
\label{fw:v33-stable-scope}
Equation~\eqref{eq:v33-stable-r20} is stable under additional
\emph{pure-gravity} curvature-squared contributions because those reduce to the
same total-stress bilinears.  It is not protected against genuinely mixed or
matter-specific dimension-eight selectors.  Such selectors remain exactly the
$\chi_*$ directions of V27 and must be tested by the mixed-species panel.
\end{firewall}

The V32 pairwise scale equalities should therefore be split correspondingly.
The completion-stable content is
\begin{equation}
 \boxed{
 b_{HH}=b_{HF}=b_{HV}=b_{FF}=b_{FV}=b_{VV},
 \qquad
 c_{HH}=c_{HF}=c_{FF}.}
 \label{eq:v33-stable-pair-equalities}
\end{equation}
The stronger statement
\begin{equation}
 \frac{b_{ab}}{14}=\frac{c_{a'b'}}{17}
 \label{eq:v33-volume-normalized-equality}
\end{equation}
for all tensor and trace pairs is a volume-sector statement until
$\Xi_{\rm add}=0$ is proved.

\subsection{Conormal-first van Vleck acquisition}
\label{sec:v33-vv-acquisition}

The detailed upstream construction identifies the missing sector as the quartic
part of the Gaussian prefactor, equivalently the same-profile average of
$\log\Delta^{1/2}$.  Define its curvature-squared coefficient by
\begin{equation}
 \mathcal V_4(p,u;\rho)
 :=[\tau^4]\int_{D_\tau(p,u)}
 \log\Delta(p,q)^{1/2}\,\rho_\tau(q)\,dq,
 \label{eq:v33-V4-def}
\end{equation}
where $[\tau^4]$ denotes the quartic term after the already-calibrated linear
Ricci response and total derivatives are removed.  After the same axis average
and covariant basis projection used for the volume term, write
\begin{equation}
 \Pi_{4\partial}\langle\mathcal V_4\rangle_u
 =\delta c_1^{\rm VV}R^2
 +\delta c_2^{\rm VV}R_{\mu\nu}R^{\mu\nu}
 +\delta c_3^{\rm VV}C^2.
 \label{eq:v33-VV-triple}
\end{equation}
The single load-bearing number for survival of the V28--V32 volume conormal is
then
\begin{equation}
 \boxed{
 \Xi_{\rm VV}
 =-21\delta c_1^{\rm VV}
 +\frac{51}{2}\delta c_2^{\rm VV}
 +65\delta c_3^{\rm VV}.}
 \label{eq:v33-Xi-VV}
\end{equation}

\begin{proposition}[Conormal-first economy]
\label{prop:v33-conormal-economy}
On the fixed geodesic-ball worked branch, deciding whether the populated volume
relation survives addition of the presently open van Vleck sector requires only
$\Xi_{\rm VV}$, not the complete triple
$(\delta c_1^{\rm VV},\delta c_2^{\rm VV},\delta c_3^{\rm VV})$.  If
\begin{equation}
 \Xi_{\rm VV}=0,
 \label{eq:v33-VV-pass}
\end{equation}
then the V28 curvature conormal, the V29 contact conormal, and the V30--V32
ratio nulls all survive the propagator completion at the retained order.  If
$\Xi_{\rm VV}\neq0$, the volume ray is not the complete worked-channel ray and
the full additive coefficient vector must be evaluated before transporting a
replacement relation.
\end{proposition}

\begin{proof}
On the fixed construction the immediate open contribution is precisely the van
Vleck sector.  Apply Theorems~\ref{thm:v33-conormal-first} and
\ref{thm:v33-contact-obstruction}, followed by
\eqref{eq:v33-Higgs-Xi} and \eqref{eq:v33-DeltaB-Xi}.
\end{proof}

This is a stronger next target than ``populate twelve amplitudes''.  Those
amplitudes remain the correct empirical/microscopic validation panel, but the
current upstream theory first owes the scalar completion gate
\eqref{eq:v33-Xi-VV}.  Populating the twelve rows from the volume ray before that
gate would merely restate the assumed sector and could not test the missing
Gaussian-prefactor physics.

\subsection{V33 correction and inheritance ledger}
\label{sec:v33-correction-ledger}

The append-only lineage keeps V28--V32 verbatim, so the scope correction must be
explicit:
\begin{enumerate}[leftmargin=2.4em]
\item The detailed source audit confirms the V28 basis
$\{R^2,R_{\mu\nu}^2,C^2\}$ and the volume-sector ray
$(15,20,-3)$; no numerical basis correction is made.
\item The V28 relation $65c_C-21c_R=0$ remains exact for the imported
geodesic-volume sector.
\item The V29 contact ray $(14,17)$, the V30 scalar-amplitude conormals, and the
V31--V32 $17{:}14$ normalization tests are inherited as exact descendants of
that same sector.
\item They are \emph{not yet} complete worked-channel predictions because the
upstream finite-window note explicitly leaves an additive van Vleck/propagator
curvature-squared contribution open.
\item The seven species-universality constraints contained in $r_2=r_0=0$ are
stable under any additional pure-gravity curvature-squared completion; only the
single ratio coordinate $\Delta_{\rm B}$ is sensitive to the missing additive
pure-gravity sector.
\end{enumerate}

\section{V33 status ledger}
\label{sec:v33-status}

\subsection*{Proved in V33}
\begin{enumerate}[leftmargin=2.2em]
\item The detailed upstream worked-channel calculation is consistent with the
V28 Stelle-type basis $\{R^2,R_{\mu\nu}^2,C^2\}$ and supplies the rational ray
$(15,20,-3)$ only in the geodesic-volume sector.
\item Every additive pure-gravity completion enters the V28 conormal through the
single scalar $\Xi_{\rm add}$ in \eqref{eq:v33-Xi-def}.
\item The same scalar controls the field-redefinition-invariant contact ratio by
\eqref{eq:v33-contact-Xi}, the V30 scalar amplitude residual by
\eqref{eq:v33-Higgs-Xi}, and the V31 Burnett-ratio residual by
\eqref{eq:v33-DeltaB-Xi}.
\item The V31--V32 species-universality tests split into seven constraints stable
under arbitrary pure-gravity curvature-squared completion and one
completion-sensitive $17{:}14$ ratio test.
\item On the fixed geodesic-ball branch, the immediate upstream completion gate
is the single van Vleck projection $\Xi_{\rm VV}$ rather than the full twelve-row
mixed-species amplitude population.
\end{enumerate}

\subsection*{Inherited}
\begin{enumerate}[leftmargin=2.2em]
\item V28: populated geodesic-volume Wilson ray and its branch-specific scope.
\item V29: exact four-dimensional field-redefinition map from the curvature
coordinates to total-stress contact operators.
\item V30: full-rank massive-scalar visibility of $(\kappa_1,\kappa_2)$.
\item V31--V32: finite nine-coordinate species-resolved contact bank and explicit
forward/nonforward tomography.
\end{enumerate}

\subsection*{Conditional}
\begin{enumerate}[leftmargin=2.2em]
\item The completion-stable equalities \eqref{eq:v33-stable-pair-equalities}
assume the added sector is pure gravity before the field-redefinition quotient.
Matter-specific or mixed dimension-eight selectors can violate them.
\item Construction changes which alter the finite-window Weyl response contribute
to $\delta c$ and must be included in $\Xi_{\rm add}$ unless the geodesic-ball
branch is kept fixed.
\item The V28--V32 volume-sector ratio becomes a complete worked-channel
prediction only after all same-order additive pure-gravity sectors satisfy the
conormal gate.
\end{enumerate}

\subsection*{Open beyond V33}
\begin{enumerate}[leftmargin=2.2em]
\item Evaluate the quartic same-profile van Vleck functional
\eqref{eq:v33-V4-def} and compute the scalar obstruction $\Xi_{\rm VV}$.
\item If $\Xi_{\rm VV}\neq0$, evaluate the complete additive triple and replace
the volume ray by the true worked-channel physical contact ray.
\item If $\Xi_{\rm VV}=0$, return to the V32 twelve-row acquisition as an
independent validation of the completed pure-gravity prediction and as a search
for mixed/matter-specific nuisance operators.
\item Only after that completion/validation split is closed should the
one-loop dimension-eight mixing problem be placed on the critical path.
\end{enumerate}

\subsection*{Exact next load-bearing theorem}
\begin{center}
\fbox{\parbox{0.92\textwidth}{
\textbf{P4q: conormal-first van Vleck completion.}
Using the same finite-window profile, reference normalization and axis averaging
as the populated geodesic-volume calculation, carry the covariant expansion of
$\log\Delta^{1/2}$ through the quartic curvature order.  Remove the already
calibrated linear Ricci response and total derivatives, project the remaining
curvature-squared term onto $\{R^2,R_{\mu\nu}^2,C^2\}$, and compute directly
\[
 \Xi_{\rm VV}
 =-21\delta c_1^{\rm VV}
 +\frac{51}{2}\delta c_2^{\rm VV}
 +65\delta c_3^{\rm VV}.
\]
If $\Xi_{\rm VV}=0$, the V28--V32 volume conormal survives the propagator
completion and the twelve-row mixed-species panel becomes the next validation
gate.  If $\Xi_{\rm VV}\neq0$, compute the full additive triple and transport
the corrected physical contact ray before any phenomenological ratio or RG
calculation.}}
\end{center}

\section*{Append-only continuation marker: V33}
\addcontentsline{toc}{section}{Append-only continuation marker: V33}
\textbf{End of V33.}  The complete V1--V32 body above is retained verbatim.
V33 goes upstream of the unpopulated twelve-row amplitude acquisition and audits
the microscopic input that generated the Burnett ray.  The detailed finite-window
calculation confirms the V28 coefficient basis and rational geodesic-volume ray,
but explicitly leaves an additive van Vleck/propagator curvature-squared sector
open.  The effect of every pure-gravity additive completion on the V28--V32
phenomenology is compressed to the single scalar $\Xi_{\rm add}$.  Seven
multi-species universality tests remain stable under such a completion; the one
$17{:}14$ ratio test is equivalent to the vanishing of the transverse completion
obstruction.  The exact next computation is therefore the conormal-first
van Vleck scalar $\Xi_{\rm VV}$, not fabricated numerical mixed-species rows.


% =====================================================================
\section{V34 --- Conormal-first van Vleck completion: the quartic direction,
nonzero obstruction, and the corrected two-ray contact family}
\label{sec:v34-vv-completion}
% =====================================================================

V33 reduced the first phenomenological completion problem to the single scalar
\(\Xi_{\rm VV}\) obtained from the quartic same-profile average of
\(\log\Delta^{1/2}\).  This section closes that directional gate.  The result is
negative for the V28 volume conormal: the intrinsic quartic van Vleck direction
is not collinear with the geodesic-volume ray.  The remaining issue is only the
common finite-window normalization relating the two additive sectors.

\subsection{Standard quartic coincidence input and the profile fourth moment}
\label{subsec:v34-vv-input}

Write
\[
 \zeta(p,q):=\log\Delta(p,q)^{1/2}.
\]
In Riemann normal coordinates about \(p\), with tangent displacement
\(\sigma^a\), the standard covariant coincidence expansion through quartic
order is
\begin{equation}
\begin{split}
 \zeta(p,q)
 ={}&\frac1{12}R_{ab}\sigma^a\sigma^b
 +\frac1{24}R_{ab;c}\sigma^a\sigma^b\sigma^c\\
 &+\left[
   \frac1{360}R^{\beta}{}_{a\gamma b}
                    R^{\gamma}{}_{c\beta d}
   +\frac1{80}R_{ab;cd}
   \right]
   \sigma^a\sigma^b\sigma^c\sigma^d
 +O(|\sigma|^5).
\end{split}
\label{eq:v34-zeta-quartic}
\end{equation}
This is the external geometric input required by P4q; the renewal-specific
input remains the calibrated finite-window profile and its same-history
normalization.

After the same Wick-rotated axis average used by the populated volume
calculation, rotational invariance implies that the fourth profile moment has
the form
\begin{equation}
 \left\langle
 \int_{D_\tau(p,u)}
 \sigma^a\sigma^b\sigma^c\sigma^d\,
 \rho_{\tau,u}(q)\,dq
 \right\rangle_u
 =m_{4,\rho}\,\tau^4
 \left(g^{ab}g^{cd}+g^{ac}g^{bd}+g^{ad}g^{bc}\right),
\label{eq:v34-profile-fourth}
\end{equation}
where \(m_{4,\rho}>0\) for every nondegenerate positive profile.  The
normalization of \(m_{4,\rho}\) is left exactly as induced by the already chosen
finite-window profile; no new profile convention is introduced here.

\begin{lemma}[Derivative part is a total divergence]
\label{lem:v34-derivative-total}
Let
\[
 S^{abcd}:=g^{ab}g^{cd}+g^{ac}g^{bd}+g^{ad}g^{bc}.
\]
Then
\begin{equation}
 S^{abcd}R_{ab;cd}=2\,\Box R.
\label{eq:v34-derivative-contraction}
\end{equation}
Consequently the quartic \(R_{ab;cd}\) term in
\eqref{eq:v34-zeta-quartic} contributes only a total derivative to the local
four-derivative action quotient used in V28--V33.
\end{lemma}

\begin{proof}
The three metric contractions give respectively
\(\Box R\), \(\nabla^a\nabla^bR_{ab}\), and
\(\nabla^b\nabla^aR_{ab}\).  The contracted Bianchi identity
\(\nabla^aR_{ab}=\frac12\nabla_bR\) turns each of the latter two terms into
\(\frac12\Box R\).  Their sum is \(2\Box R\), which is removed by the same
total-divergence quotient already used for the volume sector.
\end{proof}

\subsection{Exact isotropic contraction of the curvature-square tensor}
\label{subsec:v34-vv-contraction}

Define
\begin{equation}
 T_{abcd}:=
 R^{\beta}{}_{a\gamma b}R^{\gamma}{}_{c\beta d}.
\label{eq:v34-T-def}
\end{equation}

\begin{lemma}[Three trace contractions]
\label{lem:v34-three-traces}
With the curvature convention of \eqref{eq:v34-zeta-quartic}, the Riemann
symmetries and first Bianchi identity give
\begin{align}
 g^{ab}g^{cd}T_{abcd}
 &=R_{\mu\nu}R^{\mu\nu},
 \label{eq:v34-trace-1}\\
 g^{ac}g^{bd}T_{abcd}
 &=\frac12R_{\mu\nu\rho\sigma}R^{\mu\nu\rho\sigma},
 \label{eq:v34-trace-2}\\
 g^{ad}g^{bc}T_{abcd}
 &=R_{\mu\nu\rho\sigma}R^{\mu\nu\rho\sigma}.
 \label{eq:v34-trace-3}
\end{align}
Hence
\begin{equation}
 \boxed{
 S^{abcd}T_{abcd}
 =R_{\mu\nu}R^{\mu\nu}
 +\frac32R_{\mu\nu\rho\sigma}R^{\mu\nu\rho\sigma}.}
\label{eq:v34-isotropic-T}
\end{equation}
\end{lemma}

\begin{proof}
The first identity is the double Ricci trace.  For the second and third,
move the contracted indices using pair exchange and antisymmetry of the Riemann
tensor and apply the first Bianchi identity to the cyclic pair in the middle.
The resulting contractions are respectively one half and one full Riemann norm.
As an independent consistency check, on a four-dimensional constant-curvature
background the right-hand side of \eqref{eq:v34-isotropic-T} is
\(36K^2+\frac32(24K^2)=72K^2\); division by the unit-sphere fourth-moment
factor \(24\) gives \(3K^2\), exactly the quartic coefficient required by the
closed constant-curvature van Vleck determinant.
\end{proof}

In four dimensions,
\begin{equation}
 R_{\mu\nu\rho\sigma}R^{\mu\nu\rho\sigma}
 =C_{\mu\nu\rho\sigma}C^{\mu\nu\rho\sigma}
 +2R_{\mu\nu}R^{\mu\nu}-\frac13R^2.
\label{eq:v34-four-d-riemann}
\end{equation}
Combining \eqref{eq:v34-isotropic-T} and
\eqref{eq:v34-four-d-riemann} gives
\begin{equation}
 \boxed{
 S^{abcd}T_{abcd}
 =-\frac12R^2+4R_{\mu\nu}R^{\mu\nu}+\frac32C^2.}
\label{eq:v34-vv-direction-cont}
\end{equation}

\begin{theorem}[Quartic van Vleck curvature-square direction]
\label{thm:v34-vv-direction}
For the V33 same-profile functional \eqref{eq:v33-V4-def}, after removal of
the already calibrated quadratic Ricci response and total derivatives, the
quartic curvature-square contribution is
\begin{equation}
 \boxed{
 \Pi_{4\partial}\langle\mathcal V_4\rangle_u
 =\frac{m_{4,\rho}}{360}
 \left(
  -\frac12R^2+4R_{\mu\nu}R^{\mu\nu}+\frac32C^2
 \right).}
\label{eq:v34-V4-result}
\end{equation}
Thus, in the direct V33 profile-functional normalization,
\begin{equation}
 \boxed{
 \left(
 \delta c_1^{\rm VV},
 \delta c_2^{\rm VV},
 \delta c_3^{\rm VV}
 \right)
 =\left(
 -\frac{m_{4,\rho}}{720},
 \frac{m_{4,\rho}}{90},
 \frac{m_{4,\rho}}{240}
 \right)
 =\omega_{\rm VV}(-1,8,3),}
\label{eq:v34-vv-triple}
\end{equation}
where
\begin{equation}
 \omega_{\rm VV}:=\frac{m_{4,\rho}}{720}>0.
\label{eq:v34-omega-vv-direct}
\end{equation}
If the final common effective-action compiler carries an additional nonzero
sector normalization, it multiplies all three entries in
\eqref{eq:v34-vv-triple} by the same scalar and does not change the direction
or the nonvanishing statement below.
\end{theorem}

\begin{proof}
Insert \eqref{eq:v34-profile-fourth} into the quartic term of
\eqref{eq:v34-zeta-quartic}.  Lemma~\ref{lem:v34-derivative-total} removes the
derivative term in the physical local-action quotient.  Lemma
\ref{lem:v34-three-traces} contracts the remaining curvature-square tensor,
and \eqref{eq:v34-four-d-riemann} converts it to the V33 basis.  Reading off
the coefficients proves \eqref{eq:v34-vv-triple}.
\end{proof}

\subsection{P4q closes on the nonzero branch}
\label{subsec:v34-conormal-gate}

\begin{corollary}[Exact van Vleck conormal obstruction]
\label{cor:v34-Xi-vv}
The V33 conormal projection of the intrinsic quartic van Vleck sector is
\begin{equation}
 \boxed{
 \Xi_{\rm VV}
 =-21\delta c_1^{\rm VV}
 +\frac{51}{2}\delta c_2^{\rm VV}
 +65\delta c_3^{\rm VV}
 =\frac7{12}m_{4,\rho}
 =420\,\omega_{\rm VV}.}
\label{eq:v34-Xi-vv}
\end{equation}
For every nondegenerate positive profile,
\begin{equation}
 \boxed{\Xi_{\rm VV}\neq0.}
\label{eq:v34-Xi-vv-nonzero}
\end{equation}
Therefore the V28--V32 \(65{:}21\) curvature conormal and its
\(17{:}14\) contact descendant do \emph{not} survive the quartic van Vleck
completion as full worked-channel relations.
\end{corollary}

\begin{proof}
Substitution of \eqref{eq:v34-vv-triple} gives
\[
 \Xi_{\rm VV}
 =m_{4,\rho}
 \left(
 \frac{21}{720}+\frac{51}{180}+\frac{65}{240}
 \right)
 =\frac7{12}m_{4,\rho}.
\]
The fourth moment of a nondegenerate positive profile is strictly positive.
\end{proof}

\begin{firewall}[What this does and does not close]
\label{fire:v34-normalization}
V34 closes the \emph{directional} P4q gate and proves that the intrinsic quartic
van Vleck sector has a nonzero component transverse to the geodesic-volume
ray.  It does not yet identify the relative overall normalization with which
the volume and van Vleck pieces enter the single common finite-window effective
action.  Consequently the old volume-only ratio is ruled out as the complete
relation, but a new numerical completed ratio is not quoted until that one
relative normalization is fixed from common partition-function bookkeeping.
\end{firewall}

\subsection{Corrected pure-gravity family and physical contact quotient}
\label{subsec:v34-corrected-family}

Write the populated volume contribution in the integer normalization
\begin{equation}
 c^{\rm vol}
 =\omega_{\rm vol}(15,20,-3),
\label{eq:v34-volume-omega}
\end{equation}
and the van Vleck contribution as
\begin{equation}
 c^{\rm VV}
 =\omega_{\rm VV}(-1,8,3).
\label{eq:v34-vv-omega}
\end{equation}
Until the common normalization is fixed, the completed pure-gravity packet is
therefore the two-ray family
\begin{equation}
 \boxed{
 c^{\rm full}
 =\omega_{\rm vol}(15,20,-3)
 +\omega_{\rm VV}(-1,8,3).}
\label{eq:v34-full-two-ray}
\end{equation}
The two directions are linearly independent.

Applying the V28 map to \((c_C,c_R)\) gives
\begin{align}
 c_C^{\rm full}
 &=7(\omega_{\rm vol}+\omega_{\rm VV}),
 \label{eq:v34-cC-full}\\
 c_R^{\rm full}
 &=\frac53(13\omega_{\rm vol}+\omega_{\rm VV}).
 \label{eq:v34-cR-full}
\end{align}
The V29 physical field-redefinition quotient then gives
\begin{equation}
 \boxed{
 \kappa_1
 =\frac{14(\omega_{\rm vol}+\omega_{\rm VV})}{M_g^4},
 \qquad
 \kappa_2
 =\frac{17\omega_{\rm vol}-3\omega_{\rm VV}}{M_g^4}.}
\label{eq:v34-contact-full}
\end{equation}
In particular,
\begin{equation}
 \boxed{
 17\kappa_1-14\kappa_2
 =\frac{280\omega_{\rm VV}}{M_g^4}\neq0}
\label{eq:v34-old-contact-null-broken}
\end{equation}
on the nondegenerate van Vleck branch.

For the V30 scalar amplitude coordinates,
\begin{align}
 \Sigma
 &=\frac{31\omega_{\rm vol}+11\omega_{\rm VV}}{M_g^4},
 \label{eq:v34-Sigma-full}\\
 \Pi
 &=\frac{1188\omega_{\rm vol}-12\omega_{\rm VV}}{M_g^4},
 \label{eq:v34-Pi-full}
\end{align}
so the previous volume null becomes
\begin{equation}
 \boxed{
 31\Pi-1188\Sigma
 =-\frac{13440\omega_{\rm VV}}{M_g^4}.}
\label{eq:v34-old-amplitude-null-broken}
\end{equation}
These identities agree exactly with the V33 generic completion formula after
using \(\Xi_{\rm VV}=420\omega_{\rm VV}\).

\subsection{One scalar remains before the new completed phenomenological ray}
\label{subsec:v34-ratio-scalar}

Assume \(\omega_{\rm vol}\neq0\) and define the common-normalization ratio
\begin{equation}
 r_{\rm VV}:=\frac{\omega_{\rm VV}}{\omega_{\rm vol}}.
\label{eq:v34-rvv}
\end{equation}
Then the full curvature ray, once the common profile fixes \(r_{\rm VV}\), is
\begin{equation}
 \boxed{
 c^{\rm full}
 \propto
 (15-r_{\rm VV},\;20+8r_{\rm VV},\;-3+3r_{\rm VV}).}
\label{eq:v34-completed-ray-r}
\end{equation}
The corresponding physical contact ray is
\begin{equation}
 \boxed{
 (\kappa_1,\kappa_2)
 \propto
 \bigl(14(1+r_{\rm VV}),\;17-3r_{\rm VV}\bigr).}
\label{eq:v34-contact-ray-r}
\end{equation}
Hence its transported conormal is
\begin{equation}
 \boxed{
 (17-3r_{\rm VV})\kappa_1
 -14(1+r_{\rm VV})\kappa_2=0.}
\label{eq:v34-contact-conormal-r}
\end{equation}
Likewise the scalar-amplitude ray is
\begin{equation}
 (\Sigma,\Pi)
 \propto
 \bigl(31+11r_{\rm VV},\;1188-12r_{\rm VV}\bigr),
\label{eq:v34-amplitude-ray-r}
\end{equation}
with conormal
\begin{equation}
 \boxed{
 (31+11r_{\rm VV})\Pi
 -(1188-12r_{\rm VV})\Sigma=0.}
\label{eq:v34-amplitude-conormal-r}
\end{equation}
Thus V34 has reduced the unknown completion from a three-component Wilson triple
to a single relative normalization scalar.

\subsection{Completion-stable multi-species tests remain intact}
\label{subsec:v34-stable-tests}

Both terms in \eqref{eq:v34-full-two-ray} are pure-gravity before the V29
field-redefinition quotient.  Therefore they both descend to contacts built
from the \emph{total} stress tensor.  Consequently the seven V31--V33
species-universality conditions remain unchanged:
\begin{equation}
 \boxed{
 b_{HH}=b_{HF}=b_{HV}=b_{FF}=b_{FV}=b_{VV},
 \qquad
 c_{HH}=c_{HF}=c_{FF}.}
\label{eq:v34-stable-seven}
\end{equation}
Only the eighth transverse test---the tensor/trace ratio---is deformed by
\(r_{\rm VV}\).  Matter-specific or mixed microscopic selectors remain a
separate V26/V31 nuisance sector and can still violate
\eqref{eq:v34-stable-seven}.

\subsection{Status after V34}
\label{subsec:v34-status}

\subsection*{Proved}
\begin{enumerate}[leftmargin=2.2em]
\item The quartic covariant expansion of \(\log\Delta^{1/2}\) reduces, after
isotropic profile averaging and the physical total-derivative quotient, to the
curvature-square direction \((-1,8,3)\) in the
\((R^2,R_{\mu\nu}^2,C^2)\) basis.
\item The V33 conormal projection is nonzero:
\(\Xi_{\rm VV}=\frac7{12}m_{4,\rho}=420\omega_{\rm VV}\).
\item Therefore the V28 volume-only \(65{:}21\) relation and V29
\(17{:}14\) contact relation are not the complete worked-channel relations once
the intrinsic van Vleck quartic sector is included.
\item The complete pure-gravity correction is reduced to the two-ray family
\eqref{eq:v34-full-two-ray}; after common normalization it becomes the
one-parameter deformation \eqref{eq:v34-completed-ray-r}.
\item All seven species-universality tests remain invariant under this
pure-gravity completion.
\end{enumerate}

\subsection*{Inherited}
\begin{enumerate}[leftmargin=2.2em]
\item V28: exact geodesic-volume ray \((15,20,-3)\) for the worked isotropic
compound-Poisson branch.
\item V29: exact four-dimensional field-redefinition quotient to total-stress
contacts.
\item V30--V32: scalar and multi-species on-shell visibility/tomography.
\item V33: conormal-first reduction showing that only the transverse additive
projection is load-bearing for survival of the old ratio.
\end{enumerate}

\subsection*{Conditional}
\begin{enumerate}[leftmargin=2.2em]
\item Equations \eqref{eq:v34-completed-ray-r}--\eqref{eq:v34-amplitude-conormal-r}
become numerical predictions only after the single common finite-window ratio
\(r_{\rm VV}\) is fixed.
\item The sign \(\omega_{\rm VV}>0\) is the direct profile-functional
normalization of \eqref{eq:v33-V4-def}; a different common action convention may
multiply the entire sector by one nonzero scalar, without restoring the old
conormal.
\item The seven species-universality equalities assume that the added completion
is pure gravity before the EOM/field-redefinition quotient.
\end{enumerate}

\subsection*{Open beyond V34}
\begin{enumerate}[leftmargin=2.2em]
\item Fix the relative scale \(r_{\rm VV}\) by putting the geodesic-volume and
Gaussian-prefactor terms into one common flat-normalized diamond partition
function with one window/profile normalization.
\item After that scalar is fixed, update the V30--V32 explicit amplitude ray and
run the twelve mixed-species validation rows against the completed prediction.
\item Enumerate matter-specific or mixed fourth-order selectors and test their
transverse image against the seven completion-stable universality conditions.
\item Only after the physical bank closes should the one-loop dimension-eight
mixing calculation return to the critical path.
\end{enumerate}

\subsection*{Exact next load-bearing theorem}
\begin{center}
\fbox{\parbox{0.92\textwidth}{
\textbf{P4r: common-window normalization of the volume and van Vleck sectors.}
Start from one and the same flat-normalized finite-window renewal-diamond
partition function and one calibrated profile.  Track the numerical prefactors
with which (i) the quartic-cumulant/RNC-volume response and (ii) the quartic
\(\log\Delta^{1/2}\) response enter the local four-derivative effective action.
Determine the single ratio
\[
 r_{\rm VV}=\omega_{\rm VV}/\omega_{\rm vol}
\]
without separately refitting either Wilson packet.  Insert it into
\eqref{eq:v34-contact-ray-r} and \eqref{eq:v34-amplitude-ray-r}.  Only then
promote the deformed contact/amplitude conormal to the completed worked-channel
prediction and return to the V32 mixed-species validation panel.}}
\end{center}

\section*{Append-only continuation marker: V34}
\addcontentsline{toc}{section}{Append-only continuation marker: V34}
\textbf{End of V34.}  The complete V1--V33 body above is retained verbatim.
V34 carries the standard quartic covariant expansion of the van Vleck logarithm
through the same-profile isotropic projection required by P4q.  The derivative
part is a total divergence; the curvature-square tensor contracts to
\(-\tfrac12R^2+4R_{\mu\nu}^2+\tfrac32C^2\), giving the exact direction
\((-1,8,3)\) and the nonzero transverse obstruction
\(\Xi_{\rm VV}=\tfrac7{12}m_{4,\rho}\).  Hence the old volume-only ratio is not
the completed worked-channel relation.  The remaining completion uncertainty is
compressed to one common-normalization scalar \(r_{\rm VV}\); the exact next
step is to determine that scalar from one common finite-window partition
function before returning to mixed-species phenomenology.



% =============================================================================
% APPEND-ONLY CONTINUATION: VERSION V35
% The complete V1--V34 body above is retained verbatim.  V35 begins here.
% =============================================================================

\clearpage
\section{Version V35: Common-Window Audit, Gaussian Baseline versus Non-Gaussian Excess, and the Correct P4r Matching Tangent}
\label{sec:v35-common-window}

\subsection{Purpose and upstream correction}
\label{subsec:v35-purpose}

V34 solved a genuine geometric subproblem: the quartic coincidence expansion of
$\log\Delta^{1/2}$ was contracted against an isotropic fourth-moment tensor and
its curvature-square direction was found to be $(-1,8,3)$ in the
$(R^2,R_{\mu\nu}^2,C^2)$ basis.  The next proposed step was to determine a
single relative number $r_{\rm VV}$ and add that response directly to the
V28 geodesic-volume Wilson ray.

The common-window audit reveals that one further distinction is load-bearing.
The quartic van Vleck expression of V34 is a \emph{profile-averaged diamond
response}.  A local curvature-squared Wilson coefficient of the induced action
is instead the coefficient of the complete regularized determinant, after all
Gaussian transport, measure, recurrence/heat-kernel and subtraction terms have
been assembled in one convention.  The van Vleck factor is one factor of that
construction, not the complete coefficient.  Consequently the V34 tensor
contraction remains correct, whereas its promotion to an additive bulk Wilson
packet requires an additional matching theorem.

This version makes that correction append-only, separates the universal
Gaussian baseline from the fourth-cumulant excess, and replaces the ambiguous
ratio $r_{\rm VV}$ by the unique common-determinant tangent which is actually
load-bearing for phenomenology.

\subsection{Two expansion axes and three distinct objects}
\label{subsec:v35-three-objects}

Let
\begin{equation}
 q:=\frac{Q}{D^2}
\label{eq:v35-q}
\end{equation}
be the dimensionless isotropic fourth-cumulant coordinate, while all P1 data
(the quadratic transport $D$, reset clock, calibrated metric, screen
normalization and lower-weight physical inputs) are held fixed.  We distinguish:
\begin{enumerate}[leftmargin=2.2em]
\item the \emph{diamond response} $\mathcal D_\rho(q;\tau)$ for a selected
calibrated tomography profile $\rho$;
\item the \emph{Gaussian baseline} $c_G$, obtained from the complete common
regulated determinant at $q=0$;
\item the \emph{non-Gaussian excess tangent}
\begin{equation}
 \boxed{
 v_Q:=\left.\frac{\partial c_{\rm eff}(q)}{\partial q}\right|_{q=0}}
 \in \mathbb R^3,
\label{eq:v35-vQ-def}
\end{equation}
where $c_{\rm eff}=(c_1,c_2,c_3)$ denotes the curvature-squared coefficient
vector in one fixed off-shell matching scheme.
\end{enumerate}
The total local coefficient near the matched Gaussian reference is therefore
\begin{equation}
 c_{\rm eff}(q)=c_G+qv_Q+O(q^2).
\label{eq:v35-baseline-excess}
\end{equation}
The phenomenological selector supplied by the fourth renewal cumulant is
$v_Q$, not the profile-dependent Gaussian baseline $c_G$.

\begin{theorem}[Baseline/excess separation]
\label{thm:v35-baseline-excess}
Consider two microscopic channels with identical P1 data and identical
quadratic/Gaussian transport, differing only in the fourth-cumulant coordinate
$q$.  Match both channels in the same finite scheme and subtraction
prescription.  Then every $q$-independent Gaussian curvature-squared term
cancels in the centered difference
\begin{equation}
 \Delta_Q c(q):=c_{\rm eff}(q)-c_{\rm eff}(0)
 =qv_Q+O(q^2).
\label{eq:v35-centered-c}
\end{equation}
Under a finite differentiable scheme map $c'=F(c)$ the tangent transforms as
$v_Q'=D F(c_G)v_Q$.  Hence the image line and its transported conormal are
scheme-covariant objects.
\end{theorem}

\begin{proof}
The Taylor expansion gives \eqref{eq:v35-centered-c}.  A common finite scheme
map gives
$F(c_G+qv_Q+O(q^2))-F(c_G)=qD F(c_G)v_Q+O(q^2)$.
No $q$-independent baseline term survives the subtraction.
\end{proof}

\subsection{Profile dependence obstructs a direct Wilson interpretation of the V34 amplitude}
\label{subsec:v35-profile-obstruction}

The calibrated diamond profiles are query/tomography choices.  Their fourth
moments need not agree.  For the explicit $d=3$ polynomial family inherited
from the Ricci-tomography construction, with $\tau/2=1$, the transverse fourth
moment is
\begin{equation}
 \left\langle |z|^4\right\rangle_{m,n}
 =\frac{15(m+4)}{(m+8)(2n+5)(2n+7)}.
\label{eq:v35-profile-fourth}
\end{equation}
In particular
\begin{equation}
 \left\langle |z|^4\right\rangle_{0,0}=\frac3{14},
 \qquad
 \left\langle |z|^4\right\rangle_{0,2}=\frac5{66}.
\label{eq:v35-profile-example}
\end{equation}
Thus fourth profile moments can change while the microscopic renewal channel
and its bulk effective action are held fixed.

\begin{proposition}[Response amplitude is not automatically a bulk Wilson coefficient]
\label{prop:v35-profile-not-wilson}
The V34 coefficient
$\omega_{\rm VV}(\rho)=m_{4,\rho}/720$ is an exact normalization of the
selected quartic van Vleck \emph{response functional}.  If $\rho$ is a freely
chosen tomography profile, $\omega_{\rm VV}(\rho)$ cannot by itself be
identified with a profile-independent bulk Wilson coefficient.  Such an
identification requires either:
\begin{enumerate}[label=\textup{(\roman*)},leftmargin=2.2em]
\item a physical rule that fixes $\rho$ as part of the microscopic matching
scheme, followed by the corresponding scheme transport; or
\item a common induced-action calculation in which the profile dependence
cancels after the complete determinant is assembled.
\end{enumerate}
\end{proposition}

\begin{proof}
Equation \eqref{eq:v35-profile-example} gives two admissible profile choices
with different fourth moments.  The same microscopic bulk action cannot acquire
two different intrinsic coefficients merely because a tomography weight was
changed.  Therefore the response-to-action map is additional data unless the
profile is itself declared physical matching structure.
\end{proof}

\subsection{The van Vleck factor is not the complete Gaussian heat-kernel coefficient}
\label{subsec:v35-heat-kernel-audit}

There is an independent structural check.  The upstream renewal-gravity
analysis records the standard minimal scalar Laplacian benchmark, modulo a
total derivative,
\begin{equation}
 a_2^{\rm min}
 =\frac{13}{1080}R^2
  +\frac1{180}R_{\mu\nu}R^{\mu\nu}
  +\frac1{180}C_{\mu\nu\rho\sigma}C^{\mu\nu\rho\sigma}.
\label{eq:v35-minimal-a2}
\end{equation}
Its curvature-square direction is therefore
\begin{equation}
 b_G=(13,6,6).
\label{eq:v35-bG}
\end{equation}
By contrast, the quartic van Vleck logarithm isolated in V34 has direction
\begin{equation}
 b_{\rm VV}=(-1,8,3).
\label{eq:v35-bVV}
\end{equation}
They are not collinear:
\begin{equation}
 \boxed{
 b_G\times b_{\rm VV}
 =5(-6,-9,22)\neq0.}
\label{eq:v35-noncollinear}
\end{equation}

\begin{theorem}[Van Vleck incompleteness diagnostic]
\label{thm:v35-vv-incomplete}
A quartic coincidence coefficient extracted from
$\log\Delta^{1/2}$ cannot, in general, be identified with the complete
curvature-squared coefficient of a Gaussian induced action.  The minimal
Laplacian benchmark \eqref{eq:v35-noncollinear} is an explicit diagnostic.
The full coefficient also contains the remaining heat-kernel/transport
recurrence and measure contributions of the common operator.
\end{theorem}

\begin{proof}
For a Laplace-type operator the off-diagonal heat kernel factors into the
van Vleck half-density times the world-function exponential and a nontrivial
coefficient series.  If the quartic van Vleck term alone determined the local
curvature-square coefficient, its direction would have to be proportional to
the complete $a_2$ direction in the minimal benchmark.  Equation
\eqref{eq:v35-noncollinear} excludes this.  The benchmark is not asserted to be
the renewal operator; it is used only to falsify the stronger identification
``quartic van Vleck coefficient = complete Gaussian Wilson coefficient''.
\end{proof}

\begin{firewall}[Append-only corrigendum to the V34 interpretation]
\label{fire:v35-v34-corrigendum}
The following V34 statements remain proved exactly as written at the
\emph{response} level:
\begin{enumerate}[leftmargin=2.2em]
\item the quartic covariant van Vleck contraction has direction $(-1,8,3)$;
\item its V33 conormal projection is nonzero for a nondegenerate selected
profile.
\end{enumerate}
What is \emph{not} yet proved is the additional identification of that selected
response amplitude with an additive bulk Wilson packet in the same
normalization as the V28 geodesic-volume packet.  Accordingly the V34 two-ray
Wilson family and the scalar $r_{\rm VV}$ are retained only as a conditional
response-model parametrization.  They are removed from the critical
phenomenology path until a common induced-action map is populated.
\end{firewall}

\subsection{Exact common-determinant tangent}
\label{subsec:v35-common-determinant}

The correct P4r object can be written without any profile ambiguity.  Let
$L=-\Box$ denote the positive spatial Laplacian in the chosen Euclidean
matching chart and write the isotropic Kramers--Moyal operator as
\begin{equation}
 \mathcal R(q)
 =\frac D2L+\frac{D^2q}{8}L^2+O(q^2,\nabla D).
\label{eq:v35-Rq}
\end{equation}
Let
\begin{equation}
 \mathcal A(q)=\partial_\tau+\mathcal R(q)
\label{eq:v35-Aq}
\end{equation}
with the same boundary conditions, proper-time regulator and subtraction
prescription for all $q$, and define
\begin{equation}
 W(q)=\frac12\Tr_{\rm reg}\log\mathcal A(q).
\label{eq:v35-Wq}
\end{equation}

\begin{theorem}[Common-determinant fourth-cumulant insertion]
\label{thm:v35-common-det-derivative}
At the matched Gaussian point $q=0$,
\begin{equation}
 \boxed{
 \dot W_Q
 :=\left.\partial_qW(q)\right|_{0}
 =\frac{D^2}{16}\Tr_{\rm reg}
 \left[\mathcal A_0^{-1}L^2\right].}
\label{eq:v35-dotW}
\end{equation}
If the temporal/spatial split used by the inherited induced-action theorem is
valid and the regulator admits a proper-time representation, then, before the
standard regulator-boundary terms are evaluated,
\begin{equation}
 \boxed{
 \dot W_Q
 =\frac14\int_0^\infty ds\,\chi(s)\,
 \Tr_\tau(e^{-s\partial_\tau})
 \partial_s^2\Tr_{\rm sp}
 \left(e^{-s(D/2)L}\right)
 +\mathcal B_\chi.}
\label{eq:v35-propertime-insertion}
\end{equation}
Here $\chi$ is the declared common regulator and $\mathcal B_\chi$ contains
only the regulator/end-point terms generated when derivatives are moved through
the proper-time integral.  The curvature-square coefficient of
\eqref{eq:v35-dotW} is precisely the non-Gaussian excess tangent $v_Q$ of
\eqref{eq:v35-vQ-def}.
\end{theorem}

\begin{proof}
Differentiate $\frac12\Tr\log\mathcal A(q)$ using
$\partial_q\log\det\mathcal A=\Tr(\mathcal A^{-1}\partial_q\mathcal A)$.
Equation \eqref{eq:v35-Rq} gives
$\partial_q\mathcal A|_0=(D^2/8)L^2$, proving
\eqref{eq:v35-dotW}.  With
$\mathcal A_0^{-1}=\int_0^\infty ds\,\chi(s)e^{-s\mathcal A_0}$ and the split
heat trace, use
\[
 L^2e^{-s(D/2)L}=\frac4{D^2}\partial_s^2e^{-s(D/2)L}.
\]
The prefactor becomes $1/4$.  Any integrations by parts against the regulator
produce the explicitly retained term $\mathcal B_\chi$.  Extracting the local
curvature-square coefficient gives $v_Q$.
\end{proof}

\begin{remark}[Relation to the inherited induced-action factorization]
\label{rem:v35-Aren}
The upstream induced-action theorem factorizes the fourth-cumulant contribution
as a common temporal moment times a spatial Seeley coefficient.  Equation
\eqref{eq:v35-propertime-insertion} is the infinitesimal common-window version of
that statement.  It also shows why a numerical coefficient cannot be obtained
from the quartic van Vleck coincidence term alone: the result depends on the
complete heat trace and the common regulator/subtraction prescription.
\end{remark}

\subsection{The corrected one-scalar phenomenology gate}
\label{subsec:v35-corrected-gate}

The V28 volume contribution per unit $q$ is
\begin{equation}
 v_{\rm vol}
 =G_{\rm ren}\lambda^2
 \left(\frac54,\frac53,-\frac14\right).
\label{eq:v35-vvol}
\end{equation}
Introduce the Gauss--Bonnet-invariant projection to the V28 physical
curvature coordinates,
\begin{equation}
 P_{\rm phys}
 \begin{pmatrix}c_1\\c_2\\c_3\end{pmatrix}
 =
 \begin{pmatrix}
 c_C\\ c_R
 \end{pmatrix}
 =
 \begin{pmatrix}
 0&\frac12&1\\
 1&\frac13&0
 \end{pmatrix}
 \begin{pmatrix}c_1\\c_2\\c_3\end{pmatrix}.
\label{eq:v35-Pphys}
\end{equation}
The volume ray is annihilated by
\begin{equation}
 \ell_{\rm vol}^{\rm phys}=(65,-21).
\label{eq:v35-lvol-phys}
\end{equation}

\begin{definition}[Common-determinant transverse excess]
\label{def:v35-XiQ}
Define
\begin{equation}
 \boxed{
 \Xi_Q
 :=\ell_{\rm vol}^{\rm phys}P_{\rm phys}v_Q
 =-21v_Q^{(1)}+\frac{51}{2}v_Q^{(2)}+65v_Q^{(3)}.}
\label{eq:v35-XiQ}
\end{equation}
Because $\ell_{\rm vol}^{\rm phys}P_{\rm phys}v_{\rm vol}=0$, the same number is
obtained by replacing $v_Q$ with $v_Q-v_{\rm vol}$.
\end{definition}

\begin{theorem}[Correct survival gate for the Burnett conormal]
\label{thm:v35-survival-gate}
In the centered non-Gaussian excess, the V28/V29 physical conormal survives to
first order in $q$ if and only if
\begin{equation}
 \boxed{\Xi_Q=0.}
\label{eq:v35-XiQ-zero}
\end{equation}
If $\Xi_Q\neq0$, the full common-determinant fourth-cumulant tangent leaves the
volume conormal and the contact/amplitude ray must be updated before the V32
multi-species test is interpreted as a quantitative Burnett-ratio prediction.
This statement is independent of any $q$-independent Gaussian baseline.
\end{theorem}

\begin{proof}
By Theorem \ref{thm:v35-baseline-excess}, the centered physical curvature
response is
$qP_{\rm phys}v_Q+O(q^2)$.  Applying the nonzero covector
$\ell_{\rm vol}^{\rm phys}$ gives $q\Xi_Q+O(q^2)$.  Its first-order vanishing is
therefore equivalent to \eqref{eq:v35-XiQ-zero}.  The Gaussian baseline cancels
before this test is applied.
\end{proof}

\subsection{Observable consequence: subtract the matched Gaussian channel, not a guessed van Vleck Wilson packet}
\label{subsec:v35-observable-excess}

Let $\mathcal A(q)$ be any common on-shell amplitude panel and let
$\mathcal A_G:=\mathcal A(0)$ be the amplitude of the Gaussian-matched reference
channel with identical P1 data.  Then
\begin{equation}
 \Delta_Q\mathcal A(q)
 :=\mathcal A(q)-\mathcal A_G
 =q\,J_{\mathcal A}v_Q+O(q^2).
\label{eq:v35-amplitude-excess}
\end{equation}
Every $q$-independent curvature-square baseline---including the complete
Gaussian determinant, ordinary calibrated SM/gravity loops, and finite scheme
terms common to the pair---is absent from \eqref{eq:v35-amplitude-excess}.
Thus the V30--V32 multi-species programme should be interpreted as a test of
the \emph{centered fourth-cumulant excess} after the common Gaussian reference
has been subtracted.  This is stronger operationally than attempting to assign
a scheme-independent meaning to the total off-shell curvature coefficients.

\begin{firewall}[No numerical $r_{\rm VV}$ is extracted in V35]
\label{fire:v35-no-r}
The current corpus does not contain a common regulated evaluation of
\eqref{eq:v35-dotW} sufficient to identify $v_Q$.  Historical numerical triples
proportional to $(-1,8,3)$ occur in branches where the normalized quartic
projection was assumed as input; they are not an independent derivation of the
common-window normalization.  V35 therefore does not recycle them to assign
$r_{\rm VV}$.
\end{firewall}

\subsection{Status after V35}
\label{subsec:v35-status}

\subsection*{Proved}
\begin{enumerate}[leftmargin=2.2em]
\item The V34 quartic van Vleck calculation is a valid profile-averaged diamond
response, but its amplitude is profile dependent and therefore is not by itself
a bulk Wilson coefficient.
\item The complete Gaussian heat-kernel coefficient need not be collinear with
the van Vleck-only quartic response; the inherited minimal-Laplacian benchmark
gives the explicit noncollinearity \eqref{eq:v35-noncollinear}.
\item The exact common-determinant fourth-cumulant derivative is
\eqref{eq:v35-dotW}, with the proper-time reduction
\eqref{eq:v35-propertime-insertion} under the stated split assumptions.
\item The correct first-order physical survival gate for the V28 volume
conormal is the single common-determinant scalar $\Xi_Q$ of
\eqref{eq:v35-XiQ}, not the response-level ratio $r_{\rm VV}$.
\item Every $q$-independent Gaussian baseline cancels in the matched-channel
excess \eqref{eq:v35-amplitude-excess}.
\end{enumerate}

\subsection*{Inherited}
\begin{enumerate}[leftmargin=2.2em]
\item V28: the populated geodesic-volume fourth-cumulant contribution
$v_{\rm vol}$.
\item V29--V32: the field-redefinition quotient and the scalar/multi-species
on-shell visibility maps.
\item V34: the exact quartic van Vleck response direction $(-1,8,3)$.
\item The upstream induced-action factorization: one common temporal response
moment times the spatial Seeley coefficient of the renewal operator.
\end{enumerate}

\subsection*{Conditional / corrected}
\begin{enumerate}[leftmargin=2.2em]
\item V34's two-ray \emph{Wilson} family and the numerical parameter
$r_{\rm VV}$ are conditional on an unproved response-to-induced-action
identification; they should not be used as completed phenomenological
coefficients.
\item The common determinant formula assumes the same operator domain,
regulator, subtraction prescription and P1 calibration across the $q$ family.
\item A chosen tomography profile may be promoted to part of a matching scheme,
but its profile dependence must then be transported as ordinary scheme data and
cannot be called intrinsic.
\end{enumerate}

\subsection*{Open beyond V35}
\begin{enumerate}[leftmargin=2.2em]
\item Evaluate the curvature-square coefficient of the single trace
$\Tr_{\rm reg}(\mathcal A_0^{-1}L^2)$ in the actual renewal operator and
subtraction convention.
\item Compute $\Xi_Q$.  If it vanishes, restore the V28/V29 conormal for the
\emph{centered fourth-cumulant excess}; if not, update the excess ray from the
full $v_Q$.
\item Only after that step return to the V32 twelve-row multi-species panel and
test the selector-dependent excess amplitudes.
\item Keep the total Gaussian curvature-squared baseline as a separately
calibrated EFT contribution unless the common determinant is evaluated
absolutely in the chosen scheme.
\end{enumerate}

\subsection*{Exact next load-bearing theorem}
\begin{center}
\fbox{\parbox{0.92\textwidth}{
\textbf{P4s: common induced-action fourth-cumulant insertion.}
For the same band-reduced renewal operator, P1 calibration, proper-time
regulator and subtraction convention used to define the physical continuum,
evaluate the local curvature-square coefficient of
\[
 \dot W_Q=\frac{D^2}{16}\Tr_{\rm reg}
 [\mathcal A_0^{-1}L^2].
\]
Return the three-component tangent $v_Q$ and, first of all, its single
Gauss--Bonnet-invariant transverse projection
\[
 \Xi_Q=(65,-21)P_{\rm phys}v_Q.
\]
Do not import the V34 profile amplitude or a historical assumed quartic triple
as this trace.  If $\Xi_Q=0$, the V28 contact/amplitude conormal survives for
the centered non-Gaussian excess and the programme returns immediately to the
V32 mixed-species validation panel.  If $\Xi_Q\neq0$, transport the newly
computed excess ray through V29--V32 before any RG calculation.}}
\end{center}

\section*{Append-only continuation marker: V35}
\addcontentsline{toc}{section}{Append-only continuation marker: V35}
\textbf{End of V35.}  The complete V1--V34 body above is retained verbatim.
V35 goes upstream from the attempted common scalar $r_{\rm VV}$ and separates a
profile-averaged quartic van Vleck response from a complete induced-action
Wilson coefficient.  The V34 tensor contraction remains valid; its Wilson-level
promotion is downgraded pending a common determinant calculation.  The
phenomenological selector is reformulated as the centered fourth-cumulant
tangent $v_Q=\partial_q c_{\rm eff}|_{q=0}$, whose exact common-operator
representation is $\frac{D^2}{16}\Tr_{\rm reg}(\mathcal A_0^{-1}L^2)$.  The
single load-bearing survival test is the physical conormal projection $\Xi_Q$.


% =============================================================================
% APPEND-ONLY CONTINUATION: VERSION V36
% The complete V1--V35 body above is retained verbatim.  V36 begins here.
% =============================================================================

\clearpage
\section{Version V36: Covariant Fourth-Moment Completion, the Four-Dimensional Heat-Trace Reduction, and the Local-Subtraction Obstruction}
\label{sec:v36-covariant-fourth-moment}

\subsection{Purpose: P4s must use the complete curved fourth-moment operator}
\label{subsec:v36-purpose}

V35 isolated the correct centered microscopic coordinate
$q=Q/D^2$ and replaced the profile-level van Vleck normalization by a common
regulated determinant.  Its explicit insertion,
\begin{equation}
 \frac{D^2}{16}\Tr_{\rm reg}(\mathcal A_0^{-1}L^2),
 \label{eq:v36-v35-principal-insertion}
\end{equation}
was deliberately written from the principal scalar shorthand
$\mathcal R=D L/2+D^2qL^2/8+\cdots$.  The upstream covariant
Kramers--Moyal construction, however, defines the fourth moment before the
principal-symbol contraction.  On a curved manifold, different orderings of
four covariant derivatives acting on the Hessian of a scalar differ by Ricci
terms.  Consequently \eqref{eq:v36-v35-principal-insertion} is the principal
part of P4s, not the complete canonical scalar insertion.

This version computes the missing covariant completion exactly, evaluates its
local curvature-square heat coefficient on the four-dimensional matching
branch, and then goes one step further upstream: the surviving coefficient is a
local power-sensitive counterterm coordinate.  It is therefore not yet a
scheme-independent phenomenological prediction.  The first invariant target is
instead the nonlocal/on-shell fourth-cumulant form factor after all local
counterterms have been subtracted.

\subsection{Exact isotropic fourth-moment contraction on scalars}
\label{subsec:v36-fourth-contraction}

Let $\Delta=\nabla^\mu\nabla_\mu$ and retain the V35 convention
$L=-\Delta$.  We use the curvature convention for which the scalar Bochner
identity reads
\begin{equation}
 \nabla^\nu\nabla_\mu\nabla_\nu f
 =\nabla_\mu\Delta f+R_\mu{}^\nu\nabla_\nu f.
 \label{eq:v36-bochner}
\end{equation}
Let
\begin{equation}
 S^{\mu\nu\rho\sigma}
 :=g^{\mu\nu}g^{\rho\sigma}
  +g^{\mu\rho}g^{\nu\sigma}
  +g^{\mu\sigma}g^{\nu\rho}
 \label{eq:v36-isotropic-four-tensor}
\end{equation}
be the isotropic fourth-moment tensor.

\begin{theorem}[Exact covariant isotropic fourth-moment operator]
\label{thm:v36-exact-fourth-operator}
For every smooth scalar $f$,
\begin{equation}
 \boxed{
 S^{\mu\nu\rho\sigma}
 \nabla_\mu\nabla_\nu\nabla_\rho\nabla_\sigma f
 =3\Delta^2 f
 +2R^{\mu\nu}\nabla_\mu\nabla_\nu f
 +(\nabla^\mu R)\nabla_\mu f.}
 \label{eq:v36-exact-fourth-contraction}
\end{equation}
Equivalently,
\begin{equation}
 \boxed{
 S^{\mu\nu\rho\sigma}
 \nabla_\mu\nabla_\nu\nabla_\rho\nabla_\sigma
 =3\Delta^2
 +2\nabla_\mu(R^{\mu\nu}\nabla_\nu).}
 \label{eq:v36-divergence-form}
\end{equation}
Hence, for $q=Q/D^2$, the canonical geodesic scalar Kramers--Moyal lift has
\begin{equation}
 \boxed{
 \left.\partial_q\mathcal R(q)\right|_{0}
 =\frac{D^2}{8}L^2
 +\frac{D^2}{12}\mathcal B_R,\qquad
 \mathcal B_R:=\nabla_\mu(R^{\mu\nu}\nabla_\nu).}
 \label{eq:v36-canonical-q-operator}
\end{equation}
\end{theorem}

\begin{proof}
The $g^{\mu\nu}g^{\rho\sigma}$ contraction is $\Delta^2 f$.
For either crossed contraction, apply \eqref{eq:v36-bochner} and then one
additional divergence:
\[
 \nabla^\mu\!\left(
  \nabla_\mu\Delta f+R_\mu{}^\nu\nabla_\nu f
 \right)
 =\Delta^2f+R^{\mu\nu}\nabla_\mu\nabla_\nu f
  +(\nabla^\mu R_\mu{}^\nu)\nabla_\nu f.
\]
The contracted Bianchi identity
$\nabla^\mu R_{\mu\nu}=\tfrac12\nabla_\nu R$ shows that each crossed term
is
$\Delta^2f+R^{\mu\nu}\nabla_\mu\nabla_\nu f
+\tfrac12(\nabla R)\cdot\nabla f$.
Adding the three pairings proves \eqref{eq:v36-exact-fourth-contraction}; the
Bianchi identity also gives \eqref{eq:v36-divergence-form}.  Finally the
Kramers--Moyal coefficient is $Q/4!=D^2q/24$, and $\Delta^2=L^2$.
\end{proof}

\begin{firewall}[Principal scalar shorthand versus complete curved lift]
\label{fire:v36-principal-shorthand}
The familiar expression
$\mathcal R=D L/2+D^2qL^2/8+\cdots$ is the principal scalar symbol of the
isotropic fourth-moment operator.  Equation \eqref{eq:v36-canonical-q-operator}
shows that it omits the lower-order curvature completion
$D^2\mathcal B_R/12$.  Any additional internal-band, measure, source or
subprincipal $q$-linear curved terms of the full Renewal ESM operator must be
added separately; V36 does not set them to zero.
\end{firewall}

\subsection{Corrected finite-regulator determinant tangent}
\label{subsec:v36-corrected-determinant}

Let $\mathcal A_0=\partial_\tau+(D/2)L$ denote the matched Gaussian operator
on the canonical scalar branch.  At a fixed common regulator the determinant
variation is now
\begin{equation}
 \boxed{
 \dot W_Q^{\rm can}
 =\frac{D^2}{16}\Tr_{\rm reg}(\mathcal A_0^{-1}L^2)
 +\frac{D^2}{24}\Tr_{\rm reg}(\mathcal A_0^{-1}\mathcal B_R).}
 \label{eq:v36-corrected-dotW}
\end{equation}
For the full band operator we therefore write
\begin{equation}
 \dot W_Q^{\rm full}
 =\dot W_Q^{\rm can}+\dot W_Q^{\rm extra},
 \label{eq:v36-full-dotW}
\end{equation}
where $\dot W_Q^{\rm extra}$ is a typed placeholder for every $q$-linear
curved-lift term not contained in the canonical scalar geodesic
Kramers--Moyal operator.

\begin{proposition}[Order-changing perturbation and nonuniformity of the $q$ expansion]
\label{prop:v36-order-change}
The perturbation $qL^2$ is not relatively bounded with respect to $L$ uniformly
in the ultraviolet.  With spectral cutoff $L\le\Lambda^2$, the dimensionless
relative size of the fourth-order principal term is
\begin{equation}
 \frac{(D^2|q|/8)\Lambda^4}{(D/2)\Lambda^2}
 =\frac{D|q|}{4}\Lambda^2.
 \label{eq:v36-relative-size}
\end{equation}
Thus the limits $q\to0$ and $\Lambda\to\infty$ need not commute.  The
linearized determinant \eqref{eq:v36-corrected-dotW} is therefore defined
first at a common finite regulator and only subsequently renormalized.
In particular, the heat-kernel coefficients of a pure fourth-order operator
$L^2$ characterize a different ultraviolet problem from the renormalized
$q$-tangent about the Gaussian $L$ branch.
\end{proposition}

\begin{proof}
On an $L$-eigenmode of eigenvalue $\lambda_L$, the ratio of the fourth-order
term to the Gaussian term is $(D|q|/4)\lambda_L$.  It is unbounded as
$\lambda_L\to\infty$.  The displayed cutoff bound is obtained by setting
$\lambda_L=\Lambda^2$.  Hence the usual norm-analytic perturbation argument
cannot be uniform under cutoff removal.
\end{proof}

\subsection{Four-dimensional routing: the direct $L^2$ insertion has no interior curvature-square coefficient}
\label{subsec:v36-l2-vanishing}

The V28--V35 phenomenological quotient is four-dimensional.  We therefore
first state the heat-trace result on the four-dimensional covariant Euclidean
factor used to define $R^2$, $R_{\mu\nu}^2$ and $C^2$.  If an implementation
instead uses a literal three-dimensional spatial heat kernel, this step must be
replaced by an explicit lift from that three-dimensional coefficient to the
four-dimensional covariant action; no such lift is silently assumed here.

For the minimal scalar $L=-\Delta$ in four dimensions,
\begin{equation}
 \Tr_{\rm sp}(e^{-tL})
 \sim(4\pi t)^{-2}
 \left[A_0+tA_1+t^2A_2+t^3A_3+\cdots\right],
 \label{eq:v36-heat-trace}
\end{equation}
where $A_2$ is the integrated local curvature-square coefficient modulo a total
derivative.

\begin{theorem}[Four-dimensional principal-insertion null at curvature-square order]
\label{thm:v36-l2-a2-zero}
On the four-dimensional covariant matching factor,
\begin{equation}
 \boxed{
 \left.\Tr(L^2e^{-tL})\right|_{R^2}
 =\left.\partial_t^2\Tr(e^{-tL})\right|_{R^2}=0.}
 \label{eq:v36-l2-a2-zero}
\end{equation}
Consequently the first term of \eqref{eq:v36-corrected-dotW} carries no
\emph{interior} local curvature-square coefficient.  A curvature-square number
assigned to it solely through proper-time end points or local subtraction data
is regulator/scheme data, not an interior heat coefficient.
\end{theorem}

\begin{proof}
The curvature-square part of \eqref{eq:v36-heat-trace} is
$(4\pi)^{-2}A_2$, independent of $t$.  Its second derivative is therefore zero.
The identity $L^2e^{-tL}=\partial_t^2e^{-tL}$ proves the result.
\end{proof}

\begin{remark}[Dimension firewall]
\label{rem:v36-dimension-firewall}
In general dimension $n$, the $A_2$ term scales as
$t^{2-n/2}A_2$ and two $t$ derivatives need not vanish.  The cancellation in
Theorem \ref{thm:v36-l2-a2-zero} is specifically the four-dimensional physical
matching statement.  This makes the dimension of the heat-kernel factor a
load-bearing datum rather than harmless notation.
\end{remark}

\subsection{The Ricci-divergence insertion and its local heat coefficient}
\label{subsec:v36-BR-heat}

Let $K(t;x,x')$ be the minimal scalar heat kernel.  The standard off-diagonal
parametrix gives at coincidence, through the order required here,
\begin{equation}
 \left[\nabla_\mu\nabla_{\nu'}K(t;x,x')\right]
 =(4\pi t)^{-2}
 \left[
 \frac{g_{\mu\nu}}{2t}
 -\frac16R_{\mu\nu}
 +\frac1{12}R g_{\mu\nu}
 +O(t\mathcal R^2,t\nabla^2\mathcal R)
 \right].
 \label{eq:v36-mixed-heat-derivative}
\end{equation}
Here brackets denote coincidence and $\mathcal R$ denotes curvature generically.

\begin{theorem}[Local heat trace of the canonical curvature completion]
\label{thm:v36-BR-heat}
On a closed four-dimensional matching manifold, or for compactly supported
localization away from the boundary,
\begin{equation}
 \boxed{
 \Tr\!\left(\mathcal B_Re^{-tL}\right)
 =(4\pi t)^{-2}\!\int\!\sqrt g\,
 \left[
 -\frac{R}{2t}
 +\frac16R_{\mu\nu}R^{\mu\nu}
 -\frac1{12}R^2
 +O(t\mathcal R^3,t\mathcal R\nabla^2\mathcal R)
 \right].}
 \label{eq:v36-BR-trace}
\end{equation}
The local curvature-square direction of the canonical $q$-insertion is therefore
\begin{equation}
 \boxed{(-1,2,0)}
 \label{eq:v36-canonical-local-direction}
\end{equation}
in the ordered basis $(R^2,R_{\mu\nu}^2,C^2)$.
\end{theorem}

\begin{proof}
Using an orthonormal eigenbasis of $L$ and integration by parts,
\[
 \Tr(\mathcal B_Re^{-tL})
 =-\int\sqrt g\,R^{\mu\nu}
 \left[\nabla_\mu\nabla_{\nu'}K(t;x,x')\right].
\]
Insert \eqref{eq:v36-mixed-heat-derivative}.  Contracting the first term gives
$-R/(2t)$; the remaining two displayed terms give
$+R_{\mu\nu}R^{\mu\nu}/6-R^2/12$.
As a sign check, on constant curvature
$R_{\mu\nu}=Rg_{\mu\nu}/4$ one has
$\mathcal B_R=(R/4)\Delta=-(R/4)L$; differentiating the ordinary scalar heat
trace gives exactly the same $-R^2/24$ curvature-square coefficient.
\end{proof}

\subsection{Bare local tangent and the V35 conormal}
\label{subsec:v36-local-tangent}

Put $a=D/2$ and define the common proper-time moment
\begin{equation}
 \mathfrak M_{-2}[\chi]
 :=\int_0^\infty ds\,\chi(s)
 \Tr_\tau(e^{-s\partial_\tau})\,s^{-2},
 \label{eq:v36-Mminus2}
\end{equation}
under the same regularization used for the determinant.  The word ``moment''
is formal until the ultraviolet subtraction has been specified.

\begin{corollary}[Canonical bare local fourth-cumulant tangent]
\label{cor:v36-bare-local-vQ}
On the canonical scalar geodesic branch, ignoring
$\dot W_Q^{\rm extra}$ but retaining the complete covariant fourth-moment
operator, the local curvature-square part of the fixed-regulator tangent is
\begin{equation}
 \boxed{
 v_{Q,\rm loc}^{\rm can}
 =A_Q(-1,2,0),\qquad
 A_Q:=\frac{\mathfrak M_{-2}[\chi]}{1152\pi^2}.}
 \label{eq:v36-vQ-local}
\end{equation}
Its V35 transverse projection is
\begin{equation}
 \boxed{
 \Xi_{Q,\rm loc}^{\rm can}
 =72A_Q
 =\frac{\mathfrak M_{-2}[\chi]}{16\pi^2}.}
 \label{eq:v36-Xi-local}
\end{equation}
\end{corollary}

\begin{proof}
The determinant prefactor of the $\mathcal B_R$ term is $D^2/24$.
Equation \eqref{eq:v36-BR-trace} is evaluated at $t=as$; hence
\[
 \frac{D^2}{24}(4\pi a)^{-2}
 =\frac1{96\pi^2}.
\]
Multiplying by $(-1/12,1/6,0)$ gives
$(-1,2,0)/(1152\pi^2)$ times \eqref{eq:v36-Mminus2}.  Finally
\[
 -21(-1)+\frac{51}{2}(2)+65(0)=72.
\]
\end{proof}

The Gauss--Bonnet-invariant physical projection of this direction is
\begin{equation}
 P_{\rm phys}(-1,2,0)^T
 =\begin{pmatrix}1\\-1/3\end{pmatrix},
 \label{eq:v36-physical-projection}
\end{equation}
so the V29 stress-contact descendant is proportional to
\begin{equation}
 \boxed{(\delta\kappa_1,\delta\kappa_2)\propto(2,-1).}
 \label{eq:v36-contact-direction}
\end{equation}
This direction is linearly independent of the V29 volume-sector vector
$(14,17)$.  Equivalently its intrinsic contact conormal is
$\delta\kappa_1+2\delta\kappa_2=0$, whereas
$17\delta\kappa_1-14\delta\kappa_2\neq0$.

\subsection{Why this does not yet produce a phenomenological Wilson ratio}
\label{subsec:v36-subtraction-obstruction}

The moment \eqref{eq:v36-Mminus2} is ultraviolet sensitive.  More importantly,
the operators $R^2$, $R_{\mu\nu}^2$ and $C^2$ are allowed local counterterms at
the same derivative order.  A finite common-scheme change may therefore add
$q$-linear local terms to the matching coefficients while leaving all properly
transported observables unchanged.

\begin{theorem}[Local-subtraction obstruction]
\label{thm:v36-local-subtraction-obstruction}
The numerical vector \eqref{eq:v36-vQ-local} and the bare scalar
\eqref{eq:v36-Xi-local} are not, by themselves, scheme-independent
phenomenological predictions.  They are the local fixed-regulator contribution
of the canonical scalar geodesic Kramers--Moyal completion.  A physical claim
requires one of the following additional inputs:
\begin{enumerate}[label=\textup{(\roman*)},leftmargin=2.2em]
\item an independently fixed microscopic subtraction/matching condition which
selects the finite local $q$-linear counterterms; or
\item a projection onto a nonlocal/on-shell observable component---for example a
nonanalytic momentum form factor---which cannot be changed by local finite
counterterms.
\end{enumerate}
In particular, the sign or nonvanishing of the bare
$\Xi_{Q,\rm loc}^{\rm can}$ is not sufficient to invalidate or establish a
scheme-transported amplitude conormal.
\end{theorem}

\begin{proof}
At the retained EFT order one may change finite matching convention by adding a
local functional
\[
 q\int\sqrt g\,
 (\alpha R^2+\beta R_{\mu\nu}R^{\mu\nu}+\gamma C^2)
\]
with the corresponding inverse transformation on the coefficient coordinates
and observables.  This shifts the local $q$ tangent while leaving the physical
amplitude after matching unchanged.  By contrast a local polynomial
counterterm cannot alter a genuinely nonanalytic momentum dependence or its
spectral discontinuity.  Hence only a fixed microscopic normalization or a
nonlocal/on-shell projection removes the ambiguity.
\end{proof}

\begin{firewall}[Append-only correction to the V35 P4s target]
\label{fire:v36-v35-correction}
Equation \eqref{eq:v35-dotW} remains a correct finite-regulator derivative for
the \emph{principal scalar truncation} used in V35.  It is not the complete
canonical geodesic fourth-moment derivative because
\eqref{eq:v36-canonical-q-operator} supplies the curvature term
$D^2\mathcal B_R/12$.  Moreover, because $qL^2$ changes the differential order,
the derivative at $q=0$ is not uniformly interchangeable with ultraviolet
cutoff removal.  Historical use of the Seeley coefficient of a standalone
fourth-order operator $L^2$ therefore does not by itself compute the
renormalized $q$ tangent about the Gaussian branch.
\end{firewall}

\subsection{What remains invariant from V31--V33}
\label{subsec:v36-universality-survives}

The V36 correction is still a pure-gravity local insertion before the V29
field-redefinition quotient.  Consequently, within the canonical scalar branch,
it preserves the species-universal structure of the induced stress contacts.
The seven V31--V33 species-universality diagnostics therefore remain the correct
first closure tests:
\begin{equation}
 b_{HH}=b_{HF}=b_{HV}=b_{FF}=b_{FV}=b_{VV},
 \qquad
 c_{HH}=c_{HF}=c_{FF},
 \label{eq:v36-species-universality}
\end{equation}
while the tensor/trace ratio is sensitive to the matching prescription and to
any additional $q$-linear band/matter selector contained in
$\dot W_Q^{\rm extra}$.

\subsection{Status after V36}
\label{subsec:v36-status}

\subsection*{Proved}
\begin{enumerate}[leftmargin=2.2em]
\item The canonical geodesic scalar isotropic fourth moment is not merely
$3\Delta^2$: its exact covariant completion is
$3\Delta^2+2\nabla_\mu(R^{\mu\nu}\nabla_\nu)$.
\item The V35 determinant tangent must therefore be supplemented by the second
trace in \eqref{eq:v36-corrected-dotW} on that branch.
\item The $qL^2$ perturbation changes differential order, so the $q\to0$ tangent
must be formed at a common regulator before ultraviolet renormalization.
\item On the four-dimensional covariant matching factor, the direct $L^2$
insertion has zero interior local curvature-square heat coefficient.
\item The canonical Ricci-divergence completion has local direction
$(-1,2,0)$ and fixed-regulator conormal
$\Xi_{Q,\rm loc}^{\rm can}=\mathfrak M_{-2}/(16\pi^2)$.
\item This local number is subtraction dependent; it is not by itself a
scheme-independent phenomenological relation.
\end{enumerate}

\subsection*{Inherited}
\begin{enumerate}[leftmargin=2.2em]
\item V28--V33: the geodesic-volume branch, its field-redefinition quotient,
and the multi-species stress-contact visibility programme.
\item V34: the profile-level quartic van Vleck tensor contraction.
\item V35: baseline/excess separation and the insistence on one common
regulated determinant and matched Gaussian reference.
\end{enumerate}

\subsection*{Conditional / branch-specific}
\begin{enumerate}[leftmargin=2.2em]
\item The explicit vector $(-1,2,0)$ assumes the canonical scalar geodesic
Kramers--Moyal lift and the four-dimensional covariant matching heat kernel.
\item The full Renewal ESM band may add $q$-linear matrix, measure,
source/contact or subprincipal terms collected in
$\dot W_Q^{\rm extra}$; they remain to be populated rather than set to zero.
\item Boundary contributions require the corresponding boundary heat-kernel
packet and are not included in the bulk formulas above.
\end{enumerate}

\subsection*{Open beyond V36}
\begin{enumerate}[leftmargin=2.2em]
\item Construct the full band-covariant operator
$B_Q^{\rm full}=\partial_q\mathcal A_{\rm band}|_0$, including all matrix,
measure and source descendants required by the ESM continuum.
\item Separate its local polynomial curvature/contact terms from the nonlocal
form factors in one fixed physical renormalization scheme.
\item Extract a scheme-invariant $q$-linear on-shell/nonanalytic amplitude
coefficient and test whether its microscopic image has a nontrivial conormal.
\item Only if a microscopic subtraction condition independently fixes the local
finite part should a numerical local fourth-cumulant Wilson ratio be quoted.
\end{enumerate}

\subsection*{Exact next load-bearing theorem}
\begin{center}
\fbox{\parbox{0.92\textwidth}{
\textbf{P4t: full band-covariant fourth-cumulant form factor.}
Construct the complete same-scheme operator
\[
 B_Q^{\rm full}
 :=\left.\partial_q\mathcal A_{\rm band}(q)\right|_{q=0}
\]
through the retained ESM source order, beginning with the canonical scalar term
\[
 \frac{D^2}{8}L^2
 +\frac{D^2}{12}\nabla_\mu(R^{\mu\nu}\nabla_\nu),
\]
and append every matrix/band, measure and source/contact descendant required by
the native continuum construction.  Evaluate
$\frac12\Tr_{\rm reg}(\mathcal A_0^{-1}B_Q^{\rm full})$ to second order in
external curvature/source momenta.  Subtract all local polynomial counterterms
allowed by the frozen physical scheme and return first the nonanalytic/on-shell
$q$-linear form-factor vector (logarithmic coefficient or spectral
discontinuity).  Because local finite counterterms cannot change that vector,
its left nullspace is the next genuinely scheme-independent phenomenology gate.
If the nonanalytic vector vanishes identically, prove that fact and record that
the fourth-cumulant determinant supplies no scheme-independent curvature-square
prediction at first order without an additional microscopic local matching
condition.}}
\end{center}

\section*{Append-only continuation marker: V36}
\addcontentsline{toc}{section}{Append-only continuation marker: V36}
\textbf{End of V36.}  The complete V1--V35 body above is retained verbatim.
V36 goes upstream inside P4s.  It replaces the principal-symbol shorthand
$Q\Box^2/8$ by the exact canonical covariant scalar fourth-moment operator,
whose $q$ derivative contains the additional Ricci-divergence term
$D^2\nabla_\mu(R^{\mu\nu}\nabla_\nu)/12$.  On the four-dimensional physical
matching factor the direct $L^2$ insertion has no interior local
curvature-square coefficient, whereas the curvature completion produces the
fixed-regulator local direction $(-1,2,0)$.  Its coefficient is ultraviolet and
finite-subtraction sensitive, so the raw local conormal is not promoted to a
physical prediction.  The next load-bearing object is the complete band-covariant
$q$ insertion projected onto a nonlocal/on-shell form factor that survives all
local finite counterterms.



% =============================================================================
\section{V37 --- Upstream infrared router: Hodge degree, Burnett nonlocality,
 and the scope of the local Wilson programme}
\label{sec:v37-upstream-router}
% =============================================================================

V36 ended with a deliberately scheme-invariant target: rather than promote the
fixed-regulator local vector to phenomenology, one should isolate the
nonanalytic part of the complete $q$-linear band form factor.  Before carrying
out that expensive band calculation there is a more upstream question.  The
microscopic renewal-current programme already classifies whether the quartic
(Burnett) correction has a cutoff-independent local continuum coefficient at
all.  The classification is controlled by the infrared spectral measure of the
physical reset-displacement current on the critical decorated skeleton.

This section imports that classifier into the perturbative Einstein--Standard
Model continuum line, proves the required Hodge shift self-containedly, and
uses it to place an exact gate in front of the V28--V36 local-Wilson route.  The
main correction is conceptual:
\begin{center}
\fbox{\parbox{0.91\textwidth}{
A finite-window or fixed-regulator local four-derivative coefficient need not be
the leading continuum correction.  If the physical reset-displacement current
lies in the nonlocal Burnett branch, the correct continuum object is a singular
form factor.  The local coefficients of V28--V36 then remain legitimate finite
matching data, but they are not the primary continuum phenomenology.}}
\end{center}

\subsection{Spectral measure of an exact displacement current}
\label{subsec:v37-degree-shift}

Let $G_R$ be a finite wired ball of the decorated skeleton.  Write
$d_0:C^0(G_R)\to C^1(G_R)$ for the graph coboundary and
\begin{equation}
 \Delta_{0,R}=d_0^*d_0,
 \qquad
 \Delta_{1,R}=d_0d_0^*.
 \label{eq:v37-laplacians}
\end{equation}
Let $\phi_R$ be a scalar position/embedding observable and let
\begin{equation}
 J_R=d_0\phi_R
 \label{eq:v37-current-gradient}
\end{equation}
be the corresponding reset-displacement current.  On a tree skeleton this is
not an additional hypothesis: the bridge displacement is the difference of the
endpoint embedding coordinates and hence is exact.

\begin{theorem}[Exact Hodge degree shift]
\label{thm:v37-degree-shift}
Let $d\mu_{\phi,R}$ be the spectral measure of $\phi_R$ for
$\Delta_{0,R}$ after the constant mode is removed, and let $d\mu_{J,R}$ be the
spectral measure of $J_R$ for $\Delta_{1,R}$.  Then, as measures on the positive
spectrum,
\begin{equation}
 \boxed{d\mu_{J,R}(\lambda)=\lambda\,d\mu_{\phi,R}(\lambda).}
 \label{eq:v37-measure-shift}
\end{equation}
Consequently, whenever an infrared regular-variation limit exists with
\begin{equation}
 d\mu_{\phi}(\lambda)\asymp
 \lambda^{\eta_\phi-1}\,d\lambda,
 \qquad \lambda\downarrow0,
 \label{eq:v37-phi-exponent}
\end{equation}
then
\begin{equation}
 \boxed{\eta_J=\eta_\phi+1.}
 \label{eq:v37-eta-shift}
\end{equation}
Equivalently,
$\langle J,e^{-t\Delta_1}J\rangle\asymp t^{-(\eta_\phi+1)}$.
\end{theorem}

\begin{proof}
The graph differential intertwines the two Laplacians:
\[
 d_0\Delta_{0,R}
 =d_0d_0^*d_0
 =\Delta_{1,R}d_0.
\]
If $E_{0,R}$ denotes the spectral resolution of $\Delta_{0,R}$, then on every
positive spectral band
\[
 \|d_0E_{0,R}(d\lambda)\phi_R\|^2
 =\langle E_{0,R}(d\lambda)\phi_R,
          d_0^*d_0E_{0,R}(d\lambda)\phi_R\rangle
 =\lambda\,d\mu_{\phi,R}(\lambda).
\]
Intertwining identifies this norm with the $\Delta_{1,R}$ spectral mass of
$d_0\phi_R$, proving \eqref{eq:v37-measure-shift}.  Multiplication by
$\lambda$ shifts the infrared power by one.  The time statement follows by
Laplace transform.
\end{proof}

\begin{remark}[Cohomological and displacement currents are different sectors]
\label{rem:v37-deficiency-vs-displacement}
The degree shift applies to the exact skeleton displacement current.  It does
not apply to the non-exact finite-cell deficiency class which carries the
cohomological/cosmological current.  The two currents live in different Hodge
sectors and must not be interchanged in the phenomenology router below.
\end{remark}

\subsection{Burnett locality collapses to one scalar embedding exponent}
\label{subsec:v37-locality-router}

The upstream fourth-cumulant programme expresses the potentially worst quartic
term through three reduced resolvents.  In spectral language the infrared test
is the convergence of
\begin{equation}
 \int_0 d\mu_J(\lambda)\,\lambda^{-3}.
 \label{eq:v37-three-resolvent-test}
\end{equation}
Equivalently, in the time domain one requires absolute summability of the
connected four-time cumulant over the ordered three-separation simplex.  This
is the correct existence test for a continuum Burnett coefficient; a single
one-time or two-time decay estimate is not enough by itself.

\begin{theorem}[Embedding-exponent router]
\label{thm:v37-embedding-router}
Assume the upstream three-resolvent/simplex criterion and the exact displacement
relation $J=d_0\phi$.  If the scalar embedding spectral measure has exponent
$\eta_\phi$, then the quartic renewal correction is classified by
\begin{equation}
 \boxed{
 \begin{array}{rcl}
 \eta_\phi>2
 &\Longrightarrow&
 \text{finite local four-derivative coefficient},\\[2mm]
 \eta_\phi=2
 &\Longrightarrow&
 \text{marginal/logarithmically running coefficient},\\[2mm]
 \eta_\phi<2
 &\Longrightarrow&
 \text{nonlocal form factor; no scale-independent local Burnett constant}.
 \end{array}}
 \label{eq:v37-router}
\end{equation}
Thus the locality question does not require the complete ESM band operator as
its first input.  It requires one infrared scalar: the low-energy spectral
exponent of the embedding coordinate.
\end{theorem}

\begin{proof}
By \ref{thm:v37-degree-shift}, $\eta_J=\eta_\phi+1$.  The
three-resolvent criterion is finite precisely for $\eta_J>3$, marginal at
$\eta_J=3$, and divergent for $\eta_J<3$.  Subtracting one gives
\eqref{eq:v37-router}.
\end{proof}

This version of the router is deliberately stronger in logic and weaker in
input than the often quoted density-of-states saturation value.  To exclude a
local four-derivative continuum one does \emph{not} need to prove
$\eta_\phi=d_s/2$ exactly.  It is enough to prove
\begin{equation}
 \boxed{\eta_\phi<2.}
 \label{eq:v37-weaker-nonlocal-gate}
\end{equation}
The exact density-of-states exponent is needed only to determine the sharp
nonlocal scaling law.

\subsection{Critical decorated skeleton: conditional sharp exponent}
\label{subsec:v37-critical-branch}

The finite-variance critical decorated skeleton used by the microscopic
programme has
\begin{equation}
 d_s=\frac43,
 \qquad
 d_w=3.
 \label{eq:v37-critical-exponents}
\end{equation}
The remaining upstream spectral input is the density-of-states overlap of the
position/embedding observable.

\begin{assumption}[Embedding density-of-states saturation]
\label{ass:v37-dos-saturation}
The centered embedding coordinate has nonzero overlap with the critical slow
sector and therefore
\begin{equation}
 d\mu_\phi(\lambda)\asymp
 \lambda^{d_s/2-1}\,d\lambda,
 \qquad
 \eta_\phi=\frac{d_s}{2}=\frac23.
 \label{eq:v37-dos-saturation}
\end{equation}
\end{assumption}

\begin{theorem}[Conditional sharp nonlocal renewal branch]
\label{thm:v37-sharp-nonlocal}
Under \ref{ass:v37-dos-saturation}, the reset-displacement current obeys
\begin{equation}
 \boxed{\eta_J=\frac53<3.}
 \label{eq:v37-eta-five-thirds}
\end{equation}
Hence it lies strictly in the nonlocal Burnett branch.  With the upstream
finite-size trichotomy
\begin{equation}
 \theta=\frac{d_w}{2}\frac{3-\eta_J}{3},
 \label{eq:v37-theta-general}
\end{equation}
we obtain
\begin{equation}
 \boxed{\theta=\frac23,
 \qquad
 \epsilon_4(R)\asymp R^{2/3}.}
 \label{eq:v37-epsfour-two-thirds}
\end{equation}
The continuum correction is therefore a nonlocal/fractional current form factor
rather than a cutoff-independent local four-derivative coupling.
\end{theorem}

\begin{proof}
The degree shift gives $\eta_J=2/3+1=5/3$.  Substitution into
\eqref{eq:v37-theta-general} with $d_w=3$ gives
\[
 \theta
 =\frac32\frac{3-5/3}{3}
 =\frac23.
\]
The final interpretation is exactly the nonlocal branch of
\ref{thm:v37-embedding-router}.
\end{proof}

More generally, without assuming the sharp value, if
$\eta_\phi<2$ then on this critical skeleton
\begin{equation}
 \boxed{
 \theta=\frac{2-\eta_\phi}{2}>0.}
 \label{eq:v37-theta-phi}
\end{equation}
Thus the distance to the locality threshold has a direct finite-size exponent.

\begin{firewall}[The saturation assumption is not silently promoted]
\label{fire:v37-saturation-firewall}
The exact Hodge shift \eqref{eq:v37-measure-shift} is proved.  The value
$\eta_\phi=2/3$ for the global position/embedding coordinate is a separate
spectral-overlap statement.  Generic slow-mode saturation theorems for bounded
cell-local observables do not automatically prove it for an extended embedding
coordinate.  Until the position overlap is independently certified, the sharp
numbers $5/3$ and $2/3$ remain conditional, whereas the router
\eqref{eq:v37-router} is exact once $\eta_\phi$ is known.
\end{firewall}

\subsection{Consequences for V28--V36 local phenomenology}
\label{subsec:v37-local-programme-scope}

\begin{theorem}[Nonlocal branch excludes a primary local-Wilson interpretation]
\label{thm:v37-local-wilson-no-go}
Suppose $\eta_\phi<2$.  Then the quartic Burnett coefficient does not possess a
scale-independent infinite-volume limit.  Consequently no fixed vector of local
four-derivative Wilson coefficients can represent the \emph{leading} continuum
renewal correction across arbitrarily large scales.  The local vectors computed
in V28--V36 remain valid as finite-window/fixed-regulator matching data, but the
continuum correction requires a nonanalytic operator kernel.
\end{theorem}

\begin{proof}
For $\eta_\phi<2$, \ref{thm:v37-embedding-router} places the physical current in
the divergent simplex branch.  A cutoff-independent local four-derivative
coefficient would imply convergence of the corresponding Burnett coefficient,
contradicting that branch.  Finite-window Taylor coefficients may still be
formed at every fixed cutoff; their growth is precisely the signal that they
must be resummed into a nonlocal kernel.
\end{proof}

This theorem changes the role, not the correctness, of several earlier packets.
In particular:
\begin{enumerate}[label=\textup{(V37.\arabic*)},leftmargin=4.6em]
\item V28's geodesic-volume ray and the V34--V36 fixed-regulator curvature
      vectors remain legitimate local projections of finite-window response.
\item V35--V36 correctly identified the subtraction ambiguity of those local
      vectors; V37 explains why, on the nonlocal branch, that ambiguity is not
      merely a nuisance but evidence that local Wilson coordinates are not the
      primary infrared variables.
\item The V29 local metric field redefinition and its stress-contact descendant
      apply to a local polynomial curvature bank.  They do not automatically
      quotient a singular form factor into local
      $T_{\mu\nu}T^{\mu\nu}$ and $T^2$ contacts.
\item Accordingly, the V30--V32 local contact-ray amplitude relations are
      continuum phenomenology only on the local branch of
      \eqref{eq:v37-router}, or as finite-window benchmark relations before the
      nonlocal resummation is taken.
\end{enumerate}

\begin{firewall}[No nonlocal-to-local field-redefinition shortcut]
\label{fire:v37-nonlocal-field-redefinition}
A local EFT field redefinition may remove operators proportional to the local
Einstein equations and move polynomial curvature-squared terms into local
matter contacts.  A singular kernel such as
$\mathcal O_1F(\Box)\mathcal O_2$ with nonanalytic $F$ is not covered by that
local quotient.  One must compute its on-shell amplitude or spectral
discontinuity directly.  V29's local contact map must therefore not be reused as
a proof that the nonlocal branch has the same $(\kappa_1,\kappa_2)$ ray.
\end{firewall}

\subsection{A finite spectral-overlap certificate}
\label{subsec:v37-overlap-certificate}

The upstream residual can be tested without reconstructing the complete ESM
band.  Let $Q_R$ remove the constant mode of $\Delta_{0,R}$ and define the
normalised cumulative spectral mass of the centered embedding coordinate by
\begin{equation}
 M_R(\Lambda)
 :=
 \frac{\|E_{0,R}([g(R),\Lambda])Q_R\phi_R\|^2}
      {\|Q_R\phi_R\|^2},
 \qquad
 g(R)<\Lambda<\Lambda_{\rm micro}.
 \label{eq:v37-cumulative-mass}
\end{equation}
For two cutoffs $\Lambda_1<\Lambda_2$ in a common scaling window define
\begin{equation}
 \widehat\eta_{\phi,R}(\Lambda_1,\Lambda_2)
 :=
 \frac{
  \log M_R(\Lambda_2)-\log M_R(\Lambda_1)}
  {
  \log\Lambda_2-\log\Lambda_1}.
 \label{eq:v37-local-slope}
\end{equation}

\begin{proposition}[Finite-volume locality certificate under regular variation]
\label{prop:v37-finite-volume-certificate}
Suppose nested balls and nested spectral windows are chosen so that
$g(R)\ll\Lambda_1(R)<\Lambda_2(R)\downarrow0$ and the embedding spectral measure
has a regular-variation exponent $\eta_\phi$.  Then
\begin{equation}
 \widehat\eta_{\phi,R}(\Lambda_1,\Lambda_2)
 \longrightarrow \eta_\phi.
 \label{eq:v37-slope-limit}
\end{equation}
Consequently a certified limiting bound
\begin{equation}
 \limsup_{R\to\infty}\widehat\eta_{\phi,R}
 \le 2-\delta
 \qquad(\delta>0)
 \label{eq:v37-nonlocal-certificate}
\end{equation}
is sufficient to place the physical displacement current in the nonlocal branch,
without proving exact density-of-states saturation.
\end{proposition}

\begin{proof}
For a regularly varying cumulative measure
$M(\Lambda)=\Lambda^{\eta_\phi}L(\Lambda)$ with $L$ slowly varying, the logarithmic
secant slope converges to $\eta_\phi$ on nested fixed-ratio or controlled-ratio
windows.  The second statement follows from \ref{thm:v37-embedding-router}.
\end{proof}

For the sharp critical prediction one may additionally test the ratio
\begin{equation}
 \mathcal C_{\rm sat}(R;c,\Lambda)
 :=
 \frac{M_R(c\Lambda)}{M_R(\Lambda)}-c^{2/3}.
 \label{eq:v37-saturation-residual}
\end{equation}
Vanishing of this residual over a widening infrared window is the direct
finite-volume signature of \ref{ass:v37-dos-saturation}.  It is a much smaller
upstream acquisition than the complete band-covariant form-factor tomography
proposed at the end of V36.

\subsection{What V37 does and does not determine about the nonlocal kernel}
\label{subsec:v37-form-factor-scope}

The spectral trichotomy proves that the continuum kernel is nonanalytic when
$\eta_\phi<2$ and supplies the finite-size divergence exponent.  It does
\emph{not} by itself determine the complete tensor/operator form of the ESM
nonlocal response.  In particular, $\epsilon_4(R)\asymp R^\theta$ is not enough
to assign a unique power $F(\lambda)\propto\lambda^{-\alpha}$ without fixing the
normalisation connecting the Burnett crossover scale, the current bilinear and
the physical external-momentum kernel.

Accordingly the next form-factor calculation should be performed \emph{after}
the locality branch is identified.  On the nonlocal branch its primary outputs
are:
\begin{enumerate}[label=\textup{(N\arabic*)},leftmargin=3.2em]
\item the current/tensor channel on which the singular kernel acts;
\item the small-$\lambda$ spectral discontinuity or regularly varying kernel;
\item its source and species incidence in the complete ESM band;
\item the on-shell amplitude conormal after every local polynomial ambiguity is
      subtracted.
\end{enumerate}

\subsection{Status after V37}
\label{subsec:v37-status}

\subsection*{Proved}
\begin{enumerate}[leftmargin=2.2em]
\item For an exact skeleton displacement current $J=d_0\phi$,
$d\mu_J=\lambda d\mu_\phi$ and $\eta_J=\eta_\phi+1$.
\item Combining that degree shift with the three-resolvent Burnett criterion
reduces locality to the scalar threshold $\eta_\phi=2$.
\item It is sufficient to prove $\eta_\phi<2$ to rule out a scale-independent
local four-derivative continuum coefficient.
\item On a nonlocal branch, the local V28--V36 Wilson vectors are finite-window
matching data rather than the leading continuum variables.
\item V29's local field-redefinition/contact quotient cannot be automatically
applied to a singular nonlocal kernel.
\end{enumerate}

\subsection*{Inherited}
\begin{enumerate}[leftmargin=2.2em]
\item The upstream renewal-current programme supplies the ordered-simplex / three-resolvent
criterion for quartic Burnett existence and the finite-size trichotomy.
\item The finite-variance critical decorated skeleton has the declared
spectral/walk exponents $d_s=4/3$ and $d_w=3$ on that branch.
\item V35--V36 already separated Gaussian/local subtraction data from a
scheme-independent nonanalytic target.
\end{enumerate}

\subsection*{Conditional}
\begin{enumerate}[leftmargin=2.2em]
\item The sharp values $\eta_\phi=2/3$, $\eta_J=5/3$ and
$\epsilon_4(R)\asymp R^{2/3}$ require the embedding density-of-states saturation
condition \ref{ass:v37-dos-saturation}.
\item The skeleton itself remains the declared finite-variance critical branch;
other genealogical universality classes have different $d_s,d_w$ and must be
routed separately.
\end{enumerate}

\subsection*{Open beyond V37}
\begin{enumerate}[leftmargin=2.2em]
\item Prove or falsify the embedding overlap condition directly on the actual
critical decorated skeleton, preferably by the finite spectral packet
\eqref{eq:v37-cumulative-mass}--\eqref{eq:v37-saturation-residual}.
\item Once the branch is fixed, construct the corresponding complete ESM
nonlocal current/tensor form factor and its on-shell source incidence.
\item If instead $\eta_\phi>2$ is found, return to the V36 local-Wilson programme
and complete $B_Q^{\rm full}$ in the ordinary derivative expansion.
\end{enumerate}

\subsection*{Exact next load-bearing theorem}
\begin{center}
\fbox{\parbox{0.92\textwidth}{
\textbf{P4u: embedding-overlap and nonlocality certificate.}
On nested wired balls of the actual finite-variance critical decorated skeleton,
populate the scalar spectral packet of the centered position/embedding
coordinate,
\[
 M_R(\Lambda)
 =\|E_{0,R}([g(R),\Lambda])Q_R\phi_R\|^2/
   \|Q_R\phi_R\|^2,
\]
through a controlled infrared window.  Return first a rigorous bound relative to
the threshold $\eta_\phi=2$, since this alone decides local versus nonlocal
quartic continuum behaviour.  If the data/estimate sharpen to
$\eta_\phi=2/3$, certify the conditional critical prediction
$\eta_J=5/3$ and $\epsilon_4(R)\asymp R^{2/3}$.  Only then construct the full
ESM current/tensor nonlocal form factor and its scheme-independent on-shell
amplitude.  Do not spend the primary effort completing local four-derivative
Wilson ratios before this infrared branch has been decided.}}
\end{center}

\section*{Append-only continuation marker: V37}
\addcontentsline{toc}{section}{Append-only continuation marker: V37}
\textbf{End of V37.}  The complete V1--V36 body above is retained verbatim.
V37 goes upstream from the V36 band-form-factor target to the microscopic
infrared locality classifier.  The physical reset-displacement current is an
exact graph gradient, so its spectral exponent is shifted by one relative to
the embedding coordinate.  This reduces the Burnett locality problem to the
single threshold $\eta_\phi=2$.  Under the generic critical-skeleton overlap
condition $\eta_\phi=d_s/2=2/3$, the displacement exponent is $5/3$ and the
quartic scale grows as $R^{2/3}$, forcing a nonlocal continuum response.  The
density-of-states overlap of the global embedding coordinate remains explicitly
conditional and is promoted to the next finite spectral acquisition.  Thus the
local V28--V36 Wilson/contact programme is retained as a finite-window/local
branch analysis, not assumed to be the primary continuum phenomenology before
the infrared branch is fixed.



% ============================================================================
% V38 APPEND-ONLY CONTINUATION
% ============================================================================
\section{V38 --- Upstream correction: coboundary gauge invariance and the
true Burnett carrier}
\label{sec:v38-coboundary-router}

V37 reduced the proposed critical-skeleton locality test to the infrared
spectral weight of a position/embedding observable $\phi$ and used the exact
Hodge relation $J=d_0\phi$.  An upstream audit of the current layer reveals a
more primitive issue.  The finite-volume Burnett coefficients used throughout
this lineage are derivatives of a Perron pressure of an \emph{additive}
current.  For an exact graph cochain, however, the exponential tilt is a pure
gauge conjugacy.  Thus the exact-current spectral measure may have a rich
low-frequency profile while every rate cumulant of the corresponding additive
functional vanishes identically.

This section isolates that cancellation, retains the correct content of V37,
and replaces the embedding-overlap gate by a cohomological current-selection
gate.  No previous section is rewritten.

\subsection{Exact coboundary conjugacy at finite volume}
\label{subsec:v38-coboundary-conjugacy}

Let $G_R=(V_R,E_R)$ be a finite connected wired graph carrying a primitive
predictive kernel $P_R$.  For one real component of a displacement cochain
$\xi_R$ define
\begin{equation}
 (L_{R,\xi}(k))_{xy}
 :=(P_R)_{xy}\exp\!\bigl(ik\xi_R(x,y)\bigr),
 \qquad
 \Lambda_{R,\xi}(k):=\log\rho(L_{R,\xi}(k)).
 \label{eq:v38-tilt}
\end{equation}
The same argument applies componentwise to a vector-valued displacement and,
with scalar phases tensored with the identity, to a finite matrix-valued cell
transfer.

\begin{theorem}[Coboundary conjugacy and vanishing pressure]
\label{thm:v38-coboundary-conjugacy}
Suppose
\begin{equation}
 \xi_R=d_0\phi_R,
 \qquad
 \xi_R(x,y)=\phi_R(y)-\phi_R(x).
 \label{eq:v38-exact-current}
\end{equation}
Put
\begin{equation}
 D_R(k):=\operatorname{diag}_{x\in V_R}
 \bigl(e^{ik\phi_R(x)}\bigr).
 \label{eq:v38-Dk}
\end{equation}
Then
\begin{equation}
 \boxed{
 L_{R,\xi}(k)=D_R(k)^{-1}P_RD_R(k).}
 \label{eq:v38-similarity}
\end{equation}
Consequently the complete spectrum is independent of $k$ and, on the Perron
branch normalized by $\rho(P_R)=1$,
\begin{equation}
 \boxed{
 \Lambda_{R,\xi}(k)\equiv0.}
 \label{eq:v38-zero-pressure}
\end{equation}
In particular every scaled additive cumulant vanishes:
\begin{equation}
 \boxed{
 \partial_k^m\Lambda_{R,\xi}(0)=0,
 \qquad m\ge1.}
 \label{eq:v38-all-cumulants-zero}
\end{equation}
\end{theorem}

\begin{proof}
The $(x,y)$ entry of the right-hand side of
\eqref{eq:v38-similarity} is
\[
 e^{-ik\phi_R(x)}(P_R)_{xy}e^{ik\phi_R(y)}
 =(P_R)_{xy}e^{ik(\phi_R(y)-\phi_R(x))},
\]
which is \eqref{eq:v38-tilt}.  Similar matrices have the same spectrum.  Since
$P_R$ is stochastic and primitive, its Perron eigenvalue is $1$, so the analytic
Perron branch of the tilted family is identically $1$.  Differentiation gives
\eqref{eq:v38-all-cumulants-zero}.
\end{proof}

The pathwise form is equally transparent.  For a trajectory
$X_0\to\cdots\to X_n$,
\begin{equation}
 \sum_{j=1}^{n}\xi_R(X_{j-1},X_j)
 =\phi_R(X_n)-\phi_R(X_0).
 \label{eq:v38-telescope}
\end{equation}
At fixed $R$ this is an endpoint observable rather than an extensive additive
current.  Equation~\eqref{eq:v38-zero-pressure} is the exact spectral version
of that telescoping identity.

\begin{corollary}[Cohomological invariance of the complete Perron jet]
\label{cor:v38-cohomology-invariance}
For any cochain $h_R$ and any vertex function $\phi_R$,
\begin{equation}
 \boxed{
 \Lambda_{R,h+d\phi}(k)=\Lambda_{R,h}(k).}
 \label{eq:v38-pressure-class}
\end{equation}
Hence every finite-volume Perron response tensor---mean current, Green--Kubo
Hessian, quartic Burnett tensor and all higher cumulants---depends only on the
cohomology class
\begin{equation}
 [\xi_R]\in C^1(G_R)/\operatorname{im}d_0.
 \label{eq:v38-cohomology-class}
\end{equation}
\end{corollary}

\begin{proof}
Use the same diagonal similarity with $L_{R,h}(k)$ in place of $P_R$.
\end{proof}

\subsection{Why the V37 Hodge shift does not contradict the zero pressure}
\label{subsec:v38-hodge-vs-pressure}

The exact V37 identity
\begin{equation}
 d\mu_{d\phi}(\lambda)=\lambda\,d\mu_\phi(\lambda)
 \label{eq:v38-hodge-retained}
\end{equation}
is retained without modification.  It is a Hilbert-space statement about the
spectral content of the instantaneous one-cochain $d\phi$.  By contrast,
\eqref{eq:v38-zero-pressure} is a statement about the scaled cumulants of the
\emph{time-integrated} additive functional.  These objects need not agree term
by term: for an exact current the multi-time Green--Kubo contractions cancel the
instantaneous spectral weight exactly.

\begin{proposition}[Raw current correlations need not determine the additive
pressure]
\label{prop:v38-correlation-pressure-separation}
Let $J_R=d_0\phi_R$.  The two-point quantity
\begin{equation}
 \langle J_R,e^{-t\Delta_{1,R}}J_R\rangle
 =\int e^{-t\lambda}\,d\mu_{J_R}(\lambda)
 \label{eq:v38-raw-correlation}
\end{equation}
may be nonzero and may display slow infrared decay.  Nevertheless the additive
functional generated by the same cochain satisfies
\begin{equation}
 \lim_{n\to\infty}\frac1n
 \operatorname{Cum}_m\!\left(
   \sum_{j=1}^{n}J_R(X_{j-1},X_j)
 \right)=0
 \qquad(m\ge1),
 \label{eq:v38-rate-cumulant-zero}
\end{equation}
whenever the finite-volume Perron cumulant exists.  Thus a spectral exponent
extracted from \eqref{eq:v38-raw-correlation} cannot, by itself, be inserted
into the Burnett pressure trichotomy for a pure coboundary current.
\end{proposition}

\begin{proof}
The first statement is simply the spectral theorem.  The second follows either
from the endpoint identity \eqref{eq:v38-telescope} at fixed finite volume or,
without moment estimates, directly from
\eqref{eq:v38-all-cumulants-zero}.
\end{proof}

\begin{firewall}[Append-only correction to the V37 router]
\label{fire:v38-v37-correction}
V37's degree shift and its finite spectral packet remain correct diagnostics of
the \emph{raw exact-current correlation}.  They are not, without an additional
non-coboundary or joint-scaling bridge, a load-bearing certificate for the
finite-volume Perron/Burnett coefficient.  In particular, the threshold
$\eta_\phi=2$ must not be used to infer a physical quartic Perron divergence
when the current entering the tilt is exactly $d_0\phi$ on every finite wired
ball.  The exact similarity \eqref{eq:v38-similarity} takes precedence.
\end{firewall}

There is a possible distinct mesoscopic problem in which $R$ and the observation
time tend to infinity jointly and the endpoint field itself is the observable.
That joint limit can carry nontrivial roughness or fractional scaling, but it is
not the same object as the finite-volume rate pressure used to define the
Green--Kubo metric and Burnett hierarchy in V1--V37.  V38 keeps those limits
separate.

\subsection{Metric compatibility: an exact skeleton gradient cannot carry the
Perron metric}
\label{subsec:v38-metric-compatibility}

The second derivative of the displacement pressure was used upstream as a
positive metric calibration.  The previous theorem therefore gives an immediate
consistency test.

\begin{corollary}[Pure skeleton-gradient metric no-go]
\label{cor:v38-metric-no-go}
If the displacement cochain used in the Perron metric tilt is supported only on
tree-skeleton bridges, then it is exact and
\begin{equation}
 \boxed{
 C_R=-\partial_k^2\Lambda_{R,\xi}(0)=0.}
 \label{eq:v38-zero-metric}
\end{equation}
Therefore a nondegenerate Perron/Green--Kubo metric cannot be supplied by that
pure skeleton-gradient cochain.  It must come from a non-exact finite-cell
component, a periodic/topological displacement sector, or a different local
pre-skeleton transfer observable.
\end{corollary}

This is not a contradiction with the already-closed E1--E6 continuum theorem.
Those gates assume and calibrate the nondegenerate physical principal symbol on
the finite renewal/action carrier.  The present statement only rules out a
later identification of that metric-bearing current with a pure exact bridge
current on the critical tree.

\subsection{Decorated trees: the pressure sees only the cell-cycle component}
\label{subsec:v38-cell-cycle-carrier}

Let a finite skeleton tree $T_R$ be decorated by finite connected cells $C_v$,
with cells joined only by skeleton bridges.  Every graph cycle is then contained
inside a single cell.  Choose the orthogonal Hodge decomposition
\begin{equation}
 C^1(G_R)
 =\operatorname{im}d_0\oplus\mathcal H^1_R,
 \qquad
 \mathcal H^1_R:=(\operatorname{im}d_0)^\perp,
 \label{eq:v38-hodge-decomp}
\end{equation}
and let $P_{\rm cyc,R}$ denote the second projection.

\begin{theorem}[Cell-cycle reduction of the displacement pressure]
\label{thm:v38-cell-cycle-reduction}
For every displacement cochain $\xi_R$ there is a decomposition
\begin{equation}
 \xi_R=d_0\phi_R+h_R,
 \qquad
 h_R=P_{\rm cyc,R}\xi_R.
 \label{eq:v38-xi-decomposition}
\end{equation}
The complete tilted Perron pressure satisfies
\begin{equation}
 \boxed{
 \Lambda_{R,\xi}(k)=\Lambda_{R,h}(k).}
 \label{eq:v38-pressure-harmonic}
\end{equation}
Moreover, because the skeleton edges are bridges,
\begin{equation}
 \boxed{
 \mathcal H^1_R
 \cong
 \bigoplus_{v\in V(T_R)}H^1(C_v;\mathbb R),}
 \label{eq:v38-cell-direct-sum}
\end{equation}
so $h_R$ is cell supported.  For a $K_4$ cell the local fibre has dimension
three.
\end{theorem}

\begin{proof}
The finite-dimensional Hodge decomposition gives
\eqref{eq:v38-xi-decomposition}; Corollary~\ref{cor:v38-cohomology-invariance}
gives \eqref{eq:v38-pressure-harmonic}.  Since every skeleton edge is a bridge,
no cycle contains such an edge.  Hence the cycle space is the direct sum of the
cell cycle spaces, and dual/Hodge identification gives
\eqref{eq:v38-cell-direct-sum}.
\end{proof}

\begin{corollary}[Sharp alternative for the critical-skeleton Burnett route]
\label{cor:v38-burnett-alternative}
There are two mutually exclusive cases for the actual displacement tilt:
\begin{enumerate}[label=\textup{(\roman*)},leftmargin=2.4em]
\item $P_{\rm cyc,R}\xi_R=0$ for all sufficiently large $R$.  Then the entire
      displacement Perron jet on the critical decorated skeleton is trivial;
      no spatial Burnett correction can be obtained from that tilt.
\item $P_{\rm cyc,R}\xi_R\neq0$.  Then every nontrivial Perron metric/Burnett
      response is carried by the cell-cycle component $h_R$, and its slow
      skeleton loading---not the spectral measure of $d_0\phi_R$---is the
      quantity that must be tested for locality.
\end{enumerate}
\end{corollary}

\subsection{The corrected infrared packet: cycle amplitudes over the skeleton}
\label{subsec:v38-cycle-packet}

For identical $K_4$ cells choose an orthonormal harmonic basis
$\{e_{v,\alpha}\}_{\alpha=1}^{3}$ in every cell and write
\begin{equation}
 h_R
 =\sum_{v\in V(T_R)}\sum_{\alpha=1}^{3}
   a_{R,\alpha}(v)e_{v,\alpha}.
 \label{eq:v38-cycle-coefficients}
\end{equation}
The vector field
\begin{equation}
 a_R(v):=(a_{R,1}(v),a_{R,2}(v),a_{R,3}(v))
 \label{eq:v38-amplitude-field}
\end{equation}
is the finite microscopic datum that tells us how the metric-bearing cell
cohomology is loaded along the critical skeleton.  Let $\Delta_{T,R}$ be the
skeleton scalar Laplacian, $Q_{T,R}$ remove its constant mode and $E_{T,R}$ its
spectral resolution.  Define
\begin{equation}
 \boxed{
 M_R^{\rm cyc}(\Lambda)
 :=
 \frac{
  \sum_{\alpha=1}^{3}
  \|E_{T,R}([g(R),\Lambda])Q_{T,R}a_{R,\alpha}\|^2}
 {
  \sum_{\alpha=1}^{3}\|Q_{T,R}a_{R,\alpha}\|^2}.}
 \label{eq:v38-cycle-spectral-packet}
\end{equation}
Unlike the V37 packet, \eqref{eq:v38-cycle-spectral-packet} is attached to a
cohomologically nontrivial current and therefore is not erased by diagonal
similarity.

If the cell amplitudes form a generic centered local/scenery field on the
finite-variance critical skeleton, the upstream random-scenery calculations
suggest density-of-states loading at $d_s/2=2/3$.  That is useful evidence, not a
substitute for evaluating the actual $a_R$ generated by the Renewal ESM
transfer.  Conversely, if $a_R$ is deterministic, constant, or symmetry-forced
into a fast sector, its infrared exponent can be different.  The correct
physical acquisition is therefore the actual packet
\eqref{eq:v38-cycle-spectral-packet} together with the ordered multi-time
cumulant on the enlarged skeleton--cell state.

\begin{proposition}[Corrected locality router for a non-exact carrier]
\label{prop:v38-corrected-router}
Assume the actual cell-cycle current has a well-defined infrared spectral
exponent $\eta_{\rm cyc}$ in the finite-volume scaling regime and that the
inherited ordered-simplex/three-resolvent dominance hypothesis applies to this
non-exact current.  Then
\begin{equation}
 \boxed{
 \begin{array}{rcl}
  \eta_{\rm cyc}>3 &\Longrightarrow& \text{local quartic coefficient},\\[1mm]
  \eta_{\rm cyc}=3 &\Longrightarrow& \text{marginal/logarithmic},\\[1mm]
  \eta_{\rm cyc}<3 &\Longrightarrow& \text{nonlocal form factor}.
 \end{array}}
 \label{eq:v38-corrected-router}
\end{equation}
There is no automatic $+1$ shift unless the \emph{cycle-amplitude field itself}
is obtained by a skeleton derivative.  The Hodge degree relevant to the
Perron pressure is the degree of the surviving non-exact carrier after the
coboundary quotient.
\end{proposition}

\subsection{A useful optional branch: endogenous random edge embeddings}
\label{subsec:v38-random-edge-branch}

The correction above is already sufficient to replace P4u.  There is, however,
a second diagnostic which explains why a ``generic smooth embedding'' should
not be silently identified with a microscopic path-sum embedding.

Suppose on a finite tree the bridge increments themselves are random with
centered covariance operator $K_R$ on $C^1(T_R)$ satisfying
\begin{equation}
 \kappa_- I\preceq K_R\preceq\kappa_+I,
 \qquad 0<\kappa_-\le\kappa_+<\infty,
 \label{eq:v38-edge-ellipticity}
\end{equation}
and let $d_0\phi_R=\xi_R$ with $Q_R\phi_R=\phi_R$.
If $u_j$ is a normalized nonconstant eigenfunction of
$\Delta_{0,R}$ with eigenvalue $\lambda_j$ and
$v_j=\lambda_j^{-1/2}d_0u_j$, then
\begin{equation}
 \mathbb E\,|\langle u_j,\phi_R\rangle|^2
 =\frac{1}{\lambda_j}
   \langle v_j,K_Rv_j\rangle,
 \label{eq:v38-green-field-mode}
\end{equation}
whereas
\begin{equation}
 \mathbb E\,|\langle v_j,\xi_R\rangle|^2
 =\langle v_j,K_Rv_j\rangle.
 \label{eq:v38-white-edge-mode}
\end{equation}
Thus
\begin{equation}
 \frac{\kappa_-}{\lambda_j}
 \le
 \mathbb E|\langle u_j,\phi_R\rangle|^2
 \le
 \frac{\kappa_+}{\lambda_j},
 \label{eq:v38-green-bounds}
\end{equation}
while the raw edge-current spectral weight is uniformly comparable to the
counting density of states.

\begin{proof}
Since $d_0^*d_0u_j=\lambda_ju_j$ and $Q_R\phi_R=\phi_R$,
\[
 \langle u_j,\phi_R\rangle
 =\lambda_j^{-1}\langle d_0u_j,d_0\phi_R\rangle
 =\lambda_j^{-1/2}\langle v_j,\xi_R\rangle.
\]
Taking the covariance expectation gives
\eqref{eq:v38-green-field-mode}; the bounds follow from
\eqref{eq:v38-edge-ellipticity}.  Equation~\eqref{eq:v38-white-edge-mode} is
immediate.
\end{proof}

In the isotropic case $K_R=\sigma^2I$, this becomes the exact annealed identity
\begin{equation}
 \boxed{
 \mathbb E\,d\mu_{\phi,R}(\lambda)
 =\sigma^2\lambda^{-1}dN_R(\lambda),
 \qquad
 \mathbb E\,d\mu_{\xi,R}(\lambda)
 =\sigma^2 dN_R(\lambda).}
 \label{eq:v38-green-field-measures}
\end{equation}
This shows that an endogenous path-sum embedding is a Green-field-like object,
not a generic bounded observable saturating the density of states.  For
$d_s/2<1$, its normalized $L^2$ mass is gap dominated.  This branch therefore
requires a gap-scaled embedding statistic rather than interpreting the V37
secant slope as an ordinary infinite-volume $L^2$ spectral exponent.

\begin{firewall}[Random-edge branch is diagnostic, not the selected physical
embedding]
\label{fire:v38-random-edge-scope}
Equations~\eqref{eq:v38-green-field-mode}--\eqref{eq:v38-green-field-measures}
are exact for the declared stochastic edge-increment provenance.  The present
Renewal ESM corpus has not yet proved that its physical position coordinate is
generated by such independent or uniformly elliptic edge marks.  These formulas
therefore do not replace the current-selection theorem above; they show that the
embedding provenance must be fixed before assigning a universal exponent to
$\phi$.
\end{firewall}

\subsection{Status after V38}
\label{subsec:v38-status}

\subsection*{Proved}
\begin{enumerate}[leftmargin=2.2em]
\item An exact displacement cochain produces a diagonal similarity of the
      finite-volume tilted transfer, so its complete Perron pressure and every
      rate cumulant vanish identically.
\item The entire displacement Perron jet is invariant under addition of a
      coboundary and depends only on the graph cohomology class of the
      displacement cochain.
\item On a tree-decorated-by-cells geometry, the surviving cohomology is
      cell-supported; a pure skeleton displacement therefore cannot carry a
      nondegenerate Perron metric or Burnett coefficient.
\item V37's Hodge spectral shift remains a correct raw-correlation identity but
      is not by itself a Perron/Burnett locality certificate for a pure exact
      current.
\item For a stochastic path-sum embedding with uniformly elliptic independent
      edge covariance, the embedding spectral weight is inverse-Laplacian
      weighted while the raw edge-current weight is comparable to the density
      of states.
\end{enumerate}

\subsection*{Inherited}
\begin{enumerate}[leftmargin=2.2em]
\item The finite renewal cell carries a nontrivial $K_4$ cycle fibre of rank
      three on the declared branch.
\item The upstream current programme already records the tree-current
      telescoping obstruction and identifies cell-local occupation/non-exact
      currents as the appropriate extensive alternatives.
\item The finite-variance critical skeleton has the declared
      $d_s=4/3$, $d_w=3$ exponents, subject to its existing branch hypotheses.
\end{enumerate}

\subsection*{Conditional}
\begin{enumerate}[leftmargin=2.2em]
\item The local/marginal/nonlocal trichotomy for the surviving non-exact current
      still uses the inherited ordered-simplex/no-cancellation hypothesis.
\item Density-of-states loading of the cell-cycle amplitude field is generic
      random-scenery evidence, not yet a theorem for the actual Renewal ESM
      displacement writer.
\item The random-edge Green-field formulas apply only if that microscopic
      embedding provenance is independently established.
\end{enumerate}

\subsection*{Open beyond V38}
\begin{enumerate}[leftmargin=2.2em]
\item Compute the actual cohomology class of the metric-bearing reset
      displacement on the populated $K_4$ renewal cell and its decorated
      skeleton lift.
\item If the class is nonzero, populate the three-component cell-cycle amplitude
      field $a_R(v)$ and its true multi-time Perron cumulants on nested critical
      balls.
\item If the class vanishes, terminate the critical-skeleton displacement
      Burnett route and identify the distinct non-telescoping microscopic writer
      responsible for higher continuum response.
\end{enumerate}

\subsection*{Exact next load-bearing theorem}
\begin{center}
\fbox{\parbox{0.92\textwidth}{
\textbf{P4v: cohomological displacement landing and true Burnett-current
selection.}
On the actual populated finite renewal cell, compute the physical displacement
cochain $\xi$ entering the common Perron momentum tilt and its cell Hodge
projection
\[
 h=P_{\rm cell}\xi\in H^1(K_4;\mathbb R^3).
\]
Propagate this same writer to nested decorated critical balls without replacing
it by a generic scalar proxy.  If $h=0$, prove the exact negative result that the
critical-skeleton displacement tilt is a coboundary and has zero Perron metric
and zero Burnett jet; higher continuum corrections must then come from another
non-telescoping writer.  If $h\neq0$, form its three-component skeleton loading
$a_R(v)$, populate the finite packet \eqref{eq:v38-cycle-spectral-packet}, and
compute the actual ordered four-time simplex / exact Perron fourth cumulant.
Only this non-exact carrier may be routed into the local/marginal/nonlocal
phenomenology.  Do not use the V37 embedding-gradient exponent as a substitute
for this cohomological landing.}}
\end{center}

\section*{Append-only continuation marker: V38}
\addcontentsline{toc}{section}{Append-only continuation marker: V38}
\textbf{End of V38.}  The complete V1--V37 body above is retained verbatim.
V38 goes further upstream than the V37 embedding-overlap packet.  A pure
skeleton displacement is an exact coboundary; its exponential momentum tilt is
a diagonal similarity and therefore has identically zero Perron pressure and
zero rate cumulants at every finite cutoff.  The Hodge degree shift of V37
remains valid for raw current correlations but cannot by itself feed the
Perron/Burnett hierarchy because the exact-current multi-time contractions
cancel.  The physically relevant additive carrier is the displacement
cohomology class.  On the decorated tree this class is cell supported, so the
next load-bearing object is the actual $K_4$ cell-cycle projection of the
metric-bearing displacement writer and its slow skeleton loading.  An optional
random-edge calculation further shows that a microscopic path-sum embedding is
Green-field-like rather than a generic density-of-states-saturating scalar,
reinforcing the need to fix embedding provenance before assigning an infrared
exponent.


% ================================================================
% V39 append-only continuation
% ================================================================
\section{V39: Cohomological displacement landing on $K_4$ --- the parent no-go,
period compiler, and the first oriented metric-bearing cocycle}
\label{sec:v39-cohomological-landing}

V38 proved that an exact displacement current is invisible to the complete
Perron pressure and therefore cannot carry either the Green--Kubo metric or the
Burnett hierarchy.  The load-bearing question was left in the form
\[
 h=P_{\rm cell}\xi\in H^1(K_4;\mathbb R^3).
\]
There is an additional representation-theoretic issue which must be resolved
before this projection can be interpreted.  The upstream finite-cell current
construction distinguishes two different rank-three modules on the democratic
$K_4$ parent: the reduced vertex/displacement module and the oriented harmonic
cycle module.  They have equal dimension but are not the same $S_4$ module.
This section turns that distinction into an exact displacement-selection
criterion, gives a finite period compiler for an arbitrary measured edge mark,
and exhibits the unique orientation-preserving harmonic spatial writer allowed
by tetrahedral symmetry.

\subsection{The actual landing space is nine-dimensional before symmetry
reduction}
\label{subsec:v39-nine-dimensional-landing}

Fix the orientation
\[
 E_+=\{12,13,14,23,24,34\}
\]
of the six undirected edges of $K_4$, and let
$B\in\mathbb R^{4\times 6}$ be the incidence matrix with $-1$ at the tail and
$+1$ at the head.  For unit edge weights,
\begin{equation}
 L_0:=BB^T=4I_4-\mathbf 1\mathbf 1^T,
 \qquad
 L_0^+=\frac14\left(I_4-\frac14\mathbf 1\mathbf 1^T\right).
 \label{eq:v39-k4-laplacian}
\end{equation}
Hence the orthogonal projector from oriented one-cochains to harmonic
representatives is
\begin{equation}
 \boxed{
 P_H
 =I_6-B^TL_0^+B
 =I_6-\frac14B^TB.}
 \label{eq:v39-hodge-projector}
\end{equation}
For a spatial displacement writer
\[
 \Xi=(\xi_{ij})_{ij\in E_+}\in
 C^1(K_4;\mathbb R^3)\cong\mathbb R^{6\times 3},
\]
its cohomological landing is therefore the completely finite matrix
\begin{equation}
 \boxed{H=P_H\Xi.}
 \label{eq:v39-hodge-landing}
\end{equation}

There is an important counting correction to the shorthand used at the end of
V38.  The scalar cycle fibre has rank three, but a genuinely vector-valued
spatial current has
\begin{equation}
 H^1(K_4;\mathbb R^3)
 \cong H^1(K_4;\mathbb R)\otimes\mathbb R^3,
 \qquad
 \dim H^1(K_4;\mathbb R^3)=9.
 \label{eq:v39-vector-h1-dimension}
\end{equation}
A ``three-component cycle amplitude'' is therefore justified only after an
additional equivariance or frame reduction.  Before that reduction the correct
finite datum is a $3\times3$ period matrix.

Choose the spanning tree
\[
 T_*=\{14,24,34\}.
\]
The three chord cycles give the vector periods
\begin{align}
 \Gamma_{12}(\Xi)&:=\xi_{12}+\xi_{24}-\xi_{14},
 \label{eq:v39-period12}\\
 \Gamma_{13}(\Xi)&:=\xi_{13}+\xi_{34}-\xi_{14},
 \label{eq:v39-period13}\\
 \Gamma_{23}(\Xi)&:=\xi_{23}+\xi_{34}-\xi_{24}.
 \label{eq:v39-period23}
\end{align}
Set
\begin{equation}
 \boxed{
 \mathsf P_\Xi
 :=\bigl[\Gamma_{12}(\Xi)\ \Gamma_{13}(\Xi)\ \Gamma_{23}(\Xi)\bigr]
 \in\mathbb R^{3\times3}.}
 \label{eq:v39-period-matrix}
\end{equation}

\begin{proposition}[Finite cohomology and metric-rank compiler]
\label{prop:v39-period-compiler}
For a vector-valued edge cochain $\Xi$ on connected $K_4$:
\begin{enumerate}[leftmargin=2.2em]
\item $P_H\Xi=0$ if and only if $\mathsf P_\Xi=0$;
\item for $u\in\mathbb R^3$, the scalar current $u\cdot\Xi$ is a
      coboundary if and only if $u^T\mathsf P_\Xi=0$;
\item consequently, the finite-cell Green--Kubo covariance of the displacement
      can be positive definite only if
      \begin{equation}
       \boxed{\operatorname{rank}\mathsf P_\Xi=3.}
       \label{eq:v39-full-period-rank}
      \end{equation}
      On a primitive finite cell this condition is also sufficient under the
      inherited Perron criterion that no nonzero scalar projection of the
      centered displacement be a coboundary.
\end{enumerate}
\end{proposition}

\begin{proof}
The three cycles above form a basis of $H_1(K_4;\mathbb R)$.  A one-cochain on
a connected graph is exact exactly when all of its cycle periods vanish, which
proves the first statement componentwise.  The periods of $u\cdot\Xi$ are the
three entries of $u^T\mathsf P_\Xi$, proving the second.  The inherited
finite-cell Perron theorem states that the covariance is nondegenerate exactly
when no nonzero spatial projection is a coboundary; by the second statement
this is equivalent to the period vectors spanning $\mathbb R^3$.
\end{proof}

Thus the V38 condition ``$h\neq0$'' is enough to detect a nontrivial current,
but a nondegenerate three-dimensional metric requires the stronger full-rank
condition \eqref{eq:v39-full-period-rank}.

\subsection{Canonical quotient-state positions are exactly the wrong writer}
\label{subsec:v39-state-position-no-go}

\begin{theorem}[Finite-state position differences have zero Perron jet]
\label{thm:v39-state-position-no-go}
Let $x_i\in\mathbb R^3$ be arbitrary positions assigned to the four quotient
states and define
\begin{equation}
 \xi^{\rm st}_{ij}=x_j-x_i.
 \label{eq:v39-state-position-writer}
\end{equation}
Then
\begin{equation}
 \xi^{\rm st}=d_0x,
 \qquad
 P_H\xi^{\rm st}=0,
 \qquad
 \mathsf P_{\xi^{\rm st}}=0.
 \label{eq:v39-state-position-exact}
\end{equation}
The tilted transfer is a diagonal similarity at every momentum, so its Perron
pressure is independent of momentum.  In particular its Green--Kubo covariance,
Burnett tensor, and all higher rate cumulants vanish identically.
\end{theorem}

\begin{proof}
Equation~\eqref{eq:v39-state-position-writer} is the definition of a
coboundary.  The Hodge and period conclusions then follow from
Proposition~\ref{prop:v39-period-compiler}.  The diagonal-similarity and
vanishing-pressure statements are exactly the finite-volume coboundary theorem
proved in V38.
\end{proof}

This theorem applies to the most tempting tetrahedral assignment
$x_i\propto r_i$: the fact that the four quotient states can be drawn as the
vertices of a tetrahedron does not make the corresponding edge differences an
extensive displacement current.

\subsection{The democratic $S_4$ parent has a parity obstruction}
\label{subsec:v39-s4-parity-obstruction}

The current finite-cell construction distinguishes, as oriented $S_4$ modules,
\begin{equation}
 C^0_0(K_4;\mathbb R)\cong {\rm Std},
 \qquad
 H^1(K_4;\mathbb R)\cong {\rm Std}\otimes {\rm sgn}.
 \label{eq:v39-two-triplets}
\end{equation}
Let $V_{\rm sp}\cong{\rm Std}$ denote the ordinary polar spatial-vector
representation carried by the tetrahedral displacement frame.

\begin{theorem}[$S_4$ obstruction to a parent-selected harmonic spatial writer]
\label{thm:v39-s4-no-go}
There is no nonzero $S_4$-covariant polar-vector harmonic displacement tensor on
the fully democratic $K_4$ parent:
\begin{equation}
 \boxed{
 \bigl(H^1(K_4;\mathbb R)\otimes V_{\rm sp}\bigr)^{S_4}=0.}
 \label{eq:v39-s4-invariant-zero}
\end{equation}
Equivalently,
\begin{equation}
 \operatorname{Hom}_{S_4}
 \bigl({\rm Std}\otimes{\rm sgn},{\rm Std}\bigr)=0.
 \label{eq:v39-s4-hom-zero}
\end{equation}
Therefore a nonzero cohomological spatial displacement cannot be selected from
the democratic parent by $S_4$ symmetry and linear covariance alone.
\end{theorem}

\begin{proof}
Both real modules in \eqref{eq:v39-two-triplets} are irreducible and self-dual,
but they are inequivalent because of the sign twist.  The invariant tensor
space is naturally the equivariant Hom-space displayed in
\eqref{eq:v39-s4-hom-zero}.  Schur's lemma makes that Hom-space zero.
\end{proof}

\begin{corollary}[What the rank-three coincidence does not do]
\label{cor:v39-rank-coincidence-no-map}
The numerical equality
\[
 \dim C^0_0(K_4)=\dim H^1(K_4)=3
\]
does not provide a canonical map from the displacement triplet to the cycle
triplet.  Any such nonzero map requires additional orientation/parity-breaking
structure or an independently populated edge-record incidence.
\end{corollary}

This is the precise P4v answer on the democratic parent.  The current corpus
supplies a generic reset-displacement cochain in the Perron tilt, but it does
not supply six vector edge values from which a nonzero harmonic spatial class
could be reconstructed.  If one instead interprets the displacement triplet as
four state positions, Theorem~\ref{thm:v39-state-position-no-go} forces the
harmonic class to vanish.

\subsection{Orientation-preserving symmetry opens a unique harmonic channel}
\label{subsec:v39-a4-channel}

After restriction to the proper tetrahedral group $A_4$, the sign character is
trivial and both modules in \eqref{eq:v39-two-triplets} become the same real
irreducible triplet.  Hence
\begin{equation}
 \boxed{
 \dim\operatorname{Hom}_{A_4}
 \bigl(H^1(K_4;\mathbb R),V_{\rm sp}\bigr)=1.}
 \label{eq:v39-a4-hom-one}
\end{equation}
Thus an orientation choice does not produce an arbitrary nine-parameter
spatial cycle tensor: tetrahedral covariance leaves one real scale.

Choose the normalized tetrahedral frame
\begin{equation}
 \begin{aligned}
 r_1&=\frac1{\sqrt3}(1,1,1), &
 r_2&=\frac1{\sqrt3}(1,-1,-1),\\
 r_3&=\frac1{\sqrt3}(-1,1,-1), &
 r_4&=\frac1{\sqrt3}(-1,-1,1).
 \end{aligned}
 \label{eq:v39-tetra-frame}
\end{equation}
It obeys
\begin{equation}
 \sum_i r_i=0,
 \qquad
 r_i\cdot r_j=-\frac13\quad(i\neq j).
 \label{eq:v39-tetra-identities}
\end{equation}
For a real scale $\ell$, define the oriented edge writer
\begin{equation}
 \boxed{
 h_{ij}=\ell\,r_i\times r_j,
 \qquad h_{ji}=-h_{ij}.}
 \label{eq:v39-cross-writer}
\end{equation}

\begin{theorem}[Unique $A_4$-covariant harmonic displacement candidate]
\label{thm:v39-a4-cross-writer}
The writer \eqref{eq:v39-cross-writer} is harmonic and non-exact.  It is
$A_4$-covariant as a polar spatial vector.  Under an odd tetrahedral permutation
$g\in S_4\setminus A_4$,
\begin{equation}
 h_{g i,g j}
 =-R_g h_{ij},
 \label{eq:v39-cross-parity}
\end{equation}
where $R_g\in O(3)$ is the tetrahedral spatial action.  Consequently it spans
the one-dimensional space in \eqref{eq:v39-a4-hom-one} and is forbidden by the
unbroken $S_4$ parent but allowed once an orientation has been selected.
\end{theorem}

\begin{proof}
Antisymmetry is immediate.  At each vertex,
\[
 \sum_{j\ne i}h_{ij}
 =\ell\,r_i\times\sum_{j\ne i}r_j
 =-\ell\,r_i\times r_i=0,
\]
so $d_0^*h=0$ and $h$ is the harmonic representative of its cohomology class.
It is nonzero and hence non-exact because exact and harmonic subspaces are
orthogonal.  For the tetrahedral action $r_{g i}=R_gr_i$,
\[
 (R_gu)\times(R_gv)=\det(R_g)R_g(u\times v).
\]
On the standard tetrahedral representation $\det R_g={\rm sgn}(g)$, proving
\eqref{eq:v39-cross-parity}.  The $A_4$ covariance and uniqueness up to scale
then follow from \eqref{eq:v39-a4-hom-one}.
\end{proof}

The fundamental periods of this candidate are explicit:
\begin{equation}
 \begin{split}
 \Gamma_{12}(h)&=\frac{4\ell}{3}(-1,1,-1),\\
 \Gamma_{13}(h)&=\frac{4\ell}{3}(-1,1,1),\\
 \Gamma_{23}(h)&=\frac{4\ell}{3}(1,1,1).
 \end{split}
 \label{eq:v39-cross-periods}
\end{equation}
Therefore
\begin{equation}
 \boxed{
 \det\mathsf P_h=\frac{256}{27}\ell^3.}
 \label{eq:v39-cross-period-det}
\end{equation}
For $\ell\neq0$ the period map has full rank three, so this candidate satisfies
the exact metric-rank criterion \eqref{eq:v39-full-period-rank}.

\subsection{The non-exact writer is naturally a covering-space displacement}
\label{subsec:v39-covering-interpretation}

There is no contradiction between the finite quotient having four states and a
non-exact displacement accumulating without bound.  The correct object is a
Markov-additive lift.

\begin{proposition}[Period map and translation cover]
\label{prop:v39-translation-cover}
Let $\Xi\in C^1(K_4;\mathbb R^3)$ be any antisymmetric edge displacement.  Its
cycle periods define a homomorphism
\begin{equation}
 {\rm Per}_\Xi:H_1(K_4;\mathbb Z)\longrightarrow\mathbb R^3.
 \label{eq:v39-period-homomorphism}
\end{equation}
On the universal cover $\widetilde K_4$, the pulled-back cochain is exact:
there exists a lifted position $\widetilde x$ such that
\begin{equation}
 d\widetilde x=\pi^*\Xi.
 \label{eq:v39-cover-exact}
\end{equation}
For every lifted path, the additive displacement equals the endpoint difference
of $\widetilde x$.  If the period image is a discrete rank-three lattice, the
corresponding abelian cover is a three-dimensional translation cover of the
finite renewal cell.
\end{proposition}

\begin{proof}
The universal cover is a tree, so it has trivial first cohomology.  Fix a
lifted basepoint and define $\widetilde x$ by path integration of $\pi^*\Xi$;
uniqueness of paths on the tree makes this well-defined.  A deck transformation
changes $\widetilde x$ by the period of its projected cycle, which proves the
period homomorphism and the translation-cover statement.
\end{proof}

For the candidate \eqref{eq:v39-cross-writer},
\eqref{eq:v39-cross-period-det} shows that the three fundamental periods are
independent whenever $\ell\neq0$.  Thus $h$ may be read as a genuine
three-dimensional translation cocycle on a cover even though it is not a
state-position difference on the four-state quotient.  This is the correct
finite-state interpretation of a metric-bearing reset displacement.

\subsection{Exact democratic Perron metric and quartic jet of the oriented
candidate}
\label{subsec:v39-perron-jet}

The candidate is sufficiently rigid that its finite-cell Perron jet can be
computed completely through fourth order.  Let
\begin{equation}
 P_{ij}=\frac13\mathbf 1_{i\ne j},
 \qquad
 L_h(k)_{ij}=P_{ij}e^{i k\cdot h_{ij}},
 \qquad
 \Lambda_h(k)=\log\rho_{\rm Perron}(L_h(k)).
 \label{eq:v39-tilted-democratic}
\end{equation}

\begin{proposition}[Metric and quartic Perron jet]
\label{prop:v39-perron-jet}
For \eqref{eq:v39-cross-writer},
\begin{equation}
 \boxed{
 C_{ab}
 =\frac{8\ell^2}{27}\,\delta_{ab}.}
 \label{eq:v39-cross-metric}
\end{equation}
Moreover the pressure is even in $k$ and has expansion
\begin{equation}
 \boxed{
 \Lambda_h(k)
 =-\frac{4\ell^2}{27}|k|^2
 +\frac{4\ell^4}{729}
 \left[
 |k|^4-2\sum_{a=1}^{3}k_a^4
 \right]
 +O(|k|^6).}
 \label{eq:v39-cross-pressure}
\end{equation}
Equivalently, with the metric calibration
\begin{equation}
 \epsilon^2:=\frac13\operatorname{tr}C=\frac{8\ell^2}{27},
 \label{eq:v39-epsilon-scale}
\end{equation}
one has
\begin{equation}
 \boxed{
 \Lambda_h(k)
 =-\frac12\epsilon^2|k|^2
 +\frac{\epsilon^4}{16}
 \left[
 |k|^4-2\sum_{a=1}^{3}k_a^4
 \right]
 +O(|k|^6).}
 \label{eq:v39-cross-pressure-epsilon}
\end{equation}
Thus the quadratic response is isotropic while the first finite-cell quartic
response retains the tetrahedral anisotropy permitted by $A_4$.
\end{proposition}

\begin{proof}
For every vertex $i$,
\[
 \sum_{j\ne i}P_{ij}h_{ij}=0.
\]
Hence the edge increments are martingale differences for the democratic chain,
so all nonzero-lag two-point covariances vanish.  Since
\[
 |r_i\times r_j|^2
 =1-(r_i\cdot r_j)^2
 =\frac89,
\]
and $A_4$ acts irreducibly on the spatial frame, the one-step covariance is a
scalar matrix.  Its trace is $8\ell^2/9$, which gives
\eqref{eq:v39-cross-metric}.

Reversibility and antisymmetry give
$L_h(-k)=L_h(k)^T$, so the Perron pressure is even.  At quartic order the
$A_4$-invariant polynomial space is spanned by $|k|^4$ and
$\sum_a k_a^4$.  Direct fourth-order perturbation of the $4\times4$ transfer
matrix gives
\begin{align*}
 \Lambda_h(te_1)
 &=-\frac{4\ell^2}{27}t^2
   -\frac{4\ell^4}{729}t^4+O(t^6),\\
 \Lambda_h\bigl(t(e_1+e_2)\bigr)
 &=-\frac{8\ell^2}{27}t^2+O(t^6).
\end{align*}
These two evaluations determine the two quartic invariant coefficients and
produce \eqref{eq:v39-cross-pressure}.  Substituting
\eqref{eq:v39-epsilon-scale} gives
\eqref{eq:v39-cross-pressure-epsilon}.
\end{proof}

The quartic tensor is therefore already nonzero on one cell, but it should not
be promoted to the V28--V36 gravitational Wilson packet.  It is the finite-cell
Perron cumulant of a particular orientation-sensitive displacement cocycle.
Its continuum interpretation still requires the same-writer propagation to
the decorated geometry and the appropriate shell/observable matching.

\begin{firewall}[The $A_4$ cross-product writer is selected by symmetry only up
to a channel, not as the physical Renewal--ESM writer]
\label{fire:v39-candidate-scope}
The current corpus proves that the reset displacement in the common Perron jet
may be a general bounded edge cochain and separately proves the two-triplet
module statement.  It does not yet identify the six physical edge-displacement
vectors with \eqref{eq:v39-cross-writer}.  Equations
\eqref{eq:v39-cross-writer}--\eqref{eq:v39-cross-pressure-epsilon} therefore
construct the unique $A_4$-covariant harmonic candidate, not a newly inherited
microscopic datum.  Selecting this channel requires a same-record incidence
from the signed/oriented Renewal sector to the spatial displacement writer.
Absent that incidence, the actual period matrix
\eqref{eq:v39-period-matrix} remains unpopulated.
\end{firewall}

\subsection{The P4v router after the finite-cell landing}
\label{subsec:v39-router}

The finite-cell question is now algebraic and has three distinct outcomes.
For the actual six directed spatial edge marks, compute
\begin{equation}
 \mathsf P_\xi
 =\bigl[\Gamma_{12}(\xi)\ \Gamma_{13}(\xi)\ \Gamma_{23}(\xi)\bigr].
 \label{eq:v39-actual-period-packet}
\end{equation}
Then:
\begin{equation}
 \boxed{
 \begin{array}{rcl}
 \operatorname{rank}\mathsf P_\xi=0
 &\Longrightarrow&
 \text{exact current; zero spatial Perron jet},\\[1mm]
 0<\operatorname{rank}\mathsf P_\xi<3
 &\Longrightarrow&
 \text{non-exact but metrically degenerate current},\\[1mm]
 \operatorname{rank}\mathsf P_\xi=3
 &\Longrightarrow&
 \text{cohomologically admissible nondegenerate metric carrier}.
 \end{array}}
 \label{eq:v39-period-router}
\end{equation}
On the unbroken democratic parent, a canonically $S_4$-selected polar writer
falls in the first line.  On an oriented $A_4$ branch, the candidate
\eqref{eq:v39-cross-writer} falls in the third line for every $\ell\neq0$.

\subsection{Status after V39}
\label{subsec:v39-status}

\subsection*{Proved}
\begin{enumerate}[leftmargin=2.2em]
\item The vector-valued $K_4$ cohomological displacement landing is a
      nine-dimensional object before symmetry reduction, and its exact Hodge
      projector is \eqref{eq:v39-hodge-projector}.
\item Three fundamental vector cycle periods give a finite $3\times3$ compiler.
      Full rank is exactly the cohomological condition required for a
      nondegenerate three-dimensional Green--Kubo metric.
\item Any quotient-state position-difference writer is exact and has zero
      Perron metric and zero Burnett jet.
\item The democratic $S_4$ parent cannot canonically select a nonzero
      polar-vector harmonic displacement: the displacement and oriented cycle
      triplets differ by the sign representation.
\item After orientation reduction to $A_4$, there is a unique harmonic spatial
      channel up to scale.  The explicit writer $h_{ij}=\ell r_i\times r_j$ has
      full-rank periods and therefore defines a genuine three-dimensional
      translation cover.
\item For the democratic kernel this candidate has the exact isotropic metric
      \eqref{eq:v39-cross-metric} and quartic Perron expansion
      \eqref{eq:v39-cross-pressure}.
\end{enumerate}

\subsection*{Inherited}
\begin{enumerate}[leftmargin=2.2em]
\item The common finite-cell Perron jet accepts a bounded reset-displacement
      edge cochain and gives a positive metric exactly when no nonzero scalar
      projection is a coboundary.
\item The current $K_4$ construction distinguishes the reduced-vertex
      displacement module $C^0_0\cong{\rm Std}$ from the oriented cycle module
      $H^1\cong{\rm Std}\otimes{\rm sgn}$; after restriction to $A_4$ both are
      the real tetrahedral triplet.
\item The democratic $K_4$ kernel is primitive and reversible, and the finite
      cell carries a protected rank-three cycle fibre.
\end{enumerate}

\subsection*{Conditional}
\begin{enumerate}[leftmargin=2.2em]
\item The cross-product writer is the unique symmetry-allowed $A_4$ harmonic
      spatial channel but has not yet been derived as the actual Renewal--ESM
      reset-displacement record.
\item Its scale $\ell$ can be fixed from the Green--Kubo metric calibration once
      the channel is physically selected, but the sign/orientation and the
      same-record incidence selecting the channel remain upstream data.
\item Any use of the quartic tensor in continuum phenomenology requires
      propagation of this same non-exact writer through the decorated
      skeleton and the already established continuum matching machinery.
\end{enumerate}

\subsection*{Open beyond V39}
\begin{enumerate}[leftmargin=2.2em]
\item Populate the six actual vector reset-displacement marks of the physical
      K4 momentum tilt and compute \eqref{eq:v39-actual-period-packet}.
\item Derive, rather than posit, whether the Lorentzian signed/oriented sector
      reduces the democratic $S_4$ parent to the $A_4$ channel and selects the
      cross-product intertwiner.
\item If the period matrix has rank three, propagate its three cycle-period
      fields on the decorated critical geometry and compute the true Perron
      fourth cumulant/nonlocality packet for this same writer.
\item If the period matrix vanishes, terminate the K4 quotient-displacement
      route and locate the distinct Markov-additive external translation mark
      or port-resolved record that carries the continuum metric.
\end{enumerate}

\subsection*{Exact next load-bearing theorem}
\begin{center}
\fbox{\parbox{0.92\textwidth}{
\textbf{P4w: oriented displacement incidence and covering-space landing.}
On one physical Renewal--ESM $K_4$ cell, acquire the six directed spatial
reset-displacement vectors used in the \emph{same} momentum tilt as the
Green--Kubo metric.  Compute the three vector cycle periods
\eqref{eq:v39-period12}--\eqref{eq:v39-period23} and the rank of
$\mathsf P_\xi$.  In parallel, derive the incidence of the signed Lorentzian /
anti-renewal orientation datum with this edge writer.  If the measured period
packet is $A_4$-covariant, test whether it is proportional to the unique
cross-product tensor \eqref{eq:v39-cross-writer}; only then use the metric
calibration to fix $\ell$ and propagate the same covering-space displacement
through the decorated critical geometry.  If the rank is zero, close this route
with the exact coboundary no-go and identify the separate non-telescoping
Markov-additive translation record.  Do not identify the displacement triplet
with the cycle triplet merely because both have rank three.}}
\end{center}

\section*{Append-only continuation marker: V39}
\addcontentsline{toc}{section}{Append-only continuation marker: V39}
\textbf{End of V39.}  The complete V1--V38 body above is retained verbatim.
V39 resolves the finite-cell landing hidden inside P4v.  A vector-valued
spatial displacement has a $3\times3$ period matrix, and a nondegenerate metric
requires full period rank.  Quotient-state position differences are exact and
therefore have an identically vanishing Perron jet.  More strongly, the
$S_4$-democratic parent cannot canonically select a harmonic polar-vector
spatial cocycle because its displacement and oriented cycle triplets differ by
the sign representation.  After restricting to the orientation-preserving
$A_4$ subgroup there is a unique harmonic spatial channel up to scale; the
explicit cross-product cocycle $h_{ij}=\ell r_i\times r_j$ has full-rank periods,
defines a three-dimensional translation cover, and yields an isotropic
Green--Kubo metric together with a computable tetrahedral quartic Perron tensor.
This channel is a symmetry-allowed candidate, not yet an inherited physical
writer.  The next load-bearing datum is therefore the same-record incidence
that selects the actual six displacement edge marks and determines their
cycle-period matrix.



% ============================================================================
% V40 append-only continuation
% ============================================================================
\clearpage
\section{P4w upstream correction: finite coboundaries, thermodynamic transport, and the K4 representation fork}
\label{sec:v40-p4w}

\subsection{Why the six-edge acquisition is not yet the first load-bearing datum}

V39 deliberately stopped before identifying the symmetry-allowed harmonic
$K_4$ writer with the physical Renewal--ESM displacement.  The upstream audit
required for P4w exposes an even earlier distinction.  Three objects had been
kept too close in the later phenomenology line:
\begin{enumerate}[leftmargin=2.2em]
\item a bounded edge reward on one \emph{fixed finite primitive cell}, for
      which Perron--Frobenius perturbation gives an ordinary analytic pressure;
\item the increment of an unbounded position/embedding coordinate on the
      \emph{thermodynamic critical skeleton}; and
\item the finite $K_4$ cycle fibre used downstream as an internal matter
      carrier.
\end{enumerate}
They can be represented by edge cochains, but they are not interchangeable.
In particular, the exact diagonal-similarity theorem of V38 is a finite-carrier
statement.  It cannot be exported through the thermodynamic limit without
checking boundedness of the conjugating multiplication operator and the order
of the large-time and large-volume limits.

This section records the corresponding correction before any six-vector
physical acquisition is attempted.

\subsection{Fixed finite carrier: the V38 coboundary theorem is exact}

Let $P$ be an irreducible stochastic matrix on a finite set $X$ and let
$\phi:X\to\mathbb R$.  Put
\begin{equation}
 \xi(x,y)=\phi(y)-\phi(x),
 \qquad
 (P_t)_{xy}=P_{xy}e^{t\xi(x,y)}.
 \label{eq:v40-finite-coboundary-tilt}
\end{equation}
With $D_t=\operatorname{diag}(e^{t\phi(x)})$,
\begin{equation}
 P_t=D_t^{-1} P D_t.
 \label{eq:v40-finite-similarity}
\end{equation}
Since $D_t$ is an everywhere-defined bounded invertible matrix, the complete
spectrum is unchanged.

\begin{theorem}[Fixed-finite coboundary no-go]
\label{thm:v40-fixed-finite-coboundary}
For the finite primitive chain above, the Perron pressure of an exact reward is
identically zero:
\begin{equation}
 \Lambda_{\xi}^{\rm fin}(t)
 :=\log\rho(P_t)=0.
 \label{eq:v40-finite-pressure-zero}
\end{equation}
Consequently every extensive Perron cumulant of $\xi$ vanishes.  Equivalently,
for every path,
\begin{equation}
 \sum_{j=1}^{n}\xi(X_{j-1},X_j)
 =\phi(X_n)-\phi(X_0),
 \label{eq:v40-finite-telescope}
\end{equation}
and the right-hand side is uniformly bounded on the fixed finite state space.
\end{theorem}

\begin{proof}
Equation \eqref{eq:v40-finite-similarity} preserves all eigenvalues, so the
Perron eigenvalue remains $1$.  Equation \eqref{eq:v40-finite-telescope} is the
pathwise telescoping identity.  Because $X$ is finite,
$|\phi(X_n)-\phi(X_0)|\le 2\|\phi\|_\infty$, so every cumulant divided by $n$
vanishes as $n\to\infty$.
\end{proof}

Thus V38 remains exact on every fixed finite cell or fixed finite wired ball.
The correction below concerns only its attempted extrapolation to an unbounded
thermodynamic embedding.

\subsection{The thermodynamic limit need not commute with the finite coboundary no-go}

The endpoint in \eqref{eq:v40-finite-telescope} need not remain bounded when the
state space itself grows.  The simplest example already closes the logical
point.

\begin{proposition}[Exact one-dimensional counterexample to limit interchange]
\label{prop:v40-z-counterexample}
Let $(X_n)$ be the simple random walk on $\mathbb Z$, let $\phi(x)=x$, and set
$\xi(x,y)=y-x=d\phi(x,y)$.  Then
\begin{equation}
 S_n:=\sum_{j=1}^{n}\xi(X_{j-1},X_j)=X_n-X_0
 \label{eq:v40-z-telescope}
\end{equation}
is an exact coboundary sum, but its real moment-generating pressure is
\begin{equation}
 \Lambda_\infty(t)
 :=\lim_{n\to\infty}\frac1n\log\mathbb E_0 e^{tS_n}
 =\log\cosh t,
 \label{eq:v40-z-pressure}
\end{equation}
so
\begin{equation}
 \Lambda_\infty''(0)=1.
 \label{eq:v40-z-diffusivity}
\end{equation}
By contrast, on every fixed finite reflecting interval $[-R,R]$, the same
increment is a bounded state coboundary and Theorem~\ref{thm:v40-fixed-finite-coboundary}
gives
\begin{equation}
 \Lambda_R(t)=0.
 \label{eq:v40-z-finite-pressure}
\end{equation}
Hence
\begin{equation}
 \lim_{R\to\infty}\lim_{n\to\infty}
 \frac1n\log\mathbb E_0^{(R)}e^{tS_n}=0
 \quad\hbox{while}\quad
 \lim_{n\to\infty}\lim_{R\to\infty}
 \frac1n\log\mathbb E_0^{(R)}e^{tS_n}
 =\log\cosh t.
 \label{eq:v40-noncommuting-limits}
\end{equation}
\end{proposition}

\begin{proof}
On $\mathbb Z$, the increments are independent and take the values $\pm1$ with
probability $1/2$, so
$\mathbb E e^{tS_n}=(\cosh t)^n$.  On a fixed reflecting interval,
$S_n=X_n-X_0$ is bounded by $2R$ and the finite-state similarity theorem
applies.  Taking the limits in the two displayed orders gives
\eqref{eq:v40-noncommuting-limits}.
\end{proof}

The formal identity $P_t=D_t^{-1}PD_t$ still exists on $\mathbb Z$, but for
real $t\ne0$ the multiplier $D_t f(x)=e^{t x}f(x)$ is unbounded on the natural
unweighted spaces.  There is therefore no bounded similarity theorem that
identifies the relevant thermodynamic pressure with the untitled spectrum.
For Fourier tilt $t=ik$ the multiplier is unitary on $\ell^2$, but the infinite
walk has no isolated Perron eigenvalue whose analytic branch could play the
finite-state role.  The finite Perron theorem and the infinite transport limit
are genuinely different analytic objects.

\begin{firewall}[Append-only correction to the strongest reading of V38--V39]
\label{fw:v40-v38-correction}
The statement ``an exact displacement cannot carry the physical metric'' is
valid on a fixed finite primitive carrier with a bounded potential.  It is too
strong after a thermodynamic limit in which the potential becomes unbounded.
An unbounded skeleton position increment may be exact as a graph cochain and
still have nonzero diffusive or anomalous transport because the endpoint itself
grows with time.  V38's finite similarity theorem is retained; only the
thermodynamic extrapolation is withdrawn.
\end{firewall}

\subsection{Pre-boundary transport is the correct finite-volume approximation}

The previous proposition also tells us how the finite critical balls should be
used.  Let $X$ be the infinite skeleton walk, $B_R$ a ball, and
\begin{equation}
 \tau_R:=\inf\{n\ge0:X_n\notin B_R\}.
 \label{eq:v40-exit-time}
\end{equation}
Couple a finite-volume walk $X^{(R)}$ to $X$ so that they agree up to $\tau_R$.
For a real-valued embedding $\phi$ define, for Fourier momentum $k$,
\begin{equation}
 \Psi_{R,n}(k)
 :=\mathbb E_o\!
 \left[
 e^{ik(\phi(X_n)-\phi(o))}\mathbf 1_{\{\tau_R>n\}}
 \right].
 \label{eq:v40-preboundary-characteristic}
\end{equation}
Then, without any moment assumption,
\begin{equation}
 \left|
 \mathbb E_o e^{ik(\phi(X_n)-\phi(o))}
 -\Psi_{R,n}(k)
 \right|
 \le \mathbb P_o(\tau_R\le n).
 \label{eq:v40-preboundary-error}
\end{equation}

\begin{proposition}[Thermodynamic-order finite-volume criterion]
\label{prop:v40-preboundary-window}
Assume a walk dimension $d_w$ in the sense that for every $\delta>0$ one can
choose $n(R)$ with
\begin{equation}
 n(R)\longrightarrow\infty,
 \qquad
 \frac{n(R)}{R^{d_w-\delta}}\longrightarrow0,
 \qquad
 \mathbb P_o(\tau_R\le n(R))\longrightarrow0.
 \label{eq:v40-preboundary-window}
\end{equation}
Then the characteristic functions of the finite balls in the pre-boundary
window converge to the infinite-volume displacement characteristic function.
Any transport coefficient defined from this regime must therefore be computed
before the $n\to\infty$ fixed-$R$ equilibrium limit which annihilates bounded
coboundaries.
\end{proposition}

For the critical finite-variance skeleton inherited by the renewal programme,
$d_w=3$.  Thus the natural mesoscopic window is $n\ll R^3$ up to the usual
quenched/annealed and slowly varying refinements.  No nonzero $K_4$ cycle
period is required merely to prevent the skeleton endpoint displacement from
telescoping: the endpoint can be unbounded in the thermodynamic order.

\subsection{Representation-convention firewall on the K4 edge space}

The source audit also reveals a second apparent conflict which must be resolved
before using representation theory to select a spatial writer.  Two different
$S_4$ actions on six edge coordinates occur in the programme.

Let $U\cong\mathbb R^4$ be the vertex permutation representation and write
\begin{equation}
 U\cong\mathbf 1\oplus\mathbf 3,
 \label{eq:v40-vertex-perm-decomp}
\end{equation}
where $\mathbf3$ is the reduced standard representation.

\paragraph{Unsigned edge amplitudes.}
Let $E^+$ be the permutation representation on the six unordered two-subsets
$\{i,j\}$.  Its character decomposes as
\begin{equation}
 E^+\cong\mathbf1\oplus\mathbf2\oplus\mathbf3.
 \label{eq:v40-unsigned-edge-decomp}
\end{equation}
This is the module appropriate to genuinely unsigned bond amplitudes or
conductances.

\paragraph{Oriented antisymmetric one-cochains.}
Let
\begin{equation}
 E^-:=\{\omega_{ij}:\omega_{ji}=-\omega_{ij}\}.
 \label{eq:v40-oriented-edge-def}
\end{equation}
Naturally,
\begin{equation}
 E^-\cong\bigwedge^2 U
 \cong\mathbf3\oplus\mathbf3',
 \qquad
 \mathbf3'=\mathbf3\otimes{\rm sgn}.
 \label{eq:v40-oriented-edge-decomp}
\end{equation}
The ordinary graph coboundary
\begin{equation}
 d:C^0_0(K_4)\longrightarrow E^- ,
 \qquad
 (df)_{ij}=f_j-f_i,
 \label{eq:v40-natural-coboundary}
\end{equation}
is $S_4$-equivariant and has image $\mathbf3$.  Hence the natural oriented
graph cohomology carries
\begin{equation}
 H^1_{\rm or}(K_4;\mathbb R)
 =E^-/\operatorname{im}d
 \cong\mathbf3'.
 \label{eq:v40-oriented-h1}
\end{equation}

\begin{theorem}[K4 representation-convention fork]
\label{thm:v40-representation-fork}
The decompositions
\begin{equation}
 \mathbf1\oplus\mathbf2
 \qquad\hbox{and}\qquad
 \mathbf3'
 \label{eq:v40-two-h1-looking-modules}
\end{equation}
which appear in different finite-cell constructions refer to different edge
representations.  The unsigned two-subset module contains the
$\mathbf1\oplus\mathbf2$ sector after removal of its reduced-vertex
$\mathbf3$ summand, whereas the natural antisymmetric graph-cochain quotient is
$\mathbf3'$.  In particular a path-reversal-odd displacement cochain belongs to
$E^-$ and must use \eqref{eq:v40-oriented-h1}; an unsigned edge weight belongs
to $E^+$ and must not inherit oriented-cochain conclusions merely from the
common dimension three.
\end{theorem}

\begin{proof}
Equation \eqref{eq:v40-unsigned-edge-decomp} follows from the permutation
character on two-subsets.  For oriented cochains,
$E^-\cong\wedge^2(\mathbf1\oplus\mathbf3)
\cong\mathbf3\oplus\wedge^2\mathbf3$.  Since
$\det\mathbf3={\rm sgn}$ and $\mathbf3$ is real self-dual,
$\wedge^2\mathbf3\cong\mathbf3\otimes{\rm sgn}=\mathbf3'$.  The coboundary
of a constant is zero and its restriction to the reduced vertex module is
injective, so its image is the first $\mathbf3$ summand.  Quotienting leaves
$\mathbf3'$.
\end{proof}

\begin{firewall}[Internal generation fibre versus spatial displacement]
\label{fw:v40-internal-spatial}
The upstream renewal construction explicitly separates the cell cohomology
used as an internal generation carrier from the spatial displacement observable
used in the Perron metric.  Theorem~\ref{thm:v40-representation-fork} does not
identify those objects.  It only specifies which $S_4$ module is correct once
the operational writer---unsigned amplitude or oriented path increment---has
been fixed.  Equal rank is never an incidence map.
\end{firewall}

This resolves the apparent tension between source versions without forcing one
of the two edge modules to impersonate the other.  It also preserves V39's
oriented-cochain calculation as a mathematically valid optional finite-cell
translation-cover branch.

\subsection{The three-layer displacement decomposition}

The correct P4w acquisition target is therefore not simply ``six vectors on
$K_4$''.  One must first trace the operational reset record and determine which
of the following pieces enters the physical momentum tilt:
\begin{equation}
 \xi_{\rm phys}
 =\xi_{\rm cell}^{\rm fin}
  +\xi_{\rm skel}^{\rm th}
  +\xi_{\rm ext}^{\rm MA}.
 \label{eq:v40-three-layer-writer}
\end{equation}
Here:
\begin{enumerate}[leftmargin=2.2em]
\item $\xi_{\rm cell}^{\rm fin}$ is a bounded within-cell edge reward.  Its
      exact part is gauge-removable at fixed cell; only its oriented harmonic
      class can affect the finite-cell Perron pressure.
\item $\xi_{\rm skel}^{\rm th}=d\phi$ is the increment of a thermodynamic
      skeleton embedding.  It is graph-exact but need not have zero transport
      because $\phi$ is unbounded and the thermodynamic limit precedes the
      stationary long-time finite-box limit.
\item $\xi_{\rm ext}^{\rm MA}$ denotes any genuinely Markov-additive external
      translation mark not representable as a single-valued state function on
      the chosen quotient carrier.
\end{enumerate}

\begin{proposition}[Only the finite-cell piece is classified by K4 periods]
\label{prop:v40-period-scope}
The period matrix $\mathsf P_\xi$ of V39 is a complete gauge-invariant
classifier for the oriented harmonic part of
$\xi_{\rm cell}^{\rm fin}$ on one $K_4$.  It is not a classifier for the
thermodynamic skeleton gradient $d\phi$, whose transport is determined by the
large-volume growth and spectral content of the unbounded endpoint observable.
\end{proposition}

Thus a vanishing $K_4$ period matrix no longer forces the total physical metric
to vanish.  It forces only the finite-cell harmonic contribution to the
finite-cell pressure to vanish.

\subsection{Reinstating the V37 spectral packet in the correct order of limits}

The degree-shift identity used in V37 is independent of the Perron-similarity
argument.  On a finite or infinite graph on which the relevant operators are
defined,
\begin{equation}
 d\Delta_0=\Delta_1d,
 \qquad
 d\mu_{d\phi}(\lambda)=\lambda\,d\mu_\phi(\lambda).
 \label{eq:v40-degree-shift-restated}
\end{equation}
It is therefore still the correct raw spectral statement for the skeleton
gradient.  What V38 corrected was the attempt to read a fixed-finite Perron
pressure from that correlation measure.  What V40 corrects in turn is the
attempt to use the fixed-finite pressure to erase the thermodynamic transport.

\begin{corollary}[Corrected locality router for a skeleton-gradient carrier]
\label{cor:v40-restored-router}
Assume the actual metric-bearing displacement contains a thermodynamic skeleton
component $d\phi$ and that its ordered four-time cumulant admits the inherited
three-resolvent spectral reduction.  If
\begin{equation}
 d\mu_\phi(\lambda)\sim \lambda^{\eta_\phi-1}d\lambda,
 \label{eq:v40-phi-exponent}
\end{equation}
then its gradient current has exponent
\begin{equation}
 \eta_J=\eta_\phi+1.
 \label{eq:v40-current-exponent}
\end{equation}
Consequently the inherited Burnett criterion becomes
\begin{equation}
 \boxed{
 \eta_\phi>2\Rightarrow\text{local},\qquad
 \eta_\phi=2\Rightarrow\text{marginal},\qquad
 \eta_\phi<2\Rightarrow\text{nonlocal}.}
 \label{eq:v40-corrected-router}
\end{equation}
This statement concerns the thermodynamic/pre-boundary skeleton carrier and is
not contradicted by Theorem~\ref{thm:v40-fixed-finite-coboundary}.
\end{corollary}

The density-of-states saturation value $\eta_\phi=2/3$ remains conditional,
exactly as V37 already recorded.  The point of V40 is that this spectral
question is again meaningful if the physical writer is a genuinely unbounded
thermodynamic embedding coordinate.

\subsection{What P4w actually requires from the event record}

A single same-record audit can now decide which branch is physical.  For each
accepted reset event store:
\begin{equation}
 \mathcal R_j
 =(
   c_j^{\rm in},c_j^{\rm out},
   v_j^{\rm skel,in},v_j^{\rm skel,out},
   \delta x_j^{\rm ext},
   \text{clock/source labels}
  ).
 \label{eq:v40-reset-record}
\end{equation}
From the \emph{same} records form
\begin{align}
 \xi_j^{\rm cell}
   &=\xi_{c_j^{\rm in},c_j^{\rm out}},
   \label{eq:v40-record-cell}\\
 \xi_j^{\rm skel}
   &=\phi(v_j^{\rm skel,out})-\phi(v_j^{\rm skel,in}),
   \label{eq:v40-record-skel}\\
 \xi_j^{\rm ext}
   &=\delta x_j^{\rm ext}.
   \label{eq:v40-record-ext}
\end{align}
The physical momentum tilt itself then decides which combination is present:
\begin{equation}
 e^{ik\cdot\xi_j^{\rm phys}}
 =e^{ik\cdot(
      A_{\rm cell}\xi_j^{\rm cell}
     +A_{\rm skel}\xi_j^{\rm skel}
     +A_{\rm ext}\xi_j^{\rm ext})}.
 \label{eq:v40-physical-incidence}
\end{equation}
The incidence matrices $A_{\rm cell},A_{\rm skel},A_{\rm ext}$ are now the
load-bearing observables.  They cannot be inferred from the common dimension
three of any two downstream spaces.

\begin{definition}[Two-scale displacement certificate]
\label{def:v40-two-scale-certificate}
The P4w displacement certificate is the tuple
\begin{equation}
 \mathfrak C_{\rm disp}
 =(
 A_{\rm cell},A_{\rm skel},A_{\rm ext};
 \mathsf P_{\xi^{\rm cell}};
 M_\phi^{\rm th};
 C_{\rm GK}
 ).
 \label{eq:v40-displacement-certificate}
\end{equation}
Here $\mathsf P_{\xi^{\rm cell}}$ is the V39 $K_4$ period matrix,
$M_\phi^{\rm th}$ is the pre-boundary/thermodynamic spectral packet of the
skeleton coordinate, and $C_{\rm GK}$ is the covariance extracted from the same
physical tilt.
\end{definition}

The certificate passes only when the reconstructed covariance from its selected
writer equals the already calibrated Green--Kubo metric within the common
normalisation.  This prevents the internal cycle fibre or an arbitrary covering
writer from being substituted merely because it produces a convenient positive
matrix.

\subsection{Consequences for the V39 A4 cross-product candidate}

The writer
\begin{equation}
 h_{ij}=\ell\,r_i\times r_j
 \label{eq:v40-cross-recalled}
\end{equation}
remains an exact, full-period-rank, $A_4$-covariant finite-cell harmonic
cochain.  Its V39 metric and quartic Perron jet remain valid for the democratic
kernel.  V40 changes only its logical role:
\begin{center}
\fbox{\parbox{0.88\textwidth}{\centering
$h_{ij}$ is an optional finite-cell Markov-additive branch, not a necessary
rescue of skeleton transport.}}
\end{center}
It becomes physical only if $A_{\rm cell}\neq0$ and the same-record edge marks
project onto that harmonic class.  Conversely, $A_{\rm cell}=0$ is entirely
compatible with a nonzero spacetime metric if $A_{\rm skel}$ or
$A_{\rm ext}$ supplies the physical displacement.

\subsection{P4w status ledger after the upstream correction}

\subsection*{Proved in V40}
\begin{enumerate}[leftmargin=2.2em]
\item Exact displacement rewards have zero Perron jet on every fixed finite
      primitive carrier.
\item The finite coboundary no-go does not commute in general with the
      thermodynamic limit; the simple random walk on $\mathbb Z$ gives an exact
      counterexample with pressure $\log\cosh t$ and unit variance rate.
\item The correct finite-volume approximation to thermodynamic transport is a
      pre-boundary window, not the $n\to\infty$ equilibrium limit at fixed
      volume.
\item The unsigned six-edge permutation module and the natural oriented
      one-cochain module of $K_4$ are different $S_4$ representations.  The
      oriented graph-cohomology quotient is $\mathbf3'$, while the unsigned
      amplitude module contains the distinct $\mathbf1\oplus\mathbf2$ sector.
\item The V39 period matrix classifies only the finite-cell harmonic component;
      it does not classify an unbounded skeleton-gradient displacement.
\end{enumerate}

\subsection*{Inherited and retained}
\begin{enumerate}[leftmargin=2.2em]
\item The upstream renewal construction separates the critical skeleton, which
      carries long-distance return geometry, from the finite $K_4$ cell, whose
      cohomology is used as an internal matter fibre.
\item The common Perron tilt accepts a displacement edge observable and its
      Green--Kubo Hessian is the metric calibration whenever no nonzero
      physical projection is rate-coboundary in the relevant transport
      representation.
\item The spectral Hodge identity
      $d\mu_{d\phi}=\lambda d\mu_\phi$ remains exact and supplies the
      thermodynamic locality router once the correct writer and limit order are
      fixed.
\end{enumerate}

\subsection*{Conditional}
\begin{enumerate}[leftmargin=2.2em]
\item If the physical displacement entering the continuum momentum tilt is the
      unbounded skeleton embedding gradient, the V37 exponent test
      $\eta_\phi\gtrless2$ is reinstated in the thermodynamic/pre-boundary
      order of limits.
\item If the physical tilt contains a bounded $K_4$ edge component, only its
      harmonic period class contributes to the fixed-cell Perron jet; the V39
      cross-product writer remains one symmetry-allowed candidate for that
      component.
\item Any identification of a cell-cycle triplet with spatial displacement
      requires a measured incidence map; internal generation rank alone does
      not provide it.
\end{enumerate}

\subsection*{Open beyond V40}
\begin{enumerate}[leftmargin=2.2em]
\item Populate the three incidence blocks
      $A_{\rm cell},A_{\rm skel},A_{\rm ext}$ in
      \eqref{eq:v40-physical-incidence} from the actual same-record momentum
      tilt.
\item If $A_{\rm skel}\neq0$, populate the embedding spectral packet in the
      pre-boundary thermodynamic window and decide whether
      $\eta_\phi<2$, $=2$, or $>2$.
\item If $A_{\rm cell}\neq0$, acquire the six actual oriented edge vectors and
      compute the V39 period matrix without replacing them by the cross-product
      candidate.
\item If $A_{\rm ext}\neq0$, identify the covering/port-resolved external
      translation record and include it in the same Green--Kubo and Burnett
      cumulant calculation.
\end{enumerate}

\subsection*{Exact next load-bearing theorem}
\begin{center}
\fbox{\parbox{0.92\textwidth}{
\textbf{P4x: same-writer two-scale displacement landing and thermodynamic-order metric.}
Starting from one accepted Renewal--ESM event record, identify the exact
additive displacement used by the physical momentum tilt and decompose it into
finite-cell, skeleton-gradient, and external Markov-additive pieces as in
\eqref{eq:v40-three-layer-writer}.  Populate the incidence matrices in
\eqref{eq:v40-physical-incidence}.  For every nonzero finite-cell piece compute
its $K_4$ period matrix; for every nonzero skeleton-gradient piece compute its
pre-boundary thermodynamic characteristic/spectral packet before the fixed-box
long-time limit; and for any external piece record its covering-space periods.
Reconstruct the Green--Kubo covariance from the selected same writer and match
it to the already calibrated metric.  Only after this identity passes should
the Burnett locality router be applied.  In particular, do not require a
nonzero $K_4$ period merely to rescue an unbounded skeleton displacement, and
do not use the internal generation fibre as a spatial displacement space
without a measured incidence map.}}
\end{center}

\section*{Append-only continuation marker: V40}
\addcontentsline{toc}{section}{Append-only continuation marker: V40}
\textbf{End of V40.}  The complete V1--V39 body above is retained verbatim.
V40 resolves the finite-versus-thermodynamic coboundary ambiguity hidden in
P4w.  Exact rewards have zero Perron jet on a fixed finite primitive carrier,
but the theorem does not survive an unbounded thermodynamic embedding by
simple limit interchange; the infinite simple random walk gives an explicit
counterexample.  The correct finite-volume object is therefore the
pre-boundary transport window.  V40 also separates the unsigned and oriented
$K_4$ edge representations, showing why dimension-three edge modules appearing
in different source branches must not be identified without fixing the
operational writer.  The V39 cross-product cocycle remains a valid optional
finite-cell harmonic branch, but it is no longer needed as a rescue of the
critical-skeleton metric.  The next gate is the same-record incidence that
identifies which finite-cell, skeleton-gradient, or external Markov-additive
piece actually enters the physical momentum tilt and reproduces the calibrated
Green--Kubo metric.



\end{document}
